\documentclass[11pt,openany]{book}
\setcounter{secnumdepth}{3}
\usepackage[margin=1.1in]{geometry}
% Arabic edition: build with LuaLaTeX (see the Makefile); right-to-left text, left-to-right math
\usepackage[bidi=basic]{babel}
\babelprovide[import,main]{arabic}
\babelprovide[import]{english}
\babelfont{rm}[Renderer=HarfBuzz]{Amiri}
\babelfont{sf}[Renderer=HarfBuzz]{Amiri}
\babelfont[english]{rm}{Latin Modern Roman}
\babelfont[english]{sf}{Latin Modern Sans}
\usepackage{amsmath,amssymb,amsthm}
\usepackage{graphicx}
\usepackage{tikz}
\usetikzlibrary{arrows.meta,decorations.markings,decorations.pathmorphing,decorations.pathreplacing,calc}
\input{figures/icons}
\usepackage[unicode,colorlinks=true,linkcolor=blue!60!black,citecolor=blue!60!black]{hyperref}
\usepackage{microtype}
\setlength{\emergencystretch}{3em} % long names (licences, references) in narrow lines
\usepackage{empheq}
\usepackage{array}
\usepackage{booktabs}
\usepackage{multirow}
\usepackage{longtable}
\usepackage{circuitikz}
\definecolor{keyyellow}{rgb}{1.0,0.97,0.78}
% key equations: \begin{empheq}[box=\keybox]{equation} ... \end{empheq} -- light-yellow box, number stays outside
\newcommand{\keybox}[1]{{\setlength{\fboxsep}{7pt}\colorbox{keyyellow}{#1}}}
\usepackage[most]{tcolorbox}
% key statements (propositions etc.): the same light-yellow background, breakable across pages
\newtcolorbox{keyblock}{enhanced,breakable,colback=keyyellow,colframe=keyyellow,boxrule=0pt,arc=0pt,left=6pt,right=6pt,top=2pt,bottom=2pt,boxsep=0pt}
\renewcommand{\chaptermark}[1]{\markboth{\small\chaptername~\thechapter.\ #1}{}}
\renewcommand{\sectionmark}[1]{\markright{\small\thesection.\ #1}}
\renewcommand{\topfraction}{0.95}
\renewcommand{\textfraction}{0.05}
\renewcommand{\floatpagefraction}{0.75}

% part pages: title at the top third, and text may follow on the same page (used for the part introductions)
\makeatletter
\renewcommand\part{\if@openright\cleardoublepage\else\clearpage\fi\thispagestyle{plain}\if@twocolumn\onecolumn\@tempswatrue\else\@tempswafalse\fi\null\vspace{0.18\textheight}\secdef\@part\@spart}
\renewcommand\@endpart{\par\vspace{3em}\if@tempswa\twocolumn\fi}
\makeatother
\newtheorem{theorem}{مبرهنة}[chapter]
\newtheorem{proposition}{قضية}[chapter]
\newtheorem{lemma}{تمهيدية}[chapter]
\theoremstyle{remark}
\newtheorem{remark}{ملاحظة}[chapter]
% exercises: \begin{exercise}[\lvlB] ... \end{exercise}, numbered "Exercise 4.3"; difficulty marks
\theoremstyle{definition}
\newtheorem{exercise}{تمرين}[chapter]
\newcommand{\lvlA}{$\star$}
\newcommand{\lvlB}{$\star\star$}
\newcommand{\lvlC}{$\star\star\star$}
% answer to an exercise in the Answers appendix: \answer{ex:NN:k}{text}
\newcommand{\answer}[2]{\par\noindent\textbf{\ref{#1}.}~#2\par\smallskip}
\newcommand{\dd}{\mathrm{d}}
% neuron type code: first four letters of the primary mediator
\newcommand{\nt}[1]{\foreignlanguage{english}{\textsf{\textbf{#1}}}}
\newcommand{\mV}{\,\text{mV}}
\newcommand{\ms}{\,\text{ms}}
\newcommand{\mM}{\,\text{mM}}
\newcommand{\um}{\,\text{\textmu m}}
\newcommand{\ion}[2]{\mathrm{#1}^{#2}}
\newcommand{\Na}{\ion{Na}{+}}
\newcommand{\K}{\ion{K}{+}}
\newcommand{\Cl}{\ion{Cl}{-}}
\newcommand{\Ca}{\ion{Ca}{2+}}
\newcommand{\conc}[1]{[\mathrm{#1}]}

\title{\Huge حاسوب العشرين واط\\[16pt] \Large كيف يحسب الدماغ: مقاربة كمية للمهندسين}
\hypersetup{pdftitle={حاسوب العشرين واط. كيف يحسب الدماغ: مقاربة كمية للمهندسين},pdfauthor={Konstantin Kladko}}
\author{\Large كونستانتين كلادكو (Konstantin Kladko)}
\date{}

\begin{document}
\frontmatter
\maketitle
\thispagestyle{empty}
\vspace*{\fill}
\noindent{\small \copyright~2026 Konstantin Kladko.\\[4pt]
يُنشر النص والرسوم والتمارين والإجابات بموجب رخصة المشاع الإبداعي «نسب المصنف -- غير تجاري -- الترخيص بالمثل» 4.0 الدولية (CC BY-NC-SA 4.0): \url{https://creativecommons.org/licenses/by-nc-sa/4.0/deed.ar}. يجوز نسخ الكتاب واستخدامه في التدريس وترجمته وتعديله بحرية لأغراض غير تجارية، بشرط ذكر المؤلف ونشر الأعمال المشتقة بالرخصة نفسها.\\[4pt]
تُنشر برامج النماذج والبناء بموجب رخصة MIT.\\[4pt]
المصادر والنماذج والإصدارات الجديدة: \url{https://github.com/kladkogex/20-watt-computer}.}
\clearpage

\tableofcontents
\chapter*{كيف تقرأ هذا الكتاب}
\addcontentsline{toc}{chapter}{كيف تقرأ هذا الكتاب}

لا يلزم أن يُقرأ الكتاب من أوله إلى آخره. الفصول~\babelsublr{1--4} هي الأساس المشترك: أنواع العصبونات، وما يحسبه \nt{GLUT} و\nt{GABA}، وباللغة التي يتخاطبان بها. بعد ذلك يتفرع الكتاب. تبيّن الخريطة في الشكل~\ref{fig:roadmap} أيّ الفصول يبني على أيّها: السهم من الفصل $A$ إلى الفصل $B$ يعني أن $B$ يستعمل مفاهيم $A$ وصيغه. لا تُظهر الخريطة إلا التبعيات الرئيسية؛ أما الإحالات الصغيرة إلى الوراء فلا تعيق القراءة.

\begin{figure}[h]
\centering
\begin{tikzpicture}[>=Stealth,scale=0.95,
  ch/.style={draw,thick,rounded corners=3pt,align=center,font=\scriptsize,text width=1.85cm,minimum height=0.62cm,inner sep=2pt},
  base/.style={ch,fill=blue!8}, bus/.style={ch,fill=orange!14}, sum/.style={ch,fill=gray!18},
  io/.style={ch,fill=green!10}, sort/.style={ch,fill=cyan!8}, lab/.style={ch,fill=violet!10},
  dep/.style={->,thick,gray!60!black}]
% foundation
\node[base] (c1) at (0,0) {1 أنواع العصبونات};
\node[base] (c2) at (2.3,0) {2 \nt{GLUT}};
\node[base] (c3) at (4.6,0) {3 \nt{GLUT} و\nt{GABA}};
\node[base] (c4) at (6.9,0) {4 شفرة مورس};
% broadcast channels and the synthesis
\node[bus] (c5) at (4.6,-1.5) {5 \nt{SERO}};
\node[bus] (c6) at (7.3,-1.5) {6 \nt{DOPA}};
\node[bus] (c7) at (9.2,-0.9) {7 نظريات الدوبامين};
\node[bus] (c8) at (9.2,-2.1) {8 \nt{NORA}};
\node[bus] (c9) at (11.5,-2.1) {9 \nt{ACET}};
\node[sum] (c10) at (13.8,-0.9) {10 الآلة كاملة};
% inputs and outputs
\node[io] (c12) at (2.3,-4.1) {12 التعرّف};
\node[io] (c11) at (4.6,-4.1) {11 المداخل};
\node[io] (c13) at (6.9,-4.1) {13 العضلات};
\node[io] (c14) at (9.2,-3.05) {14 الألم};
\node[io] (c15) at (9.2,-4.1) {15 تعلّم\\الحركة};
\node[io] (c16) at (9.2,-5.15) {16 كيمياء\\الجسم};
% lab, sorting, sleep
\node[lab] (c17) at (11.5,-4.1) {17 تصحيح الأخطاء};
\node[lab] (c21) at (13.8,-4.1) {21 آلات حية};
\node[lab] (c22) at (13.8,-5.15) {22 \foreignlanguage{english}{DishBrain}};
\node[sort] (c18) at (11.5,-5.15) {18 عصبونات\\أجهزةَ قياس};
\node[sort] (c19) at (13.8,-6.2) {19 عدد\\التغصنات};
\node[bus] (c20) at (11.5,-6.2) {20 النوم};
% dependencies
\draw[dep] (c1) -- (c2); \draw[dep] (c2) -- (c3); \draw[dep] (c3) -- (c4);
\draw[dep] (c1) -- (c5); \draw[dep] (c4) -- (c6); \draw[dep] (c5) -- (c6);
\draw[dep] (c6) -- (c7); \draw[dep] (c6) -- (c8); \draw[dep] (c8) -- (c9);
\draw[dep] (c7) -- (c10); \draw[dep] (c9) -- (c10);
\draw[dep] (c4) -- (c11); \draw[dep] (c11) -- (c12); \draw[dep] (c11) -- (c13); \draw[dep] (c5) -- (c13);
\draw[dep] (c13) -- (c14); \draw[dep] (c13) -- (c15); \draw[dep] (c13) -- (c16);
\draw[dep] (c15) -- (c17); \draw[dep] (c9) -- (c17); \draw[dep] (c17) -- (c21); \draw[dep] (c21) -- (c22);
\draw[dep] (c16) -- (c18); \draw[dep] (c18) -- (c19); \draw[dep] (c16) -- (c20);
% legend
\begin{scope}[yshift=-7.25cm]
\foreach \s/\t [count=\i] in {base/الأساس, bus/القنوات والأنماط, sum/الخلاصة, io/المداخل والمخارج, sort/تصنيفات العصبونات, lab/المختبر}{
  \pgfmathtruncatemacro{\col}{mod(\i-1,3)}\pgfmathtruncatemacro{\row}{(\i-1)/3}
  \node[\s,text width=0.3cm,minimum height=0.32cm] at ({\col*4.8+0.2},{-\row*0.55}) {};
  \node[anchor=west,font=\scriptsize] at ({\col*4.8+0.55},{-\row*0.55}) {\t};}
\end{scope}
\end{tikzpicture}
\caption{خريطة الفصول. يتجه السهم من الفصل إلى الفصل الذي يبني عليه. أما الإحالات الأخرى بين الفصول فهي إشارات لا تبعيات.}
\label{fig:roadmap}
\end{figure}

\section*{مسارات عبر الكتاب}

\begin{center}
\footnotesize
\setlength{\tabcolsep}{4pt}
\renewcommand{\arraystretch}{1.2}
\begin{tabular}{>{\raggedright}p{3.0cm}>{\raggedright}p{3.9cm}>{\raggedright\arraybackslash}p{6.4cm}}
\toprule
المسار & لمن & الفصول بالترتيب \\
\midrule
مقرر ربع دراسي & المدرّس، 10 أسابيع & الأسبوع 1: الفصلان~\babelsublr{1--2}؛ 2: الفصل~3؛ 3: الفصل~4؛ 4: الفصلان~\babelsublr{5--6}؛ 5: الفصلان~\babelsublr{7--8}؛ 6: الفصل~9؛ 7: الفصل~10؛ 8: الفصلان~\babelsublr{11--12}؛ 9: الفصلان~13 و15؛ 10: الفصل~17 \\
الذكاء الاصطناعي وتعلّم الآلة & من يعرف الشبكات العصبية ويريد أن يفهم كيف يتعلم الدماغ & \babelsublr{1--4}، 5 (قراءة سريعة)، 6، 7، 8، 9، 11 (قراءة سريعة)، 12، 13 (قراءة سريعة)، 15 \\
الإلكترونيات والعتاد & مصمم الدوائر، ومطوّر الرقاقات العصبية الشكل & \babelsublr{1--4}، 5، 11، 13، 16، 18، 19، 17 (والفصلان 9 و15 قراءة سريعة) \\
الواجهات العصبية والحواسيب الحية & من يبني واجهات الدماغ والحاسوب ورقاقات بعصبونات حية & 1، 2، 4، 11، 13، 17 (والفصلان 9 و15 قراءة سريعة)، 21، 22 (والفصل 12 قراءة سريعة) \\
نظرة سريعة & مهندس لديه عطلة نهاية أسبوع & 1، 2، 3، 6، 10؛ وإحالات الفصل~6 إلى الفصلين~\babelsublr{4--5} مفهومة من السياق \\
\bottomrule
\end{tabular}
\end{center}

”القراءة السريعة“ تعني أنه يكفي أن تقرأ مطلع الفصل، وقضاياه في الأطر الصفراء، وجدوله الختامي.

في نهاية كل فصل تمارين: تقديرات، واشتقاقات، ونمذجة، وتصميم؛ والإجابات المختصرة مجموعة في الملحق~\ref{sec:answers}. وكل الرموز في الملحق~\ref{sec:notation}، أما التجارب التي تستند إليها القضايا الرئيسية في كل فصل ففي الملحق~\ref{sec:supplement}. الأرقام في الكتاب ناتجة عن برامج صغيرة في المجلد \texttt{models/}؛ ويمكن تشغيلها وتعديلها.

\mainmatter

\chapter{أنواع العصبونات}
\label{sec:neurontypes}

\section{العصبون صندوقًا أسود}

لنفترض أنك أُعطيت كيسًا فيه ستة وثمانون مليار عصبون~\cite{Azevedo2009}، وطُلب منك أن توزعها أكوامًا. كيف ستفعل ذلك؟

المهندس الذي يُعطى كيسًا من دوائر متكاملة لا يعرفها لن يفرزها بحسب شكل الغلاف. بل سيسأل: كم مدخلًا للقطعة، وكم مخرجًا، وماذا تُخرج. فلنرسم العصبون بالطريقة نفسها: كتلةً في مخطط (الشكل~\ref{fig:blackbox}).

\begin{figure}[t]
\centering
\begin{tikzpicture}[>=Stealth,x=1cm,y=1cm]
% the box
\node[draw,very thick,minimum width=2.6cm,minimum height=2.6cm,fill=blue!6] (N) at (0,0) {};
\node at (0,0.55) {\small عصبون};
\node at (0,-0.15) {$\displaystyle u=\sum_i w_i x_i$};
\node at (0,-0.8) {\small $y=\Theta(u-\theta)$};
% inputs
\foreach \k/\y/\lab in {1/1.05/{x_1},2/0.55/{x_2},3/0.05/{x_3}}{
  \draw[->,thick,blue!60!black] (-4.2,\y) -- (-1.3,\y);
  \node[left] at (-4.2,\y) {$\lab$};
  \node[above,blue!60!black] at (-2.1,\y-0.06) {\scriptsize $w_{\k}$};}
\node at (-4.45,-0.4) {$\vdots$};
\node at (-2.1,-0.4) {$\vdots$};
\draw[->,thick,blue!60!black] (-4.2,-1.05) -- (-1.3,-1.05);
\node[left] at (-4.2,-1.05) {$x_n$};
\node[above,blue!60!black] at (-2.1,-1.11) {\scriptsize $w_{n}$};
\draw[decorate,decoration={brace,amplitude=5pt},gray!60!black] (-4.9,-1.2) -- (-4.9,1.2);
\node[left,align=right,gray!40!black] at (-5.1,0) {\small $n\sim10^3\text{--}10^5$\\[-2pt]\small مدخل};
% single output wire, then fan-out
\draw[very thick,red!70!black] (1.3,0) -- (3.2,0);
\foreach \x in {1.6,2.2,2.34,2.9}{\draw[thick,red!70!black] (\x,0) -- (\x,0.4);}
\node[above right,font=\tiny\itshape,gray!30!black,inner sep=1pt] at (1.42,0.45) {نقرة\dots\ نقرة-نقرة\dots};
\node[below,red!70!black,align=center] at (2.25,-0.05) {\scriptsize مخرج واحد $y$};
\fill[red!70!black] (3.2,0) circle (0.06);
\foreach \y in {1.3,0.65,0,-0.65,-1.3}{
  \draw[->,thick,red!70!black] (3.2,0) -- (5.0,\y);}
\draw[decorate,decoration={brace,amplitude=5pt},gray!60!black] (5.3,1.45) -- (5.3,-1.45);
\node[right,align=left,gray!40!black] at (5.5,0) {\small نسخ من $y$\\[-2pt]\small إلى $10^3\text{--}10^4$\\[-2pt]\small عصبون آخر};
\end{tikzpicture}
\caption{العصبون بعين المهندس. مداخل كثيرة $x_i$، لكل منها وزنه $w_i$؛ وفي الداخل مجموع موزون $u$ ومقارنة بالعتبة $\theta$ ($\Theta$ دالة الدرجة: $1$ إذا كان المتغير موجبًا، و$0$ في غير ذلك)؛ ومخرج واحد $y$ هو سلسلة من النبضات تتكاثر بالتفرع لتصل إلى آلاف المستقبِلين.}
\label{fig:blackbox}
\end{figure}

ليس في مخرج الكتلة عدد، بل سلسلة من نبضات جهد قصيرة: ”نقرة“، ثم توقف، ثم ”نقرة-نقرة“، ثم توقف طويل. كل النبضات بالارتفاع نفسه، فالمعلومة كلها في \emph{توقيت} حدوثها. وفي داخل الكتلة، في تقريب أول خشن، جامع ومقارن: تُجمع المداخل بأوزانها، فإذا تجاوز المجموع العتبة أطلقت الكتلة نبضة. في الفصل التالي نستبدل بهذا النموذج نموذجًا يعمل بأمانة في الزمن المتصل.

المخرج واحد، لكن السلك يتفرع، وكل فرع يحمل \emph{النسخة نفسها} من الإشارة. يسمى هذا في الإلكترونيات التفرع عند المخرج (\foreignlanguage{english}{fan-out}). وهو عند عصبون القشرة من رتبة بضعة آلاف، بينما تحمّل البوابة المنطقية في الدائرة المتكاملة عادة ببضع بوابات أخرى فقط.

\emph{فما الذي} ينتقل عبر سلك المخرج إلى المستقبِل؟ تبلغ النبضة نهاية كل فرع، وهناك تتحول إلى إشارة كيميائية: يقذف الفرع في الشق الضيق بين الخليتين قليلًا من مادة معينة، هي \emph{الناقل العصبي}، فتغيّر تيار الدخل عند المستقبِل. وهذا هو معيار الفرز: \emph{أي ناقل عصبي يقذفه مخرج العصبون}.

\section{مخرج واحد، نوع واحد من الإشارة}

لا معنى للفرز إلا إذا كان لكل عصبون نوع واحد من المخرج. فهل الأمر كذلك؟ تقول التجربة إنه كذلك تقريبًا.

\begin{keyblock}
\begin{proposition}[عصبون واحد، نوع مخرج واحد]
\label{prop:dale}
لكل عصبون تقريبًا ناقل عصبي رئيسي واحد، ويقذفه في جميع فروع مخرجه. لذا يُحدَّد نوع العصبون بناقله الرئيسي~\cite{Strata1999}.
\end{proposition}
\end{keyblock}

لنتفق على أن نرمز إلى نوع العصبون \emph{بالأحرف الأربعة الأولى من الاسم الإنجليزي لناقله الرئيسي}. فالعصبون الذي ناقله الرئيسي الغلوتامات (\foreignlanguage{english}{glutamate}) هو عصبون \nt{GLUT}. الأنواع كلها مجموعة في الجدول~\ref{tab:types}.

\begin{table}[t]
\centering
\footnotesize
\setlength{\tabcolsep}{3pt}
\renewcommand{\arraystretch}{1.15}
\begin{tabular}{>{\raggedright}p{1.0cm}>{\raggedright}p{2.05cm}>{\raggedright}p{1.75cm}>{\raggedright}p{1.6cm}>{\centering}p{0.7cm}>{\raggedright}p{1.55cm}>{\raggedright}p{1.8cm}>{\raggedright\arraybackslash}p{3.2cm}}
\toprule
النوع & الناقل العصبي الرئيسي & الناقل (\foreignlanguage{english}{bus}) & السرعة & العلامة & العصبونات في الدماغ & مخارج العصبون & الدور في الحساب \\
\midrule
\nt{GLUT} & غلوتامات & عناوين & سريع وبطيء & $+$ & $\sim8\cdot10^{10}$ & $\sim10^4$ & الوزن الموجب الرئيسي \\
\nt{GABA} & غابا & عناوين & سريع وبطيء & $-$ & $\sim3\cdot10^{9}$ & $10^3\text{--}10^4$ & الوزن السالب الرئيسي \\
\nt{GLYC} & غليسين & عناوين & سريع & $-$ & $10^6\text{--}10^7$ & $10^2\text{--}10^3$ & وزن سالب في الحبل الشوكي وجذع الدماغ؛ ومفتاح ثانٍ لأحد مستقبلات الغلوتامات \\
\nt{ACET} & أسيتيل كولين & كلاهما & سريع وبطيء & $\pm$ & $\sim10^6$ & العضلة: $10\text{--}10^3$؛ الدماغ: $\sim10^5$ & مخرج إلى العضلة؛ وفي الدماغ سرعة التعلم (فرضية) \\
\nt{DOPA} & دوبامين & بثّ & بطيء & $\pm$ & $\sim5\cdot10^{5}$ & $10^5\text{--}10^6$ & خطأ التنبؤ بالمكافأة $\delta$ \\
\nt{SERO} & سيروتونين & بثّ & بطيء (نوع واحد من المستقبلات سريع) & $\pm$ & $\sim3\cdot10^{5}$ & $\sim10^5$ & أفق التخطيط (فرضية) \\
\nt{HIST} & هستامين & بثّ & بطيء & $+$ & $\sim6\cdot10^{4}$ & لا تقدير & إشارة عامة: ”ابقَ مستيقظًا“ \\
\nt{NORA} & نورأدرينالين & بثّ & بطيء & $\pm$ & $\sim5\cdot10^{4}$ & $\sim10^5$ & الكسب (فرضية) \\
\nt{ADRE} & أدرينالين & بثّ & بطيء & $\pm$ & $\sim10^4$ & لا تقدير & عصبونات قليلة؛ ودورها في الدماغ غير واضح \\
\nt{ASPA} & أسبارتات & عناوين & سريع & $+$ & لا تقدير & لا تقدير & وزن موجب ضعيف؛ ودوره موضع خلاف \\
\bottomrule
\end{tabular}
\caption{أنواع العصبونات، مرتبة تنازليًا بحسب عددها في الدماغ. \emph{الناقل}: ناقل العناوين، أي يُقذف الناقل العصبي في شق ضيق نحو مستقبِل محدد؛ وناقل البثّ، أي من مجموعة صغيرة من العصبونات المصدر إلى مناطق واسعة (القسم~\ref{sec:buses}). \emph{السرعة}: سريع، أي عبر المستقبلات مؤينة التوجه، في أجزاء من الألف من الثانية؛ وبطيء، أي عبر المستقبلات أيضية التوجه، في مئات من أجزاء الألف من الثانية وأكثر. \emph{العلامة} هي العلامة المعتادة؛ أما العلامة النهائية فيحددها المستقبِل (القسم~\ref{sec:receiver})، و$\pm$ تعني أن العلامتين كلتيهما موجودتان. \emph{العصبونات في الدماغ}: عدد عصبونات هذا النوع في دماغ الإنسان كله، من حيث رتبة المقدار؛ ومجموع العصبونات نحو $8.6\cdot10^{10}$~\cite{Azevedo2009}. معظم عصبونات \nt{GLUT}، أي نحو $7\cdot10^{10}$، هي عصبونات الدخل الصغيرة في المخيخ؛ وعددا \nt{GLUT} و\nt{GABA} مستخرجان من هذا العدّ ومن نسبة عصبونات \nt{GABA} في القشرة~\cite{Hendry1987}. ومن الأنواع القليلة العدد لم يُعدّ مباشرة إلا \nt{HIST}~\cite{Airaksinen1991} والمجموعة الرئيسية من مجموعتي عصبونات \nt{DOPA}~\cite{Pakkenberg1991}، لذا فالعدد في حالة \nt{DOPA} حدّ أدنى؛ وفي حالة \nt{SERO} هو مجموع عدّين جزئيين~\cite{Baker1990,Baker1991}. أما أعداد \nt{NORA} و\nt{GLYC} و\nt{ACET} و\nt{ADRE} فتقديرات خشنة لرتبة المقدار بلا عدّ مباشر. \emph{مخارج العصبون}: عدد النهايات المشبكية للعصبون الواحد، أي مواضع قذف الناقل العصبي على فروع سلك المخرج (الشكل~\ref{fig:blackbox})، من حيث رتبة المقدار. لم تُعدّ النهايات مباشرة في أنواع البثّ: فعددها عند \nt{DOPA} مستنتج من نسبة الهدف التي يغطيها تفرع سلك مخرج واحد~\cite{Matsuda2009}، وتقديرات البقية أخشن. ولـ\nt{ACET} حالتان: العصبون الذي يتحكم في عضلة، وعصبون البثّ في الدماغ.}
\label{tab:types}
\end{table}

% the pie is a table-type float with a figure caption, so it can never float ahead of Table~\ref{tab:types}
\begin{table}[tb]
\makeatletter\def\@captype{figure}\makeatother
\centering
\definecolor{vizglut}{HTML}{2A78D6}
\definecolor{vizgaba}{HTML}{E34948}
\definecolor{vizrest}{HTML}{8A8986}
\begin{tikzpicture}[>=Stealth]
% pie: GLUT 96.4 %, GABA 3.6 % (13.0 degrees), the rest 0.006 % (a hairline at 0 degrees)
\fill[vizglut] (0,0) -- (13:1.9) arc (13:360:1.9) -- cycle;
\fill[vizgaba] (0,0) -- (0:1.9) arc (0:13:1.9) -- cycle;
\draw[white,line width=1pt] (0,0) -- (13:1.9);
\draw[vizrest,line width=1.2pt] (0,0) -- (0:1.9);
\node[white,align=center,font=\small] at (190:0.95) {\nt{GLUT}\\$96.4\,\%$};
\draw[gray!60!black] (6.5:1.75) -- (2.05,2.1);
\node[anchor=west,font=\small] at (2.05,2.1) {\nt{GABA} $3.6\,\%$};
% zoom box for the remaining types
\draw[densely dashed,vizrest] (1.9,0) -- (3.1,1.65);
\draw[densely dashed,vizrest] (1.9,0) -- (3.1,-2.2);
\draw[vizrest,rounded corners=3pt] (3.1,-2.2) rectangle (10.1,1.65);
\node[anchor=west,font=\scriptsize] at (3.2,1.4) {كل الأنواع الأخرى معًا: $\approx0.006\,\%$؛ مقياس لوغاريتمي};
\foreach \code/\lg/\y/\val/\lx in {GLYC/6.477/1.0/{$10^6\text{--}10^7$}/4.05,ACET/6.0/0.6/{$\sim10^6$}/3.0,DOPA/5.699/0.2/{$\sim5\cdot10^5$}/2.699,SERO/5.477/-0.2/{$\sim3\cdot10^5$}/2.477,HIST/4.778/-0.6/{$\sim6\cdot10^4$}/1.778,NORA/4.699/-1.0/{$\sim5\cdot10^4$}/1.699,ADRE/4.0/-1.4/{$\sim10^4$}/1.0}{
  \node[anchor=east,font=\scriptsize] at (4.35,\y) {\nt{\code}};
  \fill[vizrest] (4.45,\y-0.13) rectangle ({4.45+\lg-3},\y+0.13);
  \node[anchor=west,font=\scriptsize] at ({4.5+\lx},\y) {\val};}
\draw[vizrest,thick] (7.45,1.0) -- (8.45,1.0);
\draw[vizrest,thick] (7.45,0.92) -- (7.45,1.08);
\draw[vizrest,thick] (8.45,0.92) -- (8.45,1.08);
\draw[gray!60!black] (4.45,-1.72) -- (8.45,-1.72);
\foreach \e in {3,4,5,6,7}{
  \draw[gray!60!black] ({4.45+\e-3},-1.72) -- ({4.45+\e-3},-1.66);
  \node[below,font=\scriptsize] at ({4.45+\e-3},-1.72) {$10^{\e}$};}
\end{tikzpicture}
\caption{حصص أنواع العصبونات في الدماغ بحسب تقديرات الجدول~\ref{tab:types}. تكاد عصبونات \nt{GLUT} تملأ الدائرة كلها؛ وعصبونات \nt{GABA} قطاع ضيق، أما كل الأنواع الأخرى معًا فشعرة على الحدّ. وعلى اليمين هذه الشعرة مكبّرة، بمقياس لوغاريتمي لعدد العصبونات؛ عمود \nt{GLYC} يبلغ منتصف المجال $10^6\text{--}10^7$، والقطعة تبيّن المجال كله. لا تقدير لـ\nt{ASPA}، فهو غائب عن الشكل.}
\label{fig:typepie}
\end{table}

يُظهر الشكل~\ref{fig:typepie} مدى التفاوت بين هذه الأكوام: من كل ألف عصبون $964$ من نوع \nt{GLUT}، و$36$ من نوع \nt{GABA}، أما الأنواع الأخرى كلها معًا فنصيبها نحو ستة عصبونات من كل مئة ألف.

اسم \foreignlanguage{english}{GABA} مؤلف أصلًا من أربعة أحرف. أما المواد التي لا تكاد تكون ناقلًا رئيسيًا، كالببتيدات العصبية والغازات والدهون والبيورينات، فلا تحصل على رمز خاص بها؛ والحديث عنها في الملحق~\ref{sec:othermediators}. كلمة ”تقريبًا“ في القضية~\ref{prop:dale} مهمة. فكثير من العصبونات يقذف مع الناقل الرئيسي موادَّ مساعدة~\cite{Yao2023}. وبعضها يقذف أيضًا ناقلًا سريعًا ثانيًا، قد يكون معاكسًا في العلامة: فجزء من عصبونات \nt{DOPA} يقذف الغابا أيضًا~\cite{Tritsch2012}. وفي بعض العصبونات تخرج نواقل مختلفة من فروع مختلفة للمخرج نفسه~\cite{Zhang2015}. كل هذا مُقاس، فقاعدة ”عصبون واحد، ناقل واحد“ قاعدة لها استثناءات وليست قانونًا. لكنها تصمد تقسيمًا أول: ففي الأطلس الخلوي لدماغ الفأر، حيث صُنّفت العصبونات بحسب الجينات العاملة فيها إلى أكثر من خمسة آلاف مجموعة، ما زالت العصبونات تنقسم إلى أصناف كبرى بحسب ناقلها الرئيسي أولًا~\cite{Yao2023}.

لهذه القاعدة نتيجة تميّز الدماغ فورًا عن الشبكة العصبية الاصطناعية. في الشبكة الاصطناعية يمكن للعصبون نفسه أن يكون له وزن موجب نحو مستقبِل ووزن سالب نحو مستقبِل آخر: فالأوزان $w_{ij}$ مجرد أعداد في مصفوفة. أما في الدماغ فعلامة التأثير يحددها في الغالب الناقل العصبي، وللعصبون ناقل واحد. إذن إذا أثار العصبون $j$ مستقبِلًا واحدًا فإنه يثير الباقين أيضًا. \emph{العمود} الخارج من العصبون $j$ في مصفوفة الأوزان كله بعلامة واحدة:
\begin{empheq}[box=\keybox]{equation}
\operatorname{sign} w_{ij} \;=\; s_j \quad\text{لكل المستقبِلين } i, \qquad s_j\in\{+1,-1\}.
\label{eq:dalesign}
\end{empheq}
تنقسم العصبونات إلى مثيرة ($s_j=+1$) ومثبطة ($s_j=-1$)، وهذه خاصية العصبون نفسه لا خاصية الوصلة. فإذا احتاجت الشبكة إلى وصلة سالبة من عصبون مثير، وجب أن تمر عبر وسيط: العصبون المثير يثير عصبونًا مثبطًا، وهذا يثبط الهدف، فينتج وزن سالب متأخر بحلقة واحدة.

النواقل العصبية المعروفة أكثر من مئة، لكن عمل الدماغ كله تقريبًا يقوم على نصف دزينة منها، وهي تنقسم إلى صنفين مختلفين تمامًا.

\section{ناقلان: ناقل العناوين وناقل البثّ}
\label{sec:buses}

في شبكات الحاسوب طريقتان لإيصال رسالة. الأولى \emph{عنوانية} (\foreignlanguage{english}{unicast}): تذهب الرزمة من مرسل واحد إلى مستقبِل محدد، بسرعة، ولا يراها أحد غيره. والثانية \emph{بثّ} (\foreignlanguage{english}{broadcast}): يرسل المرسل رسالة واحدة إلى كل عقد الشبكة دفعة واحدة. تستعمل العصبونات الطريقتين كلتيهما، وتنقسم النواقل العصبية بحسب هذا المعيار بالضبط (الشكل~\ref{fig:phone_radio}).

\begin{figure}[t]
\centering
\begin{tikzpicture}[>=Stealth,scale=0.95]
% ---- unicast: point-to-point wiring
\node[above] at (2.0,3.3) {\small \textbf{ناقل العناوين}: الغلوتامات والغابا};
\foreach \i in {0,...,4}{
  \fill[blue!60!black] (0.2+0.9*\i,2.5) circle (0.1);
  \fill[blue!60!black] (0.2+0.9*\i,0.6) circle (0.1);}
\foreach \a/\b in {0/0,0/1,1/1,1/3,2/2,3/2,3/4,4/4,2/0}{
  \draw[->,blue!60!black] (0.2+0.9*\a,2.38) -- (0.2+0.9*\b,0.72);}
\fill[red!70!black] (1.1,1.55) circle (0.09);
\pic[scale=1.2] at (1.75,1.9) {letter};
\draw[->,red!70!black,thick] (1.1,1.55) -- (0.28,0.7);
\draw[->,red!70!black,thick] (1.1,1.55) -- (1.95,0.72);
\node[align=center,below] at (2.0,0.2) {\small مليارات العصبونات،\\[-2pt]\small لكلٍّ منها مستقبِلوه،\\[-2pt]\small الزمن: أجزاء من الألف من الثانية};
% ---- broadcast: one small source, global targets
\begin{scope}[xshift=7.2cm]
\node[above] at (1.8,3.3) {\small \textbf{ناقل البثّ}: الدوبامين، \dots};
\foreach \i in {0,...,8}{\fill[blue!60!black] (-0.2+0.5*\i,2.75) circle (0.08);}
\fill[orange!85!black] (1.8,0.35) ellipse (0.35 and 0.18);
\pic[scale=1.5] at (2.7,0.75) {mast};
\foreach \x in {-0.2,0.3,0.8,1.3,1.8,2.3,2.8,3.3,3.8}{
  \draw[->,orange!85!black] (1.8,0.5) .. controls (1.8,1.3) and (\x,1.6) .. (\x,2.62);}
\node[align=center,below] at (1.8,0.1) {\small من آلاف إلى مئات الآلاف من العصبونات،\\[-2pt]\small رسالة واحدة للجميع،\\[-2pt]\small الزمن: مئات من أجزاء الألف من الثانية وأكثر};
\end{scope}
\end{tikzpicture}
\caption{طريقتان للإرسال. على اليسار ناقل العناوين، وهو يشبه البريد والعنوان مكتوب على الظرف؛ وبالأحمر عصبون واحد ومستقبِلاه. وعلى اليمين ناقل البثّ، كمحطة إذاعة: مجموعة صغيرة من العصبونات المصدر، والرسالة نفسها يتلقاها الجميع.}
\label{fig:phone_radio}
\end{figure}

\emph{ناقل العناوين} هو الغلوتامات والغابا. تقذفهما الغالبية الساحقة من العصبونات. كل مخرج يقود إلى خلايا محددة، ويعمل الناقل العصبي في الشق الضيق للتماس فقط، ويعمل بسرعة: في جزء من الألف من الثانية تنفتح القنوات الأيونية لدى المستقبِل، وفي بضعة أجزاء من الألف ينتهي كل شيء. هذا ناقل \emph{البيانات}: عبره ينتقل ما يحسبه الدماغ.

\emph{ناقل البثّ} هو الدوبامين والسيروتونين والنورأدرينالين، والأسيتيل كولين جزئيًا. تقذفها مجموعات ضئيلة من العصبونات، لكن مخرج كل منها يتفرع اتساعًا مذهلًا. وكثيرًا ما يُقذف الناقل العصبي لا في شق ضيق بل في الحيز بين الخلايا، فينتشر فيه ويؤثر في كل من حوله. ويعمل ببطء: لا يفتح القنوات مباشرة، بل يطلق داخل المستقبِل سلسلة من التفاعلات الكيميائية. لا تجري عبره البيانات، بل \emph{وسائط تحكم} مشتركة بين أجزاء كبيرة من الشبكة، تغيّر نمط عملها: الكسب، وسرعة التعلم، وإشارة ”كان هذا جيدًا“. لذلك تسمى هذه النواقل \emph{المعدِّلات}. ساد طويلًا الاعتقاد أن كل قناة من هذه القنوات عدد واحد للدماغ كله. وقد دققت القياسات الحديثة الصورة: القناة ليست سلكًا واحدًا، بل حزمة صغيرة من الخطوط المتوازية. فالعصبونات المصدر لناقل واحد تنقسم إلى أنظمة فرعية، لكل منها مستقبِلوه ورسالته، كما تضبط الإشارات المحلية كمية الناقل المقذوف في موضعه~\cite{Ren2018,Poe2020}. وعدد هذه الخطوط آحاد أو عشرات لا مليارات، فتبقى القناة قناة بثّ.

إن احتجت إلى صورة واحدة من هذا الفصل فهذه هي. الدماغ شبكة ضخمة تنقل البيانات بالعناوين، وفوقها بضع قنوات بثّ تحمل إشارات التحكم. سيعرف المبرمج هذا التصميم المألوف: حساب، ومعه مجموعة صغيرة من متغيرات التحكم، كل منها مشترك بين جزء كبير من الشبكة.

\section{\nt{GLUT}: الوزن الموجب}

الغلوتامات هو الناقل المثير الرئيسي في الدماغ. عندما يقذف مخرج العصبون الغلوتامات، تنفتح لدى المستقبِل قنوات يدخل عبرها الصوديوم، فيرتفع الجهد على غشاء المستقبِل، وتزداد فرصته لإطلاق نبضة خاصة به. وبلغة الشكل~\ref{fig:blackbox} هذا مدخل بوزن \emph{موجب}، $w_i>0$.

من يُخرج الغلوتامات؟ كل من ينقل البيانات إلى مسافات بعيدة تقريبًا: من ثلاثة أرباع إلى أربعة أخماس عصبونات القشرة، ومعظم العصبونات التي تربط مناطق الدماغ المختلفة. شبكة الدماغ هي أولًا شبكة من الوصلات الغلوتاماتية.

تفصيل هندسي واحد. بعد القذف يقفز تركيز الغلوتامات في الشق آلاف المرات، إلى نحو ميلي مول، ثم يهبط بثابت زمني نحو جزء من الألف من الثانية~\cite{Clements1992}: ينتشر الغلوتامات خارج الشق، وتضخه بروتينات خاصة عائدة به إلى الخلايا. لماذا؟ للسبب نفسه الذي يُعاد لأجله الإشارة إلى الصفر بعد كل رمز في خط الاتصال. إن لم يُزَل الناقل العصبي، فستصل النبضة التالية إلى خط ”مشغول“، وستندمج النبضات التي يفصل بينها بضعة أجزاء من الألف من الثانية.

\section{\nt{GABA}: الوزن السالب}

تخيّل شبكة كل أوزانها موجبة. إذا ولّدت كل نبضة في المتوسط أكثر من نبضة لاحقة واحدة، نمت الفعالية نموًا أسيًا واجتاحت الشبكة كلها في بضع خطوات. سيعرف مهندس الإلكترونيات هنا حلقة تغذية راجعة موجبة بكسب أكبر من الواحد: تذهب الدارة إلى الإشباع. ولتجنب ذلك نحتاج إلى أوزان سالبة. ومصدرها الرئيسي حمض غاما-أمينوبيوتيريك، واختصاره \foreignlanguage{english}{GABA} (غابا).

تفتح الغابا لدى المستقبِل قنوات تمرر أيونات الكلور. جهد اتزان الكلور قريب من جهد الراحة أو أدنى منه قليلًا، لذا تبقي قنوات الكلور المفتوحة الغشاءَ قرب الراحة ولا تدع الجهد يرتفع إلى العتبة. هذا مدخل بوزن \emph{سالب}، $w_i<0$. عصبونات \nt{GABA} تشكّل من خُمس إلى ربع عصبونات القشرة~\cite{Hendry1987}. مخارجها قصيرة، وهي تثبط جيرانها الأقربين.

التثبيط في الدماغ يفعل أكثر بكثير من مجرد الطرح. إنه يعمل تغذيةً راجعة سالبة تبقي الشبكة كلها عند نقطة التشغيل. ويُعتقد أن التيارين المثير والمثبط الواردين إلى العصبون في القشرة السليمة يكادان يتوازنان بدقة: هذا النمط تنبأت به النظرية~\cite{vanVreeswijk1996} وتُظهره قياسات التيارات في القشرة الحية~\cite{Haider2006}، فيبقى العصبون طوال الوقت قرب العتبة. والدارة المستقرة بتغذية راجعة سالبة قوية حول نقطة غير مستقرة تستجيب بسرعة وحساسية للإشارات الصغيرة؛ وهذه حيلة هندسية قديمة.

\section{\nt{DOPA}: إشارة الخطأ}

أول قناة بثّ هي الدوبامين. عدد العصبونات الدوبامينية عند الإنسان من رتبة نصف مليون، أي نحو عصبون واحد من كل مئة وسبعين ألفًا؛ هذا ما عُدّ في المجموعة الرئيسية من مجموعتيها~\cite{Pakkenberg1991}، فهو حدّ أدنى. تمتد مخارجها إلى المناطق التي تختار الأفعال.

ماذا تنقل هذه القناة؟ الجواب تعطيه تجارب تُسجَّل فيها نبضات العصبونات الدوبامينية لدى حيوان ينال مكافأة~\cite{Schultz1997}. يُعطى الحيوان من حين لآخر قطرة عصير على غير توقع، فتجيب العصبونات الدوبامينية برشقة قصيرة من النبضات. ثم يُضاء مصباح قبل العصير، دائمًا قبله بثانية. وعندما يتعلم الحيوان هذا الارتباط، تبدأ العصبونات بالاستجابة برشقة \emph{للمصباح}، ولا تستجيب إطلاقًا للعصير نفسه إذا جاء في موعده تمامًا. وإذا لم يُعطَ العصير بعد المصباح، \emph{تصمت} العصبونات في اللحظة التي كان ينبغي أن يأتي فيها، ويهبط ترددها دون المعتاد.

القاعدة التي تفسر الحالات الثلاث (الشكل~\ref{fig:rpe}): لا تخبر العصبونات بمدى جودة الأمر، بل بمدى كونه \emph{أفضل أو أسوأ مما كان متوقعًا}.

\begin{figure}[t]
\centering
\begin{tikzpicture}[>=Stealth,x=3.4cm,y=1cm]
\foreach \y/\lab in {2.8/{عصير غير متوقع},1.4/{عصير بعد المصباح},0/{مصباح بلا عصير}}{
  \draw[gray!70!black] (0,\y) -- (3,\y);
  \draw[densely dotted,gray!60!black] (0,\y+0.3) -- (3,\y+0.3);
  \node[left,align=right,font=\small] at (-0.03,\y+0.35) {\lab};}
% row 1: juice without a cue -> burst at the juice
\draw[very thick,orange!85!black] plot[domain=0:3,samples=120] (\x,{2.8+0.3+0.75*exp(-((\x-2.08)/0.07)^2)});
\pic[scale=0.8] at (2,2.58) {juice};
% row 2: cue then juice -> burst at the cue, nothing at the juice
\draw[very thick,orange!85!black] plot[domain=0:3,samples=120] (\x,{1.4+0.3+0.75*exp(-((\x-1.08)/0.07)^2)});
\pic[scale=0.85] at (1,1.15) {lamp}; \pic[scale=0.85] at (2,1.15) {juice};
% row 3: cue, juice omitted -> burst at the cue, dip when the juice was due
\draw[very thick,orange!85!black] plot[domain=0:3,samples=120] (\x,{0.3+0.75*exp(-((\x-1.08)/0.07)^2)-0.28*exp(-((\x-2.1)/0.12)^2)});
\pic[scale=0.85] at (1,-0.25) {lamp}; \pic[scale=0.85] at (2,-0.25) {nojuice};
\node[right,font=\scriptsize] at (2.36,0.1) {هبوط};
\node[right,font=\scriptsize] at (2.13,3.75) {مفاجأة!};
\node[left,font=\scriptsize] at (0.97,2.25) {دفقة عند المصباح};
\node[above,font=\scriptsize] at (2,1.75) {لا شيء};
\draw[->,gray!70!black] (0,-0.75) -- (3.05,-0.75) node[right,black,font=\small] {$t$};
\draw[gray!60!black] (1,-0.8) -- (1,-0.7) node[below=2pt,font=\scriptsize] {المصباح، $1\,\text{s}$};
\draw[gray!60!black] (2,-0.8) -- (2,-0.7) node[below=2pt,font=\scriptsize] {العصير، $2\,\text{s}$};
\end{tikzpicture}
\caption{تردد نبضات عصبونات \nt{DOPA} في ثلاث تجارب (تخطيطيًا، بحسب~\cite{Schultz1997}). الخط المنقط هو التردد الخلفي الثابت. المصباح هو الإشارة، والقطرة هي العصير، والقطرة المشطوبة عصير وُعد به ولم يُعطَ. الاستجابة هي خطأ التنبؤ~\eqref{eq:rpe}.}
\label{fig:rpe}
\end{figure}

المبرمج الذي يعرف التعلم المعزز~\cite{SuttonBarto2018} سيعرف هذه الكمية فورًا. الوكيل الذي يتعلم اختيار الأفعال يحتفظ بتقدير $V(s)$: مقدار المكافأة الذي يتوقعه وهو في الحالة $s$. وبعد كل خطوة يقارن التوقع بما ناله:
\begin{empheq}[box=\keybox]{equation}
\delta_t \;=\; r_t \;+\; \gamma\,V(s_{t+1}) \;-\; V(s_t) ,
\label{eq:rpe}
\end{empheq}
ثم يصحح التقدير: $V(s_t)\leftarrow V(s_t)+\alpha\,\delta_t$. هنا $r_t$ المكافأة المنالة، و$\gamma<1$ معامل الخصم (بكم تقل قيمة المكافأة المستقبلية عن الفورية)، و$\alpha$ سرعة التعلم. وتسمى الكمية $\delta_t$ \emph{خطأ الفروق الزمنية}.

وهي تسلك تمامًا سلوك العصبونات الدوبامينية. العصير غير المتوقع: $V$ ما زالت صفرًا، و$r>0$، و$\delta>0$. العصير المتعلَّم: تنبأ المصباح بالمكافأة، فـ$V(s_t)$ قبل العصير تساوي $r$ أصلًا، و$\delta=0$. العصير لم يأتِ: $V(s_t)=r$، لكن $r_t=0$، و$\delta<0$. أما الدفقة فتنتقل إلى المصباح لأن التوقع يقفز صعودًا في لحظة المصباح بالذات: $\gamma V(s_{t+1})-V(s_t)>0$. الخطوات $t,t+1,\dots$ هنا اصطلاح خاص بالخوارزمية: لا يوجد في الجسم مؤقت يعدّ الخطوات، وسنحصل على الكمية نفسها في الزمن المتصل في الفصل~\ref{sec:dofa}.

إذن الدوبامين بثٌّ لإشارة خطأ عددية $\delta$ إلى كل من يجب أن يتعلم بها. في تقريب أول: سلك واحد للدماغ كله وعدد واحد في كل لحظة؛ أما إلى أي حدّ هذا السلك في الواقع حزمة أسلاك فيناقَش في الفصل~\ref{sec:dofaalt}.

\section{قنوات البثّ الأخرى}

للمعدِّلات الأخرى فرضيات عمل فقط، تترجم أدوارها إلى اللغة نفسها لغة الوسائط العامة~\cite{Doya2002}. تعامل معها على أنها فرضيات تحديدًا.

\emph{\nt{NORA}، النورأدرينالين: الكسب.} العصبونات النورأدرينالينية أقل عددًا حتى من الدوبامينية، ببضع عشرات من الآلاف في تقدير خشن، ومخارجها تمتد إلى الدماغ كله تقريبًا. تنشط بحدة حين يحدث أمر غير متوقع أو مهم. يغيّر النورأدرينالين ميل منحنى النقل لدى العصبونات المستقبِلة: تضعف المداخل الضعيفة وتقوى القوية. ويُقاس ذلك هكذا: تُكبت النبضات الخلفية للمستقبِل أكثر مما تُكبت استجابته للمدخل، فتزداد نسبة الإشارة إلى الضجيج~\cite{Foote1975NE}. ويوصف هذا في نموذج بسيط بعدد واحد: إذا استجاب العصبون بتردد $f=F\bigl(g\cdot u\bigr)$ لتيار الدخل $u$، فإن النورأدرينالين يزيد المعامل $g$~\cite{AstonJones2005} (التفاصيل في الفصل~\ref{sec:nora}). الشبكة ذات $g$ الكبير تعمل بـ”تباين“ أعلى وتتخذ قراراتها أسرع؛ وذات $g$ الصغير تعمل ”بضبابية“ وتجرّب خيارات أكثر، كوكيل التعلم المعزز في نمط الاستكشاف.

\emph{\nt{ACET}، الأسيتيل كولين: سرعة التعلم.} يتحكم الأسيتيل كولين في مقدار ما تصغي فيه الشبكة إلى المدخل الحالي مقابل بياناتها المخزنة. عند ارتفاع مستواه تقوى المداخل الآتية من الحواس، وتضعف الوصلات الداخلية ”المستحضِرة“، ويسهل في الوقت نفسه تغيير الأوزان~\cite{Hasselmo2006}. بعبارة تقريبية: هو مفتاح ”كتابة/قراءة“، ومعه سرعة التعلم $\alpha$.

\emph{\nt{SERO}، السيروتونين: أفق التخطيط.} ما ينقله السيروتونين هو الأقل فهمًا. تربطه إحدى الفرضيات بمعامل الخصم $\gamma$ في~\eqref{eq:rpe}: المستوى العالي من السيروتونين يقابل الاستعداد لانتظار مكافأة كبيرة بدل انتزاع صغيرة الآن. تعمل العصبونات السيروتونينية بانتظام وببطء، ببضع نبضات في الثانية أثناء اليقظة، وتكاد تصمت في إحدى مراحل النوم~\cite{JacobsAzmitia1992}.

النواقل العصبية الستة كلها مجموعة في الجدول~\ref{tab:transmitters}.

\begin{table}[t]
\centering
\small
\begin{tabular}{>{\raggedright}p{2.4cm}>{\raggedright}p{3.0cm}>{\raggedright}p{2.0cm}>{\raggedright\arraybackslash}p{5.2cm}}
\toprule
نوع العصبون & حصة العصبونات & الناقل & الدور في الحساب \\
\midrule
\nt{GLUT} (غلوتامات) & الأغلبية & عناوين & وزن موجب، $w>0$ \\
\nt{GABA} (غابا) & $20\text{--}25\,\%$ في القشرة & عناوين & وزن سالب، $w<0$؛ تثبيت الاستقرار \\
\nt{DOPA} (دوبامين) & $\sim 6\cdot10^{-6}$ & بثّ & خطأ التنبؤ $\delta$ \\
\nt{NORA} (نورأدرينالين) & $\sim 6\cdot10^{-7}$ & بثّ & الكسب $g$ (فرضية) \\
\nt{ACET} (أسيتيل كولين) & صغيرة & كلاهما & سرعة التعلم $\alpha$؛ ”كتابة/قراءة“ (فرضية) \\
\nt{SERO} (سيروتونين) & $\sim 3\cdot10^{-6}$ & بثّ & الخصم $\gamma$ (فرضية) \\
\bottomrule
\end{tabular}
\caption{النواقل العصبية الرئيسية في الدماغ بلغة الحساب. العمود الأيمن تبسيط شديد: موثوق للغلوتامات والغابا والدوبامين، أما للبقية ففرضيات.}
\label{tab:transmitters}
\end{table}

\section{المستقبِل يحدد العلامة}
\label{sec:receiver}

لقد خدعتك قليلًا حين قلت: الغلوتامات وزن موجب، والغابا وزن سالب. لا تحمل أي جزيئة علامة في ذاتها. الجزيئة مجرد مفتاح. أما ما يحدث حين تُلتقط فيقرره \emph{القفل}، أي المستقبِل على الخلية المستقبِلة.

أفضل مثال هو الدوبامين. له مستقبلات من عائلتين~\cite{Beaulieu2011}. في إحداهما يسهّل إثارة الخلية، وفي الأخرى يصعّبها. وفي المنطقة التي تختار الأفعال تحمل بعض العصبونات مستقبلات من النوع الأول، وبعضها من النوع الثاني، فتقوّي إشارة الدوبامين نفسها مجموعةً وتضعف أخرى في آن واحد. الأمر أشبه برزمة بثّ واحدة تفسرها عقد الشبكة المختلفة تفسيرات مختلفة.

\begin{keyblock}
\begin{proposition}[المستقبِل يحدد معنى الرسالة]
\label{prop:receptor}
علامة تأثير الناقل العصبي وسرعته لا يحددهما الناقل نفسه، بل المستقبِل الذي يرتبط به. والمستقبلات صنفان. \emph{المستقبلات مؤينة التوجه} هي نفسها قنوات أيونية، وتنفتح في أجزاء من جزء الألف من الثانية: هذا تحكم مباشر في التيار، كبوابة الترانزستور. و\emph{المستقبلات أيضية التوجه} تطلق داخل الخلية سلسلة من التفاعلات الكيميائية وتعمل في مئات من أجزاء الألف من الثانية وأكثر: هذا أقرب إلى الكتابة في سجل إعدادات يغيّر عمل الخلية اللاحق. ناقل العناوين يتكلم في الغالب عبر الأولى، وناقل البثّ في الغالب عبر الثانية.
\end{proposition}
\end{keyblock}

فلماذا يصح أصلًا أن نقول ”الغلوتامات يثير“؟ لأن كل مستقبلات الغلوتامات التي نصادفها في الواقع تقريبًا مثيرة، وكل مستقبلات الغابا تقريبًا مثبطة. ليس هذا قانونًا من قوانين الطبيعة، بل اصطلاح موثوق جدًا، وعليه تقوم القاعدة~\eqref{eq:dalesign}.

\section{ما تحمله معك}

الغالبية الساحقة من العصبونات ناقلُ بياناتٍ بالعناوين، أوزان كل عصبون فيه بعلامة واحدة؛ وأقلية ضئيلة قنواتُ بثّ تنقل إلى أجزاء كبيرة من الدماغ بضع إشارات تحكم عامة. نحو عصبونين من كل مئة ألف، وعليهما يتوقف كيف تتعلم بقية الشبكة كلها وفي أي نمط تعمل. وهذا لا يدهش المهندس: ففي أي حاسوب سجلات التحكم أقل بما لا يقاس من خلايا الذاكرة، ومع ذلك لا تعمل الآلة من دونها.

في الفصل التالي سنختبر ما تستطيع أن تحسبه شبكة من عصبونات \nt{GLUT} وحدها، وهي أكثر الأنواع عددًا.

\section{تمارين}

\begin{exercise}[\lvlA]
\label{ex:01:1}
\emph{الأكوام والأسلاك.} بحسب الجدول~\ref{tab:types} (منتصف المجال بمقياس لوغاريتمي؛ \nt{ACET}: $10^5$ مخرج): (أ)~لو كان كل عصبون إنسانًا، فكم ضعفًا لسكان الأرض تبلغ عصبونات \nt{GLUT}، وبحجم أي مدينة عصبونات البثّ كلها؟ (ب)~كم ضعفًا تزيد حصة نهايات \nt{DOPA} و\nt{SERO} و\nt{NORA} و\nt{ACET} على حصة عصبوناتها؟
\end{exercise}

\begin{exercise}[\lvlA]
\label{ex:01:2}
\emph{العودة إلى الصفر.} يهبط الغلوتامات في الشق كـ$c_0e^{-t/\tau_c}$، حيث $c_0=1\mM$ و$\tau_c=1\ms$؛ ويلحظ المستقبِل التركيز إذا تجاوز $10\,\mu\text{M}$. (أ)~كم يبقى الخط ”مشغولًا“ بعد قذفة واحدة؟ (ب)~المضخات محجوبة جزئيًا، $\tau_c=10\ms$. حتى أي تردد للقذفات يلحق الخط أن يتحرر؟
\end{exercise}

\begin{exercise}[\lvlB]
\label{ex:01:3}
\emph{إلى أين تنتقل الدفقة.} في تجربة الشكل~\ref{fig:rpe} تتوالى بعد المصباح الحالات $s_0\to\dots\to s_{K-1}$ بخطوة $\Delta$، و$T=K\Delta$؛ تصل المكافأة $r$ في الانتقال الأخير، ولحظة المصباح غير متوقعة. أوجد قيم $V(s_k)$ المتعلَّمة واستجابة $\delta$ للمصباح. وبوضع $\gamma=e^{-\Delta/\tau_\gamma}$ أوجد $\tau_\gamma$ عند $T=1\,\mathrm{s}$ و$\Delta=0.1\,\mathrm{s}$ و$\gamma=0.95$.
\end{exercise}

\begin{exercise}[\lvlB]
\label{ex:01:4}
\emph{نمذجة: التعلم من الخطأ.} نمذج سلسلة التمرين~\ref{ex:01:3} مع $K=10$ و$\gamma=0.95$ و$r=1$ بالقاعدة $V(s_t)\leftarrow V(s_t)+\alpha\delta_t$. في أي محاولة $n^*$ تتجاوز الاستجابة للمصباح لأول مرة الاستجابة للعصير عند $\alpha=0.05$ و$0.1$ و$0.2$ و$0.5$؟ وكيف تتعلق $n^*$ بـ$\alpha$؟
\end{exercise}

\begin{exercise}[\lvlB]
\label{ex:01:5}
\emph{تصميم: بثّ عدد واحد.} يجب إيصال العدد $\delta$ إلى كل العصبونات الـ$10^{10}$؛ وكل مستقبِل، بما في ذلك المُرحِّلات، يجب أن يتلقى $10$ نهايات من الطبقة السابقة، وكل حلقة تضيف $1\ms$. كم عصبونًا وكم حلقة في أوفر شجرة مرحِّلات عند $10^4$ مخرج للعصبون (\nt{GLUT}) وعند $10^{5.5}$ (\nt{DOPA})؟
\end{exercise}

\begin{exercise}[\lvlB]
\label{ex:01:6}
\emph{لماذا التوازن حساس.} في الشبكة $80\,\%$ من \nt{GLUT} و$20\,\%$ من \nt{GABA}، وكل عصبون متصل بكل عصبون بأوزان $J_E/N$ و$-J_I/N$؛ $\tau\dot r=-r+Wr+h$. عند $J_E=10$ القيمة الذاتية الوحيدة غير الصفرية لـ$W$ تساوي $0.5$. أوجد نسبة التيار المثبط إلى المثير، وبكم في المئة يكفي إضعاف التثبيط لتنجرف الشبكة إلى انهيار متسلسل.
\end{exercise}

\begin{exercise}[\lvlC]
\label{ex:01:7}
\emph{بحث: هل \nt{SERO} هو $\gamma$؟} بحسب التمرين~\ref{ex:01:3} استجابة عصبونات \nt{DOPA} للمصباح تساوي $re^{-T/\tau_\gamma}$. التأخيرات $T=1,\dots,5\,\mathrm{s}$، ولكل منها $n$ عرضًا، ويُقاس $\ln\delta$ بتشتت $0.5$. كم $n$ يلزم للتمييز بين $\tau_\gamma=2\,\mathrm{s}$ و$2.4\,\mathrm{s}$ (المستوى $0.05$، القوة $0.8$)؟ وأي آلية غير $\gamma$ تعطي الهبوط نفسه مع $T$؟
\end{exercise}

\chapter{ماذا تحسب عصبونات \nt{GLUT}}
\label{sec:glutcomputer}
\begin{chaptergoals}
\item كيف يصنع العصبون المكامل التسريبي ذو العتبة \foreignlanguage{english}{OR} وكاشف تزامن، أي \foreignlanguage{english}{AND} في الزمن؛
\item لماذا لا تحسب عصبونات \nt{GLUT} وحدها إلا دوالّ رتيبة، فلا تستطيع فعاليةٌ أن توقف فعالية؛
\item كيف تتيح أزواج أسلاك ”التشغيل/الإطفاء“ لعصبونات \nt{GLUT} حساب أي دالة منطقية وإعلان الجاهزية؛
\item كيف تحوّل خطوط التأخير فرق الزمن إلى موضع، وكيف تصنع التغذية الراجعة الذاكرة والمؤقتات.
\end{chaptergoals}

\section{قواعد اللعبة}

رأينا في الفصل~\ref{sec:neurontypes} أن الغالبية الساحقة من عصبونات الدماغ من نوع \nt{GLUT}: فهي تثير فقط، وأوزانها موجبة فقط. لنأخذ من هذه العصبونات ما نشاء ونسأل: ماذا تستطيع مثل هذه الشبكة أن تحسب، إذا لم نسمح لأنفسنا إلا بما يحدث في الكائن الحي؟

وليس في الكائن الحي مولّد نبضات ساعة يبدّل كل العناصر في آن واحد. كل عصبون يعمل على حدة، في الزمن المتصل: تأتي النبضات وتذهب متى اتفق، بتأخيرات مختلفة. ثمة \emph{مداخل}، هي العصبونات الحسية التي تنطلق بالضوء أو الصوت أو الضغط، وثمة \emph{مخارج}، هي العصبونات التي تتحكم في العضلات. وبينهما شبكة، لا يخبرها أحد متى تبدأ الحساب ومتى تنهيه. بالنسبة لمهندس الإلكترونيات هذه دارة \emph{لا تزامنية} عناصرها ذوات تأخيرات غير معروفة بدقة مسبقًا.

لنأخذ للعصبون أبسط النماذج التي تعمل بأمانة في الزمن المتصل. الجهد $V$ على غشاء العصبون يساوي صفرًا في الراحة. كل نبضة واردة من العصبون $j$ ترفعه قفزةً بمقدار $w_j\ge0$، ثم ينساب الجهد من تلقاء نفسه عائدًا إلى الصفر بثابت زمني $\tau$:
\begin{empheq}[box=\keybox]{equation}
\tau\,\frac{\dd V}{\dd t} \;=\; -V \;+\; \tau\sum_j w_j \sum_k \delta\bigl(t-t_j^{(k)}-d_j\bigr), \qquad w_j\ge 0 .
\label{eq:glutneuron}
\end{empheq}
هنا $t_j^{(k)}$ لحظات نبضات العصبون $j$، و$d_j$ التأخير على الطريق منه، و$\delta$ دالة دلتا، وهي التي تعني القفزة اللحظية. حين يبلغ $V$ العتبة $\theta$ يطلق العصبون نبضة، ويُصفَّر $V$، ولا يستطيع العصبون أن ينطلق ثانية مدة جزء من الألف من الثانية تقريبًا. الأرقام النموذجية: $\tau$ من أجزاء من جزء الألف من الثانية إلى عشرين جزءًا من الألف في العصبونات المختلفة، والتأخيرات من أجزاء من جزء الألف إلى عشرات الأجزاء من الألف من الثانية. هذا \emph{عصبون مكامل تسريبي}: دارة RC بعتبة. وهو تبسيط مقصود: ففي عصبونات القشرة تطلق فروع المداخل نفسها نبضات محلية~\cite{Smith2013}، ويتصرف العصبون الواحد أقرب إلى شبكة صغيرة متعددة الطبقات (الفصل~\ref{sec:synthesis}). لكن دارة RC واحدة تكفي لأسئلة هذا الفصل.

الفرق الرئيسي عن البوابة المنطقية هو التسريب: لا يتذكر العصبون النبضة إلى الأبد، بل نحو $\tau$، ولا يُجمع إلا ما وصل ضمن هذه النافذة.

\section{\foreignlanguage{english}{OR} و\foreignlanguage{english}{AND}}

لنبدأ بالبوابات. إذا رفعت نبضة واحدة الجهد فوق العتبة، $w\ge\theta$، انطلق العصبون مع كل نبضة من أي مدخل. هذه بوابة \emph{\foreignlanguage{english}{OR}}.

أما \foreignlanguage{english}{AND} فأكثر إثارة. ليكن $w<\theta<2w$: نبضة واحدة لا تكفي، ويلزم نبضتان. لكن السؤال ”هل وصلت النبضتان كلتاهما“ يصبح بلا معنى ما لم نحدد \emph{ضمن أي زمن}. لنفترض أن النبضة الأولى وصلت في اللحظة $0$، والثانية بعد زمن $\Delta$. عند وصول الثانية يكون قد بقي من الأولى $we^{-\Delta/\tau}$، وينطلق العصبون إذا كان $we^{-\Delta/\tau}+w\ge\theta$. ومن هنا نافذة التزامن:
\begin{empheq}[box=\keybox]{equation}
\Delta \;\le\; \Delta_{\max} \;=\; \tau\,\ln\frac{w}{\theta-w} .
\label{eq:window}
\end{empheq}
هذه ليست بوابة \foreignlanguage{english}{AND} منطقية، بل \emph{كاشف تزامن}: \foreignlanguage{english}{AND} في الزمن (الشكل~\ref{fig:coincidence}). عند $\theta=1.5w$ تساوي النافذة $\tau\ln2\approx0.7\tau$. وفي عصبون قشري بـ$\tau=10\ms$ هذا نحو سبعة أجزاء من الألف من الثانية. أما في عصبونات الجهاز السمعي، حيث $\tau$ أقل من جزء من الألف من الثانية، فتضيق النافذة إلى أجزاء من جزء الألف.

\begin{figure}[t]
\centering
\begin{tikzpicture}[>=Stealth,x=0.9cm,y=1.6cm]
\foreach \dx/\lab/\c [count=\i] in {0.5/{تزامنتا}/red!70!black, 2.6/{لم تتزامنا}/blue!60!black}{
 \begin{scope}[xshift={(\i-1)*7.2cm}]
  \draw[->,gray!70!black] (0,0) -- (6.3,0) node[below left,black] {\small $t$};
  \draw[->,gray!70!black] (0,0) -- (0,1.75) node[left,black] {\small $V$};
  \draw[densely dashed,gray!60!black] (0,1.5) -- (6.0,1.5) node[right,black] {\small $\theta$};
  \draw[very thick,\c] (0,0) -- (1,0) -- (1,1)
      plot[domain=1:1+\dx,samples=30] (\x,{exp(-(\x-1)/1.5)});
  \pgfmathsetmacro{\v}{exp(-\dx/1.5)+1}
  \draw[very thick,\c] (1+\dx,{exp(-\dx/1.5)}) -- (1+\dx,{min(\v,1.5)});
  \ifdim\v pt<1.5pt
    \draw[very thick,\c] plot[domain=1+\dx:6,samples=40] (\x,{\v*exp(-(\x-1-\dx)/1.5)});
  \else
    \draw[very thick,\c] (1+\dx,1.5) -- (1+\dx,0) -- (6,0);
    \draw[->,thick,\c] (1+\dx,1.55) -- (1+\dx,1.72) node[above,font=\scriptsize] {نبضة};
  \fi
  \draw[->,thick] (1,-0.35) -- (1,-0.08); \draw[->,thick] (1+\dx,-0.35) -- (1+\dx,-0.08);
  \draw[decorate,decoration={brace,amplitude=3pt,mirror}] (1,-0.42) -- (1+\dx,-0.42) node[midway,below=3pt] {\scriptsize $\Delta$};
  \node[\c] at (3.5,-0.95) {\small \lab};
 \end{scope}}
\end{tikzpicture}
\caption{كاشف التزامن، والعتبة $\theta=1.5w$. على اليسار وصلت النبضة الثانية بعد $\Delta<\Delta_{\max}$، فبلغ المجموع العتبة وانطلق العصبون. وعلى اليمين $\Delta>\Delta_{\max}$: لم يبق من النبضة الأولى شيء يُذكر، فصمت العصبون.}
\label{fig:coincidence}
\end{figure}

وثمة طريقة أخرى للحصول على \foreignlanguage{english}{AND}، عبر المدة لا عبر التوقيت الدقيق. ما دام الضوء مضاءً يطلق العصبون الحسي نبضات متكررة ومنتظمة؛ فإذا لم يكفِ العتبةَ مدخلٌ واحد كهذا وكفاها مدخلان، عمل العصبون ما دام المدخلان كلاهما نشطين. هكذا تحسب الشبكة \emph{بالمستويات}، ولا يهم التوقيت الدقيق لوصول النبضات. ومعظم ما يلي سيكون على هذا النحو.

\section{ما لا يمكن فعله}

لنحاول الآن صنع بوابة \foreignlanguage{english}{NOT}: عصبون يعمل حين يصمت المدخل. لن ننجح: من دون نبضات دخل لا توجد في~\eqref{eq:glutneuron} قفزة واحدة، ويبقى الجهد عند الصفر، تحت العتبة. وماذا عن شبكة من عصبونات كثيرة؟ لا أيضًا، ويُبرهَن ذلك دفعة واحدة لكل الشبكات.

أيسر ما يكون الحديث بالترددات. لتكن $r_i$ تردد نبضات العصبون $i$، و$I_i$ الوارد إليه من المداخل، و$f_i$ منحنى نقله: كيف يتعلق التردد بمجموع الوارد. هذا المنحنى في العصبون المكامل غير متناقص: وارد أكبر يعني نبضات أكثر تواترًا. والترددات في شبكة بلغت الاتزان تحقق المعادلات
\begin{equation}
r_i \;=\; f_i\Bigl(\sum_j w_{ij}\,r_j + I_i\Bigr), \qquad w_{ij}\ge 0 .
\label{eq:ratenet}
\end{equation}

\begin{keyblock}
\begin{proposition}[الرتابة]
\label{prop:monotone}
لتكن الشبكة~\eqref{eq:ratenet} مؤلفة من عصبونات \nt{GLUT} فقط، ولتبدأ من صمت تام. إذا قُوّيت المداخل فلن ينقص التردد المستقر لأي عصبون. كل ما تحسبه هذه الشبكة دوال \emph{رتيبة} (غير متناقصة) في المداخل.
\end{proposition}
\end{keyblock}

البرهان. لنبدأ من الصمت، $r^{(0)}=0$، ولنعوّض الترددات في الطرف الأيمن من~\eqref{eq:ratenet}: $r^{(k+1)}=F\bigl(r^{(k)}\bigr)$. الطرف الأيمن $F$ غير متناقص في أي من متغيراته، لأن الأوزان غير سالبة والدوال $f_i$ غير متناقصة. وبما أن $r^{(1)}\ge0=r^{(0)}$، فبالاستقراء $r^{(k+1)}\ge r^{(k)}$: الترددات تزداد فقط وتصل إلى أصغر حل لـ~\eqref{eq:ratenet}. وإذا استُبدلت بالمداخل $I$ مداخل أكبر $I'$، كان $r^{(k)}(I)\le r^{(k)}(I')$ في كل خطوة، وفي النهاية كذلك. الشبكة الحقيقية تبلغ الاتزان لا بخطوات بل باستمرار، لكن الأمر نفسه صحيح لها: مع الوصلات الموجبة تبقى المتباينة بين تشغيلين، إن تحققت في البداية، متحققة طوال الوقت.

أما للنبضات المفردة فالقضية لا تصح حرفيًا: قد تجعل نبضة دخل زائدة العصبونَ ينطلق أبكر قليلًا، فتضيع نبضته التالية بسبب جزء الألف من الثانية الذي لا يستجيب فيه. لكن الأثر ضعيف وغير موثوق، ولا يصلح أساسًا للحساب. أما للترددات فالرتابة صارمة.

والنتائج هي نتائج أي دارة بلا عواكس. لا \foreignlanguage{english}{NOT}. ولا \foreignlanguage{english}{XOR}: إضافة مدخل ثانٍ لا تستطيع إطفاء المخرج. والأسوأ: \emph{لا يمكن إيقاف الفعالية بفعالية}، بل يمكن فقط إزالة ما يغذيها.

\section{سلكان: التشغيل والإطفاء}

المخرج من هذا المأزق يوحي به الكائن الحي نفسه. في كثير من أعضاء الحس ينتقل المنبه عبر مجموعتين من العصبونات: ففي البصر تستجيب بعض الخلايا لاشتداد الضوء، وأخرى لخفوته~\cite{Schiller1992}. لنسمها عصبونات \emph{التشغيل} وعصبونات \emph{الإطفاء}. فبدل السلك الواحد $x$ تتلقى الشبكة الزوج $(x,\bar x)$، والغياب الذي لا تستطيع حسابه يبلغها به المستشعر نفسه.

مع زوج من الأسلاك تصبح \foreignlanguage{english}{NOT} مجانية: يكفي أن نبدّل السلكين. و\foreignlanguage{english}{AND} و\foreignlanguage{english}{OR} تُصنع كل منهما بعصبونين:
\begin{empheq}[box=\keybox]{equation}
\begin{aligned}
x\wedge y \;&\to\; \bigl(\;x\wedge y,\;\; \bar x\vee\bar y\;\bigr),\\
x\vee y \;&\to\; \bigl(\;x\vee y,\;\; \bar x\wedge\bar y\;\bigr).
\end{aligned}
\label{eq:dualrail}
\end{empheq}
كل العمليات في الطرف الأيمن رتيبة، فلا تُنتهك القضية~\ref{prop:monotone}.

وللشيفرة ثنائية السلك ميزة ثانية للدارة اللاتزامنية، أهم حتى من الأولى. للزوج ثلاث حالات: ”نعم“، و”لا“، و\emph{”لا جواب بعد“} حين يصمت السلكان. يعرف المستقبِل أن الحساب انتهى لا من ساعة، بل من الإشارة نفسها: في كل زوج اشتغل سلك واحد بالضبط. وبين منبه وآخر تصمت المستشعرات، فتعود الشبكة إلى الصمت، مستعدة للمرة التالية. وعلى هذه الحيلة تُبنى في الإلكترونيات اللاتزامنية دارات تعمل صحيحةً مهما كانت تأخيرات عناصرها~\cite{Sparso2001}.

\begin{keyblock}
\begin{proposition}[الحساب اللاتزامني]
\label{prop:dualrail}
إذا ورد كل مدخل على زوج من الأسلاك ”تشغيل/إطفاء“، استطاعت شبكة من عصبونات \nt{GLUT} أن تحسب أي دالة منطقية في المداخل. ويأتي الجواب أيضًا على أزواج من الأسلاك، ويُعرف جاهزيته من أن سلكًا واحدًا اشتغل في كل زوج. لا حاجة إلى ساعة لذلك. والثمن: عدد عصبونات أكبر بمرتين تقريبًا.
\end{proposition}
\end{keyblock}

مثال ذلك \foreignlanguage{english}{XOR}: ”تغيّر منبه واحد بالضبط من منبهين“. السلك المباشر للجواب هو $(x\wedge\bar y)\vee(\bar x\wedge y)$، والسلك المعاكس $(x\wedge y)\vee(\bar x\wedge\bar y)$. هذه ستة عصبونات، لا يحتاج أي منها إلى التثبيط (الشكل~\ref{fig:dualxor}).

\begin{figure}[t]
\centering
\begin{tikzpicture}[>=Stealth,x=1cm,y=1cm,
  nrn/.style={circle,draw,very thick,fill=blue!6,minimum size=0.75cm,inner sep=0pt,font=\small},
  inp/.style={font=\small}]
\node[inp] (X)  at (0,3.3) {$x$};
\node[inp] (Xb) at (0,2.2) {$\bar x$};
\node[inp] (Y)  at (0,1.1) {$y$};
\node[inp] (Yb) at (0,0)   {$\bar y$};
\node[nrn] (A1) at (3.2,3.3) {$2$};
\node[nrn] (A2) at (3.2,2.2) {$2$};
\node[nrn] (A3) at (3.2,1.1) {$2$};
\node[nrn] (A4) at (3.2,0)   {$2$};
\node[nrn] (O)  at (7.2,2.75) {$1$};
\node[nrn] (Ob) at (7.2,0.55) {$1$};
\foreach \s/\t in {X/A1,Yb/A1,Xb/A2,Y/A2,X/A3,Y/A3,Xb/A4,Yb/A4}{\draw[->,thick,blue!60!black] (\s) -- (\t);}
\draw[->,thick,blue!60!black] (A1) -- node[pos=0.45,above,sloped,font=\scriptsize,black] {$x\wedge\bar y$} (O);
\draw[->,thick,blue!60!black] (A2) -- node[pos=0.45,below,sloped,font=\scriptsize,black] {$\bar x\wedge y$} (O);
\draw[->,thick,blue!60!black] (A3) -- node[pos=0.45,above,sloped,font=\scriptsize,black] {$x\wedge y$} (Ob);
\draw[->,thick,blue!60!black] (A4) -- node[pos=0.45,below,sloped,font=\scriptsize,black] {$\bar x\wedge\bar y$} (Ob);
\draw[->,very thick,red!70!black] (O) -- (8.6,2.75) node[right,black] {$x\oplus y$: ”نعم“};
\draw[->,very thick,red!70!black] (Ob) -- (8.6,0.55) node[right,black] {$\overline{x\oplus y}$: ”لا“};
\node[below,font=\small,gray!40!black] at (3.2,-0.55) {\foreignlanguage{english}{AND}: العتبة $2$};
\node[below,font=\small,gray!40!black] at (7.2,-0.55) {\foreignlanguage{english}{OR}: العتبة $1$};
\node[below,font=\small,gray!40!black] at (0,-0.55) {أزواج الأسلاك};
\end{tikzpicture}
\caption{بوابة \foreignlanguage{english}{XOR} بالشيفرة ثنائية السلك، من عصبونات \nt{GLUT} وحدها. العدد في الدائرة هو العتبة بوحدات وزن مدخل واحد. أربعة عصبونات \foreignlanguage{english}{AND} تتعرف على التوليفات الأربع للمداخل، وعصبونا \foreignlanguage{english}{OR} يجمعانها في زوج أسلاك الجواب.}
\label{fig:dualxor}
\end{figure}

\section{التأخيرات بدل الساعة}

للحساب بالزمن تكفي التأخيرات في الأسلاك وكواشف التزامن. وأفضل مثال طريقة الدماغ في تحديد الجهة التي أتى منها الصوت.

يصل الصوت الآتي من مصدر جانبي إلى الأذن القريبة قبل البعيدة. ويتوقف الفرق على الاتجاه: من الصفر حين يكون المصدر أمامنا مباشرة، إلى $0.22\,\text{m}/343\,\text{m/s}\approx0.6\ms$ عند الإنسان حين يكون على الجانب تمامًا. يميّز الدماغ فرقًا من نحو عشرة \emph{أجزاء من المليون} من الثانية~\cite{KlumppEady1956,Grothe2010}، مع أن العصبونات نفسها لا تنطلق بدقة أفضل من أجزاء من جزء الألف من الثانية. كيف؟

لنمدّ من الأذن اليسرى سلكًا على طول صف من كواشف التزامن من اليسار إلى اليمين، ومن الأذن اليمنى سلكًا من اليمين إلى اليسار (الشكل~\ref{fig:itd}). تجري الإشارة في السلك بسرعة $v$. إلى الكاشف الواقع على بعد $x$ من الطرف الأيسر لصف طوله $L$، تستغرق الإشارة اليسرى زمنًا $x/v$، واليمنى $(L-x)/v$. فإذا وصل الصوت إلى الأذن اليمنى متأخرًا بـ$\delta t$، التقت الإشارتان في آن واحد عند النقطة التي
\begin{empheq}[box=\keybox]{equation}
\frac{x}{v} \;=\; \frac{L-x}{v} + \delta t
\quad\Longrightarrow\quad
x \;=\; \frac{L}{2} + \frac{v\,\delta t}{2} .
\label{eq:itd}
\end{empheq}
لن ينطلق إلا الكاشف الذي بلغته الإشارتان ضمن النافذة~\eqref{eq:window}. ورقم الكاشف المنطلق هو اتجاه المصدر. لقد تحوّل الفرق في الزمن إلى \emph{موضع}.

\begin{figure}[t]
\centering
\begin{tikzpicture}[>=Stealth,
  nrn/.style={circle,draw,very thick,fill=blue!6,minimum size=0.7cm,inner sep=0pt}]
\foreach \k in {0,...,6}{\node[nrn] (D\k) at (1.4*\k,0) {\scriptsize $2$};}
\node[nrn,fill=red!15] at (D4) {\scriptsize $2$};
\draw[->,thick,blue!60!black] (-1.4,0.9) node[left,black] {\small الأذن اليسرى} -- (9.2,0.9);
\draw[->,thick,orange!85!black] (9.8,-0.9) node[right,black] {\small الأذن اليمنى} -- (-0.8,-0.9);
\foreach \k in {0,...,6}{
  \draw[->,blue!60!black] (1.4*\k,0.9) -- (D\k.north);
  \draw[->,orange!85!black] (1.4*\k,-0.9) -- (D\k.south);}
\draw[->,thick,red!70!black] (D4.north east) ++(0.1,0.05) -- ++(0.5,0.45) node[above right,font=\small,black] {انطلق};
\node[font=\small,gray!40!black] at (4.2,-1.5) {يزداد التأخير $\longleftarrow$ \hspace{2.2cm} $\longrightarrow$ يزداد التأخير};
\pic[scale=1.4] at (-2.3,1.75) {speaker};
\node[font=\scriptsize,align=center] at (-2.3,2.35) {المصدر على اليسار};
\end{tikzpicture}
\caption{تحديد اتجاه الصوت. تجري إشارتا الأذنين إحداهما نحو الأخرى؛ وكل دائرة كاشف تزامن بعتبة $2$. إذا وصل الصوت إلى الأذن اليمنى متأخرًا، التقت الإشارتان يمين المنتصف، بحسب الصيغة~\eqref{eq:itd}. مخرج الشبكة هو رقم العصبون المنطلق. والمصدر هنا على اليسار: وصل الصوت إلى الأذن اليسرى أولًا.}
\label{fig:itd}
\end{figure}

لا ساعة هنا، ولا تثبيط، ولا حتى شيفرة ثنائية السلك: لا تجيب الشبكة بـ”نعم“ أو ”لا“، بل تشير إلى عصبون واحد من بين كثيرين. ويبدو أن دارة كهذه من خطوط التأخير وكواشف التزامن تعمل فعلًا في جذع الدماغ السمعي لدى الطيور~\cite{Carr1990}. ودقة أجزاء المليون من الثانية لا تأتي من دقة كل عصبون، بل من كثرة الكواشف ونظر كل منها إلى تأخيره الخاص. ولم يُعثر لدى الثدييات على صف كهذا من التأخيرات: هناك يبدو أن المهمة نفسها تحلها كواشف تزامن يضبطها تثبيط محسوب التوقيت بدقة من عصبونات \nt{GLYC}~\cite{Brand2002,Grothe2010}. وكيف يضيّق التثبيط نافذة التزامن سنبحثه في الفصل~\ref{sec:gaba} على مثال \nt{GABA}؛ والغليسين يثبط بالطريقة نفسها، إذ يفتح لدى المستقبِل قنوات الكلور.

\section{الذاكرة}

يحتاج الحاسوب إلى ذاكرة. والذاكرة في شبكة كهذه فعالية تغذي نفسها بنفسها. لنأخذ مجموعة من عصبونات \nt{GLUT} متصلة بعضها ببعض اتصالًا قويًا، ولنصفها بتردد واحد $r$ (الشكل~\ref{fig:recurrentgroup}أ). في الاتزان $r=f(wr+I)$، حيث $w$ قوة التغذية الراجعة للمجموعة على نفسها، و$I$ الوارد من الخارج. لنحل هذه المعادلة بيانيًا بتقاطع المنحنى $f(wr+I)$ مع المستقيم $r$ (الشكل~\ref{fig:bistable}).

\begin{figure}[t]
\centering
% (b): tau dr/dt = -r + f(w r + I), f as in fig:bistable, w=1, I=0.6; Euler dt=0.001 tau.
\begin{tikzpicture}[>=Stealth,
  nrn/.style={circle,draw,thick,fill=blue!6,minimum size=0.5cm,inner sep=0pt}]
\begin{scope}
\node[font=\small] at (-0.5,1.55) {(أ)};
\foreach \k in {1,...,6}{\node[nrn] (n\k) at ({1.4+1.0*cos(60*\k)},{1.0*sin(60*\k)}) {};}
\foreach \a/\b in {1/2,2/3,3/4,4/5,5/6,6/1,1/4,2/5,3/6}{\draw[<->,blue!60!black] (n\a) -- (n\b);}
\foreach \k in {2,3,4}{\draw[->,thick,green!45!black] ($(n\k)+(-0.95,0)$) -- (n\k);}
\node[font=\small,green!45!black] at (-0.35,-0.45) {$I$};
\node[font=\Large] at (3.1,0) {$\approx$};
\node[nrn,minimum size=0.85cm,font=\small] (R) at (4.35,-0.1) {$r$};
\draw[->,thick,blue!60!black] (R) edge[loop above,looseness=6] node[above,font=\small] {$w$} (R);
\draw[->,thick,green!45!black] (3.55,-0.95) node[left,font=\small] {$I$} -- (R);
\end{scope}
\begin{scope}[xshift=6.3cm,y=2.6cm,x=1.05cm]
\node[font=\small] at (-0.6,1.25) {(ب)};
\draw[->,gray!70!black] (0,0) -- (7.3,0) node[below left,font=\small,black] {$t/\tau$};
\draw[->,gray!70!black] (0,0) -- (0,1.12) node[left,font=\small,black] {$r$};
\foreach \t in {0,2,4,6}{\draw[gray!60!black] (\t,-0.02) -- (\t,0.02); \node[below,font=\scriptsize] at (\t,0) {\t};}
\draw[densely dotted,gray!60!black] (0,0.5) -- (7,0.5) node[right,font=\scriptsize,black,align=left] {عتبة\\الكتابة};
\fill[green!45!black,opacity=0.25] (1,-0.2) rectangle (2,-0.13);
\fill[gray!60!black,opacity=0.35] (1,-0.29) rectangle (1.5,-0.22);
\draw[very thick,blue!60!black] plot coordinates {(0.0,0.002) (0.1,0.002) (0.2,0.002) (0.3,0.002) (0.4,0.002) (0.5,0.002) (0.6,0.002) (0.7,0.002) (0.8,0.002) (0.9,0.002) (1.0,0.002) (1.1,0.083) (1.2,0.165) (1.3,0.242) (1.4,0.313) (1.5,0.378) (1.6,0.437) (1.7,0.491) (1.8,0.539) (1.9,0.583) (2.0,0.623) (2.1,0.643) (2.2,0.665) (2.3,0.688) (2.4,0.71) (2.5,0.732) (2.6,0.753) (2.7,0.773) (2.8,0.792) (2.9,0.81) (3.0,0.826) (3.2,0.855) (3.4,0.88) (3.6,0.9) (3.8,0.917) (4.0,0.931) (4.4,0.953) (4.8,0.967) (5.2,0.977) (5.6,0.984) (6.0,0.988) (6.5,0.992) (7.0,0.994)};
\draw[very thick,gray!60!black,densely dashed] plot coordinates {(0.0,0.002) (0.1,0.002) (0.2,0.002) (0.3,0.002) (0.4,0.002) (0.5,0.002) (0.6,0.002) (0.7,0.002) (0.8,0.002) (0.9,0.002) (1.0,0.002) (1.1,0.083) (1.2,0.165) (1.3,0.242) (1.4,0.313) (1.5,0.378) (1.6,0.358) (1.7,0.336) (1.8,0.313) (1.9,0.291) (2.0,0.269) (2.2,0.228) (2.4,0.191) (2.6,0.159) (2.8,0.133) (3.0,0.11) (3.2,0.091) (3.4,0.076) (3.6,0.063) (3.8,0.052) (4.0,0.043) (4.4,0.03) (4.8,0.021) (5.2,0.015) (5.6,0.011) (6.0,0.008) (6.5,0.006) (7.0,0.004)};
\node[font=\scriptsize,blue!60!black,anchor=west] at (3.3,0.8) {وارد لمدة $1\tau$: كُتب};
\node[font=\scriptsize,gray!50!black,anchor=west] at (2.3,0.3) {وارد لمدة $0.5\tau$: انطفأ};
\end{scope}
\end{tikzpicture}
\caption{مجموعة متصلة بنفسها اتصالًا قويًا. (أ)~عصبونات \nt{GLUT} كثيرة يثير بعضها بعضًا تتصرف كعصبون واحد بتردد $r$ يثير نفسه بقوة $w$؛ ويأتيه من الخارج وارد $I$. (ب)~تردد المجموعة مع الزمن عند $w=1$ وبالمنحنى $f$ نفسه الذي في الشكل~\ref{fig:bistable}. يُقدَّم الوارد $I=0.6$ طوال الزمن المعلَّم بالشريط تحت المحور. إذا دام $1\tau$، لحقت المجموعة أن تتجاوز النقطة غير المستقرة $r=0.5$، فتتقد وتبقى متقدة بعد انتهاء الوارد. وإذا دام $0.5\tau$ فقط، لم تبلغ تلك النقطة فتنطفئ عائدة: لكتابة بت يجب أن يدوم الوارد أكثر من $0.7\tau$ تقريبًا.}
\label{fig:recurrentgroup}
\end{figure}

\begin{figure}[t]
\centering
\begin{tikzpicture}[>=Stealth,x=5cm,y=3.2cm]
\draw[->,gray!70!black] (0,0) -- (1.1,0) node[below left,black] {\small $r$};
\draw[->,gray!70!black] (0,0) -- (0,1.1);
\draw[thick,gray!60!black] (0,0) -- (1.05,1.05) node[right,black] {\small $r$};
\draw[very thick,blue!60!black] plot[domain=0:1.05,samples=80] (\x,{1/(1+exp(-(\x-0.5)/0.08))});
\draw[very thick,red!70!black,densely dashed] plot[domain=0:1.05,samples=80] (\x,{1/(1+exp(-(\x+0.3-0.5)/0.08))});
\fill[blue!60!black] (0.002,0.002) circle (0.018);
\draw[blue!60!black,fill=white,thick] (0.5,0.5) circle (0.018);
\fill[blue!60!black] (0.998,0.998) circle (0.018);
\node[below right,font=\small] at (0.02,0.0) {مطفأ};
\node[right,font=\small] at (0.52,0.46) {غير مستقر};
\node[below,font=\small] at (0.99,0.96) {مشغَّل};
\node[anchor=east,font=\small,red!70!black] at (0.19,0.62) {$f(wr+I)$};
\node[anchor=west,font=\small,blue!60!black] at (0.45,0.1) {$f(wr)$};
\end{tikzpicture}
\caption{ذاكرة من مجموعة عصبونات \nt{GLUT} ذات تغذية راجعة قوية. المنحنى المتصل من دون وارد خارجي: تقاطعان مستقران مع المستقيم $r$، ”مطفأ“ و”مشغَّل“، وتقاطع غير مستقر بينهما. والوارد القصير $I$ (المنحنى المتقطع) لا يُبقي إلا التقاطع العلوي.}
\label{fig:bistable}
\end{figure}

إذا كانت التغذية الراجعة قوية كانت التقاطعات ثلاثة: السفلي هو الصمت، والعلوي هو المجموعة متقدة بكامل طاقتها، والأوسط غير مستقر. حصلنا على عنصر بحالتين مستقرتين: قلّاب (\foreignlanguage{english}{flip-flop}). وكتابة الواحد فيه سهلة: وارد قصير يرفع المنحنى، فيختفي التقاطع السفلي، وتتقد المجموعة، وتبقى متقدة بعد انتهاء الوارد (الشكل~\ref{fig:recurrentgroup}ب). أما الوارد القصير جدًا فلا يلحق أن يبلغ بالمجموعة النقطة غير المستقرة، فتنطفئ عائدة. هذه الفعالية الذاتية الاستمرار، التي تدوم ثوانيَ بعد زوال المنبه، تُرصد فعلًا في الدماغ، وتُعدّ أساس الذاكرة قصيرة الأمد~\cite{Wang2001}. لكنها ليست الطريقة الوحيدة لتذكر الثواني. يمكن أيضًا حفظ المكتوب في متغير بطيء في المشابك نفسها: بعد رشقة من النبضات يدوم التيسير، كما في المشبك ”الشَّرطة“ في الفصل~\ref{sec:code}، نحو ثانية، وقد تكاد الشبكة تصمت بين ومضات قصيرة، فهي ليست قلّابًا بل مكثفًا مشحونًا~\cite{Mongillo2008}. وهذه الفترات ”الصامتة“ أثناء الاحتفاظ بالذاكرة مرصودة~\cite{Stokes2015}، والتخزين من دون نبضات دائمة أرخص طاقةً. أما أي الطريقتين تستعمل القشرة وفي أي حالة، فموضع نقاش.

وكيف نمحو؟ بحسب القضية~\ref{prop:monotone} لا سبيل إلى ذلك بالفعالية؛ لا يبقى إلا أن \emph{نزيل} شيئًا. الطريقة الأولى إزالة الدعم: إذا كانت المجموعة لا تتقد إلا مع وارد من مجموعة أخرى، انطفأت الذاكرة معه. والثانية التعب. ففي المشابك الحقيقية ينضب مخزون الناقل العصبي عند العمل الطويل~\cite{Tsodyks1997}، وتنمو في العصبونات تيارات بطيئة تخفض الاستثارة. فتضعف التغذية الراجعة، ويختفي التقاطع العلوي، وتنطفئ المجموعة وحدها بعد ثوانٍ. والنتيجة ذاكرة تُشغَّل بإشارة ولا تُطفأ إلا بالزمن.

\section{المتتاليات}

بدل الساعة يمكن أن يكون لدينا \emph{مؤقت يطلقه حدث}. لنصل مجموعات من عصبونات \nt{GLUT} في سلسلة: الأولى تثير الثانية، والثانية الثالثة، وهكذا. تُشعل إشارة خارجية المجموعة الأولى، فتجري موجة الفعالية على طول السلسلة، كل حلقة بعد سابقتها بتأخير أو تأخيرين، ومرة واحدة فقط: فالعصبونات لا تستجيب مباشرة بعد النبضة، والموجة قد مضت.

تقيس هذه السلسلة الزمن منذ الحدث: إذا انطلقت الحلقة رقم $k$ فقد مضى منذ الإطلاق $k$ تأخيرات. ويمكن توصيل مخارجها بالعضلات للحصول على متتالية حركات تتكشف من تلقاء نفسها. لدى الطيور المغردة، أثناء الغناء، تنطلق عصبونات \nt{GLUT} مفردة في إحدى مناطق الدماغ بالتتابع الصارم، كل منها في لحظته من الأغنية، وهذا يشبه كثيرًا سلسلة كهذه~\cite{Hahnloser2002}.

وعلى خلاف الساعة، تصمت السلسلة ما لم تُطلق، وتعمل مرة واحدة، ولا تقيس الزمن إلا لمن وُصل بها. ولا يزال لا إيقاع مشتركًا للشبكة.

\section{ما الذي ينقص}

من عصبونات \nt{GLUT} وحدها، بلا تثبيط، يمكن صنع الكثير. لكن ثلاثة أشياء تنقص، وكلها بسبب القضية~\ref{prop:monotone}.

\emph{الاختيار.} يجب على الشبكة أن تختار اتجاهًا واحدًا، وفعلًا واحدًا. فإذا تساوى منبهان تقريبًا، انطلق في شبكة الشكل~\ref{fig:itd} المخرجان كلاهما، ولا شيء يكبت الأضعف لصالح الأقوى.

\emph{التصفير.} لا يمكن محو الذاكرة بإشارة.

\emph{الاستقرار.} في شبكة لا تنشأ فيها الوصلات بحسب مخطط مهندس، لا مفر من حلقات التغذية الراجعة الموجبة، وقد تتصاعد فيها الفعالية تصاعدًا متسلسلًا (الفصل~\ref{sec:neurontypes}). والكابح الوحيد هو التعب نفسه، البطيء والخشن.

والعلل الثلاث يعالجها دواء واحد: عصبون يطفئ النبضة بنبضة. هذا هو \nt{GABA}، والحاجة إليه ليست للحساب، بل للاختيار، والتصفير، وإبقاء الحساب تحت السيطرة.

\section{تمارين}

\begin{exercise}[\lvlA]
\label{ex:02:1}
\emph{صف الكواشف للصوت.} تجري إشارتا الأذنين في أسلاك الشكل~\ref{fig:itd} بسرعة $5\,\text{m/s}$، وأكبر فرق في الوصول $0.64\ms$. الكواشف عصبونات~\eqref{eq:glutneuron} بـ$\tau=0.5\ms$ و$\theta=1.5w$، مرتبة بحيث يميّز المتجاوران $10\,\mu\text{s}$. كم كاشفًا في الصف، وكم منها ينطلق لصوت واحد؟
\end{exercise}

\begin{exercise}[\lvlB]
\label{ex:02:2}
\emph{المداخل غير المتساوية تزيح النافذة.} كاشف تزامن~\eqref{eq:glutneuron} بـ$\tau=0.5\ms$ و$\theta=1.5$ يتلقى نبضة من كل من مدخلين بوزنين $1$ و$0.8$. أوجد نافذة التزامن ومركزها. بكم جزءًا من المليون من الثانية يخطئ صف الشكل~\ref{fig:itd} إذا كان مدخل إحدى الأذنين أضعف من الأخرى بـ$20\,\%$؟
\end{exercise}

\begin{exercise}[\lvlB]
\label{ex:02:3}
\emph{أي الشبكات لا تتذبذب.} الشبكة $\tau\dot r_i=-r_i+f_i\bigl(\textstyle\sum_jw_{ij}r_j+I_i\bigr)$ ذات الدوال $f_i$ غير المتناقصة و$w_{ij}\ge0$ عند $i\ne j$ رتيبة ولا تتذبذب تذبذبًا مستقرًا. برهن: بالتعويض $r_i\to\pm r_i$ تؤول إليها الشبكة التي في كل دورة فيها عدد زوجي من الوصلات المثبطة. هل يمكن أن يتذبذب التثبيط المتبادل بين مجموعتين؟ وحلقة \foreignlanguage{english}{\nt{GLUT}--\nt{GABA}}؟
\end{exercise}

\begin{exercise}[\lvlB]
\label{ex:02:4}
\emph{تصميم: جامع ثنائي السلك.} كل بت ينتقل على زوج من الأسلاك $(x,\bar x)$. ابنِ جامعًا كاملًا لثلاثة بتات يُخرج الزوجين $(S,\bar S)$ للمجموع و$(C,\bar C)$ للمنقول، من عصبونات \nt{GLUT}، مع بيان الأوزان والعتبات. اكتفِ بثمانية عصبونات. كم عصبونًا تتطلب دارة من عناصر \foreignlanguage{english}{AND} و\foreignlanguage{english}{OR} كما في الشكل~\ref{fig:dualxor}؟
\end{exercise}

\begin{exercise}[\lvlB]
\label{ex:02:5}
\emph{نمذجة: كم تدوم الكتابة.} مجموعة الشكل~\ref{fig:recurrentgroup}: $\tau\dot r=-r+f(r+I)$، $f(u)=1/\bigl(1+e^{-(u-0.5)/0.08}\bigr)$، وهي قبل الكتابة في الحالة السفلى؛ ويُقدَّم الوارد $I$ مدة $T$. أوجد عدديًا أصغر $T$ يكتب البت عند $I=0.6$، والعتبة $I_c$ التي لا يُكتب البت دونها مهما كان $T$.
\end{exercise}

\begin{exercise}[\lvlA]
\label{ex:02:6}
\emph{السلسلة مؤقتًا.} حلقة السلسلة $100$ عصبون من \nt{GLUT}، وتأخيرها $2\ms$ بتشتت مستقل $0.2\ms$. (أ)~كم عصبونًا يلزم لقياس $1\,\text{s}$، وما الخطأ النسبي؟ وكيف يتعلق بالفترة؟ (ب)~خطأ المؤقت الحقيقي $5\,\%$ لأي فترة. أي مصدر للتشتت يعطي ذلك؟
\end{exercise}

\begin{exercise}[\lvlC]
\label{ex:02:7}
\emph{بحث: قلّاب أم مكثف.} $100$ عصبون تحفظ بتًا مدة $10\,\text{s}$. القلّاب: كل عصبون يطلق $20$ نبضة في الثانية. المكثف: التيسير $u$ يهبط بـ$\tau_F=1\,\text{s}$، ويُقرأ البت عند $u\ge u_{\max}/2$، والإنعاش رشقة من ثلاث نبضات لكل عصبون. كم مرة يكون المكثف أوفر، وماذا يحدث للبت بعد إسكات المجموعة مدة $200\ms$؟
\end{exercise}

\chapter{\nt{GLUT} و\nt{GABA} معًا}
\label{sec:gaba}
\begin{chaptergoals}
\item كيف يجعل الحظر $a\wedge\bar b$ من \nt{GLUT} و\nt{GABA} منظومة تامة بسلك واحد لكل بت؛
\item لماذا يقسم التثبيط التحويلي ولا يطرح، وكيف يضيّق التثبيط نافذة التزامن؛
\item كيف يختار التثبيط المتبادل مع $\beta>1$ فائزًا واحدًا، وكيف تصفّر إشارةٌ الذاكرة؛
\item متى يُبقي التثبيط السريع الشبكة مستقرة، ومتى يحوّلها التثبيط البطيء إلى إيقاع.
\end{chaptergoals}

\section{قواعد اللعبة}

شبكة عصبونات \nt{GLUT} وحدها لا تستطيع أن تختار، ولا أن تصفّر الذاكرة، ولا أن تمنع الفعالية من التصاعد المتسلسل، وكل ذلك بسبب الرتابة (القضية~\ref{prop:monotone}). لنضف إليها النوع الثاني من حيث العدد: \nt{GABA}. القواعد كما كانت: لا شيء إلا ما يحدث في الكائن الحي، وكل شيء يجري في الزمن المتصل.

نموذج العصبون هو نفسه~\eqref{eq:glutneuron}، لكن النبضة من عصبون \nt{GABA} لا ترفع جهد المستقبِل بل تخفضه. للأوزان الآن علامتان، والمرسِل يحدد العلامة، بحسب القاعدة~\eqref{eq:dalesign}.

بعض وسائط الدارة. عصبونات \nt{GABA} في القشرة من الخُمس إلى الربع~\cite{Hendry1987}. مخارجها قصيرة: تثبط جيرانها لا المناطق البعيدة. يصل التثبيط بعد جزء من الألف من الثانية ويدوم بضعة أجزاء من الألف. ومن دون مدخل يصمت عصبون \nt{GABA}: فلكي يثبط أحدًا يجب أن يتلقى هو نفسه إثارة، عادة من عصبونات \nt{GLUT} في الشبكة نفسها.

لذا يجب أن تُصنع الوصلة السالبة من عصبون \nt{GLUT} $A$ إلى العصبون $B$ عبر وسيط: $A$ يثير عصبون \nt{GABA}، وهذا يثبط $B$. ثمن قلب العلامة عصبون واحد وتأخير واحد، نحو جزء من الألف من الثانية.

في التثبيط لا يهم فقط \emph{من أين} تأتي الوصلة، بل \emph{أين} تستقر: فموضع الاستقرار يحدد إلى حد بعيد ما تفعله الوصلة~\cite{Markram2004}. مخارج \nt{GLUT} تستقر في الغالب على فروع الدخل لدى المستقبِل، أي \emph{تغصناته}، وتحمل البيانات: كل مدخل كهذا حدّ واحد في المجموع. والمخرج على جسم العصبون يؤثر في المجموع كله: هكذا يقسم كثير من عصبونات \nt{GABA} السريعة استجابة المستقبِل ويحدد متى ينطلق. والمخرج على الموضع الذي تولد فيه نبضة المستقبِل هو الكلمة الأخيرة، حظر على المخرج كله دفعة واحدة. أما المخرج على نهاية سلك عصبون آخر فيغيّر احتمال القذف $p$ لخط دخل واحد، دون المساس بغيره~\cite{Rudomin1999}؛ وهكذا تعمل، مثلًا، بوابة الألم في الفصل~\ref{sec:pain}. وخارج الدماغ تستقر المخارج على الألياف العضلية (الفصل~\ref{sec:muscles}) وعلى الأعضاء (الفصل~\ref{sec:chemoutput})، وبعضها لا يستقر على شيء بل يقذف الناقل العصبي في حجم النسيج أو في الدم (الفصلان~\ref{sec:serocomputer} و\ref{sec:chemoutput}).

\section{\foreignlanguage{english}{NOT} والحظر}

هل نستطيع الآن صنع \foreignlanguage{english}{NOT}؟ تقريبًا. العصبون الذي يعمل حين يصمت المدخل يجب أن يأخذ إثارته من مكان ما. وهي موجودة في الدماغ الحي: فعلى كل عصبون ترد باستمرار فعالية خلفية من آلاف غيره. لتبقِ الخلفية العصبون $Y$ عاملًا، وليثبطه المدخل $x$ عبر وسيط \nt{GABA}. هذه \foreignlanguage{english}{NOT}.

لكن الأنفع هو \emph{الحظر}: ”$a$، إلا عند $b$“:
\begin{empheq}[box=\keybox]{equation}
y \;=\; a \wedge \bar b .
\label{eq:veto}
\end{empheq}
العصبون $Y$ يتلقى الإثارة من $a$ والتثبيط من عصبون \nt{GABA} يثيره $b$. لا حاجة هنا إلى الخلفية: فالإثارة يعطيها $a$ نفسه. ومن الحظر و\foreignlanguage{english}{OR} يُبنى كل شيء. مثلًا \foreignlanguage{english}{XOR}: $x\oplus y=(x\wedge\bar y)\vee(\bar x\wedge y)$، أي عصبونا حظر، ووسيطا \nt{GABA}، وعصبون \foreignlanguage{english}{OR} واحد.

\begin{keyblock}
\begin{proposition}[تمام \nt{GLUT}+\nt{GABA}]
\label{prop:gabacomplete}
شبكة من عصبونات \nt{GLUT} و\nt{GABA} تستطيع أن تحسب أي دالة منطقية في المداخل، وكل بت ينتقل على سلك واحد. ولم تعد مستشعرات ”التشغيل/الإطفاء“ من الفصل~\ref{sec:glutcomputer} لازمة لذلك. وإذا وجب أن تعطي الدالة واحدًا حين تصمت كل المداخل، لزم مصدر فعالية خلفية.
\end{proposition}
\end{keyblock}

البرهان: الحظر~\eqref{eq:veto} مع \foreignlanguage{english}{OR} تامّ وظيفيًا إذا توفر واحد ثابت، لأن \foreignlanguage{english}{NOT} $x$ هو الحظر ”واحد، إلا عند $x$“. أما الدوال التي تعطي صفرًا حين تصمت كل المداخل فتُبنى من الحظر و\foreignlanguage{english}{OR} حتى من دون الواحد.

\section{اتجاه الحركة}

الحظر في الزمن، مثل \foreignlanguage{english}{AND} في الزمن في الفصل~\ref{sec:glutcomputer}، يعطي كاشفًا لاتجاه الحركة.

ليكن مستشعران متجاوران، $A$ و$B$، على بعد $\Delta x$ أحدهما من الآخر. العصبون المخرج $Y$ يتلقى الإثارة من $B$ والتثبيط من $A$ عبر وسيط \nt{GABA}، بتأخير $d$ (الشكل~\ref{fig:direction}). بقعة الضوء الجارية بسرعة $v$ من $A$ إلى $B$ تمر بـ$A$ في اللحظة $0$، فيصل التثبيط إلى $Y$ في اللحظة $d$؛ وتمر بـ$B$ في اللحظة $\Delta x/v$، وحينها تصل الإثارة. فإذا تطابقت هاتان اللحظتان ضمن النافذة~\eqref{eq:window}، أطفأ التثبيط الإثارة وصمت $Y$. ويحدث هذا عند سرعة قريبة من
\begin{empheq}[box=\keybox]{equation}
v^* \;=\; \frac{\Delta x}{d} .
\label{eq:vstar}
\end{empheq}
أما عند الحركة من $B$ إلى $A$ فتصل الإثارة من $B$ في اللحظة $0$، ولا يصل التثبيط من $A$ إلا في اللحظة $\Delta x/v+d$، وقد انطلق $Y$ عندئذ.

\begin{figure}[t]
\centering
\begin{tikzpicture}[>=Stealth,
  nrn/.style={circle,draw,very thick,fill=blue!6,minimum size=0.9cm,inner sep=0pt},
  gab/.style={nrn,fill=red!10},
  inh/.style={-{Bar[width=7pt]},thick,red!70!black}]
\node[nrn] (A) at (0,2) {$A$};
\node[nrn] (B) at (3,2) {$B$};
\node[gab] (G) at (0.9,0.6) {\small \nt{GABA}};
\node[nrn] (Y) at (3,-0.6) {$Y$};
\draw[->,thick] (A) -- (G);
\draw[inh] (G) -- node[below left=-2pt] {\scriptsize التأخير $d$} (Y);
\draw[->,thick] (B) -- (Y);
\draw[->,very thick,red!70!black] (Y) -- (4.6,-0.6) node[right,black] {\small المخرج};
\draw[->,thick,orange!85!black] (-0.6,2.9) -- (3.6,2.9) node[right,black] {\small يصمت};
\draw[<-,thick,orange!85!black] (-0.6,3.4) -- (3.6,3.4) node[right,black] {\small يستجيب};
\foreach \yy/\dd in {2.9/-1,3.4/1}{
  \foreach \k in {1,2,3}{\draw[orange!70!yellow,line width=0.8pt] (1.5+\dd*0.12+\dd*0.13*\k,\yy+0.1-0.1*\k+0.1) -- ++(\dd*0.25,0);}
  \fill[yellow!70!orange,draw=orange!70!black] (1.5,\yy) circle (0.14);}
\node[font=\scriptsize,align=center] at (1.5,3.95) {بقعة ضوء};
\end{tikzpicture}
\caption{كاشف اتجاه الحركة. $B$ يثير $Y$، و$A$ يثبطه عبر وسيط \nt{GABA} بتأخير $d$. عند الحركة من $A$ إلى $B$ بسرعة قريبة من $\Delta x/d$ يطفئ التثبيط الإثارة؛ وعند الحركة المعاكسة يتأخر. البقعة الصفراء ذات الشُّرَط هي الضوء الجاري فوق المستشعرات.}
\label{fig:direction}
\end{figure}

في الشبكية يبدو أن الانتقائية لاتجاه الحركة تنشأ بهذه الطريقة بالذات: بتثبيط غير متناظر من الجهة التي لا ينبغي أن تستدعي الحركة منها استجابة~\cite{Barlow1965}.

\section{طرح أم قسمة}

حتى الآن كان التثبيط عندنا يطرح. والحقيقة أنه أدق من ذلك.

المشبك يفتح قناة، أي يشغّل موصلية. في المشبك المثير يقع جهد الاتزان فوق العتبة بكثير، نحو $70\mV$ فوق الراحة: القناة المفتوحة تشد الجهد إلى أعلى. وفي مشبك \nt{GABA} المثبط تمرر القناة الكلور، الذي يقع جهد اتزانه قرب الراحة، فتشد القناة المفتوحة الجهد لا إلى أسفل بل \emph{نحو الراحة}. إذا قسنا الجهد من الراحة، فالجهد المستقر لغشاء بموصلية تسريب $g_L$ وموصلية مثيرة $g_E$ ومثبطة $g_I$ يساوي
\begin{empheq}[box=\keybox]{equation}
V_\infty \;=\; \frac{g_E\,E_E}{g_L+g_E+g_I} ,
\label{eq:shunt}
\end{empheq}
حيث $E_E\approx70\mV$ جهد اتزان المشبك المثير. هذا قانون أوم لمقسّم الجهد: $g_E$ تشد نحو $E_E$، و$g_L$ و$g_I$ نحو الصفر.

لا تقف $g_I$ في البسط بإشارة سالبة، بل في المقام: التثبيط لا \emph{يطرح} من الإثارة، بل \emph{يقسمها} (الشكل~\ref{fig:shunt}). يسمى هذا التثبيط تثبيطًا تحويليًا: قنوات الكلور المفتوحة تقصّر دارة الغشاء. وبالنسبة للشبكة هذا ضبط للكسب. فإذا تناسبت $g_I$ مع مجموع فعالية الجيران، استجاب كل عصبون لا للقوة المطلقة لمدخله بل لحصتها من الفعالية الكلية. وهذه القسمة على الفعالية الكلية توجد في مناطق كثيرة من الدماغ، وتُعد من عملياته المعيارية~\cite{Carandini2012}: فالشبكة نفسها تعمل مع المدخل الضعيف ومع القوي.

تحفّظ: الصيغة~\eqref{eq:shunt} تخص عصبونًا لا يطلق نبضات. أما في العصبون المنطلق فيبقى الجهد طوال الوقت قرب العتبة، ويكاد التيار عبر الموصلية المثبطة يكون ثابتًا، فيؤثر التحويل في \emph{تردد} النبضات بالطرح أساسًا~\cite{HoltKoch1997}. ويُقسَم التردد إذا أضيفت إلى العصبون في آن واحد موصليتان خلفيتان مثيرة ومثبطة متوازنتان، كما في النمط المرتعش المتوازن للقشرة الحية~\cite{Chance2002}. فالقسمة لا يحصل عليها الدماغ من التحويل وحده، بل من التثبيط مقرونًا بالضجيج الخلفي~\cite{Carandini2012}.

\begin{figure}[t]
\centering
\begin{tikzpicture}[>=Stealth,x=1.25cm,y=0.05cm]
\draw[->,gray!70!black] (0,0) -- (5.4,0) node[right,black] {\small $g_E/g_L$};
\draw[->,gray!70!black] (0,0) -- (0,68) node[above,black] {\small $V_\infty$، \foreignlanguage{english}{mV}};
\foreach \v in {20,40,60}{\draw[gray!60!black] (-0.05,\v) -- (0.05,\v) node[left=3pt,font=\scriptsize] {\v};}
\foreach \x in {1,2,3,4,5}{\draw[gray!60!black] (\x,-1.5) -- (\x,1.5) node[below=5pt,font=\scriptsize] {\x};}
\draw[very thick,blue!60!black] plot[domain=0:5,samples=60] (\x,{70*\x/(1+\x)});
\draw[very thick,red!70!black] plot[domain=0:5,samples=60] (\x,{70*\x/(3+\x)});
\draw[thick,densely dashed,gray!50!black] plot[domain=0.56:5,samples=60] (\x,{70*\x/(1+\x)-25});
\node[right,font=\small,blue!60!black] at (5.02,58.3) {بلا تثبيط};
\node[right,font=\small,red!70!black] at (5.02,45.5) {يقسم: $g_I=2g_L$};
\node[right,font=\small,gray!40!black] at (5.02,32.5) {يطرح $25\mV$};
\end{tikzpicture}
\caption{التثبيط التحويلي يقسم الجهد ولا يطرح منه. المنحنى الأزرق هو الجهد المستقر~\eqref{eq:shunt} بلا تثبيط، والأحمر مع موصلية مثبطة $g_I=2g_L$. الأحمر هو الأزرق نفسه ممطوطًا أفقيًا ثلاث مرات: كأن المدخل قُسم على $1+g_I/g_L=3$. والمتقطع تثبيط يطرح $25\mV$ ثابتة: ينخفض المنحنى كله ولا يتغير الميل.}
\label{fig:shunt}
\end{figure}

وثمة نتيجة ثانية. ثابت زمن الغشاء $\tau_{\text{eff}}=C/(g_L+g_E+g_I)$، والتثبيط يقصّره. ونافذة التزامن~\eqref{eq:window} متناسبة مع $\tau$، فالتثبيط الذي يصل مع الإثارة \emph{يضيّق النافذة}. والدارة التي يثير فيها المدخلُ العصبونَ مباشرة ويثبطه في الوقت نفسه، بتأخير جزء من الألف من الثانية، عبر وسيط \nt{GABA}، من أكثر الدارات شيوعًا في الدماغ: لا يلحق العصبون أن يستجيب إلا لما وصل في الجزء الأول أو الثاني من الألف من الثانية~\cite{Pouille2001}.

\section{الاختيار}

والآن الأهم مما كان ينقص شبكة \nt{GLUT}: اختيار واحد من عدة. لنأخذ مجموعتين من عصبونات \nt{GLUT} بمدخلين $I_1$ و$I_2$. كل منهما تثبط الأخرى عبر وسطاء \nt{GABA} الخاصين بها بقوة $\beta$ (الشكل~\ref{fig:wta}). وللترددين $r_1,r_2$ نحصل على
\begin{equation}
\tau\,\frac{\dd r_1}{\dd t} = -r_1 + \bigl[I_1-\beta r_2\bigr]_+ ,\qquad
\tau\,\frac{\dd r_2}{\dd t} = -r_2 + \bigl[I_2-\beta r_1\bigr]_+ ,
\label{eq:wta}
\end{equation}
حيث $[u]_+=\max(u,0)$: لا يكون التردد سالبًا.

\begin{figure}[t]
\centering
\begin{tikzpicture}[>=Stealth,
  nrn/.style={circle,draw,very thick,fill=blue!6,minimum size=0.9cm,inner sep=0pt},
  gab/.style={nrn,fill=red!10,minimum size=0.75cm},
  inh/.style={-{Bar[width=7pt]},thick,red!70!black}]
\node[nrn] (E1) at (0,0) {$r_1$};
\node[nrn] (E2) at (4,0) {$r_2$};
\node[gab] (G1) at (1.4,1.1) {};
\node[gab] (G2) at (2.6,-1.1) {};
\draw[->,thick] (E1) -- (G1); \draw[inh] (G1) -- (E2);
\draw[->,thick] (E2) -- (G2); \draw[inh] (G2) -- (E1);
\draw[->,thick,blue!60!black] (-1.6,0) node[left,black] {$I_1$} -- (E1);
\draw[->,thick,blue!60!black] (5.6,0) node[right,black] {$I_2$} -- (E2);
\end{tikzpicture}
\caption{الاختيار من اثنين. مجموعتان من عصبونات \nt{GLUT} (بالأزرق) تثبط إحداهما الأخرى عبر وسطاء \nt{GABA} (بالأحمر).}
\label{fig:wta}
\end{figure}

ما دام العصبونان كلاهما يعملان، فالنظام~\eqref{eq:wta} خطي، ويحدد سلوكه قيمتان ذاتيتان: $-(1+\beta)/\tau$ للانحرافات المتماثلة في الترددين، و$-(1-\beta)/\tau$ للمتعاكسة. عند التثبيط الضعيف، $\beta<1$، القيمتان سالبتان: يعمل العصبونان كلاهما، والتثبيط لا يفعل سوى أن يخفت الأضعف. وعند التثبيط القوي، $\beta>1$، تصبح القيمة الثانية موجبة: أي فرق بين $r_1$ و$r_2$ ينمو حتى يصمت أحد العصبونين تمامًا.

\begin{keyblock}
\begin{proposition}[الفائز يأخذ كل شيء]
\label{prop:wta}
إذا كان التثبيط المتبادل أقوى من الواحد، $\beta>1$، فالحالة التي تعمل فيها المجموعتان كلتاهما غير مستقرة. تصل الشبكة إلى حالة تعمل فيها مجموعة واحدة وتصمت الأخرى. والفائز بالتردد $r_1=I_1$ يحتفظ بفوزه ما دام $I_2<\beta I_1$.
\end{proposition}
\end{keyblock}

اختبرنا ذلك على النموذج. عند $\beta=0.5$ والمدخلين $1.0$ و$0.9$ تعمل المجموعتان كلتاهما بترددين $0.73$ و$0.53$. وعند $\beta=2$ والمدخلين أنفسهما تصمت المجموعة الثانية تمامًا. ولهذا الاختيار ذاكرة: إذا تبادل المدخلان مكانيهما لاحقًا، بقي الفائز السابق فائزًا، لأن $1.0<2\cdot0.9$.

ومع مئة مجموعة يبقى الأمر كذلك: كل منها تثبط البقية، وتبقى واحدة.

\section{تصفير الذاكرة}

ذاكرة الفصل~\ref{sec:glutcomputer} لم تكن تنطفئ إلا بالتعب. لنضف إليها عصبون \nt{GABA} تثيره إشارة ”تصفير“ ويثبط المجموعة. يخفض التثبيط المنحنى $f(wr+I)$ في الشكل~\ref{fig:bistable}، فيختفي التقاطع العلوي، وتنطفئ المجموعة، وتبقى في الحالة السفلى بعد انتهاء التصفير. الآن تُكتب الذاكرة بإشارة وتُمحى بأخرى، كما في قلّاب بمدخلي ضبط وتصفير.

والأجمل أن نصل مجموعتين كهاتين بتثبيط متبادل، كما في الشكل~\ref{fig:wta}، ونعطي كلًّا منهما تغذية راجعة على نفسها. الإشارة على مدخل إحدى المجموعتين تنقل الخلية إلى حالتها، وتتذكر الخلية القرار الأخير. هذا نظير مباشر للماسك المؤلف من عنصري \foreignlanguage{english}{NOR} موصولين تقاطعيًا.

\section{الاستقرار والإيقاع}

تبقى العلة الثالثة لشبكة \nt{GLUT}: الانهيار المتسلسل. لنأخذ مجموعة من عصبونات \nt{GLUT} بتغذية راجعة على نفسها $w_{EE}$، ومجموعة من عصبونات \nt{GABA} تثيرها الأولى بقوة $w_{IE}$ وتثبط الأولى بقوة $w_{EI}$. وللترددين $E$ و$I$ نحصل على
\begin{equation}
\tau_E\frac{\dd E}{\dd t} = -E + w_{EE}E - w_{EI}I + h, \qquad
\tau_I\frac{\dd I}{\dd t} = -I + w_{IE}E ,
\label{eq:ei}
\end{equation}
حيث $h$ المدخل الخارجي~\cite{WilsonCowan1972}. من دون تثبيط، عند $w_{EE}>1$، تنمو الفعالية بلا حد: كل نبضة تولّد أكثر من نبضة لاحقة. ومع التثبيط يكون التردد المستقر
\begin{equation}
E^* \;=\; \frac{h}{1-w_{EE}+w_{EI}w_{IE}}
\label{eq:eistar}
\end{equation}
منتهيًا إذا كان المقام موجبًا. لكن هذا لا يكفي: يجب أيضًا أن يكون الاتزان جاذبًا. ومن التحليل الخطي لـ~\eqref{eq:ei} نحصل على شرطين.

\begin{keyblock}
\begin{proposition}[شروط الاستقرار]
\label{prop:ei}
الاتزان~\eqref{eq:eistar} مستقر إذا كان التثبيط
\begin{enumerate}
\item[(i)] \emph{قويًا بما يكفي}: $w_{EI}\,w_{IE} > w_{EE}-1$؛
\item[(ii)] \emph{سريعًا بما يكفي}: $\tau_I < \tau_E/(w_{EE}-1)$.
\end{enumerate}
إذا انتُهك (i) تصاعدت الفعالية تصاعدًا متسلسلًا. وإذا تحقق (i) وانتُهك (ii) تأخر التثبيط، وبدأت الفعالية تتذبذب.
\end{proposition}
\end{keyblock}

الشرط~(i) هو محدد مصفوفة النظام~\eqref{eq:ei}، والشرط~(ii) أثرها. ومعنى الثاني واضح لكل من ضبط منظِّمًا: التغذية الراجعة المتأخرة لا تخمد الانحراف بل تؤرجحه.

\begin{figure}[t]
\centering
\begin{tikzpicture}[>=Stealth]
\draw[->,gray!70!black] (0,0) -- (10.9,0) node[right,black] {\small $t$، \foreignlanguage{english}{ms}};
\draw[->,gray!70!black] (0,0) -- (0,3.3) node[left,black] {\small $E$};
\foreach \t in {0,50,100,150}{\draw[gray!70!black] (\t*0.07,0.06) -- (\t*0.07,-0.06) node[below,font=\scriptsize] {\t};}
\input{figures/ei_traces}
\node[font=\small,gray!40!black,anchor=west] at (1.3,3.45) {بلا تثبيط: انهيار متسلسل};
\node[font=\small,blue!60!black,anchor=west] at (7.4,1.25) {تثبيط سريع};
\node[font=\small,red!70!black,anchor=west] at (7.4,2.55) {تثبيط بطيء};
\end{tikzpicture}
\caption{تردد مجموعة \nt{GLUT} ذات تغذية راجعة $w_{EE}=1.5$ و$\tau_E=10\ms$، بعد تشغيل المدخل. بلا تثبيط (المتقطع) تنمو الفعالية بلا حد. ومع تثبيط بقوة $w_{EI}w_{IE}=2$ و$\tau_I=2\ms$ (المنحنى الأزرق) تستقر عند المستوى $E^*=h/(1-1.5+2)=0.67h$. ومع التثبيط نفسه لكن بطيئًا، $\tau_I=30\ms$ (المنحنى الأحمر)، تتذبذب الفعالية.}
\label{fig:ei}
\end{figure}

يُعتقد أن القشرة تعمل في هذا النمط بالذات: تغذيتها الراجعة الذاتية قوية بما يكفي لتنجرف إلى انهيار متسلسل لولا التثبيط، ولا يبقيها عند نقطة التشغيل إلا التثبيط السريع~\cite{Tsodyks1997isn}.

ولهذا النمط بصمة متناقضة ظاهريًا يمكن التعرف عليه بها من الخارج~\cite{Tsodyks1997isn}. لنضف إلى~\eqref{eq:ei} مدخلًا خارجيًا $h_I$ مباشرة إلى مجموعة \nt{GABA}. تصبح الترددات المستقرة
\begin{equation}
E^* \;=\; \frac{h-w_{EI}\,h_I}{D},\qquad
I^* \;=\; \frac{w_{IE}\,h+(1-w_{EE})\,h_I}{D},\qquad
D \;=\; 1-w_{EE}+w_{EI}w_{IE} .
\label{eq:isn}
\end{equation}
إذا كان $w_{EE}>1$ فمعامل $h_I$ في الصيغة الثانية سالب: ادفع عصبونات \nt{GABA} فـ\emph{يهبط} ترددها. والسبب في الحلقة. التثبيط الزائد يخفت مجموعة \nt{GLUT}، فتقل إثارتها لعصبونات \nt{GABA}، فتخسر هذه أكثر مما كسبت. وبأرقام الشكل~\ref{fig:ei} ($w_{EE}=1.5$، $w_{EI}w_{IE}=2$، $D=1.5$) كل وحدة من المدخل المضاف تخفض $I^*$ بثلث وحدة. وقد وُجدت هذه البصمة في القشرة: حين يُكبت استجابة عصبونات القشرة البصرية منبهٌ محيط، فإن التيار المثبط الذي تتلقاه لا يزداد بل يهبط مع المثير~\cite{Ozeki2009}. ولاحقًا وُجدت البصمة نفسها في مناطق وطبقات مختلفة من القشرة~\cite{Sanzeni2020}.

والتذبذبات عند انتهاك الشرط~(ii) موجودة في الدماغ أيضًا. الحلقة ”\nt{GLUT} يثير \nt{GABA}، و\nt{GABA} يثبط \nt{GLUT}“ بتأخير بضعة أجزاء من الألف من الثانية تولّد إيقاعًا دوره من رتبة $12\text{--}50\ms$ (التردد $20\text{--}80\,\text{Hz}$)، وهذه الإيقاعات السريعة في القشرة معروفة جيدًا~\cite{Cardin2009,Buzsaki2012}. هذه ليست ساعة الحاسوب: الإيقاع محلي ولا ينشأ إلا حين تكون الشبكة نشطة. لكنه على مدى بضع دورات يقسم الزمن إلى نوافذ يمكن للعصبونات أن تنطلق فيها.

\section{الخلاصة}

\begin{table}[t]
\centering
\small
\begin{tabular}{>{\raggedright}p{3.3cm}>{\raggedright}p{4.4cm}>{\raggedright\arraybackslash}p{4.4cm}}
\toprule
& \nt{GLUT} فقط & \nt{GLUT} + \nt{GABA} \\
\midrule
\foreignlanguage{english}{NOT} & عبر مستشعرات ”التشغيل/الإطفاء“ فقط & الحظر $a\wedge\bar b$؛ و\foreignlanguage{english}{NOT} مع فعالية خلفية \\
الأسلاك لكل بت & اثنان & واحد \\
اتجاه الحركة & لا & حظر مع تأخير \\
ضبط الكسب & لا & القسمة: التحويل مع الضجيج الخلفي \\
نافذة التزامن & $\sim\tau$ & يضيّقها التثبيط \\
اختيار واحد من كثيرين & لا & تثبيط متبادل، $\beta>1$ \\
تصفير الذاكرة & بالتعب فقط & بإشارة عبر \nt{GABA} \\
الاستقرار & بالتعب فقط & تغذية راجعة سالبة سريعة \\
الإيقاع & غير لازم & ينشأ مع التثبيط البطيء \\
\bottomrule
\end{tabular}
\caption{ما يضيفه \nt{GABA} إلى شبكة من عصبونات \nt{GLUT}.}
\label{tab:gaba}
\end{table}

لا يكاد \nt{GABA} يضيف إلى الشبكة \emph{حسابات} جديدة، فكل ذلك كان يمكن حسابه بالمستشعرات ثنائية السلك، بل يضيف \emph{التحكم} (الجدول~\ref{tab:gaba}): الاختيار، والتصفير، وضبط الكسب، والاستقرار. الإثارة في الدماغ تحمل البيانات، والتثبيط يقرر أي البيانات تمر، ومتى، وبأي علوّ. ولعصبون من كل أربعة أو خمسة في القشرة هذا تقسيم معقول جدًا للعمل.

\section{تمارين}

\begin{exercise}[\lvlA]
\label{ex:03:1}
\emph{التحويل بالأرقام.} $E_E=70\mV$. بحسب~\eqref{eq:shunt} أوجد $V_\infty$ عند $g_E=g_L$ و$4g_L$، بلا تثبيط ومع تحويل $g_I=2g_L$. أي $V_\infty$ عند $g_E=4g_L$ يعطيه طرحٌ يُنقص عند $g_E=g_L$ قدر ما ينقصه التحويل؟ كم مرة يضيّق التحويل نافذة التزامن؟
\end{exercise}

\begin{exercise}[\lvlB]
\label{ex:03:2}
\emph{استنتاج شروط الاستقرار.} أوجد أثر مصفوفة النظام الخطي~\eqref{eq:ei} ومحددها، واستنتج منهما شرطي القضية~\ref{prop:ei}. على حدّ الشرط~(ii) يفقد الاتزان استقراره تذبذبيًا. أوجد دور هذه التذبذبات بأرقام الشكل~\ref{fig:ei}.
\end{exercise}

\begin{exercise}[\lvlB]
\label{ex:03:3}
\emph{إيقاع من التأخير.} لنستبدل بالتثبيط البطيء في النموذج~\eqref{eq:ei} تثبيطًا سريعًا بتأخير $d$: $\tau_E\dot E=-E(t)-GE(t-d)+h$. عند $\tau_E=10\ms$ أوجد أصغر تأخير $d_c$ يعطي تذبذبات، ودورها عند $G=2$ و$G=5$. هل تقع ضمن إيقاعات القشرة السريعة، $12\text{--}50\ms$، خلافًا للتمرين~\ref{ex:03:2}؟
\end{exercise}

\begin{exercise}[\lvlB]
\label{ex:03:4}
\emph{نمذجة: اختيار بذاكرة.} كامل~\eqref{eq:wta} عند $\beta=2$ و$\tau=1$. (أ)~عند $I_1=1$ ارفع $I_2$ ببطء من $0$ إلى $3$ ثم أعده: عند أي قيم $I_2$ تنقلب الشبكة؟ (ب)~بكم يزداد زمن القرار انطلاقًا من الصمت حين يصغر فرق المدخلين $\varepsilon$ عشر مرات؟
\end{exercise}

\begin{exercise}[\lvlB]
\label{ex:03:5}
\emph{تصميم: كاشف لنطاق من السرعات.} يجب أن يصمت كاشف الشكل~\ref{fig:direction} عند الحركة من $A$ إلى $B$ بزمن عبور $5\text{--}20\ms$. وسيط \nt{GABA} بتأخير $d_k$ يحظر $Y$ على المجال $[d_k,\,d_k+5\ms]$ بعد نبضة $A$؛ والإثارة من $B$ تصل فورًا. كم وسيطًا يلزم، وبأي تأخيرات؟
\end{exercise}

\begin{exercise}[\lvlB]
\label{ex:03:6}
\emph{الحظر والخلفية.} كم عصبونًا تتطلب الزوجية $x\oplus y\oplus z$ في المخطط العام ”وسيط \nt{GABA} لكل مدخل، وحظر لكل توليفة فيها $f=1$، و\foreignlanguage{english}{OR} مشتركة“؟ ابنها من عصبونين، \nt{GABA} و\nt{GLUT}، يتلقيان المداخل مباشرة، مع بيان الأوزان والعتبات.
\end{exercise}

\begin{exercise}[\lvlC]
\label{ex:03:7}
\emph{تصميم: الاختيار من مئة.} يجب أن تثبط كل واحدة من $100$ مجموعة \nt{GLUT} كل المجموعات الأخرى بالتساوي، لا نفسها. وسطاء \nt{GABA} خطيون تثيرهم المجموعات ويثبطونها؛ ويمكن للمجموعات أن يثير بعضها بعضًا، لا نفسها. ما أصغر عدد من الوسطاء؟ وماذا لو لم تكن ثمة وصلات إثارة مباشرة أصلًا؟
\end{exercise}

\begin{exercise}[\lvlC]
\label{ex:03:8}
\emph{بحث: متى يكون التثبيط متناقضًا.} في نموذج الشكل~\ref{fig:ei} ($w_{EI}=2$، $w_{IE}=1$) صارت دالة نقل مجموعة \nt{GLUT} هي $f(u)=[u]_+^2$، وبقيت مجموعة \nt{GABA} خطية. فوق أي مدخل $h_c$ تخفض الإضافة إلى مجموعة \nt{GABA} ترددها هي نفسها؟ تحقق بالنمذجة.
\end{exercise}

\chapter{شفرة مورس: ما نعرفه عن اللغة التي تتكلم بها عصبونات \nt{GLUT} و\nt{GABA}}
\chaptermark{شفرة مورس}
\label{sec:code}

\section{أبجدية من رمز واحد}

في شفرة مورس رمزان، النقطة والشَّرطة، وفواصل بثلاثة أطوال بينهما. أما أبجدية العصبون، كما رأينا في الفصل~\ref{sec:neurontypes}، فأفقر: رمز واحد. تدوم النبضة نحو جزء من الألف من الثانية، وكل نبضات العصبون الواحد متساوية في الارتفاع والشكل. لذلك فالرسالة كلها في الفواصل، أي في متتالية اللحظات $t_1,t_2,t_3,\dots$ إنها شفرة بلا شُّرَط، لا يحمل المعنى فيها إلا إيقاع النقاط (الشكل~\ref{fig:morse}).

\begin{figure}[t]
\centering
\begin{tikzpicture}[>=Stealth,x=0.08cm,y=1cm]
% Morse letter: dot, dash, dot
\node[left,font=\small] at (-6,2.55) {مورس:};
\fill[black] (0,2.45) rectangle (6,2.65);
\fill[black] (14,2.45) rectangle (32,2.65);
\fill[black] (40,2.45) rectangle (46,2.65);
\node[right,font=\small,gray!40!black] at (52,2.55) {نقطة، شَرطة، نقطة: المعنى في المُدد والفواصل};
% stimulus onset
\node[left,font=\small] at (-6,0.4) {العصبون:};
\draw[->,thick,green!45!black] (0,-0.55) -- (0,-0.05);
\node[below,font=\scriptsize,green!45!black] at (0,-0.55) {المنبه};
\draw[gray!70!black] (0,0) -- (152,0) node[right,black] {\small $t$};
% spikes: single ones blue, the burst red
\foreach \s in {14,26,41,88,112,131}{\draw[very thick,blue!60!black] (\s,0) -- (\s,0.8);}
\foreach \s in {60,63,66,69}{\draw[very thick,red!70!black] (\s,0) -- (\s,0.8);}
% latency
\draw[<->,thick] (0,1.1) -- node[above,font=\scriptsize] {$t_1$: زمن النبضة الأولى} (14,1.1);
% burst
\draw[decorate,decoration={brace,amplitude=4pt},red!70!black] (58,0.95) -- (71,0.95) node[midway,above=4pt,font=\scriptsize] {الرشقة هي ”الشَّرطة“};
% counting windows
\foreach \a/\b/\n in {0/50/3,50/100/5,100/150/2}{
  \draw[decorate,decoration={brace,amplitude=4pt,mirror},gray!50!black] (\a+1,-0.2) -- (\b-1,-0.2) node[midway,below=4pt,font=\small] {$n=\n$};}
\node[font=\scriptsize,gray!40!black] at (75,-1.05) {ترميز التردد: عدد النبضات $n$ في كل نافذة};
\foreach \x in {0,50,100,150}{\draw[gray!60!black] (\x,-0.06) -- (\x,0.06);}
\end{tikzpicture}
\caption{شفرة مورس وشفرة العصبون. يمكن قراءة متتالية النبضات نفسها بطرق مختلفة: عدّ النبضات في نوافذ من $50\ms$ (القسم~\ref{sec:rate})، أو قياس التأخير $t_1$~\eqref{eq:latency}، أو ملاحظة الرشقة، أي ”الشَّرطة“~\eqref{eq:burst}.}
\label{fig:morse}
\end{figure}

ما الفواصل الممكنة أصلًا؟ من الأسفل يحدها عدم الاستجابة: بعد النبضة لا يستطيع العصبون أن ينطلق ثانية مدة جزء من الألف من الثانية تقريبًا، فلا يتجاوز العدد بضع مئات من النبضات في الثانية. ومن الأعلى لا يحدها شيء: قد يصمت العصبون دقائق. ترددات عصبونات \nt{GLUT} في القشرة تتراوح بين أجزاء من النبضة وبضع عشرات من النبضات في الثانية، ومعظم العصبونات تعمل معظم الوقت قرب الحد الأدنى.

في الفصلين~\ref{sec:glutcomputer} و\ref{sec:gaba} اعتمدنا عرضًا حينًا طريقة وحينًا أخرى لكتابة الأعداد بالنبضات. والآن نسأل مباشرة: بأي لغة تتخاطب عصبونات \nt{GLUT} و\nt{GABA}؟ لا جواب كاملًا؛ فهذه اللغة لم تُفكّ رموزها. لكن بعض ما يُعرف عنها ثابت، ويمكن الوصول إلى معظمه بتفكير مهندس الاتصالات: سعة القناة، والضجيج، والمستقبِل، وكلفة الإرسال.

\section{كم يمكن أن ننقل}

لنبدأ بالحد الأعلى. ليكن المستقبِل يميز لحظات النبضات بدقة $\Delta t$: الإزاحات الأصغر من $\Delta t$ لا يميزها. لنقطّع الزمن إلى خانات طولها $\Delta t$، أصغر من زمن عدم الاستجابة، بحيث تحوي كل خانة نبضة واحدة أو لا شيء (الشكل~\ref{fig:cells}). عند تردد متوسط $r$ تقع النبضة في الخانة باحتمال $p=r\Delta t$. يحمل التدفق أكبر قدر من المعلومات حين تكون الخانات مستقلة، وعندئذ تعطي كل خانة إنتروبيا $-p\log_2p-(1-p)\log_2(1-p)\approx p\log_2(e/p)$ بت عند $p\ll1$. وبالقسمة على $\Delta t$ نحصل على السرعة القصوى للنقل وحصة النبضة الواحدة منها~\cite{MacKay1952}:
\begin{empheq}[box=\keybox]{equation}
R_{\max} \;\approx\; r\,\log_2\frac{e}{r\,\Delta t}\ \ \text{bit/s},
\qquad
\frac{R_{\max}}{r} \;\approx\; \log_2\frac{e}{r\,\Delta t}\ \ \text{bit/spike}.
\label{eq:spikeentropy}
\end{empheq}
عند $r=10$ نبضات في الثانية و$\Delta t=1\ms$ هذا نحو ثمانية بتات للنبضة وثمانين بتًا في الثانية.

\begin{figure}[t]
\centering
\begin{tikzpicture}[>=Stealth,x=0.5cm,y=1cm]
% time cut into cells of width Delta t; a cell holds one impulse or none
\foreach \k in {0,...,23}{\draw[gray!55] (\k,0) rectangle (\k+1,0.7);}
\foreach \k in {2,7,8,15,21}{
  \fill[red!10] (\k,0) rectangle (\k+1,0.7);
  \draw[very thick,red!70!black] (\k+0.5,0.05) -- (\k+0.5,0.65);}
\foreach \k in {0,...,23}{
  \pgfmathtruncatemacro{\bit}{(\k==2)||(\k==7)||(\k==8)||(\k==15)||(\k==21)}
  \ifnum\bit=1 \def\bitcol{red!70!black}\else \def\bitcol{gray!60!black}\fi
  \node[font=\scriptsize,text=\bitcol] at (\k+0.5,-0.25) {\bit};}
\draw[<->,gray!60!black] (0,0.95) -- (1,0.95) node[midway,above,font=\scriptsize,black] {$\Delta t$};
\node[font=\small,anchor=east] at (-0.3,0.35) {النبضات};
\node[font=\small,anchor=east] at (-0.3,-0.25) {البتات};
\draw[decorate,decoration={brace,amplitude=5pt,mirror},gray!60!black] (0,-0.55) -- (24,-0.55)
  node[midway,below=6pt,font=\scriptsize,black,align=center] {المستقبِل الذي يعدّ فقط: ”$n=5$“؛\\ والمستقبِل الذي يرى الخانات: أي $5$ من $24$ بالضبط، أي $\log_2\binom{24}{5}\approx15$ بتًا};
\end{tikzpicture}
\caption{من أين تأتي السعة~\eqref{eq:spikeentropy}. قُطّع الزمن إلى خانات طولها $\Delta t$، وفي كل خانة إما نبضة ($1$) أو لا ($0$): تدفق النبضات سلسلة ثنائية. تقع النبضة في الخانة باحتمال $p=r\Delta t$. كلما صغرت الخانات طالت السلسلة وكثرت البتات؛ وعدّاد النبضات يفقد تقريبًا كل ما هو مكتوب في \emph{مواضع} الآحاد.}
\label{fig:cells}
\end{figure}

البتات لكل نبضة لا تتعلق بدقة الزمن إلا لوغاريتميًا: بدقة $10\ms$ يبقى نحو خمسة بتات للنبضة. لكن إذا كان المستقبِل يعدّ النبضات في الثانية فقط، فعند $r=10$ لن يحصل على أكثر من بتّين في الثانية كلها (القسم~\ref{sec:rate}).

\begin{keyblock}
\begin{proposition}[سعة القناة النبضية]
\label{prop:capacity}
تدفق النبضات بتردد متوسط $r$، مقروءًا بدقة زمنية $\Delta t$، لا يحمل أكثر من~\eqref{eq:spikeentropy}. وتكاد السعة كلها تقبع في \emph{توقيت} النبضات: المستقبِل الذي يعدّ النبضات فقط يستخرج من التدفق نفسه أقل بعشرات المرات مما يستخرجه من يميز لحظاتها بدقة جزء من الألف من الثانية.
\end{proposition}
\end{keyblock}

هذا حد أعلى لا قياس. القياسات على العصبونات الحسية، التي يُعرف منبهها، تعطي من بت واحد إلى بضعة بتات للنبضة؛ وفي أفضل الحالات يبلغ هذا نصف الحد المحسوب لدقتها~\cite{Rieke1997}. إذن، على الأقل عند مدخل الجهاز العصبي، يُستعمل توقيت النبضات ولا يضيع.

\section{القناة الضجيجية}

حُسبت السعة~\eqref{eq:spikeentropy} بلا ضجيج. أما أين ينشأ الضجيج، فثلاث حقائق ثابتة.

\emph{مولّد النبضات نفسه دقيق.} إذا حُقن في عصبون قشري، مأخوذ من الدماغ مع قطعة من النسيج، التيارُ نفسه المتغير بسرعة في الزمن مرات كثيرة، تكررت النبضات بدقة نحو جزء من الألف من الثانية~\cite{Mainen1995}. وإذا كان التيار ثابتًا، تتبعثر لحظات النبضات من مرة إلى أخرى: فالتوقيت الدقيق تولّده التغيرات السريعة في المدخل، لا العصبون نفسه.

\emph{المشبك غير موثوق.} النبضة التي تبلغ مشبك \nt{GLUT} تقذف دفعة من الناقل العصبي باحتمال ما $p$ فقط. في مشابك الحُصين لا يصل أكثر من نصف النبضات إلى المستقبِل أصلًا، والسبب هو عشوائية القذف بالذات، لا أعطال التوصيل~\cite{AllenStevens1994}. وبالنسبة لمهندس الاتصالات هذه قناة محو؛ وعدة مشابك بين عصبونين تنقذ الموقف جزئيًا.

\emph{في القشرة الحية يبدو تدفق النبضات عشوائيًا.} الفترات بين النبضات مبعثرة تقريبًا كما في تدفق بواسون العشوائي، وعند تكرار المنبه نفسه يتغير عدد النبضات في النافذة بحيث يقارب تباينه متوسطه~\cite{Softky1993}.

كيف يولّد عنصر دقيق تدفقًا عشوائيًا؟ يتلقى العصبون آلاف المداخل، كل منها غير موثوق، والإثارة والتثبيط متوازنان كما في الفصل~\ref{sec:gaba}. يرتعش جهد الغشاء قرب العتبة، وهذه الارتعاشات هي التي تحدد لحظات النبضات. أهو ضجيج أم رسالة لا نحسن قراءتها؟ هذا إلى حد بعيد سؤال مفتوح. وقد تبيّن أن جزءًا من ”الضجيج“ رسالة فعلًا: ففي القشرة البصرية تتبع حصة كبيرة من التشتت حركات الحيوان نفسه~\cite{Stringer2019}. والصعوبة هنا مبدئية: الرسالة المضغوطة جيدًا لا تتميز عن تدفق عشوائي عند من لا يعرف الشيفرة.

\section{التردد: بطيء لكنه موثوق}
\label{sec:rate}

أبسط الشيفرات \emph{ترميز التردد}: العدد مكتوب في عدد النبضات التي يطلقها العصبون في النافذة $T$. لا المحو ولا ارتعاش اللحظات يضر بهذه الشيفرة. لكن لها ثمنًا. إذا كان التدفق بواسونيًا فعدد النبضات $n$ في النافذة متوسطه $rT$ وتشتته $\sqrt{rT}$:
\begin{empheq}[box=\keybox]{equation}
\frac{\sigma_n}{\bar n} \;=\; \frac{1}{\sqrt{r\,T}} .
\label{eq:rateerror}
\end{empheq}
لمعرفة التردد بدقة عشرة في المئة تلزم مئة نبضة، أي خمس ثوانٍ عند عشرين نبضة في الثانية. والمستويات المتجاورة تتمايز إذا تباعدت بنحو تشتتين، لذا يتمايز حتى التردد $r$ في الزمن $T$ نحو $\sqrt{rT}$ مستوى: عند $r=10$ و$T=1\,\text{s}$ هذا ثلاثة أو أربعة مستويات، أقل من بتين.
والدماغ لا يعمل بهذا البطء. يتعرف الإنسان على حيوان في صورة تُعرض عليه لحظة، وتظهر الاستجابة في الدماغ بعد نحو $150\ms$~\cite{Thorpe1996}. وبتقدير خشن تمر الإشارة في هذا الزمن بنحو عشر مراحل معالجة متتالية، بمعدل $10\text{--}20\ms$ للمرحلة. وبالترددات المعتادة في القشرة لا يلحق كل عصبون أن يطلق في المرحلة إلا صفرًا أو نبضة واحدة، فتردده في نافذة كهذه غير معرّف. وثمة مخرجان: أخذ عصبونات كثيرة، أو النظر إلى الزمن.

\emph{عصبونات كثيرة.} لتكن $N$ عصبونات تحمل العدد نفسه، والمستقبِل يجمع نبضاتها. إذا كانت ضجيجاتها مستقلة حلّ في~\eqref{eq:rateerror} $NrT$ محل $rT$: مئة عصبون عند $r=20$ في $15\ms$ تعطي ثلاثين نبضة ودقة نحو عشرين في المئة. المكان يحل محل الزمن. غير أن ضجيجات العصبونات المتجاورة في القشرة ليست مستقلة: معامل الارتباط بين أعداد النبضات لزوج من الجيران نحو $0.1$ عادة~\cite{Zohary1994}. فإذا كان يساوي $c$، فتباين المتوسط على $N$ عصبونًا يساوي
\begin{empheq}[box=\keybox]{equation}
\sigma^2_N \;=\; \sigma^2_1\,\frac{1+(N-1)\,c}{N}
\;\xrightarrow[N\to\infty]{}\; c\,\sigma^2_1 .
\label{eq:correlated}
\end{empheq}

\begin{keyblock}
\begin{proposition}[حد التوسيط]
\label{prop:pooling}
مع الضجيج المترابط لا يقلل التوسيط على مجموعة الخطأ بأكثر من $1/\sqrt{c}$ مرة. عند $c=0.1$ هذا ثلاث مرات، ويُبلغ هذا الربح ببضع عشرات من العصبونات؛ وإضافة المزيد لا تجدي.
\end{proposition}
\end{keyblock}

ليس كل ارتباط ضارًا بالقدر نفسه. لا يحد الدقة إلا ذلك الجزء من الضجيج المشترك الذي \emph{يشبه الإشارة نفسها}، أي يزيح المجموعة كلها كما كان سيزيحها تغيّر الكمية المقيسة~\cite{MorenoBote2014}. وكم من هذا الضجيج في الدماغ فعلًا، ما زال قيد البحث.

وثمة طريقة أخرى لاستعمال عصبونات كثيرة: كتابة العدد في \emph{أي} عصبون انطلق، كما في محدد اتجاه الصوت (الشكل~\ref{fig:itd}). هذا ترميز \emph{بالموضع}، وهو أوثق المعروف: في العصبونات الحسية يدل رقم السلك على أي نقطة من الجلد لُمست أو على طبقة الصوت، وهذا ثابت مرارًا. وهو مكلف في عدد العصبونات، لكنه سريع: تستغرق الرسالة تأخيرًا واحدًا.

\section{الزمن: سريع، لكن من أين نبدأ العدّ}

المخرج الثاني كتابة العدد في زمن النبضة. لنأخذ العصبون~\eqref{eq:glutneuron} ونقدّم إليه من اللحظة $t=0$ مدخلًا ثابتًا كان سيرفع الجهد إلى المستوى $U$. ينمو الجهد كـ$V=U\bigl(1-e^{-t/\tau}\bigr)$ ويبلغ العتبة $\theta$ في اللحظة
\begin{empheq}[box=\keybox]{equation}
t_1 \;=\; \tau\,\ln\frac{U}{U-\theta}, \qquad U>\theta .
\label{eq:latency}
\end{empheq}
المدخل القوي نبضة مبكرة، والضعيف نبضة متأخرة، وما دون العتبة لا نبضة. عند $\tau=10\ms$ يعطي المدخل $U=4\theta$ النبضة الأولى بعد $2.9\ms$، و$U=2\theta$ بعد $6.9\ms$، و$U=1.2\theta$ بعد $17.9\ms$. نبضة واحدة وحيدة تحمل عددًا.

لكن من أي لحظة نعدّ $t_1$؟ لا ساعة، والمستقبِل لا يعرف متى بدأ المنبه. ثمة ثلاثة أجوبة.

\emph{التأخير النسبي.} عصبونان يريان المنبه نفسه، وفرق تأخيريهما $t_1^A-t_1^B$ لا يتوقف إلا على مدخليهما، لا على لحظة البدء. وتقارن اللحظاتِ كواشفُ التزامن ذات التأخيرات (الشكل~\ref{fig:itd}). إذا أطلق $N$ عصبونًا نبضة واحدة لكل منها، فإن \emph{ترتيب} وصولها وحده يحمل حتى $\log_2N!$ بت: لعشرة عصبونات نحو اثنين وعشرين بتًا، بينما الجواب ”انطلق أم لا“ لا يعطي أكثر من عشرة. وقد بُيّن في الشبكية أن الصورة يمكن استعادتها بسرعة وموثوقية من التأخيرات النسبية للنبضات الأولى لعصبوناتها المخرجة~\cite{Gollisch2008}.

\emph{الرشقة الشبكية بداية للكلمة.} التغير الحاد في المنبه يجعل عددًا كبيرًا من العصبونات ينطلق في آن واحد تقريبًا. هذه الرشقة الشبكية نفسها علامة على البداية، كالفاصل الطويل بين الكلمات في شفرة مورس، ويعدّ المستقبِل التأخيرات منها.

\emph{الإيقاع المحلي.} في الفصل~\ref{sec:gaba} ولّدت حلقة \foreignlanguage{english}{\nt{GLUT}--\nt{GABA}} إيقاعًا محليًا دوره $12\text{--}50\ms$. ليس ساعة، لكن داخل دورته يلحق العصبون ذو المدخل القوي أن ينطلق قبل العصبون ذي المدخل الضعيف، فيتحول~\eqref{eq:latency} إلى ترميز \emph{بالطور}: العدد مكتوب في الجزء من الدورة الذي وصلت فيه النبضة. وقد رُصد هذا الترميز: العصبونات المسؤولة عن موضع معين في المكان تنطلق أبكر فأبكر في دورة إيقاع محلي بطيء كلما عبر الحيوان موضعها ”الخاص“~\cite{OKeefe1993}. أما إلى أي حد يستعمل الدماغ الطور، فغير معروف بعد.

\section{المستقبِل يختار الشيفرة}

في متتالية النبضات تردد، وأزمنة، وترتيب. أيها يُقرأ يقرره المستقبِل، ويقرره بوسيط واحد: ثابته الزمني $\tau$.

لنأخذ متوسط المعادلة~\eqref{eq:glutneuron} على الزمن. إذا وردت على العصبون تدفقات مستقلة بترددات $r_j$، فالجهد المتوسط وارتعاشه النسبي يساويان
\begin{empheq}[box=\keybox]{equation}
\bar V \;=\; \tau\sum_j w_j\,r_j ,
\qquad
\frac{\sigma_V}{\bar V} \;\approx\; \frac{1}{\sqrt{2\,r_{\text{in}}\,\tau}} ,
\label{eq:reader}
\end{empheq}
حيث المساواة الثانية مكتوبة لأوزان متساوية، و$r_{\text{in}}$ مجموع ترددات كل المداخل. حين يكون $\tau$ كبيرًا، لا يكاد الجهد يرتعش ويتبع المجموع الموزون للترددات: المستقبِل \emph{مكامل}، يقرأ الترددات وينسى الأزمنة. وحين يكون $\tau$ صغيرًا، يلحق الجهد أن يهبط بين النبضات، ولا تُبلغ العتبة إلا إذا وصلت عدة نبضات ضمن النافذة~\eqref{eq:window}: المستقبِل \emph{كاشف تزامن}، يقرأ الأزمنة (الشكل~\ref{fig:reader}). ألف مدخل بخمس نبضات في الثانية لكل منها عند $\tau=10\ms$ تعطي ارتعاشًا نحو عشرة في المئة، أي تكاملًا شبه خالص. وإذا قصّر التثبيط، كما في الفصل~\ref{sec:gaba}، $\tau$ إلى $2\ms$، نما الارتعاش إلى أكثر من عشرين في المئة، وبدأ العصبون يلاحظ التزامنات.

\begin{figure}[t]
\centering
\begin{tikzpicture}[>=Stealth,x=0.115cm,y=1cm]
% common input: the same spike train for both receivers
\foreach \s in {6,21,22,38,52,55.5,59,62.5,66,69.5,73,76.5,80,83.5,87,90.5,94,97.5}{\draw[thick,gray!40!black] (\s,5.25) -- (\s,5.65);}
\draw[gray!70!black] (0,5.25) -- (105,5.25);
\node[left,font=\small] at (-1,5.45) {المدخل};
\node[above,font=\scriptsize,gray!40!black] at (13.5,5.65) {نادرًا};
\node[above,font=\scriptsize,gray!40!black] at (21.5,6.0) {زوج};
\node[above,font=\scriptsize,gray!40!black] at (75,5.65) {كثيرًا};
% axes and thresholds
\foreach \y/\th in {2.7/2.5,0/1.5}{
  \draw[gray!70!black] (0,\y) -- (106,\y) node[right,black] {\small $t$};
  \draw[densely dashed,gray!60!black] (0,\y+0.6*\th) -- (104,\y+0.6*\th) node[right,black,font=\small] {$\theta$};
}
\node[anchor=south west,font=\small] at (0,4.3) {$\tau=10\ms$: مكامل};
\node[anchor=south east,font=\small] at (104,1.0) {$\tau=2\ms$: كاشف تزامن};
% integrator: tau=10 ms, theta=2.5w
\begin{scope}[yshift=2.7cm]
\foreach \a/\b/\v in {0/6/0.000,6/21/1.000,21/22/1.223,22/38/2.107,38/52/1.425,52/55.5/1.351,55.5/59/1.952,59/62.5/2.376,62.5/66/0.000,66/69.5/1.000,69.5/73/1.705,73/76.5/2.201,76.5/80/0.000,80/83.5/1.000,83.5/87/1.705,87/90.5/2.201,90.5/94/0.000,94/97.5/1.000,97.5/104/1.705}{\draw[very thick,blue!60!black] plot[domain=\a:\b,samples=14] (\x,{0.6*\v*exp(-(\x-\a)/10)});}
\foreach \s/\p/\q in {6/0.000/1.000,21/0.223/1.223,22/1.107/2.107,38/0.425/1.425,52/0.351/1.351,55.5/0.952/1.952,59/1.376/2.376,62.5/1.674/2.500,62.5/2.500/0.000,66/0.000/1.000,69.5/0.705/1.705,73/1.201/2.201,76.5/1.551/2.500,76.5/2.500/0.000,80/0.000/1.000,83.5/0.705/1.705,87/1.201/2.201,90.5/1.551/2.500,90.5/2.500/0.000,94/0.000/1.000,97.5/0.705/1.705}{\draw[very thick,blue!60!black] (\s,0.6*\p) -- (\s,0.6*\q);}
\foreach \s in {62.5,76.5,90.5}{\draw[->,thick,red!70!black] (\s,1.58) -- (\s,1.95);}
\end{scope}
% coincidence: tau=2 ms, theta=1.5w
\begin{scope}[yshift=0cm]
\foreach \a/\b/\v in {0/6/0.000,6/21/1.000,21/22/1.001,22/38/0.000,38/52/1.000,52/55.5/1.001,55.5/59/1.174,59/62.5/1.204,62.5/66/1.209,66/69.5/1.210,69.5/73/1.210,73/76.5/1.210,76.5/80/1.210,80/83.5/1.210,83.5/87/1.210,87/90.5/1.210,90.5/94/1.210,94/97.5/1.210,97.5/104/1.210}{\draw[very thick,blue!60!black] plot[domain=\a:\b,samples=14] (\x,{0.6*\v*exp(-(\x-\a)/2)});}
\foreach \s/\p/\q in {6/0.000/1.000,21/0.001/1.001,22/0.607/1.500,22/1.500/0.000,38/0.000/1.000,52/0.001/1.001,55.5/0.174/1.174,59/0.204/1.204,62.5/0.209/1.209,66/0.210/1.210,69.5/0.210/1.210,73/0.210/1.210,76.5/0.210/1.210,80/0.210/1.210,83.5/0.210/1.210,87/0.210/1.210,90.5/0.210/1.210,94/0.210/1.210,97.5/0.210/1.210}{\draw[very thick,blue!60!black] (\s,0.6*\p) -- (\s,0.6*\q);}
\foreach \s in {22}{\draw[->,thick,red!70!black] (\s,0.98) -- (\s,1.35);}
\end{scope}
\foreach \x in {0,50,100}{\draw[gray!60!black] (\x,-0.06) -- (\x,0.06) node[below,font=\scriptsize] {$\x\,\text{ms}$};}
\end{tikzpicture}
\caption{المدخل نفسه، ومستقبِلان مختلفان. كل نبضة ترفع الجهد بمقدار $w$، ثم ينساب نازلًا كما في نموذج العصبون~\eqref{eq:glutneuron}؛ والأسهم الحمراء نبضات المستقبِل. المستقبِل البطيء ($\tau=10\ms$، $\theta=2.5w$) ينطلق حيث النبضات \emph{كثيرة}، ولا يلاحظ الزوج. والسريع ($\tau=2\ms$، $\theta=1.5w$) لا ينطلق إلا للزوج الواقع ضمن النافذة~\eqref{eq:window}، ويُفلت التدفق المتواتر غير المتزامن.}
\label{fig:reader}
\end{figure}

تُظهر القياسات أن موصلية الغشاء في القشرة النشطة تزداد عدة مرات مقارنة بالراحة~\cite{Destexhe2003}. إذن $\tau$ هناك أقصر، وعصبونات القشرة في نمط العمل أقرب إلى كواشف التزامن منها إلى المكاملات.

\begin{keyblock}
\begin{proposition}[المستقبِل يحدد الشيفرة]
\label{prop:reader}
متتالية النبضات نفسها تحمل للمستقبِل البطيء ترددًا، وللسريع أزمنة، وللحجم الذي قُذف فيه الناقل العصبي مستوى تركيز منعّمًا على مدى ثوانٍ (الفصل~\ref{sec:serocomputer}). أي شيفرة تُستعمل لا يحدده المرسِل، بل الثابت الزمني للمستقبِل، والتثبيط هو الذي يضبطه.
\end{proposition}
\end{keyblock}
\section{النقاط والشُّرَط: المشبك مرشِّحًا}

حتى الآن كان المشبك عندنا سلكًا ذا وزن. والحقيقة أن استجابته تتوقف على النبضات السابقة، وهذا يعطي لغة العصبونات رمزًا ثانيًا.

\emph{النضوب: المشبك ينقل التغيرات.} كل قذفة تستهلك حصة $U$ من مخزون الناقل العصبي الجاهز للقذف~\cite{Tsodyks1997}، ويُستعاد المخزون بثابت زمني $\tau_{\text{rec}}$ من رتبة مئات من أجزاء الألف من الثانية. عند التردد $r$ تستقر سعة الاستجابة للنبضة تقريبًا عند المستوى $A(r)=A_0/(1+U r\tau_{\text{rec}})$، حيث $A_0$ استجابة المشبك المرتاح. والتيار الذي يتلقاه المستقبِل في وحدة الزمن يساوي
\begin{empheq}[box=\keybox]{equation}
r\,A(r) \;=\; \frac{A_0\,r}{1+U r\,\tau_{\text{rec}}}
\;\xrightarrow[r\to\infty]{}\; \frac{A_0}{U\tau_{\text{rec}}} .
\label{eq:depression}
\end{empheq}
عند $U=0.5$ و$\tau_{\text{rec}}=0.8\,\text{s}$ يبدأ الإشباع منذ $1/(U\tau_{\text{rec}})=2.5$ نبضة في الثانية. فإذا ارتفع التردد من خمس إلى عشرين، لم يزد التيار المستقر إلا بالثلث. لكن النبضات الأولى على التردد الجديد تصل والمخزون ما زال على حاله، فيقفز التيار لحظةً أربعة أضعاف، ثم يهبط خلال $\tau_{\text{rec}}/(1+Ur\tau_{\text{rec}})\approx90\ms$ (الشكل~\ref{fig:depression}).

هذا المشبك لا ينقل التردد بل \emph{تغيّره}، وفي جوهر الأمر التغير النسبي $\Delta r/r$~\cite{Abbott1997depression}. وبالنسبة لمهندس الإلكترونيات هذا مرشِّح تمرير عالٍ موضوع في السلك مباشرة.

\begin{figure}[t]
\centering
\begin{tikzpicture}[>=Stealth,x=12.5cm,y=1cm]
% input rate
\draw[->,gray!70!black] (0,2.6) -- (0.9,2.6) node[right,black] {\small $t$};
\draw[->,gray!70!black] (0,2.6) -- (0,4.3) node[above,black] {\small $r$};
\draw[very thick,blue!60!black] (0,2.9) -- (0.2,2.9) -- (0.2,3.8) -- (0.8,3.8);
\node[left,font=\scriptsize] at (0,2.9) {$5$};
\node[left,font=\scriptsize] at (0,3.8) {$20$};
\node[right,font=\small,blue!60!black] at (0.4,4.1) {التردد عند مدخل المشبك};
% output current
\draw[->,gray!70!black] (0,0) -- (0.9,0) node[right,black] {\small $t$};
\draw[->,gray!70!black] (0,0) -- (0,2.45) node[left,black] {\small $rA$};
\draw[densely dashed,gray!60!black] (0,0.667) -- (0.8,0.667);
\draw[very thick,red!70!black] (0,0.5) -- (0.2,0.5) -- (0.2,2.0)
   plot[domain=0.2:0.8,samples=60] (\x,{0.667+(2.0-0.667)*exp(-(\x-0.2)/0.089)});
\node[left,font=\scriptsize] at (0,0.5) {$1.7$};
\node[left,font=\scriptsize] at (0,2.0) {$6.7$};
\node[right,font=\scriptsize] at (0.8,0.667) {$2.2$};
\node[right,font=\small,red!70!black] at (0.28,1.6) {التيار إلى المستقبِل: قفزة ثم هبوط خلال $\approx90\ms$};
\draw[gray!60!black] (0.2,-0.08) -- (0.2,0.08); \node[below,font=\scriptsize] at (0.2,0) {$0.2\,\text{s}$};
\draw[gray!60!black] (0.8,-0.08) -- (0.8,0.08); \node[below,font=\scriptsize] at (0.8,0) {$0.8\,\text{s}$};
\end{tikzpicture}
\caption{مشبك ناضب بحسب الصيغة~\eqref{eq:depression} عند $U=0.5$ و$\tau_{\text{rec}}=0.8\,\text{s}$؛ التيار بوحدات $A_0$ في الثانية، والخط المتقطع هو التيار المستقر: لم يرتفع إلا من $1.7$ إلى $2.2$، أما لحظة القفزة فإلى $6.7$. يبلغ المشبك المستقبِلَ أن التردد تغيّر، لا ما هو.}
\label{fig:depression}
\end{figure}

\emph{التيسير: الرشقة شَرطة.} في مشابك أخرى احتمال القذف صغير، لكنه يزداد بعد كل نبضة لعشرات من أجزاء الألف من الثانية. النبضة المفردة عبر مشبك كهذا تضيع في الغالب. أما \emph{الرشقة}، أي بضع نبضات متتالية تفصل بينها بضعة أجزاء من الألف من الثانية، فتمر شبه مؤكدة. حتى من دون تيسير، احتمال أن تضيع نبضات الرشقة الـ$k$ كلها يساوي
\begin{empheq}[box=\keybox]{equation}
P_{\text{loss}} \;=\; (1-p)^k ,
\label{eq:burst}
\end{empheq}
وعند $p=0.2$ هذا $0.8$ للنبضة المفردة و$0.41$ لرشقة من أربع؛ والتيسير يخفضه أكثر. هكذا يظهر لدى العصبون رمز ثانٍ: النبضة المفردة نقطة، والرشقة شَرطة. المشبك الناضب أفضل ما ينقل النبضة الأولى بعد توقف؛ والمشبك التيسيري يمرر الرشقات ويخمد النقاط المفردة. أما أن الرشقات تنتقل بموثوقية أكبر فمُقاس؛ وأما أن الدماغ يستعملها فعلًا رمزًا مستقلًا ففرضية معقولة~\cite{Lisman1997}.

\section{كيف تتكلم عصبونات \nt{GABA}}

أكثرها عددًا في القشرة هي السريعة~\cite{Tremblay2016}. نبضاتها أقصر من نبضات عصبونات \nt{GLUT}، وتستطيع أن تطلقها بانتظام ولمدة طويلة بتردد مئات في الثانية، بلا تعب يُذكر. ومشابكها أوثق: فالوصلة مع كل مستقبِل مؤلفة من مواضع قذف كثيرة، ولا تكاد النبضة تضيع أبدًا. وتستقر مخارجها قرب الموضع الذي تولد فيه نبضة المستقبِل، فهي لا تتحكم في كيفية جمع المستقبِل لمداخله، بل في هل ينطلق ومتى.

وهي نفسها تتلقى الإثارة من عصبونات \nt{GLUT} متجاورة كثيرة وتُبلغ عن الفعالية الكلية من حولها، وهذا بالضبط ما تحتاجه القسمة في الفصل~\ref{sec:gaba}. وأخيرًا، عصبونات \nt{GABA} السريعة متصلة بعضها ببعض وتنطلق معًا بسهولة؛ وهي بالذات التي تحدد الإيقاع المحلي الذي تحدثنا عنه أعلاه~\cite{Cardin2009}.

بلغة شفرة مورس، عصبونات \nt{GABA} ليست مسؤولة عن الحروف، بل عن \emph{الفواصل}. تقطّع التدفق المتصل لنبضات \nt{GLUT} إلى نوافذ يمكن للمستقبِل أن ينطلق داخلها، وتضيّق نوافذ التزامن، وتخفت التدفق كله حين يعلو أكثر مما ينبغي.

\begin{keyblock}
\begin{proposition}[تقسيم الأدوار]
\label{prop:punctuation}
عصبونات \nt{GLUT} تحمل مضمون الرسالة: \emph{ماذا} و\emph{كم}. وعصبونات \nt{GABA} السريعة تحمل علاماتها: \emph{متى} يجوز الكلام و\emph{بأي علوّ}؛ وصنف آخر، بتثبيطه للمثبِّطات، يقرر \emph{لمن} تُفتح البوابة. بعض أجزاء هذا التقسيم مُقاس؛ أما بوصفه قاعدة عامة ففرضية عمل.
\end{proposition}
\end{keyblock}

وثمة عصبونات \nt{GABA} أخرى أبطأ، تستقر مخارجها على مداخل منفردة للمستقبِل. تعمل كالحظر~\eqref{eq:veto} في الفصل~\ref{sec:gaba}: تغلق تدفقًا واحدًا من نبضات \nt{GLUT} دون المساس بغيره.

تقسيم عصبونات \nt{GABA} إلى سريعة وبطيئة صورة قديمة، وهي ناقصة. فاليوم يُميَّز في القشرة ثلاثة أصناف كبرى على الأقل~\cite{Tremblay2016}، والثالث مصمم على نحو غير متوقع: لا يثبط عصبونات \nt{GLUT}، بل عصبونات \nt{GABA} أخرى، ولا سيما تلك التي تغلق المداخل~\cite{Pfeffer2013}. تثبيط التثبيط نفي مزدوج. عصبون كهذا \emph{يفتح} البوابة دون أن يرسل إلى المستقبِل نبضة مثيرة واحدة. تشغّله إشارات من مناطق أخرى في القشرة ومن قنوات البثّ، فيعمل مدخلَ تمكين لجزء كامل من الشبكة: ما دام نشطًا فالحظور مرفوعة، والتدفق يمر~\cite{Pi2013}. وفي الملحق~\ref{sec:othermediators} سُميت هذه الحيلة رفع التثبيط.
\section{لغة مقتصدة}

آخر ما يُعرف بثبات عن لغة العصبونات أنها مكلفة. النبضة، مع عمل المشابك الذي تستدعيه، تكلف نحو $10^9$ جزيء من ATP (تقدير لقشرة القوارض~\cite{Attwell2001})، والجزيء الواحد يعطي نحو $10^{-19}$ جول. إذن تكلف النبضة من رتبة $E_{\text{spk}}\sim10^{-10}$ جول. يستهلك الدماغ كله نحو $20$ واط، وإذا ذهب نصفها إلى نقل الإشارات، $P\sim10$ واط، فالتردد المتوسط للعصبون الواحد
\begin{empheq}[box=\keybox]{equation}
\bar r \;\approx\; \frac{P}{N\,E_{\text{spk}}}
\;\sim\; \frac{10\ \text{W}}{8.6\cdot10^{10}\cdot10^{-10}\ \text{J}}
\;\approx\; 1\ \text{spike/s}.
\label{eq:energy}
\end{empheq}
والتقديرات الأكثر تفصيلًا للقشرة تعطي أقل من ذلك~\cite{Lennie2003}. شيفرة الدماغ مضطرة إلى أن تكون \emph{متناثرة}: لا يستطيع أن يعمل بسرعة إلا جزء صغير من العصبونات.

ومن~\eqref{eq:spikeentropy} يتبين أن هذا ليس سيئًا إلى هذا الحد. بدقة $1\ms$ لا تحمل نبضة عصبون يعمل بتردد $100$ في الثانية أكثر من خمسة بتات، أما نبضة عصبون يعمل مرة في الثانية فتحمل حتى أحد عشر. وإذا كان الدفع على النبضات، فالأربح أن تتكلم نادرًا. والقشرة فعلًا تقايض الدقة بالطاقة: عند شح الغذاء تحافظ العصبونات على تردد نبضاتها، لكنها تنفق أقل على التيارات المشبكية، فتصبح استجاباتها أقل دقة~\cite{Padamsey2022}.

في شفرة مورس الحروف الشائعة قصيرة: أشيع حرف في النص الإنجليزي، \foreignlanguage{english}{E}، نقطة واحدة. وشيفرة العصبونات، بحسب ما هو معروف، تتبع القاعدة نفسها في صورة أخرى: المتوقع يكلف نبضات قليلة~\cite{Barlow1961}. العصبون يخفض تردده تدريجيًا مع المنبه الثابت، ومستشعرات الفصل~\ref{sec:glutcomputer} تستجيب للتغيرات لا للمستويات، وعصبونات \nt{DOPA} في الفصل~\ref{sec:neurontypes} لا تبلغ إلا عن خطأ التنبؤ بالمكافأة.

\begin{keyblock}
\begin{proposition}[الاقتصاد]
\label{prop:economy}
تحدّ الطاقة التردد المتوسط للعصبون بقيمة من رتبة نبضة واحدة في الثانية. لذلك فلغة عصبونات \nt{GLUT} متناثرة، وتنفق نبضاتها على الأخبار: على ما تغيّر أو لم يُتنبأ به. الحد الطاقي مُقاس؛ أما أن شيفرة الدماغ مضبوطة على أكبر عدد من البتات لكل جول ففرضية حسنة الأسس.
\end{proposition}
\end{keyblock}

\section{الخلاصة}

\begin{table}[t]
\centering
\footnotesize
\setlength{\tabcolsep}{4pt}
\renewcommand{\arraystretch}{1.15}
\begin{tabular}{>{\raggedright}p{2.6cm}>{\raggedright}p{2.3cm}>{\raggedright}p{3.0cm}>{\raggedright}p{2.0cm}>{\raggedright\arraybackslash}p{3.6cm}}
\toprule
طريقة الكتابة & ماذا تحمل & من يقرؤها & زمن الرسالة & ما هو معروف \\
\midrule
رقم العصبون المنطلق (الترميز بالموضع) & أي قيمة & أي مستقبِل؛ كواشف التزامن & تأخير واحد & ثابت في العصبونات الحسية وفي تحديد اتجاه الصوت \\
تردد عصبون واحد & المقدار & مكامل بـ$\tau$ كبير؛ الحجم (الفصل~\ref{sec:serocomputer}) & من مئات الأجزاء من الألف من الثانية إلى ثوانٍ & ثابت؛ وبطيء جدًا لمرحلة معالجة من $10\text{--}20\ms$ \\
تردد مجموعة & المقدار & عصبون بآلاف المداخل & $10\text{--}50\ms$ & ثابت؛ والربح تحده الارتباطات~\eqref{eq:correlated} \\
زمن النبضة الأولى، الترتيب & المقدار، الترتيب & كواشف التزامن ذات التأخيرات & تأخير واحد & ثابت عند مدخل الجهاز العصبي؛ وفي القشرة موضع نقاش \\
الطور في الإيقاع المحلي & المقدار، الموضع & عصبونات ذات إيقاع مشترك & دورة $12\text{--}50\ms$ وأبطأ & مرصود في مناطق بعينها؛ ودوره العام فرضية \\
الرشقة (”الشَّرطة“) & ”مهم“: إيصال موثوق & المشبك التيسيري & $\sim10\ms$ & الموثوقية مقيسة؛ والرمز المستقل فرضية \\
تغيّر التردد & خبر & المشبك الناضب & $\sim0.1\,\text{s}$ & ثابت \\
نبضات عصبونات \nt{GABA} السريعة & متى وبأي علوّ & كل عصبونات \nt{GLUT} المجاورة & $1\text{--}30\ms$ & الإيقاع والقسمة مقيسان؛ و”علامات الترقيم“ قاعدةً فرضية \\
\bottomrule
\end{tabular}
\caption{الطرق التي تكتب بها عصبونات \nt{GLUT} و\nt{GABA} الرسائل بالنبضات. العمود الأيسر يفصل المقيس عن المفترض.}
\label{tab:code}
\end{table}

يجمع الجدول~\ref{tab:code} ما نعرفه؛ ويظهر منه أيضًا ما لا نعرفه. لا نعرف إلى أي حد تستعمل القشرة التوقيت الدقيق للنبضات، لا ترددات المجموعات فقط. ولا نعرف هل للعصبونات ”كلمات“، أي توليفات مستقرة من نبضات عصبونات كثيرة تُقرأ كلًّا واحدًا. ولا نعرف هل اللغة واحدة في أجزاء الدماغ المختلفة. يمكن تعلم شفرة مورس في أمسية، لأن جدولها معروف. أما جدول لغة العصبونات فما زال ينتظر من يضعه، والأرجح أن من سيضعه مهندسو الاتصالات.

\section{تمارين}

\begin{exercise}[\lvlA]
\label{ex:04:1}
\emph{السعة بالأرقام.} $\Delta t=1\ms$. (أ)~كم بتًا لكل جول يعطي عصبون تردده $1$ نبضة في الثانية، بكلفة النبضة من~\eqref{eq:energy}؟ (ب)~عند $r=10$، كم مرة أقل من الحد~\eqref{eq:spikeentropy} يستخرج مستقبِل لا يفعل سوى عدّ النبضات في الثانية؟
\end{exercise}

\begin{exercise}[\lvlB]
\label{ex:04:2}
\emph{التردد الأمثل.} ينفق العصبون قدرة $P_0+rE_{\text{spk}}$، حيث $P_0$ هو النصف الآخر من $20\,\text{W}$ مقسومًا على كل العصبونات، و$E_{\text{spk}}=10^{-10}\,\text{J}$. عند أي تردد $r$ يبلغ عدد البتات لكل جول $R_{\max}(r)/(P_0+rE_{\text{spk}})$ بحسب~\eqref{eq:spikeentropy} مع $\Delta t=1\ms$ أقصاه؟
\end{exercise}

\begin{exercise}[\lvlB]
\label{ex:04:3}
\emph{أي ارتباط ضار.} $N=1000$ عصبون، التباين $\sigma^2=1$، والارتباط الثنائي $c=0.1$. احسب معلومات فيشر $\mathcal I=f'^{\mathsf T}\Sigma^{-1}f'$ ($\Sigma$ مصفوفة التغاير) لميول متساوية $f'=\mathbf 1$ ولميول $\pm1$ مع $\mathbf 1^{\mathsf T}f'=0$. كم مرة تختلفان؟
\end{exercise}

\begin{exercise}[\lvlB]
\label{ex:04:4}
\emph{نمذجة: كاشف التزامن.} العصبون~\eqref{eq:glutneuron} يتلقى $1000$ مدخل بواسوني بـ$5$ نبضات في الثانية لكل منها، $w=1$؛ والعتبة $\theta$ بحيث يعبرها الجهد الحر مرة في الثانية. أوجد بالنمذجة كم مدخلًا متزامنًا $m$ يُطلق العصبون باحتمال $1/2$ عند $\tau=10$ و$2\ms$.
\end{exercise}

\begin{exercise}[\lvlB]
\label{ex:04:5}
\emph{تصميم: كاشف التغيرات النسبية.} اختر مشبكًا~\eqref{eq:depression}: التيار المستقر عند $40$ نبضة في الثانية لا يزيد بأكثر من $10\,\%$ عليه عند $10$، وبعد القفزة إلى $40$ يبلغ التيار المستوى الجديد بثابت زمني لا يتجاوز $50\ms$. ما أصغر $\tau_{\text{rec}}$ وعند أي $U$؟
\end{exercise}

\begin{exercise}[\lvlA]
\label{ex:04:6}
\emph{زمن النبضة الأولى والرشقات.} (أ)~بحسب~\eqref{eq:latency} أوجد المدخل الذي تأتي عنده النبضة الأولى بعد ثابت زمني واحد بالضبط. علامَ يتوقف؟ (ب)~عند احتمال قذف $p=0.2$، كم نبضة تلزم في الرشقة ليكون احتمال ضياعها كلها أقل من $5\,\%$؟
\end{exercise}

\begin{exercise}[\lvlC]
\label{ex:04:7}
\emph{بحث: كم بتًا في الترتيب المكاني.} العصبون خانات مستقلة~\eqref{eq:spikeentropy}، $r=10$، $\Delta t=1\ms$. (أ)~فكّك إنتروبيا ”كلمة“ من $20$ خانة إلى إنتروبيا عدد النبضات والباقي. (ب)~نمذج تقدير الإنتروبيا من ترددات الكلمات بـ$N=30$ و$10^4$ تسجيل. بكم يكون التقدير منخفضًا؟
\end{exercise}

\chapter{ماذا تحسب عصبونات \nt{SERO}}
\label{sec:serocomputer}

\section{أول قناة بثّ}

حتى الآن تخاطبت كل عصبونات الكتاب عبر ناقل العناوين: كان \nt{GLUT} و\nt{GABA} يرسلان النبضات إلى متلقّين محددين، وعرفنا ماذا تحسب شبكة كهذه وبأي لغة تتكلم. ننتقل الآن إلى الناقل الثاني من القسم~\ref{sec:buses}، وهو ناقل البثّ. أولى قنواته \nt{SERO}. وكما من قبل، لا نسمح لأنفسنا إلا بما يوجد في الكائن الحي.

في هذا الفصل تظهر ثلاث أدوات ستستعملها بعد ذلك كل قنوات البثّ: عصبون يعمل من تلقاء نفسه ما دام يصله دفق خلفي ثابت، إلى أن يوقفه شيء؛ وسلك هو في الحقيقة حجم من النسيج؛ وتركيز بطيء بدل النبضات المنفردة. وستُبنى \nt{DOPA} و\nt{NORA} و\nt{ACET} في الفصول~\babelsublr{\ref{sec:dofa}--\ref{sec:acet}} من القطع نفسها.

من وجهة نظر المهندس، يختلف عصبون \nt{SERO} في أربعة أمور مهمة.

\emph{الأول: المتلقي هو الذي يختار الإشارة.} للسيروتونين مستقبلات بالإشارتين، والقاعدة~\eqref{eq:dalesign} لا تنطبق هنا: لا يحدد الإشارةَ المرسِلُ، بل مستقبِلُ المتلقي (القضية~\ref{prop:receptor})~\cite{BarnesSharp1999}.

\emph{الثاني: عصبون \nt{SERO} يعمل من تلقاء نفسه إذا توفّر دفق خلفي ثابت.} في اليقظة يطلق النبضات بانتظام، بضع مرات في الثانية، ولا يحتاج إلى دفعة لكل نبضة: إنه ناظمة إيقاع. لكنه يحتاج إلى دفق خلفي دائم. في شريحة الدماغ، حيث لا يوجد هذا الدفق، تصمت معظم هذه الخلايا؛ ويعود الإيقاع المنتظم إذا أُعطي النورأدرينالين عبر مستقبلات $\alpha_1$، ويخمده حجبُها~\cite{VandermaelenAghajanian1983}. في الكائن الحي يأتي هذا الدفق الخلفي من \nt{NORA}. وبلغة الدارات هذا مذبذب يحتاج إلى تغذية كهربائية: ما دامت موجودة يعمل وحده، إلى أن يوقفه التثبيط.

\emph{الثالث: إنه بطيء.} تكاد كل مستقبلات السيروتونين تكون أيضية التوجه: لا تفتح قناة بنفسها، بل تطلق داخل الخلية سلسلة من التفاعلات. والاستجابة بطيئة: يرتفع التيار المثبِّط عبر المستقبل الرئيسي ثم يهبط في غضون ثانية تقريبًا~\cite{CourtneyFord2016}، أي أطول بعشرات المرات من استجابة مشبك \nt{GLUT} السريع.

\emph{الرابع: يحرّر الناقل لا في سلك بل في حجم.} في المناطق التي تصل إليها مخارج عصبونات \nt{SERO} كثيرًا ما يخرج السيروتونين لا إلى شق ضيق، بل إلى الحيز بين الخلايا، وينتشر حوله~\cite{Bunin1998}. يحسّ المتلقي بالتركيز، ولا يستطيع أن يعرف من أي مرسِل جاءت الجزيئة.

مع هذه الأزمنة لم تعد النبضات المنفردة مهمة، لذلك سنصف الشبكة بالترددات $r_j$ وبالتركيزات. يزداد تركيز السيروتونين $c$ في الحجم الذي تحرّر فيه العصبونات $j$ الناقل بفعل التحرير، وينقص لأن الناقل يُضَخّ عائدًا إلى الخلايا:
\begin{empheq}[box=\keybox]{equation}
\tau_c\,\frac{\dd c}{\dd t} \;=\; -c \;+\; \kappa\sum_j r_j ,
\qquad
r_i \;=\; f_i\bigl(b_i + s_i\,g\,c\bigr), \quad s_i=\pm1 .
\label{eq:seroneuron}
\end{empheq}
هذه طريقة أخرى لقراءة سيل النبضات تكمّل القضية~\ref{prop:reader}: الحجم لا يرى النبضات المنفردة ولا ترتيبها، بل التردد فقط، متوسَّطًا على الزمن $\tau_c$. المعادلة الأولى هي دارة \foreignlanguage{english}{RC} نفسها التي رأيناها في الغشاء: $\tau_c$ عمر الناقل في الحجم، ولم يُقَس مباشرة، فنفترضه من رتبة الثانية؛ و$\kappa$ مقدار الناقل الذي تعطيه نبضة واحدة. والمعادلة الثانية تصف المتلقي: $b_i$ دفقه الخاص، وهو موجب عند ناظمة الإيقاع (ويدخل فيه الدخل الخلفي الثابت من \nt{NORA})؛ و$g$ حساسية المستقبل؛ والإشارة $s_i$ يختارها المتلقي؛ و$f_i$ خاصية نقل غير متناقصة.

\section{بوابة \foreignlanguage{english}{NOT} من عصبون واحد}

لنأخذ ناظمة إيقاع ذات مستقبل مثبِّط، $s=-1$. ما دام لا سيروتونين حولها فهي تعمل على دفقها الخاص $b$. وعندما يظهر السيروتونين ينقص الدفق بمقدار $gc$، وإذا كان $gc>b$ يصمت العصبون. هذه بوابة \foreignlanguage{english}{NOT}: يوجد خرج حين لا يوجد دخل. وقد كلّفتنا عصبونًا واحدًا (الشكل~\ref{fig:serogates}). والقضية~\ref{prop:monotone} لا تنطبق هنا: فقد اشترطت أوزانًا موجبة.

إذا أثّر في ناظمة الإيقاع نفسها سيروتونين من حجمين مختلفين، فإنها تصمت متى وُجد السيروتونين في أحدهما على الأقل. هذه بوابة \foreignlanguage{english}{NOR}، ومن بوابات \foreignlanguage{english}{NOR} يمكن بناء أي دالة منطقية. ولا حاجة هنا إلى الترميز ثنائي الأسلاك الذي تستعمله شبكة \nt{GLUT}: غياب الإشارة تحسبه عصبونات تعمل إلى أن يوقفها شيء.

\begin{figure}[t]
\centering
\begin{tikzpicture}[>=Stealth,
  nrn/.style={circle,draw,very thick,fill=blue!6,minimum size=0.95cm,inner sep=0pt},
  pace/.style={nrn,double,double distance=1.2pt},
  inh/.style={-{Bar[width=7pt]},thick,red!70!black}]
\node[pace] (N1) at (0,0) {};
\draw[inh] (-1.6,0) node[left,black] {$c$} -- (N1);
\draw[->,very thick,red!70!black] (N1) -- (1.3,0) node[right,black] {$\bar c$};
\node[below=0.55cm] at (0,0) {\small \foreignlanguage{english}{NOT}};
\node[pace] (N2) at (5,0) {};
\draw[inh] (3.4,0.45) node[left,black] {$c_1$} -- (N2);
\draw[inh] (3.4,-0.45) node[left,black] {$c_2$} -- (N2);
\draw[->,very thick,red!70!black] (N2) -- (6.3,0) node[right,black] {$\overline{c_1\vee c_2}$};
\node[below=0.55cm] at (5,0) {\small \foreignlanguage{english}{NOR}};
\end{tikzpicture}
\caption{المنطق بعصبونات \nt{SERO}. الدائرة المزدوجة ناظمة إيقاع: مع دفق خلفي ثابت تعمل وحدها إلى أن تُثبَّط. الخط المنتهي بشرطة مستقبل مثبِّط، $s=-1$؛ و$c$ و$c_1$ و$c_2$ تركيزات السيروتونين في الأحجام المحيطة بالمتلقي. على اليسار \foreignlanguage{english}{NOT}، وعلى اليمين \foreignlanguage{english}{NOR}.}
\label{fig:serogates}
\end{figure}

\begin{keyblock}
\begin{proposition}[اكتمال منطق \nt{SERO}]
\label{prop:serocomplete}
إذا كانت في الشبكة نواظم إيقاع ذات مستقبلات مثبِّطة، ولم تختلط التحريرات في الأحجام المختلفة، استطاعت الشبكة~\eqref{eq:seroneuron} أن تحسب أي دالة منطقية، ويُنقل كل بت بسلك واحد.
\end{proposition}
\end{keyblock}

من الناحية المنطقية شبكة \nt{SERO} أقوى من \nt{GLUT}، لكن الشرط ”لا تختلط التحريرات في الأحجام المختلفة“ لا يتحقق في الدماغ (القسم~\ref{sec:volume}).

\section{التغذية الراجعة السالبة}

لعصبونات السيروتونين مستقبلات مثبِّطة عليها هي نفسها~\cite{Sprouse1987}. فالسيروتونين الذي حرّرته يعود إليها ويكبحها. هذه دارة مألوفة للمهندس: مضخّم بتغذية راجعة سالبة (الشكل~\ref{fig:serofeedback}).

ما المقيس هنا وما النموذج؟ في نواة الرفاء نفسها، حيث تقع عصبونات \nt{SERO}، يثبّط السيروتونين الجيران لا عبر حجم مشترك، بل من نقطة إلى نقطة عبر المشابك: يزيله بروتين النقل بسرعة ولا يدعه يتراكم خارج الشق. والتحرير المنفرد لا يعطي إلا توقفًا قصيرًا في التفريغ، أما التردد المتوسط فلا يتغير~\cite{CourtneyFord2016}. لذلك فالحلقة أدناه، المغلقة عبر حجم مشترك، نموذج: نحسب ما كانت ستفعله تغذية راجعة كهذه لو عملت على مستوى الترددات المتوسطة.

\begin{figure}[t]
\centering
\begin{tikzpicture}[>=Stealth,
  nrn/.style={circle,draw,very thick,fill=blue!6,minimum size=0.95cm,inner sep=0pt},
  pace/.style={nrn,double,double distance=1.2pt},
  blk/.style={draw,very thick,rounded corners=2pt,fill=orange!10,minimum height=0.9cm,minimum width=2.2cm,align=center,font=\small},
  inh/.style={-{Bar[width=7pt]},thick,red!70!black}]
\node[pace] (N) at (0,0) {};
\node[blk] (V) at (3.6,0) {الحجم\\$\tau_c$};
\draw[->,very thick] (N) -- node[above] {\small $r$} (V);
\draw[->,very thick] (V) -- node[above] {\small $c$} (6.0,0) node[right] {\small إلى كل المتلقين};
\draw[inh] (V.south) -- ++(0,-0.9) -| node[pos=0.25,below] {\small $-g\,c$} (N.south);
\draw[->,thick,blue!60!black] (-1.7,0) node[left,black] {\small $b$} -- (N);
\end{tikzpicture}
\caption{عصبون \nt{SERO} بتغذية راجعة عبر ناقله نفسه: يذهب التركيز $c$ إلى كل المتلقين ويثبّط العصبون نفسه. في النموذج هذا منظِّم مستوى؛ أما أن للحلقة هذا الدور في الدماغ ففرضية.}
\label{fig:serofeedback}
\end{figure}

لنحسب ما تفعله هذه الحلقة. لتكن خاصية النقل خطية، $f(u)=au$ عند $u>0$. في التوازن $c=\kappa r$ و$r=a(b-gc)$. وبتعويض إحداهما في الأخرى نحصل على
\begin{empheq}[box=\keybox]{equation}
r \;=\; \frac{a\,b}{1+G}, \qquad G \;=\; a\,g\,\kappa .
\label{eq:feedback}
\end{empheq}
المقدار $G$ هو كسب الحلقة. إنها الصيغة نفسها لأي مضخّم بتغذية راجعة سالبة، وتعطي المكسبين نفسيهما.

\emph{المناعة من التشتت.} لنفترض أن استثارية العصبون $a$ تغيّرت بنسبة $\varepsilon$. من دون تغذية راجعة كان التردد سيتغير بالنسبة نفسها. أما مع التغذية الراجعة، عند $G\gg1$، فيكاد التردد $r\approx b/(g\kappa)$ لا يعتمد على $a$: تغيّره أصغر بـ$1+G$ مرة. عند $G=9$ يُكبَت التشتت عشر مرات.

\emph{السرعة.} نعوّض $r=a(b-gc)$ في المعادلة الأولى من~\eqref{eq:seroneuron} فنحصل على $\tau_c\,\dd c/\dd t=-(1+G)\,c+\kappa ab$. يقترب التركيز من مستواه بثابت زمني $\tau_c/(1+G)$، أي أسرع بـ$1+G$ مرة مما هو من دون تغذية راجعة.

هذا أول حساب يمكن أن يؤديه نظام \nt{SERO}: أن يُبقي المستوى العام للسيروتونين في الدماغ عند قيمة مضبوطة، كما يُبقي الترموستات درجة الحرارة. هذه فرضية: في التجارب لا يظهر التثبيط الذاتي حتى الآن إلا توقفاتٍ قصيرة لا تغيّر التردد المتوسط.

\section{من النبضات إلى المستوى}

الحساب الثاني مخبوء في المعادلة الأولى من~\eqref{eq:seroneuron}. تطلق العصبونات نبضات منفردة، أما المتلقي فيحسّ بالتركيز، الذي ينعّمها على الزمن $\tau_c$. وإذا تغيّر تردد المصدر تبعه التركيز متأخرًا:
\begin{equation}
c(t) \;=\; \frac{\kappa}{\tau_c}\int_{-\infty}^{t} e^{-(t-t')/\tau_c}\,\sum_j r_j(t')\,\dd t' .
\label{eq:lowpass}
\end{equation}
هذا متوسط متحرك لنشاط المصادر على آخر بضعة $\tau_c$، أي مرشّح تمرير منخفض (الشكل~\ref{fig:lowpass}). وبهذه الطريقة يمرّر أبسط محوّل رقمي-تماثلي النبضاتِ عبر دارة \foreignlanguage{english}{RC} ويحوّل ترددها إلى مستوى. ونظام \nt{SERO} يفعل الشيء نفسه بمئات الآلاف من العصبونات دفعة واحدة.

وهذا المقدار ذاكرة أيضًا. بعد أن تصمت المصادر يبقى التركيز قائمًا زمنًا من رتبة $\tau_c$. إنه ليس بتًا، بل أثر يذوب تدريجيًا.

\begin{figure}[t]
\centering
\begin{tikzpicture}[>=Stealth,x=1.45cm,y=1cm]
\draw[gray!70!black] (0,3.2) -- (8.1,3.2);
\node[left,font=\small] at (-0.05,3.4) {النبضات};
\node[above,font=\scriptsize,gray!40!black] at (1,3.65) {$2$ في الثانية};
\node[above,font=\scriptsize,gray!40!black] at (3.5,3.65) {$8$ في الثانية};
\node[above,font=\scriptsize,gray!40!black] at (6.5,3.65) {$2$ في الثانية};
\draw[->,gray!70!black] (0,0) -- (8.3,0) node[right,black] {\small $t$};
\draw[->,gray!70!black] (0,0) -- (0,2.75) node[left,black] {\small $c$};
\foreach \s in {0.136,0.529,0.760,1.223,1.714,1.748,1.755,2.663,2.701,2.734,3.414,3.493,3.719,3.800,3.928,3.948,4.074,4.327,4.420,4.589,4.728,4.736,4.914,5.025,5.205,5.220,6.224,6.544,7.178}{\draw[thick,blue!60!black] (\s,3.2) -- (\s,3.6);}
\foreach \a/\b/\v in {0.000/0.136/0.000,0.136/0.529/1.000,0.529/0.760/1.675,0.760/1.223/2.330,1.223/1.714/2.466,1.714/1.748/2.509,1.748/1.755/3.425,1.755/2.663/4.402,2.663/2.701/2.775,2.701/2.734/3.672,2.734/3.414/4.553,3.414/3.493/3.306,3.493/3.719/4.055,3.719/3.800/4.235,3.800/3.928/4.905,3.928/3.948/5.316,3.948/4.074/6.211,4.074/4.327/6.476,4.327/4.420/6.028,4.420/4.589/6.493,4.589/4.728/6.483,4.728/4.736/6.642,4.736/4.914/7.589,4.914/5.025/7.351,5.025/5.205/7.579,5.205/5.220/7.331,5.220/6.224/8.221,6.224/6.544/4.012,6.544/7.178/3.914,7.178/8.000/3.076}{\draw[very thick,orange!85!black] plot[domain=\a:\b,samples=10] (\x,{0.25*\v*exp(-(\x-\a))});}
\foreach \s/\p/\q in {0.136/0.000/1.000,0.529/0.675/1.675,0.760/1.330/2.330,1.223/1.466/2.466,1.714/1.509/2.509,1.748/2.425/3.425,1.755/3.402/4.402,2.663/1.775/2.775,2.701/2.672/3.672,2.734/3.553/4.553,3.414/2.306/3.306,3.493/3.055/4.055,3.719/3.235/4.235,3.800/3.905/4.905,3.928/4.316/5.316,3.948/5.211/6.211,4.074/5.476/6.476,4.327/5.028/6.028,4.420/5.493/6.493,4.589/5.483/6.483,4.728/5.642/6.642,4.736/6.589/7.589,4.914/6.351/7.351,5.025/6.579/7.579,5.205/6.331/7.331,5.220/7.221/8.221,6.224/3.012/4.012,6.544/2.914/3.914,7.178/2.076/3.076}{\draw[very thick,orange!85!black] (\s,0.25*\p) -- (\s,0.25*\q);}
\draw[thick,densely dashed,black] plot[domain=0:2,samples=30] (\x,{0.25*2*(1-exp(-\x))})
  plot[domain=2:5,samples=40] (\x,{0.25*(8-6.271*exp(-(\x-2)))})
  plot[domain=5:8,samples=40] (\x,{0.25*(2+5.688*exp(-(\x-5)))});
\foreach \x in {0,2,5,8}{\draw[gray!60!black] (\x,-0.05) -- (\x,0.05) node[below=3pt,font=\scriptsize] {$\x\,\text{s}$};}
\end{tikzpicture}
\caption{من النبضات إلى المستوى. في الأعلى نبضات عصبونات \nt{SERO}: متفرقة ثانيتين، وكثيفة ثلاث ثوانٍ، ثم متفرقة من جديد. في الأسفل تركيز السيروتونين~\eqref{eq:lowpass} عند $\tau_c=1\,\text{s}$. الخط المتقطع هو التركيز نفسه لو جاءت النبضات بانتظام تام.}
\label{fig:lowpass}
\end{figure}

\section{السلك هو الحجم}
\label{sec:volume}

لنعد إلى شرط القضية~\ref{prop:serocomplete}. السيروتونين الذي يحرّره مرسِلون مختلفون متجاورون يتجمع في تركيز واحد.

\emph{السلك هنا ليس ليفًا، بل حجم من النسيج.} لا تبقى إشارتان متمايزتين إلا إذا حُرّرتا في منطقتين تفصل بينهما مسافة أكبر مما يتسع الوقت لانتشار الناقل فيه. خلال الزمن $\tau$ تبتعد جزيئة معامل انتشارها $D$ مسافة
\begin{empheq}[box=\keybox]{equation}
\ell \;\sim\; \sqrt{D\,\tau}\, .
\label{eq:diffusionlength}
\end{empheq}
يُبطئ تعرّجُ الشقوق بين الخلايا انتشارَ الجزيئة الصغيرة نحو $2.6$ مرة~\cite{Sykova2008}، أما معامل انتشار السيروتونين نفسه في النسيج فلم يُقَس؛ ومن حجم الجزيئة نقدّره ببضع وحدات من $10^{-6}\,\text{cm}^2/\text{s}$. إذا عاشت الإشارة $\tau\sim1\,\text{s}$ (وهذا تقدير أيضًا)، فإن $\ell\sim\sqrt{3\cdot10^{-6}}\,\text{cm}\approx20\um$. العنوان في شبكة كهذه هو \emph{المكان}: لا يعرف المتلقي مصدر الإشارة إلا من الموضع الذي يوجد هو فيه.

عصبونات \nt{SERO} الحقيقية مبنية بحيث لا يكاد يبقى شيء من هذه العناوين المنفصلة. خرج كل منها موزع على مساحات هائلة من الدماغ: فروع الخلية الواحدة تبلغ مناطق كثيرة متباعدة~\cite{GagnonParent2014}. يُقدَّر عدد مواقع التحرير في الخلية بنحو $10^5$ (الجدول~\ref{tab:types})، ولم يعدّها أحد مباشرة. ”أسلاك“ عصبونات \nt{SERO} المختلفة تتداخل وتختلط. ومع ذلك لا تندمج كليًا في سلك واحد: تنقسم عصبونات \nt{SERO} إلى بضعة أنظمة فرعية ترسل خرجها إلى مناطق مختلفة وتستجيب للحدث نفسه استجابات مختلفة~\cite{Ren2018}. لكن هذه الأنظمة الفرعية آحاد، لا مئات آلاف.

\begin{keyblock}
\begin{proposition}[عدد الأسلاك المستقلة]
\label{prop:volumewires}
في شبكة ذات نقل حجمي، لا يحدّ عددَ الإشارات المستقلة عددُ العصبونات، بل عددُ المناطق غير المتداخلة ذات الحجم $\ell\sim\sqrt{D\tau}$ التي تحرّر فيها الناقل. وإذا كان خرج كل عصبون موزعًا على حجم كبير، فلا تنقل الشبكة كلها إلا بضعة متغيرات مشتركة.
\end{proposition}
\end{keyblock}

لذلك لا يوجد في الدماغ ما تُبنى منه الدارات المنطقية للقضية~\ref{prop:serocomplete}. أما المنظِّم~\eqref{eq:feedback} والمرشّح~\eqref{eq:lowpass} فلا يحتاجان إلى أسلاك منفصلة: يكفيهما مقدار مشترك واحد. المرشّح ينتج مباشرة من التحرير في الحجم، وهو مقيس في مناطق الهدف~\cite{Bunin1998}؛ أما منظِّم المستوى ففرضية.

\section{أفق التخطيط}
\label{sec:horizon}

فيمَ قد ينفع مقدار مشترك بطيء واحد؟ أنسب ما يخطر هنا معامِل يجب أن يكون واحدًا للشبكة كلها، وألا يتغير أسرع من ثوانٍ أو دقائق. في الفصل~\ref{sec:neurontypes} ظهر معامل كهذا: المعامل $\gamma$ في صيغة الخطأ~\eqref{eq:rpe}، الذي يقول بكم تقلّ قيمة المكافأة المقبلة عن المكافأة الفورية. في الزمن المتصل يحل محله \emph{الأفق} $\tau_\gamma$: المكافأة $R$ التي يجب انتظارها زمنًا $D$ تساوي الآن $R\,e^{-D/\tau_\gamma}$.

الأفق يقرر هل يستحق الأمر الانتظار. لنفترض أنه يمكن أخذ مكافأة صغيرة $r_0$ فورًا أو انتظار مكافأة كبيرة $R$. يكون الانتظار مجديًا ما دام
\begin{empheq}[box=\keybox]{equation}
D \;<\; \tau_\gamma\,\ln\frac{R}{r_0} .
\label{eq:patience}
\end{empheq}
عند $\tau_\gamma=2\,\text{s}$ ومكافأة مؤجلة أكبر بثلاث مرات يستحق الأمر الانتظار حتى $2.2\,\text{s}$؛ وإذا كان الأفق أطول بمرتين، فحتى $4.4\,\text{s}$. الصبر يتناسب مع الأفق.

تستند الصيغة~\eqref{eq:patience} إلى خصم أسّي. أما الخصم المقيس على تأخيرات من رتبة الثواني فأقرب إلى المنحنى الزائدي $R/(1+D/\tau_\gamma)$: هكذا تنخفض عند القردة قيمةُ المكافأة المؤجلة في الاختيار، واستجابةُ عصبونات \nt{DOPA} للإشارة التي تُنبئ بها~\cite{KobayashiSchultz2008}. ومع المنحنى الزائدي يحل محل الشرط~\eqref{eq:patience} الشرطُ $D<\tau_\gamma\,(R/r_0-1)$: لم يعد الاعتماد على نسبة المكافأتين لوغاريتميًا بل خطيًا، ومع الأرقام نفسها تنزاح حدود الانتظار إلى الضعف تقريبًا. ويُستخرج المنحنى الزائدي من الأسّية نفسها إذا كان التأخير عشوائيًا (التمرين~\ref{ex:05:7})، ويبقى الصبر متناسبًا مع مقياس الزمن.

وفق فرضيةٍ ما، \nt{SERO} هو الذي يضبط الأفق~\cite{Doya2002}. ولديه كل ما يلزم لهذا الدور: مستوى مشترك واحد للشبكة كلها، يتغير خلال ثوانٍ. وثمة تجارب مباشرة أيضًا: إذا شُغّلت عصبونات السيروتونين اصطناعيًا، انتظر الحيوان المكافأة المؤجلة مدة أطول قبل أن يكفّ عن الانتظار~\cite{Miyazaki2014}، وواصل محاولة الحصول على المكافأة مدة أطول في المكان الذي صارت فيه نادرة~\cite{Lottem2018}. أما كيف يتحول تركيز السيروتونين إلى $\tau_\gamma$ داخل الشبكة فغير معروف. في الفصل التالي سيدخل الأفق في الخطأ الذي يبثّه \nt{DOPA}.

\section{الخلاصة}

\begin{table}[t]
\centering
\small
\begin{tabular}{>{\raggedright}p{3.6cm}>{\raggedright}p{4.3cm}>{\raggedright\arraybackslash}p{4.3cm}}
\toprule
& \nt{GLUT} & \nt{SERO} \\
\midrule
زمن الاستجابة & ميلي ثوانٍ & من أجزاء الثانية إلى ثانية \\
إشارة الأثر & $+$ فقط & $+$ و$-$، يختارها المتلقي \\
من دون دخل & يصمت & يعمل وحده مع دفق خلفي ثابت \\
العنونة & من نقطة إلى نقطة & حجم، والعنوان هو المكان \\
\foreignlanguage{english}{NOT} & عبر مستشعرات ”الإطفاء“ فقط & عصبون واحد \\
الذاكرة & نشاط ذاتي الاستمرار، يخبو بالإجهاد & تركيز، يذوب خلال $\tau_c$ \\
الحسابات الرئيسية & التزامن، والتأخيرات، والمنطق، والتسلسلات & التنعيم؛ وضبط المستوى والأفق $\tau_\gamma$ (فرضيتان) \\
\bottomrule
\end{tabular}
\caption{ما تحسبه شبكات من عصبونات \nt{GLUT} وحدها ومن عصبونات \nt{SERO} وحدها.}
\label{tab:glutvssero}
\end{table}

بوصفه عنصرًا منطقيًا (الجدول~\ref{tab:glutvssero}) عصبون \nt{SERO} أغنى من \nt{GLUT}، لكنه بوصفه نظام اتصال بثٌّ عام: بطيء ويكاد يخلو من العناوين. لذلك لا يحاول أن يكون معالجًا، بل ينعّم ويبثّ بضعة مقادير مشتركة تحدثنا عنها في الفصل~\ref{sec:neurontypes}، ومنها، وفق فرضيةٍ، أفق التخطيط. وسنلتقي ناظمة الإيقاع والحجم والتركيز البطيء ثلاث مرات أخرى: عند \nt{DOPA} و\nt{NORA} و\nt{ACET}.

\section{تمارين}

\begin{exercise}[\lvlA]
\label{ex:05:1}
\emph{السلك هو الحجم.} خذ للسيروتونين $D=3\cdot10^{-6}\,\text{cm}^2/\text{s}$ (القسم~\ref{sec:volume}). جد من~\eqref{eq:diffusionlength} قيمة $\ell$ لإشارة تعيش $1\,\text{s}$، والزمن الذي ينتشر فيه الناقل مسافة $5\,\text{cm}$. فبمَ إذن يتحقق البثّ في \nt{SERO}: بالانتشار أم بتفرّع سلك له $\sim10^5$ موقع تحرير؟
\end{exercise}

\begin{exercise}[\lvlA]
\label{ex:05:2}
\emph{منظِّم المستوى.} بيّن أنه في~\eqref{eq:seroneuron} مع $f(u)=au$ تساوي الحساسية النسبية $S=(a/r)\,\partial r/\partial a$ القيمةَ $1/(1+G)$ تمامًا، لا عند $G\gg1$ فقط. عند $G=9$ ازدادت الاستثارية $a$ بنسبة $20\,\%$. بكم في المئة يتغير $r$ من دون تقريب خطي؟
\end{exercise}

\begin{exercise}[\lvlB]
\label{ex:05:3}
\emph{مستقبل بطيء في الحلقة.} يستجيب مستقبل \nt{SERO} بتأخير: $r(t)=a\bigl(b-gc(t-T)\bigr)$، والحجم وفق~\eqref{eq:seroneuron}. بيّن أنه عند $G>1$ لا تكون الحلقة مستقرة إلا إذا كان $T<T^*=\tau_c\arccos(-1/G)/\sqrt{G^2-1}$. جد $T^*$ عند $\tau_c=1\,\text{s}$ لـ$G=2$ و$G=9$. ما كبت التشتت المتاح مع تأخير من مئات الميلي ثانية؟
\end{exercise}

\begin{exercise}[\lvlB]
\label{ex:05:4}
\emph{ضجيج الطلقات في المستوى.} كل نبضة تضيف إلى $c$ وفق~\eqref{eq:lowpass} دالة أسّية بـ$\tau_c=1\,\text{s}$. جد معامل الاختلاف لـ$c$ إذا كانت كل عصبونات \nt{SERO}، وعددها $3\cdot10^5$، تحرّر في حجم واحد نبضات بواسونية بمعدل $2$ في الثانية. أي $\tau_c$ يعطي $\mathrm{CV}=10\,\%$ لعصبون واحد تردده $4$ نبضات في الثانية؟
\end{exercise}

\begin{exercise}[\lvlB]
\label{ex:05:5}
\emph{نمذجة: التغذية الراجعة ضد الضجيج.} حاكِ $N=50$ ناظمة إيقاع بواسونية بترددات $r_i=a_i[b-gc]_+$، حيث $a_i$ موزعة بانتظام في $[2.8,\,5.2]$، والحجمَ~\eqref{eq:seroneuron} مع $\kappa=1/N$ و$\tau_c=1\,\text{s}$. كم مرة تخفّض التغذية الراجعة ($b=1$، $g=2.25$) معامل الاختلاف \foreignlanguage{english}{CV} للمستوى $c$ مقارنةً بـ$b=0.1$، $g=0$؟ وهل تسوّي ترددات العصبونات؟
\end{exercise}

\begin{exercise}[\lvlB]
\label{ex:05:6}
\emph{تصميم: \foreignlanguage{english}{XOR} بمنطق \nt{SERO}.} البوابة ناظمة إيقاع تعطي $r_0$ إذا كان $b-g\sum_kc_k>0$، و$0$ خلاف ذلك، في حجمها الخاص $\tau_c\,\dd c/\dd t=-c+\kappa r$؛ وهي تحسب \foreignlanguage{english}{NOR}. ابنِ \foreignlanguage{english}{XOR} من خمس بوابات كهذه، وجد زمن استقراره عند $\tau_c=1\,\text{s}$ و$G'=g\kappa r_0/b=2$.
\end{exercise}

\begin{exercise}[\lvlB]
\label{ex:05:7}
\emph{الصبر مع تأخير مجهول.} (أ)~يقبل الحيوان انتظار مكافأة أكبر بثلاث مرات إذا لم يتجاوز الانتظار $6\,\text{s}$. أي أفق $\tau_\gamma$ يعني هذا؟ (ب)~ليكن التأخير $D$ موزعًا أسّيًا بمتوسط $\bar D$. بيّن أن الانتظار مجدٍ عند $\bar D<\tau_\gamma(R/r_0-1)$، وقارن بالتأخير الثابت عند $\tau_\gamma=2\,\text{s}$ و$R=3r_0$.
\end{exercise}

\begin{exercise}[\lvlC]
\label{ex:05:8}
\emph{كيف يضبط التركيز الأفق.} يتلقى عصبون \nt{DOPA} المقدار $V$ مباشرة بوزن $k_1$، وعبر وسيط \nt{GABA} ينعّمه بـ$\tau_d=0.1\,\text{s}$ بوزن $-k_2$؛ وللمقادير $V$ البطيئة يساوي المجموع $(k_1-k_2)V+k_2\tau_d\dot V$. اختر $k_1$ و$k_2$ لتطابقا~\eqref{eq:ctd} مع $\tau_\gamma=2\,\text{s}$. يضرب \nt{SERO} المعامل $k_1$ في $m$: كم يجب أن يتغير $m$ لمضاعفة الأفق؟
\end{exercise}

\chapter{شبكات عصبية تستطيع أن تتعلم: نضيف \nt{DOPA}}
\chaptermark{شبكات تتعلم}
\label{sec:dofa}
\begin{chaptergoals}
\item لماذا لا يجوز للمشبك أن يستعمل إلا إشارات محلية، ولماذا يستبعد ذلك الانتشار العكسي؛
\item ماذا تتعلم قاعدة هيب: الاتجاه الذي تتغير فيه المداخل أكثر من غيره؛
\item كيف يحوّل أثر الأهلية مع إشارة البثّ $\delta$ الضجيجَ إلى خطوات صعودًا على تدرج المكافأة؛
\item كيف يعطي الناقد $\delta$ في الزمن المتصل، ولماذا يتعلم السلك المشترك الواحد ببطء.
\end{chaptergoals}

\section{قواعد اللعبة}

في كل دارات الفصول السابقة كان المهندس هو من يضع الأوزان. وفي الكائن الحي لا يوجد مهندس. يجب أن تستقر الأوزان وحدها، أثناء العمل، وبحيث تفعل الشبكة شيئًا مفيدًا. وهذا ما يسمى التعلّم.

إلى القواعد السابقة تضاف قاعدة أخرى. المشبك يغيّر وزنه بنفسه، ولا يستطيع أن يعرف إلا ما يجري \emph{بقربه}: نشاط المرسِل $x$، ونشاط المتلقي $y$، والمواد التي تصل إليه. لا أحد يستطيع أن يرسل إلى المشبك تصحيحًا خاصًا به.

وهذا يستبعد الطريقة الرئيسية في الشبكات العصبية الاصطناعية: الانتشار العكسي للخطأ. هناك يعبر خطأ الخرج الشبكة في الاتجاه المعاكس عبر الأوزان نفسها، في طور مستقل بعد التمرير الأمامي، ويتلقى كل وزن تصحيحه. ولهذا تلزم أسلاك عكسية تكرر الأمامية بدقة، ونبضة ساعة مشتركة. ولم يُعثر في الدماغ على أي منهما، على الأقل بالشكل الذي هما عليه في الشبكات الاصطناعية~\cite{Lillicrap2020,WhittingtonBogacz2019}.

سنكتب تغيّر الوزن عملية متصلة بطيئة:
\begin{equation}
\frac{\dd w}{\dd t} \;=\; \eta\,\Phi(\text{ما يعرفه المشبك}),
\label{eq:plasticity}
\end{equation}
حيث $\eta$ معدل التعلّم، وهو صغير إلى حد أن الوزن يكاد لا يتغير خلال نبضة واحدة. والفصل كله عن الشكل الذي يمكن أن تتخذه الدالة $\Phi$.

\section{أين يُخزَّن الوزن}
\label{sec:weightstore}

ما المشبك الذي ”يغيّر وزنه“؟ للوصلة جانبان: نهاية سلك المرسِل (الجانب \emph{قبل المشبكي}) وقطعة من غشاء المتلقي عليها المستقبلات (الجانب \emph{بعد المشبكي})، وبينهما شق ضيق. في وصلات \nt{GLUT} يقع غشاء المتلقي عادةً على \emph{شويكة}، وهي نتوء صغير على فرع الدخل عند المتلقي. من المناسب أن نفكك وزن الوصلة إلى ثلاثة عوامل~\cite{delCastillo1954}:
\begin{empheq}[box=\keybox]{equation}
w \;=\; N_s\,p\,A_1 .
\label{eq:quantal}
\end{empheq}
هنا $N_s$ عدد مواقع التحرير بين العصبونين، و$p$ احتمال أن تحرّر النبضة دفعة من الناقل في موقع واحد (وهي $p$ نفسها من الفصل~\ref{sec:code})، و$A_1$ استجابة المتلقي لدفعة واحدة، أي عدد المستقبلات التي تلتقطها. العامل $p$ يقيم عند المرسِل، و$A_1$ عند المتلقي، أما $N_s$ فيتطلب الجانبين: مواقع تحرير جديدة تنمو، وقديمة تختفي، ويستمر هذا طوال الحياة.

\emph{القرار للمتلقي.} في مشبك \nt{GLUT} النموذجي لا يمرّر أحد مستقبلات الغلوتامات التيار إلا إذا اجتمع شرطان: وصل الغلوتامات من المرسِل، وغشاء المتلقي مستثار أصلًا~\cite{Nowak1984}. هذه قاعدة هيب مدمجة في جزيئة واحدة، أي كاشف التزامن من الفصل~\ref{sec:glutcomputer} وقد وُضع على مدخل المشبك نفسه. وحين يعمل المستقبل يُدخل الكالسيوم إلى المتلقي، فيطلق الكالسيوم تغيّر الوزن. المتلقي هو الموضع الوحيد الذي يُرى فيه الدخل والخرج معًا، لذلك يُتخذ القرار عنده.

\emph{التنفيذ يمكن أن يتولاه أي جانب.} أكثر أنواع التقوية الطويلة الأمد دراسةً تجري عند المتلقي: تُدمج في الغشاء مستقبلات جديدة~\cite{Malinow2002}، وتكبر الشويكة~\cite{Matsuzaki2004}، أي يكبر $A_1$. لكن المتلقي يستطيع أن يغيّر $p$ أيضًا: فهو يحرّر مادة تعبر الشق عائدة إلى المرسِل، وهي القناة العكسية من الملحق~\ref{sec:othermediators}، فيبدأ المرسِل بتحرير ناقل أقل~\cite{WilsonNicoll2001}. أما التغيرات القصيرة الأمد في الوزن، أي الاستنزاف والتيسير من الفصل~\ref{sec:code}، فهي تغيرات في $p$ يديرها المرسِل وحده بحسب نشاطه القريب.

بالنسبة إلى المهندس، سجل الوزن مقسوم بين رقاقتين. قاعدة التعلّم ينفذها المتلقي، ويستطيع أن يكتب النتيجة عنده أو عند المرسِل عبر القناة العكسية. لذلك فكلمة ”محلي“ في قواعد هذا الفصل لا تعني جانبًا واحدًا من الوصلة، بل الوصلة كلها.

\section{التعلّم بلا معلّم}

أبسط ما يمكن أن يفعله المشبك أن يقوى حين يكون المرسِل والمتلقي نشطين معًا:
\begin{equation}
\frac{\dd w}{\dd t} \;=\; \eta\,x\,y .
\label{eq:hebb}
\end{equation}
هذه قاعدة هيب~\cite{Hebb1949}. إنها محلية ولا تحتاج إلى ساعة. لكن فيها العيب نفسه الذي في شبكة \nt{GLUT} في الفصل~\ref{sec:glutcomputer}: التغذية الراجعة الموجبة. الوزن القوي يكبّر $y$، و$y$ الكبير يقوّي الوزن أكثر، فتنمو الأوزان بلا حد. والعلاج تغذية راجعة سالبة على الوزن: ليتناقص الوزن أيضًا بما يتناسب مع $y^2$. لعصبون مداخله $\mathbf x$ وخرجه خطي $y=\mathbf w\cdot\mathbf x$ نحصل على
\begin{empheq}[box=\keybox]{equation}
\frac{\dd\mathbf w}{\dd t} \;=\; \eta\,y\,\bigl(\mathbf x - y\,\mathbf w\bigr) .
\label{eq:oja}
\end{empheq}
ما زالت القاعدة محلية: يعرف المشبك دخله $x_i$ والخرج $y$ ووزنه $w_i$.

\begin{keyblock}
\begin{proposition}[ماذا تجد قاعدة هيب]
\label{prop:oja}
تقود القاعدة~\eqref{eq:oja} متجه الأوزان إلى طول الوحدة على امتداد الاتجاه الذي تتغير فيه المداخل أكثر من غيره، أي إلى المتجه الذاتي الرئيسي لمصفوفة ارتباطات المداخل $C=\langle\mathbf x\mathbf x^{\mathsf T}\rangle$~\cite{Oja1982}. يتعلم العصبون أن يخرج أكثر مقدار واحد تعبيرًا يمكن أن تُضغط فيه مداخله.
\end{proposition}
\end{keyblock}

البرهان. نأخذ متوسط~\eqref{eq:oja} على المداخل: $\dd\mathbf w/\dd t=\eta\bigl(C\mathbf w-(\mathbf w^{\mathsf T}C\mathbf w)\,\mathbf w\bigr)$. النقاط الثابتة هي المتجهات الذاتية لـ$C$ ذات طول الوحدة. إذا كان $\mathbf w$ قريبًا من متجه ذاتي $\mathbf e_k$ قيمته الذاتية $\lambda_k$، فإن إضافة صغيرة على امتداد متجه آخر $\mathbf e_j$ تنمو مثل $e^{\eta(\lambda_j-\lambda_k)t}$. ولا تستقر إلا النقطة التي تكون عندها كل $\lambda_j<\lambda_k$، أي المتجه الرئيسي.

وثمة صيغة من قاعدة هيب تنظر إلى توقيت النبضات. إذا وصلت نبضة المرسِل \emph{قبل} نبضة المتلقي بـ$s$ ميلي ثانية، ازداد الوزن بمقدار $A_+e^{-s/\tau_+}$؛ وإذا وصلت \emph{بعدها} بـ$s$ ميلي ثانية، نقص بمقدار $A_-e^{-s/\tau_-}$، حيث $\tau_\pm$ من رتبة عشرين ميلي ثانية. هذه القاعدة مقيسة على مشابك عصبونات \nt{GLUT}~\cite{BiPoo1998}. إنها تقوّي الوصلات التي \emph{يمكن أن تكون سببًا} في الاستجابة، وتضعف المتأخرة (الشكل~\ref{fig:stdp}). مرّر عبر الشبكة سلسلة الإثارات نفسها مرات كثيرة، فتنمو فيها وحدها وصلات من كل حلقة إلى التي تليها: السلسلة من الفصل~\ref{sec:glutcomputer} وقد جُمعت بلا مهندس.

\begin{figure}[t]
\centering
\begin{tikzpicture}[>=Stealth,x=0.07cm,y=1.3cm]
\draw[->,gray!70!black] (-66,0) -- (68,0) node[right,black] {\small $s$};
\draw[->,gray!70!black] (0,-1.2) -- (0,1.35) node[above,black] {\small تغيّر الوزن};
\draw[very thick,blue!60!black] plot[domain=0.5:60,samples=50] (\x,{exp(-\x/20)});
\draw[very thick,red!70!black] plot[domain=-60:-0.5,samples=50] (\x,{-0.6*exp(\x/20)});
\draw[blue!60!black,thick] (0.5,0) -- (0.5,1);
\draw[red!70!black,thick] (-0.5,0) -- (-0.5,-0.6);
\foreach \x in {-40,-20,20,40}{\draw[gray!60!black] (\x,-0.04) -- (\x,0.04);}
\node[below,font=\scriptsize] at (20,-0.04) {$20\ms$};
\node[above,font=\scriptsize] at (-20,0.04) {$-20\ms$};
\node[align=left,font=\small,blue!60!black,anchor=west] at (22,0.95) {المرسِل أسبق:\\ الوصلة تقوى};
\node[align=right,font=\small,red!70!black,anchor=east] at (-12,-0.95) {المرسِل متأخر:\\ الوصلة تضعف};
\end{tikzpicture}
\caption{قاعدة هيب بحسب التوقيت. على المحور الأفقي $s=t_{\text{post}}-t_{\text{pre}}$، أي بكم تأخرت نبضة المتلقي عن نبضة المرسِل؛ $\tau_\pm=20\ms$، $A_-=0.6A_+$. عرض النافذة بضع عشرات من الميلي ثواني، أي من رتبة نافذة التزامن~\eqref{eq:window} نفسها.}
\label{fig:stdp}
\end{figure}

لكن كل قواعد هذا القسم تعلّم ما هو \emph{متكرر}، لا ما هو \emph{مفيد}. المشبك الذي لا يعرف إلا $x$ و$y$ لا يدري أجلب الفعلُ عصيرًا أم ألمًا. ولكي يتعلم المشبك من النتيجة يحتاج إلى إشارة ثالثة.

\section{العامل الثالث}

الإشارة الثالثة يأتي بها \nt{DOPA}: عصبونات الدوبامين تبثّ عددًا واحدًا، هو خطأ التنبؤ بالمكافأة $\delta$~\eqref{eq:rpe} (الفصل~\ref{sec:neurontypes}). ليكن تغيّر الوزن متناسبًا مع جداء ثلاثة عوامل: نشاط المرسِل، ونشاط المتلقي، و$\delta$.

الصعوبة أن المكافأة لا تأتي حين كانت الشبكة تختار الفعل، بل بعد مئات الميلي ثواني أو ثوانٍ. وحين تصل $\delta$ يكون الجداء $xy$ قد صار صفرًا منذ زمن، ويحتاج المشبك إلى ذاكرة تفيد أنه شارك. وأبسط ذاكرة هي دارة \foreignlanguage{english}{RC} التي نعرفها:
\begin{empheq}[box=\keybox]{equation}
\tau_e\,\frac{\dd e}{\dd t} \;=\; -e \;+\; x\,y ,
\qquad
\frac{\dd w}{\dd t} \;=\; \eta\,\delta(t)\,e(t) .
\label{eq:threefactor}
\end{empheq}
يسمى المقدار $e$ \emph{أثر الأهلية}~\cite{Fremaux2016}. يترك النشاط المشترك في المشبك علامة تذوب خلال الزمن $\tau_e$. إذا وصل الدوبامين خلال هذا الزمن تغيّر الوزن بإشارة $\delta$؛ وإلا اختفت العلامة بلا أثر (الشكل~\ref{fig:trace}). وقد قيس على مشابك \nt{GLUT} أن الدوبامين الذي يصل خلال ثانية أو ثانيتين بعد النشاط المشترك يحوّله إلى تقوية للوصلة، أما الذي يصل بعد ذلك فلا~\cite{Yagishita2014}. ووُجدت نوافذ لدونة طولها ثوانٍ أيضًا حيث لا تكون الإشارة الثالثة الدوبامين، بل الإثارة القوية للمتلقي نفسه~\cite{Bittner2017}.

\begin{figure}[t]
\centering
\begin{tikzpicture}[>=Stealth,x=7cm,y=1cm]
\fill[orange!12] (0.7,-0.2) rectangle (0.8,5.2);
% rows: x*y, delta, traces, weights
\foreach \y/\lab in {4.4/{$x\,y$},3.1/{$\delta$},1.4/{الأثر $e$},0/{الوزن $w$}}{
  \draw[gray!70!black] (0,\y) -- (1.42,\y);
  \node[left,font=\small] at (-0.02,\y+0.3) {\lab};}
\draw[very thick,blue!60!black] (0,4.4) -- (0,4.9) -- (0.2,4.9) -- (0.2,4.4);
\node[right,font=\scriptsize,blue!60!black] at (0.21,4.8) {المرسِل والمتلقي نشطان معًا};
\draw[very thick,orange!85!black] (0,3.1) -- (0.7,3.1) -- (0.7,3.7) -- (0.8,3.7) -- (0.8,3.1) -- (1.4,3.1);
\node[right,font=\scriptsize,orange!85!black] at (0.81,3.6) {وصل الدوبامين};
% traces, each in units of its own maximum
\draw[very thick,blue!60!black] plot[domain=0:0.2,samples=20] (\x,{1.4+1.1*(1-exp(-\x))/(1-exp(-0.2))})
  plot[domain=0.2:1.4,samples=60] (\x,{1.4+1.1*exp(-(\x-0.2))});
\draw[thick,densely dashed,gray!50!black] plot[domain=0:0.2,samples=30] (\x,{1.4+1.1*(1-exp(-\x/0.05))/(1-exp(-4))})
  plot[domain=0.2:1.4,samples=80] (\x,{1.4+1.1*exp(-(\x-0.2)/0.05)});
\node[right,font=\scriptsize,blue!60!black] at (1.0,2.15) {$\tau_e=1\,\text{s}$};
\node[right,font=\scriptsize,gray!40!black] at (0.3,1.65) {$\tau_e=50\ms$};
% weights
\draw[very thick,blue!60!black] (0,0.25) -- (0.7,0.25) -- (0.8,0.8) -- (1.4,0.8) node[right,font=\scriptsize] {ازداد الوزن عند $\tau_e=1\,\text{s}$};
\draw[thick,densely dashed,gray!50!black] (0,0.2) -- (1.4,0.2) node[right,font=\scriptsize,gray!40!black] {لا تغيير عند $\tau_e=50\ms$};
% time axis marks
\foreach \x/\l in {0/0,0.2/0.2,0.7/0.7,1.4/1.4}{\draw[gray!60!black] (\x,-0.05) -- (\x,0.05) node[below=3pt,font=\scriptsize] {$\l\,\text{s}$};}
\draw[<->,thick,gray!50!black] (0.2,5.35) -- node[above,font=\scriptsize,black] {تأخر المكافأة $0.5\,\text{s}$} (0.7,5.35);
\end{tikzpicture}
\caption{أثر الأهلية وفق القاعدة~\eqref{eq:threefactor}. يدوم النشاط المشترك $0.2\,\text{s}$، ويصل الدوبامين بعد نصف ثانية من انتهائه. الآثار معروضة كنسب من قيمتها العظمى. الأثر البطيء ($\tau_e=1\,\text{s}$) ما زال حيًا في تلك اللحظة، أما السريع ($\tau_e=50\ms$) فقد ذاب منذ زمن.}
\label{fig:trace}
\end{figure}

\begin{keyblock}
\begin{proposition}[التعلّم بإشارة بثّ]
\label{prop:threefactor}
لا تغيّر القاعدة~\eqref{eq:threefactor} إلا المشابك التي شاركت في عمل الشبكة خلال آخر $\tau_e$، وفي الاتجاه الذي تحدده الإشارة المشتركة $\delta$. إشارة البثّ مع الأثر المحلي تعطي تصحيحًا موجّهًا إلى عنوان. والتأخر بين الفعل والنتيجة مقبول ما دام لا يتجاوز $\tau_e$.
\end{proposition}
\end{keyblock}

\section{لماذا ينجح هذا: الضجيج بحثًا}
\label{sec:dofagradient}

لماذا تقود القاعدة~\eqref{eq:threefactor} إلى الأفضل؟ الجواب في الضجيج من الفصل~\ref{sec:code}. ليكن خرج العصبون $y=f(u)+\xi$، حيث $u=\sum_jw_jx_j$، و$\xi$ ارتعاش عشوائي متوسطه صفر وتباينه $\sigma^2$. المكافأة $R$ تعتمد على $y$. لنأخذ أثرًا لا يساوي $x_jy$ ببساطة، بل جزأه العشوائي $x_j\xi$، ولنأخذ عاملًا ثالثًا هو الخطأ $\delta=R-\bar R$، حيث $\bar R$ المكافأة المتوقعة. عند ضجيج صغير $R\approx R\bigl(f(u)\bigr)+R'\,\xi$، لذلك
\begin{empheq}[box=\keybox]{equation}
\bigl\langle \delta\,\xi\,x_j \bigr\rangle \;=\; \sigma^2\,R'\,x_j \;=\; \frac{\sigma^2}{f'(u)}\,\frac{\partial\langle R\rangle}{\partial w_j} .
\label{eq:perturbation}
\end{empheq}
الخطوة المتوسطة تسير على تدرج المكافأة المتوقعة~\cite{Williams1992}. لم يحسب أحد أي مشتقة: انحرفت الشبكة انحرافًا عشوائيًا، وأخبر الدوبامين هل صار الوضع أفضل، فتحركت المشابك المشاركة في الاتجاه المفيد (الشكل~\ref{fig:perturbation}). الضجيج الذي أعاق نقل الرسائل في الفصل~\ref{sec:code} يصير هنا بحثًا.

\begin{figure}[t]
\centering
\begin{tikzpicture}[>=Stealth,x=2cm,y=3cm]
\draw[->,gray!70!black] (0,0) -- (4.2,0) node[below left,font=\small,black] {الخرج $y$};
\draw[->,gray!70!black] (0,0) -- (0,1.2) node[left,font=\small,black] {المكافأة $R$};
% reward hill
\draw[very thick,blue!60!black] plot[domain=0:4,samples=60] (\x,{max(0,1-((\x-3)/2.2)^2)});
\node[font=\scriptsize,blue!60!black,anchor=south] at (3,1.02) {أفضل خرج};
% noise around f(u)
\draw[thick,gray!60!black] plot[domain=0.6:2.4,samples=50] (\x,{0.28*exp(-((\x-1.5)/0.3)^2/2)});
\draw[densely dashed,gray!60!black] (1.5,0) -- (1.5,0.535);
\node[below,font=\small] at (1.5,0) {$f(u)$};
\node[font=\scriptsize,gray!50!black] at (1.5,-0.2) {الارتعاش $\xi$};
% expected reward
\draw[densely dotted,gray!60!black] (0.5,0.535) node[left,font=\scriptsize,black] {$\bar R$} -- (2.1,0.535);
% two random deviations
\fill[green!45!black] (2.1,0.833) circle (0.035);
\draw[very thick,green!45!black] (2.1,0.535) -- (2.1,0.833);
\node[font=\scriptsize,green!45!black,anchor=west,align=left] at (2.18,0.6) {$\xi>0$: أفضل،\\ $\delta>0$};
\fill[red!70!black] (0.9,0.089) circle (0.035);
\draw[very thick,red!70!black] (0.9,0.535) -- (0.9,0.089);
\node[font=\scriptsize,red!70!black,anchor=east,align=right] at (0.83,0.27) {$\xi<0$: أسوأ،\\ $\delta<0$};
% resulting step
\draw[->,very thick,orange!85!black] (1.55,0.4) -- (1.95,0.4) node[right,font=\scriptsize,black] {خطوة};
\end{tikzpicture}
\caption{الضجيج بحثًا~\eqref{eq:perturbation}. خرج العصبون يرتعش حول $f(u)$. الانحراف العشوائي إلى اليمين (الأخضر) أعطى أكثر من المتوقع، $\delta>0$؛ والانحراف إلى اليسار (الأحمر) أعطى أقل، $\delta<0$. في الحالتين الجداء $\delta\,\xi$ موجب، فيتحرك الوزن إلى اليمين، حيث تزداد المكافأة. في المتوسط تتناسب الخطوة مع الميل $R'$: على قمة التلّ يكون الانحراف في الاتجاهين سيئًا بالقدر نفسه، فتلغي الخطوات بعضها بعضًا.}
\label{fig:perturbation}
\end{figure}

\begin{keyblock}
\begin{proposition}[النزول على التدرج بلا أسلاك عكسية]
\label{prop:gradient}
إذا كانت في الشبكة انحرافات عشوائية في النشاط، وكانت إشارة البثّ تساوي خطأ التنبؤ بالمكافأة ومتوسطها صفر في كل موقف، فإن القاعدة~\eqref{eq:threefactor} تغيّر الأوزان في المتوسط على تدرج المكافأة المتوقعة. ولا حاجة لذلك إلى أسلاك عكسية ولا إلى طور تعلّم مستقل.
\end{proposition}
\end{keyblock}

الآن يتضح لماذا يُخبر الدوبامين لا عن المكافأة، بل عن \emph{خطأ} التنبؤ بها. لو كان العامل الثالث المكافأةَ نفسها $R\ge0$ لما تغيّر المتوسط في~\eqref{eq:perturbation}، لكن لظهرت مشكلتان. الأولى أن الجزء المفيد سيُضاف إليه الحد $\bar R\,\xi x_j$: متوسطه صفر، لكنه ضجيج قوي. والثانية، وهي أسوأ، أن المشبك الحقيقي لا يستطيع أن يفصل في الأثر $x_jy$ الجزء العشوائي عن غير العشوائي. عند مداخل معينة يكون $x_jf(u)$ العدد نفسه من محاولة إلى أخرى؛ فإذا كان متوسط $\delta$ في هذا الموقف صفرًا فإن هذا الجزء من الأثر لا يغيّر شيئًا، أما عند $\delta\ge0$ فإنه يقوّي مرة بعد مرة كل ما فعلته الشبكة، صوابًا كان أم خطأً. لذلك يجب أن يكون $\bar R$ التوقعَ \emph{في الموقف المعين}، لا المتوسط على الحياة كلها. وحسابه مهمة الناقد الذي نتحدث عنه أدناه.

تحققنا من ذلك على نموذج. مجموعتان من عصبونات \nt{GLUT} تختاران أحد فعلين عبر التثبيط المتبادل، كما في الشكل~\ref{fig:wta}؛ الأوزان من إشارة ”ابدأ“ إلى المجموعتين متساوية في البداية، والضجيج يقرر من يفوز. الفعل $A$ يجلب العصير باحتمال $0.8$، والفعل $B$ باحتمال $0.2$، ويصل العصير بعد نصف ثانية من الاختيار. عند $\tau_e=1\,\text{s}$ و$\delta=R-\bar R$ ارتفعت نسبة اختيار $A$ على خمسمئة محاولة من $0.65\text{--}0.8$ في المئة الأولى إلى $0.96\text{--}1.0$ في الأخيرة، وبقيت الأوزان ضمن $0.85\text{--}1.35$. وعند $\tau_e=50\ms$، وهو أقل من تأخر المكافأة، لم تتعلم الشبكة شيئًا: بقيت نسبة اختيار $A$ نحو النصف. وعند $\delta=R$ دون طرح التوقع صارت الشبكة تختار $A$ أيضًا، لكن وزن الفائز ازداد خلال الخمسمئة محاولة نفسها ألف مرة واستمر في الازدياد؛ ومع $\delta=R$ نفسها ومعدل تعلّم أكبر عشر مرات ثبّتت الشبكة في واحدة من ثلاث محاكاة اختيارها العشوائي الأول إلى الأبد، وكان الأسوأ.

\section{ثمن السلك الواحد}

الإشارة $\delta$ واحدة للجميع، أما الضجيج الذي تقيّمه فهو مجموع ارتعاشات كل العصبونات $N$ المؤثرة في النتيجة. يرى كل مشبك في $\delta$ جزأه المفيد وضجيج الجيران الآخرين $N-1$. ولكي يستخلص نصيبه عليه أن يأخذ المتوسط على محاولات كثيرة، فيزداد زمن التعلّم متناسبًا مع $N$~\cite{Werfel2005}. أما الانتشار العكسي فلا يتطلب هذا الثمن.

\begin{keyblock}
\begin{proposition}[ثمن البثّ]
\label{prop:broadcastcost}
زمن التعلّم بإشارة مشتركة واحدة يزداد متناسبًا مع عدد العصبونات التي يؤثر ضجيجها في النتيجة. هذا التعلّم يصلح لدارات صغيرة تختار من بين خيارات قليلة، لا لضبط مليارات الأوزان.
\end{proposition}
\end{keyblock}

يبدو أن الدماغ يقسّم العمل هكذا بالضبط: تذهب مخارج عصبونات \nt{DOPA} قبل كل شيء إلى المناطق التي تختار الأفعال (الفصل~\ref{sec:neurontypes})، أما شبكة \nt{GLUT} الهائلة في القشرة، حيث الغالبية الساحقة من الأوزان، فتتعلم أساسًا بقواعد محلية، بلا معلّم خارجي. كانت هذه القواعد تُردّ سابقًا إلى قواعد هيب~\eqref{eq:oja}. أما الآن فتتزايد المعطيات على أن للقشرة هدفًا خاصًا بها، هو أن \emph{تتنبأ} بدخلها التالي، وعندئذ يُحسب الخطأ في مكانه، فرقًا بين المتنبَّأ به والواصل~\cite{Keller2018}: المعلّم مدمج في الشبكة نفسها. ونوافذ اللدونة التي طولها ثوانٍ، حيث تكون الإشارة الثالثة الإثارةَ القوية للمتلقي نفسه~\cite{Bittner2017}، تبيّن أن إشارة خطأ محلية تصل إلى المشابك فعلًا. أما إلى أي حد هذه قاعدة عامة في القشرة ففرضية.

\section{من أين يأتي الخطأ: الناقد في الزمن المتصل}

يبقى السؤال: كيف تحسب عصبونات \nt{DOPA} نفسها $\delta$؟ في برامج التعلّم المعزَّز يُحسب خطأ التنبؤ بخطوات $t,t+1,\dots$ وفي الكائن الحي لا توجد خطوات، لذلك نكتبه في الزمن المتصل~\cite{Doya2000}. ليكن $r(t)$ دفق المكافأة، و$V(t)$ توقع كل المكافأة المقبلة، حيث تُقدَّر المكافأة الأبعد بقيمة أقل:
\begin{equation}
V(t) \;=\; \int_t^{\infty} e^{-(s-t)/\tau_\gamma}\,r(s)\,\dd s .
\end{equation}
بالاشتقاق نحصل على $\dd V/\dd t=V/\tau_\gamma-r$. والفرق الذي يخالف به الواقعُ هذه المساواة هو خطأ التنبؤ:
\begin{empheq}[box=\keybox]{equation}
\delta(t) \;=\; r(t) \;+\; \frac{\dd V}{\dd t} \;-\; \frac{V}{\tau_\gamma} .
\label{eq:ctd}
\end{empheq}
الثابت $\tau_\gamma$ هو أفق التخطيط، أي النظير المتصل للمعامل $\gamma$؛ ووفق فرضيةٍ يضبطه \nt{SERO} (القسم~\ref{sec:horizon})، وعندئذ تكون~\eqref{eq:patience} قاعدةً لمقدار ما يستحق الانتظار. لنتحقق من ثلاث حالات (الشكل~\ref{fig:ctdcases}). أضاء مصباح يعد بالعصير: قفز $V$، و$\dd V/\dd t$ كبير، فتظهر ذروة. وصل العصير الموعود: $r>0$، لكن التوقع نفد في اللحظة نفسها، $\dd V/\dd t\approx-r$، فلا ذروة. لم يصل العصير: اختفى التوقع بلا مكافأة، $\dd V/\dd t<0$، فيحدث هبوط.

\begin{figure}[t]
\centering
\begin{tikzpicture}[>=Stealth,x=1.15cm,y=1cm]
\foreach \col/\ttl/\lampon/\juice in {0/{عصير غير متوقع}/0/1, 1/{عصير موعود}/1/1, 2/{لم يصل العصير}/1/0}{
 \begin{scope}[xshift=\col*4.6cm]
  \node[font=\small] at (1.5,3.55) {\ttl};
  \foreach \yy in {2.4,1.2,0}{\draw[gray!60!black] (0,\yy) -- (3.1,\yy);}
  % reward r
  \ifnum\juice=1
    \fill[green!45!black] (1.95,2.4) rectangle (2.05,2.85);
    \pic[scale=0.55] at (2.0,3.1) {juice};
  \else
    \pic[scale=0.55] at (2.0,3.1) {nojuice};
  \fi
  % expectation V and error delta
  \ifnum\lampon=1
    \pic[scale=0.55] at (1.0,3.1) {lamp};
    \draw[very thick,blue!60!black] (0,1.2) -- (1,1.2) -- (1,1.45)
      plot[domain=1:2,samples=20] (\x,{1.2+0.25*exp((\x-1)/1.5)}) -- (2,{1.2+0.25*exp(1/1.5)}) -- (2,1.2) -- (3.1,1.2);
    \draw[very thick,red!70!black] (0,0) -- (0.95,0) -- (0.95,0.5) -- (1.05,0.5) -- (1.05,0) -- (3.1,0);
    \ifnum\juice=0
      \draw[very thick,red!70!black] (1.95,0) -- (1.95,-0.45) -- (2.05,-0.45) -- (2.05,0);
      \node[font=\scriptsize,red!70!black,anchor=west] at (2.1,-0.35) {هبوط};
    \else
      \node[font=\scriptsize,gray!50!black,anchor=north,align=center] at (2.0,-0.08) {$r$ وانخفاض $V$\\ يلغي أحدهما الآخر};
    \fi
  \else
    \draw[very thick,blue!60!black] (0,1.2) -- (3.1,1.2);
    \draw[very thick,red!70!black] (0,0) -- (1.95,0) -- (1.95,0.5) -- (2.05,0.5) -- (2.05,0) -- (3.1,0);
  \fi
  \ifnum\col=0
    \node[font=\small,anchor=east] at (-0.1,2.55) {$r$};
    \node[font=\small,anchor=east] at (-0.1,1.35) {$V$};
    \node[font=\small,anchor=east] at (-0.1,0.1) {$\delta$};
  \fi
 \end{scope}}
\end{tikzpicture}
\caption{ثلاث حالات للخطأ~\eqref{eq:ctd}؛ من الأعلى إلى الأسفل: المكافأة $r$، والتوقع $V$، والخطأ $\delta$. على اليسار لا يوجد توقع، فالعصير مفاجأة كله. في الوسط يرفع المصباح $V$ قفزةً واحدة: هذه قفزة في $\dd V/\dd t$، فتقع ذروة $\delta$ عند المصباح. بعد ذلك ينمو $V$ تمامًا كما يقتضي الأفق، و$\delta=0$؛ ولحظة العصير تلغي المكافأة $r$ وانخفاض $V$ أحدهما الآخر. على اليمين ينخفض $V$ بلا مكافأة، فيحدث هبوط. هذه هي الذروة والانتقال والهبوط أنفسها التي في الشكل~\ref{fig:rpe}، لكننا نرى الآن من أين تأتي.}
\label{fig:ctdcases}
\end{figure}

ماذا يلزم لـ~\eqref{eq:ctd} من الدارات التي لدينا؟ ثلاثة أشياء.

\emph{التقدير $V$ نفسه.} تخرجه مجموعة من عصبونات \nt{GLUT}، هي الناقد، تتعلم أوزانها بالقاعدة نفسها~\eqref{eq:threefactor} وبـ$\delta$ نفسها: إذا كان $\delta>0$ فالتوقع كان أقل من اللازم، فتقوى الوصلات التي كانت نشطة قبل ذلك.

\emph{المشتقة $\dd V/\dd t$.} تحسبها أي دارة تستجيب للتغيرات لا للمستوى: إثارة مباشرة مع تثبيط متأخر عبر وسيط \nt{GABA}، كما في كاشف الاتجاه في الفصل~\ref{sec:gaba}، أو مشبك مستنزَف~\eqref{eq:depression} من الفصل~\ref{sec:code}.

\emph{الإشارة في سلك واحد.} الخطأ قد يكون موجبًا أو سالبًا، أما تردد النبضات فموجب فقط. والحل هو نفسه عند عصبون \nt{SERO} في الفصل~\ref{sec:serocomputer}: عصبون \nt{DOPA} ناظمة إيقاع. نبضاته المنتظمة القليلة في الثانية تعني $\delta=0$، والرشقة تعني $\delta>0$، والتوقف يعني $\delta<0$. ويبقى للخطأ السالب مجال أضيق: لا يمكن خفض التردد تحت الصفر. وعند عصبونات \nt{DOPA} الحقيقية تكون الهبوطات فعلًا أضحل من الذرى~\cite{Bayer2005}.

أن الدوبامين يسلك سلوك $\delta$ أمر مقيس. أما كيف يحسب الدماغ $\dd V/\dd t$ بالضبط ففرضية.

\section{الفاعل والناقد}

لنجمع كل شيء (الشكل~\ref{fig:actorcritic}). عصبونات الموقف تثير مجموعات الأفعال التي تختار الفائز عبر \nt{GABA} (القضية~\ref{prop:wta})، وهذا هو \emph{الفاعل}، وتثير كذلك الناقد الذي يخرج $V$~\cite{SuttonBarto2018}. يتلقى عصبون \nt{DOPA} المكافأة $r$ و$V$، ويحسب~\eqref{eq:ctd}، ويبثّ $\delta$ إلى كل مشابك الفاعل والناقد.

\begin{figure}[t]
\centering
\begin{tikzpicture}[>=Stealth,
  nrn/.style={circle,draw,very thick,fill=blue!6,minimum size=0.9cm,inner sep=0pt,font=\small},
  gab/.style={circle,draw,very thick,fill=red!10,minimum size=0.6cm,inner sep=0pt},
  dofa/.style={circle,draw,very thick,double,double distance=1.2pt,fill=orange!15,minimum size=1.35cm,inner sep=0pt,font=\scriptsize},
  inh/.style={-{Bar[width=7pt]},thick,red!70!black},
  bc/.style={->,thick,densely dashed,orange!85!black}]
\node[nrn] (S1) at (0,2.4) {$x_1$};
\node[nrn] (S2) at (0,1.2) {$x_2$};
\node[nrn] (S3) at (0,0) {$x_3$};
\node[left=0.15cm,align=right,font=\small] at (-0.5,1.2) {الموقف};
\node[nrn] (A) at (4,2.6) {$A$};
\node[nrn] (B) at (4,1.0) {$B$};
\node[gab] (G) at (5.2,1.8) {};
\draw[->,thick] (A) -- (G); \draw[->,thick] (B) -- (G);
\draw[inh] (G) to[bend left=35] (A);
\draw[inh] (G) to[bend right=35] (B);
\node[above,font=\small] at (4.6,3.05) {الفاعل: اختيار الفعل};
\node[nrn] (V) at (4,-1.1) {$V$};
\node[below,font=\small] at (4,-1.6) {الناقد};
\foreach \s in {S1,S2,S3}{\foreach \t in {A,B,V}{\draw[->,gray!60!black] (\s) -- (\t);}}
\node[dofa] (D) at (8.0,-1.1) {\nt{DOPA}};
\draw[->,thick] (V) -- node[above,font=\scriptsize] {$V,\ \dd V/\dd t$} (D);
\draw[->,thick,green!45!black] (10.0,-1.1) node[right,black,font=\small] {المكافأة $r$} -- (D);
\pic[scale=1.1] at (10.6,-0.5) {juice};
\draw[bc] (D.north) -- (8.0,4.0) -- node[above,font=\small,black] {$\delta$ إلى كل مشابك الفاعل والناقد} (2.2,4.0) -- (2.2,3.0);
\draw[bc] (D.south) -- (8.0,-2.4) -- (2.2,-2.4) -- (2.2,-0.6);
\end{tikzpicture}
\caption{شبكة تتعلم اختيار الأفعال. الأسهم الرمادية مشابك \nt{GLUT} بأوزان متغيرة؛ الدائرة الحمراء عصبون \nt{GABA}؛ الدائرة المزدوجة ناظمة إيقاع \nt{DOPA} تحسب $\delta$ وفق~\eqref{eq:ctd}؛ الأسهم المتقطعة بثّ $\delta$.}
\label{fig:actorcritic}
\end{figure}

إشارة أثر الناقل يختارها المتلقي (القضية~\ref{prop:receptor})، وفي منطقة اختيار الأفعال يحمل جزء من العصبونات مستقبلات دوبامين من عائلة، وجزء آخر من عائلة أخرى. في تجارب على الشرائح يقوّي الدوبامين مع النشاط المشترك المشابكَ عند الأولى، ويضعفها عند الثانية~\cite{Shen2008}؛ وفي النموذج يعني هذا أن إشارة $\delta$ في~\eqref{eq:threefactor} معكوسة عند الثانية. الأولى تتعلم ما \emph{يستحق الفعل}، والثانية ما \emph{لا يستحق الفعل}.

\section{الخلاصة}

\begin{table}[t]
\centering
\footnotesize
\setlength{\tabcolsep}{4pt}
\renewcommand{\arraystretch}{1.15}
\begin{tabular}{>{\raggedright}p{2.6cm}>{\raggedright}p{2.7cm}>{\raggedright}p{3.1cm}>{\raggedright\arraybackslash}p{5.0cm}}
\toprule
القاعدة & ما يعرفه المشبك & ما تتعلمه الشبكة & ما هو معروف \\
\midrule
هيب بالترددات~\eqref{eq:oja} & $x$، $y$، وزنه & ضغط المداخل: الاتجاهات الرئيسية & التقوية عند النشاط المشترك مقيسة؛ وآلية تقييد الأوزان موضوع بحث \\
هيب بالتوقيت & ترتيب نبضات $x$ و$y$ ضمن $\sim20\ms$ & الروابط السببية، والتسلسلات & مقيسة على مشابك \nt{GLUT} \\
العوامل الثلاثة~\eqref{eq:threefactor} & $x$، $y$، الأثر $e$، الإشارة المشتركة $\delta$ & الأفعال التي تجلب المكافأة & تأثير الدوبامين خلال ثوانٍ بعد النشاط مقيس؛ وكونها الآلية الرئيسية لتعلّم الاختيار فرضية قوية الأساس \\
الناقد~\eqref{eq:ctd} & كذلك & توقع المكافأة $V$ & تشابه الدوبامين مع $\delta$ مقيس؛ والدارة التي تحسب $\dd V/\dd t$ فرضية \\
الانتشار العكسي للخطأ & تصحيحًا خاصًا لكل وزن & كل شيء، لكنه يحتاج إلى أسلاك عكسية وساعة & لم يوجد في الدماغ بهذا الشكل؛ ويُبحث عن تقريبات له \\
\bottomrule
\end{tabular}
\caption{قواعد التعلّم في شبكة من عصبونات \nt{GLUT} و\nt{GABA} و\nt{DOPA}.}
\label{tab:learning}
\end{table}

قواعد التعلّم مجموعة في الجدول~\ref{tab:learning}. \nt{GLUT} يحمل البيانات، و\nt{GABA} يدير التدفق، و\nt{DOPA} يقرر أي التغييرات في الشبكة تبقى.

\section{تمارين}

\begin{exercise}[\lvlA]
\label{ex:06:1}
\emph{كم يعيش الأثر.} يدوم النشاط المشترك $xy=1$ مدة $T_a=0.2\,\text{s}$ بدءًا من $e=0$، ويصل الدوبامين بعد $d=0.5\,\text{s}$ من نهايته. (أ)~كم مرة يكون الأثر~\eqref{eq:threefactor} في تلك اللحظة أكبر عند $\tau_e=1\,\text{s}$ منه عند $\tau_e=50\ms$؟ (ب)~عند أي $\tau_e$ يكون أكبر ما يمكن؟
\end{exercise}

\begin{exercise}[\lvlA]
\label{ex:06:2}
\emph{انجراف قاعدة هيب.} المرسِل والمتلقي يخرجان دفقين بواسونيين مستقلين بمعدل $10$ نبضات في الثانية، وكل زوج من النبضات يغيّر الوزن وفق نافذة الشكل~\ref{fig:stdp} مع $\tau_\pm=20\ms$ و$A_-=0.6A_+$. بكم $A_+$ يزداد الوزن خلال ساعة من نشاط عشوائي تمامًا؟ عند أي $\tau_-$ يختفي الانجراف؟
\end{exercise}

\begin{exercise}[\lvlB]
\label{ex:06:3}
\emph{ارتباطات لا تغايرات.} قاعدة القضية~\ref{prop:oja} تجد المتجه الذاتي الرئيسي لـ$C=\langle\mathbf x\mathbf x^{\mathsf T}\rangle$. الترددات موجبة: $\mathbf x=\mathbf m+\mathbf n$، $\mathbf m=(\mu,0)$، وتغاير الضجيج $\operatorname{diag}(s_1^2,s_2^2)$. بأي شرط يتعلم العصبون $\mathbf e_1$؟ وماذا يتعلم عند $\mu=1$، $s_1=0.1$، $s_2=0.5$؟
\end{exercise}

\begin{exercise}[\lvlA]
\label{ex:06:4}
\emph{الأفق في الذروة.} يضيء مصباح في لحظة لا يمكن التنبؤ بها، ويصل عصير حجمه $R$ بعد $L=1\,\text{s}$. بيّن أنه للتوقع المتعلَّم وفق~\eqref{eq:ctd} يكون $\delta\equiv0$ بين المصباح والعصير، وجد مساحة الذروة عند المصباح عند $\tau_\gamma=2\,\text{s}$ و$\tau_\gamma=4\,\text{s}$.
\end{exercise}

\begin{exercise}[\lvlB]
\label{ex:06:5}
\emph{من أين يأتي ثمن البثّ.} $N$ عصبونًا ترتعش باستقلال، $\xi_i\sim\mathcal N(0,\sigma^2)$، والخطأ $\delta=a\sum_i\xi_i$، ويقدّر المشبك $j$ التدرج بـ$\delta\,\xi_jx_j$. جد نسبة الإشارة إلى الضجيج في محاولة واحدة. عصبون واحد يتعلم في $100$ محاولة مدة كل منها $5\,\text{s}$: كم يستغرق التعلّم عند $N=10^4$ و$N=10^{10}$؟
\end{exercise}

\begin{exercise}[\lvlB]
\label{ex:06:6}
\emph{تصميم: دارة تحسب $\delta$.} يتلقى عصبون \nt{DOPA} المقدار $r$، و$V$ بوزن $k_1$، ونسخة من $V$ منعَّمة بـ$\tau_d$ بوزن $-k_2$. (أ)~اختر $k_1$ و$k_2$ بحيث يكون الخرج للمقادير $V$ البطيئة هو~\eqref{eq:ctd}. (ب)~عند $V=V_0e^{t/\tau_\gamma}$ تكون $\delta$ الحقيقية صفرًا. عند أي $\tau_d$ لا يتجاوز الخطأ الزائف للدارة $5\,\%$ من $V/\tau_\gamma$ إذا كان $\tau_\gamma=2\,\text{s}$؟
\end{exercise}

\begin{exercise}[\lvlB]
\label{ex:06:7}
\emph{نمذجة: أقصى تأخر للمكافأة.} أضف إلى \texttt{run()} في نسخة من \texttt{models/\allowbreak dofa\_learning.py} تأخرًا للمكافأة $d$. جد نسبة اختيار $A$ في المئة الأخيرة من $500$ محاولة عند $d=0.5\,\text{s}$ و$\tau_e=0.05$ و$0.2$ و$1$ و$4\,\text{s}$، وعند $\tau_e=1\,\text{s}$ و$d=3\,\text{s}$. عند أي نسبة من الأثر الباقي حتى المكافأة (التمرين~\ref{ex:06:1}) يختفي التعلّم؟
\end{exercise}

\begin{exercise}[\lvlC]
\label{ex:06:8}
\emph{هل يمكن تفادي ثمن البثّ؟} $N$ عصبونًا: $y_i=w_i+\xi_i$، $\xi_i\sim\mathcal N(0,\sigma^2)$، والمكافأة $R=-\sum_i(y_i-w_i^*)^2$، والقاعدة $\Delta w_i=\eta\,(R-\langle R\rangle)\,\xi_i$. في كم محاولة، عند أفضل $\eta$، يتناقص الخطأ $\sum_i(w_i-w_i^*)^2$ بمعامل $e$؟ وماذا لو ارتعشت $m$ عصبونات عشوائية فقط من $N$؟ هل يساعد البحث المتناثر؟
\end{exercise}

\chapter{نظريات حديثة بديلة للتعلّم القائم على الدوبامين}
\chaptermark{نظريات أخرى في \nt{DOPA}}
\label{sec:dofaalt}
\begin{chaptergoals}
\item كيف تسجل عصبونات \nt{DOPA} ذات الاستجابات اللامتناظرة للأخطاء توزيعَ المكافأة كله؛
\item كيف تنقل بعض عصبونات \nt{DOPA} الأهمية أو الخطر بدل علامة الخطأ؛
\item كيف يحمل سلك واحد خطأً سريعًا يعلّم ومستوى بطيئًا يحدد سرعة الأفعال؛
\item لماذا إشارة \nt{DOPA} حزمة لا عدد واحد، وأي نظرية تنافس الخطأ نفسه.
\end{chaptergoals}

\section{ما الذي يسري في السلك حقًا}

في الفصل~\ref{sec:dofa} كان نظام \nt{DOPA} عندنا سلكًا واحدًا يسري فيه عدد واحد، هو خطأ التنبؤ بالمكافأة $\delta$~\eqref{eq:ctd}. هذه الصورة تفسر الكثير، لكن كلما قيس الدوبامين بدقة أكبر وُجد ما لا ينسجم معها. بعض العصبونات يستجيب للمزعج كما يستجيب للمكافأة؛ وتركيز الدوبامين يرتفع ما دام الحيوان يركض نحو الهدف، مع أنه لا يحدث شيء غير متوقع؛ وعصبونات \nt{DOPA} المختلفة تسلك سلوكًا مختلفًا.

بالنسبة إلى المهندس، كل هذه المكتشفات سؤال واحد: \emph{ما صيغة الإشارة على ناقل البثّ؟} عدد واحد أم متجه؟ إشارة واحدة في السلك أم عدة إشارات موزعة على الترددات؟ هل تتحدث عن المستقبل، مثل $\delta$، أم عن الماضي؟ فيما يلي ستة أجوبة حديثة؛ ولكل منها سنفصل المقيس عن المفترض.

\section{لا عدد واحد بل توزيع}

الخطأ~\eqref{eq:ctd} يعلّم الناقد \emph{متوسط} توقع المكافأة. لكن نصف قطرة عصير مضمونة وقطرة كاملة باحتمال النصف لهما المتوسط نفسه. ولكي نميّز بينهما نحتاج إلى توزيع المكافأة كله.

لنأخذ أبسط مسألة: بعد الإشارة المنبئة تأتي مكافأة حجمها عشوائي $r$. لتكن عصبونات \nt{DOPA} كثيرة، ولكلٍّ منها، وليكن العصبون $i$، تقديره الخاص $V_i$، بحيث يكون خطؤه لحظة المكافأة $\delta_i=r-V_i$. وليأخذ كلٌّ منها الخطأ الموجب بوزن $\alpha_i^+$، والسالب بوزن $\alpha_i^-$. بعد كل مكافأة يتحرك التقدير بمقدار
\begin{equation}
\Delta V_i \;=\; \eta\,\bigl(\alpha_i^+\,[\delta_i]_+ \;-\; \alpha_i^-\,[-\delta_i]_+\bigr).
\label{eq:asymmetric}
\end{equation}
يكفّ التقدير عن التغير حين تتوازن التصحيحات الموجبة والسالبة في المتوسط:
\begin{empheq}[box=\keybox]{equation}
\alpha_i^+\,\bigl\langle[r-V_i]_+\bigr\rangle \;=\; \alpha_i^-\,\bigl\langle[V_i-r]_+\bigr\rangle ,
\qquad
\kappa_i \;=\; \frac{\alpha_i^+}{\alpha_i^++\alpha_i^-} .
\label{eq:expectile}
\end{empheq}
عند $\kappa_i=1/2$ هذا هو المتوسط العادي. وعند $\kappa_i>1/2$ يكون العصبون ”متفائلًا“: تقديره مزاح نحو الطرف الأعلى من التوزيع، وعند $\kappa_i<1/2$ يكون ”متشائمًا“.

مثال: يأتي العصير باحتمال $1/2$، وحجم القطرة $1$. عندئذ تعطي~\eqref{eq:expectile} $\kappa_i\cdot\tfrac12(1-V_i)=(1-\kappa_i)\cdot\tfrac12V_i$، أي $V_i=\kappa_i$ (الشكل~\ref{fig:expectiles}). مجموعة عصبونات بقيم $\kappa_i$ مختلفة تسجّل توزيع المكافأة كما تسجّل مجموعة من المئينات توزيعًا في الإحصاء.

\begin{figure}[t]
\centering
\begin{tikzpicture}[>=Stealth,x=7cm,y=3cm]
\draw[->,gray!70!black] (-0.08,0) -- (1.15,0) node[right,black] {\small المكافأة};
\draw[->,gray!70!black] (-0.08,0) -- (-0.08,0.75) node[left,black] {\small الاحتمال};
\fill[blue!25] (-0.03,0) rectangle (0.03,0.5);
\fill[blue!25] (0.97,0) rectangle (1.03,0.5);
\node[below,font=\small] at (0,0) {$0$};
\node[below,font=\small] at (1,0) {$1$};
\node[left,font=\scriptsize] at (-0.08,0.5) {$1/2$};
\foreach \k/\c/\lab/\h in {0.2/blue!60!black/{متشائم، $\kappa=0.2$}/0.62, 0.5/gray!50!black/{المتوسط، $\kappa=0.5$}/0.42, 0.8/red!70!black/{متفائل، $\kappa=0.8$}/0.22}{
  \draw[very thick,\c] (\k,0) -- (\k,\h);
  \node[above,font=\scriptsize,\c] at (\k,\h) {\lab};}
\end{tikzpicture}
\caption{توزيع المكافأة مسجَّلًا بمجموعة من عصبونات \nt{DOPA}. المكافأة $0$ أو $1$ باحتمال $1/2$ (العمودان الأزرقان)؛ والعصبون ذو النسبة $\kappa$ يصل إلى التقدير $V=\kappa$~\eqref{eq:expectile}.}
\label{fig:expectiles}
\end{figure}

المقيس: لعصبونات \nt{DOPA} المختلفة ”نقاط انقلاب“ مختلفة، أي حجم المكافأة الذي تتحول عنده الاستجابة من ذروة إلى هبوط، والعصبونات ذات اللاتناظر الأكبر بين استجابتيها للخطأ الموجب والسالب تنقلب عند مكافآت أكبر، كما تقتضي~\eqref{eq:expectile}~\cite{Dabney2020}. ومن استجابات مجموعة من العصبونات يمكن استعادة شكل التوزيع. أما أن هذا هو الوظيفة الرئيسية للتشتت، لا أثرًا جانبيًا، ففرضية.

\begin{keyblock}
\begin{proposition}[التوزيع من اللاتناظر]
\label{prop:expectile}
إذا اختلفت عصبونات \nt{DOPA} في النسبة $\kappa_i$ التي تأخذ بها الأخطاء الموجبة، تقاربت تقديراتها إلى نقاط مختلفة~\eqref{eq:expectile} من توزيع المكافأة. فتنقل مجموعة العصبونات لا المتوسط، بل التوزيع، ومعه المخاطرة.
\end{proposition}
\end{keyblock}

\section{لا الخطأ وحده بل الأهمية أيضًا}

وفق صيغة الخطأ~\eqref{eq:ctd} يجب أن يسبب الحدث المزعج، كصدمة أو طعم مرّ، هبوطًا: جاءت النتيجة أسوأ من المتوقع. جزء من عصبونات \nt{DOPA} يفعل ذلك فعلًا، وجزء آخر يستجيب للمزعج \emph{بذروة} كما يستجيب للمكافأة~\cite{Matsumoto2009}. هذه العصبونات لا تخبر بإشارة الخطأ، بل بأن شيئًا \emph{مهمًا} حدث: غير متوقع، قويًا، يستدعي الانتباه~\cite{BrombergMartin2010}.

من المناسب للمهندس أن يفكر في هذا كإشارتين. الأولى الخطأ ذو الإشارة $\delta$: في أي اتجاه تُغيَّر الأوزان. والثانية قيمته المطلقة $|\delta|$ أو جِدّة الحدث: \emph{بأي قوة} تُغيَّر، أي معدل التعلّم؛ فبعد حدث غير متوقع ينبغي التعلّم أسرع. وهي في معناها قريبة من معامل الكسب النورأدرينالي من الفصل~\ref{sec:neurontypes}.

وثمة خيار ثالث: سلك مستقل لمحور مستقل. عصبونات \nt{DOPA} الذاهبة إلى إحدى مناطق اختيار الأفعال لا تستجيب للمكافأة، بل لجِدّة المنبّه وشدته، وهي ضرورية لكي يتعلم الحيوان تجنّب الخطر~\cite{Menegas2018}: ما يسري في السلك ليس ”أفضل أم أسوأ“، بل ”خطر أم لا“.

المقيس: مجموعات مختلفة من عصبونات \nt{DOPA} تستجيب للمزعج وللجديد استجابات مختلفة، والمخارج المختلفة تعلّم أشياء مختلفة. أما كيف يستعمل الدماغ إشارة الأهمية بالضبط، معدلًا للتعلّم، أم إشارة انتباه، أم خطأً مستقلًا على محور الخطر، فمسألة قيد النقاش.

\section{لا التعليم وحده بل الحثّ أيضًا}

للدوبامين أثر فوري أيضًا: كلما زاد، أقبل الحيوان على الفعل برغبة وسرعة أكبر. كيف يتفق هذا مع التعلّم؟

الجواب الهندسي هو \emph{الفصل بالتردد}. الذرى والهبوطات السريعة التي تدوم مئات الميلي ثواني تحمل $\delta$ وتعلّم. والمستوى البطيء الذي يتغير خلال دقائق يحمل مقدارًا آخر، هو متوسط المكافأة في وحدة الزمن $\bar R$~\cite{Niv2007}. ليكن الفعل المنجز في زمن $\tau_a$ يتطلب جهدًا $C/\tau_a$: كلما كان أسرع كان أغلى. لكن كل ثانية تأخير تكلّف $\bar R$، أي قدر المكافأة التي كان يمكن نيلها فيها. الكلفة الكلية $C/\tau_a+\bar R\,\tau_a$ أصغر ما يمكن عند
\begin{empheq}[box=\keybox]{equation}
\tau_a^{*} \;=\; \sqrt{\frac{C}{\bar R}} .
\label{eq:vigor}
\end{empheq}
كلما كانت البيئة أغنى كان الوقت أغلى، وصار الأجدى أن يُعمل بسرعة أكبر: إذا تضاعف $\bar R$ قصر زمن الفعل بـ$\sqrt2\approx1.4$ مرة. ولاحظ أن $\bar R$ هو المكافأة المتوقعة نفسها التي طرحناها في الفصل~\ref{sec:dofa}.

قد يتغير تحرير الدوبامين في موضع الهدف دون أن يتغير تردد النبضات: فالإشارات المحلية عند المخارج نفسها تضبطه~\cite{Mohebi2019}. لكن المكونات البطيئة موجودة في النبضات نفسها أيضًا: في تجارب القسم التالي يُرى الارتفاع التدريجي عند الاقتراب من المكافأة في تردد عصبونات \nt{DOPA} كذلك~\cite{Kim2020}. إذن يتألف المستوى البطيء من مصدرين، تردد النبضات والضبط المحلي للتحرير، ويبدو أن أيهما الغالب يعتمد على المخرج وعلى المهمة.

\begin{keyblock}
\begin{proposition}[إشارتان على سلك واحد]
\label{prop:multiplex}
ينقل نظام \nt{DOPA} مقدارين على الأقل مفصولين في الزمن: خطأً سريعًا $\delta$ يعلّم، ومستوى بطيئًا يحدد ثمن الوقت~\eqref{eq:vigor} ومعه سرعة الأفعال. المقيس أن المكوّنين السريع والبطيء يسلكان سلوكين مختلفين؛ أما أن البطيء يساوي $\bar R$ تحديدًا ففرضية.
\end{proposition}
\end{keyblock}

\section{الارتفاع عند الاقتراب}

إذا ركض الحيوان في ممر نحو مكافأة بعيدة، يرتفع تركيز الدوبامين في منطقة اختيار الأفعال تدريجيًا على مدى ثوانٍ ما دام الهدف يقترب~\cite{Hamid2016}. ووفق الفصل~\ref{sec:dofa} لا ينبغي أن يحدث هذا: كل شيء يسير وفق الخطة، ويجب أن تكون $\delta$ صفرًا.

التفسير الأول: هذا الارتفاع ليس $\delta$، بل التوقع $V$ نفسه، أي إشارة ”الهدف قريب، يستحق الجهد“ من القسم السابق~\cite{Hamid2016}. والتفسير الثاني يحتفظ بـ$\delta$~\cite{Kim2020}. إذا كان التوقع صحيحًا، فإنه عند الأفق $\tau_\gamma$ ينمو في الطريق إلى المكافأة $R$ مثل $V=Re^{-(T-t)/\tau_\gamma}$، ووفق صيغة الخطأ~\eqref{eq:ctd} $\dd V/\dd t-V/\tau_\gamma=0$. لكن لنفترض أن التوقع من بعيد أقل من اللازم، وأن التقدير يتدقق مع الاقتراب. عندئذ ينمو التوقع بانحدار أشد، مثلًا $V=Re^{-(T-t)/\tau_v}$ مع $\tau_v<\tau_\gamma$، و
\begin{empheq}[box=\keybox]{equation}
\delta(t) \;=\; \frac{\dd V}{\dd t}-\frac{V}{\tau_\gamma} \;=\; V\Bigl(\frac{1}{\tau_v}-\frac{1}{\tau_\gamma}\Bigr) \;>\; 0 .
\label{eq:ramp}
\end{empheq}
الخطأ موجب ويزداد مع $V$، وهذا هو الارتفاع التدريجي نفسه (الشكل~\ref{fig:ramp}). عند $\tau_\gamma=4\,\text{s}$ و$\tau_v=1\,\text{s}$ نحصل على $\delta=0.75\,V$ في الثانية. اختُبرت هذه الصيغة في ممر افتراضي بنقل الحيوان لحظيًا إلى نقطة أقرب من الهدف: استجاب الدوبامين بذروة، كما يليق بالمشتقة، لا بمستوى متدرج~\cite{Kim2020}.

\begin{figure}[t]
\centering
\begin{tikzpicture}[>=Stealth,x=1.6cm,y=2.6cm]
\draw[->,gray!70!black] (-4.2,0) -- (0.5,0) node[below left,black] {\small $t$};
\draw[->,gray!70!black] (-4.2,0) -- (-4.2,1.2);
\foreach \x/\l in {-4/{$T-4\,\text{s}$},-2/{$T-2\,\text{s}$},0/{$T\,\text{s}$}}{\draw[gray!60!black] (\x,-0.03) -- (\x,0.03); \node[below,font=\scriptsize] at (\x,0) {\l};}
\draw[thick,densely dashed,gray!60!black] plot[domain=-4:0,samples=40] (\x,{exp(\x/4)});
\draw[very thick,blue!60!black] plot[domain=-4:0,samples=60] (\x,{exp(\x)});
\draw[very thick,red!70!black] plot[domain=-4:0,samples=60] (\x,{0.75*exp(\x)});
\node[anchor=south east,font=\small,gray!50!black] at (-1.3,0.77) {$V$ عند $\tau_v=\tau_\gamma$: $\delta=0$};
\node[right,font=\small,blue!60!black] at (0.05,1.0) {$V$ عند $\tau_v=1\,\text{s}$};
\node[right,font=\small,red!70!black] at (0.05,0.72) {$\delta$ وفق~\eqref{eq:ramp}};
\end{tikzpicture}
\caption{ارتفاع الدوبامين عند الاقتراب من المكافأة بوصفه خطأ تنبؤ، $\tau_\gamma=4\,\text{s}$. الخط المتقطع توقع ينمو تمامًا كما يقتضي الأفق ($\delta=0$)؛ والأزرق توقع ينمو بانحدار أشد؛ والأحمر الخطأ~\eqref{eq:ramp}.}
\label{fig:ramp}
\end{figure}

المقيس: الارتفاع التدريجي موجود، في تركيز الدوبامين وفي تردد نبضات عصبونات \nt{DOPA} معًا~\cite{Kim2020}، والقفزات اللحظية في الموقف تسبب قفزات في الدوبامين. أما أي التفسيرين صحيح، أو هل كلاهما صحيح على مخارج مختلفة، فمسألة قيد النقاش.

\section{النظر إلى الوراء}

كل النظريات السابقة تنظر إلى \emph{الأمام}. وإحدى النظريات الحديثة تعكس الاتجاه~\cite{Jeong2022}. لنفترض أن الدماغ لا يتعلم التنبؤ بالمكافأة من الإشارة المنبئة، بل \emph{يلتفت إلى الوراء} بعد أن ينال المكافأة: أي حدث سبقها أكثر من غيره؟ ومقياس هذه الصلة فرق بين احتمالين:
\begin{empheq}[box=\keybox]{equation}
M \;=\; P(\text{إشارة منبئة قبل المكافأة}) \;-\; P(\text{إشارة منبئة في نافذة عشوائية}) .
\label{eq:retro}
\end{empheq}
إذا ظهرت الإشارة المنبئة كثيرًا دون مكافأة كان الحد الثاني كبيرًا، والصلة ضعيفة. في هذه النظرية لا يُخبر الدوبامين بخطأ التنبؤ، بل بمدى كون الحدث المعيّن \emph{سببًا} محتملًا للمكافأة.

آلية الالتفات إلى الوراء تتطلب ما تتطلبه قاعدة العوامل الثلاثة في الفصل~\ref{sec:dofa}: يجب أن يترك كل حدث أثرًا يذوب، لكي يمكن لحظة المكافأة قراءة ما سبقها. والفرق ليس في المشبك، بل فيما ينقله \nt{DOPA}.

في التجارب التي تتنبأ فيها النظرية~\eqref{eq:retro} والخطأ~\eqref{eq:ctd} بنتائج مختلفة، تبع تحرير الدوبامين في عدد من الحالات الأولى~\cite{Jeong2022}. لكن في تجارب أخرى انزاحت استجابة الدوبامين أثناء التعلّم تدريجيًا من لحظة المكافأة إلى لحظة الإشارة المنبئة، كما يقتضي التعلّم بالخطأ~\eqref{eq:ctd}~\cite{Amo2022}. أي النظريتين تصف الدماغ، لا يُعرف بعد.

\section{الخطأ لا يخص المكافأة وحدها}

التوسيع الأخير يتعلق \emph{بموضوع} الخطأ. إذا استُبدل بالعصير المتوقع عصير آخر بالقيمة نفسها لكن بطعم مختلف، لم تتغير القيمة، والخطأ~\eqref{eq:ctd} يساوي صفرًا. ومع ذلك تستجيب عصبونات \nt{DOPA} لهذا الاستبدال~\cite{Takahashi2017}. والذرى المستحثّة اصطناعيًا للدوبامين تستطيع أن تعلّم الحيوان صلة بين منبّهين محايدين، دون أي مكافأة~\cite{Sharpe2017}. يبدو أن \nt{DOPA} يخبر بخطأ التنبؤ لا بالمكافأة وحدها، بل \emph{بما} سيحدث أيضًا.

صوريًا هذا تعميم بسيط لـ~\eqref{eq:ctd}: بدل دفق مكافأة واحد $r$ تؤخذ مجموعة من سمات ما يجري $\varphi_k$، كالطعم والمكان والصوت، ولكل منها توقعها $M_k$ وخطؤها~\cite{Gardner2018}:
\begin{empheq}[box=\keybox]{equation}
\delta_k(t) \;=\; \varphi_k(t) \;+\; \frac{\dd M_k}{\dd t} \;-\; \frac{M_k}{\tau_\gamma} ,
\qquad k=1,\dots,K .
\label{eq:vectortd}
\end{empheq}
المكافأة ليست إلا واحدة من السمات، ومتجه الأخطاء لا يعلّم ما هو جيد فحسب، بل ما يتبع ماذا أيضًا، أي نموذجًا للعالم.

ويتفق مع هذا عدم تجانس عصبونات \nt{DOPA}: في المهمة الواحدة تحمل عصبونات مختلفة مقادير مختلفة، بعضها مرتبط بالمكافأة، وبعضها بالحركة، وبعضها بموضع الانتباه~\cite{Engelhard2019}.

\begin{keyblock}
\begin{proposition}[حزمة أسلاك بدل سلك]
\label{prop:bundle}
إشارة نظام \nt{DOPA} ليست عددًا واحدًا. إنها موزعة على العصبونات (التوزيع~\eqref{eq:expectile}، والسمات المختلفة~\eqref{eq:vectortd}، ومحور الخطر المستقل) وعلى الزمن (الخطأ السريع والمستوى البطيء~\eqref{eq:vigor}). والخطأ العددي~\eqref{eq:ctd} تقريب أول لهذه الإشارة، كما أن الجامع ذا العتبة تقريب أول للعصبون.
\end{proposition}
\end{keyblock}

\section{الخلاصة}

\begin{table}[t]
\centering
\footnotesize
\setlength{\tabcolsep}{4pt}
\renewcommand{\arraystretch}{1.15}
\begin{tabular}{>{\raggedright}p{2.5cm}>{\raggedright}p{3.2cm}>{\raggedright}p{3.6cm}>{\raggedright\arraybackslash}p{4.2cm}}
\toprule
النظرية & ما يسري في السلك & القياس الرئيسي & ما هو معروف \\
\midrule
خطأ التنبؤ (الفصل~\ref{sec:dofa}) & العدد $\delta$~\eqref{eq:ctd} & ذروة، وانتقال إلى الإشارة المنبئة، وهبوط & مقيس؛ تقريب أول \\
التوزيع & مجموعة $\delta_i$ بلاتناظر مختلف & نقاط انقلاب مختلفة لعصبونات مختلفة & يتفق مع التجربة؛ ودور التشتت فرضية \\
الأهمية & $|\delta|$، الجِدّة؛ محور الخطر & ذرى للمزعج عند جزء من العصبونات & مقيس؛ وكيف يُستعمل قيد النقاش \\
ثمن الوقت & المستوى البطيء $\bar R$ & الدوبامين يسرّع الأفعال؛ والمكوّن البطيء في التحرير عند المخارج وفي تردد النبضات معًا & الفصل بين السريع والبطيء مقيس؛ و$\bar R$ فرضية \\
الارتفاع عند الاقتراب & $V$ أو $\delta$ مع تدقيق التقدير~\eqref{eq:ramp} & ارتفاع تدريجي، وقفزات عند النقل في الممر & الارتفاع مقيس؛ والتفسير قيد النقاش \\
النظر إلى الوراء & مقياس الصلة السببية~\eqref{eq:retro} & تجارب تختلف فيها تنبؤات النظريتين & النتائج متناقضة \\
متجه الأخطاء & $\delta_k$ بحسب السمات~\eqref{eq:vectortd} & استجابة لتغيّر الطعم؛ تعلّم بلا مكافأة & مقيس؛ والصيغة المتجهية فرضية \\
\bottomrule
\end{tabular}
\caption{ما ينقله نظام \nt{DOPA}: الصورة المعيارية في الفصل~\ref{sec:dofa} وتوسيعاتها الحديثة.}
\label{tab:dofaalt}
\end{table}

نظريات الجدول~\ref{tab:dofaalt} لا ينفي بعضها بعضًا: كلها تقريبًا تحتفظ بخطأ التنبؤ أساسًا وتوسّع صيغته. ولا ينافسه جديًا إلا النظر إلى الوراء، والتجارب هي التي ستحسم هذا الخلاف. والخلاصة للمهندس بسيطة: ثمن البثّ في القضية~\ref{prop:broadcastcost} ينخفض إذا كان السلك في الحقيقة أسلاكًا كثيرة بمرسَل إليهم مختلفين.

\section{تمارين}

\begin{exercise}[\lvlA]
\label{ex:07:1}
\emph{نقاط التوزيع لقطرة واحدة.} يأتي عصير حجمه $1$ باحتمال $p=0.2$، وإلا فالمكافأة $0$. جد من~\eqref{eq:expectile} قيم $V_i$ عند $\kappa_i=0.2$ و$0.5$ و$0.8$. كم وسيط هذه المكافأة، وبمَ تختلف النقطة~\eqref{eq:expectile} عن المئين؟
\end{exercise}

\begin{exercise}[\lvlB]
\label{ex:07:2}
\emph{التوزيع من عصبونين.} المكافأة $0$ أو $x$، واحتمال $x$ هو $p$. أفاد عصبونان يعملان بالقاعدة~\eqref{eq:asymmetric} بأن $V(1/2)=0.25$ و$V(0.8)=0.5$. استعد $p$ و$x$ والانحراف المعياري للمكافأة. بمَ سيفيد عصبون عند $\kappa=0.2$؟
\end{exercise}

\begin{exercise}[\lvlB]
\label{ex:07:3}
\emph{نمذجة: التوزيع من اللاتناظر.} حاكِ تسعة عصبونات \nt{DOPA} بقيم $\kappa_i=0.1,\dots,0.9$ والقاعدة~\eqref{eq:asymmetric}، $\alpha_i^+=2\kappa_i$، $\alpha_i^-=2(1-\kappa_i)$، ومكافأة موزعة بانتظام على $[0,1]$. كيف يعتمد تشتت $V_i$ على $\eta$، وأي $\eta$ يلزم لتمييز هذا التوزيع من قطرتين متساويتي الاحتمال $0.25$ و$0.75$؟
\end{exercise}

\begin{exercise}[\lvlA]
\label{ex:07:4}
\emph{ثمن الوقت.} (أ)~اشتق~\eqref{eq:vigor} وجد الكلفة الكلية عند الأمثل. (ب)~قدّر الدماغ $\bar R$ بخطأ قدره $k$ مرة. جد الزيادة النسبية في الكلفة عند $k=2$ و$k=4$. ما الدقة التي يجب أن يخبر بها المستوى البطيء للدوبامين عن $\bar R$؟
\end{exercise}

\begin{exercise}[\lvlB]
\label{ex:07:5}
\emph{الذروة عند النقل.} ينمو التوقع وفق~\eqref{eq:ramp} مع $\tau_v=1\,\text{s}$ و$\tau_\gamma=4\,\text{s}$؛ ويُنقل الحيوان لحظيًا من النقطة $T-2\,\text{s}$ إلى $T-1\,\text{s}$. جد مساحة ذروة $\delta$ كنسبة من $R$، وكم ثانية من الارتفاع التدريجي قبل النقل تحل محلها. ماذا سيُظهر النقل إذا كان الدوبامين يخبر بـ$V$ نفسه؟
\end{exercise}

\begin{exercise}[\lvlB]
\label{ex:07:6}
\emph{تصميم: إشارتان على السلك.} يسري في سلك \nt{DOPA} مستوى ينمو بمعدل $s=1/60$ نبضة في الثانية كل ثانية، وأحداث من $A=3$ نبضات. مشبك أثره $\tau_e=1\,\text{s}$ يطرح من التردد نسخته المنعّمة بـ$\tau_f$. عند أي $\tau_f$ لا يضيع أكثر من $10\,\%$ من الحدث، وتكون مساهمة المستوى خلال زمن الأثر أقل من $10\,\%$ من مساهمة الحدث؟
\end{exercise}

\begin{exercise}[\lvlB]
\label{ex:07:7}
\emph{تصميم تجربة: النظر إلى الوراء ضد الخطأ.} في النافذة تظهر الإشارة المنبئة باحتمال $0.1$، وبعدها المكافأة باحتمال $0.8$. قارن $\Delta P=P(\text{مكافأة}\mid\text{إشارة})-P(\text{مكافأة}\mid\text{لا إشارة})$ و$M$ من~\eqref{eq:retro} إذا (أ)~أُضيفت مكافآت بلا إشارة منبئة، $0.08$ لكل نافذة؛ (ب)~ضوعف تواتر الإشارة المنبئة دون مكافآت.
\end{exercise}

\begin{exercise}[\lvlC]
\label{ex:07:8}
\emph{حزمة الأسلاك وثمن البثّ.} في نموذج التمرين~\ref{ex:06:8} قُسّم $N$ عصبونًا إلى $K$ مجموعة لكل منها سلكها، ويتسرب إلى كل سلك جزء $\rho$ من أخطاء غيره. بيّن أن ثابت التعلّم يساوي $N/K+2+\rho^2N(1-1/K)$. كم محاولة يعطي عند $K\to\infty$ و$N=10^{10}$ و$\rho=0.01$؟
\end{exercise}

\chapter{متى نستكشف ومتى نقرر: نضيف \nt{NORA}}
\chaptermark{نضيف \nt{NORA}}
\label{sec:nora}
\begin{chaptergoals}
\item كيف ينقل مقبض كسب واحد العصبون من مضخّم تماثلي إلى مقارِن؛
\item لماذا لا يفيد رفع الكسب إلا ضد الضجيج الذي ينشأ بعد المضخّم؛
\item كيف ينقل الكسب شبكة الاختيار من التروّي إلى القرار دون أن يمسّ أي وزن؛
\item لماذا الكسب مقلوب درجة الحرارة، ولماذا تتبع المكافأة حرف U مقلوبًا بدلالة الكسب.
\end{chaptergoals}

\section{ما الذي ينقص}

في الفصل~\ref{sec:dofa} تعلّمت الشبكة بفضل الضجيج: كانت الانحرافات العشوائية في النشاط بحثًا، وكان \nt{DOPA} يخبر هل صار الوضع بها أفضل. لكن بقي هناك معامل واحد غير محدد: \emph{كم} نبحث. إذا كان الضجيج قويًا جدًا تخبطت الشبكة ولم تنتفع بما تعرفه. وإذا كان ضعيفًا جدًا اختارت الشيء نفسه مرة بعد مرة، ولم تلاحظ أن خيارًا قريبًا صار أفضل. في التعلّم المعزَّز يسمى هذا الاختيارَ بين الاستغلال والاستكشاف، ولكل خوارزمية مقبض له. وفي الفصل~\ref{sec:neurontypes} افترضنا أن \nt{NORA} هو الذي يدير هذا المقبض في الدماغ.

عصبونات النورأدرينالين عند الإنسان بضع عشرات الآلاف (الجدول~\ref{tab:types})، وكلها تقريبًا مجموعة في تجمّع صغير واحد، أما مخارجها فتتفرق في الدماغ كله تقريبًا، وإن تبيّن أن الأجزاء المختلفة من هذا التجمع تخدم جزئيًا مناطق مختلفة~\cite{Poe2020}. هذه قناة بثّ أضيق حتى من \nt{DOPA}. وهي تعمل في وضعين~\cite{AstonJones2005}. في الوضع \emph{الخلفي} تطلق العصبونات النبضات بانتظام، بتردد من أجزاء النبضة إلى بضع نبضات في الثانية، ويتغير هذا التردد ببطء مع درجة اليقظة. وفي الوضع \emph{الطوري} تستجيب برشقة قصيرة لحدث مهم للمهمة الجارية، نحو لحظة اتخاذ القرار. ونقطة الاختبار المريحة لكن غير الدقيقة هي قطر الحدقة: فهو يتغير مع نشاط هذه العصبونات، لكنه يتغير أيضًا مع نشاط مناطق أخرى~\cite{Joshi2016} ومع المخارج الكولينية في القشرة~\cite{Reimer2016}. الحدقة تبيّن المستوى العام للإثارة، لا \nt{NORA} وحده.

\section{معامل الكسب}

المقيس هو التالي: النورأدرينالين يكبت النشاط الخلفي للعصبونات أكثر مما يكبت استجابتها للمنبّه، فترتفع نسبة الإشارة إلى الخلفية~\cite{Foote1975NE}. أما أن انحدار خاصية العصبون المفرد يُضرب حرفيًا في معامل، فلا يلزم من التجربة. وكما في الفصل~\ref{sec:neurontypes}، نأخذ نموذجًا بسيطًا ينقل هذا الأثر: النورأدرينالين يغيّر انحدار خاصية النقل عند المتلقي~\cite{ServanSchreiber1990}. عندئذ تعطي المداخل الضعيفة تحت العتبة (الخلفية) أقل، والقوية أكثر. ليستجب العصبون بتردد
\begin{empheq}[box=\keybox]{equation}
r \;=\; F\bigl(g\,(u-\theta)\bigr), \qquad F(z)=\frac{r_{\max}}{1+e^{-z}} ,
\label{eq:gain}
\end{empheq}
حيث $u$ تيار الدخل، و$\theta$ العتبة، و$g$ معامل الكسب الذي يتحكم فيه \nt{NORA} في النموذج. عند $g$ صغير يتبع التردد الدخلَ بسلاسة في مدى واسع: العصبون \emph{مضخّم تماثلي}. وعند $g$ كبير يعطي أي دخل فوق العتبة تقريبًا التردد الكامل، ويعطي ما تحتها صفرًا: العصبون \emph{مقارِن}، أي عنصر شبه رقمي (الشكل~\ref{fig:gaincurves}). مقبض واحد ينقل العصبون نفسه بين وضعين تنفذهما في الإلكترونيات دارات مختلفة.

\begin{figure}[t]
\centering
\begin{tikzpicture}[>=Stealth,x=1.1cm,y=2.4cm]
\draw[->,gray!70!black] (-3.3,0) -- (3.4,0) node[below left,black] {\small $u$};
\draw[->,gray!70!black] (-3.3,0) -- (-3.3,1.2) node[left,black] {\small $r$};
\draw[densely dashed,gray!60!black] (0,0) -- (0,1.1);
\node[below,font=\small] at (0,0) {$\theta$};
\draw[densely dotted,gray!60!black] (-3.3,1) -- (3.2,1) node[right,black,font=\small] {$r_{\max}$};
\draw[very thick,blue!60!black] plot[domain=-3.2:3.2,samples=60] (\x,{1/(1+exp(-0.8*\x))});
\draw[very thick,orange!85!black] plot[domain=-3.2:3.2,samples=60] (\x,{1/(1+exp(-2*\x))});
\draw[very thick,red!70!black] plot[domain=-3.2:3.2,samples=120] (\x,{1/(1+exp(min(-8*\x,9)))});
\node[blue!60!black,font=\small,anchor=west] at (0.5,0.12) {$g$ صغير: مضخّم};
\node[red!70!black,font=\small,anchor=east] at (-0.3,0.85) {$g$ كبير: مقارِن};
\end{tikzpicture}
\caption{خاصية النقل~\eqref{eq:gain} عند ثلاث قيم لمعامل الكسب $g$.}
\label{fig:gaincurves}
\end{figure}

هل يساعد الكسب الكبير على مقاومة الضجيج؟ ليس دائمًا. ليختلف دخلان بمقدار $\Delta u$، والمطلوب أن نقول أيهما أكبر. إذا أُضيف الضجيج $\sigma_{\text{pre}}$ إلى الدخل \emph{قبل} المضخّم، تضخّمت الإشارة والضجيج معًا، ولم تتغير نسبتهما. أما إذا أُضيف الضجيج $\sigma_{\text{post}}$ \emph{بعد} المضخّم، من المشابك غير الموثوقة ومن النبضات العشوائية للشبكة نفسها، كما في الفصل~\ref{sec:code}، فإن نسبة الإشارة إلى الضجيج
\begin{empheq}[box=\keybox]{equation}
d' \;=\; \frac{g\,\Delta u}{\sqrt{g^2\sigma_{\text{pre}}^2+\sigma_{\text{post}}^2}}
\label{eq:gainsnr}
\end{empheq}
تزداد مع $g$ إلى أن يتساوى $g\sigma_{\text{pre}}$ مع $\sigma_{\text{post}}$~\cite{ServanSchreiber1990}.

\begin{keyblock}
\begin{proposition}[متى يساعد الكسب]
\label{prop:gainnoise}
حين يرفع \nt{NORA} معامل الكسب، فإنه يحسّن التمييز بين المداخل ضد الضجيج الذي ينشأ \emph{بعد} المضخّم، في الشبكة نفسها، فقط. أما الضجيج الذي جاء مع الدخل فلا يزيله الكسب.
\end{proposition}
\end{keyblock}

\section{الكسب في الاختيار بين اثنين}

لنعد إلى الاختيار من الفصل~\ref{sec:gaba}: مجموعتان من عصبونات \nt{GLUT} مدخلاهما $I_1$ و$I_2$ تثبّط إحداهما الأخرى بقوة $\beta$. الآن لكل مجموعة معامل كسب $g$: في معادلات الاختيار~\eqref{eq:wta} يُضرب ما بين القوسين في $g$. ما دامت المجموعتان تعملان، تُستخرج الترددات المستقرة من $r_1=g(I_1-\beta r_2)$ و$r_2=g(I_2-\beta r_1)$. بجمع هاتين المساواتين وطرحهما نحصل على المجموع $r_1+r_2=g(I_1+I_2)/(1+g\beta)$ والفرق $r_1-r_2=g(I_1-I_2)/(1-g\beta)$. ونسبتهما تبيّن مدى حدة تفضيل الشبكة لدخل على آخر:
\begin{empheq}[box=\keybox]{equation}
\frac{r_1-r_2}{r_1+r_2} \;=\; \varepsilon\,\frac{1+g\beta}{1-g\beta}, \qquad \varepsilon=\frac{I_1-I_2}{I_1+I_2}, \quad g\beta<1 .
\label{eq:wtagain}
\end{empheq}
الفرق الصغير بين المدخلين $\varepsilon$ يتضخم $(1+g\beta)/(1-g\beta)$ مرة، وهذا العامل يكبر بلا حد حين يقترب $g\beta$ من الواحد. وعند $g\beta>1$ تكون الحالة التي تعمل فيها المجموعتان غير مستقرة، كما في القضية~\ref{prop:wta}: تفوز واحدة وتصمت الأخرى (الشكل~\ref{fig:pitchfork}).

\begin{figure}[t]
\centering
\begin{tikzpicture}[>=Stealth,x=3.2cm,y=1.5cm]
\draw[->,gray!70!black] (0,0) -- (2.15,0) node[below left,black] {\small $g\beta$};
\draw[->,gray!70!black] (0,-1.25) -- (0,1.35) node[above,black] {\small $(r_1-r_2)/(r_1+r_2)$};
\draw[gray!60!black] (1,-0.04) -- (1,0.04); \node[below right,font=\small] at (1,0) {$1$};
\node[left,font=\scriptsize] at (0,1) {$1$}; \node[left,font=\scriptsize] at (0,-1) {$-1$};
\draw[very thick,blue!60!black] plot[domain=0:0.905,samples=60] (\x,{0.05*(1+\x)/(1-\x)}) -- (1,1) -- (2.05,1);
\draw[very thick,blue!60!black] (1.105,-1) -- (2.05,-1);
\draw[thick,densely dashed,gray!60!black] (1,0) -- (2.05,0);
\node[font=\small,align=center] at (0.45,-0.62) {تتروّى:\\المجموعتان\\تعملان};
\node[font=\small,align=center] at (1.55,0.45) {قرّرت:\\تعمل واحدة};
\node[font=\small,blue!60!black,anchor=south west] at (1.05,-0.97) {فاز الدخل الأضعف أولًا، فيبقى فائزًا};
\end{tikzpicture}
\caption{الاختيار بين اثنين عند معاملات كسب مختلفة، والمدخلان يختلفان بـ$\varepsilon=0.05$. عند $g\beta>1$ يكون القراران كلاهما مستقرين (الخطان المتصلان؛ وفوز الدخل الأضعف مستقر عند $g\beta>(1+\varepsilon)/(1-\varepsilon)\approx1.1$)، وتتمسك الشبكة بما اتخذته منهما؛ أما الحالة المتناظرة (الخط المتقطع) فغير مستقرة.}
\label{fig:pitchfork}
\end{figure}

\begin{keyblock}
\begin{proposition}[الكسب يبدّل وضع الاختيار]
\label{prop:gainwta}
في شبكة اختيار بتثبيط متبادل $\beta$ ينقل معامل الكسب $g$ الشبكة عبر الحد $g\beta=1$. تحته ”تتروّى“ الشبكة: تعمل المجموعتان، ويتضخم فرق المدخلين وفق~\eqref{eq:wtagain}. وفوقه ”قرّرت“ الشبكة: تعمل مجموعة واحدة، ويثبت القرار. وبتغيير $g$ يستطيع \nt{NORA} أن ينقل الشبكة بين هذين الوضعين دون أن يمسّ أي وزن.
\end{proposition}
\end{keyblock}

\section{الكسب بوصفه درجة حرارة}

لنضف الآن إلى الاختيار ضجيجًا ينشأ في الشبكة نفسها، بعد المضخّم. أي المجموعتين تفوز تحدده إشارة المقدار $g(u_A-u_B)+\eta$، حيث $u_A-u_B$ فرق المدخلين، أي فرق الأوزان المتعلَّمة من الفصل~\ref{sec:dofa}، و$\eta$ ضجيج مقياسه $\sigma$. إذا كان الضجيج موزعًا وفق القانون اللوجستي (والمنحنى يكاد يكون نفسه للتوزيع الغاوسي)، فإن احتمال اختيار $A$ هو
\begin{empheq}[box=\keybox]{equation}
P(A) \;=\; \frac{1}{1+e^{-g\,(u_A-u_B)/\sigma}} .
\label{eq:softmax}
\end{empheq}
سيعرف المبرمج هنا الاختيار بدالة \foreignlanguage{english}{softmax} بدرجة حرارة $T=\sigma/g$. معامل الكسب هو مقلوب درجة الحرارة. عند $g$ صغير لا يُختار حتى الخيار الأفضل بوضوح أكثر من الأسوأ إلا قليلًا: الشبكة تستكشف كثيرًا. وعند $g$ كبير يُختار أفضل خيار معروف دائمًا تقريبًا: الشبكة تستغل ما تعرفه.

لنوضح ما هو $g$ في~\eqref{eq:wtagain} و~\eqref{eq:softmax}: إنه الكسب \emph{الفعّال} لفرق مدخلات المهمة الجارية لحظة الاختيار. ووضعا \nt{NORA} يؤثران فيه تأثيرين متعاكسين. الرشقة الطورية لحظة القرار ترفعه، فتثبّت الشبكة اختيارها. أما النشاط الخلفي المرتفع فيرافق الاستكشاف على العكس: عند البشر تكون الحدقة القاعدية أوسع قبل الاختيار العشوائي~\cite{JepmaNieuwenhuis2011}، وعند الجرذان ينقل دخل \nt{NORA} إلى القشرة الحزامية الأمامية الاختيارَ إلى وضع عشوائي~\cite{Tervo2014stochastic}. إذن يقابل \nt{NORA} الخلفيَّ المرتفع $g$ فعّال \emph{صغير}، ويقابل الرشقةَ $g$ كبير.

\section{كم نستكشف}

أي كسب أفضل؟ تحققنا من ذلك على نموذج. تختار الشبكة أحد فعلين وفق~\eqref{eq:softmax}؛ وكل فعل يجلب مكافأة باحتمال ما، والشبكة تقدّره من مكافآت الفعل المختار، كما في الفصل~\ref{sec:dofa}. ومن حين لآخر يتغير العالم: يُسحب الاحتمالان من جديد. قد يتغير الفعل المختار، وقد يتغير الفعل الذي لم تجربه الشبكة منذ زمن؛ ولا تعرف الشبكة بالثاني إلا بالاستكشاف. يبيّن الشكل~\ref{fig:invertedU} نسبة أفضل مكافأة ممكنة التي تنالها الشبكة عند قيم ثابتة مختلفة لـ$g$.

\begin{figure}[t]
\centering
\begin{tikzpicture}[>=Stealth,x=1.2cm,y=1cm]
\draw[->,gray!70!black] (0,0) -- (7.6,0) node[below left,black] {\small $g$};
\draw[->,gray!70!black] (0,0) -- (0,4.4) node[above,black,font=\small] {نسبة أفضل مكافأة};
\foreach \x/\l in {0/0.5,1/1,2/2,3/4,4/8,5/16,6/32,7/64}{\draw[gray!60!black] (\x,-0.06) -- (\x,0.06); \node[below,font=\scriptsize] at (\x,0) {\l};}
\foreach \v/\y in {0.8/0.8,0.9/2.4,1.0/4.0}{\draw[gray!60!black] (-0.06,\y) -- (0.06,\y); \node[left,font=\scriptsize] at (0,\y) {\v};}
\draw[very thick,blue!60!black] plot[mark=*,mark size=1.3pt] coordinates {(0,0.368) (1,0.880) (2,1.728) (3,2.784) (4,3.456) (5,3.616) (6,3.488) (7,3.264)};
\draw[very thick,red!70!black] plot[mark=*,mark size=1.3pt] coordinates {(0,0.368) (1,0.672) (2,1.184) (3,1.840) (4,2.176) (5,1.920) (6,1.456) (7,1.104)};
\node[blue!60!black,font=\small,anchor=south] at (5,3.7) {عالم هادئ};
\node[red!70!black,font=\small,anchor=north] at (5.1,1.25) {عالم سريع التغير};
\draw[->,thick,blue!60!black] (5,3.25) -- (5,3.55);
\draw[->,thick,red!70!black] (4,1.8) -- (4,2.1);
\node[font=\small\itshape,gray!35!black,anchor=west] at (0.15,2.3) {تبحث عشوائيًا};
\node[font=\small\itshape,gray!35!black] at (6.45,2.55) {لا تلاحظ التغيرات};
\end{tikzpicture}
\caption{نسبة أفضل مكافأة ممكنة عند معامل كسب ثابت $g$ (المقياس على محور $g$ لوغاريتمي). المنحنى الأزرق: يتغير العالم في المتوسط مرة كل ألف محاولة، والأحمر: مرة كل عشرين. تشير السهام إلى أفضل $g$: في العالم السريع التغير هو أصغر بمرتين. المتوسط على 60 محاكاة في كل منها 4000 محاولة.}
\label{fig:invertedU}
\end{figure}

عند $g=0.5$ تكاد الشبكة لا تستعمل معرفتها، فتنال نحو ثلاثة أرباع الممكن، وعند $g$ كبير جدًا تفوتها التغيرات. أفضل قيمة: $g=16$ حين يتغير العالم مرة كل ألف محاولة (النسبة $0.976$)، و$g=8$ حين يتغير مرة كل عشرين ($0.886$).

\begin{keyblock}
\begin{proposition}[حرف U مقلوب]
\label{prop:explore}
المكافأة التي تجمعها الشبكة تعتمد على معامل الكسب اعتمادًا على شكل حرف U مقلوب: $g$ الصغير جدًا يهدر المحاولات في البحث، والكبير جدًا لا يلاحظ التغيرات. وأفضل $g$ يصغر كلما تغيّر العالم أسرع.
\end{proposition}
\end{keyblock}

ويُلاحَظ حرف U المقلوب في التجارب أيضًا: تحل الحيوانات المهمة على أفضل وجه عند نشاط خلفي معتدل لعصبونات \nt{NORA} مع استجابات طورية واضحة، أما عند النشاط الخلفي المرتفع فتتشتت وتتنقل بين المهام~\cite{AstonJones2005}. وهذا يتفق مع الفرق بين الوضعين في القسم السابق: الرشقة الطورية لحظة القرار تعطي $g$ فعالًا كبيرًا في الوقت الذي يجب فيه الاختيار تحديدًا، أما النشاط الخلفي المرتفع بلا رشقات فيضخّم كل شيء دون تمييز، بما في ذلك المداخل الدخيلة، فيهبط $g$ الفعّال للمهمة الجارية، وهذا هو الاستكشاف.

\section{من يدير المقبض}

ما دام أفضل كسب يعتمد على سرعة تغيّر العالم، فمن الحسن أن يُضبط. لكن من أين تعرف عصبونات \nt{NORA} أن العالم تغيّر؟ العلامة الطبيعية هي \emph{المفاجأة}: صارت أخطاء التنبؤ أكبر من المعتاد. من المهم التمييز بين نوعين من عدم اليقين. إذا كانت المكافأة عشوائية ببساطة، فالأخطاء كبيرة دائمًا، وهذا متوقع. أما إذا ازدادت الأخطاء فجأة مقارنة بحجمها المعتاد، فقد تغيّر شيء ما. ووفق إحدى الفرضيات، يخبر \nt{NORA} بعدم اليقين غير المتوقع هذا تحديدًا، أما المتوقع فيخبر به \nt{ACET}~\cite{YuDayan2005}.

ماذا نفعل بهذه الإشارة؟ يمكن خفض الكسب والاستكشاف أكثر. ويمكن رفع معدل التعلّم لنسيان التقديرات القديمة أسرع: تبيّن التجارب أن الناس يتعلمون أسرع بعد تغيّر حاد، وأن الحدقة تتسع عندئذ~\cite{Nassar2012}؛ ونذكّر بأن الحدقة مؤشر لا على \nt{NORA} وحده، بل على \nt{ACET} أيضًا~\cite{Reimer2016}. ويمكن فعل الأمرين معًا: وفق فرضية أخرى تعمل رشقة \nt{NORA} الطورية \emph{مقاطعةً}، أي إشارة تعيد بها الشبكة ضبط حالتها الجارية وتتكيف مع الموقف الجديد~\cite{BouretSara2005}، كخط المقاطعة العام في المعالج.

لنقل بصراحة ما لم نجده. جربنا في النموذج طريقتي الضبط البسيطتين كلتيهما: خفض $g$ أو رفع معدل التعلّم حين يتجاوز الخطأ المنعَّم متوسطه الطويل الأمد. في عالم يثبت مدة طويلة تارة ويتغير كثيرًا تارة أخرى، أعطت أفضل إعدادات الطريقتين $0.926\text{--}0.927$ من أفضل مكافأة، أي ما يكاد يساوي أفضل $g$ ثابت ($0.925$ عند $g=8$) أو معدل تعلّم ثابتًا لكنه كبير بما يكفي ($0.928$)، وأعطت الإعدادات غير الموفقة أقل. المنحنيات في الشكل~\ref{fig:invertedU} مسطحة قرب القمة، وفي عالم كهذا لا يكاد يوجد ما يُضبط. ينبغي أن يؤتي الضبط ثماره حيث تتطلب الأوضاع المختلفة إعدادات مختلفة جدًا. أما قانون التحكم الذي ينفذه \nt{NORA} بالضبط فلم يُحدَّد بعد.

\section{الخلاصة}

\begin{table}[t]
\centering
\footnotesize
\setlength{\tabcolsep}{4pt}
\renewcommand{\arraystretch}{1.15}
\begin{tabular}{>{\raggedright}p{3.0cm}>{\raggedright}p{4.4cm}>{\raggedright\arraybackslash}p{6.0cm}}
\toprule
ما يفعله \nt{NORA} & كيف & ما هو معروف \\
\midrule
يغيّر انحدار استجابة العصبونات & معامل الكسب $g$ في~\eqref{eq:gain} & ارتفاع نسبة الإشارة إلى الضجيج مقيس؛ والنموذج الضربي تبسيط \\
يبدّل وضع الاختيار & ينقل الشبكة عبر $g\beta=1$ & ينتج من النموذج~\eqref{eq:wtagain}؛ ولا قياسات مباشرة للحد \\
يحدد كم نستكشف & $g$ مقلوب درجة الحرارة~\eqref{eq:softmax}؛ الخلفي $g$ فعّال صغير (استكشاف)، والرشقة كبير (قرار) & حرف U المقلوب ووضعا العمل مقيسة؛ والصلة بالاستكشاف فرضية قوية الأساس \\
يخبر بالتغيرات & عدم اليقين غير المتوقع & اتساع الحدقة وارتفاع معدل التعلّم بعد التغيرات مقيسان؛ وأنه إشارة \nt{NORA} فرضية \\
يقاطع ويعيد الضبط & رشقة طورية عند حدث غير متوقع & الرشقات عند الأحداث المهمة مقيسة؛ و”المقاطعة“ فرضية \\
\bottomrule
\end{tabular}
\caption{ما يضيفه \nt{NORA} إلى شبكة من \nt{GLUT} و\nt{GABA} و\nt{DOPA}.}
\label{tab:nora}
\end{table}

\nt{DOPA} يقول للشبكة \emph{إلى أين} تغيّر الأوزان. أما \nt{NORA} فلا يغيّر أي وزن، لكنه يقرر \emph{كيف} تستعمل الشبكة ما تعرفه (الجدول~\ref{tab:nora}): مقبض واحد، هو معامل الكسب، يحوّل العصبون من مضخّم إلى مقارِن، وشبكة الاختيار من ”متروّية“ إلى ”قررت“، والاختيار من استكشاف إلى استغلال. أما كيف يعرف \nt{NORA} سرعة تغيّر العالم فسؤال مفتوح.

\section{تمارين}

\begin{exercise}[\lvlA]
\label{ex:08:1}
\emph{مقبض واحد للدماغ كله.} (أ)~قدّر من الجدول~\ref{tab:types} كم عصبونًا في الدماغ يقابل عصبون \nt{NORA} واحدًا. (ب)~الضجيج قبل المضخّم $\sigma_{\text{pre}}=0.5$، وبعده $\sigma_{\text{post}}=4$. عند أي $g$ يبلغ $d'$ من~\eqref{eq:gainsnr} $90\,\%$ من حدّه؟
\end{exercise}

\begin{exercise}[\lvlB]
\label{ex:08:2}
\emph{الاختيار مع معامل كسب.} $\tau\,\dd r_1/\dd t=-r_1+g[I_1-\beta r_2]_+$، $\tau\,\dd r_2/\dd t=-r_2+g[I_2-\beta r_1]_+$، $\tau=20\ms$. (أ)~في كم من الزمن يخمد $r_1-r_2$ عند $g\beta=0.5$ و$0.9$؟ (ب)~عند $\varepsilon=0.05$ جد قيمة $g\beta$ التي تعمل تحتها المجموعتان، وتلك التي يستقر فوقها فوز الدخل الأضعف.
\end{exercise}

\begin{exercise}[\lvlA]
\label{ex:08:3}
\emph{\foreignlanguage{english}{Softmax} من الضجيج.} تتلقى كل مجموعة من $K$ مجموعة ضجيج غامبل خاصًا بها، $P(\eta_k<z)=\exp(-e^{-z/\sigma})$، ويفوز أكبر $gu_k+\eta_k$. جد $P(k)$. عند أي $g$ يُختار الأفضل من خيارين، مع $u_A-u_B=0.1$ و$\sigma=1$، في $90\,\%$ من الحالات؟
\end{exercise}

\begin{exercise}[\lvlB]
\label{ex:08:4}
\emph{تقدير أفضل كسب.} في نموذج الشكل~\ref{fig:invertedU} تكون احتمالات المكافآت بعد التغيّر موزعة بانتظام على $[0,1]$. تجرب الشبكة الفعل الأسوأ مرة كل $e^{g\Delta}$ محاولة؛ وبمساواة ذلك بـ$1/h$ نحصل على $g^*\approx\ln(1/h)/\bar\Delta$. جد $\bar\Delta=\langle|p_A-p_B|\rangle$ و$g^*$ عند $h=0.001$ و$0.01$ و$0.05$.
\end{exercise}

\begin{exercise}[\lvlB]
\label{ex:08:5}
\emph{نمذجة: الكسب وسرعة التغيرات.} جد $g^*(h)$ عند $h$ من $0.001$ إلى $0.2$ بنسخة من \texttt{models/\allowbreak nora\_explore.py} (وهي تكتب ملفات في المجلد الحالي). أين يصح تقدير التمرين~\ref{ex:08:4}، ولماذا يكفّ $g^*$ عن التناقص عند التغيرات السريعة؟
\end{exercise}

\begin{exercise}[\lvlB]
\label{ex:08:6}
\emph{تصميم: مقاطعة لشبكة الاختيار.} شبكة التمرين~\ref{ex:08:2} مع $\beta=2$ و$\tau=20\ms$ و$g=1$ تثبت على $r_1=0.9$، $r_2=0$، مع أن المدخلين صارا $I_1=0.9$ و$I_2=1.0$. يخفض \nt{NORA} مدة $T_p$ قيمة $g$ إلى $g_{\text{lo}}$. جد أصغر $T_p$ تحليليًا عند $g_{\text{lo}}=0$، وبالنمذجة عند $g_{\text{lo}}=0.25$.
\end{exercise}

\begin{exercise}[\lvlC]
\label{ex:08:7}
\emph{متى يؤتي الضبط ثماره.} يعرف عرّاف هل الحقبة هادئة أم متقلبة، ويضع في كل منها $g$ خاصًا بها. (أ)~جد بالنمذجة مكسبه على أفضل $g$ ثابت في عالم الفصل (حقب من $1000$ محاولة، $h=0.001$ و$0.05$). (ب)~ابنِ عالمًا يكون فيه المكسب أكبر، وجد الجزء الذي يستعيده منه الضبط بحسب المفاجأة.
\end{exercise}

\chapter{متى نكتب ومتى نقرأ: نضيف \nt{ACET}}
\chaptermark{نضيف \nt{ACET}}
\label{sec:acet}

\section{ما الذي ينقص}

لرقاقة الذاكرة، إلى جانب أسلاك العنوان والبيانات، خط آخر: تمكين الكتابة. ما دام مطفأً تُقرأ الذاكرة فقط؛ وحين يُشغَّل يُكتب ما على خطوط البيانات. ومن دون هذا الخط لا تعمل الذاكرة: إذا كانت كل قراءة تكتب شيئًا في الوقت نفسه، فإن المقروء، ومنه المقروء خطأً، يُكتب فورًا من جديد.

في الدماغ تبدو هذه المسألة أحدّ منها في الرقاقة. في الفصل~\ref{sec:glutcomputer} كانت الذاكرة مجموعة من عصبونات \nt{GLUT} تُبقيها وصلاتها الخاصة مشتعلة (الشكل~\ref{fig:bistable}). وفي الفصل~\ref{sec:dofa} رأينا أن هذه الوصلات تتعلم بقاعدة هيب أثناء العمل مباشرة. إذن المشابك نفسها تحفظ القديم وتستقبل الجديد، ولا يوجد طور كتابة مستقل، كما لا توجد ساعة. فبمَ تميّز الشبكة اللحظة التي يجب فيها حفظ ما تراه من اللحظة التي يجب فيها تذكّر ما تعرفه؟ في الفصل~\ref{sec:neurontypes} افترضنا أن \nt{ACET} هو الذي يمسك هذا الخط. وفي هذا الفصل نتحقق من الكيفية التي يجب أن يعمل بها خط كهذا، ولماذا لا غنى عنه.

\section{ثلاثة أفعال لسلك واحد}

عصبونات الأسيتيل كولين العاملة داخل الدماغ قليلة، مجموعة في بضع مجموعات صغيرة، أما مخارجها فتتفرق في القشرة كلها. هذه قناة بثّ، مثل \nt{DOPA} و\nt{NORA}. مستوى الأسيتيل كولين في القشرة مرتفع في اليقظة المنتبهة ومنخفض في النوم البطيء~\cite{Hasselmo1999}. وكما في الدوبامين، يحدد معنى الإشارة مستقبِلُ المتلقي (القضية~\ref{prop:receptor}): للأسيتيل كولين مستقبلات سريعة تفتح قناة، ومستقبلات بطيئة تطلق تفاعلات داخل الخلية. ويهمنا هنا ثلاثة أفعال.

\emph{الأول: إضعاف الوصلات الداخلية.} بيّنت التجارب على شرائح القشرة الشمية أن الأسيتيل كولين يكبت بشدة النقل عبر الوصلات بين عصبونات المنطقة الواحدة، ويكاد لا يمس المداخل الواصلة إليها من الخارج~\cite{HasselmoBower1992}.

\emph{الثاني: تقوية المداخل.} يرفع الأسيتيل كولين استثارية العصبونات وييسّر النقل عبر بعض مسارات الدخل؛ فتصير الإشارة الآتية من أعضاء الحس أعلى صوتًا من نشاط الشبكة نفسها~\cite{Hasselmo2006}.

\emph{الثالث: تيسير تغيّر الأوزان.} عند مستوى مرتفع من الأسيتيل كولين تتغير الوصلات بسهولة أكبر~\cite{Hasselmo2006}.

بالنسبة إلى المهندس، الأفعال الثلاثة عملية واحدة: التبديل من وضع ”اقرأ“ إلى وضع ”اكتب“. تكفّ الشبكة عن الإصغاء إلى نفسها، وتبدأ بالإصغاء إلى الدخل، وتحفظ ما تسمعه.

\section{شبكة تُكمل}

لنأخذ $N$ عصبونًا من عصبونات \nt{GLUT} موصولة بعضها ببعض. ولنحفظ في وصلاتها بضعة أنماط، أي مجموعات من $pN$ عصبونًا نشطة معًا، بقاعدة هيب~\cite{Hebb1949}. إذا قُدّم إلى هذه الشبكة نصف نمط، أثارت الوصلات النصف الناقص: الشبكة \emph{تُكمل} النمط من جزئه، كما كانت المجموعة في الشكل~\ref{fig:bistable} تُبقي نفسها. هذه ذاكرة ترابطية: العنوان فيها جزء من المحتوى. و\nt{GABA}، كما في الفصل~\ref{sec:gaba}، يُبقي العدد الكلي للعصبونات النشطة قرب $pN$، لكيلا يشتعل نمطان معًا.

لنرمز لمستوى الأسيتيل كولين بـ$a$، من $0$ إلى $1$، ولنكتب أفعاله الثلاثة في النموذج مباشرة:
\begin{empheq}[box=\keybox]{equation}
\tau\frac{\dd r_i}{\dd t} = -r_i + F\Bigl(\tfrac{1+a}{2}\,x_i + (1-a)\sum_j W_{ij}\,r_j - g\Bigr),
\qquad
\frac{\dd W_{ij}}{\dd t} = \eta\,a\,(r_i-p)(r_j-p) .
\label{eq:acetnet}
\end{empheq}
هنا $r_i$ تردد العصبون $i$، و$x_i$ دخله من أعضاء الحس، و$F$ خاصية نقل ذات عتبة، و$g$ التثبيط العام من \nt{GABA}، الذي يزداد حين يتجاوز عدد العصبونات النشطة $pN$. العامل $\frac{1+a}{2}$ تقوية الدخل، و$1-a$ إضعاف الوصلات الداخلية، و$\eta a$ معدل التعلّم. والمعادلة الثانية قاعدة هيب~\eqref{eq:hebb} بصيغة مناسبة للأنماط المتناثرة~\cite{Tsodyks1988}: يزداد الوزن حين يكون العصبونان كلاهما أنشط من المعتاد، وينقص حين يكون أحدهما فقط نشطًا. والجزء السالب من الأوزان، كالعادة، تقف وراءه وسائط \nt{GABA}. لا توجد ساعة: العصبونات والأوزان تتغير كلها باتصال.

\begin{figure}[t]
\centering
\begin{tikzpicture}[>=Stealth,
  nrn/.style={circle,draw,very thick,fill=blue!6,minimum size=0.8cm,inner sep=0pt,font=\small},
  acet/.style={circle,draw,very thick,double,double distance=1.2pt,fill=orange!15,minimum size=1.35cm,inner sep=0pt,font=\scriptsize},
  bc/.style={->,thick,densely dashed,orange!85!black}]
\foreach \i/\y in {1/2.4,2/1.2,3/0}{
  \node[nrn] (N\i) at (3,\y) {$r_\i$};
  \draw[->,thick,blue!60!black] (0,\y) node[left,black] {\small $x_\i$} -- (N\i);}
\node[left,align=right,font=\small] at (-0.5,1.2) {من أعضاء\\الحس};
\draw[->,thick,gray!60!black] (N1) to[bend left=30] (N2);
\draw[->,thick,gray!60!black] (N2) to[bend left=30] (N1);
\draw[->,thick,gray!60!black] (N2) to[bend left=30] (N3);
\draw[->,thick,gray!60!black] (N3) to[bend left=30] (N2);
\draw[->,thick,gray!60!black] (N1.east) .. controls (4.6,2.0) and (4.6,0.4) .. (N3.east);
\node[right,font=\small,gray!40!black,align=left] at (4.7,1.2) {الوصلات\\الداخلية $W$};
\node[acet] (A) at (1.2,-1.9) {\nt{ACET}};
\draw[bc] (A) -- (1.6,-0.15);
\node[font=\small,orange!60!black,anchor=east] at (0.95,-0.9) {$+$ الدخل};
\draw[bc] (A) -- (4.3,0.6);
\node[font=\small,orange!60!black,anchor=west] at (4.3,0.1) {$-$ الوصلات الداخلية، $+$ معدل التعلّم};
\node[right,font=\small,align=left] at (2.2,-1.9) {مستوى مرتفع: ”اكتب“\\مستوى منخفض: ”اقرأ“};
\end{tikzpicture}
\caption{\nt{ACET} خطًا لتمكين الكتابة. العصبونات $r_i$ تتلقى الدخل $x_i$ من أعضاء الحس (الأسهم الزرقاء)، ويثير بعضها بعضًا عبر الوصلات الداخلية $W$ (الرمادية) التي تُحفظ فيها الأنماط. الأسيتيل كولين (الأسهم المتقطعة) يقوّي الدخل، ويضعف الوصلات الداخلية، وييسّر تغيّر الأوزان، كما في~\eqref{eq:acetnet}.}
\label{fig:acetnet}
\end{figure}

\section{الكتابة والقراءة يعيق أحدهما الآخر}

لنختبر النموذج~\eqref{eq:acetnet} على مهمة تتعارض فيها الكتابة والقراءة فعلًا. شبكة من $200$ عصبون تحفظ خمسة أنماط، في كل منها $20$ عصبونًا نشطًا. والمطلوب الآن حفظ نمط سادس يتطابق نصفه مع أحد الأنماط القديمة: عشرة عصبونات مشتركة بينهما. وهذا يحدث دائمًا: الغرفة الجديدة تشبه غرفة مألوفة، والوجه الجديد يشبه وجهًا قديمًا.

\emph{الكتابة عند $a$ منخفض.} لنفصل أولًا الفعلين الأولين للأسيتيل كولين عن الثالث: نُبقي معدل التعلّم كاملًا، ولا يتغير إلا تقوية الدخل وإضعاف الوصلات الداخلية. عند $a=0$ يُشعل الدخل العصبونات العشرة المشتركة، وتُشعل هذه عبر الوصلات الداخلية النصف الباقي من النمط \emph{القديم}. لا يدع \nt{GABA} النمطين يشتعلان معًا، فيطرد النمطُ القديم، المدعوم بوصلاته، النمطَ الجديد الذي لا يستند إلا إلى الدخل. الشبكة لا ترى ما أمامها، بل ما تتذكره، وقاعدة هيب تكتب المرئي بإخلاص.

والنتيجة أن النمط الجديد لم يُحفظ أصلًا: من نصفه المميز لا تتذكر الشبكة شيئًا. أما القديم فقد فسد: إذا استُدعي بنصفه المميز جرّ معه في المتوسط ستة عصبونات غريبة من عشرة. عند $a=0.1$ يُحفظ النمط الجديد، لكنه يندمج بالقديم، ويكون الفساد أشد: كل منهما، إذا استُدعي بنصفه، يجرّ معه نحو ثمانية عصبونات غريبة؛ وبدءًا من $a=0.2$ تصير الكتابة نظيفة: العصبونات الزائدة عند التذكّر أقل من واحد (بالمتوسط على عشر محاكاة).

\emph{الكتابة بمعدل التعلّم الحقيقي.} لنجعل الآن الأسيتيل كولين يحدد معدل التعلّم أيضًا، كما في~\eqref{eq:acetnet}. في عرض واحد لا يُحفظ النمط الجديد كاملًا إلا عند $a\ge0.9$؛ وعند $a=0.75$ يُحفظ منه ثمانية أعشار، وعند $a\le0.5$ لا يُحفظ أبدًا.

\emph{القراءة.} لنقدّم إلى شبكة فيها خمسة أنماط مكتوبة كتابة نظيفة نصفَ نمط، ثم نزيله. عند $a\le0.2$ تُكمل الشبكة النمط كله وتُبقيه وحدها؛ وعند $a=0.3$ تُكمله كله تقريبًا؛ وعند $a\ge0.4$ تكون الوصلات الداخلية قد ضعفت أكثر مما ينبغي، فينطفئ النشاط بعد اختفاء التلميح (الشكل~\ref{fig:acetsim}).

\begin{figure}[t]
\centering
\begin{tikzpicture}[>=Stealth,x=9cm,y=3.2cm]
\draw[->,gray!70!black] (0,0) -- (1.1,0) node[right,black] {\small $a$};
\draw[->,gray!70!black] (0,0) -- (0,1.2);
\foreach \x in {0,0.2,0.4,0.6,0.8,1.0}{\draw[gray!60!black] (\x,-0.02) -- (\x,0.02); \node[below,font=\scriptsize] at (\x,0) {$\x$};}
\foreach \y in {0,0.5,1}{\draw[gray!60!black] (-0.01,\y) -- (0.01,\y); \node[left,font=\scriptsize] at (0,\y) {$\y$};}
\node[rotate=90,font=\small] at (-0.1,0.6) {نسبة النمط المستعاد};
\draw[very thick,blue!60!black] plot[mark=*,mark size=1.6pt] coordinates {(0,1.00) (0.1,1.00) (0.2,1.00) (0.3,0.96) (0.4,0.04) (0.5,0.00) (0.75,0.00) (1.0,0.00)};
\draw[very thick,red!70!black] plot[mark=*,mark size=1.6pt] coordinates {(0.1,0.00) (0.2,0.00) (0.3,0.00) (0.5,0.00) (0.75,0.80) (0.9,1.00) (1.0,1.00)};
\node[font=\small,blue!60!black,anchor=west,align=left] at (0.03,0.5) {القراءة:\\النمط\\من نصفه};
\node[font=\small,red!70!black,anchor=east,align=right] at (1.0,0.35) {الكتابة:\\نمط جديد\\دفعة واحدة};
\end{tikzpicture}
\caption{النموذج~\eqref{eq:acetnet}: $200$ عصبون، وفي كل نمط $20$ عصبونًا نشطًا، والمتوسط على عشر محاكاة. النقاط الزرقاء القراءة: أي نسبة من النمط تُستعاد من نصفه بعد إزالة التلميح. النقاط الحمراء الكتابة: أي نسبة من نمط جديد يشبه نصفه نمطًا قديمًا تُتذكَّر بعد عرض واحد، إذا كان الأسيتيل كولين يحدد معدل التعلّم أيضًا. القراءة تنجح عند $a\le0.3$، والكتابة عند $a\ge0.9$.}
\label{fig:acetsim}
\end{figure}

\begin{keyblock}
\begin{proposition}[الكتابة والقراءة تتطلبان وضعين مختلفين]
\label{prop:writeread}
في النموذج~\eqref{eq:acetnet}، حيث الوصلات نفسها تحفظ القديم وتتعلم الجديد، ويحدد الأسيتيل كولين معدل التعلّم أيضًا، لا يوجد مستوى $a$ تنجح عنده الكتابة والقراءة بالقدر نفسه. للكتابة يجب إضعاف الوصلات الداخلية، وإلا كتبت الشبكة ما تذكّرته لا الدخل؛ وللقراءة يجب تقويتها، وإلا لم تُكمل النمط من جزئه. يلزم مبدِّل، ويمكن لسلك بثّ واحد أن يؤدي هذا الدور.
\end{proposition}
\end{keyblock}

لو لم يغيّر الأسيتيل كولين معدل التعلّم، لصلح للعملين شريط ضيق $0.2\le a\le0.3$. لكن الشبكة عندئذ ستتعلم مع كل قراءة، أي ستعيد كتابة المقروء بأخطائه. أما الفكرة نفسها، أي كبت الوصلات الداخلية أثناء التعلّم، فمأخوذة من نماذج بُنيت على قياسات في الشرائح~\cite{HasselmoAndersonBower1992}. والأرقام أعلاه تخص نموذجنا المبسّط؛ أما أين تقع حدود الأوضاع في الدماغ بالضبط فغير معروف.

\section{إلى أي حد نصدّق العينين}

المبدِّل ”اكتب/اقرأ“ حالة قصوى لسؤال أعم: إلى أي حد نصدّق الدخل، وإلى أي حد نصدّق الذاكرة؟ لدى المهندسين جواب دقيق عن هذا السؤال، وهو يوضح لماذا يتحكم سلك واحد في الوضع ومعدل التعلّم معًا.

لتتابع الشبكة مقدارًا $s(t)$ يهيم ببطء: في كل ثانية يزداد تباينه بمقدار $q$. تُخبر أعضاء الحس بـ$y(t)=s(t)+$ضجيج شدته $R$. وأفضل تقدير $\hat s$ في الزمن المتصل يعطيه مرشح كالمان-بيوسي~\cite{KalmanBucy1961}:
\begin{empheq}[box=\keybox]{equation}
\frac{\dd\hat s}{\dd t} \;=\; \frac{P}{R}\,\bigl(y-\hat s\bigr),
\qquad
\frac{\dd P}{\dd t} \;=\; q-\frac{P^2}{R} ,
\label{eq:kalman}
\end{empheq}
حيث $P$ تباين خطأ التقدير، أي مدى عدم يقين الشبكة نفسها بما تعرفه. المعادلة الأولى هي دارة \foreignlanguage{english}{RC} نفسها التي في الغشاء في نموذج العصبون~\eqref{eq:glutneuron}، لكن بثابت زمني $R/P$ متغير. حين لا تكون الشبكة واثقة يكون $P$ كبيرًا، والثابت الزمني صغيرًا، فيتبع التقدير الدخل بسرعة: هذا ”اكتب“. وحين تكون واثقة يكون $P$ صغيرًا، ولا يكاد التقدير يصغي إلى الدخل، متمسكًا بالذاكرة: هذا ”اقرأ“.

في الحالة المستقرة $P=\sqrt{qR}$، والثابت الزمني للذاكرة $\sqrt{R/q}$: إذا انزاح العالم بمقدار واحد في مئة ثانية، أي $q=0.01$، وكان ضجيج أعضاء الحس $R=1$، فعلى الذاكرة أن تتذكر نحو عشر ثوانٍ.

الأهم هنا أن عددًا واحدًا $P/R$ يحدد شيئين معًا: وزن الدخل مقارنةً بالذاكرة، وسرعة تغيّر التقدير المحفوظ. وهذا بالضبط ما يفعله الأسيتيل كولين في~\eqref{eq:acetnet}. ووفق فرضية~\cite{YuDayan2005}، يخبر الأسيتيل كولين الشبكة بمقدار شبيه بـ$P$، هو عدم اليقين \emph{المتوقع}، الذي تعرفه الشبكة مسبقًا. ولا تكون هذه القناة بطيئة دائمًا: جزء من عصبونات \nt{ACET} يستجيب للمكافأة والعقاب في نحو عشرين ميلي ثانية، وبقوة تزداد كلما كان التعزيز أقل توقعًا~\cite{Hangya2015}.

\begin{keyblock}
\begin{proposition}[مقبض واحد لوزن الدخل ومعدل التعلّم]
\label{prop:kalman}
في المرشح الأمثل~\eqref{eq:kalman} يكون الوزن الذي يُمزج به الدخل مع الذاكرة، ومعدل التعلّم، مقدارًا واحدًا هو $P/R$. لذلك من المعقول التحكم فيهما بإشارة واحدة. ومع ثبات خصائص العالم تستقر هذه الإشارة عند مستوى ثابت $\sqrt{q/R}$، ولا حاجة إلى ضبطها إلا حين تتغير خصائص العالم.
\end{proposition}
\end{keyblock}

الجملة الثانية من القضية تفسر نتيجة الفصل~\ref{sec:nora}، حيث لم يتفوق ضبط معدل التعلّم بحسب المفاجأة على أفضل معدل ثابت: في عالم خصائصُ تغيّره ثابتة يكون أفضل معدل تعلّم ثابتًا فعلًا. ويظهر المكسب حين يصير العالم فجأة مختلفًا. عندئذ يجد التقدير نفسه فجأة بعيدًا عن الدخل، و$P$ لا يعلم بذلك. والفعل الصحيح أن يُرفع $P$ بحدة، فيُنتقل بذلك إلى وضع الكتابة. ووفق الفرضية التي ناقشناها في الفصل~\ref{sec:nora}، يخبر \nt{NORA} بتغيّر كهذا \emph{غير متوقع}~\cite{YuDayan2005}. ويأتي تقسيم العمل طبيعيًا: \nt{NORA} يلاحظ أن العالم تغيّر، و\nt{ACET} يُبقي الشبكة في وضع الكتابة إلى أن يُتعلَّم الموقف الجديد.

\section{النوم وضعًا للقراءة}

إذا كان المستوى المرتفع للأسيتيل كولين وضعَ الكتابة، فماذا تفعل الشبكة عند المستوى المنخفض؟ في النوم البطيء ينخفض مستوى الأسيتيل كولين، فتتحرر الوصلات الداخلية من الكبت، ولا يكاد يوجد دخل من أعضاء الحس. الشبكة في وضع القراءة بلا تلميح تعيد تشغيل ما هو مكتوب فيها. ويُظهر تسجيل النشاط أن العصبونات التي عملت معًا في النهار تنطلق معًا من جديد في النوم البطيء~\cite{WilsonMcNaughton1994}، بل بالترتيب نفسه~\cite{Skaggs1996}. ووفق فرضية~\cite{Hasselmo1999}، تنقل إعادات التشغيل هذه ما تعلّمته الشبكة بسرعة في النهار إلى الذاكرة الطويلة الأمد (وإطالة ومضات إعادة التشغيل القصيرة اصطناعيًا تحسّن الذاكرة~\cite{FernandezRuiz2019}): في النهار كتابة من الدخل، وفي الليل إعادة كتابة من الداخل. وبالنسبة إلى المهندس يشبه هذا الكتابة في مخزن مؤقت سريع نهارًا، وتفريغه في ذاكرة دائمة بطيئة ليلًا.

\section{الخلاصة}

\begin{table}[t]
\centering
\footnotesize
\setlength{\tabcolsep}{4pt}
\renewcommand{\arraystretch}{1.15}
\begin{tabular}{>{\raggedright}p{3.2cm}>{\raggedright}p{4.2cm}>{\raggedright\arraybackslash}p{6.0cm}}
\toprule
ما يفعله \nt{ACET} & كيف & ما هو معروف \\
\midrule
يضعف الوصلات الداخلية & العامل $1-a$ في~\eqref{eq:acetnet} & مقيس في الشرائح \\
يقوّي الدخل & العامل $\frac{1+a}{2}$ & رفع الاستثارية والنقل عبر مسارات الدخل مقيس \\
ييسّر التعلّم & المعدل $\eta a$ & مقيس \\
يبدّل ”اكتب/اقرأ“ & القضية~\ref{prop:writeread} & ينتج من النماذج؛ ولا قياس مباشرًا للوضعين \\
يخبر بعدم اليقين المتوقع & $P$ في~\eqref{eq:kalman} & فرضية \\
يحرر الشبكة لإعادة التشغيل في النوم & $a$ منخفض في النوم البطيء & المستوى المنخفض في النوم وإعادات التشغيل مقيسة؛ والنقل إلى الذاكرة الطويلة الأمد فرضية \\
\bottomrule
\end{tabular}
\caption{ما يضيفه \nt{ACET} إلى شبكة من \nt{GLUT} و\nt{GABA} و\nt{DOPA} و\nt{NORA}.}
\label{tab:acet}
\end{table}

\nt{DOPA} يقول للشبكة إلى أين تغيّر الأوزان، و\nt{NORA} بأي حزم تستعمل ما تعرفه، و\nt{ACET} هل تصغي إلى العالم أم إلى نفسها (الجدول~\ref{tab:acet}). هذا آخر خطوط التحكم بالبثّ في الجدول~\ref{tab:types}: تمكين الكتابة، الذي لولاه لأفسدت الذاكرةُ، التي تحفظ الأنماط في الوصلات نفسها التي تقرؤها بها، نفسَها مع كل قراءة.

\section{تمارين}

\begin{exercise}[\lvlA]
\label{ex:09:1}
\emph{كم نتذكر.} المرشح~\eqref{eq:kalman} مع $q=0.01$ و$R=1$ يتذكر نحو $10\,\text{s}$. جد $P_\infty$ والذاكرة $\sqrt{R/q}$ إذا (أ)~تضاعف ضجيج المستشعر في السعة؛ (ب)~صار العالم يهيم أسرع بمئة مرة. (ج)~كم مرة يجب أن يزداد الضجيج لكي تتذكر الشبكة دقيقة؟
\end{exercise}

\begin{exercise}[\lvlB]
\label{ex:09:2}
\emph{وضع الكتابة بعد التغيّر.} تغيّر العالم، فقفز عدم يقين المرشح~\eqref{eq:kalman} مع $q=0.01$ و$R=1$ إلى $P_0\to\infty$. (أ)~بأي ثابت زمني يعود $P$ إلى $P_\infty$؟ (ب)~بعد كم ثانية لا يتجاوز معدل التعلّم قيمته المستقرة بأكثر من $10\,\%$؟
\end{exercise}

\begin{exercise}[\lvlB]
\label{ex:09:3}
\emph{معدل تعلّم ثابت.} تتعلم الشبكة وفق $\dd\hat s/\dd t=(y-\hat s)/T$؛ والعالم يهيم، $\dd s/\dd t=w$، $y=s+n$، حيث $w$ و$n$ ضجيجان أبيضان شدتاهما $q$ و$R$. جد أفضل $T$ وتباين الخطأ عنده. كم مرة يزداد إذا أخطأنا في $T$ بمعامل $2$ و$10$؟
\end{exercise}

\begin{exercise}[\lvlB]
\label{ex:09:4}
\emph{أين حد الكتابة.} في \texttt{models/\allowbreak acet\_memory.py} العتبة $\theta=0.35$، والوزن داخل النمط $w=0.041$ (المثالي $0.045$). النمط الجديد يتقاسم $m=10$ عصبونات مع القديم. عند أي $a$ في~\eqref{eq:acetnet} لا تشتعل بقية عصبونات النمط القديم من هذه الـ$m$ (العتبة حادة)؟
\end{exercise}

\begin{exercise}[\lvlB]
\label{ex:09:5}
\emph{نمذجة: سعة الذاكرة.} عدّل في \texttt{models/\allowbreak acet\_memory.py} الدالة \texttt{read\_sweep}: عند $a=1$ اكتب $M$ نمطًا ($M$ من $5$ إلى $100$)، ثم عند $a=0$ اقرأ العشرة الأولى من أنصافها. عند أي $M$ تظهر الأخطاء، وكيف يعتمد الحد $M_{\max}$ على $N$ و$p$؟
\end{exercise}

\begin{exercise}[\lvlB]
\label{ex:09:6}
\emph{تصميم: كتابة بلا تسرّب.} مع معدل التعلّم $\eta a$ تتعلم الشبكة أثناء القراءة أيضًا. اقترح دالة رتيبة $\eta(a)\le\eta$ تساوي صفرًا عند $a\le0.3$ ولا تقل عن $\eta a$ عند $a\ge0.9$. تحقق: بعد $20$ قراءة عند $a=0.2$ بتلميح فيه ثلاثة عصبونات غريبة، كم منها دخل النمط؟
\end{exercise}

\begin{exercise}[\lvlC]
\label{ex:09:7}
\emph{كيف يُنقل المخزن المؤقت ليلًا.} (أ)~في النوم $a=0.05$، وتدوم إعادة التشغيل $0.1\,\text{s}$ مقابل $0.2\,\text{s}$ نهارًا عند $a=1$. كم إعادة تشغيل تراكم تغيّر الأوزان الناتج عن عرض نهاري واحد، بالمعدل $\eta a$، وبالدالة $\eta(a)$ من التمرين~\ref{ex:09:6}؟ (ب)~المخزن الدائم يتعلم أبطأ من المؤقت بمئة مرة ويسمع النمط $20$ مرة في الليلة، مدة كل منها $0.1\,\text{s}$. كم ليلة يلزم؟
\end{exercise}

\chapter{الآلة كاملة}
\chaptermark{الآلة كاملة}
\label{sec:synthesis}

\section{ما الذي جمعناه}

أخذنا أنواع العصبونات في الجدول~\ref{tab:types} واحدًا تلو الآخر، وسألنا: ماذا يضيف كلٌّ منها إلى آلة حاسبة؟ وكانت القواعد في كل مرة واحدة: لا شيء إلا ما يوجد في كائن حي، وكل شيء يجري في زمن متصل. والآن نجمع الأجزاء في آلة واحدة ونرى كيف تعمل معًا.

الصورة التي رسمناها في الفصل~\ref{sec:neurontypes} صمدت وازدادت دقة. شبكة معنونة هائلة من عصبونات \nt{GLUT} تحمل البيانات. وعصبونات \nt{GABA} تتحكم في تدفق هذه البيانات. وفوق الشبكة أربع قنوات بثّ، لكلٍّ منها مقبضها: \nt{DOPA} يقول أيّ تغييرات الأوزان تُحفظ، و\nt{SERO} يمسك المستوى العام ويحدد، بحسب فرضية، أفق التخطيط، و\nt{NORA} يختار بين البحث والحسم، و\nt{ACET} بين الكتابة والقراءة (الشكل~\ref{fig:machine}).

\begin{figure}[t]
\centering
\begin{tikzpicture}[>=Stealth,
  blk/.style={draw,very thick,rounded corners=3pt,align=center,font=\small},
  reg/.style={draw,very thick,double,double distance=1.2pt,circle,minimum size=1.25cm,inner sep=0pt,font=\scriptsize,fill=orange!15},
  bc/.style={->,thick,densely dashed,orange!85!black}]
% sensors, bus, actions
\node[blk,fill=green!8,text width=1.7cm] (S) at (0.9,0) {المستشعرات\\\scriptsize تشغيل/إطفاء};
\draw[very thick,rounded corners=4pt,fill=blue!5] (2.6,-1.35) rectangle (9.0,1.35);
\node[font=\small] at (5.8,0.95) {ناقل \nt{GLUT}: البيانات};
\node[font=\scriptsize,align=center] at (4.3,0.25) {التزامنات، التأخيرات،\\المنطق (الفصل~\ref{sec:glutcomputer})};
\node[font=\scriptsize,align=center] at (7.4,0.25) {الذاكرة، الشفرة\\(الفصلان~\ref{sec:code} و\ref{sec:acet})};
\draw[thick,red!70!black,fill=red!8,rounded corners=2pt] (2.8,-1.15) rectangle (8.8,-0.55);
\node[font=\scriptsize,red!60!black] at (5.8,-0.85) {\nt{GABA} (الفصل~\ref{sec:gaba}): منع، اختيار، قسمة، إيقاع};
\node[blk,fill=blue!5,text width=2.0cm] (A) at (10.9,0) {اختيار\\الفعل};
\draw[->,very thick] (S) -- (2.6,0);
\draw[->,very thick] (9.0,0) -- (A);
\draw[->,very thick] (A) -- (12.6,0) node[right,font=\small] {العضلات};
\pic[scale=1.1] at (13.2,-0.5) {spring};
\pic[scale=1.15] at (0.4,0.85) {eyepic};
\pic[scale=1.05] at (1.35,0.85) {speaker};
% registers
\node[reg] (D) at (3.4,3.2) {\nt{DOPA}};
\node[reg] (C) at (5.6,3.2) {\nt{SERO}};
\node[reg] (N) at (7.8,3.2) {\nt{NORA}};
\node[reg] (H) at (10.0,3.2) {\nt{ACET}};
\node[font=\scriptsize,align=center,above=1pt] at (D.north) {الخطأ $\delta$};
\node[font=\scriptsize,align=center,above=1pt] at (C.north) {المستوى، $\tau_\gamma$};
\node[font=\scriptsize,align=center,above=1pt] at (N.north) {الكسب $g$};
\node[font=\scriptsize,align=center,above=1pt] at (H.north) {الكتابة $a$};
\draw[bc] (D.south) -- (3.4,1.35);
\draw[bc] (C.south) -- (5.6,1.35);
\draw[bc] (N.south) -- (7.8,1.35);
\draw[bc] (H.south) -- (10.0,2.1) -- (8.8,1.35);
\draw[bc] (H.south east) to[bend left=20] (A.north east);
\draw[->,thick] (2.85,1.35) -- (2.85,2.75) -- (D.west) node[pos=0.25,left,font=\scriptsize] {$V$};
\draw[->,thick,green!45!black] (0.4,3.2) node[left,black,font=\small] {المكافأة $r$} -- (2.2,3.2) -- (D.west);
\pic[scale=0.9] at (-0.45,3.7) {juice};
\draw[->,thick,densely dotted,gray!50!black] ([yshift=0.45cm]D.north) to[bend left=18] node[above,font=\scriptsize,black] {المفاجأة} ([yshift=0.45cm]N.north);
\end{tikzpicture}
\caption{الآلة بكاملها. ينقل ناقل \nt{GLUT} البيانات من المستشعرات إلى اختيار الفعل، ويتحكم \nt{GABA} في التدفق داخله. الدوائر المزدوجة قنوات البثّ؛ والأسهم المتقطعة بثُّ إشاراتها إلى الشبكة كلها، ولكل قناة مقبضها. يتلقى \nt{DOPA} المكافأة $r$ والتوقع $V$ من الناقد داخل الناقل (الفصل~\ref{sec:dofa})؛ والسهم المنقّط فرضيةُ أن \nt{NORA} يعرف بالمفاجأة من أخطاء كبيرة على غير العادة (الفصل~\ref{sec:nora}).}
\label{fig:machine}
\end{figure}

\section{صيغة واحدة لكل المقابض}

يمكن جمع المقابض في صيغة تخطيطية واحدة. ليست هذه نموذجًا جديدًا، بل خلاصة نماذج الفصول من~\ref{sec:glutcomputer} إلى~\ref{sec:acet}: كل رمز مأخوذ من فصله.
\begin{empheq}[box=\keybox]{equation}
\begin{aligned}
\tau\,\frac{\dd r_i}{\dd t} \;&=\; -r_i + F\Bigl(g\,\Bigl[\tfrac{1+a}{2}\,x_i + (1-a)\sum_j W_{ij}\,r_j - \sum_k M_{ik}\,q_k - \theta\Bigr]\Bigr),\\
\frac{\dd W_{ij}}{\dd t} \;&=\; \eta\,a\,\delta(t)\,e_{ij}(t),
\qquad \tau_e\,\frac{\dd e_{ij}}{\dd t} = -e_{ij} + r_i\,r_j,
\qquad \delta = r + \frac{\dd V}{\dd t} - \frac{V}{\tau_\gamma} .
\end{aligned}
\label{eq:machine}
\end{empheq}
هنا $r_i$ ترددات عصبونات \nt{GLUT}، و$x_i$ الدخل من المستشعرات، و$W_{ij}\ge0$ أوزانها المتبادلة (الفصل~\ref{sec:glutcomputer})؛ و$q_k$ ترددات عصبونات \nt{GABA}، و$M_{ik}\ge0$ شدة تثبيطها (الفصل~\ref{sec:gaba})؛ و$e_{ij}$ أثر الأهلية، و$\delta$ خطأ \nt{DOPA} (الفصل~\ref{sec:dofa})؛ و$\tau_\gamma$ الأفق الذي يحدده \nt{SERO} بحسب الفرضية؛ و$g$ كسب \nt{NORA} (الفصل~\ref{sec:nora})؛ و$a$ مستوى \nt{ACET} (الفصل~\ref{sec:acet}).

انظر إلى ما ليس موجودًا في~\eqref{eq:machine}. لا نبضة ساعة: كل شيء مكتوب بمعادلات تفاضلية، وكل عصبون يتغير من تلقاء نفسه. لا ذاكرة منفصلة: تتذكر الأوزان $W_{ij}$ نفسها التي يجري عليها الحساب، ويتذكر النشاط $r_i$ نفسه. لا تصحيحات فردية للغالبية العظمى من المشابك: تعلّمها حاصل ضرب أثر محلي في عدد واحد مشترك؛ ولا يوجد سلك خطأ مستقل إلى كل مخرج إلا في الدارة التي تتعلم الحركات (الفصل~\ref{sec:motorlearning}). وإشارات التحكم أربعة أنواع فقط، $\delta$ و$\tau_\gamma$ و$g$ و$a$، لثمانين مليار عصبون. وكلٌّ منها في الحقيقة ليس عددًا واحدًا، بل حزمة صغيرة من الخطوط المتوازية (القضية~\ref{prop:bundle})، لكن خطوط الحزمة تُعدّ بالآحاد.

\begin{keyblock}
\begin{proposition}[البنية]
\label{prop:architecture}
يمكن وصف آلة الدماغ، بقدر ما نفهمها اليوم، بثلاث طبقات. تسير البيانات على ناقل العناوين \nt{GLUT}. وتوجّه عصبونات \nt{GABA} المعنونة تدفقَ البيانات وتحدّه وتطبّعه. ويحدد نمطَ العمل عددٌ قليل من قنوات البث، كلٌّ منها حزمة صغيرة من الخطوط المتوازية المشتركة بين مناطق واسعة من الشبكة. الطبقتان الأوليان راسختان؛ أما توزيع الأدوار بين قنوات البث ففرضية في معظمه.
\end{proposition}
\end{keyblock}

\section{كل الأنواع في جدول واحد}

يجمع الجدول~\ref{tab:synthesis} كل أنواع العصبونات في الكتاب. الأنواع الأربعة التي لم يُخصَّص لها فصل نال كلٌّ منها سطرًا فيه؛ ووجد اثنان منها دورًا في فصول مخارج الآلة: \nt{GLYC} في حلقات الحبل الشوكي (الفصل~\ref{sec:muscles})، و\nt{ADRE} في المخرج الكيميائي (الفصل~\ref{sec:chemoutput}). أما دور الاثنين الآخرين في الحساب فإما بسيط وإما مجهول. والمواد التي لا تحدد نوع العصبون مجموعة في الملحق~\ref{sec:othermediators}.

\begin{table}[t]
\centering
\footnotesize
\setlength{\tabcolsep}{3pt}
\renewcommand{\arraystretch}{1.15}
\begin{tabular}{>{\raggedright}p{1.3cm}>{\raggedright}p{1.0cm}>{\raggedright}p{3.9cm}>{\raggedright}p{1.5cm}>{\raggedright}p{1.8cm}>{\raggedright\arraybackslash}p{3.8cm}}
\toprule
النوع & الفصل & ماذا يحسب & الزمن & الناقل & ما هو معروف \\
\midrule
\nt{GLUT} & \ref{sec:glutcomputer}, \ref{sec:code}, \ref{sec:sensors}, \ref{sec:motorlearning}, \ref{sec:dendrites} & البيانات: التزامنات، التأخيرات، المنطق، الذاكرة، التتابعات & \foreignlanguage{english}{ms} & معنون & راسخ \\
\nt{GABA} & \ref{sec:gaba}, \ref{sec:code}, \ref{sec:pain}, \ref{sec:motorlearning} & التحكم في التدفق: المنع، الاختيار، القسمة، الاستقرار، الإيقاع؛ بوابة الألم & \foreignlanguage{english}{ms} & معنون & راسخ؛ قسمة التردد بمعونة الضجيج الخلفي \\
\nt{SERO} & \ref{sec:serocomputer}, \ref{sec:pain}, \ref{sec:sleep} & منظّم المستوى، التنعيم؛ الأفق $\tau_\gamma$؛ مقبض صوت الألم؛ أنماط النوم & ثوانٍ & حجمي & التثبيط الذاتي مُقاس؛ الأفق فرضية \\
\nt{DOPA} & \ref{sec:dofa}, \ref{sec:dofaalt} & خطأ التنبؤ $\delta$؛ المستوى البطيء ثمنُ الوقت & $0.1\,\text{s}$ ودقائق & بثّ & $\delta$ مُقاس؛ التوسعات في الفصل~\ref{sec:dofaalt} \\
\nt{NORA} & \ref{sec:nora}, \ref{sec:pain}, \ref{sec:chemoutput}, \ref{sec:sleep} & الكسب $g$: البحث أم الحسم؛ المفاجأة؛ ”دوّاسة الوقود“ للأعضاء & ثوانٍ & بثّ & ارتفاع نسبة الإشارة إلى الضجيج مُقاس؛ الدور في البحث فرضية \\
\nt{ACET} & \ref{sec:acet}, \ref{sec:muscles}, \ref{sec:chemoutput}, \ref{sec:sleep} & الكتابة أو القراءة؛ سرعة التعلم؛ المخرج إلى العضلة؛ ”المكابح“ للأعضاء & ثوانٍ & كلاهما & ثلاثة تأثيرات مُقاسة؛ تبديل الأنماط فرضية \\
\nt{GLYC} & \ref{sec:muscles}, \ref{sec:pain} & وزن سالب في الحبل الشوكي وجذع الدماغ: منع العضلة المضادة & \foreignlanguage{english}{ms} & معنون & راسخ \\
\nt{HIST} & --- & إشارة عامة ”ابقَ يقظًا“ & بطيء & بثّ & دوره في الحساب لم يُدرس \\
\nt{ADRE} & \ref{sec:chemoutput} & ”دوّاسة الوقود“ للجسم كله عبر الدم؛ في الدماغ مجهول & بطيء & بثّ & خارج الدماغ راسخ؛ دوره في الدماغ غير واضح \\
\nt{ASPA} & --- & وزن موجب ضعيف & \foreignlanguage{english}{ms} & معنون & دوره موضع خلاف \\
\bottomrule
\end{tabular}
\caption{كل أنواع العصبونات في الكتاب: ماذا يحسب كلٌّ منها، وإلى أي حد نعرف ذلك.}
\label{tab:synthesis}
\end{table}

\section{كيف تعمل المقابض معًا}

حتى الآن كان كل فصل يدير مقبضًا واحدًا. فلنتتبع الآن حدثًا واحدًا عبر الآلة كلها (الشكل~\ref{fig:timeline}). تعلّم الحيوان منذ زمن أن الفعل $A$ يجلب العصير. وفي لحظة ما يتغير العالم: يكفّ $A$ عن جلب العصير، ويجلبه بدلًا منه الفعل $B$ الذي لا يكاد الحيوان يجربه.

\emph{الخطأ.} لم يأتِ العصير بعد $A$، فتجيب عصبونات \nt{DOPA} بهبوط، $\delta<0$ بحسب صيغة الخطأ~\eqref{eq:ctd}. والمشابك التي اختارت $A$ لتوّها تحمل أثرًا، فتضعف.

\emph{المفاجأة.} هبوط واحد مصادفة عادية. لكن الهبوطات تتكرر، وتصبح الأخطاء أكبر من المعتاد. وبحسب فرضية الفصل~\ref{sec:nora}، هذه بالذات إشارة \nt{NORA}: لقد تغيّر العالم. ينخفض الكسب $g$، ويصبح الاختيار ليّنًا، ويقود الضجيج إلى تجربة $B$ أكثر.

\emph{الكتابة.} بحسب فرضية الفصل~\ref{sec:acet}، يرتفع \nt{ACET} بعد التغيير وينقل الشبكة إلى نمط الكتابة: تضعف الوصلات الداخلية التي تحفظ العادة القديمة، ويقوى الدخل من المستشعرات، وتتغير الأوزان بسهولة.

\emph{عادة جديدة.} تجربة $B$ تجلب عصيرًا لم يتوقعه أحد: دفقة، $\delta>0$، فتقوى المشابك التي اختارت $B$. بضع محاولات كهذه، ويلحق التوقع $V$ بالعالم الجديد، وتعود الأخطاء صغيرة.

\emph{العودة.} مرّت المفاجأة، فيعيد \nt{NORA} الكسب العالي، ويعود الاختيار حادًّا؛ وينخفض \nt{ACET}، فتستعمل الشبكة في نمط القراءة ما تعلمته. وطوال هذا الوقت يحدد \nt{SERO}، بحسب الفرضية، الأفق $\tau_\gamma$: إلى أي مدى يستحق عصيرٌ بعيد أن يُنتظر.

\begin{figure}[t]
\centering
\begin{tikzpicture}[>=Stealth,x=1.1cm,y=0.95cm]
\foreach \y/\lab in {4/{$\delta$ (\nt{DOPA})},3/{$g$ (\nt{NORA})},2/{$a$ (\nt{ACET})},1/{اختيار $B$}}{
  \draw[gray!60!black] (0,\y-0.35) -- (10,\y-0.35);
  \node[left,font=\scriptsize] at (0,\y) {\lab};}
\draw[densely dashed,gray!60!black] (2,0.4) -- (2,4.55) node[above,font=\scriptsize,black] {تغيّر العالم};
\draw[densely dashed,gray!60!black] (5,0.4) -- (5,4.55) node[above,font=\scriptsize,black] {أول عصير مقابل $B$};
\pic[scale=0.85] at (2,5.25) {nojuice};
\pic[scale=0.85] at (5,5.25) {juice};
% delta: dips after the change, burst at the first juice for B, then back to zero
\draw[thick,red!70!black] (0,3.65) -- (2.3,3.65) -- (2.3,3.3) -- (2.5,3.3) -- (2.5,3.65) -- (3.2,3.65) -- (3.2,3.3) -- (3.4,3.3) -- (3.4,3.65) -- (4.1,3.65) -- (4.1,3.35) -- (4.3,3.35) -- (4.3,3.65) -- (5.0,3.65) -- (5.0,4.35) -- (5.2,4.35) -- (5.2,3.65) -- (5.8,3.65) -- (5.8,4.1) -- (6.0,4.1) -- (6.0,3.65) -- (6.6,3.65) -- (6.6,3.85) -- (6.8,3.85) -- (6.8,3.65) -- (10,3.65);
% gain: high, drops after surprise, recovers
\draw[thick,blue!60!black] (0,3.2) -- (3.0,3.2) -- (3.4,2.75) -- (5.8,2.75) -- (7.2,3.2) -- (10,3.2);
% ACh: low, rises after the change, falls after learning
\draw[thick,green!45!black] (0,1.75) -- (2.8,1.75) -- (3.3,2.25) -- (6.8,2.25) -- (7.8,1.75) -- (10,1.75);
% choice of B
\draw[thick,black] (0,0.7) -- (3.0,0.7) plot[domain=3:10,samples=40] (\x,{0.7+0.6/(1+exp(-(\x-5.3)*1.8))});
\draw[->,gray!70!black] (0,0.2) -- (10.2,0.2) node[below left,font=\scriptsize,black] {الزمن};
\end{tikzpicture}
\caption{حدث واحد متتبَّع عبر الآلة كلها. هذا مخطط نوعي مبني على نماذج الفصول من~\ref{sec:dofa} إلى~\ref{sec:acet} وفرضياتها، لا نتيجة حساب. بعد التغيير يجيب \nt{DOPA} بهبوطات؛ والهبوطات المتكررة، بحسب الفرضية، ترفع إشارة المفاجأة، فيخفض \nt{NORA} الكسب، ويشغّل \nt{ACET} الكتابة؛ ويعطي أول عصير مقابل $B$ دفقة، فينتقل الاختيار إلى $B$، وتعود المقابض إلى مواضعها.}
\label{fig:timeline}
\end{figure}

لاحظ أن أيّ مقبض لا يعرف ما تفعله المقابض الأخرى. كلٌّ منها يستجيب لعلاماته الخاصة، أي الخطأ وحجمه غير المعتاد وعدم اليقين، ويتكوّن ترتيب الأفعال من تلقاء نفسه، كما في دارة غير متزامنة ينطلق فيها كل مكوّن حين تجهز مداخله. أما هذا العمل المنسَّق بكامله فلم يُختبر بعد على حد علمنا: المقابض مقيسة كلٌّ على حدة، وعملها المشترك فرضية.

\section{بمَ تختلف هذه الآلة عن الحاسوب}

\begin{table}[t]
\centering
\footnotesize
\setlength{\tabcolsep}{4pt}
\renewcommand{\arraystretch}{1.15}
\begin{tabular}{>{\raggedright}p{2.2cm}>{\raggedright}p{4.2cm}>{\raggedright\arraybackslash}p{6.6cm}}
\toprule
& الحاسوب العادي & الدماغ \\
\midrule
الزمن & ساعة مشتركة، نحو جيجاهرتز & لا ساعة؛ كل عصبون يتغير بنفسه، والنوافذ تحددها الأحداث والإيقاعات المحلية (الفصول~\ref{sec:glutcomputer}، \ref{sec:gaba}، \ref{sec:code}) \\
العناصر & بوابات ثنائية موثوقة & عناصر عتبية تماثلية؛ يفقد المشبك معظم النبضات (الفصل~\ref{sec:code}) \\
علامة الوصلة & أيّ علامة & يحددها العصبون المرسل (الفصل~\ref{sec:neurontypes}) \\
الذاكرة & منفصلة عن المعالج & في الوصلات نفسها التي يجري فيها الحساب، وفي النشاط الذي يديم نفسه (الفصلان~\ref{sec:glutcomputer} و\ref{sec:acet}) \\
التعلّم & الانتشار العكسي للخطأ & قواعد محلية، وعدد بثّ واحد $\delta$ (الفصل~\ref{sec:dofa})، وأسلاك خطأ إلى مخارج دارة الحركة (الفصل~\ref{sec:motorlearning}) \\
التحكم & سجلات وناقل أوامر & أربع قنوات بثّ، كلٌّ منها حزمة من بضعة خطوط (الفصول~\ref{sec:serocomputer}، \babelsublr{\ref{sec:dofa}--\ref{sec:acet}}) \\
الطاقة & عشرات إلى مئات الواطات للمعالج الواحد & نحو $20$ واطًا للدماغ كله، بمعدل نبضة في الثانية تقريبًا لكل عصبون (الفصل~\ref{sec:code}) \\
\bottomrule
\end{tabular}
\caption{الدماغ والحاسوب العادي.}
\label{tab:computer}
\end{table}

يلخّص الجدول~\ref{tab:computer} الفروق الرئيسية. وليس أيٌّ منها عيبًا لم تجد الطبيعة وقتًا لإصلاحه، بل جزءًا من حل واحد. من دون ساعة مشتركة تلزم مستشعرات ”التشغيل والإطفاء“ وكواشف التزامن. والعناصر غير الموثوقة تحتاج إلى فائض من العصبونات الكثيرة. والذاكرة المدمجة في الحساب تحتاج إلى \nt{ACET} كي لا تفسد الكتابةُ القراءةَ. والتعلم بلا أسلاك راجعة يحتاج إلى \nt{DOPA} وإلى الضجيج. وحدّ الطاقة البالغ نبضة واحدة في الثانية يجبر اللغة كلها على التقتير وعلى إنفاق النبضات على الجديد وحده.

\section{ما لا نعرفه}

أصدق ما تُختم به هذه الخلاصة قائمةٌ بما لا نعرفه. وكل بند فيها مسألة لمهندس.

\emph{الشفرة.} نعرف سعة قناة النبضات، ونعرف أن المتلقي يختار أيّ جزء من الرسالة يقرأ (الفصل~\ref{sec:code}). لكن إلى أي حد تستعمل القشرة التوقيت الدقيق للنبضات، وهل للعصبونات ”كلمات“، فأمر مجهول.

\emph{التعلم عبر طبقات كثيرة.} إشارة البث تعلّم ببطء، وتزداد بطئًا مع كل عصبون يؤثر في النتيجة (القضية~\ref{prop:broadcastcost}). فكيف تضبط القشرة إذن مليارات الأوزان في دارات متعددة الطبقات؟ لا نعرف؛ وثمة نماذج تنقل فيها رشقاتُ النبضات الخطأَ بين الطبقات~\cite{Payeur2021}. وربما يكمن جزء من الجواب في الحسابات داخل فروع العصبون نفسه، حيث يتصرف عصبون واحد كشبكة صغيرة متعددة الطبقات: لتكرار استجابة نموذج عصبون قشري واحد تحتاج الشبكة الاصطناعية إلى خمس إلى ثماني طبقات~\cite{Beniaguev2021}، وفروع عصبونات الإنسان قادرة حتى على حساب ”أو الحصرية“~\cite{Gidon2020}. ومرشح آخر هو التعلم الذاتي: إذا تعلمت القشرة أن تتنبأ بمداخلها، نشأ الخطأ في موضعه، في كل طبقة، ولم تعد هناك حاجة إلى إشارة مشتركة (الفصل~\ref{sec:dofa})~\cite{Keller2018}.

\emph{ما الذي تنقله قنوات البث حقًّا.} لـ\nt{DOPA} صورة معيارية وستة توسعات لها (الفصل~\ref{sec:dofaalt})؛ ولـ\nt{NORA} و\nt{ACET} فرضيات معقولة (الفصلان~\ref{sec:nora} و\ref{sec:acet})؛ أما \nt{SERO} فمعظم ما لدينا عنه تخمينات.

\emph{التنسيق في الزمن.} رأينا كيف نحسب دالة، وكيف نقيس الزمن من حدث، وكيف نقطّع التدفق إلى نوافذ بإيقاع محلي. لكن كيف تنسّق شبكة هائلة بلا ساعة مشتركة عمل مناطق متباعدة، فأمر لا نفهمه إلا جزئيًا.

\emph{الطاقة.} عشرون واطًا لكل شيء. نفهم لماذا يجب أن تكون الشفرة متناثرة (القضية~\ref{prop:economy})، لكن لا أحد يعرف حتى الآن كيف يبني آلة تفعل الشيء نفسه بالقدرة نفسها~\cite{MehonicKenyon2022}.

الدماغ آلة حاسبة لم يجمعها مهندس، لكنها مبنية وفق قوانين يعرفها أي مهندس: دارات \foreignlanguage{english}{RC}، والعتبات، والتغذية الراجعة، والمرشحات، والنواقل، وسجلات البث. لم نقرأ مخططها إلا جزئيًا. وسيكمل قراءته من يجيد قراءة المخططات.

\section{ما التالي في الكتاب}

جمع هذا الفصل النواة الحسابية للآلة. أما بقية الفصول فتخرج عن حدودها. الفصلان~\ref{sec:sensors} و\ref{sec:recognition} عن المداخل: كيف تحوّل المستشعرات الضوء والصوت والجزيئات إلى نبضات، وكيف تتعرّف الآلة فيها على الأشياء. والفصول من~\ref{sec:muscles} إلى~\ref{sec:chemoutput} عن المخارج وما يخدمها: العضلات، والألم بوصفه إشارة إنذار، وتعلّم الحركات عبر أسلاك خطأ مستقلة، والمخرج الكيميائي البطيء إلى الجسم. والفصل~\ref{sec:debugging} عن تصحيح أخطاء الآلة من الخارج، ومن ذلك واجهات الدماغ والحاسوب. والفصلان~\ref{sec:eetypes} و\ref{sec:dendrites} يصنّفان العصبونات مرة أخرى، هذه المرة بحسب وظيفتها الكهربائية وعدد مداخلها. والفصل~\ref{sec:sleep} عمّا تفعله الآلة حين تنفصل عن العالم، والفصلان~\ref{sec:livingmachine} و\ref{sec:dishbrain} عن آلات مجمّعة من عصبونات حية على رقاقة.

\section{تمارين}

تمارين هذا الفصل تربط نتائج فصول مختلفة: كل تمرين يستند إلى فصلين على الأقل.

\begin{exercise}[\lvlA]
\label{ex:10:1}
\emph{الناقل والسجلات.} قنوات البث الأربع حزم من خمسة خطوط تقريبًا لكلٍّ منها (القضية~\ref{prop:bundle})؛ كل خط يتغير مرة في الثانية ويمكن تمييز أربعة مستويات عليه. كم سنة يحتاج التحكم لينقل ما ينقله ناقل \nt{GLUT} ($8.6\cdot10^{10}$ عصبون بالتردد~\eqref{eq:energy}، ودقة $1\ms$ بحسب~\eqref{eq:spikeentropy}) في ثانية واحدة؟
\end{exercise}

\begin{exercise}[\lvlB]
\label{ex:10:2}
\emph{مقبضان على حلقة واحدة.} في الحلقة~\eqref{eq:ei} تعمل مقابض~\eqref{eq:machine} على الإثارة: $\tau_E\dot E=-E+g[(1-a)w_{EE}E-w_{EI}I+h]$، ولا تتحكم في \nt{GABA}. بأرقام الشكل~\ref{fig:ei} ترفع رشقةُ \nt{NORA} الكسبَ $g$ إلى $6$. ما أصغر $a$ تكون عنده الحلقة مستقرة؟ وهل يمكن إدارة مقبضي \nt{NORA} و\nt{ACET} كلٍّ على حدة؟
\end{exercise}

\begin{exercise}[\lvlB]
\label{ex:10:3}
\emph{من أين يأتي ثمن البث.} مخارج $N$ عصبونًا $y_k=f(u_k)+\xi_k$ بضجيج غاوسي مستقل تباينه $\sigma^2$؛ و$R\approx R_0+\sum_kR'_k\xi_k$ بقيم $R'_k$ متساوية، ويبث \nt{DOPA} القيمة $\delta=R-R_0$. أوجد نسبة الإشارة إلى الضجيج للتصحيح $\delta\,\xi_jx_j$ في محاولة واحدة. كيف يزداد عدد المحاولات مع $N$؟
\end{exercise}

\begin{exercise}[\lvlB]
\label{ex:10:4}
\emph{ربط الصوت بالومضة.} يصل الصوت إلى الناقل بعد $5\ms$، والومضة بعد $30\ms$ مضافًا إليها كمون النبضة الأولى~\eqref{eq:latency} ($\tau=10\ms$، و$U$ من $1.2\theta$ إلى $4\theta$)؛ والخلفية على كل مدخل $5$ نبضات في الثانية. اختر تأخير خط الصوت وأصغر نافذة تزامن تصلح لكل التباينات. كم ربطًا كاذبًا في الثانية ستعطي الخلفية؟
\end{exercise}

\begin{exercise}[\lvlB]
\label{ex:10:5}
\emph{نمذجة: من يلاحظ التغيير.} يرى الناقد $y=s+\text{ضجيج}$ ($R=1$)؛ وباحتمال $1/200$ تقفز $s$ إلى قيمة جديدة تباينها $J^2=4$. نمذج مرشح كالمان بـ$q=0$ يرفع $P$ إلى $J^2$ حين تتجاوز $E\leftarrow0.9E+\delta^2/(P+R)$ القيمة $20$. بكم يقل خطؤه عن خطأ أفضل $\alpha$ ثابت؟
\end{exercise}

\begin{exercise}[\lvlA]
\label{ex:10:6}
\emph{كم محاولة لتغيير العادة.} تعلّمت الشبكة $Q_A=0.8$ و$Q_B=0.2$، وتختار بحسب~\eqref{eq:softmax} مع $\sigma=1$ و$g=8$. الآن يعطي $A$ العصير باحتمال $0.2$؛ $Q_A\leftarrow Q_A+\alpha(r-Q_A)$، و$Q_B$ لا يتغير. بعد كم اختيارًا لـ$A$ ينخفض احتمال اختيار $A$ إلى $0.6$ عند $\alpha=0.1$ و$0.5$؟ وماذا لو خفّض \nt{NORA} الكسب $g$ إلى $2$؟
\end{exercise}

\begin{exercise}[\lvlC]
\label{ex:10:7}
\emph{حزمة أسلاك بدل سلك.} المخارج $y_k=w_k+\xi_k$، والمكافأة $R=-\sum_ky_k^2$؛ يبث كل خط من خطوط الحزمة إلى مجموعة من $G$ عصبونًا خطأَ مخارجها، $w_k\leftarrow w_k+\eta\,\delta_l\xi_k$. بيّن أنه عند أفضل $\eta$ يلزم $\approx(G+2)\ln(1/\varepsilon)$ محاولة للوصول إلى $\sum w_k^2=\varepsilon\sum w_k^2(0)$. كم تتعلم دارة من $10^6$ عصبون على خمسة خطوط عند $\varepsilon=0.01$، بمحاولة كل ثلاث ثوانٍ؟
\end{exercise}

\chapter{مداخل الآلة: كيف توصَل بها العيون والآذان والمستشعرات الأخرى}
\chaptermark{مداخل الآلة}
\label{sec:sensors}
\begin{chaptergoals}
\item كيف تعمل كل حاسة محوِّلًا من مستقبِل بوابة وضبط للكسب ومرمّز للنبضات؛
\item لماذا يحدّ ضجيج الطلقات للفوتونات الرؤيةَ في الظلام، ولماذا نرى في الغسق أبطأ؛
\item كيف يحشر الضبط التلقائي للكسب مدى هائلًا في النبضات، فتنقل المستشعرات التغيرات؛
\item كيف تضغط العين صورتها نحو مئة مرة، وكيف تعمل الأذن بنكَ مرشحات.
\end{chaptergoals}

\section{المستشعر محوِّلًا}

في الشكل~\ref{fig:machine} كانت البيانات تدخل الآلة من اليسار، من كتلة ”المستشعرات“. فلنفتح الآن هذه الكتلة. كل عضو حسّ عند المهندس \emph{محوِّل} له طور دخل تماثلي وفي آخره محوِّل تماثلي رقمي، إلا أن الأرقام في المخرج ليست أعدادًا، بل نبضات.

المخطط واحد في كل أعضاء الحس (الشكل~\ref{fig:transducer}). كمية فيزيائية، ضوءًا كانت أو ضغطًا أو اهتزازًا أو جزيئًا، تؤثر في بروتين مستقبِل في غشاء الخلية الحسية. يعمل هذا البروتين عمل البوابة: يفتح قناة أيونية بنفسه أو عبر سلسلة قصيرة من التفاعلات. فينشأ \emph{تيار المستقبِل}، ومعه جهد تماثلي على الغشاء. ثم يمر الجهد بضبط الكسب ويصل إلى مرمّز النبضات، وهو العصبون المكامل المعروف لنا~\eqref{eq:glutneuron}، الذي يحوّل المستوى إلى تردد وإلى لحظات نبضات. والمخرج واحد في الغالب: العصبونات الحسية في العين والأذن والشم والجلد تطلق الغلوتامات، فتتصل المداخل مباشرة بناقل \nt{GLUT}.

\begin{figure}[t]
\centering
\begin{tikzpicture}[>=Stealth,
  blk/.style={draw,very thick,rounded corners=2pt,align=center,font=\footnotesize,minimum height=1.0cm,text width=1.9cm,fill=blue!5}]
\node[align=center,font=\small] (P) at (0.1,0) {ضوء، صوت،\\ضغط،\\جزيء};
\node[blk] (R) at (3.0,0) {المستقبِل\\البوابة};
\node[blk] (G) at (6.5,0) {ضبط\\الكسب};
\node[blk] (E) at (10.0,0) {مرمّز\\النبضات};
\node[align=center,font=\small] (B) at (13.4,0) {ناقل\\\nt{GLUT}};
\draw[->,very thick] (P) -- (R);
\foreach \ic/\xx in {sun/-1.1,speaker/-0.3,press/0.5,molecule/1.3}{\pic[scale=0.85] at (\xx,1.05) {\ic};}
\draw[->,very thick] (R) -- node[above,font=\scriptsize] {تيار} (G);
\draw[->,very thick] (G) -- node[above,font=\scriptsize] {مستوى} (E);
\draw[->,very thick,red!70!black] (E) -- node[above,font=\scriptsize,black] {نبضات} (B);
\draw[decorate,decoration={brace,amplitude=5pt,mirror},gray!60!black] (1.8,-0.8) -- (7.7,-0.8) node[midway,below=5pt,font=\scriptsize,black] {الجزء التماثلي};
\draw[decorate,decoration={brace,amplitude=5pt,mirror},gray!60!black] (8.8,-0.8) -- (11.2,-0.8) node[midway,below=5pt,font=\scriptsize,black] {نبضات};
\end{tikzpicture}
\caption{أي عضو حسّ سلسلةً من التحويلات. يفتح البروتين المستقبِل قناة فيعطي تيارًا؛ ويكيّفه ضبط الكسب مع الظروف؛ ويحوّل المرمّز~\eqref{eq:glutneuron} المستوى إلى نبضات تخرج إلى ناقل \nt{GLUT}. في العين والأذن لا تُصدر الأطوار الأولى نبضات أصلًا: تبقى الإشارة تماثلية إلى أن يلزم نقلها بعيدًا.}
\label{fig:transducer}
\end{figure}

تفصيل يعرفه مهندس الإلكترونيات جيدًا. في العين والأذن لا تُصدر الخلايا الأولى نبضات إطلاقًا: إنها تنقل إلى جاراتها جهدًا متدرجًا، كطور تماثلي على لوحة دارة. ولا تظهر النبضات إلا حيث يجب نقل الإشارة بعيدًا. فالإشارة التماثلية تخمد وتتلوث بالضجيج على مسافة ملليمترات، أما النبضة، كما رأينا في الفصل~\ref{sec:neurontypes}، فتصل إلى نهاية السلك بالارتفاع نفسه. والرقمنة تجري حيث يبدأ الخط الطويل.

\section{عند حدود الفيزياء}

تعمل مستشعرات الدماغ عند الحدود الفيزيائية للحساسية. مستشعر الضوء في الظلام يستجيب \emph{لفوتون واحد}: امتصاص جسيم ضوء واحد يعطي تيارًا نحو بيكوأمبير، يمكن تمييزه من ضجيج المستشعر نفسه~\cite{Baylor1979}. ومستشعر الصوت يلاحظ إزاحة حزمته الحساسة بأقل من نانومتر~\cite{Hudspeth1989}.

حين يقل الضوء لا يحدّ الدقةَ المستشعرُ، بل الضوء نفسه: تصل الفوتونات عشوائيًا في تدفق بواسوني. إذا دخل المستشعر خلال زمن التجميع $T$ في المتوسط $N=\Phi T$ فوتونًا، تذبذب عددها بمقدار $\sqrt N$. ولكي تُلاحظ زيادة $\Delta N$ يجب أن تكون أكبر من هذا التذبذب بعدة مرات، فتكون أصغر نسبة لتغير السطوع يمكن تمييزها
\begin{empheq}[box=\keybox]{equation}
\frac{\Delta I}{I} \;\approx\; \frac{k}{\sqrt{\Phi\,T}} .
\label{eq:photonnoise}
\end{empheq}
هذه هي الصيغة نفسها التي رأيناها لعدد النبضات~\eqref{eq:rateerror}: ضجيج الطلقات للفوتونات وضجيج الطلقات للنبضات يخضعان لقانون واحد. وفي الظلام لا تملك العين إلا طريقة واحدة لتحسين الدقة، هي أن تجمع الفوتونات مدة أطول، وزمن التجميع يطول فعلًا حتى أعشار الثانية. لهذا نرى في الغسق أبطأ.

\section{المدى الديناميكي: ضبط الكسب}

أصعب مسألة هندسية في المستشعرات هي المدى. تتغير الإضاءة من ليلة مرصعة بالنجوم إلى ثلج تحت الشمس بنحو عشر رتب عشرية، وشدة الصوت من عتبة السمع إلى عتبة الألم باثنتي عشرة رتبة. أما تردد النبضات فيتغير من الصفر إلى بضع مئات في الثانية، أي بأقل من ثلاث رتب. ولا يمكن حشر دخل كهذا في مخرج كهذا خطيًا.

والحل هو حل أي مستقبِل لاسلكي: \emph{الضبط التلقائي للكسب}. توصف استجابات الخلايا الحسية في العين وصفًا جيدًا بالصيغة
\begin{empheq}[box=\keybox]{equation}
r \;=\; r_{\max}\,\frac{I}{I+\sigma} ,
\label{eq:nakarushton}
\end{empheq}
حيث $\sigma$ السطوع الذي تبلغ عنده الاستجابة نصف قيمتها العظمى~\cite{NakaRushton1966}. والصيغة~\eqref{eq:nakarushton} وحدها لا تغطي إلا رتبتين أو ثلاثًا. لكن $\sigma$ ليست ثابتة: عند العمل طويلًا على خلفية ساطعة تزداد تبعًا لمتوسط السطوع. فينزاح منحنى الاستجابة على المحور اللوغاريتمي للسطوع، ويضع جزأه الشديد الانحدار في كل مرة حيث يقع الدخل الآن (الشكل~\ref{fig:adaptation}). هذه هي القسمة نفسها على النشاط الكلي التي رأيناها في التثبيط التحويلي في الفصل~\ref{sec:gaba}~\cite{Carandini2012}، إلا أن المقام هنا متوسط السطوع خلال الثواني الأخيرة.

\begin{figure}[t]
\centering
\begin{tikzpicture}[>=Stealth,x=1.25cm,y=2.8cm]
\draw[->,gray!70!black] (-2,0) -- (6.4,0) node[below left,font=\small,black] {$\log_{10} I$};
\draw[->,gray!70!black] (-2,0) -- (-2,1.15) node[left,font=\small,black] {$r/r_{\max}$};
\foreach \u in {-2,0,2,4,6}{\draw[gray!60!black] (\u,-0.02) -- (\u,0.02); \node[below,font=\scriptsize] at (\u,0) {$\u$};}
\draw[densely dashed,gray!60!black] (-2,0.5) -- (6.2,0.5);
\foreach \s/\c/\lab in {0/blue!60!black/{خلفية مظلمة},2/green!45!black/{متوسطة},4/red!70!black/{ساطعة}}{
  \pgfmathsetmacro{\lo}{max(-2,\s-3)}
  \draw[very thick,\c] (-2,0) -- (\lo,0) plot[domain=\lo:6.2,samples=80] (\x,{1/(1+10^(\s-\x))});
  \node[font=\scriptsize,\c,anchor=west] at (\s+0.15,0.42) {\lab};}
\foreach \s/\ic in {0/moon,2/cloud,4/sun}{\pic at (\s+0.6,0.64) {\ic};}
\end{tikzpicture}
\caption{ضبط كسب مستشعر الضوء بحسب~\eqref{eq:nakarushton}. كل منحنى يغطي رتبتين أو ثلاثًا من السطوع؛ وعند الانتقال إلى خلفية أسطع تزداد $\sigma$ فينزاح المنحنى على المحور اللوغاريتمي. والمنحنيات المنزاحة معًا تغطي المدى كله.}
\label{fig:adaptation}
\end{figure}

لهذا الضبط ثمن نعرفه من الفصل~\ref{sec:code}: لا يُبلغ المستشعر عن السطوع، بل عن السطوع \emph{بالنسبة إلى الخلفية}. الدخل الثابت يكفّ تدريجيًا عن استدعاء الاستجابة، أما كل تغيير فيستدعيها. فمداخل الآلة تتحدث منذ البداية عن التغيرات لا عن المستويات، وفي هذا الاقتصاد نفسه في النبضات الذي في القضية~\ref{prop:economy}~\cite{Barlow1961}. ومن هنا أيضًا مستشعرات التشغيل والإطفاء في الفصل~\ref{sec:glutcomputer}~\cite{Schiller1992}: بعضها يُبلغ أن الضوء ازداد، وبعضها أنه خفت.

وقد كرر المهندسون هذه الحيلة في \emph{كاميرات الأحداث}. لا تلتقط هذه الكاميرا إطارات وفق ساعة: كل بكسل فيها يُصدر بنفسه، بلا ساعة مشتركة، حدثًا ”أسطع“ أو ”أعتم“ حين يتغير لوغاريتم السطوع بنسبة محددة~\cite{Lichtsteiner2008}. هذه بالضبط شفرة السلكين للتشغيل والإطفاء من الفصل~\ref{sec:glutcomputer} مجسَّدة في السيليكون، وهي تمنح الكاميرا مدى وسرعة لا تبلغهما المصفوفة العادية~\cite{Gallego2022}.

\section{العين: ضغط حتى يتسع له السلك}

في العين نحو $10^8$ مستشعر ضوء~\cite{Curcio1990}، أما عصبونات المخرج التي تنقل الصورة إلى الدماغ فنحو مليون. وقبل أن تغادر الإشارة العين تُضغط نحو مئة مرة. وأهم حيل الضغط معروفة لنا. كل عصبون مخرج يستجيب للفرق بين السطوع في بقعة صغيرة والسطوع حولها، وهذا طرح الجيران عبر التثبيط كما في الفصل~\ref{sec:gaba}. المساحات المستوية من الصورة لا تكاد تكلف شيئًا، أما الحواف والتغيرات فتُنقل. وعصبونات المخرج المختلفة متخصصة: بعضها يُبلغ عن الإضاءة، وبعضها عن الإعتام، وبعضها عن الحركة في اتجاه معيّن (الشكل~\ref{fig:direction})؛ ولدى الفأر أُحصي أكثر من ثلاثين قناة مخرج كهذه~\cite{Baden2016}.

فما عرض نطاق الكابل الناتج؟ لدى خنزير غينيا قيس من تسجيلات عصبونات المخرج نحو $875$ ألف بت في الثانية لنحو $\sim10^5$ من هذه العصبونات؛ والتحويل إلى مليون عصبون مخرج لدى الإنسان يعطي رتبة $10^7$ بت في الثانية~\cite{Koch2006}، أي ما يعادل كابل شبكة قديمًا بسرعة عشرة ميغابت. وبتردد متوسط يبلغ بضع نبضات في الثانية لكل عصبون يكون ذلك بضع بتات لكل نبضة، كما خرج لنا في الفصل~\ref{sec:code}.

\section{الأذن: بنك مرشحات}

تحل الأذن مسألة أخرى: يجب تحليل الصوت إلى تردداته. كان المهندس سيضع بنكًا من مرشحات تمرير النطاق، والأذن تفعل ذلك بالضبط ميكانيكيًا. يجري الاهتزاز الصوتي على حاجز مرن تتغير خصائصه تدريجيًا على طوله، فيرنّ كل جزء منه على تردده. والمستشعرات الجالسة على جزء ما لا تستجيب إلا لنطاقها (الشكل~\ref{fig:filterbank}). وتتوزع الترددات على المواضع توزيعًا لوغاريتميًا تقريبًا: كل أوكتاف ينال حصة متقاربة من المستشعرات. ورقم المستشعر الذي أطلق هو التردد، وهذه شفرة الموضع من الفصل~\ref{sec:code}.

\begin{figure}[t]
\centering
\begin{tikzpicture}[>=Stealth,x=1.35cm,y=2.2cm]
\draw[->,gray!70!black] (-0.3,0) -- (7.6,0) node[above left,font=\small,black] {التردد، \foreignlanguage{english}{Hz}};
\draw[->,gray!70!black] (-0.3,0) -- (-0.3,1.15) node[left,font=\small,black] {الاستجابة};
\foreach \k/\f in {0/125,1/250,2/500,3/1000,4/2000,5/4000,6/8000}{
  \draw[gray!60!black] (\k,-0.02) -- (\k,0.02); \node[below,font=\scriptsize] at (\k,0) {\f};
  \draw[thick,blue!60!black] plot[domain={max(\k-1.2,-0.3)}:{\k+0.7},samples=60] (\x,{1/(1+((\x-\k)/(\x<\k?0.45:0.2))^4)});}
% a piano keyboard: seven white keys per octave, like the octaves on the axis
\foreach \i in {0,...,48}{\draw[gray!50!black,fill=white,line width=0.3pt] ({\i/7},-0.44) rectangle ({(\i+1)/7},-0.26);}
\foreach \o in {0,...,6}{\foreach \j in {0,1,3,4,5}{\fill[black] ({\o+(\j+1)/7-0.035},-0.35) rectangle ({\o+(\j+1)/7+0.035},-0.26);}}
\end{tikzpicture}
\caption{الأذن بنكًا من مرشحات تمرير النطاق، مخطط. كل منحنى استجابةُ مستشعرات جزء واحد لصوت بترددات مختلفة؛ والترددات المركزية متباعدة بأوكتاف، كمفاتيح على مقياس لوغاريتمي. للمرشحات الحقيقية منحدر شديد في جهة الترددات العالية ومنحدر لطيف في جهة المنخفضة. وفي الأسفل لوحة مفاتيح: عليها، كما في الأذن، ينال كل أوكتاف المساحة نفسها.}
\label{fig:filterbank}
\end{figure}

مرنانات الأذن ليست منفعلة. فبعض المستشعرات يهزّ الحاجز بنفسه بإيقاع الصوت، ويضخّم الاهتزازات الضعيفة آلاف المرات في السعة (بمقدار $66\text{--}76\,\text{dB}$)، ومع ازدياد الشدة ينخفض الكسب~\cite{Ruggero1997,Robles2001}. هذا عند المهندس مضخّم بتغذية راجعة موجبة موضوع عند حافة الإثارة الذاتية تمامًا، وهي الحياة نفسها على الحافة التي رأيناها في شبكة الفصل~\ref{sec:gaba}: عند الحافة يكون النظام حساسًا للإشارات الصغيرة على نحو خاص. وضغط الشدة يؤديه المضخّم نفسه: الهادئ يُضخَّم كثيرًا، والصاخب قليلًا.

وللأذن شفرة ثانية. في الترددات المنخفضة تسير نبضات العصبونات في طور مع الصوت: يطلق العصبون عند جزء معيّن من كل موجة، وإن لم يطلق عند كل موجة. لدى خنزير غينيا يبدأ التقيد بالطور بالضعف قرب $600\,\text{Hz}$، ويبقى ملحوظًا حتى نحو $3.5\,\text{kHz}$~\cite{PalmerRussell1986}. هكذا يُحفظ في لحظات النبضات التوقيت الدقيق للصوت، ومنه الفرق بين وصول الصوت إلى الأذنين، الذي تجد به شبكة الفصل~\ref{sec:glutcomputer} الاتجاه~\cite{Grothe2010}. وفي الترددات العالية لا تلحق النبضات بالموجة، فلا تبقى إلا شفرة الموضع. فالأذن الواحدة تستعمل شفرتي الفصل~\ref{sec:code} كلتيهما: التوقيت حيث تلحق العصبونات، والموضع حيث لا تلحق.

\section{بقية المستشعرات}

\begin{table}[t]
\centering
\footnotesize
\setlength{\tabcolsep}{3pt}
\renewcommand{\arraystretch}{1.15}
\begin{tabular}{>{\raggedright}p{1.9cm}>{\raggedright}p{2.6cm}>{\raggedright}p{4.0cm}>{\raggedright\arraybackslash}p{4.6cm}}
\toprule
الحاسة & ماذا تقيس & كيف بُني المدخل & الحيلة من الكتاب \\
\midrule
البصر & الضوء & $\sim10^8$ مستشعر، ضغط $\sim100$ مرة، $\sim10^7$ بت في الثانية & ضبط الكسب، طرح الجيران، سلكان، اتجاه الحركة \\
السمع & ضغط الهواء & بنك مرشحات ميكانيكي، مضخّم فعّال & شفرة الموضع وشفرة التوقيت، التأخيرات \\
اللمس & الضغط، الاهتزاز & مستشعرات سريعة التعوّد وبطيئته & زوج ”مرشح تمرير عالٍ ومرشح تمرير منخفض“ \\
الحرارة & الدفء & مستشعرات منفصلة للدفء وللبرد & شفرة السلكين \\
الشم & الجزيئات & $\sim400$ نوع من المستقبِلات، ولكل مستشعر نوع واحد & شفرة تركيبية: الرائحة مجموعة الأنواع التي أطلقت \\
الذوق & مواد الطعام & خمس صفات: حلو، مرّ، مالح، حامض، أومامي & خطوط موسومة؛ بعض الخلايا يطلق \foreignlanguage{english}{ATP} أو السيروتونين لا الغلوتامات \\
وضع الجسم & تمدد العضلات & مستشعرات الطول وسرعة التمدد & تغذية راجعة لمنظّم الحركة \\
\bottomrule
\end{tabular}
\caption{كيف توصَل المستشعرات. كل الأجهزة في العمود الثالث مقيسة؛ والعمود الأخير يربطها بحيل الفصول السابقة.}
\label{tab:sensors}
\end{table}

بقية الحواس مجموعة في الجدول~\ref{tab:sensors}، وفي كل سطر منه تقريبًا نتعرف على حيلة من الفصول السابقة. يستعمل اللمس زوجًا من المرشحات: بعض مستشعرات الضغط يستجيب للمسة ثم يصمت فورًا، وبعضها يُبقي استجابته ما دام الضغط قائمًا. وتنقل الحرارةَ سلكان، أحدهما للدفء والآخر للبرد. وأغرب المداخل الشم. لدى الإنسان نحو أربعمئة نوع من المستقبِلات الشمية، ولا يحمل كل عصبون شمي إلا نوعًا واحدًا~\cite{Mombaerts2004}. والرائحة لا تنشّط نوعًا واحدًا بل مجموعة، والرسالة هي \emph{أيّ} الأنواع أطلقت. هذا عند المهندس متجه ثنائي طوله نحو أربعمئة، وعدد قيمه الممكنة فلكي.

\begin{keyblock}
\begin{proposition}[مداخل الآلة]
\label{prop:sensors}
كل حاسة محوِّل يعمل عند الحد الفيزيائي للحساسية، ويضغط مدى ديناميكيًا هائلًا بالضبط التلقائي للكسب، وينقل على ناقل \nt{GLUT} \emph{التغيرات} قبل كل شيء. ولهذا تستعمل المستشعرات حيل الآلة نفسها: شفرة السلكين، وشفرة الموضع، وشفرة التوقيت، والقسمة على الخلفية. بنية المستشعرات مقيسة؛ أما أن شفراتها مضبوطة على أكبر قدر من المعلومات لكل نبضة ففرضية راسخة الأسس.
\end{proposition}
\end{keyblock}

تعمل المداخل، ككل شيء آخر، بلا تزامن. كل مستشعر يُصدر نبضة حين يكون لديه ما يبلّغ عنه، والحواس المختلفة تصل بتأخيرات مختلفة: الصوت بعد ملّي ثوانٍ، والضوء بعد عشرات الملّي ثانية. أما كيف تجمعها الآلة في صورة واحدة للعالم من غير ساعة مشتركة، فهي إحدى مسائل التنسيق في الزمن من الفصل~\ref{sec:synthesis}.

\section{تمارين}

\begin{exercise}[\lvlA]
\label{ex:11:1}
\emph{كم نجمع من الفوتونات.} في الغسق يدخل المستشعرَ $100$ فوتون في الثانية؛ والزيادة ملحوظة إذا كانت أكبر من التذبذب~\eqref{eq:photonnoise} بـ$3$ مرات. (أ)~كم يحتاج مستشعر واحد من الوقت ليلاحظ تغير السطوع بـ$10\,\%$؟ (ب)~كم مستشعرًا يجب جمعها لينجز ذلك في $0.2\,\text{s}$، وبكم مرة تنخفض الدقة المكانية على المحور؟
\end{exercise}

\begin{exercise}[\lvlA]
\label{ex:11:2}
\emph{كم بتًا في النبضة.} (أ)~تنقل العين $\sim10^7$ بت في الثانية عبر $10^6$ عصبون مخرج بتردد من $3$ إلى $5$ نبضات في الثانية. ما النسبة التي تستعملها من السعة~\eqref{eq:spikeentropy} عند $\Delta t=1\ms$؟ (ب)~رائحة تشغّل $20$ من $\approx400$ نوع مستقبِل. كم بتًا تحمل هذه المجموعة، وكم نوعًا مشتركًا في المتوسط بين رائحتين عشوائيتين؟
\end{exercise}

\begin{exercise}[\lvlB]
\label{ex:11:3}
\emph{ما الذي يقدر عليه منحنى واحد~\eqref{eq:nakarushton}.} (أ)~في أي مدى من السطوع تزداد الاستجابة من $10\,\%$ إلى $90\,\%$ من $r_{\max}$؟ كم رتبة هذا؟ (ب)~كم موضعًا للمنحنى المنزاح يلزم لتغطية عشر رتب من السطوع؟ واثنتي عشرة رتبة من الشدة الصوتية؟
\end{exercise}

\begin{exercise}[\lvlB]
\label{ex:11:4}
\emph{حد شفرة التوقيت في الأذن.} الفرق بين وصول الصوت إلى الأذنين يصل إلى $0.22\,\text{m}/343\,\text{m/s}$. (أ)~النبضات تسير في طور مع النغمة: حتى أي تردد يتحدد الاتجاه بها بلا لبس؟ (ب)~ترتجف النبضة بمقدار $0.5\ms$، والمطلوب تمييز $20\,\text{\textmu s}$. كم عصبونًا بمعدل $100$ نبضة في الثانية يجب أن نأخذ متوسطها خلال $50\ms$؟
\end{exercise}

\begin{exercise}[\lvlB]
\label{ex:11:5}
\emph{كاميرا أحداث على كابل العين.} كاميرا من $10^6$ بكسل بمخرج $10^7$ بت في الثانية ترسل حدثًا (العنوان والإشارة) حين يتغير $\ln I$ في البكسل بمقدار $\ln1.2$. ما أكبر سرعة يمكن أن تجري بها حافة طولها $500$ بكسل وفرق السطوع عبرها $4$ أضعاف؟ وكم إطارًا في الثانية تعطي كاميرا عادية عبر المخرج نفسه بـ$8$ بتات لكل بكسل؟
\end{exercise}

\begin{exercise}[\lvlB]
\label{ex:11:6}
\emph{نمذجة: ضبط الكسب.} يضبط المستشعر~\eqref{eq:nakarushton} القيمة $\sigma$ بـ$\tau=1\,\text{s}$ لوغاريتميًا، $\tau\,\dd\ln\sigma/\dd t=\ln I-\ln\sigma$، أو خطيًا، $\tau\,\dd\sigma/\dd t=I-\sigma$؛ والخلفية تقفز $1\to100\to1$. بعد كم ثانية تعود الاستجابة لاختبار $+10\,\%$ إلى نصف الاستجابة المتكيفة بعد القفزة صعودًا وهبوطًا، في كل من القانونين؟
\end{exercise}

\begin{exercise}[\lvlC]
\label{ex:11:7}
\emph{لأي عالم ضُبط المنحنى.} عند ضجيج مخرج ثابت صغير تكون المعلومات عن السطوع أكبر ما يمكن حين $r(I)=r_{\max}\Phi_I(I)$، حيث $\Phi_I$ دالة توزيع السطوع. (أ)~لأي توزيع يكون المنحنى~\eqref{eq:nakarushton} أمثل، وما انتشار $\log_{10}I$ فيه؟ (ب)~أي منحنى هو الأمثل عند ضجيج مخرج بواسوني؟
\end{exercise}

\chapter{التعرّف على الصور والفيديو}
\chaptermark{التعرّف}
\label{sec:recognition}
\begin{chaptergoals}
\item لماذا يجب أن يتجاهل التعرّف الإزاحة والمقياس والإضاءة، ولماذا تفشل المقارنة بكسلًا ببكسل؛
\item كيف يتعرّف خط أنابيب يناوب بين ”و“ على الأجزاء و”أو“ على المواضع في $150\ms$؛
\item كيف تتعلم قاعدة هيب مع أثر اللاتغيّر من فيديو بلا أسماء؛
\item بمَ يختلف الدماغ عن الشبكة الالتفافية، وماذا تكشف الخدع البصرية عنه.
\end{chaptergoals}

\section{المسألة: الشيء نفسه ببكسلات مختلفة}

أوصل الفصل~\ref{sec:sensors} الصورة إلى مدخل الآلة: نحو مليون عصبون مخرج في العين، وتدفق مضغوط لا يُنقل فيه في الغالب إلا الحواف والتغيرات. لكن بين المستشعر والفعل خطوة سكتنا عنها حتى الآن. القطة في الصورة ليست مجموعة من البكسلات. أزِحها نصف إطار، أو أدِرها، أو أبعِدها، أو أطفئ نصف الضوء، فتتغير كل البكسلات تقريبًا، ومع ذلك يجب أن يبقى الجواب ”قطة“ كما هو. يسمي المهندس هذا \emph{اللاتغيّر}: يجب ألا يتغير جواب المصنِّف مع الإزاحة والمقياس والدوران والإضاءة.

تظهر الصعوبة جليًا في تجربة بسيطة. لنأخذ ”شبكية“ أحادية البعد من $24$ نقطة وشكلين من خمس نقاط يمكن أن يقعا في أي مكان. المصنِّف الذي يقارن الدخل بنموذج محفوظ في موضع واحد يصيب في $49\text{--}52\%$ من الحالات، أي بقدر رمي عملة. فالإزاحة تحطم المقارنة بكسلًا ببكسل تمامًا. لا بد من مخطط آخر.

\section{خط أنابيب من مراحل}

نعرف من الفصل~\ref{sec:code} مدى سرعة الدماغ في ذلك: يُتعرَّف على حيوان في صورة معروضة لحظيًا في نحو $150\ms$، وتمر الإشارة بنحو عشر مراحل، من $10$ إلى $20\ms$ لكل منها~\cite{Thorpe1996}. وفي كل مرحلة لا يلحق العصبون أن يُصدر إلا صفرًا من النبضات أو نبضة واحدة. إذن التعرّف السريع مرور واحد عبر خط الأنابيب، بلا تفكير طويل ولا تكرار. والسؤال: ماذا تفعل كل مرحلة؟

الجواب الذي أوحت به القياسات في المراحل الأولى للبصر~\cite{HubelWiesel1962} يتألف من عمليتين متناوبتين (الشكل~\ref{fig:conveyor}).

\emph{الانتقائية.} ينظر عصبون المرحلة $S$ إلى جزء صغير من الدخل، ولا يطلق إلا إذا وُجد فيه تركيب معيّن من الأجزاء: هذه الحافة، وبجوارها تلك. هذا ”و“ على الأجزاء، وهو العصبون العتبي من الفصل~\ref{sec:glutcomputer} الذي عتبته أعلى مما يعطيه أي جزء وحده.

\emph{اللاتغيّر.} يجمع عصبون المرحلة $C$ مخارج عصبونات $S$ كثيرة تبحث عن التركيب نفسه في مواضع مختلفة، ويطلق إذا وُجد التركيب في أي مكان. هذا ”أو“ على المواضع، أو بتعبير أدق: اختيار الأكبر.

وفي النموذج الأحادي البعد تُكتب العمليتان هكذا:
\begin{empheq}[box=\keybox]{equation}
s_{f,p} \;=\; \Bigl[\sum_i t_{f,i}\,x_{p+i} - \theta\Bigr]_+ ,
\qquad
c_f \;=\; \max_p\, s_{f,p} ,
\label{eq:scstage}
\end{empheq}
حيث $x$ الدخل، و$t_{f,i}$ نموذج السمة $f$، و$p$ الموضع، و$\theta$ العتبة. التعبير الأول يجعل الجواب انتقائيًا، والثاني يجعله مستقلًا عن الموضع. وتحصل الشبكة على أكبر المداخل بأدوات الفصل~\ref{sec:gaba}: التثبيط المتبادل يُبقي أقوى الدخل (القضية~\ref{prop:wta})، والقسمة على النشاط الكلي تسوّي الاستجابات عند سطوع مختلف~\cite{Carandini2012}.

\begin{figure}[t]
\centering
\begin{tikzpicture}[>=Stealth,
  px/.style={draw,minimum size=0.36cm,inner sep=0pt},
  on/.style={px,fill=blue!55!black},
  snode/.style={circle,draw,very thick,minimum size=0.62cm,inner sep=0pt,font=\scriptsize,fill=blue!6},
  cnode/.style={circle,draw,very thick,minimum size=0.8cm,inner sep=0pt,font=\scriptsize,fill=red!10}]
% two inputs: the same shape at two positions
\foreach \row/\sh/\lab in {0/1/{الإطار 1},1/7/{الإطار 2}}{
  \begin{scope}[yshift=-\row*1.0cm]
  \node[left,font=\scriptsize] at (-0.25,0) {\lab};
  \foreach \i in {0,...,13}{\node[px] at (\i*0.4,0) {};}
  \foreach \d in {0,1,3,4}{\node[on] at ({(\sh+\d)*0.4},0) {};}
  \end{scope}}
% S stage: detectors of the shape at several positions
\foreach \k in {0,...,5}{
  \node[snode] (S\k) at ({0.8+0.8*\k},1.7) {$S$};}
\node[font=\scriptsize,left] at (-0.25,1.7) {المرحلة $S$: ”و“ على الأجزاء};
\draw[->,thick,blue!60!black] (1.2,0.25) -- (S1.south);
\draw[->,thick,orange!85!black] (3.6,0.25) -- (S4.south);
\node[font=\scriptsize,blue!60!black,anchor=east] at (1.25,0.85) {الإطار 1};
\node[font=\scriptsize,orange!85!black,anchor=west] at (3.9,0.85) {الإطار 2};
% C stage
\node[cnode] (C) at (2.8,3.25) {$C$};
\node[font=\scriptsize,left] at (-0.25,3.25) {المرحلة $C$: ”أو“ على المواضع};
\foreach \k in {0,...,5}{\draw[->,gray!60!black] (S\k.north) -- (C);}
\draw[->,very thick,red!70!black] (C.north) -- ++(0,0.6) node[above,font=\scriptsize,black] {”الشكل موجود“ في الإطارين كليهما};
% next stages
\draw[decorate,decoration={brace,amplitude=5pt},gray!60!black] (6.3,1.4) -- (6.3,-1.3);
\node[right,font=\scriptsize,align=left] at (6.5,0.05) {الشيء نفسه،\\وبكسلات\\أخرى};
\draw[->,thick,densely dashed,gray!60!black] (3.6,3.25) -- (5.9,3.25) node[right,font=\scriptsize,black,align=left] {الأزواج التالية\\$S$، $C$: أكبر،\\وأعقد، وأكثر لاتغيّرًا};
\end{tikzpicture}
\caption{زوج واحد من مراحل خط الأنابيب. تبحث عصبونات $S$ عن تركيب الأجزاء نفسه في مواضع مختلفة؛ ويستجيب العصبون $C$ إذا وُجد التركيب في أي مكان~\eqref{eq:scstage}. الشكل في الإطار 1 وفي الإطار 2 يثير عصبونات $S$ مختلفة، لكنه يثير العصبون $C$ نفسه. والأزواج التالية من المراحل تكرر الشيء نفسه على تراكيب أكبر فأكبر وأعقد فأعقد.}
\label{fig:conveyor}
\end{figure}

في النموذج الأحادي البعد نفسه يتعرف زوج المراحل~\eqref{eq:scstage} على الشكل في أي مكان بلا أخطاء، وإذا قُلبت كل نقطة من الدخل عشوائيًا باحتمال $0.1$ فإنه يصيب في $86\%$ من الحالات، وباحتمال $0.2$ في $70\%$. أما رمي العملة فيعطي النصف كما من قبل. ضع عدة أزواج كهذه بعضها فوق بعض: الأول يبحث عن الحواف، والتالي عن تراكيب الحواف، وأعلى منه عن أجزاء الأشياء، وفي القمة عن الأشياء كاملة. ومع كل زوج تكبر التراكيب، ويقل اعتماد الجواب على مكان الشيء وهيئته~\cite{Riesenhuber1999}.

\begin{keyblock}
\begin{proposition}[خط أنابيب الانتقائية واللاتغيّر]
\label{prop:conveyor}
التعرّف السريع مرور واحد عبر نحو عشر مراحل. في كل زوج من المراحل يجعل ”و“ على الأجزاء~\eqref{eq:scstage} الجواب انتقائيًا، ويجعله ”أو“ على المواضع مستقلًا عن الإزاحة. وتناوب هاتين العمليتين يعطي جوابًا يميّز الأشياء ولا يكاد يتوقف على مكانها وهيئة رؤيتها. بنية المراحل الأولى وسرعة المرور مقيستان؛ أما أن السلسلة كلها تُختزل إلى هاتين العمليتين فتبسيط.
\end{proposition}
\end{keyblock}

إذا بدت هذه البنية مألوفة، فهذا طبيعي. فتناوبها هذا بالذات، أي البحث عن السمة في كل المواضع ثم الدمج على المواضع المتجاورة، هو ما نقله المهندسون إلى الشبكات الاصطناعية \emph{الالتفافية}~\cite{Fukushima1980}. ثم ساروا مؤخرًا في الاتجاه المعاكس: تبيّن أن الشبكات الالتفافية المدرَّبة على التعرف على الأشياء أفضل نموذج معروف لاستجابات العصبونات في المراحل العليا للبصر~\cite{Yamins2014}. والنتيجة في القمة تُقرأ ببساطة: من الترددات الكلية لنحو مئة عصبون في المرحلة الأخيرة تتعرف وحدة قرار خطية على الشيء بعد عُشر ثانية من عرضه~\cite{Hung2005}. هذه شفرة التردد الجماعية من الفصل~\ref{sec:code}.

\section{المعلّم الغائب: التعلم من الفيديو}

تُدرَّب الشبكة الالتفافية على ملايين الصور المعنونة بأسمائها. أما الدماغ فلا يكاد أحد يعطيه أسماء: يرى الطفل الأشياء ساعات، ولا يسمع كلمة ”فنجان“ إلا نادرًا. فمن أين تعرف مراحل $C$ أيّ عصبونات $S$ تدمج؟ فلكي ندمج ”هذا الشكل هنا“ مع ”هذا الشكل هناك“ يجب أن نعرف مسبقًا أنه الشكل نفسه.

والتلميح يأتي من الفيديو نفسه. العالم يتغير باستمرار: الشيء في مجال الرؤية يبقى الشيء نفسه وهو يتزحزح ويدور وتتغير إضاءته، ولا ينتقل البصر إلى غيره إلا أحيانًا. فكل ما يتغير \emph{ببطء} يرجح أنه يخص الشيء، وكل ما يتغير بسرعة يخص موضعه وهيئته~\cite{WiskottSejnowski2002}. إذن يكفي العصبون $C$ قاعدة هيب من الفصل~\ref{sec:dofa} مع أثر: أن يربط ما جاء متتابعًا، لا ما جاء متزامنًا فقط~\cite{Foldiak1991}:
\begin{empheq}[box=\keybox]{equation}
\tau_{\text{tr}}\,\frac{\dd\bar y_k}{\dd t} \;=\; -\bar y_k + y_k ,
\qquad
\frac{\dd\mathbf w_k}{\dd t} \;=\; \eta\,\bar y_k\,\bigl(\mathbf x - \mathbf w_k\bigr) .
\label{eq:traceinvariance}
\end{empheq}
هنا $y_k$ نشاط العصبون $C$ رقم $k$ الذي فاز في التنافس عبر التثبيط، و$\bar y_k$ أثره، وهو دارة \foreignlanguage{english}{RC} نفسها التي في القاعدة~\eqref{eq:threefactor}، و$\mathbf x$ مخارج عصبونات $S$. العصبون الذي فاز على الهيئة الأولى للشيء لا يزال يتذكر ذلك حين تأتي الهيئة الثانية، فيجذبها إليه هي أيضًا.

\begin{figure}[t]
\centering
\begin{tikzpicture}[>=Stealth,x=1.0cm,y=1.0cm]
\foreach \i/\lab in {0/1,1/2,2/3,3/4}{
  \node[draw,thick,fill=blue!8,minimum width=0.9cm,minimum height=0.55cm,font=\scriptsize] at (\i+0.5,2.2) {$A$ في \lab};}
\foreach \i/\lab in {0/4,1/3,2/2,3/1}{
  \node[draw,thick,fill=orange!12,minimum width=0.9cm,minimum height=0.55cm,font=\scriptsize] at (\i+6.5,2.2) {$B$ في \lab};}
\node[font=\scriptsize,gray!40!black] at (5.0,2.2) {توقف};
\draw[->,gray!70!black] (0,0) -- (10.4,0) node[below left,font=\scriptsize,black] {الزمن};
% traces
\draw[thick,blue!60!black] (0,0.05) .. controls (0.4,1.2) and (0.8,1.3) .. (1.2,1.35) -- (4,1.4) .. controls (4.6,0.5) and (5.2,0.15) .. (6,0.08) -- (10,0.05);
\draw[thick,orange!85!black] (0,0.05) -- (6,0.05) .. controls (6.4,1.2) and (6.8,1.3) .. (7.2,1.35) -- (10,1.4);
\node[font=\scriptsize,blue!60!black,anchor=west] at (1.4,1.65) {الأثر $\bar y_1$ يدوم طوال مرور $A$};
\node[font=\scriptsize,orange!85!black,anchor=west] at (7.2,1.65) {الأثر $\bar y_2$};
\draw[decorate,decoration={brace,amplitude=4pt},gray!60!black] (4.0,2.6) -- (0,2.6);
\draw[decorate,decoration={brace,amplitude=4pt,mirror},gray!60!black] (0,-0.35) -- (4,-0.35) node[midway,below=4pt,font=\scriptsize,black] {كل هيئات $A$ ترتبط بالعصبون 1};
\end{tikzpicture}
\caption{كيف يعلّم الفيديو اللاتغيّر، مخطط. يمر الشيء $A$ بأربعة مواضع، ثم توقف، ثم الشيء $B$. أثر~\eqref{eq:traceinvariance} العصبون الذي فاز على الهيئة الأولى لـ$A$ يدوم طوال المرور، فترتبط كل هيئات $A$ بعصبون واحد. والتوقف يمحو الأثر، فيذهب $B$ إلى عصبون آخر.}
\label{fig:tracevideo}
\end{figure}

اختبرنا ذلك على نموذج (الشكل~\ref{fig:tracevideo}). عصبونا $C$ يتنافسان على مداخل من $16$ عصبون $S$: شكلان في ثمانية مواضع. يمر الشكل عبر كل المواضع، $50\ms$ لكل موضع، ثم توقف $0.3\,\text{s}$، ثم الشكل العشوائي التالي. بلا أي أسماء. إذا كان الأثر قصيرًا، $\tau_{\text{tr}}=5\ms$، اقتسم العصبونان الدخل بحسب \emph{الموضع}، ولم يطابق الجوابُ الشكلَ بأفضل من رمي عملة: $52\%$ في المتوسط على عشرة تشغيلات. وإذا كان الأثر مماثلًا لزمن عرض موضع واحد، ابتداءً من $\tau_{\text{tr}}\approx20\ms$، جمع كل عصبون \emph{كل} مواضع شكله (عند $20$ و$50$ و$200\ms$ في التشغيلات العشرة كلها): تعلّم اللاتغيّر من الفيديو وحده. والتوقف بين الأشياء مهم: بدونه ينساب الأثر إلى الشيء التالي فيخلط بينهما.

والشيء نفسه يُرى في التجربة. إذا استُبدل الشيء خفيةً في التجربة بينما تقفز العين من مكان إلى آخر، فإن عصبونات المراحل العليا تبدأ في غضون ساعة أو ساعتين بالاستجابة للشيء البديل كأنه الشيء نفسه: إنها تربط ما جاء متتابعًا~\cite{LiDiCarlo2008}. فاللاتغيّر في الدماغ يُتعلَّم فعلًا من استمرارية الزمن. أما إلى أي حد تفسر هذه القاعدة كل تعلم البصر، ففرضية.

\section{الفيديو حركة}

ليس الفيديو مجرد تدفق من الإطارات للتعلم، بل هو المسألة نفسها أيضًا: ما الذي يتحرك، وإلى أين، وبأي سرعة. الخطوة الأولى معروفة لنا: كاشف الاتجاه من الفصل~\ref{sec:gaba} (الشكل~\ref{fig:direction})، الذي يُطفئ فيه التثبيطُ المتأخر الحركةَ في اتجاه واحد. تنتشر هذه الكواشف في كل مجال الرؤية، ثم يُعامَل معها كما يُعامَل مع سمات الشكل: المراحل التي تجمع كواشف كثيرة للاتجاه نفسه في مواضع مختلفة تستجيب لحركة شيء كبير أينما كان. هذا زوج ”و“ و”أو“ نفسه~\eqref{eq:scstage}، إلا أن السمة ليست ”حافة“، بل ”حافة انزاحت خلال زمن $d$“.

والفيديو فوق ذلك قابل للتنبؤ: الإطار التالي يكاد يطابق السابق. بحسب فرضية \emph{الترميز التنبؤي} ترسل المراحل العليا إلى الدنيا تنبؤًا، ولا يصعد إلى الأعلى في الغالب إلا الخطأ، أي ما لم يحسب التنبؤ حسابه~\cite{RaoBallard1999}. هذا هو الاقتصاد نفسه في النبضات الذي في القضية~\ref{prop:economy}: ادفع ثمن الجديد، لا ثمن المعروف أصلًا. بعض المشاهدات تتفق مع هذه الفرضية، لكنها لم تُثبت بوصفها البنية العامة لخط الأنابيب.

\section{الدماغ والشبكة الالتفافية}

إلى أي مدى يذهب التشابه؟ في التعرّف السريع على شيء واحد هو عميق حقًا~\cite{Yamins2014}. لكن في الدماغ ما ليس في الشبكة الالتفافية البسيطة.

\emph{الوصلات الراجعة.} خط أنابيب الدماغ ليس أماميًا فقط: لكل مرحلة وصلات إلى الخلف وإلى الجانب. في الصور الصعبة، كالمتداخلة أو ضعيفة الإضاءة، تتكوّن استجابة المراحل العليا متأخرة، وهذه الاستجابات المتأخرة تصفها الشبكات ذات الوصلات الراجعة أفضل من الشبكات الأمامية~\cite{Kar2019}. المرور الواحد يحل الحالات السهلة، والصعبة تتطلب عدة دورات.

\emph{المتانة.} يمكن خداع الشبكة الاصطناعية بتزييف في البكسلات لا تلحظه العين~\cite{Szegedy2014}: تغيير لا يراه الإنسان يحوّل ”الباندا“ إلى ”جِبّون“~\cite{Goodfellow2015}. والبصر البشري أمتن بكثير أمام هذا التزييف، لكنه ليس محصنًا: الصور التي تخدع شبكات كثيرة معًا تزيح اختيار الإنسان أيضًا عند العرض القصير~\cite{Elsayed2018}. لماذا تختلف المتانة إلى هذا الحد؟ سؤال مفتوح؛ والمرشحون: الضجيج، والقسمة على النشاط الكلي، والوصلات الراجعة.

\emph{المعلّم.} تحتاج الشبكة إلى ملايين الأمثلة المعنونة، أما الدماغ فيحتاج في الغالب إلى فيديو بلا أسماء~\eqref{eq:traceinvariance} وقليل من التلميحات من \nt{DOPA} عما هو مهم.

\emph{الطاقة.} الدماغ كله نحو $20$ واطًا (القضية~\ref{prop:economy})، والتعرّف لا يشغل إلا جزءًا من ذلك؛ أما الشبكة الالتفافية المدرَّبة على مسرّع رسوميات فتستهلك مئات الواطات.

\section{الخداع البصري: أين تخطئ الآلة ولماذا}
\label{sec:illusions}

المهندس الذي يريد أن يفهم دارة غيره يغذيها بإشارات غير مألوفة، وينظر أين تخطئ. وللبصر إشارات كهذه ابتُكرت منذ زمن بعيد، هي الخدع البصرية. الخدعة ليست عطلًا. كل خدعة تشير إلى آلية بعينها في الآلة، تعمل صحيحًا في الظروف العادية، وتكشف نفسها في صورة مختارة خصيصًا.

\emph{توقعات معقولة.} الدخل مشوَّش بالضجيج، والآلة تكمله بما يحدث عادةً في العالم. الحركات البطيئة أكثر من السريعة؛ وحين يكون الدخل غير موثوق، كأن يكون تباين الصورة ضعيفًا، ينزاح تقدير السرعة نحو البطء، فيبدو الشيء الباهت أبطأ من الساطع~\cite{Weiss2002}. وبالطريقة نفسها تُفسَّر خدع كثيرة في السطوع: نحن لا نرى سطوع البكسل، بل السطح الذي طابقه هذا السطوع في أغلب الأحيان في الماضي~\cite{Purves2011}. وصورة الفستان الذي يراه بعضهم أبيض وذهبيًا ويراه آخرون أزرق وأسود هي الآلية نفسها: الصورة ملتبسة، ويختلف الناس في الإضاءة التي يفترضونها دون وعي~\cite{LaferSousa2015,Wallisch2017}. الفروق بين الناس مقيسة؛ أما أن الدماغ يجمع هنا الدخل والتوقع وفق قواعد نظرية الاحتمالات ففرضية راسخة الأسس.

\emph{ضبط الكسب.} انظر دقيقة إلى شلال، ثم إلى صخور ساكنة: ستبدو الصخور كأنها تطفو إلى أعلى. قناتا ”إلى أسفل“ و”إلى أعلى“ زوج أسلاك كما في~\eqref{eq:dualrail}؛ وقد خفّض ضبط الكسب~\eqref{eq:nakarushton} القناة المتعبة، فاختل توازن الزوج. وخدع الميل والتباين التي يغيّر فيها المحيطُ مظهرَ المركز تُفسَّر بالقسمة على نشاط الجيران~\cite{Schwartz2009}، كما في الفصل~\ref{sec:gaba}. وهناك مثال مفيد على الحذر: البقع الداكنة عند تقاطعات الخطوط البيضاء في الشبكة المربعة فُسرت طويلًا بتثبيط الجيران البسيط، لكن التجارب على شبكة مربعة معدّلة أظهرت أن هذا التفسير لا يصح~\cite{SchillerCarvey2005}.

\emph{تعويض التأخيرات.} يستغرق المسار من العين إلى المراحل العليا عشرات الملّي ثانية، والشيء المتحرك يلحق أن ينزاح خلالها. إذا ومضت نقطة ساكنة بجانبه، بدا الشيء متقدمًا على الومضة مع أنهما كانا متطابقين~\cite{Nijhawan1994}. وبحسب أحد التفسيرات تمدّ الآلة الحركة إلى الأمام لتعوّض التأخير، كما في الفصل~\ref{sec:muscles}: مع تغذية راجعة متأخرة لا بد من التنبؤ. ولهذه الخدعة تفسيرات عدة، والخلاف لم يُحسم.

\emph{خطأ في شفرة التوقيت.} صورة ساكنة من حلقات ذات انتقالات لونية حادة تبدو كأنها تدور. فالمناطق عالية التباين تُعالَج أسرع من الباهتة، وتقرأ كواشفُ الاتجاه فرقَ التأخيرات~\eqref{eq:latency} على أنه حركة~\cite{Conway2005}. ويطلق الخدعة قفزات العين الصغيرة اللاإرادية والرمش: كل قفزة تجدد الدخل، فترى الكواشف ”حركة“ من جديد~\cite{OteroMillan2012}. هذه شفرة التوقيت من الفصل~\ref{sec:code} وقد قُرئت خطأً.

\emph{صورتان في صورة.} هيكل مكعب يواجهنا تارة بوجه وتارة بآخر، أو صورتان مختلفتان في العينين: يقفز الإدراك كل بضع ثوانٍ. أفضل النماذج تثبيط متبادل بين التأويلين، وتعب بطيء، وضجيج~\cite{MorenoBote2007}، أي المخطط نفسه~\eqref{eq:halfcentre} الذي ولّد في الفصل~\ref{sec:muscles} إيقاع الخطو. يحدد زمنَ التبدل الضجيجُ والتعب؛ وبحسب النموذج~\cite{MorenoBote2007} الدور الأكبر للضجيج، والتعب يساعد فقط.

وماذا عن الشبكات الاصطناعية؟ الشبكة المدرَّبة على التنبؤ بالإطار التالي من الفيديو ترى الدوران في الحلقات الساكنة نفسها، ولا تراه في الصور الضابطة~\cite{Watanabe2018}. والشبكات المدرَّبة على تنقية الصور من الضجيج وزيادة حدتها تنخدع بخدع اللون والسطوع نفسها التي ينخدع بها الناس~\cite{GomezVilla2020}. إذن بعض الخدع نتيجة للمسألة نفسها التي تتعلمها الآلة، لا لخصائص النسيج العصبي. وأي الخدع يختص بها الدماغ وحده؟ هذا اختبار جيد لأي نموذج للبصر.

\begin{keyblock}
\begin{proposition}[الخدع بصمات للآليات]
\label{prop:illusions}
تنشأ الخدعة حين تتلقى آلية نافعة في العالم العادي دخلًا غير مألوف: التوقع يغلب الدخل غير الموثوق، وضبط الكسب يزيح زوج القنوات، وتعويض التأخير يقدّم الشيء إلى الأمام، وفرق التأخيرات يُقرأ حركةً، والتثبيط المتبادل مع الضجيج والتعب يقفز. كل خدعة إشارة اختبار تشير إلى آليتها. الخدع نفسها مقيسة؛ وتفسيراتها تتراوح بين الراسخ الأسس والمختلف فيه.
\end{proposition}
\end{keyblock}

\section{الخلاصة}

\begin{table}[t]
\centering
\footnotesize
\setlength{\tabcolsep}{4pt}
\renewcommand{\arraystretch}{1.15}
\begin{tabular}{>{\raggedright}p{3.0cm}>{\raggedright}p{4.6cm}>{\raggedright\arraybackslash}p{5.2cm}}
\toprule
ماذا تفعل الآلة & كيف & ما هو معروف \\
\midrule
تتعرف في $150\ms$ & مرور واحد عبر نحو عشر مراحل & السرعة مقيسة \\
تجعل الجواب انتقائيًا & ”و“ على الأجزاء، المرحلة $S$~\eqref{eq:scstage} & مقيس في المراحل الأولى \\
تجعل الجواب لا متغيرًا & ”أو“ على المواضع، المرحلة $C$؛ التثبيط والقسمة & مقيس في المراحل الأولى؛ وفي العليا تبسيط \\
تُصدر الجواب & ترددات نحو مئة عصبون في المرحلة الأخيرة، قراءة خطية & مقيس \\
تتعلم بلا أسماء & قاعدة هيب مع أثر~\eqref{eq:traceinvariance} على فيديو متصل & استبدال الشيء يغيّر الاستجابات، وهذا مقيس؛ وكونها القاعدة الرئيسية فرضية \\
ترى الحركة & كواشف الاتجاه، ثم زوج ”و“/”أو“ نفسه & مقيس \\
تقتصد في المتوقَّع & يصعد خطأ التنبؤ إلى الأعلى & فرضية \\
تحل الحالات الصعبة & وصلات راجعة، عدة دورات & الاستجابات المتأخرة ودور الوصلات الراجعة مقيسة \\
تخطئ على نحو متوقَّع & الخدع: التوقعات، ضبط الكسب، التأخيرات، التثبيط المتبادل مع الضجيج والتعب & الخدع مقيسة؛ والتفسيرات بين الراسخ والمختلف فيه \\
\bottomrule
\end{tabular}
\caption{كيف تتعرف الآلة على الصور والفيديو.}
\label{tab:recognition}
\end{table}

التعرّف مجمَّع من قطع كانت لدينا من قبل (الجدول~\ref{tab:recognition}): العتبة تعطي ”و“ على الأجزاء، والتثبيط المتبادل ”أو“ على المواضع، والقسمة الاستقلال عن السطوع، وقاعدة هيب مع أثر تعطي المعلّم الغائب. أخذ المهندسون من البصر تناوب الانتقائية واللاتغيّر وبنوا عليه الشبكات الالتفافية؛ أما الدماغ فلم يسلّمنا بعدُ أهم ما لديه: القدرة على التعلم من فيديو بلا أسماء، بلا طاقة تقريبًا.

\section{تمارين}

\begin{exercise}[\lvlA]
\label{ex:12:1}
\emph{التردد في مرحلة واحدة.} تدوم مرحلة خط الأنابيب $15\ms$، والعصبونات تعمل بتردد $20$ نبضة في الثانية، والضجيج مترابط بـ$c=0.1$~\eqref{eq:correlated}. ما الخطأ النسبي لتردد مجموعة من مئة عصبون، وما حده لمجموعة كبيرة بلا حد؟ كم مستوى من التردد يمكن تمييزه في المرحلة؟
\end{exercise}

\begin{exercise}[\lvlA]
\label{ex:12:2}
\emph{طاقة تعرّف واحد.} في مرور واحد خلال $150\ms$ تُصدر عشر مراحل من $10^7$ عصبون لكل منها نبضة واحدة على الأكثر، كلفتها $10^{-10}\,\text{J}$. (أ)~ما نسبة طاقة الدماغ كله خلال هذه الـ$150\ms$ التي يكلفها التعرّف؟ (ب)~كم مرة أكثر يستهلك مسرّع قدرته $300\,\text{W}$ يتعرف على الصورة في $10\ms$؟
\end{exercise}

\begin{exercise}[\lvlB]
\label{ex:12:3}
\emph{عتبات المرحلة $S$.} ”شبكية“ من $24$ نقطة، والشكلان $A=11011$ و$B=10101$، والنموذج~\eqref{eq:scstage} يساوي $+1$ على آحاد الشكل و$-1$ على أصفاره. (أ)~عند أي عتبات يتسامح عصبون $S$ مع نقطة مقلوبة واحدة، ويصمت على الشكل الآخر؟ (ب)~كم عصبون $S$ يلزم لعشر سمات $5\times5$ في صورة $100\times100$؟
\end{exercise}

\begin{exercise}[\lvlB]
\label{ex:12:4}
\emph{كم تدوم إحدى الصورتين.} في مخطط القفزات~\eqref{eq:halfcentre} عند $\tau\ll\tau_a$، ما دامت المجموعة~1 فائزة، يكون $r_2=0$ و$r_1=I-a_1$. (أ)~أوجد مدة الفوز عند $I=1$ و$\beta=2$ و$b=2$ و$\tau_a=0.4\,\text{s}$. (ب)~أي $\tau_a$ يعطي تبدّل الصورة كل ثلاث ثوانٍ؟
\end{exercise}

\begin{exercise}[\lvlB]
\label{ex:12:5}
\emph{نمذجة: ثمن اللاتغيّر.} بالدالة \foreignlanguage{english}{\texttt{two\_stage}} من \foreignlanguage{english}{\texttt{models/\allowbreak recognition\_models.py}} ($4000$ محاولة) أوجد عند أي احتمال لقلب النقطة $p$ يخطئ زوج المراحل $S$، $C$ عند $N=24$ في حالة من كل أربع. وكيف تتغير نسبة الأجوبة الصحيحة عند $p=0.1$ من $N=8$ إلى $N=192$؟
\end{exercise}

\begin{exercise}[\lvlB]
\label{ex:12:6}
\emph{نمذجة: نافذة الأثر.} عشرة تشغيلات لـ\foreignlanguage{english}{\texttt{trace\_learning}} من \foreignlanguage{english}{\texttt{models/\allowbreak recognition\_models.py}} بتوقف $0.3\,\text{s}$: أوجد عند أي قيم $\tau_{\text{tr}}$ بين $5\ms$ و$2\,\text{s}$ تتعلم كل التشغيلات اللاتغيّر بلا خطأ. ومن أين يأتي الحد الأعلى لهذه النافذة؟
\end{exercise}

\begin{exercise}[\lvlB]
\label{ex:12:7}
\emph{الحركة في أي مكان.} على نقاط متجاورة من صف بخطوة $0.1^\circ$ توجد كواشف اتجاه من الفصل~\ref{sec:gaba}: التثبيط من $A$ بتأخير $d$ يُطفئ الإثارة من $B$ إذا لم يتباعدا بأكثر من $5\ms$؛ وفوقها المرحلة $C$. ما التأخيرات اللازمة لكي تصمت $C$ عند الحركة المعاكسة بسرعة من $5$ إلى $20^\circ$ في الثانية؟
\end{exercise}

\begin{exercise}[\lvlC]
\label{ex:12:8}
\emph{تزييف البكسلات.} كم نقطة مقلوبة تلزم لتغيير جواب النموذج \foreignlanguage{english}{\texttt{two\_stage}} ($N=24$) على شكل نظيف؟ وكم تلزم للمصنِّف ”أكثر من $3.5$ آحاد في الدخل يعني $A$“، وهو أيضًا لا يتغير بالإزاحة؟ وما نسبة أجوبته الصحيحة عند $p=0.1$ مقابل $0.86$ لزوج المراحل؟
\end{exercise}

\chapter{مخارج الآلة: كيف يتحكم الدماغ في العضلات}
\chaptermark{مخارج الآلة}
\label{sec:muscles}

\section{المخرج السريع}

كل ما تفعله الآلة في العالم بسرعة تفعله بالعضلات. الكلمة، والنظرة، والخطوة، والابتسامة كلها انقباضات عضلات. أما المخرج الثاني البطيء، أي كيمياء الجسم، فنتناوله في الفصل~\ref{sec:chemoutput}. في الشكل~\ref{fig:machine} كان المخرج سهمًا عنوانه ”العضلات“؛ فلننظر الآن إلى ما وراء هذا السهم.

الحلقة الأخيرة في السلسلة عصبونات \nt{ACET}، وهي نفسها التي كُتب عنها في الجدول~\ref{tab:types} ”المخرج إلى العضلة“. كل ليف عضلي يتلقى دخله من عصبون واحد بالضبط من هذه العصبونات، والعصبون الواحد يتحكم في مجموعة من الألياف، من بضع مئات إلى نحو ألفين~\cite{Feinstein1955mu}. وفي العضلات التي تحتاج إلى ضبط دقيق تكون الوحدات أصغر، وفي عضلات الساقين الكبيرة أكبر. ويسمى العصبون مع أليافه \emph{الوحدة الحركية}: إنها أصغر قطعة من العضلة يستطيع الدماغ تشغيلها وحدها.

مشبك العصبون على العضلة مبني على نحو مختلف تمامًا عن مشابك القشرة في الفصل~\ref{sec:code}. هناك كان معظم النبضات يضيع. أما هنا فكل نبضة تطلق من الأستيل كولين أضعاف ما يلزم لإطلاق الليف~\cite{WoodSlater2001}. إنه مشبك ذو \emph{هامش أمان}: نبضة واحدة من العصبون تعني انقباضة واحدة بالضبط للألياف. داخل الآلة يمكن تحمّل وصلات غير موثوقة وتصحيح الأخطاء بالفائض؛ أما في المخرج فالخطأ يكلف حركة، والموثوقية مشتراة في العتاد مباشرة.

\section{القوة: التردد وعدد الوحدات}

النبضة الواحدة تعطي انقباضة مفردة، أي نفضة: تتصاعد القوة خلال بضع عشرات من الملّي ثانية ثم تهبط. وتقريب جيد لها~\cite{Fuglevand1993}
\begin{equation}
h(t) \;=\; F_1\,\frac{t}{T_c}\,e^{\,1-t/T_c},
\label{eq:twitch}
\end{equation}
حيث $T_c$ زمن التصاعد، من رتبة $50\ms$، و$F_1$ قوة النفضة الواحدة. وإذا تقاربت النبضات تجمّعت النفضات:
\begin{empheq}[box=\keybox]{equation}
F(t) \;=\; \sum_k h\bigl(t-t_k\bigr) .
\label{eq:forcesum}
\end{empheq}
هذا هو مرشح التمرير المنخفض نفسه الذي حوّل النبضات إلى تركيز في الفصل~\ref{sec:serocomputer}، إلا أن المخرج الآن ليس تركيزًا بل قوة. العضلة محوِّل رقمي تماثلي من النبضات إلى القوة (الشكل~\ref{fig:tetanus}). في نموذجنا، عند خمس نبضات في الثانية لا تكاد النفضات تتداخل، فتقفز القوة بين $0.2F_1$ و$1.1F_1$؛ وعند خمس عشرة تبقى بين $1.8F_1$ و$2.2F_1$؛ وعند أربعين تصبح شبه مستوية، نحو $5.4F_1$. والعضلة الحقيقية تتشبع مع النبضات المتقاربة، فلا يصح الجمع الخطي~\eqref{eq:forcesum} إلا حتى بضع عشرات من النبضات في الثانية.

\begin{figure}[t]
\centering
\begin{tikzpicture}[>=Stealth,x=20cm,y=0.55cm]
\draw[->,gray!70!black] (0,0) -- (0.66,0) node[right,font=\small,black] {$t,\ \text{s}$};
\draw[->,gray!70!black] (0,0) -- (0,6.4) node[left,font=\small,black] {$F/F_1$};
\foreach \t in {0,0.2,0.4,0.6}{\draw[gray!60!black] (\t,-0.1) -- (\t,0.1); \node[below,font=\scriptsize] at (\t,0) {\t};}
\foreach \f in {1,3,5}{\draw[gray!60!black] (-0.004,\f) -- (0.004,\f); \node[left,font=\scriptsize] at (0,\f) {\f};}
\draw[thick,blue!60!black] plot file {figures/muscle_twitch_5.dat};
\draw[thick,green!45!black] plot file {figures/muscle_twitch_15.dat};
\draw[thick,red!70!black] plot file {figures/muscle_twitch_40.dat};
\node[font=\scriptsize,blue!60!black,anchor=west] at (0.605,0.9) {5 نبضات/ث};
\node[font=\scriptsize,green!45!black,anchor=west] at (0.605,2.1) {15 نبضة/ث};
\node[font=\scriptsize,red!70!black,anchor=west] at (0.605,5.4) {40 نبضة/ث};
\end{tikzpicture}
\caption{العضلة محوِّلًا رقميًا تماثليًا: القوة~\eqref{eq:forcesum} عند نبضات منتظمة من وحدة حركية واحدة، $T_c=50\ms$. النبضات المتباعدة تعطي نفضات منفصلة، والمتقاربة تندمج في قوة مستوية، كما تندمج النبضات المتقاربة في تركيز مستوٍ في الشكل~\ref{fig:lowpass}.}
\label{fig:tetanus}
\end{figure}

للدماغ مقبضان للقوة: تردد نبضات كل وحدة، وعدد الوحدات المشغَّلة. والمقبض الثاني مبني ببراعة كبيرة. الوحدات في العضلة مختلفة: أضعفها أضعف من أقواها مئة مرة. وحين تلزم قوة متزايدة تُشغَّل الوحدات \emph{بترتيب الحجم}، من الصغيرة إلى الكبيرة~\cite{Henneman1957}. لا أحد يحسب هذا الترتيب: العصبونات الصغيرة ببساطة أسهل استثارة بالدخل المشترك نفسه، فترتيب التشغيل يحدده العتاد نفسه.

ماذا يعطي ذلك؟ لنفترض أن كل وحدة تالية في الترتيب أقوى من سابقتها بـ$q$ مرة، حيث $q$ أكبر من الواحد بقليل. قوة كل الوحدات المشغَّلة مجموع متوالية هندسية، ومعظمها يأتي من الوحدات الأخيرة. عندئذ تضيف الوحدة التالية المشغَّلة
\begin{empheq}[box=\keybox]{equation}
\frac{\Delta F}{F} \;\approx\; q-1 \;=\; \text{const} ,
\label{eq:sizeprinciple}
\end{empheq}
أي أن خطوة القوة تتناسب مع القوة نفسها. في الجهد الضعيف يجرعه الدماغ بحصص صغيرة، وفي القوي بحصص كبيرة. هذا \emph{عدد بفاصلة عائمة}: الأس يحدده عدد الوحدات المشغَّلة، والجزء العشري (المانتيسا) يحدده تردد نبضاتها. وهو المقياس اللوغاريتمي نفسه الذي لمستشعرات الفصل~\ref{sec:sensors}: مدخل الآلة ومخرجها يقيسان العالم بمقاييس نسبية.

وللفاصلة العائمة وجه آخر. فانتشار القوة يزداد هو أيضًا بما يتناسب معها: الحركة القوية أقل دقة حتمًا من الضعيفة. وبحسب فرضية مؤثرة، فهذا الضجيج المعتمد على الإشارة بالذات هو ما يفسر نعومة حركاتنا: الأمر الناعم يعطي أصغر انتشار في النقطة النهائية~\cite{HarrisWolpert1998}.

\section{تسحب ولا تدفع: سلكان لكل مفصل}

العضلة لا تستطيع إلا السحب. لا تستطيع الدفع، كما لا يستطيع عصبون \nt{GLUT} أن يُصدر إشارة سالبة. والحل معروف من الفصل~\ref{sec:glutcomputer}: \emph{سلكان}. كل مفصل يخدمه زوج من العضلات يسحبان في اتجاهين متعاكسين (الشكل~\ref{fig:antagonists}). إذا كانت قوتاهما $F_1$ و$F_2$، وذراع القوة $\ell$، فإن عزم الدوران وصلابة المفصل
\begin{empheq}[box=\keybox]{equation}
M \;=\; \ell\,(F_1-F_2), \qquad K \;\approx\; \kappa_m\,(F_1+F_2),
\label{eq:antagonists}
\end{empheq}
حيث $\kappa_m$ معامل يربط قوة العضلة بصلابتها: العضلة المشدودة أشد نابضية من المرتخية. الفرق بين القوتين يحدد العزم، ومجموعهما يحدد الصلابة~\cite{Hogan1984}. بسلكين يتحكم الدماغ في كميتين مستقلتين: إلى أين يُدار المفصل، وبأي إحكام يُمسك. وحين يلزم أن نمسك فنجانًا دون أن ينسكب، نشد العضلتين معًا: العزم نفسه، والصلابة أعلى، والدفعة العارضة تزيح اليد أقل.

\begin{figure}[t]
\centering
\begin{tikzpicture}[>=Stealth,
  nrn/.style={circle,draw,very thick,minimum size=0.95cm,inner sep=0pt,font=\scriptsize},
  inh/.style={-{Bar[width=7pt]},thick,red!70!black}]
% the joint: a bar on a pivot, two springs pulling down on either side
\fill[gray!30] (4.6,-0.25) -- (5.4,-0.25) -- (5,0.15) -- cycle;
\draw[line width=3pt,gray!50!black] (2.4,0.2) -- (7.6,0.2);
\fill[black] (5,0.2) circle (0.08);
\draw[decorate,decoration={coil,segment length=5pt,amplitude=4pt},very thick,blue!60!black] (3.0,0.2) -- (3.0,-1.6);
\draw[decorate,decoration={coil,segment length=5pt,amplitude=4pt},very thick,orange!85!black] (7.0,0.2) -- (7.0,-1.6);
\draw[very thick] (2.5,-1.6) -- (3.5,-1.6); \draw[very thick] (6.5,-1.6) -- (7.5,-1.6);
\node[font=\small,blue!60!black] at (2.35,-0.75) {$F_1$};
\node[font=\small,orange!85!black] at (7.65,-0.75) {$F_2$};
\draw[->,thick] (5.9,0.75) arc (60:120:1.8) node[above,font=\scriptsize,pos=0.5] {$M=\ell(F_1-F_2)$};
% motor neurons and reflex circuit
\node[nrn,fill=green!12] (M1) at (1.6,-3.0) {\nt{ACET}};
\node[nrn,fill=green!12] (M2) at (8.4,-3.0) {\nt{ACET}};
\node[nrn,fill=red!10] (G) at (5.0,-3.0) {\nt{GLYC}};
\node[nrn,fill=blue!6] (S) at (1.6,-5.0) {\nt{GLUT}};
\draw[->,very thick] (M1) -- (3.0,-1.7);
\draw[->,very thick] (M2) -- (7.0,-1.7);
\draw[->,thick] (S) -- (M1);
\draw[->,thick] (S) -- (G);
\draw[inh] (G) -- (M2);
\draw[->,thick,densely dashed,gray!60!black] (3.15,-1.2) to[bend left=25] node[pos=0.35,right,font=\scriptsize,black] {التمدد} (S.north east);
\draw[->,thick,blue!60!black] (-0.3,-3.0) node[left,font=\scriptsize,black,align=right] {أمر\\الدماغ} -- (M1);
\draw[->,thick,blue!60!black] (10.3,-3.0) node[right,font=\scriptsize,black,align=left] {أمر\\الدماغ} -- (M2);
\end{tikzpicture}
\caption{عضلتان لمفصل واحد، وأسرع حلقة تغذية راجعة. تتلقى عصبونات \nt{ACET} أوامر الدماغ وتسحب المفصل في اتجاهين متعاكسين؛ الفرق بين القوتين يديره، ومجموعهما يزيده صلابة. مستشعر التمدد (\nt{GLUT}) في العضلة اليسرى يثير عصبون \nt{ACET} الخاص بها، ويثبّط عبر وسيط \nt{GLYC} عصبون العضلة المضادة، وهذا هو المنع من الفصل~\ref{sec:gaba}.}
\label{fig:antagonists}
\end{figure}

\section{تغذية راجعة متأخرة}

لا تعمل العضلات عمياء. ففي كل منها مستشعرات تمدد، ذُكرت في الجدول~\ref{tab:sensors}، وأقصر حلقة تُغلق مباشرة في الحبل الشوكي (الشكل~\ref{fig:antagonists}). إذا مُدّت العضلة فجأة، أثار مستشعر التمدد عصبون \nt{ACET} الخاص بها، فتنقبض العضلة وتعيد المفصل. وفي الوقت نفسه تُطفأ العضلة المضادة عبر عصبون وسيط مثبِّط، كي لا تعيق. ويعمل وسيطًا هنا \nt{GLYC}، وهو النوع نفسه من الجدول~\ref{tab:types} الذي لم يُخصَّص له فصل: مكانه هنا بالذات، في حلقات الحبل الشوكي السريعة.

لكل حلقة تأخير $d$: للحلقة الشوكية نحو $25\text{--}30\ms$ في الذراع، وللحلقة عبر القشرة أطول بمرة ونصف إلى ثلاث مرات، وللحلقة عبر العين $100\text{--}200\ms$. والتغذية الراجعة المتأخرة، كما رأينا في القضية~\ref{prop:ei}، تُحدث تذبذبًا في النظام. ولأبسط منظِّم يصحح الانحراف $x$ بسرعة $G$ استنادًا إلى معلومات عمرها $d$ ثانية، $\dd x/\dd t=-G\,x(t-d)$، يكون شرط الاستقرار~\cite{Hayes1950}:
\begin{empheq}[box=\keybox]{equation}
G\,d \;<\; \frac{\pi}{2} ,
\label{eq:delaystable}
\end{empheq}
وعلى الحد يتذبذب النظام بتردد $1/(4d)$. اختبرنا ذلك على نموذج: عند $d=30\ms$ يخمد الانحراف عند $G=40$ في الثانية، أما عند $G=55$ في الثانية، فوق الحد $52$ بقليل، فيتذبذب متناميًا.

ومن هنا نتيجتان. الأولى: كلما طالت الحلقة وجب أن تصحح أبطأ. فعبر العين بتأخير مئة وخمسين ملّي ثانية لا يمكن تصحيح اليد أسرع من نحو عُشر ثانية، وإلا ارتجفت. والثانية: حد استقرار الحلقة الشوكية، $1/(4\cdot30\ms)\approx8$ ذبذبات في الثانية، قريب من تردد الرعشة الدقيقة الموجودة لدى كل إنسان سليم، أي من ثماني إلى اثنتي عشرة ذبذبة في الثانية. وحلقة التمدد أحد المصادر المحتملة لهذه الرعشة~\cite{McAuleyMarsden2000}.

فكيف يتحرك الدماغ إذن بسرعة ودقة مع أن التغذية الراجعة متأخرة؟ بحسب فرضية شائعة: \emph{بالتنبؤ}. فهو يحتفظ بنموذج داخلي للجسم ويحسب نتيجة الأمر دون انتظار المستشعرات~\cite{WolpertGhahramani2000}. والمستشعرات لازمة لتصحيح النموذج، لا لتصحيح كل حركة. وكان المهندس سيسمي هذا تحكمًا تنبؤيًا مع مراقِب للحالة.

\section{مولّدات إيقاع الحركات}

المشي والتنفس والمضغ حركات إيقاعية. فهل تحتاج إلى ساعة؟ لا. في الحبل الشوكي وجذع الدماغ شبكات صغيرة تولّد الإيقاع بنفسها، حتى لو لم يصلها أي شيء إيقاعي~\cite{MarderBucher2001}. وأبسط شبكة كهذه مجمَّعة من قطع لدينا من قبل: مجموعتان من عصبونات \nt{GLUT} تثبّط كل منهما الأخرى عبر وسطاء \nt{GABA} أو \nt{GLYC}، كما في الشكل~\ref{fig:wta}، وتتعبان ببطء، كذاكرة الفصل~\ref{sec:glutcomputer}:
\begin{empheq}[box=\keybox]{equation}
\tau\,\frac{\dd r_i}{\dd t} = -r_i + \bigl[I - \beta\,r_j - a_i\bigr]_+ ,
\qquad
\tau_a\,\frac{\dd a_i}{\dd t} = -a_i + b\,r_i ,
\label{eq:halfcentre}
\end{empheq}
حيث $a_i$ تعب المجموعة $i$، و$j$ المجموعة الأخرى. التثبيط المتبادل مع $\beta>1$ يختار فائزًا (القضية~\ref{prop:wta}). يعمل الفائز ويتعب؛ وحين يأكل التعب دخله، ينطلق الخاسر، وقد استراح، إلى المقدمة ويبدأ هو بالتعب. تعمل المجموعتان بالتناوب، وتتحكم كل منهما في عضلتها من الزوج (الشكل~\ref{fig:halfcentre}).

\begin{figure}[t]
\centering
\begin{tikzpicture}[>=Stealth,x=3.2cm,y=2.4cm]
\draw[->,gray!70!black] (0,0) -- (4.2,0) node[right,font=\small,black] {$t,\ \text{s}$};
\draw[->,gray!70!black] (0,0) -- (0,1.2) node[left,font=\small,black] {$r$};
\foreach \t in {0,1,2,3,4}{\draw[gray!60!black] (\t,-0.03) -- (\t,0.03); \node[below,font=\scriptsize] at (\t,0) {\t};}
\draw[thick,blue!60!black] plot file {figures/halfcentre1.dat};
\draw[thick,orange!85!black] plot file {figures/halfcentre2.dat};
\node[font=\scriptsize,blue!60!black,anchor=west] at (1.6,1.28) {المجموعة 1: ”إلى الأمام“};
\node[font=\scriptsize,orange!85!black,anchor=west] at (2.7,1.28) {المجموعة 2: ”إلى الخلف“};
\end{tikzpicture}
\caption{مولّد الإيقاع بحسب النموذج~\eqref{eq:halfcentre}: $I=1$، $\beta=2$، $b=2$، $\tau=20\ms$، $\tau_a=0.4\,\text{s}$. مجموعتان بتثبيط متبادل وتعب تعملان بالتناوب بدور $0.84\,\text{s}$، مع أن الدخل إشارة ثابتة.}
\label{fig:halfcentre}
\end{figure}

في نموذجنا، مع دخل ثابت، تتناوب المجموعتان بدور $0.84\,\text{s}$، أي نحو ضعف زمن التعب $\tau_a$. الأمر الثابت من الأعلى، ”امشِ“، يتحول إلى إيقاع في موضعه. وهذا إيقاع محلي آخر، كإيقاع حلقة \foreignlanguage{english}{\nt{GLUT}--\nt{GABA}} في الفصل~\ref{sec:gaba}: لا ينشأ إلا حين تعمل الشبكة، ولا يضبط الإيقاع إلا للعضلات التي تتحكم فيها.

\begin{keyblock}
\begin{proposition}[مخرج الآلة]
\label{prop:muscles}
يتحكم الدماغ في العضلات كهرمية من المنظِّمات. في الأعلى يُختار الفعل (الفصل~\ref{sec:dofa})؛ وأدنى منه تحوّل المولّدات المحلية الأوامرَ الثابتة إلى إيقاع، وتمسك الحلقاتُ السريعة في الحبل الشوكي المفاصلَ. وتُحدَّد القوة بفاصلة عائمة: الأس بعدد الوحدات المشغَّلة بحسب الحجم، والجزء العشري بتردد النبضات. وكل مفصل يتحكم فيه زوج من العضلات الساحبة، يحدد الفرق بين قوتيهما العزمَ، ومجموعهما الصلابةَ. هذه الأجهزة مقيسة؛ أما اعتماد الدماغ على نموذج داخلي للجسم ففرضية راسخة الأسس.
\end{proposition}
\end{keyblock}

\section{الخلاصة}

\begin{table}[t]
\centering
\footnotesize
\setlength{\tabcolsep}{4pt}
\renewcommand{\arraystretch}{1.15}
\begin{tabular}{>{\raggedright}p{3.0cm}>{\raggedright}p{4.6cm}>{\raggedright\arraybackslash}p{5.2cm}}
\toprule
ماذا يفعل المخرج & كيف & ما هو معروف \\
\midrule
يحوّل النبضات إلى قوة & مجموع النفضات~\eqref{eq:forcesum}، مرشح تمرير منخفض & مقيس؛ والجمع الخطي تبسيط \\
يجرّع القوة & تشغيل الوحدات بحسب الحجم، فاصلة عائمة~\eqref{eq:sizeprinciple} & ترتيب التشغيل مقيس \\
يدفع وهو لا يستطيع إلا السحب & زوج عضلات: العزم فرق، والصلابة مجموع~\eqref{eq:antagonists} & مقيس \\
يمسك المفصل & حلقة التمدد، منع المضادة عبر \nt{GLYC} & مقيس \\
لا يرتجف & $Gd<\pi/2$~\eqref{eq:delaystable} & ينتج من نظرية التحكم؛ والإسهام في الرعشة سبب محتمل \\
يتحرك بسرعة & التنبؤ بالنموذج الداخلي & فرضية راسخة الأسس \\
يمشي ويتنفس بإيقاع & مولّد الإيقاع~\eqref{eq:halfcentre} & المولّدات مقيسة؛ والمخطط الأبسط نموذج \\
\bottomrule
\end{tabular}
\caption{كيف تتحكم الآلة في العضلات.}
\label{tab:muscles}
\end{table}

مخرج الآلة مبني بالحيل نفسها التي بُني بها كل شيء آخر (الجدول~\ref{tab:muscles}): سلكان بدل العلامة، ومرشح تمرير منخفض بدل المحوِّل الرقمي التماثلي، ومقياس لوغاريتمي بدل الخطي، وتثبيط متبادل وتعب بدل مولّد نبضات الساعة. المدخل (الفصل~\ref{sec:sensors}) والمخرج يقيسان العالم بمقاييس نسبية، وبينهما تعمل الآلة التي خُصص لها هذا الكتاب.

\section{تمارين}

\begin{exercise}[\lvlA]
\label{ex:13:1}
\emph{الفاصلة العائمة بالأرقام.} في العضلة $N=100$ وحدة حركية قواها $f_n=f_1q^{\,n-1}$؛ أقواها أقوى من أضعفها $100$ مرة. (أ)~أوجد $q$، والخطوة النسبية~\eqref{eq:sizeprinciple}، والمدى الديناميكي بالديسيبل. (ب)~ما الخطوة عند القوة $10f_1$ في ”محوِّل رقمي تماثلي خطي“ من $100$ وحدة متماثلة بالقوة الكلية نفسها؟
\end{exercise}

\begin{exercise}[\lvlA]
\label{ex:13:2}
\emph{ثمن هامش الأمان.} مع كل نبضة يطلق المشبك على الليف العضلي عددًا بواسونيًا من الحزم متوسطه $m$، والليف يطلق عند $n_{\text{th}}=20$ حزمة فأكثر؛ وهامش الأمان $S=m/n_{\text{th}}$. كم إخفاقًا في اليوم تعطي وحدة تعمل بتردد $20$ نبضة في الثانية عند $S=2$ وعند $S=4$؟
\end{exercise}

\begin{exercise}[\lvlB]
\label{ex:13:3}
\emph{العضلة مرشحًا.} النفضة~\eqref{eq:twitch} مع $T_c=50\ms$ هي الاستجابة النبضية لنظام خطي. (أ)~أوجد دالة تحويله $H(s)$ وتردد القطع. (ب)~عند أي تردد للنبضات المنتظمة يبلغ مدى تموجات القوة $10\%$ من متوسط القوة؟
\end{exercise}

\begin{exercise}[\lvlB]
\label{ex:13:4}
\emph{لا تصحيح أسرع من التأخير.} المنظِّم $\dd x/\dd t=-G\,x(t-d)$. (أ)~بيّن أن الانحراف يخمد بلا تذبذب بأسرع ما يمكن عند $Gd=1/e$، بمعدل $1/d$. (ب)~أوجد هذا $G$ وثابت زمن الخمود للحلقة الشوكية، $d=30\ms$، وللحلقة عبر العين، $d=150\ms$.
\end{exercise}

\begin{exercise}[\lvlB]
\label{ex:13:5}
\emph{دور مولّد الإيقاع.} عند $\tau\ll\tau_a$ يُحدَّد الدور $T$ للنموذج~\eqref{eq:halfcentre} بالمعادلة $(\beta-1)(1-E^{2+b})/(\beta-E)=b\,(1-E^{1+b})/(1+b)$، حيث $E=e^{-T/(2\tau_a)}$. (أ)~أوجد $T$ عند $\beta=b=2$ و$\tau_a=0.4\,\text{s}$. (ب)~بيّن أن $T$ لا يتوقف على الدخل $I$، وأن الإيقاع ممكن فقط عند $b>\beta-1$.
\end{exercise}

\begin{exercise}[\lvlB]
\label{ex:13:6}
\emph{نمذجة المولّد.} استورد الدالة \foreignlanguage{english}{\texttt{halfcentre}} من \foreignlanguage{english}{\texttt{models/\allowbreak muscle\_models.py}} (لا تشغّل الملف نفسه)، وقس الدور، مهملًا الدورتين الأوليين، عند $b=0.8$، $1.05$، $1.2$، $1.5$، $2$، $4$ وعند $I=0.5$، $1$، $2$. هل يستطيع الأمر ”امشِ أسرع“ أن يغيّر الوتيرة عبر $I$؟
\end{exercise}

\begin{exercise}[\lvlB]
\label{ex:13:7}
\emph{الصلابة في مواجهة المنعكس.} يحيط بالمفصل حلقة تمدد: $\dd\theta/\dd t=-a\theta(t)-G\theta(t-d)$، $a=K/B$، $d=30\ms$. ردّها إلى التمرين~\ref{ex:13:4}، وأوجد أصغر $a$ يخمد عنده الانحراف بلا تذبذب بثابت زمن $15\ms$، و$G$ اللازم. كيف يجب تغيير $G$ عند زيادة الشد المشترك للعضلات؟
\end{exercise}

\begin{exercise}[\lvlC]
\label{ex:13:8}
\emph{مقبض الوتيرة.} بحسب التمرينين~\ref{ex:13:5} و\ref{ex:13:6} لا يغيّر الدخل $I$ للمولّد~\eqref{eq:halfcentre} إلا السعة، لا الوتيرة. اقترح أدنى تعديل للنموذج من قطع الكتاب، بحيث لا يوجد إيقاع دون $I_{\min}$ ما، ويتناقص الدور فوقها مع ازدياد $I$، إلى النصف على الأقل عند مضاعفة $I$ ثلاث مرات. أوجد $I_{\min}$ عند $\tau\ll\tau_a$ وتحقق بالنمذجة.
\end{exercise}

\chapter{الألم: إشارة إنذار}
\chaptermark{الألم}
\label{sec:pain}

\section{إنذار لا قياس}

كل مستشعرات الفصل~\ref{sec:sensors} كانت تقيس العالم: كم من الضوء، وما حدة الصوت، وما الضغط. أما الألم فمبني على نحو آخر. إنه لا يُبلغ ”ماذا هناك“، بل ”خطر، اترك كل شيء“. وهو عند المهندس ليس قناة قياس، بل \emph{إشارة إنذار}: خط ذو أولوية عليا يجب أن يخترق أي عمل تؤديه الآلة في تلك اللحظة.

يختلف نظام الإنذار عن قناة القياس بثلاث خصائص، وكلها في الألم. أولًا، يصعب إخماده بالتعوّد: مستشعرات الألم تتعوّد أضعف بكثير من مستشعرات اللمس، بل تصبح بعد إصابة خفيفة أشد حساسية~\cite{LaMotte1982hyper}؛ ولم يُقَس ضعف الاستجابة بعد تسخين قوي إلا لدى نوع واحد منها~\cite{Treede1995heat}. مستشعر الضوء يصمت على خلفية ثابتة (الشكل~\ref{fig:adaptation})، أما جهاز إنذار يصمت بعد دقيقة فلا فائدة منه. ثانيًا، عتبته عالية: لا تطلق مستشعرات الألم إلا حين يقترب المؤثر من حد الخطر، فللتسخين مثلًا تتراوح عتبات المستشعرات المختلفة من $41^\circ$ إلى $46\text{--}48^\circ$ بحسب النوع~\cite{Treede1995heat, Weidner1999cfib}. ثالثًا، وهذا هو الأهم في هذا الفصل، لا يقرر علوَّ صوت الإنذار المستشعرُ وحده، بل الآلة نفسها أيضًا. فالجرح نفسه يؤلم على نحو مختلف في ظروف مختلفة. فلنرَ من أي قطع مألوفة لنا جُمع ذلك.

مستشعرات الألم عصبونات \nt{GLUT}، كبقية العصبونات الحسية في الفصل~\ref{sec:sensors}. وبعضها يطلق مع الغلوتامات المادة~\foreignlanguage{english}{P} من الملحق~\ref{sec:othermediators}، التي تقوّي النقل ببطء.

\section{سلكان: ألم سريع وألم بطيء}

اضرب إصبعك، وستشعر بالألم مرتين: أولًا ألم قصير حاد، ثم بعد ثانية ألم كليل طويل. هذان سلكان مختلفان. السلك السريع معزول، ويوصل النبضات بسرعة $v$ من $5$ إلى $30$ مترًا في الثانية؛ والبطيء رفيع وعارٍ، $0.5\text{--}2$ متر في الثانية. تقطع الإشارة من القدم نحو متر، وزمن الوصول
\begin{empheq}[box=\keybox]{equation}
t \;=\; \frac{L}{v}
\label{eq:paintime}
\end{empheq}
للسلك السريع عشرات الملّي ثانية، وللبطيء نحو ثانية. يحمل السلكان رسالتين. السريع يقول \emph{أين} و\emph{متى}: نبضاته تصل أبكر وأدق، كما في شفرة التوقيت من الفصل~\ref{sec:code}. والبطيء يقول \emph{ما مدى السوء}، وألمه الكليل الطويل يُبقي الآلة في نمط الإنذار حتى تزول الإصابة.

\section{البوابة في الحبل الشوكي}

أول ما تصل إليه إشارة الألم الحبل الشوكي، وهناك تمر عبر \emph{بوابة}~\cite{MelzackWall1965}؛ وقد وُجدت العصبونات المحددة لهذه البوابة حديثًا نسبيًا~\cite{Duan2014}. على عصبون \nt{GLUT} المخرج، الذي ينقل الألم إلى الدماغ، يلتقي سلك الألم وسلك اللمس. وبجانبه عصبون وسيط مثبِّط، \nt{GABA} أو \nt{GLYC}، يثيره اللمس (الشكل~\ref{fig:gate}). فاللمس يغلق البوابة. وفي نظرية البوابة الأصلية~\cite{MelzackWall1965} يُطفئ الألمُ على العكس هذا الوسيط، فيفتح الألم القوي البوابة؛ وهذا افتراض في النظرية، لم يُبيَّن مباشرة على عصبونات البوابة المكتشفة.

\begin{figure}[t]
\centering
\begin{tikzpicture}[>=Stealth,
  nrn/.style={circle,draw,very thick,minimum size=1.0cm,inner sep=0pt,font=\scriptsize},
  inh/.style={-{Bar[width=7pt]},thick,red!70!black},
  bc/.style={->,thick,densely dashed,orange!85!black}]
\node[nrn,fill=blue!6] (Z) at (6,0) {\nt{GLUT}};
\node[nrn,fill=red!10] (Q) at (3.6,-1.6) {\nt{GABA}};
\draw[->,very thick,red!70!black] (0,0.6) node[left,font=\small,black,align=right] {الألم $\nu$} -- (5.45,0.15);
\draw[->,very thick,blue!60!black] (0,-2.6) node[left,font=\small,black,align=right] {اللمس $\mu$} -- (2.9,-2.0);
\draw[->,thick,blue!60!black] (1.8,-2.35) -- (5.5,-0.35) node[pos=0.8,below right=-2pt,font=\scriptsize] {$w_{z\mu}$};
\draw[inh] (1.6,0.5) -- (3.35,-1.1) node[pos=0.5,left=1pt,font=\scriptsize] {$w_{q\nu}$};
\draw[inh] (Q) -- (Z) node[pos=0.5,above left=-1pt,font=\scriptsize] {$\beta$};
\draw[->,very thick] (Z) -- (8.2,0) node[right,font=\small,align=left] {إلى الدماغ: $z$};
\draw[bc] (6,2.1) node[above,font=\scriptsize,black,align=center] {من الأعلى: \nt{SERO}، \nt{NORA}، الأفيونيات} -- (Z.north) node[pos=0.45,right,font=\scriptsize,black] {الكسب $g$};
\end{tikzpicture}
\caption{بوابة الألم في الحبل الشوكي بحسب النموذج~\eqref{eq:gate}. يتلقى عصبون \nt{GLUT} المخرج الألم $\nu$ وإثارة ضعيفة من اللمس $\mu$. الوسيط المثبِّط (\nt{GABA} أو \nt{GLYC}) يثيره اللمس، ويطفئه الألم بحسب افتراض نظرية البوابة. ومن الأعلى، عبر قنوات البث، يأتي الكسب $g$.}
\label{fig:gate}
\end{figure}

لنكتب ذلك للترددات، كما في الفصل~\ref{sec:gaba}. تردد الوسيط $q$ وتردد المخرج $z$ هما
\begin{empheq}[box=\keybox]{equation}
q \;=\; \bigl[w_{q\mu}\,\mu + q_0 - w_{q\nu}\,\nu\bigr]_+ ,
\qquad
z \;=\; \bigl[\,g\,(\nu + w_{z\mu}\,\mu) - \beta\,q\,\bigr]_+ ,
\label{eq:gate}
\end{empheq}
حيث $\nu$ دخل الألم، و$\mu$ دخل اللمس، و$w$ أوزان الوصلات (الدليل الأول: إلى مَن، والثاني: من مَن)، و$\beta$ شدة التثبيط، و$q_0$ النشاط الخلفي الذاتي للوسيط، و$g$ الكسب الذي يُحدَّد من الأعلى (عنه أدناه). هذا هو المنع~\eqref{eq:veto} من الفصل~\ref{sec:gaba}: ”الألم، ولكن ليس مع اللمس“، مع تعديل واحد مأخوذ من نظرية البوابة افتراضًا: الألم القوي يطفئ بنفسه العصبون المانع.

حسبنا النموذج عند $w_{q\mu}=1$، $q_0=0.2$، $w_{q\nu}=0.5$، $w_{z\mu}=0.3$، $\beta=1.5$ (الشكل~\ref{fig:gatecurves}). بلا لمس يبدأ الألم بالمرور من $\nu\approx0.18$. ومع اللمس ترتفع العتبة إلى $0.86$، وعند $\nu=1$ يهبط المخرج أربع مرات، من $1$ إلى $0.25$. أما مع الألم القوي، $\nu=2$، فلا يغيّر اللمس شيئًا: البوابة مفتوحة على مصراعيها. وهكذا في الحياة: فرك الكدمة يساعد، أما في الإصابة الخطيرة فلا. واللمس نفسه، بلا ألم، لا يمر عبر البوابة: الإنذار لا ينطلق من اللمس.

\begin{figure}[t]
\centering
\begin{tikzpicture}[>=Stealth,x=5.2cm,y=1.35cm]
\draw[->,gray!70!black] (0,0) -- (2.12,0) node[right,font=\small,black] {الألم $\nu$};
\draw[->,gray!70!black] (0,0) -- (0,4.3) node[left,font=\small,black] {$z$};
\foreach \t in {0,0.5,1,1.5,2}{\draw[gray!60!black] (\t,-0.05) -- (\t,0.05); \node[below,font=\scriptsize] at (\t,0) {\t};}
\foreach \f in {1,2,3,4}{\draw[gray!60!black] (-0.01,\f) -- (0.01,\f); \node[left,font=\scriptsize] at (0,\f) {\f};}
\draw[very thick,red!70!black] plot file {figures/pain_gate_notouch.dat};
\draw[very thick,blue!60!black] plot file {figures/pain_gate_touch.dat};
\draw[very thick,orange!85!black,densely dashed] plot file {figures/pain_gate_sens.dat};
\draw[very thick,red!70!black] (0.08,3.9) -- (0.2,3.9) node[right,font=\scriptsize,black] {بلا لمس};
\draw[very thick,blue!60!black] (0.08,3.45) -- (0.2,3.45) node[right,font=\scriptsize,black] {مع اللمس};
\draw[very thick,orange!85!black,densely dashed] (0.08,3.0) -- (0.2,3.0) node[right,font=\scriptsize,black] {بعد التحسّس، $g=2$};
\end{tikzpicture}
\caption{مخرج البوابة~\eqref{eq:gate} بحسب شدة الألم. اللمس يرفع العتبة ويخفت الألم الضعيف والمتوسط، لا القوي. والتحسّس (الخط المتقطع) يضاعف الكسب: الألم نفسه يُنقل بصوت أعلى مرتين، وتنخفض العتبة.}
\label{fig:gatecurves}
\end{figure}

\section{مقبض الصوت من الأعلى}

لا تُدار البوابة من الأسفل وحده. فمن جذع الدماغ تنزل إليها مسارات هابطة تغيّر الكسب $g$ في~\eqref{eq:gate}: \nt{SERO} و\nt{NORA} والببتيدات الأفيونية من الملحق~\ref{sec:othermediators}~\cite{Fields2004}. هذه هي قنوات البث نفسها التي في الفصول من~\ref{sec:serocomputer} إلى~\ref{sec:acet}، إلا أن مقبضها هنا علوّ صوت الإنذار. وهي تستطيع إدارته في الاتجاهين: الدارات الهابطة نفسها قادرة على إخفات الألم وعلى تضخيمه.

لماذا تخفض الآلة الإنذار؟ لأن الإنذار مقاطعة، والمقاطعة ليست مناسبة دائمًا. الحيوان الهارب من مفترس يجب ألا يتوقف بسبب قائمة جريحة: الركض أولًا، ثم لعق الجرح. وتوقع الراحة أيضًا يخفت الألم، ويزول جزء من هذا الأثر إذا حُجبت المستقبِلات الأفيونية~\cite{Levine1978}: فالتوقع المحسوب في الدماغ ينزل عبر القناة الأفيونية إلى البوابة. هذه الصورة بحد ذاتها مقيسة؛ أما كيف يقرر الدماغ بالضبط كم يجب أن يكون علو الصوت ففرضية.

\section{التحسّس: إنذار يتعلم}

بعد الإصابة تصبح عصبونات المخرج في البوابة أشد حساسية: الدخل نفسه يستدعي مخرجًا أكبر، والعتبة تنخفض~\cite{Woolf1983}. في النموذج~\eqref{eq:gate} هذا ازدياد في الكسب $g$، وأبسط وصف له دارة \foreignlanguage{english}{RC} بطيئة يشحنها الألم نفسه:
\begin{empheq}[box=\keybox]{equation}
\tau_s\,\frac{\dd g}{\dd t} \;=\; -(g-1) + k_s\,z ,
\label{eq:sensitization}
\end{empheq}
حيث $\tau_s$ الزمن الذي تعود فيه الحساسية إلى طبيعتها، وهو أطول من مدة الألم نفسه، لأن التحسّس يدوم حتى بعد توقف المنبّه المؤذي~\cite{Woolf1983}، و$k_s$ قوة الشحن. ما دام النسيج مصابًا يُبقي مخرج البوابة $g$ فوق الواحد، فينتقل كل ما يحدث قرب الجرح بصوت أعلى (الخط المتقطع في الشكل~\ref{fig:gatecurves}: عند $g=2$ تنخفض العتبة إلى $0.11$، ويتضاعف المخرج). هذا ضبط معقول لجهاز الإنذار: ما دامت الإصابة لم تلتئم، فالاحتياط أفضل.

هذه عند المهندس تغذية راجعة موجبة بتسريب بطيء: الألم يرفع الكسب، والكسب يرفع الألم. ما دامت الحلقة ضعيفة يعود $g$ إلى طبيعته حين يلتئم الجرح. أما إذا كانت الحلقة قوية، فقد يعلق جهاز الإنذار في حالة التشغيل حتى بعد زوال الإصابة، كالمجموعة ذات التغذية الراجعة القوية في الشكل~\ref{fig:bistable}. النموذج~\eqref{eq:sensitization} تبسيط؛ أما وجود التحسّس المركزي فمقيس.

\section{الألم معلّمًا ومقاطعةً}

للألم في الآلة عملان إلى جانب الإنذار.

\emph{المعلّم.} الألم مكافأة سالبة: يدخل في صيغة الخطأ~\eqref{eq:ctd} بوصفه $r<0$. وكل ما قيل في الفصل~\ref{sec:dofa} يصح هنا أيضًا: نذير الألم يبدأ باستدعاء الخطأ مسبقًا، فتتعلم الشبكة تجنب ما يقود إلى الألم. والتجارب على البشر تُظهر أن نشاط الدماغ في التعلم من الألم يتبع خطأ الفروق الزمنية بالذات~\cite{Seymour2004}، وأن بعض عصبونات \nt{DOPA} تستجيب للأحداث المزعجة أيضًا (الفصل~\ref{sec:dofaalt}).

\emph{المقاطعة.} الألم يصرف الانتباه عن المهمة الجارية، وهذه وظيفته لا أثر جانبي له~\cite{EcclestonCrombez1999}. وفي الفصل~\ref{sec:nora} نوقش دور ”المقاطعة“ نفسه للرشقة الطورية من \nt{NORA}. أما أسرع استجابة فلا تصل إلى الدماغ أصلًا: سحب اليد عن الساخن يُغلق مباشرة في الحبل الشوكي بزوج العضلات نفسه والمنع نفسه للعضلة المضادة كما في الشكل~\ref{fig:antagonists}.

\begin{keyblock}
\begin{proposition}[إشارة الإنذار]
\label{prop:pain}
الألم ليس قياسًا، بل إشارة إنذار ذات أولوية عليا. يتعوّد أضعف بكثير من قنوات القياس، ويسير في سلكين (السريع: ”أين“، والبطيء: ”ما مدى السوء“)، ويمر عبر بوابة يغلقها اللمس (ويفتحها الألم القوي بحسب افتراض نظرية البوابة)، وتحدد علوَّ صوته من الأعلى قنواتُ البث. هذه الأجهزة مقيسة، عدا فتح البوابة بالألم؛ والنموذجان البسيطان~\eqref{eq:gate} و\eqref{eq:sensitization} تبسيطان.
\end{proposition}
\end{keyblock}

\section{الخلاصة}

\begin{table}[t]
\centering
\footnotesize
\setlength{\tabcolsep}{4pt}
\renewcommand{\arraystretch}{1.15}
\begin{tabular}{>{\raggedright}p{3.0cm}>{\raggedright}p{4.9cm}>{\raggedright\arraybackslash}p{4.9cm}}
\toprule
ماذا يفعل الألم & كيف & ما هو معروف \\
\midrule
يطلق الإنذار & عتبة عالية، وتعوّد أضعف من اللمس & العتبة مقيسة؛ ومقارنة التعوّد تقدير \\
يُبلغ ”أين“ و”ما مدى السوء“ & سلكان بسرعتين مختلفتين~\eqref{eq:paintime} & مقيس \\
يمر عبر بوابة & منع عبر \nt{GABA}/\nt{GLYC}~\eqref{eq:gate} & البوابة مقيسة؛ وفتحها بالألم افتراض النظرية؛ والنموذج تبسيط \\
يغيّر علو الصوت & \nt{SERO} و\nt{NORA} والأفيونيات الهابطة & مقيس؛ وقاعدة اختيار علو الصوت فرضية \\
يتعلم بعد الإصابة & ازدياد الكسب~\eqref{eq:sensitization} & التحسّس مقيس؛ والنموذج تبسيط \\
يعلّم التجنب & مكافأة سالبة في~\eqref{eq:ctd} & التشابه مع خطأ الفروق الزمنية مقيس \\
يقاطع العمل & الاستيلاء على الانتباه؛ منعكس السحب & مقيس \\
\bottomrule
\end{tabular}
\caption{الألم إشارةَ إنذار.}
\label{tab:pain}
\end{table}

الألم مجمَّع من قطع رأيناها من قبل (الجدول~\ref{tab:pain}): مستشعرات \nt{GLUT}، وسلكان، ومنع عبر وسيط مثبِّط، ومقبض صوت يُدار بالبث، ودارة \foreignlanguage{english}{RC} بطيئة، وخطأ التنبؤ، والمقاطعة. والجديد فيه أمر واحد: إنه المدخل الوحيد للآلة الذي تحدد الآلة نفسها علوّ صوته، فهو جهاز إنذار يستطيع صاحبه أن يخفته وأن يعيد ضبطه.

\section{تمارين}

\begin{exercise}[\lvlA]
\label{ex:14:1}
\emph{سلكان.} تأتي الإشارة من القدم، $L=1\,\text{m}$. (أ)~تسير الدفعة في حزمة من الأسلاك من نوع واحد، تملأ سرعاتها المدى~\eqref{eq:paintime}. بكم تتمدد حين تصل إلى الحبل الشوكي، للنوع السريع وللبطيء؟ (ب)~موجتا الألم في سلكين سرعتاهما $15$ و$1\,\text{m/s}$ يمكن تمييزهما إذا كان الفرق $0.1\,\text{s}$ فأكثر. من أي مسافة تندمجان؟
\end{exercise}

\begin{exercise}[\lvlB]
\label{ex:14:2}
\emph{متى يؤلم اللمس.} البوابة~\eqref{eq:gate} بمعاملات النص، والكسب $g$ يزداد مع التحسّس. حتى أي $g$ لا يمر اللمس بلا ألم عبر البوابة مهما كانت شدة $\mu$؟ وعند أي $g$ يبدأ اللمس $\mu=1$ بالمرور؟
\end{exercise}

\begin{exercise}[\lvlB]
\label{ex:14:3}
\emph{استقرار التحسّس.} المداخل ثابتة، ومخرج البوابة خطي في $g$: $z=gA-C$، $A=\nu+w_{z\mu}\mu$، $C=\beta q\ge0$. (أ)~أوجد التوازن $g^*$ للنموذج~\eqref{eq:sensitization} وشرط استقراره. (ب)~عند $k_s=0.5$ أوجد شدة الألم التي ينفلت النموذج فوقها، بلا لمس ومع اللمس $\mu=1$.
\end{exercise}

\begin{exercise}[\lvlA]
\label{ex:14:4}
\emph{نذير الألم.} إشارة في اللحظة $0$ تتنبأ بألم في اللحظة $T=2\,\text{s}$، وفي~\eqref{eq:ctd} $r(t)=-P\,\delta(t-T)$، $\tau_\gamma=2\,\text{s}$. (أ)~أوجد مساحة دفقة $\delta$ لحظة الإشارة. (ب)~فعل في اللحظة $t_e=1\,\text{s}$ يلغي الألم. ما دفقة $\delta$ في هذه اللحظة، وماذا ستعلّم الشبكة؟
\end{exercise}

\begin{exercise}[\lvlB]
\label{ex:14:5}
\emph{إنذار ينغلق على نفسه.} أضف إلى البوابة من \foreignlanguage{english}{\texttt{models/\allowbreak pain\_gate.py}} ديناميكا~\eqref{eq:sensitization} وتشبعًا $z\le4$. اللمس $\mu=1$ موجود دائمًا، والإصابة $\nu=2$ تدوم زمنًا $T$، ثم $\nu=0$. أوجد بالنمذجة أصغر $T$ (بوحدات $\tau_s$) يبقى الإنذار بعده مشغَّلًا، عند $k_s=4$، $4.6$، $5$، $6$، $8$.
\end{exercise}

\begin{exercise}[\lvlB]
\label{ex:14:6}
\emph{تصميم البوابة.} عند $\beta=1.5$، $w_{z\mu}=0.3$، $g=1$ اختر في~\eqref{eq:gate} القيم $q_0$ و$w_{q\nu}$ و$w_{q\mu}$ بحيث تكون عتبة الألم $0.2$ بلا لمس و$1.0$ مع اللمس $\mu=1$، ويكفّ اللمس عن إخفات الألم عند $\nu_\times=2.5$. حتى أي كسب $g$ لا يمر اللمس بلا ألم عبر البوابة؟
\end{exercise}

\begin{exercise}[\lvlC]
\label{ex:14:7}
\emph{مقبض الصوت.} العمل يجلب قيمة $V$ في وحدة الزمن، والعمل مع إصابة $\nu$ يكلف $c\nu$، فتكون المقاطعة مجدية عند $\nu>V/c$؛ والألم يقاطع عند $z>\theta=0.5$. (أ)~أي كسب $g^*$ للبوابة~\eqref{eq:gate} بلا لمس يعطي هذه العتبة عند $V/c=0.25$، $0.5$، $2$؟ (ب)~من أي إشارات الكتاب يمكن للآلة أن تقدّر $V$ و$c$؟
\end{exercise}

\chapter{كيف تتعلم الآلة الحركة}
\chaptermark{كيف تتعلم الآلة الحركة}
\label{sec:motorlearning}
\begin{chaptergoals}
\item لماذا تحتاج الحركة إلى معلّم ثالث، سلك خطأ إلى كل مخرج، وما ثمنه؛
\item كيف تعطي طبقة توسيع هائلة من \nt{GLUT} وسلك خطأ لكل مخرج قاعدة المربعات الصغرى؛
\item كيف تصير الدارة بعد التعلّم نموذجًا داخليًا، مرشحًا تكيّفيًا يتنبأ بالجسم؛
\item لماذا تفسر عمليتا تعلّم، سريعة وبطيئة، الأثر اللاحق والاستعادة التلقائية.
\end{chaptergoals}

\section{ثلاثة معلّمين}

رأينا في الفصل~\ref{sec:muscles} كيف تحرك الآلة العضلات، وتوقفنا عند فرضية: الحركات السريعة الدقيقة تتطلب نموذجًا داخليًا للجسم يتنبأ بنتيجة الأمر~\cite{WolpertGhahramani2000}. فمن أين يأتي هذا النموذج؟ يجب أن يُتعلَّم، وأن يُتعلَّم على غير الطريقة التي تعلمنا بها حتى الآن.

مرّ في الكتاب معلّمان. الأول \emph{لا أحد}: قاعدة هيب (الفصل~\ref{sec:dofa}) تقوّي ما يحدث معًا كثيرًا، وتضغط المداخل. والثاني \emph{المكافأة}: \nt{DOPA} يبث إلى الجميع عددًا واحدًا، ”أفضل أو أسوأ مما كان متوقعًا“. أما الحركة فتحتاج إلى معلّم ثالث هو \emph{الخطأ}. أخطأت اليد الفنجان بسنتيمترين إلى اليسار: هذا ليس ”أسوأ“، بل متجه يقول لكل مخرج أين أخطأ وبكم. وكان المهندس سيسمي هذا التعلم الموجَّه.

الفرق بين المعلّم الثاني والثالث هو الفرق بين سلك واحد للجميع وسلك مستقل لكل مخرج. في القضية~\ref{prop:broadcastcost} كان التعلم بعدد مشترك واحد يبطؤ مع كل عصبون يؤثر ضجيجه في النتيجة. أما إذا تلقى كل مخرج $i$ خطأه الخاص $e_i$، فيمكن تغيير الأوزان بأبسط قاعدة:
\begin{empheq}[box=\keybox]{equation}
\frac{\dd w_{ij}}{\dd t} \;=\; -\eta\,e_i(t)\,x_j(t) .
\label{eq:lms}
\end{empheq}
هذه قاعدة المربعات الصغرى المعروفة منذ زمن في تقنية المرشحات التكيّفية~\cite{WidrowHoff1960}: يتغير الوزن بما يتناسب مع حاصل ضرب دخله في خطأ مخرجه، ويسير في المتوسط مع تدرج مربع الخطأ. ولا حاجة إلى ضجيج للبحث كما في الفصل~\ref{sec:dofa}: اتجاه التصحيح يأتي مباشرة عبر سلك الخطأ. ولا تنخفض السرعة مع عدد المخارج، لأن أخطاء الآخرين لا تصل إلى المشبك.

\begin{keyblock}
\begin{proposition}[ثلاثة معلّمين]
\label{prop:teachers}
قاعدة هيب تعلّم بلا معلّم ما هو متكرر؛ وإشارة \nt{DOPA} تعلّم بعدد واحد للشبكة كلها ما هو مجدٍ، وتبطؤ مع كل عصبون؛ وسلك الخطأ إلى كل مخرج يعلّم بالقاعدة~\eqref{eq:lms} ما هو دقيق، وسرعته لا تتوقف على عدد المخارج. وثمن الطريقة الثالثة سلك مستقل لكل مخرج، ومن يعرف الجواب الصحيح.
\end{proposition}
\end{keyblock}

\section{سلك الخطأ}

من يستطيع أن يعرف الجواب الصحيح للحركة؟ المستشعرات نفسها. إذا انحرفت اليد يسار الهدف أبلغت عن ذلك العينُ ومستشعرات التمدد من الفصل~\ref{sec:sensors}. والمطلوب دارة توصل هذا الخطأ إلى المخارج المعنية.

في الدماغ منطقة يبدو أنها مبنية لهذا بالضبط~\cite{Marr1969,Albus1971}. دارتها منتظمة على نحو غير مألوف، تكاد تكون بلورية، ومن السهل أن نتعرف فيها على القاعدة~\eqref{eq:lms} (الشكل~\ref{fig:errorwire}).

\emph{طبقة توسيع الدخل.} تصل الإشارات عن الأمر وعن حالة الجسم إلى طبقة هائلة من عصبونات \nt{GLUT} الصغيرة. عددها نحو سبعين مليارًا، أي أكثر مما في بقية الدماغ مجتمعة~\cite{Azevedo2009}. كل منها يمزج بضعة مداخل، فتنبسط الإشارة في متجه طويل جدًا. والمهندس يعرف هذه الحيلة: إذا رُفع بُعد الدخل، أمكن في الفضاء الجديد الحصول على أي علاقة مطلوبة تقريبًا بمجموع موزون بسيط~\cite{Cover1965}.

\emph{عصبونات المخرج.} يحسب المجموع الموزون نحو ثلاثين مليون عصبون مخرج كبير~\cite{Andersen1992cereb}، لكل منها نحو مئة ألف مدخل من طبقة التوسيع. وهي عصبونات \nt{GABA}: نتيجة التعلم تُنقل إلى الأمام لا بالإثارة بل بالتثبيط، كما في المنع من الفصل~\ref{sec:gaba}.

\emph{سلك الخطأ.} يتلقى كل عصبون مخرج مدخلًا آخر خاصًا: سلك \nt{GLUT} واحد بالضبط، يطلق نادرًا، نحو مرة في الثانية، لكنه يستدعي في كل مرة ومضة قوية في عصبون المخرج~\cite{Ito2001}. هذا هو $e_i$: احتمال الومضة يتوقف على اتجاه خطأ الحركة~\cite{Herzfeld2018}. وتزامن ومضة الخطأ مع نشاط المدخل يُضعف هذا المدخل، وهذا بالضبط الحد $-\eta\,e_i x_j$ في~\eqref{eq:lms}.

\begin{figure}[t]
\centering
\begin{tikzpicture}[>=Stealth,
  gran/.style={circle,draw,thick,fill=blue!6,minimum size=0.32cm,inner sep=0pt},
  outn/.style={circle,draw,very thick,fill=red!10,minimum size=1.1cm,inner sep=0pt,font=\scriptsize},
  inh/.style={-{Bar[width=7pt]},very thick,red!70!black}]
% inputs
\foreach \y/\lab in {1.6/{الأمر},0.4/{حالة الجسم}}{
  \draw[->,thick,blue!60!black] (-0.6,\y) node[left,font=\scriptsize,black] {\lab} -- (0.35,\y);}
% expansion layer
\foreach \i in {0,...,11}{\node[gran] (g\i) at ({0.8+0.28*mod(\i,2)},{-0.3+0.23*\i}) {};}
\foreach \i in {0,...,11}{\pgfmathsetmacro{\yy}{mod(\i,3)>0 ? 1.6 : 0.4}\draw[gray!60!black] (0.35,\yy) -- (g\i);}
\node[font=\scriptsize,align=center] at (0.95,-1.05) {طبقة التوسيع\\$\sim7\cdot10^{10}$ \nt{GLUT}};
% parallel wires to the output neuron
\foreach \i in {0,...,11}{\draw[->,gray!70!black] (g\i.east) -- (4.1,{1.0+0.07*(\i-5.5)});}
\node[outn] (P) at (4.6,1.0) {\nt{GABA}};
\node[font=\scriptsize,align=center,above=2pt] at (P.north) {عصبون المخرج\\$\sim10^5$ مدخل $x_j$، أوزان $w_{ij}$};
\draw[inh] (P) -- (6.8,1.0) node[right,font=\scriptsize,align=left] {تصحيح\\للحركة};
% error wire
\draw[->,very thick,red!70!black,densely dashed] (4.6,-1.4) node[below,font=\scriptsize,black,align=center] {سلك خطأ واحد $e_i$\\\nt{GLUT}، $\sim1$ مرة/ث} -- (P.south);
\end{tikzpicture}
\caption{دارة التعلم الموجَّه كما يبدو أن الدماغ ينفذها. تبسط طبقة هائلة من عصبونات \nt{GLUT} الأمرَ وحالةَ الجسم في متجه طويل $x_j$؛ ويجمعه عصبون \nt{GABA} المخرج بالأوزان $w_{ij}$. ويأتي سلك الخطأ الوحيد بـ$e_i$؛ وتزامن الخطأ مع مدخل نشط يُضعف وزن هذا المدخل، كما في القاعدة~\eqref{eq:lms}.}
\label{fig:errorwire}
\end{figure}

وهناك صعوبة نعرفها من الفصل~\ref{sec:dofa}: الخطأ يصل متأخرًا عن الدخل الذي قاد إليه. فالإخفاق يظهر بعد عشرات إلى مئات الملّي ثانية من الأمر. إذن يجب أن يتذكر المشبك أنه كان نشطًا، أي يلزم أثر، كما في القاعدة~\eqref{eq:threefactor}. وطول هذا الأثر، بحسب التجارب، مضبوط على تأخير التغذية الراجعة في كل مهمة: حيث يصل الخطأ بعد مئة ملّي ثانية يضعف أكثر ما يضعف المدخلُ الذي كان نشطًا قبله بمئة ملّي ثانية~\cite{Suvrathan2016}.

\begin{keyblock}
\begin{proposition}[سلك الخطأ]
\label{prop:errorwire}
في الدارة المبيّنة في الشكل~\ref{fig:errorwire} يتلقى كل عصبون مخرج سلك خطأ واحدًا بالضبط، وتزامن الخطأ مع المدخل يُضعف هذا المدخل. هذه قاعدة المربعات الصغرى~\eqref{eq:lms} مطبَّقة على دخل وُسِّع إلى بُعد هائل. بنية الدارة والإضعاف عند التزامن مقيسان؛ أما أن الدارة كلها تعمل تعلمًا موجَّهًا ففرضية رائدة.
\end{proposition}
\end{keyblock}

\section{النموذج الداخلي مرشحًا تكيّفيًا}

ماذا تتعلم دارة كهذه؟ أقرب جواب: التنبؤ. ليصل إلى الدخل الأمر $u(t)$ مارًّا بمجموعة من المرشحات بثوابت زمنية مختلفة، كما في طبقة التوسيع: $x_j(t)=(g_j*u)(t)$. والمخرج مجموعها الموزون، أي التنبؤ بما ستُبلغ عنه المستشعرات:
\begin{empheq}[box=\keybox]{equation}
\hat y(t) \;=\; \sum_j w_j\,(g_j*u)(t), \qquad e(t) \;=\; \hat y(t) - y(t) ,
\label{eq:adaptivefilter}
\end{empheq}
والأوزان تتعلم بحسب~\eqref{eq:lms}. هذا \emph{مرشح تكيّفي}: دارة تختار بنفسها دالة تحويلها لتكرر نظامًا مجهولًا~\cite{Fujita1982}. والنظام المجهول هنا هو الجسم نفسه بكتلته ومرونته وتأخيراته. والمرشح المتعلَّم هو النموذج الداخلي من الفصل~\ref{sec:muscles}: يقول ما ستُبلغ عنه المستشعرات دون انتظارها.

ونموذج كهذا مفيد مرتين. أولًا، يمكن تقديم تنبؤه بدل التغذية الراجعة المتأخرة، فلا يعود شرط الاستقرار~\eqref{eq:delaystable} يقيّد اليدين: يمكن التصحيح بسرعة، وفق النموذج. وثانيًا، الفرق بين المتنبَّأ به والمتلقَّى إشارة عن \emph{غير المتوقع}. فكل ما فعلته الآلة بنفسها تنبأ به النموذج وطرحه؛ ويبقى ما فعله العالم. ولهذا، بحسب إحدى الفرضيات، لا نستطيع أن ندغدغ أنفسنا~\cite{Blakemore1998}.

\section{حين يتغير العالم: السريع والبطيء}

لنختبر الدارة بتجربة سهلة الإعداد. يمدّ الإنسان يده نحو هدف، فيتغير العالم فجأة: مثلًا، تؤثر في اليد قوة جانبية. تخطئ الحركات الأولى، ثم يتناقص الخطأ. وإذا أُزيلت القوة، أخطأت اليد أولًا في الاتجاه \emph{الآخر}: لقد تعلم النموذج القوة ويستمر في تعويضها. وهذا دليل مباشر على أن المتعلَّم نموذج، لا مجرد أن الأمر ”صار ينجح“.

وتُظهر التجارب أمرًا أدق: تتعلم عمليتان في آن واحد، سريعة تتعلم بسرعة وتنسى بسرعة، وبطيئة تتعلم ببطء وتتذكر طويلًا~\cite{Smith2006}. ويُحدَّث التصحيحان بعد كل حركة، عند كل حدث:
\begin{empheq}[box=\keybox]{equation}
e = p - (x_f + x_s), \qquad x_f \leftarrow A_f\,x_f + B_f\,e, \qquad x_s \leftarrow A_s\,x_s + B_s\,e ,
\label{eq:tworate}
\end{empheq}
حيث $p$ الاضطراب، و$x_f$ و$x_s$ التصحيحان المتعلَّمان للعمليتين، و$A$ مقدار ما يُحفظ من التصحيح حتى الحركة التالية، و$B$ مدى قوة تعلمه من الخطأ. للعملية السريعة $B$ كبير و$A$ أصغر من الواحد بوضوح؛ وللبطيئة العكس.

\begin{figure}[t]
\centering
\begin{tikzpicture}[>=Stealth,x=0.034cm,y=1.9cm]
\draw[->,gray!70!black] (0,0) -- (355,0) node[right,font=\small,black] {الحركة};
\draw[->,gray!70!black] (0,-1.15) -- (0,1.25);
\foreach \v in {-1,1}{\draw[gray!60!black] (-3,\v) -- (3,\v); \node[left,font=\scriptsize] at (0,\v) {$\v$};}
\foreach \n in {100,200,300}{\draw[gray!60!black] (\n,-0.04) -- (\n,0.04); \node[below,font=\scriptsize] at (\n,0) {\n};}
\draw[thin,gray!70!black] plot file {figures/tworate_p.dat};
\draw[thick,orange!85!black] plot file {figures/tworate_fast.dat};
\draw[thick,blue!60!black] plot file {figures/tworate_slow.dat};
\draw[very thick,black] plot file {figures/tworate_net.dat};
\node[font=\scriptsize,gray!50!black] at (120,1.12) {الاضطراب $p$};
\node[font=\scriptsize,black,anchor=west] at (238,0.52) {المجموع: عودة};
\node[font=\scriptsize,blue!60!black,anchor=west] at (150,0.39) {البطيئة $x_s$};
\node[font=\scriptsize,orange!85!black,anchor=west] at (60,0.08) {السريعة $x_f$};
\draw[densely dashed,gray!60!black] (226,-1.15) -- (226,1.25);
\node[font=\scriptsize,align=center,anchor=south west] at (228,-1.1) {لا أخطاء\\بعد الآن};
\end{tikzpicture}
\caption{عمليتا التعلم بحسب النموذج~\eqref{eq:tworate}، $A_f=0.92$، $B_f=0.1$، $A_s=0.996$، $B_s=0.02$. مئتا حركة مع الاضطراب $p=1$، ثم ست مع الاضطراب المعاكس $p=-1$، حتى يعود المجموع إلى الصفر. بعد الخط المتقطع يتوقف وصول الأخطاء، فيعود المجموع بنفسه إلى التصحيح القديم: العملية السريعة نسيت الاضطراب المعاكس، والبطيئة تتذكر الأصلي.}
\label{fig:tworate}
\end{figure}

يفسر النموذج~\eqref{eq:tworate} تجربة يصعب فهمها بدونه (الشكل~\ref{fig:tworate}). في حسابنا يبلغ مجموع التصحيحين بعد مئتي حركة مع الاضطراب $0.84$، وتحمل العملية البطيئة معظمه، $0.63$. ثم تعيد ست حركات مع الاضطراب المعاكس المجموع إلى الصفر: السريعة ذهبت إلى السالب، $-0.53$، والبطيئة لا تزال موجبة، $0.46$. فإذا توقف الإبلاغ عن الأخطاء الآن، نسيت العملية السريعة تصحيحها خلال عشرات الحركات، والبطيئة خلال مدة أطول بكثير، فيرتفع المجموع \emph{بنفسه} إلى $0.37$ خلال أربعين حركة: تعود المهارة القديمة بلا أي تدريب. وهذه الاستعادة التلقائية مقيسة لدى البشر؛ أما النموذج ذو العملية الواحدة فلا يستطيع إعطاءها، إذ يخمد بلا أخطاء إلى الصفر ببساطة.

لماذا عمليتان؟ يعطي المرشح الأمثل من الفصل~\ref{sec:acet} جوابًا هندسيًا معقولًا. يتغير الجسم لأسباب مختلفة وبسرعات مختلفة: تتعب العضلات في دقائق، وتنمو وتضعف في أشهر. وإذا كانت الاضطرابات سريعة وبطيئة، فأفضل تقدير يتألف من مرشحين بثابتي زمن مختلفين، يلتقط كل منهما اضطرابه~\cite{Kording2007}. العملية السريعة تصحيح مؤقت، ”اليد تعبت“؛ والبطيئة ”اليد صارت أخرى“. هذا لا يزال فرضية، لكنه يفسر لماذا تضيع المهارة المتعلَّمة في يوم جزئيًا في اليوم التالي، ولماذا تُتعلَّم أسرع في المرة الثانية.

\section{الخلاصة}

\begin{table}[t]
\centering
\footnotesize
\setlength{\tabcolsep}{4pt}
\renewcommand{\arraystretch}{1.15}
\begin{tabular}{>{\raggedright}p{3.0cm}>{\raggedright}p{4.6cm}>{\raggedright\arraybackslash}p{5.2cm}}
\toprule
ماذا تفعل الدارة & كيف & ما هو معروف \\
\midrule
تتعلم بمعلّم & سلك خطأ إلى كل مخرج، القاعدة~\eqref{eq:lms} & بنية الدارة والإضعاف عند التزامن مقيسان؛ والتعلم الموجَّه فرضية رائدة \\
توسّع الدخل & سبعون مليار عصبون \nt{GLUT} & عدد العصبونات مقيس؛ ودور التوسيع نظرية \\
تراعي تأخير الخطأ & أثر مضبوط على التأخير & مقيس \\
تبني نموذجًا داخليًا & مرشح تكيّفي~\eqref{eq:adaptivefilter} & فرضية راسخة الأسس \\
تتعلم بوتيرتين & عمليتان سريعة وبطيئة~\eqref{eq:tworate} & الأثر اللاحق والاستعادة التلقائية مقيسان؛ والعمليتان نموذج \\
\bottomrule
\end{tabular}
\caption{كيف تتعلم الآلة الحركة.}
\label{tab:motorlearning}
\end{table}

صار للآلة الآن ثلاثة معلّمين (الجدول~\ref{tab:motorlearning}). هيب يضغط ما هو متكرر؛ و\nt{DOPA} يعلّم بعدد واحد اختيار المجدي؛ وسلك الخطأ يعلّم الدقة، ويعلّمها بسرعة، لأن كل مخرج يصله خطؤه الخاص. ويبدو أن الدماغ يستعمل الثلاثة، كلًّا حيث يكون أرخص: إشارة البث للقرارات القليلة، والأسلاك المستقلة للضبط الدقيق للجسم، حيث تُبلغ المستشعرات نفسها عن الجواب الصحيح.

\section{تمارين}

\begin{exercise}[\lvlA]
\label{ex:15:1}
\emph{سعة الدارة ذات سلك الخطأ.} في الدارة $3\cdot10^7$ عصبون مخرج لكل منها $10^5$ مدخل، ولكل منها سلك خطأ خاص (نبضة واحدة في الثانية). العصبون ذو $n$ مدخلًا يتعلم أي تصنيف لـ$\approx2n$ متجهًا~\cite{Cover1965}. (أ)~كم ارتباطًا تخزن الدارة؟ (ب)~قدّر بحسب~\eqref{eq:spikeentropy} عند $\Delta t=10\ms$ أقصر زمن لملء سعة العصبون.
\end{exercise}

\begin{exercise}[\lvlB]
\label{ex:15:2}
\emph{سرعة قاعدة المربعات الصغرى.} أنماط القاعدة~\eqref{eq:lms} تخمد بسرعات $\eta\lambda_k(C)$. لخمسة مرشحات~\eqref{eq:adaptivefilter} على ضجيج أبيض $C_{jk}=2\sqrt{\tau_j\tau_k}/(\tau_j+\tau_k)$. أوجد نسبة أسرع السرعات إلى أبطئها إذا تزايدت $\tau_j$ بعامل $2$ ($10$ إلى $160\ms$) وبعامل $4$.
\end{exercise}

\begin{exercise}[\lvlB]
\label{ex:15:3}
\emph{تحليل العمليتين.} اكتب النموذج~\eqref{eq:tworate} بمعاملات الشكل~\ref{fig:tworate} على الصورة $\mathbf x_{n+1}=M\mathbf x_n+\mathbf b\,p$. (أ)~أوجد من القيم الذاتية لـ$M$ الثوابت الزمنية مقيسة بالحركات. (ب)~أوجد قيمتي $x_f$ و$x_s$ المستقرتين ومجموعهما عند $p=1$ ثابت.
\end{exercise}

\begin{exercise}[\lvlB]
\label{ex:15:4}
\emph{نمذجة: المرة الثانية أسرع.} بالنموذج~\eqref{eq:tworate} بمعاملات من \foreignlanguage{english}{\texttt{models/\allowbreak motor\_learning.py}} أوجد عدد الحركات حتى $x_f+x_s\ge0.8$ عند $p=1$: (أ)~لآلة غير مدرَّبة؛ (ب)~بعد $200$ حركة مع $p=1$ واضطراب معاكس حتى صفر المجموع، كما في الشكل~\ref{fig:tworate}. هل يمكن أن يتسارع هكذا نموذج بعملية واحدة؟
\end{exercise}

\begin{exercise}[\lvlB]
\label{ex:15:5}
\emph{منظِّم بنموذج داخلي.} الجسم المتحكَّم فيه $\dd x/\dd t=ku$ يُرى بتأخير $d$؛ والمنظِّم $u=-G\hat x$ يتنبأ بـ$\hat x(t)=x(t-d)+\hat k\int_{t-d}^{t}u\,\dd s$. (أ)~عند $\hat k=k=1$ و$d=150\ms$ أوجد $G$ لثابت زمن $20\ms$ وقارنه بالحد~\eqref{eq:delaystable}. (ب)~هل الدارة مستقرة عند $\hat k=0.4k$؟
\end{exercise}

\begin{exercise}[\lvlB]
\label{ex:15:6}
\emph{المرشح التكيّفي يتعلم العضلة.} الجسم المتحكَّم فيه نفضة~\eqref{eq:twitch} مساحتها واحد، $T_c=50\ms$؛ والمرشحات~\eqref{eq:adaptivefilter} دارات \foreignlanguage{english}{RC} بـ$\tau_j=12.5$، $25$، $50$، $100$، $200\ms$؛ والدخل عدد طبيعي التوزيع جديد كل $10\ms$. في كم ثانية تخفض القاعدة~\eqref{eq:lms} الخطأ إلى $0.5\%$ عند $\eta=50$ و$200$؟ ولماذا تبقى الأوزان بعيدة عن الأفضل حتى بعد $300\,\text{s}$؟
\end{exercise}

\begin{exercise}[\lvlC]
\label{ex:15:7}
\emph{العمليتان مرشحًا أمثل.} الاضطراب مجموع عمليتين $p^{f,s}_{n+1}=A_{f,s}\,p^{f,s}_n+w^{f,s}_n$ بتباينَي ضجيج $q_f$ و$q_s$، والخطأ يُرصد بضجيج تباينه $R$؛ ومرشح كالمان يعطي~\eqref{eq:tworate}. عند أي $q_f/R$ و$q_s/R$ ينتج من $A_f=0.92$ و$A_s=0.996$ أن $B_f=0.1$ و$B_s=0.02$؟ وكيف يتغير $B_f$ و$B_s$ إذا ضوعف $q_s$ أربع مرات؟
\end{exercise}

\chapter{المخرج الثاني: كيف يتحكم الدماغ في كيمياء الجسم}
\chaptermark{المخرج الكيميائي}
\label{sec:chemoutput}

\section{مخرج سريع ومخرج بطيء}

رأينا في الفصل~\ref{sec:muscles} المخرج السريع للآلة: العضلات. لكن للآلة مخرجًا ثانيًا بطيئًا. فالدماغ يُبقي حرارة الجسم ومعدل ضربات القلب وكمية الماء والسكر ضمن الحدود المطلوبة، ويهيّئ الجسم للركض أو للنوم. ولا يحرّك لذلك المفاصل، بل يطلق إشارات كيميائية. وللإشارات طريقان. الأول أعصاب تمتد مباشرة إلى الأعضاء الداخلية: إلى القلب والأوعية الدموية والأمعاء والغدد. والثاني الدم: بعض العصبونات والغدد لا تطلق مادتها في الشق نحو الجار، بل في مجرى الدم مباشرة.

الدم أوسع نواقل البث بين كل ما مرّ في هذا الكتاب. فقنوات البث في الفصول~\babelsublr{\ref{sec:serocomputer}--\ref{sec:acet}} كانت ترسل الإشارة إلى مناطق واسعة من الدماغ؛ أما الدم فيرسلها إلى كل خلية في الجسم. يدور الدم في الجسم في نحو دقيقة، فتصل الرسالة إلى الجميع خلال عشرات الثواني وتبقى دقائق وساعات حتى تتفكك المادة. وليس لهذا البث عنوان على الإطلاق. وكما في القضية~\ref{prop:receptor}، يختار المتلقّي معنى الرسالة: لا تستجيب إلا الخلايا التي تملك مستقبلًا لهذه المادة.

\section{المسرِّع والمكبح: سلكان إلى الأعضاء}

خير مثال القلب. في القلب ناظمة إيقاع خاصة به، كما في عصبون \nt{SERO} في الفصل~\ref{sec:serocomputer}: فهي تضع الإيقاع بنفسها، ولو قُطعت كل الأعصاب لنبض القلب نحو مئة مرة في الدقيقة. الدماغ لا يولّد الإيقاع، بل يضبطه فقط، وذلك عبر سلكين متعاكسي الإشارة (الشكل~\ref{fig:gasbrake}). في أحدهما يسري \nt{NORA}: وهو \emph{المسرِّع}، يرفع الإيقاع. وفي الآخر \nt{ACET}: وهو \emph{المكبح}، يخفضه. تقريبًا،
\begin{empheq}[box=\keybox]{equation}
f \;\approx\; f_0 \;+\; a_S\,S \;-\; a_P\,P ,
\label{eq:gasbrake}
\end{empheq}
حيث $f_0\approx100$ ضربة في الدقيقة هو الإيقاع الذاتي لناظمة الإيقاع، و$S$ و$P$ نشاط المسرِّع والمكبح. في الراحة يكون المكبح مضغوطًا، فينبض القلب عادةً نحو ستين إلى سبعين مرة في الدقيقة؛ وعند الجهد يُرفع المكبح ويُضغط المسرِّع.

\begin{figure}[t]
\centering
\begin{tikzpicture}[>=Stealth,
  blk/.style={draw,very thick,rounded corners=2pt,align=center,font=\small},
  pace/.style={circle,draw,very thick,double,double distance=1.2pt,minimum size=1.8cm,inner sep=0pt,font=\scriptsize,fill=red!8,align=center},
  inh/.style={-{Bar[width=7pt]},thick,red!70!black}]
\node[blk,fill=blue!5,text width=1.6cm] (B) at (0,0) {أمر\\الدماغ};
\node[pace] (H) at (6.2,0) {ناظمة\\الإيقاع};
\draw[->,very thick,orange!85!black] (B.north east) to[bend left=20] node[above,font=\scriptsize,black,align=center] {\nt{NORA}: مسرِّع، ثوانٍ} (H.north west);
\draw[inh,very thick,blue!60!black] (B.south east) to[bend right=20] node[below,font=\scriptsize,black,align=center] {\nt{ACET}: مكبح، أجزاء من الثانية} (H.south west);
\draw[->,very thick] (H) -- (8.4,0) node[right,font=\small] {التردد $f$};
\node[blk,fill=orange!10,text width=1.6cm] (A) at (0,-2.6) {\nt{ADRE}\\إلى الدم};
\draw[->,thick,orange!85!black] (B) -- (A);
\foreach \y/\lab in {-2.0/{القلب},-2.6/{الأوعية},-3.2/{الكبد، العضلات}}{
  \draw[->,thick,densely dashed,orange!85!black] (A.east) -- (4.3,\y) node[right,font=\scriptsize,black] {\lab};}
\end{tikzpicture}
\caption{سلكان متعاكسا الإشارة إلى ناظمة إيقاع القلب: المسرِّع \nt{NORA} بطيء، والمكبح \nt{ACET} سريع. في الأسفل المسرِّع نفسه عبر ناقل البث: تطلق خلايا \nt{ADRE} الأدرينالين في الدم، فتتلقى الإشارة كل الأعضاء التي تملك المستقبل المناسب (الأسهم المتقطعة).}
\label{fig:gasbrake}
\end{figure}

هذه هي الشفرة ثنائية السلك في الفصل~\ref{sec:glutcomputer} وزوج العضلات في الفصل~\ref{sec:muscles}، غير أن السلكين الآن لا يتحكمان في مفصل، بل في سرعة ناظمة الإيقاع. وللسلكين سرعتان مختلفتان. كلاهما يعمل مرشحَ تمرير منخفض، لكن المكبح يمرّر التغيرات أسرع من المسرِّع بعدة أضعاف~\cite{Berger1989}: طريق \nt{ACET} قصير، إذ يفتح البروتين المرتبط بالمستقبل قناة البوتاسيوم بنفسه، دون مرسال ثانٍ داخل الخلية~\cite{Logothetis1987gbg}، أما \nt{NORA} فيعمل عبر سلسلة من التفاعلات داخل الخلية. يعرف المهندس هنا حيلة المشغّلَين المختلفَي السرعة: السريع يتولى الضبط الدقيق الفوري، والبطيء يتولى التغيرات الكبيرة الطويلة.

وهنا يجد \nt{ADRE} دوره، وقد كُتب عنه في الجدول~\ref{tab:types} أن دوره في الدماغ غير واضح. أما خارج الدماغ فدوره واضح. بعض خلايا طريق المسرِّع لا ترسل سلكًا إلى عضو، بل تطلق الأدرينالين في الدم مباشرة. إنه المسرِّع نفسه، لكن عبر ناقل البث: يصل خلال عشرات الثواني إلى القلب والأوعية والكبد والعضلات معًا ويبقى دقائق. المسرِّع المعنوَن \nt{NORA} يضبط عضوًا واحدًا؛ والمسرِّع المبثوث \nt{ADRE} ينقل الجسم كله إلى وضع الركض.

\section{سلسلة مراحل بتغذية راجعة}

كثير من الهرمونات لا تطلقها غدة واحدة، بل سلسلة من المراحل. تطلق مجموعة صغيرة من العصبونات في قاعدة الدماغ الإشارة الأولى في الدم؛ فتدفع غدةً وسيطة إلى إطلاق إشارة ثانية؛ وهذه تدفع الغدة الأخيرة إلى إطلاق الهرمون الذي يؤثر في الجسم. ثم يعود الهرمون الأخير إلى المراحل الأولى فيثبطها. فينتج مضخم من ثلاث مراحل تحيط به تغذية راجعة سالبة، أي منظِّم مستوى، كما في منظومة \nt{SERO} في الصيغة~\eqref{eq:feedback}.

لنصف كل مرحلة بدارة RC ثابتها الزمني $\tau$ وكسبها $g$، والتغذية الراجعة بالمعامل $k$:
\begin{equation}
\tau\,\frac{\dd x_1}{\dd t} = -x_1 + g\bigl[u - k\,x_3\bigr]_+ ,\qquad
\tau\,\frac{\dd x_2}{\dd t} = -x_2 + g\,x_1 ,\qquad
\tau\,\frac{\dd x_3}{\dd t} = -x_3 + g\,x_2 ,
\label{eq:cascade}
\end{equation}
حيث $u$ أمر الدماغ، و$x_3$ مستوى الهرمون الأخير. كسب الحلقة $K=g^3k$. في التوازن $x_3=g^3u/(1+K)$، وكما في أي منظِّم بتغذية راجعة، تجعل الحلقةُ الكبيرة المستوى غير حساس لتفاوت كسب المراحل. لكن كل مرحلة تزيح الطور: عند التردد $\omega=\sqrt3/\tau$ تعطي المراحل الثلاث المتماثلة معًا إزاحة نصف دورة وتوهينًا بثمانية أضعاف. ومن هنا النتيجة الكلاسيكية في نظرية التحكم:
\begin{empheq}[box=\keybox]{equation}
K \;<\; 8 ,
\label{eq:cascadestable}
\end{empheq}
وإلا بدأ المنظِّم يتذبذب بدور يقارب $2\pi\tau/\sqrt3\approx3.6\,\tau$.

\begin{keyblock}
\begin{proposition}[الدقة مقابل الاستقرار]
\label{prop:cascade}
في سلسلة من ثلاث مراحل بتغذية راجعة تناسبية، خطأ المستوى يساوي $1/(1+K)$، ولا تستقر السلسلة إلا عندما $K<8$. لذلك لا يستطيع هذا المنظِّم أن يثبّت المستوى بدقة أفضل من نحو عشرة في المئة؛ ومحاولة رفع الكسب تحوّله إلى مولّد تذبذبات.
\end{proposition}
\end{keyblock}

تحققنا من ذلك على النموذج~\eqref{eq:cascade} (الشكل~\ref{fig:cascade}). عند $K=4$ يتذبذب المستوى قليلًا بعد التشغيل ثم يستقر. وعند $K=7$ يستقر أيضًا، لكن ببطء. وعند $K=12$ لا تخمد التذبذبات ولا تنمو بلا حد: فالمرحلة لا تستطيع إخراج إشارة سالبة، فتدخل السلسلة في نبضات منتظمة بين $0.83$ و$1.24$ من المستوى المتوسط بدور $3.7\text{--}3.9\,\tau$.

\begin{figure}[t]
\centering
\begin{tikzpicture}[>=Stealth,x=0.3cm,y=1.2cm]
\draw[->,gray!70!black] (0,0) -- (41.5,0) node[right,font=\small,black] {$t/\tau$};
\draw[->,gray!70!black] (0,0) -- (0,2.8) node[left,font=\small,black] {$x_3/x_3^*$};
\foreach \t in {0,10,20,30,40}{\draw[gray!60!black] (\t,-0.05) -- (\t,0.05); \node[below,font=\scriptsize] at (\t,0) {\t};}
\foreach \f in {1,2}{\draw[gray!60!black] (-0.3,\f) -- (0.3,\f); \node[left,font=\scriptsize] at (0,\f) {\f};}
\draw[densely dashed,gray!60!black] (0,1) -- (40,1);
\draw[thick,blue!60!black] plot file {figures/cascade_K4.dat};
\draw[thick,red!70!black] plot file {figures/cascade_K12.dat};
\node[font=\scriptsize,blue!60!black,anchor=west] at (16,0.55) {$K=4$: يستقر};
\node[font=\scriptsize,red!70!black,anchor=west] at (16,1.75) {$K=12$: ينبض};
\end{tikzpicture}
\caption{سلسلة من ثلاث مراحل بتغذية راجعة وفق~\eqref{eq:cascade}؛ مستوى الهرمون الأخير بوصفه كسرًا من قيمة التوازن $x_3^*$. حين يكون كسب الحلقة أقل من ثمانية يستقر المستوى؛ وحين يزيد عليها تولّد السلسلة بنفسها نبضات منتظمة.}
\label{fig:cascade}
\end{figure}

والهرمونات الحقيقية في هذه السلاسل تخرج فعلًا على شكل نبضات: فمستوى هرمون الإجهاد مثلًا يتذبذب بدور يقارب الساعة. وبحسب إحدى الفرضيات، تولّد هذه النبضاتِ التغذيةُ الراجعة نفسها بين المراحل مع تأخرها، ولا حاجة إلى مولّد إيقاع مستقل~\cite{Walker2010}. وهذا هو سبب التذبذب نفسه في حلقة \foreignlanguage{english}{\nt{GLUT}--\nt{GABA}} في القضية~\ref{prop:ei} وفي حلقة شدّ العضلة في الشرط~\eqref{eq:delaystable}: تغذية راجعة سالبة متأخرة.

\section{الشفرة النبضية للهرمونات}

النبضات ليست أثرًا جانبيًا فحسب، بل قد تكون شفرة. العصبونات التي تتحكم في التكاثر تطلق إشارتها في الدم دفعاتٍ قصيرة نحو مرة كل ساعة. فإذا أُعطيت الغدة المتلقية الإشارةَ نفسها باستمرار لا دفعاتٍ، فإنها لا تعمل بكامل طاقتها كما مع النبضات، بل تصمت؛ وإذا عادت الإشارة دفعاتٍ مرة كل ساعة، عاد العمل~\cite{Belchetz1978}. إذن المتلقّي لا يقرأ المستوى، بل \emph{تردد} الدفعات.

لا لغز في ذلك عند المهندس. المستقبل الذي يتعب ويكفّ عن الاستجابة تحت تأثير المادة الطويل هو مرشح تمرير عالٍ، مثل المشبك المستنفَد~\eqref{eq:depression} في الفصل~\ref{sec:code}. فهو يخمد الإشارة المستمرة ويمرّر التغيرات. ولا يمكن إبلاغ مثل هذا المتلقّي شيئًا إلا بالنبضات، فتنتقل المعلومة من المستوى إلى تردد الدفعات وشكلها. ثم إن الترددات المختلفة للدفعات تؤثر تأثيرات مختلفة في خلايا مختلفة من الغدة نفسها، فيمكن نقل أكثر من كمية واحدة عبر خط كيميائي واحد، كما سرت مركّبتان سريعة وبطيئة عبر سلك \nt{DOPA} واحد في الفصل~\ref{sec:dofaalt}.

\section{الإيقاع اليومي: مولّد بطيء للأوضاع}

أبطأ إيقاعات الجسم هو الإيقاع اليومي. تولّده تغذية راجعة سالبة متأخرة داخل الخلايا: تنتج الجينات بروتينات تثبّط إنتاجها هي نفسها بتأخر عدة ساعات~\cite{TakahashiJS2017}. مرة أخرى تغذية راجعة سالبة متأخرة، للمرة الرابعة في هذا الكتاب، ومرة أخرى إيقاع، دوره هذه المرة نحو يوم. المولّد الرئيسي مجموعة صغيرة من عشرات آلاف العصبونات في قاعدة الدماغ، يزامن إيقاعُها بقية خلايا الجسم.

الدور الذاتي لهذا المولّد عند الإنسان ليس يومًا بالضبط، بل $24.18$ ساعة في المتوسط~\cite{Czeisler1999}. ويضبطه كل يوم الضوءُ الآتي من مستشعرات العين (الفصل~\ref{sec:sensors}). وهذا عند مهندس الإلكترونيات \emph{حلقة مقفلة الطور}: مولّد بتردده الذاتي $\omega_0$ يتزامن مع إشارة خارجية ترددها $\omega$، وفرق الطور $\varphi$ يخضع للمعادلة
\begin{empheq}[box=\keybox]{equation}
\frac{\dd\varphi}{\dd t} \;=\; (\omega-\omega_0) \;-\; K_\varphi\,\sin\varphi .
\label{eq:pll}
\end{empheq}
يبقى التزامن قائمًا إذا كان $|\omega-\omega_0|<K_\varphi$. والمولّد الذي دوره $24.18$ ساعة يحتاج إلى أن ينزاح كل يوم نحو إحدى عشرة دقيقة؛ ويكفي ضوء الصباح عادةً لهذه الإزاحة. أما في الليل القطبي أو في كهف بلا ضوء فلا تزامن، وينجرف الإيقاع إلى دوره الذاتي.

المولّد اليومي ليس مولّد نبضات الساعة في الحاسوب. فهو لا يعدّ خطوات الحساب ولا يزامن نبضات العصبونات. إنه يبدّل \emph{الأوضاع}: النوم واليقظة، ومستويات الهرمونات، وحرارة الجسم، والاستعداد للطعام. إنه سجل بث بطيء من جنس سجلات الفصول~\babelsublr{\ref{sec:serocomputer}--\ref{sec:acet}}، غير أن قيمته تتغير وفق جدول مضبوط على الشمس.

\begin{keyblock}
\begin{proposition}[المخرج الكيميائي]
\label{prop:chemoutput}
المخرج البطيء للآلة مبني من القطع نفسها التي بُني منها المخرج السريع. تتلقى الأعضاء سلكين متعاكسي الإشارة مختلفي السرعة، المسرِّع \nt{NORA} والمكبح \nt{ACET}، ويتلقى الجسم كله معًا إشارات بث عبر الدم، ومنها الأدرينالين \nt{ADRE}. وتثبّت مستوياتِ الهرمونات سلاسلُ مراحل بتغذية راجعة سالبة، دقتها محدودة بالاستقرار؛ وبعض الهرمونات مرمَّز بتردد الدفعات؛ والإيقاع اليومي مولّد بطيء بحلقة مقفلة الطور يضبطها الضوء. بنية الطرق والتجارب على الدفعات وعلى الدور مُقاسة؛ أما تفسير نبضات السلسلة بالتغذية الراجعة نفسها ففرضية.
\end{proposition}
\end{keyblock}

\section{الخلاصة}

\begin{table}[t]
\centering
\footnotesize
\setlength{\tabcolsep}{4pt}
\renewcommand{\arraystretch}{1.15}
\begin{tabular}{>{\raggedright}p{3.1cm}>{\raggedright}p{4.8cm}>{\raggedright\arraybackslash}p{4.9cm}}
\toprule
ماذا يفعل المخرج الكيميائي & كيف & ما المعروف \\
\midrule
يضبط الأعضاء & المسرِّع \nt{NORA} والمكبح \nt{ACET}~\eqref{eq:gasbrake} & مُقاس؛ المكبح أسرع من المسرِّع \\
ينقل الجسم كله إلى وضع الركض & الأدرينالين \nt{ADRE} في الدم & مُقاس \\
يثبّت مستويات الهرمونات & سلسلة مراحل بتغذية راجعة، $K<8$~\eqref{eq:cascadestable} & السلاسل والتغذية الراجعة مُقاسة؛ نشوء النبضات من الحلقة نفسها فرضية \\
ينقل بتردد الدفعات & المستقبل المُتعَب مرشحَ تمرير عالٍ & الاعتماد على الدفعات مُقاس \\
يبدّل الأوضاع على مدار اليوم & مولّد بحلقة مقفلة الطور يضبطها الضوء~\eqref{eq:pll} & الدور والضبط بالضوء مُقاسان \\
\bottomrule
\end{tabular}
\caption{كيف تتحكم الآلة في كيمياء الجسم.}
\label{tab:chemoutput}
\end{table}

للآلة مخرجان (الجدول~\ref{tab:chemoutput}). السريع هو العضلات: أجزاء من الألف من الثانية، معنوَن، إلى كل مفصل. والبطيء هو كيمياء الجسم: ثوانٍ ودقائق وأيام، عبر سلكين إلى الأعضاء وعبر الدم إلى الجميع معًا. والقطع في المخرجين واحدة: سلكان متعاكسا الإشارة، وناظمات إيقاع، وتغذية راجعة وتأخرها، وهي نفسها القطع التي جُمعت منها الآلة من الداخل.

\section{تمارين}

\begin{exercise}[\lvlA]
\label{ex:16:1}
\emph{الدم مرشحًا.} يُطرح هرمون من الدم بعمر نصف $t_{1/2}=10$ دقائق ويُطلق دفعاتٍ قصيرة كل $T=60$ دقيقة. كم ضعفًا تزيد القيمة العظمى للتركيز على الصغرى، وكم ضعفًا تُوهَّن التوافقية الأساسية للدفعات؟ عند أي $t_{1/2}$ تزيد العظمى على الصغرى ضعفين فقط؟
\end{exercise}

\begin{exercise}[\lvlA]
\label{ex:16:2}
\emph{التزامن اليومي.} الدور الذاتي للمولّد~\eqref{eq:pll} يساوي $24.18$ ساعة، والضوء يزيح الطور ساعة في اليوم على الأكثر: $K_\varphi=2\pi/24$ راديان/يوم. أوجد فرق الطور المستقر $\varphi^*$ بالدقائق وزمن الاسترخاء. بعد كم يومًا يصبح عدم التطابق أقل من ساعة إثر قفزة في الإشارة الخارجية مقدارها $\pm6$ ساعات؟
\end{exercise}

\begin{exercise}[\lvlB]
\label{ex:16:3}
\emph{الدقة مقابل الاستقرار.} مُدّدت السلسلة الخطية~\eqref{eq:cascade} إلى $n$ مرحلة متماثلة. أوجد الكسب الحرج $K_c(n)$ وأصغر خطأ ممكن $1/(1+K_c)$ عند $n=3$ و$n=6$. إلامَ يؤول هذا الخطأ عندما $n\to\infty$؟
\end{exercise}

\begin{exercise}[\lvlB]
\label{ex:16:4}
\emph{المسرِّع والمكبح.} في~\eqref{eq:gasbrake} مساهمتا المسرِّع والمكبح مخرجا دارتي RC بـ$\tau_S=2$ ثانية و$\tau_P=0.4$ ثانية، والأمران لا يتجاوزان $80$ و$60$ ضربة في الدقيقة. الإيقاع $f_0=100$؛ في الراحة المكبح $35$ والمسرِّع $0$ و$f=65$. في كم من الزمن يمكن بلوغ $f=120$؟ وماذا لو لم يُرفع المكبح وعمل المسرِّع وحده؟
\end{exercise}

\begin{exercise}[\lvlB]
\label{ex:16:5}
\emph{نمذجة: سلسلة بتأخير.} أضف تأخيرًا $d$ إلى التغذية الراجعة في الدالة \texttt{cascade} من \texttt{models/\allowbreak chem\_models.py} (انسخها ولا تشغّل الملف): ترى المرحلة الأولى $x_3(t-d)$. أوجد $K_c$ والدور عند الحد عند $d/\tau=0$ و$0.5$ و$1$ و$2$، وتحقق بالمحاكاة عند $0.9K_c$ و$1.1K_c$.
\end{exercise}

\begin{exercise}[\lvlB]
\label{ex:16:6}
\emph{متلقٍّ يقرأ الدفعات.} المستقبل يتعب: $y=[c-a]_+$، $\tau_a\,\dd a/\dd t=c-a$، $\tau_a=30$ دقيقة. يصل الهرمون دفعاتٍ مدة كل منها $6$ دقائق بدور $T$ مع الجرعة المتوسطة نفسها $\bar c$. أوجد متوسط الاستجابة عند $T=15$ و$60$ و$120$ دقيقة وعندما $T\to\infty$. وكم يكون عند تركيز ثابت $\bar c$؟
\end{exercise}

\begin{exercise}[\lvlC]
\label{ex:16:7}
\emph{مولّد أم تغذية راجعة؟} اقترح تجربة تميّز النبضات التي تولّدها التغذية الراجعة المتأخرة في السلسلة~\eqref{eq:cascade} من ناظمة إيقاع مستقلة. هل يمكن التمييز بين الفرضيتين بإزاحة الدور إذا لم يمكن تغيير $K$ إلا مرة ونصفًا، وكان الدور يُقاس بدقة $5\%$؟
\end{exercise}

\chapter{تصحيح أخطاء الآلة}
\chaptermark{تصحيح أخطاء الآلة}
\label{sec:debugging}
\begin{chaptergoals}
\item لماذا يكفي مسبار يسمع ألف عصبون من أصل $8.6\cdot10^{10}$ لفكّ شفرة الحركة؛
\item لماذا يجب أن يستخرج الجهاز المزروع النبضات على الرقاقة ولا يرسل إلا الأحداث، لا التدفق الخام؛
\item كيف يقرأ مفكِّك الشفرة ترددات المجموعات، ولماذا لا تستخرج الواجهات إلا بضعة بتات في الثانية؛
\item لماذا الكتابة في الدماغ أصعب من القراءة، وما الذي لا يراه المسبار.
\end{chaptergoals}

\section{كيف تُصحَّح أخطاء آلة بلا مخطط}

المهندس الذي يُعطى لوحة تعمل بلا مخطط يصحح أخطاءها بأربع حيل. يضع \emph{مسبارًا} وينظر ما الذي يسري في السلك. \emph{ويحقن إشارة اختبار} وينظر ما الذي تغيّر. \emph{ويفصل} جزءًا من الدارة وينظر ما الذي توقف عن العمل. \emph{ويقرأ سجلات التحكم} ليعرف في أي وضع تعمل اللوحة. هذا الكتاب كله محاولة لقراءة مخطط الدماغ، وفي هذا الفصل سننظر في أيّ الحيل الأربع يُحسن المهندسون استعمالها اليوم، وما الذي تقوله لهم الفصول السابقة عمّا ينبغي أن يتوقعوه.

أبرز أمثلة هذا التصحيح اليوم واجهات الدماغ والحاسوب: أجهزة تقرأ نبضات العصبونات وتحوّلها إلى أوامر لآلات خارجية، أو تحقن في الدماغ إشارات من الخارج بالعكس. ولها يُخصَّص معظم الفصل، وأولًا لما تفعله شركة \foreignlanguage{english}{Neuralink}.

\section{المسبار: ماذا يسمع القطب}

مسبار الدماغ قطب رفيع مغروس في النسيج. يسمع نبضات العصبونات الواقعة في نصف قطر من رتبة مئة ميكرومتر، وهي عادةً من عصبون واحد إلى بضعة عصبونات. النبضة على القطب ومضة من عشرات إلى مئات الميكروفولتات تدوم نحو جزء من الألف من الثانية، ومن شكلها يمكن تمييز العصبونات المتجاورة بعضها من بعض. وأوضح ما يُسمع عصبونات \nt{GLUT} الكبيرة: فهي الأكثرية في القشرة (الجدول~\ref{tab:types})، وتياراتها أقوى.

والصعوبة الرئيسية ظاهرة على الفور. المسبار الجيد اليوم مئات أو آلاف الأقطاب، أما عدد العصبونات في الدماغ فـ$8.6\cdot10^{10}$. نحن نتنصّت على نحو جزء من مئة مليون من الآلة. فلماذا يمكن أن نفهم شيئًا أصلًا من عيّنة ضئيلة كهذه؟ أجاب الفصل~\ref{sec:code}. العصبونات المتجاورة تتشارك ضجيجها، فيصطدم المتوسط على المجموعة بحد (القضية~\ref{prop:pooling}): مئة عصبون تحمل تقريبًا ما تحمله ألف. ثم إن المسألة التي تحلها القشرة الحركية منخفضة الأبعاد: للذراع مفاصل قليلة (الفصل~\ref{sec:muscles})، ونشاط مئات العصبونات التي تتحكم فيها ينحصر في فضاء من عشرة أو عشرين متغيرًا مشتركًا فقط~\cite{Gallego2017}. والمسبار الذي يسمع مئة عصبون يسمع تقريبًا كل هذه المتغيرات. ولو كان الدماغ يخزن في كل عصبون رسالة مستقلة، لما أجدى هذا التنصّت شيئًا.

\section{تدفق البيانات: الضغط على المسبار نفسه}

يُخرج المسبار تدفقًا هائلًا. لرؤية شكل النبضة يجب رقمنة كل قطب نحو $20$ ألف مرة في الثانية بدقة نحو عشرة بتات. وألف قطب تعطي
\begin{empheq}[box=\keybox]{equation}
\begin{aligned}
R_{\text{raw}} \;&=\; N\,f_s\,b \;\approx\; 10^3\cdot2\cdot10^4\cdot10 \;\approx\; 2\cdot10^8\ \text{bit/s},\\
R_{\text{spk}} \;&\approx\; N\,r\,\log_2\frac{e}{r\,\Delta t} \;\approx\; 8\cdot10^4\ \text{bit/s}.
\end{aligned}
\label{eq:probedata}
\end{empheq}
التقدير الثاني هو سعة تدفق النبضات~\eqref{eq:spikeentropy} عند $r=10$ نبضات في الثانية ودقة $\Delta t=1\ms$: كل ما فيه فعلًا يشغل حيزًا أصغر بألفين وخمسمئة مرة. ولجهاز مزروع هذا فرق حاسم: لا يمكن إرسال التدفق الخام لاسلكيًا عبر الجلد بالطاقة التي يُسمح بها داخل الرأس دون تسخين النسيج. لذلك يجب على الجهاز المزروع أن يستخرج النبضات على الرقاقة نفسها بجوار الأقطاب، وألا يرسل إلى الخارج إلا الأحداث: رقم القناة ولحظة النبضة. في الوصف الأول لمنظومة \foreignlanguage{english}{Neuralink} لم يكن الأمر قد بلغ ذلك بعد: الرقاقات تضخّم الإشارات وترقمنها، والتدفق الخام كله يخرج عبر كابل، وبرنامج في الخارج يستخرج النبضات؛ أما الضغط على الرقاقة والاتصال اللاسلكي فقد ذُكرا هناك خطةً~\cite{Musk2019}.

هذه حيلة مألوفة. العين تضغط الإشارة مئة مرة قبل أن ترسلها عبر كابل طويل (الفصل~\ref{sec:sensors})؛ وكاميرا الأحداث لا ترسل إطارات بل تغيرات. والمسبار، كي يتسع له السلك، مضطر إلى أن يفعل ما يفعله الدماغ نفسه: أن يتكلم بالنبضات وألا يتكلم إلا عن الجديد.

\section{ماذا تفعل \foreignlanguage{english}{Neuralink}}

جهاز \foreignlanguage{english}{Neuralink} مسبار مصمم لهذه القيود~\cite{Musk2019}. بدل مصفوفة صلبة من الإبر، خيوط مرنة كثيرة أرفع من الشعرة، على طول كل منها أقطاب: في الوصف الأول للمنظومة حتى $3072$ قطبًا على $96$ خيطًا، $32$ على كل خيط، وفي الجهاز المزروع في البشر، بحسب ما أعلنته الشركة، نحو ألف قناة على $64$ خيطًا. الخيوط المرنة تؤذي النسيج أقل، وتتحمل أفضل أن الدماغ داخل الجمجمة يتحرك قليلًا طوال الوقت. يغرس الخيوطَ روبوتٌ يتجنب الأوعية. وفي الوصف الأول، كما سبق، كان التدفق الخام~\eqref{eq:probedata} يخرج عبر كابل. أما الجهاز المخصص للبشر، بحسب ما أعلنته الشركة، فمبني بطريقة أخرى: الغلاف مغطى بالجلد كليًا، والنبضات تُستخرج في الداخل، والبيانات تخرج لاسلكيًا، والطاقة تصل بلا أسلاك.

مهمة الجهاز الأول هي القراءة. توضع الخيوط في جزء القشرة الذي يخطط حركات الذراع، ويحوّل مفكِّكُ الشفرة النبضاتِ إلى حركة مؤشر على الشاشة. فيتحكم إنسان مشلول اليدين في الحاسوب بتخيّل الحركة. الفكرة نفسها ليست جديدة: الواجهات القائمة على مصفوفات صلبة من مئة قطب سمحت منذ زمن بالتحكم في المؤشر~\cite{Hochberg2006}، وبكتابة النص بتخيّل حركات اليد في الكتابة~\cite{Willett2021}، بل بتحويل محاولات الكلام إلى نص على الشاشة بسرعة نحو ستين كلمة في الدقيقة~\cite{Willett2023}. الجديد عند \foreignlanguage{english}{Neuralink} هو الهندسة: عدد القنوات، واللاسلكية، والغلاف المغلق، والغرس الروبوتي، أي محاولة صنع مسبار يمكن حمله سنوات خارج المختبر. وبحسب علمنا، نُشرت نتائج التجارب على البشر في الغالب من الشركة نفسها لا في مقالات محكَّمة؛ وقد أعلنت الشركة، من بين أمور أخرى، أن بعض الخيوط انزاح بعد الغرس فنقص عدد القنوات العاملة. وهذا إزعاج عادي في تصحيح الأخطاء: المسبار الذي يتحرك بالنسبة إلى اللوحة يسمع أسلاكًا أخرى.

الاتجاه الثاني للشركة هو الكتابة: جهاز للبصر يحقن إشارات في القشرة البصرية لإنسان فقد بصره. وعنه نتحدث أدناه، في قسم الكتابة.

\section{القراءة: مفكِّك الشفرة متلقّيًا}

مفكِّك شفرة الواجهة متلقٍّ بمعنى القضية~\ref{prop:reader}: فهو يقرر بنفسه أي جزء من الرسالة يقرأ. عمليًا، تقرأ كل الواجهات تقريبًا \emph{ترددات المجموعات}: تعدّ نبضات كل قناة في نوافذ من بضع عشرات من أجزاء الألف من الثانية، وتجمعها بأوزان مختارة بحيث تنتج الحركة المقصودة. هذا ترميز التردد للمجموعة من القسم~\ref{sec:rate}، ولم يُختر مصادفة: مع بضع نبضات في النافذة لا يتحدد تردد عصبون واحد، أما المجموع على مئة فيعمل.

كم من المعلومات يمكن استخراجه؟ الكتابة ”باليد المتخيَّلة“ سارت بسرعة نحو تسعين حرفًا في الدقيقة~\cite{Willett2021}، أي حرف ونصف في الثانية: وحتى بخمسة بتات للحرف يكون ذلك أقل من ثمانية بتات في الثانية. والتحكم في المؤشر، بحسب ما أعلنته \foreignlanguage{english}{Neuralink}، يعطي نحو عشرة بتات في الثانية. أما الحد الأعلى~\eqref{eq:spikeentropy} لعصبون \emph{واحد} فعشرات البتات في الثانية، ولمئة عصبون آلاف. الواجهة تستخرج من العصبونات المتنصَّت عليها جزءًا صغيرًا مما يمكن أن تحمله. وعنق الزجاجة ليس السلك، بل عدد المتغيرات المستقلة التي تحملها المجموعة (القضية~\ref{prop:pooling})، ومدى فهم المفكِّك للشفرة التي، كما رأينا في الفصل~\ref{sec:code}، لم تُفكّ كاملة بعد.

وثمة خاصية ثانية مألوفة من الفصل~\ref{sec:motorlearning}. المفكِّك والدماغ يتعلم كل منهما من الآخر. يرى الدماغ إلى أين ذهب المؤشر، فيتلقى خطأً ويعدّل أوامره، وهو التعلم الحركي نفسه عبر سلك الخطأ. والمهندسون يعيدون من حين إلى آخر اختيار أوزان المفكِّك، لأن العصبونات تحت الخيوط تتغير والوصلات في الشبكة يُعاد تعلّمها. فينتج متعلمان يتكيف كل منهما مع الآخر؛ ولكي لا يفلت هذا الزوج من الضبط، يحسن الفصل بين سرعتي تعلمهما: ليتعلم أحدهما بسرعة والآخر ببطء.

\begin{keyblock}
\begin{proposition}[لماذا يعمل مسبار على ألف عصبون]
\label{prop:probe}
الواجهة التي تتنصت على جزء ضئيل من العصبونات تستطيع التحكم في الحركة لأن نشاط المجموعة منخفض الأبعاد وضجيجها متشارك: بضع مئات من العصبونات تحمل تقريبًا كل المتغيرات المشتركة. وللسبب نفسه لا تستخرج الواجهة إلا بضعة بتات في الثانية، بقدر المتغيرات المستقلة في المسألة، لا بقدر ما يتسع له تدفق النبضات. عمل الواجهات مُقاس؛ أما مقدار تحسّن المفكِّك حين تُفهم الشفرة فسؤال مفتوح.
\end{proposition}
\end{keyblock}

\section{الكتابة: أصعب من القراءة}

حقن إشارة في الآلة أصعب من القراءة منها. القطب الذي يمرّ عبره تيار لا يثير عصبونًا واحدًا، بل كل العصبونات والأسلاك حوله، دون تمييز للنوع أو الإشارة. ولا يمكن به كتابة شفرة يهم فيها \emph{أيّ} عصبون أطلق و\emph{متى} (الفصل~\ref{sec:code}): لا يمكن إلا تحديد \emph{أين} و\emph{بأي شدة} تحديدًا خشنًا.

وهذا يكفي للبصر جزئيًا. يرى الإنسان التيار عبر قطب في القشرة البصرية نقطةً مضيئة في موضع محدد من مجال الرؤية؛ وحين أُضيئت النقاط معًا، حتى ست عشرة نقطة في آن، تعرّف إنسان فقد بصره على بعض الحروف وميّز حدود الأشياء~\cite{Fernandez2021}. هذه شفرة بالموضع، ويمكن كتابتها. لكن تذكّر الفصل~\ref{sec:sensors}: العين لا ترسل إلى الدماغ بكسلات، بل وصفًا مضغوطًا: الحواف، والتغيرات، والإضاءة والإعتام، على أسلاك منفصلة. كتابة ”نقاط الضوء“ في القشرة كتابةُ بكسلات في آلة تنتظر عند مدخلها شفرة مختلفة تمامًا. لذلك ما زال البصر الاصطناعي أخشن مما يمكن توقعه من عدد الأقطاب. وثمة إشارات اختبار لا تحتاج إلى قطب: الخداع البصري يحقن عبر العين مدخلًا غير مألوف ويُخرج الآلة من وضعها المعتاد، وكل خطأ يدل على آليته (القسم~\ref{sec:illusions}).

\section{ما لا يراه المسبار}

القطب يسمع ناقل العناوين: نبضات عصبونات \nt{GLUT}، وبدرجة أقل عصبونات \nt{GABA} الصغيرة. لكننا رأينا في الفصول~\babelsublr{\ref{sec:serocomputer}--\ref{sec:acet}} أن وضع عمل الآلة كلها تحدده سجلات البث: تراكيز \nt{SERO} و\nt{DOPA} و\nt{NORA} و\nt{ACET}. وهذه لا يسمعها القطب: العصبونات المصدر قليلة إلى حد الضآلة، وإشارتها ليست نبضة بجوار المسبار بل تركيز في حجم. ومصحح الأخطاء الذي يرى ناقل البيانات ولا يرى سجلات التحكم سيبقى متعجبًا دائمًا: المدخل الواحد سيعطي استجابات مختلفة، لأن $g$ أو $a$ أو $\tau_\gamma$ تغيّر في مكان ما. يمكن قياس التراكيز بمستشعرات كيميائية داخل النسيج مباشرة~\cite{Bunin1998}، لكن هذه المستشعرات لا تُستعمل بعد في الواجهات. ويبقى خارج مجال رؤية المسبار تمامًا المخرجُ الثاني البطيء للآلة: الهرمونات في الدم (الفصل~\ref{sec:chemoutput})، وهي تقاس في عينات الدم لا بالقطب.

ولا يرى المسبار الأوزان، وفيها بالذات تُخزَّن الذاكرة كلها تقريبًا وكل ما تعلّمته الآلة. لا يمكن إلا تخمينها من الطريقة التي تتغير بها الاستجابات. ولا يرى الألم كما يشعر به الإنسان: رأينا في الفصل~\ref{sec:pain} أن شدة إشارة الإنذار تحددها الآلة نفسها، ببوابات في الحبل الشوكي ومنظّمات من الأعلى، فلا يمكن أن يُقرأ من المستشعر مقدار الألم.

\section{الخلاصة: التصحيح والأسئلة المفتوحة}

\begin{table}[t]
\centering
\footnotesize
\setlength{\tabcolsep}{4pt}
\renewcommand{\arraystretch}{1.15}
\begin{tabular}{>{\raggedright}p{2.2cm}>{\raggedright}p{3.6cm}>{\raggedright}p{3.4cm}>{\raggedright\arraybackslash}p{4.0cm}}
\toprule
حيلة التصحيح & ماذا تفعل الواجهة & ماذا يفسر الكتاب & ما المعروف \\
\midrule
المسبار & يسمع مئات إلى آلاف العصبونات & انخفاض الأبعاد والضجيج المتشارك (الفصل~\ref{sec:code}) & مُقاس \\
الضغط على المسبار & يرسل أحداثًا بدل الإشارة الخام & سعة تدفق النبضات~\eqref{eq:spikeentropy} & محلول هندسيًا \\
القراءة & من ترددات المجموعات إلى حركة ونص وكلام & المفكِّك متلقٍّ، ترميز التردد للمجموعة & يعمل؛ بضعة بتات في الثانية \\
التعلم المشترك & الدماغ والمفكِّك يتكيف كل منهما مع الآخر & سلك الخطأ (الفصل~\ref{sec:motorlearning}) & مُلاحَظ؛ قواعد التكيف موضوع بحث \\
الكتابة & التيار يضيء نقاطًا في مجال الرؤية & ترميز الموضع يمكن كتابته، وترميز التوقيت لا (الفصل~\ref{sec:sensors}) & النقاط مُقاسة؛ البصر المفيد ليس بعد \\
قراءة السجلات & لا تكاد تُجرى & قنوات البث (الفصول~\babelsublr{\ref{sec:serocomputer}--\ref{sec:acet}}) & المستشعرات الكيميائية موجودة في التجارب \\
\bottomrule
\end{tabular}
\caption{تصحيح أخطاء الآلة: ما تستطيعه واجهات الدماغ والحاسوب وما تقوله فصول الكتاب عنها.}
\label{tab:debugging}
\end{table}

يلخص الجدول~\ref{tab:debugging} الفصل. تستطيع واجهات الدماغ والحاسوب اليوم أن تضع مسبارًا على ناقل العناوين وأن تقرأ منه بضعة بتات في الثانية، وكون ذلك يعمل أصلًا يفسره انخفاض الأبعاد والضجيج المتشارك من الفصل~\ref{sec:code}. أما الكتابة في الآلة فلا تستطيعها إلا خشنةً، لأن شفرة مدخلها مضغوطة وموجهة إلى أسلاك منفصلة. وسجلات التحكم والأوزان، وهي ما يجعل الآلة قابلة للتعلم والضبط، لا يكاد المسبار يراها بعد.

كل سؤال مفتوح في الفصل~\ref{sec:synthesis} هو أيضًا مسألة لمصحح الأخطاء. أيّ شفرة تُقرأ، سيتضح حين تصير المسابير أكثف ويتعلم المفكِّكون استعمال توقيت النبضات لا تردداتها فقط. وكيف تتعلم شبكة متعددة الطبقات، يمكن فحصه بوضع مسابير في عدة طبقات معًا ومراقبة تغير استجاباتها أثناء التعلم. وما الذي تنقله قنوات البث، سيتضح حين تضاف المسابير الكيميائية إلى الكهربائية. هذا الكتاب محاولة لقراءة مخطط الآلة من سلوكها ومن القوانين التي يعرفها كل مهندس. ومصححات الأخطاء التي تُبنى اليوم ستتيح فحص هذا المخطط بالمسبار.

\section{تمارين}

\begin{exercise}[\lvlA]
\label{ex:17:1}
\emph{ميزانية القناة اللاسلكية.} الإرسال عبر الجلد يكلف $1\,\text{nJ}$ للبت، ولا يجوز للمرسل أن ينفق أكثر من $10\,\text{mW}$. ما القدرة اللازمة للتدفق الخام~\eqref{eq:probedata} عند $N=1000$ قطب، ولتدفق النبضات؟ وأي طاقة للبت كان التدفق الخام سيتطلب؟
\end{exercise}

\begin{exercise}[\lvlA]
\label{ex:17:2}
\emph{كم بتًا في مئة عصبون.} يجمع المفكِّك نبضات $N$ عصبونًا في نوافذ $50\ms$ عند $r=20$ نبضة/ثانية، والارتباط الزوجي للضجيج $c=0.1$، والإشارة تغيّر تردد المجموعة بـ$30\%$ (جذر متوسط المربعات). قدّر المعدل $\tfrac12\log_2(1+\text{SNR})$ في النافذة عند $N=10$ و$100$ و$1000$ و$N\to\infty$. وعند $c=0$ و$N=1000$؟
\end{exercise}

\begin{exercise}[\lvlB]
\label{ex:17:3}
\emph{سرعة الواجهة-لوحة المفاتيح.} تختار الواجهة مرة ونصفًا في الثانية حرفًا من $N=32$: المقصود باحتمال $P$، والأخطاء متساوية الاحتمال. كم بتًا في الثانية تنقل عند $P=0.99$ و$0.94$ و$0.8$؟ وكم حرفًا في الثانية يلزم عند $P=0.94$ لبلوغ $10$ بتات/ثانية؟
\end{exercise}

\begin{exercise}[\lvlB]
\label{ex:17:4}
\emph{نمذجة: مفكِّك المجموعة.} $N$ عصبونًا بـ$\lambda_i=[b+m\,\mathbf v\cdot\mathbf p_i]_+$، $b=20$، $m=15$ نبضة/ثانية، نوافذ $50\ms$، و$\mathbf v$ بانحراف معياري $0.5$ لكل مركّبة؛ يرى الجميع ضجيجًا مشتركًا، $\mathbf v+\boldsymbol\xi$، $\sigma_\xi=0.2$. حاكِ مفكِّكًا غير متحيز بمتجه المجموعة. عند أي $N$ يتساوى الضجيجان، وكم بتًا للمركّبة في النافذة يعطي $N\to\infty$؟
\end{exercise}

\begin{exercise}[\lvlB]
\label{ex:17:5}
\emph{متعلمان.} يُخرج الدماغ الأمر $m=b\,v^*$، والمفكِّك $v=d\,m$، $\langle v^{*2}\rangle=1$. كلاهما بعد كل محاولة يخطو خطوة تدرج على $\tfrac12(v-v^*)^2$ بمعدل $\eta$. حاكِ بدءًا من $b=1.2$ و$d=1$ عند $\eta=0.8$ و$1.1$ و$1.5$ و$2$. كل منهما وحده مستقر عند $\eta<2$؛ فعند أي $\eta$ يستقر الزوج؟
\end{exercise}

\begin{exercise}[\lvlB]
\label{ex:17:6}
\emph{تصميم: خط معالجة المسبار.} $1024$ قناة بـ$20$ ألف عيّنة في الثانية، والعصبونات $10$ نبضات/ثانية، والإرسال اللاسلكي $1\,\text{nJ}$ للبت. أي عتبة على الضجيج الأبيض للعيّنات تعطي $0.1\,\text{Hz}$ على الأكثر من النبضات الكاذبة لكل قناة؟ نرسل أعداد النبضات في $20\ms$ بـ$2$ بت: ما القدرة، وكم مرة هي أصغر مما تتطلبه العيّنات الخام بـ$10$ بتات؟
\end{exercise}

\begin{exercise}[\lvlC]
\label{ex:17:7}
\emph{حد سرعة الواجهة.} قدّر $R\le DW\log_2(1+\text{SNR})$ لـ$D=10$ متغيرات بعرض نطاق $W=5\,\text{Hz}$ عند $\text{SNR}\approx0.9$ (ضجيج يشبه الإشارة) و$90$ (ضجيج مستقل لـ$1000$ عصبون). المُقاس نحو $10$ بتات/ثانية. أي تجربة تميّز بين تفسيري الفجوة: ضجيج يشبه الإشارة، أم الجهل بالشفرة؟
\end{exercise}

\chapter{أنواع العصبونات بحسب وظيفتها الكهربائية}
\chaptermark{العصبونات أجهزةً}
\label{sec:eetypes}

صنّفنا العصبونات في الفصل~\ref{sec:neurontypes} بحسب الناقل العصبي الذي يطلقه مخرجها. أما مهندس الإلكترونيات فكان سيصنّفها بطريقة أخرى: بحسب ما يفعله العصبون بإشارته، أي بحسب الجهاز الذي يعمل عمله. الجهاز الرئيسي في الكتاب نعرفه منذ الفصل~\ref{sec:glutcomputer}: العصبون المكامل التسريبي~\eqref{eq:glutneuron}، أي دارة RC مع مقارِن. لكن في الدماغ أجهزة أخرى. جُمعت في الجدول~\ref{tab:eetypes} في أربع مجموعات، وقد مرّ كل منها تقريبًا في الكتاب.

\begin{center}
\footnotesize
\setlength{\tabcolsep}{3pt}
\renewcommand{\arraystretch}{1.15}
\begin{longtable}{>{\raggedright}p{3.1cm}>{\raggedright}p{3.3cm}>{\raggedright}p{4.7cm}>{\raggedright\arraybackslash}p{2.4cm}}
\caption{العصبونات أجهزةً: الاسم، والنظير في الإلكترونيات، وما يفعله الجهاز، وأين يرد في الكتاب.}\label{tab:eetypes}\\
\toprule
العصبون & الجهاز & ماذا يفعل & أين في الكتاب \\
\midrule
\endfirsthead
\toprule
العصبون & الجهاز & ماذا يفعل & أين في الكتاب \\
\midrule
\endhead
\bottomrule
\endfoot
\multicolumn{4}{l}{\emph{المرشحات والمكاملات}} \\
العصبون المكامل التسريبي & مكامل RC مع مقارِن & يجمع المداخل في نافذة $\tau$ ويطلق عند العتبة & الفصل~\ref{sec:glutcomputer}، \eqref{eq:glutneuron} \\
كاشف التزامن & بوابة AND بنافذة زمنية & لا يطلق إلا إذا وصلت المداخل معًا تقريبًا؛ $\tau$ صغيرة جدًا & الفصول~\ref{sec:glutcomputer}، \ref{sec:code}، \eqref{eq:window} \\
العصبون الطوري & مرشح تمرير عالٍ، كاشف حافة & يستجيب لتغير المدخل ثم يصمت & الفصل~\ref{sec:sensors} \\
العصبون المتكيّف & مرشح تمرير عالٍ مع تحكم آلي في الكسب & يهبط التردد مع المدخل الثابت & الفصل~\ref{sec:sensors}، \eqref{eq:nakarushton} \\
الرنّان & دارة رنين، مرشح تمرير نطاقي & أقوى استجابته لمدخل بتردده الخاص & القسم~\ref{sec:resonator} \\
المكامل غير التسريبي & مكامل بمضخم عمليات & يحوّل السرعة إلى موضع ويحفظه & القسم~\ref{sec:perfectint} \\
\midrule
\multicolumn{4}{l}{\emph{المحوّلات}} \\
العصبون التوتري & محوّل جهد إلى تردد & يتبع التردد المدخلَ خطيًا ولا يكاد يتعب & الفصل~\ref{sec:code} \\
العصبون المتدرج (بلا نبضات) & مضخم تماثلي، عازل & ينقل جهدًا متصلًا إلى مسافة قصيرة & الفصل~\ref{sec:sensors} \\
العصبون المقوِّم & مقوّم نصف موجة & التردد لا يكون سالبًا، $[u]_+$ & الفصلان~\ref{sec:glutcomputer}، \ref{sec:gaba} \\
زوج ”تشغيل/إطفاء“ & زوج مقوّمات، شفرة ثنائية السلك & أحدهما يستجيب لـ”أكثر“ والآخر لـ”أقل“ & الفصل~\ref{sec:glutcomputer}، \eqref{eq:dualrail} \\
عصبون بتثبيط تحويلي & مقسِّم تماثلي، ضبط الكسب & يقسم المدخل بدل أن يطرح منه & الفصل~\ref{sec:gaba}، \eqref{eq:shunt} \\
\midrule
\multicolumn{4}{l}{\emph{المولّدات والزمن}} \\
ناظمة الإيقاع & مذبذب استرخائي، هزّاز متعدد & يُصدر بنفسه نبضات بتردد ثابت & الفصلان~\ref{sec:serocomputer}، \ref{sec:dofa} \\
ناظمة إيقاع لها مدخل & مذبذب متحكَّم فيه بالجهد & إيقاع خاص يسرّعه المدخل أو يبطئه & الفصل~\ref{sec:chemoutput} \\
عصبون الرشقات & مولّد رشقات مبوَّب & يُصدر رشقات من نبضات متقاربة، أي ”شَرطات“ & الفصل~\ref{sec:code}، \eqref{eq:burst} \\
عصبون مقيَّد بالطور & كاشف طور، حلقة مقفلة الطور & يطلق عند طور محدد من الموجة الداخلة & الفصلان~\ref{sec:sensors}، \ref{sec:chemoutput} \\
محوار بتأخير & خط تأخير & يؤخر الإشارة زمنًا محددًا & الفصل~\ref{sec:glutcomputer}، الشكل~\ref{fig:itd} \\
\midrule
\multicolumn{4}{l}{\emph{الذاكرة والتبديل}} \\
العصبون ثنائي الاستقرار & قادح شميت، ماسك & مدخل قصير ينقله إلى حالة طويلة الأمد & القسم~\ref{sec:bistable} \\
عصبون بفروع تغصنية & مُضاعِف إرسال، شبكة صغيرة متعددة الطبقات & الفرع لا يمرّر المدخل إلا إذا وُجد مدخل سماح & الفصل~\ref{sec:synthesis} \\
\end{longtable}
\end{center}

تحفّظ واحد أهم من الجدول كله: هذه ليست خلايا مختلفة، بل \emph{أوضاع} مختلفة. العصبون نفسه قد يكون مكاملًا مع المدخل البطيء وكاشف تزامن حين يقصّر التثبيطُ $\tau$ الخاصة به (الفصل~\ref{sec:code})؛ وناظمة إيقاع في اليقظة ومتلقيًا صامتًا في النوم. أيّ جهاز سينتج تقرره القنوات الأيونية في الغشاء وقنوات البث التي تضبطها.

\section{أربعة أجهزة لم ترد في الكتاب بعد}

\subsection{الرنّان}
\label{sec:resonator}

لبعض العصبونات تيار بطيء يعاكس أي تغير في الجهد: إذا ارتفع الجهد شدّه التيار إلى الأسفل، وإذا انخفض شدّه إلى الأعلى، لكن بتأخر. هذا التيار يتصرف عند مهندس الإلكترونيات كالمحاثة، ويشكّل مع سعة الغشاء دارة رنين. يستجيب العصبون أقوى ما يستجيب لمدخل بتردد محدد، عادةً من بضعة هرتزات إلى عشرة أو عشرين، ويستجيب أضعف للأبطأ وللأسرع~\cite{HutcheonYarom2000}. إنه مرشح تمرير نطاقي في خلية واحدة: يختار من المدخل المضجوج إيقاعًا بتردده، كالإيقاع المحلي في الفصل~\ref{sec:gaba} مثلًا.

\subsection{المكامل غير التسريبي}
\label{sec:perfectint}

المكامل التسريبي لا يتذكر المدخل إلا زمنًا قدره $\tau$، أي عشرات الأجزاء من الألف من الثانية. لكن العين، كي تثبت، تحتاج إلى أن تتذكر موضعها ثوانيَ: أمر ”الدوران“ يأتي سرعةً قصيرة الأمد، والعين يجب أن تُمسَك في الموضع الجديد. تحل الشبكة ذلك بالتغذية الراجعة نفسها التي بُنيت بها الذاكرة في الفصل~\ref{sec:glutcomputer}. إذا كانت مجموعة من العصبونات ثابتها الزمني $\tau$ تثير نفسها بوزن $w$، فإن $\tau\,\dd r/\dd t=-r+w\,r+I$، والثابت الزمني للمجموعة يساوي
\begin{empheq}[box=\keybox]{equation}
\tau_{\text{eff}} \;=\; \frac{\tau}{1-w} .
\label{eq:perfectint}
\end{empheq}
عند $\tau=0.1\,\text{s}$ و$w=0.995$ نحصل على $20\,\text{s}$: مكامل شبه مثالي من عصبونات ينسى كل منها وحده في عُشر ثانية~\cite{Seung1996}. والثمن هو الضبط الدقيق: إذا زاد $w$ قليلًا على الواحد انجرف المكامل إلى الإشباع، وإذا نقص قليلًا نسي بسرعة. لذلك يحتاج هذا المكامل إلى ضبط دائم بالتعلم، كالذي وُصف في الفصل~\ref{sec:motorlearning}.

\subsection{العصبون ثنائي الاستقرار}
\label{sec:bistable}

في الفصلين~\ref{sec:glutcomputer} و\ref{sec:gaba} جُمع القادح من مجموعة عصبونات. لكن القادح قد يكون في خلية واحدة أيضًا. لبعض العصبونات تيار إذا اشتغل مرة حافظ على الإثارة بنفسه: مدخل قصير ينقل العصبون إلى عمل طويل، هو \emph{الهضبة}، والتثبيط يعيده إلى الراحة. هذا قادح شميت: للعصبون حالتان مستقرتان وبينهما تخلّف. هكذا تتصرف مثلًا عصبونات \nt{ACET} التي تتحكم في العضلات، وتشغّل هذه الخاصيةَ قناةُ بث: تظهر الهضبة حين يصل السيروتونين إلى العصبون~\cite{Hounsgaard1988}. الخلية نفسها ماسك مع \nt{SERO}، ومكامل عادي من دونه.

\subsection{المكاملات والرنّانات}

تقسم النظرية الخلايا القابلة للإثارة إلى صنفين~\cite{Izhikevich2007}. \emph{المكاملات} تشبه مذبذبًا متحكمًا فيه بالجهد يبدأ من الصفر: فوق العتبة بقليل يعمل العصبون ببطء كيفما شئنا، ويرتفع التردد بسلاسة مع المدخل. مثل هذا العصبون يجمع كل ما وصل، وينقل مقدار المدخل بالتردد نقلًا جيدًا. \emph{والرنّانات} تشبه مذبذبًا له تردد أدنى: لا تستطيع العمل ببطء، فعند المدخل الحدّي تُصدر فورًا نبضات بتردد منتهٍ وتفضّل المدخل الذي بترددها. الأولى هي الأجهزة الطبيعية لترميز التردد في الفصل~\ref{sec:code}، والثانية لترميز التوقيت.

\begin{keyblock}
\begin{proposition}[عصبون واحد، أجهزة كثيرة]
\label{prop:eetypes}
تعمل العصبونات، بحسب وظيفتها الكهربائية، مكاملاتٍ تسريبية وغير تسريبية، وكواشف تزامن، ومرشحات تمرير عالٍ ونطاقي، ومحوّلات جهد إلى تردد، ومقوّمات، ومقسِّمات، ومولّدات، وخطوط تأخير، وقوادح. وأيّ جهاز ينتج من خلية بعينها تحدده قنواتها الأيونية وقنوات البث التي تضبطها. الأوضاع نفسها مُقاسة؛ أما مدى استعمال الدماغ لهذه إعادة التهيئة في الحساب ففرضية في الغالب.
\end{proposition}
\end{keyblock}

يشبه هذا عند المهندس المصفوفة المنطقية القابلة للبرمجة: العتاد واحد، والجهاز تحدده التهيئة. غير أن التهيئة هنا لا تتغير عند البرمجة، بل أثناء العمل، بقنوات البث البطيئة التي تحدثنا عنها في الفصول~\babelsublr{\ref{sec:serocomputer}--\ref{sec:acet}}.

\section{تمارين}

\begin{exercise}[\lvlA]
\label{ex:18:1}
\emph{تفاوت المكامل غير التسريبي.} عصبونات بـ$\tau=0.1$ ثانية تحفظ موضع العين وفق~\eqref{eq:perfectint} مدة $T=1$ ثانية بخطأ لا يتجاوز $5\%$ في الاتجاهين. (أ)~أوجد المدى المسموح لـ$w$. (ب)~$w$ مجموع أوزان $K$ مشبكًا، لكل منها خطأ مستقل قدره $10\%$. عند أي $K$ يقع تشتت المجموع ضمن التفاوت؟
\end{exercise}

\begin{exercise}[\lvlB]
\label{ex:18:2}
\emph{الرنّان دارةً.} لنصف التيار البطيء في القسم~\ref{sec:resonator} بالمتغير $W$: $C\,\dd V/\dd t=-g_LV-g_wW+I$، $\tau_w\,\dd W/\dd t=V-W$. بيّن أن هذه دارة $RC$ مع فرع موازٍ $R_w=1/g_w$، $L_w=\tau_w/g_w$، وأوجد عند $C/g_L=10\ms$ و$\tau_w=100\ms$ و$\alpha=g_w/g_L=4$ تردد الرنين والمكسب $|Z|$ عنده على $|Z(0)|$.
\end{exercise}

\begin{exercise}[\lvlB]
\label{ex:18:3}
\emph{كيف يبدأ المكامل العمل.} (أ)~عصبون $\tau\,\dd V/\dd t=-V+\mu$، $\tau=10\ms$، بعتبة $\theta$ وإعادة ضبط إلى الصفر. عند أي $\mu/\theta-1$ يساوي التردد $1\,\text{Hz}$؟ (ب)~ما التردد الذي يعطيه النموذج $\tau\,\dd v/\dd t=v^2+\mu$ (النبضة ذهاب $v$ إلى $+\infty$، وإعادة الضبط إلى $-\infty$) عند $\mu=0.01$؟
\end{exercise}

\begin{exercise}[\lvlB]
\label{ex:18:4}
\emph{نمذجة: المكامل والرنّان.} نموذج التمرين~\ref{ex:18:2} مع $g_L=1$، $\tau_m=10\ms$، $\tau_w=100\ms$، عتبة $1$ وإعادة ضبط $V\to0$ ($W$ لا يتغير) يتلقى $\mu=1.5\sin2\pi ft$، $f=1\text{--}40\,\text{Hz}$. في أي نطاق ترددي يُصدر العصبون نبضات عند $\alpha=0$ وعند $\alpha=4$؟ قارن بالشرط $1.5\,|Z(f)|>1$.
\end{exercise}

\begin{exercise}[\lvlB]
\label{ex:18:5}
\emph{ماسك من خلية واحدة.} عصبون بتيار هضبة: $\tau\,\dd r/\dd t=-r+F(wr+I)$، $F(x)=\min(1,\max(0,x))$، $\tau=10\ms$، $w=1.5$، والخلفية $I=-0.2$. (أ)~ما أقصر نبضة $I=0.3$ تنقل العصبون من $r=0$ إلى الهضبة؟ (ب)~ما أصغر سعة $B$ لنبضة طويلة $I=-0.2-B$ تعيده إلى الراحة؟
\end{exercise}

\begin{exercise}[\lvlB]
\label{ex:18:6}
\emph{الضبط الذاتي للمكامل.} عند تثبيت النظر يكون المدخل $I=0$، ويعطي مستشعر الانزلاق سرعة الانجراف $e=\dd r/\dd t$، و$\dd w/\dd t=-\eta\,e\,r$. عند $w=0.95$، $\tau=0.1$ ثانية، $\langle r^2\rangle=1$ أوجد $\eta$ الذي يُدخل $|1-w|$ في تفاوت التمرين~\ref{ex:18:1} خلال $60$ ثانية من التثبيت. ماذا يفسد لو جرى التعلم أثناء دوران العين أيضًا؟
\end{exercise}

\begin{exercise}[\lvlC]
\label{ex:18:7}
\emph{لماذا تُعاد تهيئة الجهاز؟} قناة بث تبدّل $\alpha$ في نموذج التمرين~\ref{ex:18:2} بين $0$ و$4$. كم ضعفًا يتفوق الرنّان على المكامل في نسبة الإشارة إلى الضجيج لجيب ضعيف في ضجيج تيار أبيض عند $11\,\text{Hz}$ وعند $1\,\text{Hz}$؟ اقترح مسألة يكون فيها تبديل الوضع نفسه هو المفيد.
\end{exercise}

\chapter{العصبونات بحسب عدد التغصنات}
\chaptermark{عدد التغصنات}
\label{sec:dendrites}

\section{عدد المداخل معاملًا في الدارة}

صنّفنا العصبونات في الفصل~\ref{sec:neurontypes} بحسب ما يطلقه مخرجها، وفي الفصل~\ref{sec:eetypes} بحسب الجهاز الذي تعمل عمله. وثمة معيار ثالث يلاحظه المهندس على الفور: \emph{كم مدخلًا للعصبون}. تجمع المداخلَ التغصناتُ، وهي الفروع التي تحط عليها محاوير الخلايا الأخرى (الفصل~\ref{sec:dofa}، القسم~\ref{sec:weightstore}). بعض العصبونات بلا تغصنات أصلًا، وبعضها له تغصن أو اثنان، وبعضها له شجرة تجمع مئة ألف مدخل. هذا عند مهندس الإلكترونيات عامل التحميل عند المدخل، أي \foreignlanguage{english}{fan-in}، وهو يكاد يوافق المهمة دائمًا.

لماذا يهم عدد التغصنات وحجمها إلى هذا الحد؟ يتضح ذلك من تقدير بسيط. الغشاء مكثف سعته النوعية $c_m\approx1\,\text{\textmu F/cm}^2$، وكلما كبرت مساحة العصبون $A$ مع تغصناته، كبرت الشحنة التي يجب جلبها لرفع الجهد بمقدار $\Delta V$. والزمن الذي يفعل فيه تيار الدخل $I$ ذلك، إن أهملنا التسريب، هو
\begin{empheq}[box=\keybox]{equation}
t \;\approx\; \frac{c_m\,A\,\Delta V}{I} .
\label{eq:dendriteload}
\end{empheq}
العصبون الصغير ذو التغصنات القصيرة القليلة مساحته نحو $2\cdot10^{-5}\,\text{cm}^2$ وسعته نحو $20\,\text{pF}$. ومشبك واحد كبير بتيار بضعة نانوأمبيرات يرفعه $20\mV$ في أقل من عُشر جزء من الألف من الثانية. أما العصبون ذو الشجرة الكبيرة فمساحته أكبر بنحو خمس عشرة مرة وسعته نحو $300\,\text{pF}$؛ والمشبك العادي بتيار نحو $0.1\,\text{nA}$ كان سيشحنه في $60\ms$، وهو أطول بكثير من الثابت الزمني للتسريب، أي أنه لن يشحنه أبدًا. مثل هذا العصبون يحتاج إلى عشرات المداخل المتزامنة. الأعداد هنا رتب مقادير، لكن النتيجة لا تعتمد عليها: \emph{مداخل قليلة، سرعة ودقة؛ مداخل كثيرة، بطء وعمل بالمجموع}.

\begin{figure}[t]
\centering
\begin{tikzpicture}[>=Stealth,
  soma/.style={circle,draw,very thick,fill=blue!6,minimum size=0.55cm,inner sep=0pt},
  ax/.style={->,very thick,red!70!black},
  den/.style={very thick,blue!60!black}]
% 0 dendrites: sensor on a line
\begin{scope}[xshift=0cm]
\draw[den,decorate,decoration={snake,amplitude=1pt,segment length=4pt}] (-0.9,0) -- (0,0);
\node[font=\scriptsize] at (-0.9,0.3) {مستشعر};
\draw[ax] (0,0) -- (1.0,0);
\draw[thick] (0.1,0) -- (0.1,0.45);
\node[soma,minimum size=0.4cm] at (0.1,0.65) {};
\node[font=\scriptsize,align=center] at (0.05,-0.75) {$0$\\خط بمستشعر};
\end{scope}
% 1 dendrite
\begin{scope}[xshift=3.3cm]
\draw[den] (-0.9,0) -- (-0.28,0);
\node[soma] at (0,0) {};
\draw[ax] (0.28,0) -- (0.9,0);
\node[font=\scriptsize,align=center] at (0,-0.75) {$1$\\عازل، مكرِّر};
\end{scope}
% 2 dendrites: one per ear
\begin{scope}[xshift=6.2cm]
\draw[den] (-0.28,0) -- (-0.9,0); \draw[den] (0.28,0) -- (0.9,0);
\node[soma] at (0,0) {};
\draw[ax] (0,-0.28) -- (0,-0.6);
\node[font=\scriptsize] at (-0.9,0.25) {يسار}; \node[font=\scriptsize] at (0.9,0.25) {يمين};
\node[font=\scriptsize,align=center] at (0,-1.0) {$2$\\AND زمنية};
\end{scope}
% ~4 short dendrites: mixer
\begin{scope}[xshift=9.0cm]
\foreach \a in {45,135,225,315}{\draw[den] (\a:0.28) -- (\a:0.7); \fill[blue!60!black] (\a:0.72) circle (0.06);}
\node[soma] at (0,0) {};
\draw[ax] (0.28,0) -- (1.0,0);
\node[font=\scriptsize,align=center] at (0,-1.0) {$\sim4$\\مازج};
\end{scope}
% many: summing tree
\begin{scope}[xshift=12.0cm]
\foreach \a in {60,90,120}{\draw[den] (0,0.28) -- ++(\a:0.55) coordinate (b); \foreach \d in {-20,20}{\draw[den] (b) -- ++(\a+\d:0.35);}}
\foreach \a in {-120,-60}{\draw[den] (0,-0.28) -- ++(\a:0.45);}
\node[soma] at (0,0) {};
\draw[ax] (0.28,0) -- (1.0,0);
\node[font=\scriptsize,align=center] at (0,-1.0) {$10^4\text{--}10^5$\\جامع};
\end{scope}
\end{tikzpicture}
\caption{خمسة أصناف بحسب عدد المداخل. بالأزرق التغصنات (المداخل)، وبالأحمر المحوار (المخرج). كلما قلّت المداخل، نقل العصبون الإشارة المفردة أسرع وأدق؛ وكلما كثرت، أحسن الجمع والتصنيف. هذا مخطط، لا رسم لأشكال الخلايا.}
\label{fig:dendriteclasses}
\end{figure}

\section{صفر تغصنات: مستشعر في طرف الخط}

العصبونات الحسية للمس والألم وتمدد العضلات (الفصول~\ref{sec:sensors} و\ref{sec:pain} و\ref{sec:muscles}) ليس على جسم خليتها تغصن واحد. يخرج المحوار من الجسم فرعًا واحدًا وينقسم فورًا إلى اثنين: فرع يذهب إلى الجلد أو العضلة، وطرفه نفسه هو المستشعر؛ وفرع آخر يذهب إلى الحبل الشوكي. تجري النبضة من المستشعر إلى الحبل الشوكي مباشرة، متجاوزة جسم الخلية. هذا عند المهندس خط اتصال مستشعرُه في الطرف البعيد، وجسم الخلية معلّق جانبًا كوحدة تغذية: يمدّ الخط بالطاقة، لكنه لا يشارك في النقل. والمكسب هو السرعة والموثوقية: لا يوجد على الطريق موضع واحد يلزم فيه جمع الإشارة وتقرير إرسالها من جديد.

ومستشعرات الضوء والصوت حالة أكثر تطرفًا: فهي ليست عصبونات جامعة، بل محوّلات بحد ذاتها، مدخلها ليس مشبكًا بل بروتين مستقبِل (الشكل~\ref{fig:transducer}).

\section{تغصن واحد: عازل}

العصبون ذو التغصن الواحد والمحوار الواحد يتلقى خط دخل واحدًا وينقله إلى الأمام. هكذا بُنيت عصبونات الشم: تغصنها الوحيد يخرج إلى سطح الأنف ويحمل نوعًا واحدًا من المستقبلات~\cite{Mombaerts2004}؛ وهكذا بُنيت خلايا النقل في الشبكية بين مستشعرات الضوء والعصبونات المخرجة للعين؛ وهكذا بُنيت العصبونات التي تصل مستشعرات الأذن بالدماغ. كان المهندس سيسميها عازلًا أو مكرِّرًا: فهي لا تحسب، بل توصل إشارة واحدة، وقد تغيّر شكلها أحيانًا، كأن تحوّل الجهد المتصل للمستشعر إلى نبضات، كما في الفصل~\ref{sec:sensors}.

\section{تغصنان: بوابة AND زمنية}

كواشف التزامن التي تحدد اتجاه الصوت (الشكل~\ref{fig:itd}) كثيرًا ما يكون لها عند الثدييات تغصنان بالضبط، متجهان في اتجاهين متعاكسين: إلى أحدهما تصل المداخل من الأذن اليسرى، وإلى الآخر من اليمنى~\cite{Grothe2010}. لماذا تُفرَّق المداخل؟ لنتذكر الصيغة~\eqref{eq:shunt}: حين تنفتح قنوات مثيرة كثيرة في موضع واحد من الغشاء، يقترب الجهد هناك من جهد التوازن، ويضيف كل مدخل تالٍ أقل فأقل. لذلك فالمداخل الكثيرة من أذن واحدة، المجموعة على تغصن واحد، تُشبعه ولا تستطيع وحدها بلوغ العصبون العتبة. أما مدخل واحد من كل تغصن فيُجمعان خطيًا تقريبًا. التغصنان يحوّلان العصبون إلى بوابة AND أكثر صرامة: أطلِق فقط إذا جاءت الإشارتان من الجانبين معًا. وقد بيّنت النماذج أن هذا التفريق يحسّن فعلًا كشف التزامن~\cite{AgmonSnir1998}.

\section{تغصنات قصيرة قليلة: المازج ومُعيد الإرسال}

\emph{المازجات.} في دارة التعلم الحركي في الفصل~\ref{sec:motorlearning} يوسّع نحو سبعين مليار عصبون \nt{GLUT} صغير المدخلَ إلى متجه طويل جدًا. لكل منها نحو أربعة تغصنات قصيرة فقط، وينتهي كل تغصن بـ”مخلب“ يطوّق نهاية مدخل واحد. وهكذا لا يرى كل مازج إلا نحو أربعة مداخل من مداخل كثيرة، ويعمل بوابة AND صغيرة على رباعيته. من $M$ خط دخل يمكن اختيار $\binom{M}{4}$ رباعية مختلفة: عند $M=100$ يقارب ذلك أربعة ملايين. لماذا هذا العدد كله؟ العنصر العتبي ذو $n$ مدخلًا يستطيع أن يفصل فصلًا صحيحًا إلى صنفين ما لا يزيد على
\begin{empheq}[box=\keybox]{equation}
P_{\max} \;\approx\; 2n
\label{eq:cover}
\end{empheq}
نمطًا عشوائيًا~\cite{Cover1965}. العصبون المخرج في الدارة نفسها يتلقى نحو $10^5$ مدخل من المازجات، والحد الأعلى له وفق~\eqref{eq:cover} من رتبة $2\cdot10^5$ تقابل مختلف. هذا حد العنصر المثالي: إذا كانت كل الأوزان موجبة، كما في المشابك المثيرة، ولزم هامش للضجيج، فالتقدير للعصبون نفسه نحو $4\cdot10^4$ تقابل~\cite{Brunel2004}. المازجات ذات المداخل القليلة لازمة تحديدًا كي يكون للجامع الكبير مداخل كثيرة مختلفة شبه مستقلة~\cite{Marr1969,Albus1971}.

\emph{مُعيدات الإرسال.} على الطريق من الأذن إلى كواشف الاتجاه عصبونات بتغصنات قصيرة قليلة، تتلقى بضعة مشابك ضخمة فقط، وبعضها في الجوهر مشبكًا واحدًا. ووفق~\eqref{eq:dendriteload} هذا بالضبط ما يلزم: سعة صغيرة وتيار كبير، فيطلق العصبون بعد أجزاء من جزء الألف من الثانية من كل نبضة داخلة، دون أن يفقد تقريبًا شيئًا من دقة التوقيت. أحد مُعيدات الإرسال هذه يتلقى \nt{GLUT} ويُخرج \nt{GLYC}: إنه عاكس بدقة عشرات الميكروثواني، يزوّد كواشف الاتجاه بتثبيط سريع~\cite{BorstSoria2012}.

\section{آلاف المداخل: الجامع والمصنِّف}

في الطرف الآخر من السلّم عصبونات \nt{GLUT} القشرية ذات العشرة آلاف مدخل، وعصبونات \nt{GABA} المخرجة في دارة التعلم الحركي ذات المئة ألف. هذا العصبون بطيء وفق~\eqref{eq:dendriteload} ولا يستطيع الاستجابة لمدخل مفرد: إنه يجمع. إنه العصبون المكامل في الفصل~\ref{sec:glutcomputer}، والجامع الذي تُبنى منه المصنِّفات. والشجرة الكبيرة تضيف شيئًا جديدًا أيضًا: تعمل فروعها دارات فرعية مستقلة لكل منها عتبتها، فيتصرف العصبون الواحد كشبكة صغيرة متعددة الطبقات~\cite{Beniaguev2021,Gidon2020}.

\section{تغصنات هي المخرج أيضًا}

ثمة عصبونات لا محوار لها أصلًا، وتغصناتها هي المدخل والمخرج معًا: تتلقى الإشارات وتطلق الناقل العصبي بنفسها. وأكثر الأمثلة دراسةً عصبونات \nt{GABA} في الشبكية المشاركة في تحديد اتجاه الحركة (الفصل~\ref{sec:gaba}، الشكل~\ref{fig:direction}). في هذا العصبون يستجيب كل فرع بمفرده للحركة من جسم الخلية نحو طرفه، ويُصدر التثبيط من هذا الطرف~\cite{Euler2002}. العصبون الواحد عدة كواشف اتجاه مستقلة، واحد لكل فرع. وهو عند المهندس ليس عنصرًا واحدًا، بل دارة متكاملة بعدة قنوات في غلاف واحد.

\section{الخلاصة}

\begin{table}[t]
\centering
\footnotesize
\setlength{\tabcolsep}{3pt}
\renewcommand{\arraystretch}{1.15}
\begin{tabular}{>{\raggedright}p{1.9cm}>{\raggedright}p{3.4cm}>{\raggedright}p{2.6cm}>{\raggedright}p{1.8cm}>{\raggedright\arraybackslash}p{3.8cm}}
\toprule
التغصنات & مثال & الجهاز & المداخل & ما المعروف \\
\midrule
$0$ & مستشعرات اللمس والألم والتمدد & خط بمستشعر في طرفه & --- & ثابت \\
$1$ & عصبونات الشم، خلايا النقل في الشبكية، عصبونات الأذن & عازل، مكرِّر & من $1$ إلى بضعة & ثابت \\
$2$ & كواشف اتجاه الصوت & بوابة AND زمنية & حزمتان & البنية ثابتة؛ ومكسب تفريق المداخل مبيَّن في النماذج \\
$\sim4$ & مازجات دارة التعلم الحركي & بوابة AND صغيرة & $\sim4$ & ثابت؛ ودور توسيع المدخل فرضية جيدة الأساس \\
قليلة، قصيرة & مُعيدات الإرسال على الطريق من الأذن & مُعيد إرسال، عاكس & من $1$ إلى بضعة ضخمة & ثابت \\
آلاف & عصبونات \nt{GLUT} القشرية، العصبونات المخرجة للتعلم الحركي & جامع، مصنِّف & $10^4\text{--}10^5$ & ثابت؛ والفروع دارات فرعية مُقاسة في أمثلة منفردة \\
تغصنات مُخرِجة & عصبونات \nt{GABA} للاتجاه في الشبكية & دارة متكاملة متعددة القنوات & كثيرة & مُقاس \\
\bottomrule
\end{tabular}
\caption{العصبونات بحسب عدد التغصنات.}
\label{tab:dendrites}
\end{table}

\begin{keyblock}
\begin{proposition}[عدد المداخل يتبع المهمة]
\label{prop:dendrites}
حيث تلزم سرعة الإشارة المفردة ودقتها، في المستشعرات ومُعيدات الإرسال وكواشف التزامن، يكون للعصبونات تغصنات قليلة ومداخل قليلة ومشابك كبيرة: السعة الصغيرة~\eqref{eq:dendriteload} تُشحن بسرعة. وحيث يلزم الجمع والتصنيف، يكون للعصبونات أشجار كبيرة بآلاف المداخل. بنية الأصناف مُقاسة؛ والتفسير الحسابي جيد الأساس، لكنه يبقى فرضية لأصناف كثيرة.
\end{proposition}
\end{keyblock}

يكمّل الجدول~\ref{tab:dendrites} التصنيفين السابقين: بحسب الناقل العصبي (الجدول~\ref{tab:types}) وبحسب الجهاز (الجدول~\ref{tab:eetypes}). والثلاثة تنظر إلى العنصر نفسه من جهات مختلفة: ماذا يُخرج، وماذا يحسب، وإلى كم يستمع.

\section{تمارين}

\begin{exercise}[\lvlA]
\label{ex:19:1}
\emph{كم مدخلًا يحتاج العصبون الكبير.} عصبون بـ$C=300\,\text{pF}$ و$\tau_m=20\ms$ يتلقى مشابك هي مصادر تيار كل منها $0.1\,\text{nA}$، والعتبة أعلى من الراحة بـ$20\mV$. (أ)~كم مشبكًا يجب أن يعمل معًا لبلوغ العتبة أصلًا؛ وفي $10\ms$؛ وفي $5\ms$؟ (ب)~في كم يبلغ العتبةَ عصبونٌ بـ$C=20\,\text{pF}$ من مشبك $4\,\text{nA}$؟
\end{exercise}

\begin{exercise}[\lvlA]
\label{ex:19:2}
\emph{لماذا أربعة.} المازج بوابة AND على $k$ خطًا من $M=100$، ونصف الخطوط نشط في النمط، والمازجات $7\cdot10^{10}$. (أ)~كم منها يستجيب للنمط عند $k=2$؛ $4$؛ $40$؟ (ب)~كم ضعفًا تزيد المازجات عدد التقابلات~\eqref{eq:cover} للعصبون المخرج ذي $10^5$ مدخل مقارنة بـ$M=100$ خط مباشر؟
\end{exercise}

\begin{exercise}[\lvlB]
\label{ex:19:3}
\emph{من أين يأتي $2n$.} المستوى الفائق المار بالأصل يقسم $P$ نقطة في وضع عام في فضاء ذي $n$ بعدًا إلى صنفين بـ$C(P,n)=2\sum_{k=0}^{n-1}\binom{P-1}{k}$ طريقة. (أ)~أثبت أنه عند $P=2n$ يكون نصف التقسيمات الـ$2^P$ بالضبط قابلًا للفصل. (ب)~ما النسبة القابلة للفصل عند $n=100$ و$P=180$؛ $220$؟
\end{exercise}

\begin{exercise}[\lvlB]
\label{ex:19:4}
\emph{التغصنان بوابةَ AND صارمة.} التغصن حجيرة بتسريب $g_L$، وكل مدخل يفتح عليه موصلية $g_0=0.5\,g_L$ بجهد $E_E$، والجسم يرى $\kappa(V_1+V_2)$، والعتبة $\theta=1.1\,\kappa E_E$. لماذا لا يطلق العصبونَ أيُّ عدد من المداخل على تغصن واحد، وكم مدخلًا يلزم على كل تغصن؟
\end{exercise}

\begin{exercise}[\lvlB]
\label{ex:19:5}
\emph{تصميم مُعيد إرسال.} ارتعاش النبضة $\sigma_t=\sigma_V/\dot V$، حيث $\dot V$ ميل الجهد عند العتبة. المطلوب $\sigma_t\le20\,\text{\textmu s}$ عند ضجيج $\sigma_V=1\mV$؛ أهمل التسريب. كم مشبكًا متزامنًا بـ$0.1\,\text{nA}$ يلزم لذلك عصبونًا بـ$C=300\,\text{pF}$، وما ارتعاش عصبون بـ$C=20\,\text{pF}$ ومشبك واحد $4\,\text{nA}$؟
\end{exercise}

\begin{exercise}[\lvlB]
\label{ex:19:6}
\emph{نمذجة: شحن تغصن طويل.} حاكِ كابلًا منفعلًا $\tau_m\partial_tV=\lambda^2\partial_x^2V-V+i/g$ طوله $L=\ell/\lambda$ بطرفين مغلقين ($40$ حجيرة، أويلر). تُحقن في الطرف البعيد نبضة تيار قصيرة: متى، وإلى أي نسبة من الجهد الناتج عن الشحنة نفسها موزعة على التغصن كله، يرتفع الطرف القريب عند $L=0.2$؛ $1$؛ $2$؟
\end{exercise}

\begin{exercise}[\lvlC]
\label{ex:19:7}
\emph{بحث: العدد الأمثل للمداخل.} السعة $C=C_0+nc_s$، ويعمل معًا $m$ مدخلًا تيار كل منها $I_s$، ويحفظ العصبون $P\approx2n$ تقابلًا~\eqref{eq:cover}. كيف يعتمد زمن الاستجابة على $P$ عند $m$ ثابت وعند نسبة ثابتة $m=\varphi n$؟ اقترح معيار كلفة وأوجد $n$ الأمثل لمُعيد الإرسال وللمصنِّف.
\end{exercise}

\chapter{ماذا تفعل العصبونات في النوم}
\chaptermark{النوم}
\label{sec:sleep}
\begin{chaptergoals}
\item لماذا النوم مجموعة أوضاع يبدّلها \nt{ACET} و\nt{NORA} و\nt{SERO}، لا إطفاء؛
\item كيف يكوّن ضغط النوم وعتبتان مُرحِّلًا بتخلّف يضبطه الإيقاع اليومي؛
\item كيف يحوّل التعب القشرة إلى مذبذب استرخائي، وكيف يوقف \nt{ACET} الأرجحة؛
\item كيف يُعاد تشغيل اليوم مسرَّعًا، ولماذا قد تضعف الأوزان كلها بالنسبة نفسها في النوم.
\end{chaptergoals}

\section{النوم ليس إطفاءً، بل وضع آخر}

المهندس الذي يرى حاسوبًا مطفأً سيقول: إنه لا يفعل شيئًا. أما الدماغ النائم فليس كذلك. العصبونات لا تصمت في النوم: في النوم العميق تُصدر القشرة نبضات، وفي النوم السريع تُصدر منها قرابة ما تُصدره نهارًا. الذي يتغير ليس \emph{هل} تعمل العصبونات، بل \emph{كيف} تعمل. النوم مجموعة من أوضاع الآلة، وتبدّلها قنوات البث نفسها التي تعمل نهارًا.

لكل وضع يمكن كتابة ”حالة السجلات“ (الجدول~\ref{tab:sleepmodes}). في النهار تعمل \nt{ACET} و\nt{NORA} و\nt{SERO}، والمستشعرات موصولة، والعضلات تطيع. في النوم البطيء العميق تهدأ القنوات الثلاث كلها، ويضعف المدخل من أعضاء الحس~\cite{Hasselmo1999}: يهبط إطلاق الأسيتيل كولين في القشرة، وتُصدر عصبونات \nt{NORA} و\nt{SERO} نبضات أقل مما في النهار~\cite{Marrosu1995,AstonJonesBloom1981,McGintyHarper1976}. وفي النوم السريع صورة غريبة: يرتفع \nt{ACET} من جديد إلى مستوى النهار~\cite{Marrosu1995}، بينما يصمت \nt{NORA} و\nt{SERO} صمتًا شبه تام~\cite{AstonJonesBloom1981,McGintyHarper1976,JacobsAzmitia1992}. وعضلات الجسم مطفأة في النوم السريع: عصبونات \nt{ACET} التي تتحكم فيها تُثبَّط بقوة عبر \nt{GABA} و\nt{GLYC}~\cite{BrooksPeever2012}، بالمنع نفسه الذي رأيناه في الفصل~\ref{sec:gaba}، لكنه مفروض على المخرج كله دفعة واحدة.

\begin{table}[t]
\centering
\footnotesize
\setlength{\tabcolsep}{4pt}
\renewcommand{\arraystretch}{1.15}
\begin{tabular}{>{\raggedright}p{3.1cm}>{\raggedright}p{2.9cm}>{\raggedright}p{3.2cm}>{\raggedright\arraybackslash}p{3.4cm}}
\toprule
& اليقظة & النوم البطيء & النوم السريع \\
\midrule
\nt{ACET} (الكتابة، الفصل~\ref{sec:acet}) & مرتفع & منخفض & مرتفع \\
\nt{NORA} (الكسب، الفصل~\ref{sec:nora}) & متوسط، مع ومضات & منخفض & شبه صامت \\
\nt{SERO} (الفصل~\ref{sec:serocomputer}) & متوسط & منخفض & شبه صامت \\
المدخل من المستشعرات & موصول & مُضعَف & مُضعَف \\
المخرج إلى العضلات & موصول & مُضعَف & مطفأ بالتثبيط \\
نشاط القشرة & منتظم & أرجحة بطيئة: عمل، صمت & منتظم، كما في النهار \\
\bottomrule
\end{tabular}
\caption{حالة سجلات الآلة في الأوضاع الثلاثة. كل الأسطر مُقاسة؛ ومستويات \nt{ACET} و\nt{NORA} و\nt{SERO} بتسجيلات مباشرة بحسب مراحل النوم~\cite{Marrosu1995,AstonJonesBloom1981,McGintyHarper1976}.}
\label{tab:sleepmodes}
\end{table}

\section{متى ننام: مُرحِّل بتخلّف}

ما الذي يقرر متى تنتقل الآلة إلى النوم؟ يعمل هنا جيدًا مخطط بسيط من جزأين~\cite{Borbely1982}. الجزء الأول \emph{ضغط النوم} $S$، يتراكم ما دمنا مستيقظين ويُفرَّغ ما دمنا نائمين. إنه مكثف يُشحن نهارًا ويُفرَّغ ليلًا:
\begin{empheq}[box=\keybox]{equation}
\frac{\dd S}{\dd t} \;=\; \frac{1-S}{\tau_r}\ \ \text{(اليقظة)},\qquad
\frac{\dd S}{\dd t} \;=\; -\frac{S}{\tau_d}\ \ \text{(النوم)} .
\label{eq:sleeppressure}
\end{empheq}
الجزء الثاني عتبتان: تنام الآلة حين يبلغ $S$ العتبة العليا $H(t)$، وتستيقظ حين ينخفض $S$ إلى الدنيا $L(t)$. وكلتا العتبتين تتأرجح بسلاسة مع الإيقاع اليومي من الفصل~\ref{sec:chemoutput}. فنحصل على مُرحِّل بتخلّف، أي قادح شميت من القسم~\ref{sec:bistable}، ويعطي مع المكثف مذبذبًا استرخائيًا يضبط الإيقاعُ اليومي طورَه~\eqref{eq:pll}.

\begin{figure}[t]
\centering
\begin{tikzpicture}[>=Stealth,x=0.16cm,y=4.2cm]
\foreach \a/\b in {0/4.3,21.2/29.3,45.4/53.4,69.5/72}{\fill[blue!8] (\a,0) rectangle (\b,1);}
\draw[->,gray!70!black] (0,0) -- (75,0) node[right,font=\small,black] {$t$، ساعة};
\draw[->,gray!70!black] (0,0) -- (0,1.08) node[left,font=\small,black] {$S$};
\foreach \t in {0,24,48,72}{\draw[gray!60!black] (\t,-0.02) -- (\t,0.02); \node[below,font=\scriptsize] at (\t,0) {\t};}
\draw[thick,densely dashed,red!70!black] plot file {figures/sleep_H.dat};
\draw[thick,densely dashed,green!45!black] plot file {figures/sleep_L.dat};
\draw[very thick,blue!60!black] plot file {figures/sleep_S.dat};
\node[font=\scriptsize,red!70!black,anchor=west] at (73,0.72) {$H$: النوم};
\node[font=\scriptsize,green!45!black,anchor=west] at (73,0.24) {$L$: الاستيقاظ};
\node[font=\scriptsize,blue!60!black] at (25.25,0.93) {نوم};
\node[font=\scriptsize,blue!60!black] at (49.4,0.93) {نوم};
\end{tikzpicture}
\caption{ضغط النوم $S$~\eqref{eq:sleeppressure} عند $\tau_r=18.2$ ساعة و$\tau_d=4.2$ ساعة. في النهار يُشحن $S$ حتى العتبة العليا $H$، وفي الليل (المظلَّل) يُفرَّغ حتى الدنيا $L$؛ والعتبتان تتأرجحان مع الإيقاع اليومي. تصل الآلة بنفسها إلى نظام من نحو $16$ ساعة يقظة و$8$ ساعات نوم.}
\label{fig:twoprocess}
\end{figure}

في نموذجنا، بالقيم المعتادة $\tau_r=18.2$ ساعة و$\tau_d=4.2$ ساعة، تصل الآلة بنفسها إلى $16.1$ ساعة يقظة و$8.0$ ساعات نوم (الشكل~\ref{fig:twoprocess}). ويفسر النموذج أيضًا ما يبدو غريبًا: بعد أربعين ساعة بلا نوم يبلغ الضغط $0.90$، لكن نوم التعافي في النموذج لا يدوم إلا $8.8$ ساعة، لا يومًا كاملًا زيادة على المعتاد. فالتفريغ أُسّي، والشحنة العالية تزول بسرعة؛ و”الدَّين“ يُسدَّد بعمق النوم لا بطوله.

ما الذي يؤدي دور $S$ فيزيائيًا؟ المرشح القوي هو الأدينوزين، ”عدّاد الحمل“ من الملحق~\ref{sec:othermediators}: يرتفع تركيزه خلال اليقظة وينخفض في النوم~\cite{PorkkaHeiskanen1997,Dunwiddie2001}. أما أن ضغط النوم هو الأدينوزين بالذات ولا شيء غيره ففرضية.

\section{النوم البطيء: الشبكة كلها مذبذب استرخائي}

في النوم العميق تعمل عصبونات القشرة بالتناوب مجموعةً كاملة: أقل من مرة في الثانية تشتغل معًا لبضع مئات من أجزاء الألف من الثانية (\emph{حالة العمل}) وتصمت معًا (\emph{حالة الصمت}): تردد هذه الأرجحة أقل من $1\,\text{Hz}$، وتحت التخدير غالبًا $0.3\text{--}0.4\,\text{Hz}$~\cite{Steriade1993}. يمكننا تجميع هذه الأرجحة البطيئة من القطع المتاحة لدينا. لنأخذ مجموعة من عصبونات \nt{GLUT} بتغذية راجعة قوية على نفسها، أي الذاكرة من الفصل~\ref{sec:glutcomputer} التي لها حالتان مستقرتان، ولنضف تعبًا بطيئًا، كما في مولّد الإيقاع في الفصل~\ref{sec:muscles}:
\begin{empheq}[box=\keybox]{equation}
\tau\,\frac{\dd r}{\dd t} = -r + F\bigl(w\,r + I - a\bigr),
\qquad
\tau_a\,\frac{\dd a}{\dd t} = -a + b\,r .
\label{eq:updown}
\end{empheq}
تشتغل المجموعة وتعمل وتتعب؛ ويأكل التعبُ $a$ الحالةَ العليا، فتسقط المجموعة في الصمت؛ وفي الصمت يزول التعب، وتختفي الحالة الدنيا، فتشتغل المجموعة من جديد. فنحصل على المذبذب الاسترخائي نفسه الذي رأيناه في ضغط النوم، لكنه أسرع بنحو ثمانين ألف مرة.

\begin{figure}[t]
\centering
\begin{tikzpicture}[>=Stealth,x=2.9cm,y=1.0cm]
\begin{scope}[yshift=1.9cm]
\draw[->,gray!70!black] (0,0) -- (4.1,0);
\draw[->,gray!70!black] (0,0) -- (0,1.25) node[left,font=\small,black] {$r$};
\draw[thick,blue!60!black] plot file {figures/updown_sleep.dat};
\node[font=\scriptsize,anchor=west] at (4.15,0.5) {نوم: $b=5$};
\end{scope}
\draw[->,gray!70!black] (0,0) -- (4.1,0) node[right,font=\small,black] {$t$، ثانية};
\draw[->,gray!70!black] (0,0) -- (0,1.25) node[left,font=\small,black] {$r$};
\draw[thick,red!70!black] plot file {figures/updown_wake.dat};
\node[font=\scriptsize,anchor=west] at (4.15,0.5) {نهار: $b=1$};
\foreach \t in {0,1,2,3,4}{\draw[gray!60!black] (\t,-0.05) -- (\t,0.05); \node[below,font=\scriptsize] at (\t,0) {\t};}
\end{tikzpicture}
\caption{المجموعة نفسها~\eqref{eq:updown} في وضعين: $w=5$، $I=2$، $\theta=2.5$، $\tau=10\ms$، $\tau_a=0.3$ ثانية. مع تعب قوي ($b=5$، كما في النوم البطيء) تتأرجح المجموعة بين العمل والصمت بدور $1.1$ ثانية (قيمة النموذج؛ وفي القشرة الدور أطول من ثانية، وتحت التخدير نحو $3$ ثوانٍ). وإذا أضعفنا التعب ($b=1$)، وهذا ما يفعله الأسيتيل كولين، عملت المجموعة نفسها بانتظام.}
\label{fig:updown}
\end{figure}

فما الذي ينقل الشبكة من الأرجحة إلى العمل النهاري المنتظم؟ الأسيتيل كولين يثبّط في العصبونات التيارات البطيئة التي تصنع التعب~\cite{McCormick1992}. في نموذجنا، مع التعب القوي $b=5$ تتأرجح المجموعة بدور $1.1$ ثانية، وهو أسرع من الأرجحة المُقاسة لكنه من رتبتها، وتعمل نحو $35\%$ من الوقت إذا أُخذ المتوسط على عدد صحيح من الأدوار؛ ومع التعب الضعيف $b=1$ تعمل المجموعة نفسها بانتظام (الشكل~\ref{fig:updown}). سجل واحد، هو مستوى \nt{ACET}، يحوّل المذبذب إلى مضخم. وفوق الأرجحة البطيئة تسري في النوم ومضات أسرع من إيقاع $7\text{--}15$ ذبذبة في الثانية؛ تولّدها حلقة \foreignlanguage{english}{\nt{GLUT}--\nt{GABA}} بين القشرة وعقدة الترحيل الوسطى~\cite{SteriadeMcCormickSejnowski1993}، وهي قريبة الإيقاع الذي رأيناه في الفصل~\ref{sec:gaba}.

\section{الإعادات: تشغيل اليوم مسرَّعًا}

رأينا في الفصل~\ref{sec:acet} أن العصبونات التي عملت معًا نهارًا تطلق في النوم البطيء معًا من جديد وبالترتيب نفسه. ولنضف تفصيلًا هندسيًا واحدًا: الإعادات تجري \emph{مسرَّعة}. التسلسل الذي استغرق ثواني في النهار يُعاد في النوم في أعشار من الثانية~\cite{Nadasdy1999}: في الحُصين أسرع بنحو عشرين مرة~\cite{LeeWilson2002}، وفي القشرة بست أو سبع مرات~\cite{Euston2007}. ولا يُعاد في أي وقت اتفق: الومضات القصيرة للإعادات في منطقة واحدة تتزامن مع أرجحة العمل والصمت ومع الومضات السريعة للإيقاع في القشرة. وإذا عُزّز هذا التزامن اصطناعيًا تحسنت الذاكرة~\cite{Maingret2016}، وهذه علاقة سببية لا مجرد تزامن.

لماذا تشغّل الآلة يومها مسرَّعًا؟ يعرف المهندس صعوبة تسمى \emph{النسيان الكارثي}: الشبكة التي تتعلم الجديد متتابعًا، شيئًا بعد شيء، تفسد القديم. والمنقذ هو \emph{الخلط}: أن يُتعلم الجديد مخلوطًا بالقديم، على دفعات صغيرة، مرات كثيرة. وفرضية نظامي التعلم~\cite{McClelland1995} تقول إن الذاكرة السريعة تسجل اليوم كاملًا، ثم تعيد تشغيله في النوم للقشرة البطيئة شيئًا فشيئًا، مخلوطًا ومرات كثيرة. هذا هو الزوج نفسه ”مخزن مؤقت سريع، مستودع بطيء“ في الفصل~\ref{sec:acet}، والزوج نفسه من المتعلم السريع والمتعلم البطيء في الفصل~\ref{sec:motorlearning}. الإعادات وتزامنها مُقاسة؛ أما أن غرضها هو التعلم المخلوط للذاكرة البطيئة ففرضية جيدة الأساس.

\section{إعادة تطبيع الأوزان}

في النهار تقوّي قواعد هب من الفصل~\ref{sec:dofa} الوصلات في الغالب. ولو اقتصر الأمر على التقوية لنمت الأوزان، ونمت معها كلفة الطاقة وهبطت الحساسية: العصبون الذي كل مداخله قوية يكفّ عن التمييز بينها. وبحسب فرضية \emph{الاتزان الداخلي المشبكي}~\cite{TononiCirelli2014}، يلزم النوم، ضمن ما يلزم له، لإعادة الأوزان إلى مستوى العمل: تضعف الوصلات كلها بالنسبة نفسها، فتبقى \emph{نسبها}، أي ما تعلّمته الآلة. وهذا عند المهندس تطبيع، كما في القاعدة~\eqref{eq:oja}، غير أنه لا يجري أثناء العمل، بل في وقت منفصل.

والقياسات المباشرة لا تناقض ذلك: عند الفئران تكون مساحة التماس في مشابك القشرة بعد النوم أصغر بنحو $18\%$ مما بعد اليقظة، والتناقص متناسب مع الحجم، أما أكبر التماسات وأثبتها فلا تكاد تتغير~\cite{deVivo2017}. تحجيم الأوزان مُقاس؛ أما أنه المهمة الرئيسية للنوم ففرضية.

\section{النوم السريع: آلة مفصولة عن مداخلها ومخارجها}

في النوم السريع تعمل القشرة تقريبًا كما في النهار، لكن المدخل من أعضاء الحس مُضعَف، والمخرج إلى العضلات مطفأ بالتثبيط. هذا عند المهندس وضع مألوف: \emph{الاختبار على منصة}. الدارة تعمل بكامل قدرتها، لكن مداخلها موصولة بمولّد داخلي، ومخارجها مفصولة عن المشغّلات كي لا ينكسر شيء. والأحلام، بحسب إحدى الفرضيات، هي بالضبط هذا التشغيل الداخلي لنموذج العالم الذي جُمع خلال النهار؛ أما ما تفعله الشبكة حينئذ بالضبط وهل تتعلم، فغير معروف. المُقاس هنا هو فقط ما في الجدول~\ref{tab:sleepmodes}: أي قناة مشغّلة وأيها صامتة، وأن المخرج محجوب.

\section{الصيانة والنوم الموضعي}

ساد طويلًا أن الفراغات بين الخلايا تتسع في النوم فيُغسل الدماغ أسرع من فضلات الأيض~\cite{Xie2013}. لكن تجارب حديثة قاست التصريف بطريقة أخرى فوجدت العكس: في النوم يكون أبطأ~\cite{Miao2024}. السؤال مفتوح الآن، وسنتركه مفتوحًا.

وثمة حقيقة أوثق: النوم خاصية للشبكات الموضعية، لا للآلة كلها دفعة واحدة. إذا حُرم حيوان من النوم طويلًا، فإن مناطق منفردة من قشرته، والحيوان مستيقظ يتحرك، تسقط لحظات قصيرة في الصمت كما في النوم البطيء، ويخطئ الحيوان في تلك اللحظات أكثر~\cite{Vyazovskiy2011}. كل شبكة تتعب بنفسها وتتبدل بنفسها؛ ويتكون الوضع العام حين تنقل قنوات البث الشبكات كلها معًا.

\begin{keyblock}
\begin{proposition}[النوم مجموعةَ أوضاع]
\label{prop:sleep}
في النوم لا تنطفئ العصبونات: قنوات البث تنقل الآلة إلى أوضاع أخرى. ومتى ننام يقرره مُرحِّل بتخلّف~\eqref{eq:sleeppressure} يضبطه الإيقاع اليومي. في النوم البطيء تصير الشبكة مذبذبًا استرخائيًا~\eqref{eq:updown} وتعيد تشغيل اليوم مسرَّعًا؛ وفي النوم السريع تعمل مفصولةً عن المداخل والمخارج. الأوضاع والإعادات وتحجيم الأوزان مُقاسة؛ أما غرضها، أي التعلم المخلوط وإعادة التطبيع، ففرضيات جيدة الأساس.
\end{proposition}
\end{keyblock}

\section{الخلاصة}

\begin{table}[t]
\centering
\footnotesize
\setlength{\tabcolsep}{4pt}
\renewcommand{\arraystretch}{1.15}
\begin{tabular}{>{\raggedright}p{3.2cm}>{\raggedright}p{4.4cm}>{\raggedright\arraybackslash}p{5.2cm}}
\toprule
ماذا يحدث & كيف & ما المعروف \\
\midrule
تبدّل الأوضاع & السجلات \nt{ACET} و\nt{NORA} و\nt{SERO} (الجدول~\ref{tab:sleepmodes}) & مُقاس \\
متى ننام & المكثف $S$ ومُرحِّل بتخلّف~\eqref{eq:sleeppressure} & النموذج يصف التجارب جيدًا؛ الأدينوزين بوصفه $S$ فرضية \\
الأرجحة البطيئة & مجموعة بتغذية راجعة وتعب~\eqref{eq:updown} & الأرجحة مُقاسة؛ والنموذج أبسط مخطط \\
تشغيل اليوم مسرَّعًا & إعادات أسرع بـ$6\text{--}20$ مرة، متزامنة مع الأرجحة & مُقاس؛ وتعزيز التزامن يحسّن الذاكرة \\
لماذا الإعادات & التعلم المخلوط للذاكرة البطيئة & فرضية جيدة الأساس \\
إعادة تطبيع الأوزان & تحجيم الوصلات في النوم & تناقص التماسات مُقاس؛ وكونه الدور الرئيسي للنوم فرضية \\
النوم السريع & العمل مع مداخل ومخرج مفصولة & الوضع مُقاس؛ والغرض غير معروف \\
التصريف & اتساع الفراغات بين الخلايا & نتائج متناقضة \\
النوم الموضعي & مناطق منفردة تسقط في الصمت & مُقاس \\
\bottomrule
\end{tabular}
\caption{ماذا تفعل العصبونات في النوم.}
\label{tab:sleep}
\end{table}

يلخص الجدول~\ref{tab:sleep} الفصل. تبيّن أن النوم ليس توقفًا في عمل الآلة، بل صيانتها أثناء العمل: تبديل الأوضاع بقنوات البث نفسها، ومولّد من القطع نفسها التي يُبنى منها إيقاع الخطوات، وتشغيل اليوم مسرَّعًا للذاكرة البطيئة، وإعادة تطبيع الأوزان. في النهار تكتب الآلة، وفي الليل تعيد كتابة ما كُتب وترتّبه.

\section{تمارين}

\begin{exercise}[\lvlA]
\label{ex:20:1}
\emph{مُرحِّل بلا إيقاع يومي.} لتكن العتبتان ثابتتين: $H=0.67$، $L=0.17$ (متوسطا عتبتي نموذج الشكل~\ref{fig:twoprocess}). استخرج من~\eqref{eq:sleeppressure} مدتي اليقظة والنوم ودور المذبذب عند $\tau_r=18.2$ ساعة و$\tau_d=4.2$ ساعة. لماذا يصل النموذج ذو العتبات المتأرجحة مع ذلك إلى $16.1+8.0=24$ ساعة؟ وإلى أي جهاز في الفصل~\ref{sec:chemoutput} يعود هذا الأثر؟
\end{exercise}

\begin{exercise}[\lvlA]
\label{ex:20:2}
\emph{دَين النوم.} النوم الذي يبدأ عند ضغط $S_0$ يدوم $t_s=\tau_d\ln(S_0/L)$، $\tau_d=4.2$ ساعة، و$L$ العتبة الدنيا عند الاستيقاظ. (أ)~بعد $40$ ساعة بلا نوم $S_0=0.90$، ودام النوم $8.8$ ساعة. كم كانت $L$؟ وكم كان النوم سيدوم مع $L\approx0.092$ المعتادة؟ (ب)~خلال الليل تنقص مساحة التماسات المشبكية $18\%$. بكم يجب أن تنمو الأوزان خلال النهار؟
\end{exercise}

\begin{exercise}[\lvlB]
\label{ex:20:3}
\emph{تصميم العتبات.} اختر عتبتين ثابتتين $H$ و$L$ يعطي عندهما المُرحِّل~\eqref{eq:sleeppressure} مع $\tau_r=18.2$ ساعة و$\tau_d=4.2$ ساعة $16$ ساعة يقظة و$8$ ساعات نوم بالضبط. بمَ تكون هذه الآلة أسوأ من الآلة ذات العتبات المتأرجحة إذا أُوقظت ذات ليلة بعد $4$ ساعات؟
\end{exercise}

\begin{exercise}[\lvlB]
\label{ex:20:4}
\emph{دور الأرجحة البطيئة.} في النموذج~\eqref{eq:updown} $F(x)=1/\bigl(1+e^{-(x-\theta)/k}\bigr)$، $w=5$، $I=2$، $\theta=2.5$، $k=0.25$، $b=5$، $\tau\ll\tau_a=0.3$ ثانية. (أ)~أوجد نقطتي الانعطاف $a_{\text{hi}}$ و$a_{\text{lo}}$ للمنحنى $r=F(wr+I-a)$، حيث يختفي العمل والصمت. (ب)~باعتبار $r\approx1$ في العمل و$r\approx0$ في الصمت، أوجد الدور ونسبة العمل.
\end{exercise}

\begin{exercise}[\lvlB]
\label{ex:20:5}
\emph{متى تنطفئ الأرجحة.} في نموذج التمرين~\ref{ex:20:4} تقع النقطة الثابتة عند تقاطع المنحنى $a=wr+I-F^{-1}(r)$ مع المستقيم $a=br$. تستمر التذبذبات ما دام التقاطع على الفرع الأوسط، بين نقطتي الانعطاف. عند أي $b$ يكون ذلك؟ وكيف يحوّل \nt{ACET} المذبذب إلى مضخم بتغيير $b$؟
\end{exercise}

\begin{exercise}[\lvlB]
\label{ex:20:6}
\emph{نمذجة: دَين أم طور.} في \texttt{sleep\_models.py} خذ الاستيقاظ الأول بعد $48$ ساعة وأبقِ الآلة مستيقظة $D=24$ و$32$ و$40$ و$48$ و$64$ ساعة (\texttt{stay\_awake}، \texttt{days=8}). كم يدوم نوم التعافي عند كل $D$؟ وما الذي يحدد مدته؟
\end{exercise}

\begin{exercise}[\lvlB]
\label{ex:20:7}
\emph{نمذجة: النسيان والخلط.} عصبون خطي $y=\mathbf w\cdot\mathbf x$، $n=40$، يتعلم بالقاعدة $\mathbf w\leftarrow\mathbf w-0.5\,(y-y^*)\,\mathbf x$. المهمتان A وB، لكل منهما $10$ أنماط، $x_j\sim\mathcal N(0,1/n)$، $y^*=\pm1$. ما جذر متوسط مربع الخطأ على A بعد $200$ حقبة من A ثم $200$ حقبة من B؟ وبعد $200$ حقبة من A وB مخلوطتين؟
\end{exercise}

\begin{exercise}[\lvlC]
\label{ex:20:8}
\emph{بحث: الخلط أم إعادة التطبيع؟} عصبون التمرين~\ref{ex:20:7} مع عتبة؛ وإجراءان ليليان: (i)~تُضرب كل الأوزان في $0.82$؛ (ii)~تُعاد الأنماط القديمة مخلوطة بالجديدة. ماذا يفعل كل منهما بنسب الأوزان وباستجابات العصبون؟ وكيف يمكن التمييز بين الفرضيتين بأحجام المشابك بعد النوم؟
\end{exercise}

\chapter{آلة من عصبونات حية}
\chaptermark{آلة من عصبونات حية}
\label{sec:livingmachine}
\begin{chaptergoals}
\item لماذا تطلق مزرعة بلا قنوات بثّ رشقات شبكية، وما الذي يحدد الفاصل بينها؛
\item كيف تغيّر الشبكات الحية استجاباتها تحت التغذية الراجعة، وما الذي لم يثبت بعد؛
\item كيف تعمل العضيّة خزّانًا يُدرَّب قارئه الخطي في الخارج، في السيليكون؛
\item كيف عُلِّمت المزارع حتى الآن، ولماذا ما زال المعلّم خارج الطبق.
\end{chaptergoals}

\section{منصة اختبار بمخطط مكشوف}

في الفصل~\ref{sec:debugging} صحّحنا أخطاء آلة لا مخطط لها: يسمع المسبار ألف عصبون من أصل ثمانين مليارًا، وكل ما عدا ذلك علينا أن نخمّنه. والمهندس في هذا الموقف يفعل شيئًا بسيطًا: يجمّع قطعة صغيرة من الدارة على لوحة تجارب، حيث يرى كل مكوّن. ويمكن فعل ذلك مع العصبونات أيضًا.

تُفصل عصبونات القشرة بعضها عن بعض وتُزرع على شريحة فيها مصفوفة أقطاب، من بضع عشرات من الأقطاب إلى عشرات الآلاف. وفي غضون أسابيع قليلة تمدّ الخلايا زوائدها وتتصل في شبكة من آلاف العصبونات أو مئات الآلاف، أغلبها \nt{GLUT} ونسبة أصغر \nt{GABA} تتوقف على المستحضر: ففي مزارع الحُصين تبلغ نحو $6\%$ من العصبونات~\cite{Benson1994}. وكل قطب يستطيع أن يستمع، كمسبار الفصل~\ref{sec:debugging}، وأن يحقن تيارًا، كإشارة اختبار. والنوع الثاني من المنصات هو \emph{العضيّات} (\foreignlanguage{english}{organoids}): كُتل ثلاثية الأبعاد من النسيج العصبي قطرها نحو مليمتر، تُستنبت من خلايا جذعية بشرية~\cite{Lancaster2013}. على هذه المنصات تعمل الشبكات الحية شهورًا، وهناك مرافق يمكن فيها إجراء تجارب على العضيّات عن بُعد، عبر الشبكة، أكثر من مئة يوم متواصلة~\cite{Jordan2024}. ويسمى هذا المجال \emph{الحوسبة البيولوجية} أو \emph{ذكاء العضيّات}~\cite{Smirnova2023}.

تختلف المنصة عن الدماغ في أمرين مهمين. فهي ضئيلة: عصبوناتها أقل بمئة ألف مرة. وليس فيها قنوات بثّ: مزرعة القشرة لا تحوي عصبونات \nt{DOPA} ولا \nt{SERO} ولا \nt{NORA} ولا \nt{ACET}. إنها آلة لها ناقل بيانات وتحكم في التدفق، لكن لا سجلات تحكم لها. فلننظر ماذا تستطيع.

\section{ما تفعله الشبكة من تلقاء نفسها}

المزرعة المتروكة وشأنها لا تصمت، ولا تُصدر ضجيجًا منتظمًا. فمن حين لآخر تشتعل الشبكة كلها تقريبًا دفعة واحدة، في \emph{رشقة شبكية}، ثم تهدأ من جديد؛ والفواصل بين الرشقات الشبكية من ثانية إلى دقائق، بحسب المزرعة~\cite{Wagenaar2006}. والصورة مألوفة للمهندس: مضخِّم ذو تغذية راجعة موجبة قوية وبلا حِمل يتحول إلى مذبذب.

الآلية نعرفها. الروابط المتبادلة القوية بين عصبونات \nt{GLUT} تُطلق انهيارًا متسلسلًا، كما في الفصل~\ref{sec:gaba} عند الإخلال بشروط القضية~\ref{prop:ei}. والانهيار يستنزف بسرعة مخزون الناقل العصبي في المشابك~\eqref{eq:depression}، فتصمت الشبكة. ثم يتعافى المخزون بثابت زمني $\tau_{\text{rec}}$، وحين تتعافى النسبة $x^*$ التي تكفي لانهيار جديد تأتي رشقة شبكية جديدة. والفاصل بين الرشقات الشبكية
\begin{empheq}[box=\keybox]{equation}
T_{\text{burst}} \;\approx\; \tau_{\text{rec}}\,\ln\frac{1}{1-x^*} .
\label{eq:burstinterval}
\end{empheq}
مع $\tau_{\text{rec}}=0.8\,\text{s}$ من الفصل~\ref{sec:code} و$x^*=0.9$ نحصل على نحو ثانيتين، ومع $x^*=0.99$ على نحو أربع. هذا تقدير النموذج بقيمة $\tau_{\text{rec}}$ المعتمدة في الكتاب، وهو يقع عند الطرف القصير للفواصل المقيسة. أما القياس المباشر في المزارع الدقيقة فيُظهر أن مخزون الناقل يتعافى بعد الرشقة الشبكية في عشرات الثواني~\cite{CohenSegal2011}، أي أن $\tau_{\text{rec}}$ هناك أكبر؛ ومع قيمة كهذه تعطي الصيغة~\eqref{eq:burstinterval} عشرات الثواني والدقائق، كما في المزارع البطيئة. هذا مذبذب من نوع الاسترخاء، قريب لمولّد الإيقاع في الفصل~\ref{sec:muscles}: هناك كان يعيد شحنَه تعبُ مجموعات العصبونات، وهنا مخزونُ الناقل.

الرشقات الشبكية هي ما تفعله الشبكة حين لا يكون لديها ما تفعله. وفي الدماغ الحي تظهر صورة مشابهة في النوم العميق (الفصل~\ref{sec:sleep}) وفي الدماغ النامي، قبل أن تصله مداخل حقيقية. والتشابه يوحي بفكرة، لكن تطابق الآليات فرضية.

\section{كيف نعلّم الشبكة بلا معلّم دوباميني}

ما دام \nt{DOPA} غائبًا عن المزرعة، فعلينا نحن أن نقدّم إشارة ”أفضل أو أسوأ“، عبر الأقطاب. وتُظهر التجارب أن الشبكة على المنصة تغيّر استجاباتها فعلًا، وأطرف ما فيها هو \emph{بماذا} تُعلَّم.

\emph{معلّم يصمت.} تُحفَّز الشبكة عبر قطب واحد مرة كل ثانية إلى ثلاث ثوانٍ، إلى أن تظهر على قطب آخر الاستجابة المطلوبة بعد $50\pm10\ms$ من المنبه. وما إن تظهر حتى يُوقَف التحفيز خمس دقائق. وبعد عدة دورات كهذه تظهر الاستجابة المطلوبة على المنبه فورًا~\cite{Shahaf2001}. المكافأة هنا هي الصمت: المنبه يخلخل الشبكة، وإيقافه يتركها في الحالة التي أعطت الجواب الصحيح. هذا هو النزول على التدرج عبر اضطرابات عشوائية من الفصل~\ref{sec:dofa} (القضية~\ref{prop:gradient})، حيث يقوم الهزّ نفسه مقام الخطأ: نهزّ حتى يصير الجواب صحيحًا.

\emph{شبكة تلعب التنس.} في تجربة على منصة بمصفوفة أقطاب كثيفة، وُصل نحو مليون عصبون بلعبة تنس حاسوبية بمضرب واحد~\cite{Kagan2022}. كان موضع الكرة يُقدَّم تيارًا إلى أقطاب ”حسية“: أين الكرة، بترميز المكان؛ وكم تبعد، بالتردد. وكانت الفعالية في منطقتين ”حركيتين“ تحرّك المضرب إلى أعلى أو إلى أسفل. وكانت قاعدة التعلّم غير مألوفة. إذا صدّ المضرب الكرة تلقّت الشبكة منبهًا قصيرًا \emph{متوقَّعًا}؛ وإذا أخفق تلقّت بضع ثوانٍ من تيار عشوائي \emph{غير متوقَّع}. وفي عشرين دقيقة من اللعب طالت سلاسل الإصابات. لم يظهر تحسّن ذو دلالة داخل الجلسة إلا مع التغذية الراجعة الكاملة؛ ومع الصمت بدل المنبه لعبت المزرعة أيضًا في آخر الجلسة أفضل مما بلا تغذية راجعة، لكن بأثر أصغر، ولم يكن هناك تعلّم ثابت من جلسة إلى أخرى على مدى أيام. والتحليل المفصّل لهذه التجربة في الفصل~\ref{sec:dishbrain}.

فسّر الباحثون النتيجة بفرضية أن الشبكة تتصرف بحيث يصير دخلها متوقَّعًا~\cite{Friston2010}. وهي الفكرة نفسها التي رأيناها في الاقتصاد في النبضات في الفصل~\ref{sec:code} وفي التنبؤ بالإطار في الفصل~\ref{sec:recognition}: الآلة تنفق أقل حين يكون الدخل متوقَّعًا. لكن التجربة نفسها، وخاصة كلام الباحثين عن ”وعي“ المزرعة، لقيت نقدًا حادًا: الأثر صغير، ويمكن تفسيره بما هو أبسط~\cite{Balci2023}. والتقييم المنصف هو هذا: المزرعة على المنصة تغيّر سلوكها تحت تأثير التغذية الراجعة، وهذا مُقاس؛ أما أنها تتعلم تحديدًا التنبؤ بدخلها ففرضية.

\emph{شبكة تقود جسمًا.} وبطريقة مشابهة وُصلت مزرعة بـ”جسم“ افتراضي بسيط وعُلّمت أن تحرّكه نحو هدف، باختيار النمط المكاني الزماني لمنبهات التدريب~\cite{Bakkum2008}. لا دوبامين في أيٍّ من هذه التجارب: المعلّم نمط تيار يختاره المهندس. أما ما يتغير حين يصير المعلّم برنامجًا أو الدوبامين نفسه، ففي الأقسام~\babelsublr{\ref{sec:dishdopamine}--\ref{sec:cartpole}}.

\begin{figure}[t]
\centering
\begin{tikzpicture}[>=Stealth,
  blk/.style={draw,very thick,rounded corners=2pt,align=center,font=\footnotesize,minimum height=0.8cm,fill=blue!5}]
% (a) closed loop
\node[font=\small] at (1.5,4.3) {(أ) حلقة مغلقة};
\node[blk,text width=2.6cm] (G) at (1.5,3.3) {اللعبة: كرة ومضرب};
\draw[very thick,fill=orange!10,rounded corners=3pt] (0.5,0.2) rectangle (2.5,1.6);
\foreach \x in {0.7,1.0,...,2.35}{\foreach \y in {0.4,0.7,1.0,1.3}{\fill[gray!60] (\x,\y) circle (0.04);}}
\draw[->,thick,blue!60!black] (G.west) -| (-0.5,0.9) -- (0.5,0.9);
\node[font=\scriptsize,blue!60!black,anchor=east] at (-0.6,2.1) {الكرة $\leftarrow$ تيار};
\draw[->,thick,red!70!black] (2.5,0.9) -- (3.5,0.9) |- (G.east);
\node[font=\scriptsize,red!70!black,anchor=west] at (3.6,2.1) {النبضات $\leftarrow$ المضرب};
\node[font=\scriptsize] at (1.5,-0.2) {عصبونات على مصفوفة أقطاب};
\node[font=\scriptsize,align=center] at (1.5,-0.85) {إصابة: منبه متوقَّع\\إخفاق: تيار عشوائي};
% (b) reservoir
\begin{scope}[xshift=7.9cm]
\node[font=\small] at (1.6,4.3) {(ب) خزّان};
\node[blk,text width=1.1cm] (X) at (-0.4,1.0) {الدخل\\$u(t)$};
\draw[very thick,fill=orange!10] (1.6,1.0) circle (0.8);
\foreach \a/\r in {20/0.35,80/0.5,150/0.3,210/0.55,280/0.4,330/0.2,110/0.15}{\fill[red!60!black] ({1.6+\r*cos(\a)},{1.0+\r*sin(\a)}) circle (0.05);}
\node[font=\scriptsize] at (1.6,-0.2) {عضيّة: بلا معلّم};
\node[blk,text width=1.8cm,fill=green!8] (W) at (4.0,1.0) {القراءة\\$W_{\text{out}}$};
\draw[->,thick] (X) -- (0.8,1.0);
\draw[->,thick] (2.4,1.0) -- node[above,font=\scriptsize] {$\mathbf r(t)$} (W);
\draw[->,thick] (W) -- (4.0,2.3) node[above,font=\scriptsize] {$y(t)$};
\node[font=\scriptsize,green!40!black] at (4.0,0.3) {تتعلم بمعلّم};
\end{scope}
\end{tikzpicture}
\caption{طريقتان لجعل شبكة حية تحسب. (أ) حلقة مغلقة: موضع الكرة يُقدَّم تيارًا، ونبضات المنطقتين الحركيتين تحرّك المضرب، والمنبه المتوقَّع أو العشوائي بعد الضربة هو المعلّم. (ب) خزّان~\eqref{eq:reservoir}: لا تتلقى العضيّة إشارة تعليم، بل تفكّك الدخل إلى استجابات لاخطية كثيرة ذات ذاكرة؛ ولا يتعلم بالخطأ إلا قارئ خطي واحد في الخارج. وروابط العضيّة نفسها تتغير أيضًا في أثناء ذلك، بلا معلّم.}
\label{fig:livingmachine}
\end{figure}

\section{خزّان حي}

وهناك طريقة أخرى لاستعمال الشبكة الحية: ألّا نعلّمها أصلًا. تسمى هذه الطريقة في التقنية \emph{الحوسبة بالخزّان}~\cite{Maass2002,JaegerHaas2004}. نأخذ نظامًا لاخطيًا جاهزًا ذا ذاكرة، أيَّ نظام، بشرط أن تكون استجابته للدخل غنية وأن تتوقف على الماضي القريب. الدخل $u(t)$ يهزّه، فنحصل على متجه استجابات $\mathbf r(t)$، ولا يُدرَّب إلا القارئ الخطي
\begin{empheq}[box=\keybox]{equation}
y(t) \;=\; W_{\text{out}}\,\mathbf r(t) ,
\label{eq:reservoir}
\end{empheq}
الذي تُضبط أوزانه بقاعدة المربعات الصغرى~\eqref{eq:lms} من الفصل~\ref{sec:motorlearning}. يفعل الخزّان ما فعلته عصبونات \nt{GLUT} الصغيرة في الفصل~\ref{sec:motorlearning}: يفكّك الدخل إلى متجه طويل تصير فيه الأصناف المطلوبة قابلة للفصل بمستوٍ (الفصل~\ref{sec:dendrites}).

وقد صلحت عضيّة على مصفوفة أقطاب لدور هذا الخزّان. فأصوات الكلام، حين قُدّمت إليها أنماطًا من التيار، ميّزها قارئ خطي واحد بعد تدريبه بحسب المتكلم بدقة نحو $78\%$~\cite{Cai2023}. ودقة $78\%$ رقم يذكره الباحثون أنفسهم وتكرره الصحافة. النتيجة مفيدة، لكن من المهم أن نفهم أين ”الذكاء“ فيها: العضيّة تقدّم اللاخطية والذاكرة المتلاشية، أما التعلم بالخطأ فيجري في الخارج، في السيليكون (الشكل~\ref{fig:livingmachine}). ومع ذلك لا تبقى العضيّة على حالها: فبحسب الباحثين، تغيّرت في أثناء التدريب اتصاليتها الوظيفية نفسها، وهم يسمّون ذلك تعلّمًا بلا معلّم في العضيّة ذاتها~\cite{Cai2023}. ولم يُفصل بعدُ أيّ جزء من الدقة يعود إلى هذا التغير وأيّ جزء إلى القارئ.

\section{ما تقوله المنصة عن الآلة}

\emph{الطاقة.} بحسب التقدير~\eqref{eq:energy} تكلّف النبضة نحو $10^{-10}\,\text{J}$. فمزرعة من مليون عصبون بتردد نبضة واحدة في الثانية تنفق على نقل الإشارات
\begin{empheq}[box=\keybox]{equation}
P \;\approx\; 10^{6}\times1\,\text{s}^{-1}\times10^{-10}\,\text{J} \;=\; 10^{-4}\,\text{W},
\label{eq:dishpower}
\end{empheq}
أي عُشر ميلي واط. أما الإلكترونيات التي تحفّز هذه النبضات وتسجلها وتعالجها فتنفق أكثر بمراتب. وكما في الفصل~\ref{sec:debugging}، عنق الزجاجة ليس الحساب، بل الاتصال به.

\emph{القطع الناقصة.} تُظهر المنصة كم يستطيع ناقل \nt{GLUT} مع تحكم \nt{GABA} في التدفق، بلا سجلات تحكم: أن ينتظم ذاتيًا، ويولّد الرشقات الشبكية، ويغيّر استجاباته تحت تأثير التغذية الراجعة، ويكون خزّانًا جيدًا. وما لا يستطيعه بدونها هو أن يقرر بنفسه ما الذي يُعدّ نجاحًا. هذه الإشارة يقدّمها على المنصة المهندس، وفي الدماغ، كما رأينا في الفصول~\babelsublr{\ref{sec:dofa}--\ref{sec:acet}}، قنوات البثّ.

\emph{الأخلاقيات.} تُستنبت العضيّات من خلايا بشرية، وكلما ازدادت تعقيدًا صار السؤال أجدّ: هل يمكن لنسيج كهذا أن يشعر بشيء؟ والنظرة الهندسية توحي بجواب جزئي: الألم في الفصل~\ref{sec:pain} ليس إشارة من مستشعر واحد، بل منظومة إنذار كاملة ببوابات ومقبض لشدة الصوت وقنوات بثّ، وكل ذلك غائب عن المنصة. لكن هذا استدلال لا برهان، والمجال نفسه يقترح أن يجري التقييم الأخلاقي جنبًا إلى جنب مع التجارب، منذ البداية~\cite{Smirnova2023}.

\begin{keyblock}
\begin{proposition}[شبكة حية على المنصة]
\label{prop:livingmachine}
شبكة من عصبونات \nt{GLUT} و\nt{GABA} على مصفوفة أقطاب تولّد الرشقات الشبكية من تلقاء نفسها~\eqref{eq:burstinterval}، وتغيّر استجاباتها تحت تأثير التغذية الراجعة، ويمكن أن تكون خزّانًا~\eqref{eq:reservoir} لقارئ مدرَّب في الخارج. كل هذا مُقاس. أما أن هذه الشبكة تتعلم بقاعدة عامة ما، كأن تتنبأ بدخلها، ففرضية، والعبارات الرنانة عن ”وعيها“ لا يقبلها أغلب المختصين.
\end{proposition}
\end{keyblock}

\section{دوبامين بومضة ضوء}
\label{sec:dishdopamine}

التجربة المباشرة الوحيدة التي نعرفها، حيث أُعطيت المزرعة الدوبامين معلّمًا، أُجريت على منصة عضيّات يتصل بها الباحثون عن بُعد~\cite{Jordan2024}. أُضيف إلى السائل الذي يغمر العضيّات دوبامين داخل ”قفص“ حساس للضوء: فالجزيء المقيَّد لا يؤثر في المستقبلات. وتركيز الدوبامين المقيَّد $1\mM$. وحين يلزم أن تُكافأ الشبكة يُسلَّط عليها ضوء فوق بنفسجي: فينفتح القفص ويخرج الدوبامين إلى الحجم.

والحلقة مبنية هكذا. يُقدَّم المنبه نفسه مرة بعد مرة؛ فإذا أثار نبضات أكثر من ذي قبل أُطلقت الومضة وتحرر الدوبامين. وعلى مدى $240$ تكرارًا ازداد عدد النبضات المستثارة إجمالًا. ويكتب الباحثون أنفسهم أنه لا يمكن من تجربة واحدة كهذه أن يُستنتج أن الزيادة سببها الدوبامين~\cite{Jordan2024}.

هذه، في عين المهندس، أول محاولة لتجميع قاعدة العوامل الثلاثة~\eqref{eq:threefactor} في طبق: فعالية المرسل والمستقبِل تترك أثرًا، والإشارة العامة تقرر هل يُثبَّت التغيير. لكن الإشارة هنا موجبة فقط: المكافأة إما موجودة وإما غائبة، والجهاز لا يُصدر خطأً سالبًا $\delta<0$. ثم إن الدوبامين ينتشر ويُزال ببطء، كالتركيز في~\eqref{eq:lowpass}، فقد ترفع الومضات المتكررة مستواه العام لا غير. ولتمييز التعلم من ذلك تلزم تجربتان ضابطتان: العدد نفسه من الومضات في لحظات عشوائية لا صلة لها باستجابة الشبكة، والومضات نفسها بلا دوبامين مقيَّد. ويلزم أيضًا اختبار بعد الراحة: هل بقي التغيير بعد أن زال الدوبامين؟

\section{الدوبامين يعلّم السيليكون}
\label{sec:biohybrid}

في التجربة المعاكسة تعلّم الخلايا الحية الإلكترونيات~\cite{Keene2020}. تُستنبت خلايا تفرز الدوبامين في قناة دقيقة فوق عنصر عصبي الشكل عضوي، أي ترانزستور تؤدي موصلية قناته دور وزن المشبك. والدوبامين الخارج من الخلايا يتأكسد عند قطب العنصر، فيغيّر ذلك الموصلية. ويحاكي الجهاز أيضًا آلية إزالة الناقل من الشق، فيمكن أن يُغيَّر الوزن تغييرًا دائمًا وأن يُعاد أيضًا إلى ما كان عليه.

من حيث الدارة هذا مكامل تسريبي بدخل كيميائي: الوزن يراكم تدفق الدوبامين، و”الإزالة“ تحدد سرعة نسيانه، وهي دارة \foreignlanguage{english}{RC} نفسها التي في~\eqref{eq:lowpass}. والمهم هنا اتجاه الاتصال. فمسبار الفصل~\ref{sec:debugging} لا يرى سجلات البثّ لأنها كيميائية. وهذا العنصر يقرأ سجلًا كيميائيًا بالذات ويحوّله إلى وزن كهربائي. وهو، بالنسبة إلى واجهات الدماغ والحاسوب، طريق إلى مستشعر يسمع لا ناقل البيانات وحده، بل سجلات التحكم أيضًا.

\section{عضيّات بعصبوناتها الدوبامينية}
\label{sec:midbrainorg}

صار في الإمكان أن تُستنبت من الخلايا الجذعية البشرية عضيّات تشبه منطقة الدماغ التي توجد فيها عصبونات \nt{DOPA}. وفيها عصبونات دوبامينية عاملة تفرز الدوبامين~\cite{Jo2016}. وتُستنبت هذه العضيّات في الغالب لدراسة الأمراض. ولم نجد تجارب استُعمل فيها دوبامينها معلّمًا لشبكة أخرى.

وهذه، بالنسبة إلى الآلة المصنوعة من عصبونات حية، هي القطعة الناقصة. فمنصة الفصل~\ref{sec:livingmachine} ناقل بيانات وتحكم في التدفق بلا سجلات تحكم؛ وهنا مصدر إشارة البثّ موجود. ما ينقص هو الناقد: شبكة تحسب خطأ التنبؤ~\eqref{eq:ctd} وتتحكم في هذه العصبونات. ووصل عضيّة قشرية بعضيّة كهذه بحيث يُفرز الدوبامين بحسب الخطأ مسألة مفتوحة، وهو حتى الآن فرضية لا تجربة.

\section{عضيّة توازن عمودًا بلا دوبامين}
\label{sec:cartpole}

أكمل التجارب الحديثة يستغني عن الدوبامين تمامًا~\cite{Robbins2026}. وُصلت عضيّات قشرة الفأر في حلقة مغلقة بمسألة ”العمود على العربة“: فعالية العضيّة تحرّك العربة، وموضع العمود يعود إليها بالتحفيز. وبين المحاولات تتلقى العضيّة منبهات تعليم قصيرة عالية التردد. أيّها بالضبط، يختاره برنامج تعلّم معزَّز.

النتائج مقيسة:
\begin{itemize}
\item في أغلب العضيّات أعطت إشارات التعليم التي اختارها البرنامج نتيجة أفضل من إشارات مختارة عشوائيًا أو من غياب الإشارة؛
\item بعد $45$ دقيقة من الراحة لم يبقَ التحسن؛
\item حجب مستقبلات الغلوتامات بنوعيها، \foreignlanguage{english}{AMPA} و\foreignlanguage{english}{NMDA}، ألغى التحسن.
\end{itemize}

البند الأخير يدل على مكان ما تعلّمته العضيّة: في مشابك \nt{GLUT}، وعبر المستقبِل الذي يعمل في القسم~\ref{sec:weightstore} كاشفًا لتزامن الدخل والخرج. والبند الثاني يبيّن ما ينقص: التغير السريع موجود، والتثبيت غائب، كالمتعلم السريع في الفصل~\ref{sec:motorlearning} بلا المتعلم البطيء. أما دور المعلّم فتغيّر: حلّ برنامج محل المهندس الذي يختار نمط التيار، لكن المعلّم لا يزال خارج الطبق.

ولم نجد كذلك بين التجارب الأقدم على تعليم المزارع فوق مصفوفات الأقطاب تجربة واحدة استُعمل فيها الدوبامين: إشارة التعليم في كل مكان كانت تحفيزًا كهربائيًا.

\section{من يعلّم المزرعة}

\begin{table}[t]
\centering
\footnotesize
\setlength{\tabcolsep}{4pt}
\renewcommand{\arraystretch}{1.15}
\begin{tabular}{>{\raggedright}p{2.7cm}>{\raggedright}p{2.5cm}>{\raggedright}p{3.2cm}>{\raggedright\arraybackslash}p{4.4cm}}
\toprule
التجربة & مَن يعلّم & إشارة التعليم & ما هو معروف \\
\midrule
إيقاف التحفيز عند الاستجابة المطلوبة & المهندس & انتهاء التحفيز & تثبت الاستجابة، وهذا مُقاس \\
\foreignlanguage{english}{DishBrain} (الفصل~\ref{sec:dishbrain}) & المهندس & تيار متوقَّع أو عشوائي & نمو ذو دلالة في الجولات داخل الجلسة مع التغذية الراجعة الكاملة فقط، وأثر أصغر مع الصمت، وهذا مُقاس؛ غير ثابت من جلسة إلى أخرى؛ والسبب موضع نقاش \\
العمود على العربة & برنامج تعلّم معزَّز & منبهات عالية التردد يختارها البرنامج & أفضل من العشوائية، لكنه لا يبقى بعد الراحة، وهذا مُقاس \\
دوبامين بالومضة & قاعدة ”نبضات أكثر = مكافأة“ & دوبامين يحرّره الضوء & زيادة النبضات ملاحظة؛ دور الدوبامين غير مثبت \\
مشبك هجين حيوي & خلايا حية & دوبامين من الخلايا & تغيّر طويل لوزن العنصر وإعادته، وهذا مُقاس \\
عضيّات بعصبونات \nt{DOPA} & --- & --- & مصدر الدوبامين موجود؛ ولم يُستعمل معلّمًا \\
\bottomrule
\end{tabular}
\caption{من يعلّم العصبونات الحية على الشريحة، وبماذا.}
\label{tab:dishteachers}
\end{table}

في كل التجارب التي تتعلم فيها المزرعة نفسها، المعلّم حتى الآن في الخارج (الجدول~\ref{tab:dishteachers}): مهندس، أو برنامج، أو قاعدة وضعها مهندس. ولم يدخل الدوبامين الطبق إلا مرة واحدة، ودوره لم يُثبت بعد. أما تجميع قناة بثّ حقيقية في الطبق، بناقد وخطأ وأثر كما في الفصل~\ref{sec:dofa}، فهو الخطوة التالية التي لم يخطُها أحد بعد.


\section{خلاصة}

\begin{table}[t]
\centering
\footnotesize
\setlength{\tabcolsep}{4pt}
\renewcommand{\arraystretch}{1.15}
\begin{tabular}{>{\raggedright}p{2.8cm}>{\raggedright}p{4.2cm}>{\raggedright\arraybackslash}p{5.8cm}}
\toprule
التجربة & ما تفعله الشبكة الحية & ما هو معروف \\
\midrule
شبكة بلا دخل & رشقات شبكية بفواصل من ثانية إلى دقائق~\eqref{eq:burstinterval} & مُقاس؛ والآلية عبر استنزاف المشابك نموذج مقبول عمومًا \\
معلّم يصمت & تثبت الاستجابة المطلوبة حين يُوقَف التحفيز & مُقاس \\
لعبة التنس & تطول سلاسل الإصابات؛ والتحسن ذو الدلالة داخل الجلسة لا يظهر إلا مع التغذية الراجعة الكاملة & تغيّر السلوك مُقاس؛ التعلم من جلسة إلى أخرى غير ثابت؛ حجم الأثر وتفسيره بالتنبؤ بالدخل موضع خلاف \\
جسم افتراضي & حركة نحو الهدف بنمط منبهات مختار & مُقاس \\
خزّان & لاخطية وذاكرة من أجل القراءة~\eqref{eq:reservoir} & مُقاس؛ التعلم بالخطأ في الخارج، في السيليكون؛ وروابط العضيّة تتغير أيضًا \\
الطاقة & نحو $10^{-4}\,\text{W}$ لكل مليون عصبون~\eqref{eq:dishpower} & تقدير وفق~\eqref{eq:energy} \\
\bottomrule
\end{tabular}
\caption{ما أظهرته الآلات المصنوعة من عصبونات حية.}
\label{tab:livingmachine}
\end{table}

الآلة المصنوعة من عصبونات حية (الجدول~\ref{tab:livingmachine}) منصة اختبار نرى عليها ما يستطيعه ناقل البيانات والتحكم في التدفق بلا سجلات تحكم. يستطيعان الكثير، لكنهما لا يستطيعان أن يقررا بنفسيهما ما هو الجيد. على المنصة يؤدي هذا الدور المهندس بنمط التيار الذي يختاره؛ وفي الدماغ تؤديه أسلاك البثّ القليلة التي خُصّص لها وسط الكتاب.

\section{تمارين}

\begin{exercise}[\lvlA]
\label{ex:21:1}
\emph{الفاصل بين الرشقات الشبكية.} تستنزف الرشقة الشبكية المخزون حتى الصفر، ثم $\tau_{\text{rec}}\,\dd x/\dd t=1-x$، وتبدأ رشقة شبكية جديدة عند $x=x^*$؛ $\tau_{\text{rec}}=0.8\,\text{s}$. (أ)~أوجد الفاصل عند $x^*=0.9$ و$0.99$. (ب)~العتبة $x^*$ تتذبذب عشوائيًا بمقدار $\pm0.005$ حول $0.99$. في أي مدى تقع الفواصل؟
\end{exercise}

\begin{exercise}[\lvlA]
\label{ex:21:2}
\emph{طاقة المنصة.} (أ)~ما القدرة التي يعطيها التقدير~\eqref{eq:dishpower} للدماغ كله ($8.6\cdot10^{10}$ عصبون)؟ (ب)~تسجّل المنصة $1024$ قطبًا بتردد $20\,\text{kHz}$ و$10$ بتات للعينة. كم مرة يفوق نقلُ هذا التدفق بكلفة $1\,\text{nJ}$ للبت قدرةَ المزرعة نفسها؟
\end{exercise}

\begin{exercise}[\lvlB]
\label{ex:21:3}
\emph{تصميم: إخماد الرشقات الشبكية.} تحفيز متباعد عبر أقطاب كثيرة يجعل كل مشبك يطلق الناقل بتردد $r_b$؛ والمخزون $\tau_{\text{rec}}\,\dd x/\dd t=1-x-Ur_b\tau_{\text{rec}}x$، $U=0.5$، $\tau_{\text{rec}}=0.8\,\text{s}$. ما أصغر $r_b$ لا يبلغ عندها المخزون $x^*=0.9$؛ $0.99$، وكم يضعف المشبك عندئذ؟
\end{exercise}

\begin{exercise}[\lvlB]
\label{ex:21:4}
\emph{ذاكرة الخزّان لا تزيد على عدد العصبونات.} خزّان له $N$ مخرجًا $\mathbf r(t)$ يتلقى دخلًا أبيض $u(t)$ متوسطه صفر وتباينه واحد؛ و$m_k$ أكبر مربع لارتباط القراءة $\mathbf w_k\cdot\mathbf r(t)$ مع $u(t-k)$. برهن أن سعة الذاكرة $\mathrm{MC}=\sum_{k\ge0}m_k\le N$.
\end{exercise}

\begin{exercise}[\lvlB]
\label{ex:21:5}
\emph{محاكاة: خزّان في السيليكون.} الخزّان $\mathbf r(t+1)=\tanh(W\mathbf r(t)+\mathbf v\,u(t)+\mathbf c)$، $N=30$، و$W$ عشوائية نصف قطرها الطيفي $0.9$، $v_i\sim U(-0.5,0.5)$، $c_i\sim U(-0.2,0.2)$، $u\sim U(-1,1)$، والقراءة انحدار حَرْفي. أوجد سعة ذاكرة التمرين~\ref{ex:21:4} مع $\tanh$ وبدونها. أيّ الخزّانين يتذكر أكثر؟
\end{exercise}

\begin{exercise}[\lvlB]
\label{ex:21:6}
\emph{قراءة لخزّان حي.} تعطي العضيّة $1024$ قناة و$P=240$ تسجيلًا تدريبيًا من صنفين. (أ)~كم سمة $n$ يمكن إبقاؤها كي يكون التصنيف العشوائي قابلًا للفصل، بحسب صيغة التمرين~\ref{ex:19:3}، باحتمال لا يزيد على $1\%$؟ (ب)~بكم انحرافًا معياريًا تعلو دقة $78\%$ على $100$ تسجيل اختبار فوق التخمين؟
\end{exercise}

\begin{exercise}[\lvlC]
\label{ex:21:7}
\emph{بحث: معلّم بلا قناة بثّ.} اقترح بروتوكول تحفيز يطبّق على المنصة قاعدة العوامل الثلاثة من الفصل~\ref{sec:dofa} عبر $8\text{--}1024$ قطبًا، وتجربة ضابطة بقراءة مجمَّدة تميّز تعلّم الأوزان من انزياح حالة الشبكة. أيّ نتيجة تدحض فرضية التعلم؟
\end{exercise}

\chapter{\foreignlanguage{english}{DishBrain}}
\chaptermark{\foreignlanguage{english}{DishBrain}}
\label{sec:dishbrain}

\section{شبكة في طبق تلعب التنس}

\foreignlanguage{english}{DishBrain} (دماغ في طبق) هو الاسم الذي أطلقه الباحثون على منصة تلعب فيها مزرعة من عصبونات حية، مستنبتة على شريحة بأقطاب، لعبة تنس مبسّطة بمضرب واحد~\cite{Kagan2022}. وهي أكثر التجارب على الآلات المصنوعة من عصبونات حية إثارةً للنقاش (وعرض هذه الآلات في الفصل~\ref{sec:livingmachine})، وهي مناسبة لكتابنا خاصة. فالشبكة الصغيرة هنا متاحة كلها: كل دخل يحدده المجرِّب، وكل خرج يُسجَّل. وكل أسئلة الفصول السابقة، أي كيف تُرمَّز الإشارة، ومن يعلّم، وماذا يرى المسبار، يمكن أن تُطرح على شبكة لا عيون لها ولا عضلات ولا عصبونات \nt{DOPA}. فلنحلّل التجربة كما يحلّل المهندس منصة غيره: المخطط، والدخل، والخرج، والمعلّم، والنتيجة، والخلاف.

\section{المنصة}

تؤخذ العصبونات من قشرة جنين الفأر، نحو $800\,000$ خلية للشريحة، أو تُستنبت من خلايا جذعية بشرية، نحو مليون للشريحة. وتنمو على الشريحة بضعة أسابيع، فتتشابك زوائدها وتبدأ من تلقاء نفسها بإطلاق النبضات والرشقات. وتحتها، على مساحة $8\,\text{mm}^2$، $26\,000$ قطب من البلاتين؛ يمكن التسجيل من $1024$ منها في آن واحد، أما التحفيز ففي الممارسة عبر ثمانية.

وحول المزرعة حلقة مغلقة (الشكل~\ref{fig:dishloop}). يحاكي الحاسوب اللعبة: الكرة تطير وترتد عن الجدران، والمضرب يتحرك إلى أعلى وإلى أسفل. ويتحول موضع الكرة إلى تيار على ثمانية أقطاب ”حسية“. وتُعدّ النبضات من منطقتين ”حركيتين“ في نوافذ مدة كل منها $10\ms$، فتحرّك المضرب. وبعد كل إصابة أو إخفاق تتلقى المزرعة منبهًا آخر، هو التغذية الراجعة. ونوافذ $10\ms$ هي نبضة ساعة حاسوب المنصة لا المزرعة: فالشبكة نفسها تبقى تعمل عملًا لا متزامنًا، ككل شبكات هذا الكتاب.

\begin{figure}[t]
\centering
\begin{tikzpicture}[>=Stealth,
  blk/.style={draw,very thick,rounded corners=3pt,align=center,font=\small,minimum height=1.2cm},
  lab/.style={font=\scriptsize,align=center}]
\node[blk,fill=blue!6,text width=3.4cm] (C) at (0,0) {مزرعة على شريحة\\\scriptsize $\sim10^6$ عصبون، $26\,000$ قطب};
\node[blk,fill=orange!10,text width=3.0cm] (G) at (7.2,0) {الحاسوب:\\لعبة التنس};
\draw[->,very thick,blue!60!black] (G.north west) to[bend right=25] node[lab,above] {الكرة: \emph{المكان}: أيّ الأقطاب الثمانية،\\\emph{التردد}: $4\text{--}40$ نبضة/ث} (C.north east);
\draw[->,very thick,red!70!black] (C.south east) to[bend right=25] node[lab,below] {المضرب: عدّ نبضات المنطقتين\\”أعلى“ و”أسفل“ في $10\ms$} (G.south west);
\draw[->,thick,densely dashed,green!45!black] (G.east) -- ++(0.9,0) |- ++(0,-2.6) -| node[lab,pos=0.25,below] {إصابة: منبه متوقَّع؛ إخفاق: منبه غير متوقَّع} (C.south);
\end{tikzpicture}
\caption{حلقة \foreignlanguage{english}{DishBrain}. السهم الأزرق هو الدخل: موضع الكرة مرمَّز بترميز المكان وبترميز التردد (الفصل~\ref{sec:code}). والأحمر هو الخرج: الفرق بين عدّي المنطقتين يحرّك المضرب. والأخضر المتقطع هو التغذية الراجعة بعد كل جولة، وهي ”المعلّم“ الوحيد في المنصة.}
\label{fig:dishloop}
\end{figure}

\section{الدخل والخرج}

\emph{الدخل} مرمَّز بترميزين من الفصل~\ref{sec:code} معًا. أيّ الأقطاب الثمانية يُحفَّز يخبر على أي ارتفاع توجد الكرة: هذا ترميز المكان، كما في مستشعرات الفصل~\ref{sec:sensors}. وكم مرة تأتي نبضات التحفيز يخبر كم تبعد الكرة: $4$ في الثانية حين تكون عند الجدار البعيد، وأكثر خطيًا حتى $40$ في الثانية حين تكون عند جدار المضرب. وهذا ترميز التردد؛ وسعة نبضات الكرة $75\mV$. والمزرعة لا ”تعرف“ أيًّا من هذين الترميزين مسبقًا: الترميز حدده المجرِّب، وعلى الشبكة أن تتعلم قراءته.

\emph{الخرج} يُقرأ كما في واجهات الفصل~\ref{sec:debugging}. نبضات إحدى المنطقتين الحركيتين تدفع المضرب إلى أعلى، ونبضات الأخرى إلى أسفل؛ وفي أبسط صورة، في كل نافذة
\begin{empheq}[box=\keybox]{equation}
\Delta p \;=\; c\,\bigl(n_\uparrow - n_\downarrow\bigr),
\label{eq:dishreadout}
\end{empheq}
حيث $n_\uparrow$ و$n_\downarrow$ عددا نبضات المنطقتين في $10\ms$، و$c$ معامل المنصة. هذا هو الترميز بسلكين من الفصل~\ref{sec:glutcomputer}: اتجاه الحركة لا تحمله علامة الإشارة، بل أيّ السلكين أعلى صوتًا.

لننظر في ضجيج هذا الخرج. إذا أطلقت المنطقة ما مجموعه $R$ نبضة في الثانية، ففي النافذة $T=10\ms$ يكون متوسطها $RT$، وبحسب~\eqref{eq:rateerror} يكون التشتت النسبي $1/\sqrt{RT}$. عند $R=1000$ نبضة في الثانية يبلغ هذا نحو $30$ في المئة في كل نافذة؛ وعند $R=100$ أكثر من مئة. فالمضرب يرتجف حتمًا، ولا تستطيع الشبكة أن تتعلم إلا بطريقة واحدة: أن تزيح \emph{متوسط} الفرق بين العدّين في الاتجاه المطلوب. إنها القوانين نفسها التي تحدّ كل واجهة في الفصل~\ref{sec:debugging}، لكنها الآن على طرفي الحلقة كليهما.

\section{معلّم بلا دوبامين}

أطرف ما في المنصة معلّمها. فمزرعة عصبونات القشرة لا عصبونات \nt{DOPA} فيها، والمنصة لا ترسل أي إشارة من نوع ”أفضل مما توقعنا“. التغذية الراجعة من نوع آخر. إذا صدّت الشبكة الكرة، قُدّم إلى الأقطاب الحسية الثمانية كلها معًا منبه قصير \emph{متوقَّع}: نبضات بسعة $75\mV$، $100$ في الثانية مدة $0.1\,\text{s}$. وإذا أخفقت، استمر أربع ثوانٍ منبه يسميه الباحثون \emph{غير متوقَّع}: نبضات بسعة $150\mV$، $5$ في الثانية، في لحظات عشوائية وعلى أقطاب مختارة عشوائيًا؛ ثم أربع ثوانٍ من التوقف بلا تحفيز، وتنطلق الكرة من جديد~\cite{Kagan2022}.

انطلق الباحثون من فرضية أن الشبكة تتصرف بحيث يصير دخلها متوقَّعًا~\cite{Friston2010}. والمعنى عند المهندس بسيط: ليس لدى المزرعة إلا طريقة واحدة للتخلص من أربع ثوانٍ من الضجيج العشوائي، وهي أن تصدّ الكرة. وقد صادفنا الفكرة نفسها في الاقتصاد في النبضات في الفصل~\ref{sec:code} وفي التنبؤ بالإطار في الفصل~\ref{sec:recognition}.

لكن هل يلزم هنا التوقُّع بعينه؟ هناك آلية أخشن ممكنة أيضًا. لنفترض أن المنبه غير المتوقَّع بعد الإخفاق \emph{يخلط} فحسب التغيرات الأخيرة في الشبكة، وأن المنبه الهادئ بعد الإصابة لا يمسّها. لنصف ضبط الشبكة بمتجه واحد $\boldsymbol\psi$، ولتكن احتمالية إخفاق الشبكة في الضبط $\boldsymbol\psi$ هي $P_{\text{miss}}(\boldsymbol\psi)$. إذا كانت الدفعات متناظرة، أي أن الانتقال من $\boldsymbol\psi$ إلى $\boldsymbol\psi'$ محتمل احتمال الانتقال المعاكس، ففي النظام المستقر يتوازن تدفق المغادرة من كل ضبط مع تدفق القدوم إليه، ونسبة الوقت الذي تقضيه الشبكة في الضبط $\boldsymbol\psi$ هي
\begin{empheq}[box=\keybox]{equation}
\rho(\boldsymbol\psi) \;\propto\; \frac{1}{P_{\text{miss}}(\boldsymbol\psi)} .
\label{eq:dishdensity}
\end{empheq}
الضبط الذي يخفق أقل بعشر مرات يدوم أطول بعشر مرات. لا أحد يخبر الشبكة في أي اتجاه تغيّر أوزانها: الضبوط الجيدة تُهدم أقل فحسب.

اختبرنا ذلك على نموذج يختزل المزرعة في عددين: يصل المضرب إلى النقطة $a\,y+b$ حين تصل الكرة إلى الارتفاع $y$، ويصدّها إذا كان خطؤه أقل من $0.1$. بعد الإخفاق يتلقى العددان دفعة عشوائية، وبعد الإصابة يبقيان على حالهما. فارتفعت نسبة الكرات المصدودة في $2000$ جولة من $0.21$ في أول مئتين إلى $0.38$ في آخر خمسمئة (أربعون ”مزرعة“، الشكل~\ref{fig:dishmodel}). وإذا أُعطيت الدفعة بعد \emph{كل} جولة، فلا معلومة في التغذية الراجعة، وتبقى النسبة نحو $0.12$؛ وإذا لم تكن دفعات أصلًا ثبتت عند $0.19$. والتجربة على المنصة تتفق مع هذا: لم يظهر نمو ذو دلالة داخل الجلسة إلا مع التغذية الراجعة الكاملة. ومع الصمت بدل المنبهات لعبت المزرعة في آخر الجلسة مع ذلك أفضل مما بلا تغذية راجعة، لكن بأثر أصغر~\cite{Kagan2022}؛ ونلاحظ أن في هذا الشرط أيضًا يأتي بعد الإخفاق توقف طويل وبعد الإصابة توقف قصير، فلا يختفي الفرق بين النتيجتين كليًا.

\begin{figure}[t]
\centering
\begin{tikzpicture}[x=1cm,y=5cm]
\draw[->,gray!70!black] (0,0) -- (10.6,0);
\draw[->,gray!70!black] (0,0) -- (0,0.5) node[above,font=\small,black] {نسبة المصدود};
\foreach \v in {0.1,0.2,0.3,0.4}{\draw[gray!60!black] (-0.06,\v) -- (0.06,\v); \node[left,font=\scriptsize] at (0,\v) {\v};}
\foreach \x/\f/\l/\lab in {1.8/0.21/0.38/{ضجيج بعد\\الإخفاق},5.2/0.13/0.11/{ضجيج بعد\\كل جولة},8.6/0.19/0.19/{بلا ضجيج}}{
  \fill[blue!35] (\x-0.8,0) rectangle (\x-0.05,\f);
  \fill[red!60!black] (\x+0.05,0) rectangle (\x+0.8,\l);
  \node[below,font=\scriptsize,align=center] at (\x,-0.01) {\lab};
  \node[above,font=\scriptsize] at (\x-0.425,\f) {\f};
  \node[above,font=\scriptsize] at (\x+0.425,\l) {\l};}
\fill[blue!35] (7.4,0.47) rectangle (7.7,0.495); \node[right,font=\scriptsize] at (7.75,0.482) {أول 200};
\fill[red!60!black] (7.4,0.41) rectangle (7.7,0.435); \node[right,font=\scriptsize] at (7.75,0.422) {آخر 500};
\end{tikzpicture}
\caption{نموذج ”الخلط عند الإخفاق“ (\foreignlanguage{english}{\texttt{dishbrain\_model.py}}): نسبة الكرات المصدودة في بداية $2000$ جولة وفي نهايتها، بالمتوسط على أربعين مزرعة نموذجية. لا يحدث التعلم إلا إذا جاءت الدفعة بعد الإخفاق لا بعد الإصابة، كما في التجربة، حيث لم يظهر نمو ذو دلالة داخل الجلسة إلا مع التغذية الراجعة الكاملة.}
\label{fig:dishmodel}
\end{figure}

\begin{keyblock}
\begin{proposition}[المعلّم الخالط]
\label{prop:dishscramble}
إذا كان الفشل يهزّ الشبكة والنجاح يتركها وشأنها، انجرفت الشبكة من تلقاء نفسها نحو الضبوط التي تخطئ أقل: الوقت الذي تقضيه في ضبط ما يتناسب عكسيًا مع احتمالية الفشل~\eqref{eq:dishdensity}. ولا يحتاج هذا التعلم إلى إشارة مكافأة ولا إلى علامتها ولا إلى التوقُّع بحد ذاته؛ يكفي أن تختلف التغذية الراجعة بعد النجاح عنها بعد الفشل. تغيّر سلوك المزرعة مُقاس؛ أما الآلية العاملة في الطبق، أهي التنبؤ بالدخل أم الخلط، فلم تُحدَّد.
\end{proposition}
\end{keyblock}

قارن بالفصل~\ref{sec:dofa}. هناك أيضًا كان الضجيج أداة بحث، لكن الدوبامين كانت له علامة: الخطأ $\delta$ كان يخبر في أي اتجاه تُزاح الأوزان، وكانت الخطوة تسير على التدرج~\eqref{eq:perturbation}. أما هنا فلا علامة: ليس إلا ”أبقِ“ أو ”اهزز“. وهذا البحث أخشن وأبطأ، لكنه يطلب من الشبكة أقل: لا عصبونات \nt{DOPA}، ولا ناقد، ولا أثر أهلية ممتد على ثوانٍ.

\section{ما قيس وما الخلاف عليه}

دامت كل لعبة $20$ دقيقة؛ وقورنت أول $5$ دقائق بآخر $15$. في المزارع ذات التغذية الراجعة الكاملة طالت سلاسل الكرات المصدودة وقلّت الإخفاقات الفورية؛ ولم يُرَ شيء من هذا في التجارب الضابطة بلا خلايا، وبلا لعبة، وبـ”مضرب“ يحرّكه الحاسوب. ومزارع الخلايا البشرية صدّت في المتوسط مددًا أطول من مزارع الفأر. والتحسن متين إحصائيًا، على تسع مزارع من الفأر وإحدى عشرة بشرية، في أكثر من مئة لعبة لكل مجموعة، لكنه صغير: لم تصر الشبكة لاعبًا جيدًا، بل صارت أفضل قليلًا من العشوائية. وهو متين داخل الجلسة؛ أما التعلم الثابت من جلسة إلى أخرى على مدى أيام فلم يحصل عليه الباحثون~\cite{Kagan2022}. ووجدت الدراسة التالية للمجموعة نفسها أن فعالية المزرعة في أثناء اللعب تقترب من النظام \emph{الحرج}، أي الحد بين الخمود والانهيار المتسلسل، وهو العيش على الحافة نفسه الذي في الفصل~\ref{sec:gaba}، وكلما اقتربت كان اللعب أفضل~\cite{Habibollahi2023}.

الخلاف الرئيسي لم تُثِره الأرقام، بل الكلمات. فعنوان المقالة زعم أن العصبونات في الطبق ”تُظهر وعيًا“ (\foreignlanguage{english}{sentience})، واعترضت مجموعة من الباحثين: تغيّر السلوك في حلقة مغلقة ليس ذكاءً ولا إحساسًا بعد، والاستنتاجات يجب أن تُصاغ بتواضع أكبر~\cite{Balci2023}. ومن الزاوية الهندسية سؤالان مفتوحان، وكلاهما عن الآلية. الأول: هل تتعلم الشبكة، أي هل تتغير أوزانها تغيرًا دائمًا، أم أن التغذية الراجعة تزيح حالتها الراهنة فحسب؟ والثاني: هل يلزم التوقُّع بعينه، أم يكفي أي فرق بين الاستجابة للنجاح والاستجابة للفشل، كما في القضية~\ref{prop:dishscramble}؟ والمنصة التي يُتاح فيها كل قطب تسمح بالإجابة عن السؤالين معًا، ولعل هذه أهم قيمها.

\section{مواصفات المنصة: كيف نعيد بناءها}
\label{sec:dishspec}

جمعنا أدناه كل ما يلزم لبناء منصة كهذه من جديد: العتاد، والحلقة، وترميز الدخل، وقراءة الخرج، والتغذية الراجعة، وبروتوكول التجربة. وكل معامل موسوم بمصدره. ”المقالة“: القيمة مأخوذة من نص الدراسة~\cite{Kagan2022}. ”اختيار“: القيمة غير موجودة في المقالة، وعلى من يعيد التجربة أن يحددها بنفسه؛ ونقترح قيمة افتراضية ونبيّن كيف تُختبر.

\subsection*{العتاد والمزرعة}

\begin{center}
\footnotesize
\setlength{\tabcolsep}{4pt}
\renewcommand{\arraystretch}{1.15}
\begin{tabular}{>{\raggedright}p{4.0cm}>{\raggedright}p{6.9cm}>{\raggedright\arraybackslash}p{1.6cm}}
\toprule
المعامل & القيمة & المصدر \\
\midrule
الشريحة & مصفوفة \foreignlanguage{english}{MaxOne}: $26\,000$ قطب من البلاتين على مساحة $8\,\text{mm}^2$ & المقالة \\
التسجيل & حتى $1024$ قناة في آن واحد، $20\,000$ عينة في الثانية & المقالة \\
التحفيز & حتى $32$ قطبًا تقنيًا، وفي التجربة $8$ مستقلة & المقالة \\
الخلايا & قشرة جنين الفأر؛ عصبونات بشرية من خلايا جذعية (طريقتا استنبات)؛ خلايا \foreignlanguage{english}{HEK293T} (ليست عصبونات) للمقارنة الضابطة & المقالة \\
الزرع & نحو $10^6$ خلية للشريحة & المقالة \\
عمر المزرعة عند التجربة & الفأر: من $\sim10$ أيام؛ الإنسان: $14\text{--}24$ يومًا أو $\sim80$ يومًا بحسب طريقة الاستنبات & المقالة \\
\bottomrule
\end{tabular}
\end{center}

\subsection*{الحلقة والزمن}

تعمل الحلقة بنوافذ مدة كل منها $10\ms$ ($200$ عينة): في كل نافذة تُعدّ النبضات على أقطاب المنطقتين الحركيتين، وتُحدَّث اللعبة، ويُصدر التحفيز التالي. والتأخير من النبضة في المزرعة إلى المنبه المجيب نحو $5\ms$. وفي أثناء اللعب تمرّ إشارة كل قطب بمرشح بِسل تمرير عالٍ من الرتبة الثانية بتردد قطع $100\,\text{Hz}$؛ وتُنعَّم القيمة المطلقة للإشارة بمرشح بِسل تمرير منخفض من الرتبة الأولى بتردد $1\,\text{Hz}$، وعتبة النبضة متناسبة مع هذه القيمة المنعَّمة. أما عتبة ستة انحرافات معيارية للخلفية فتذكرها المقالة لقياسات منفصلة للفعالية، لا للعب. وكي لا يُحسب التحفيز استجابةً من المزرعة، تُخمد كل الإشارات حين يصل في آن واحد أكثر من $15$ نبضة كبيرة، تزيد على $75\mV$. كل هذا من المقالة.

\subsection*{ترميز الكرة}

\begin{center}
\footnotesize
\setlength{\tabcolsep}{4pt}
\renewcommand{\arraystretch}{1.15}
\begin{tabular}{>{\raggedright}p{3.4cm}>{\raggedright}p{7.5cm}>{\raggedright\arraybackslash}p{1.6cm}}
\toprule
المعامل & القيمة & المصدر \\
\midrule
المنطقة الحسية & $8$ أقطاب مرتبة بحيث تعني الأقطاب المتجاورة مواضع متجاورة للكرة & المقالة \\
ترميز المكان & رقم القطب $k=1,\dots,8$ يحدد موضع الكرة بالنسبة إلى مركز المضرب & المقالة \\
أي إحداثي & الموضع الرأسي للكرة؛ وللإعادة: بالنسبة إلى المضرب، $k=\lceil 8(y-p+\tfrac12)\rceil$ عند $y-p\in[-\tfrac12,\tfrac12]$ & المقالة؛ والصيغة اختيار \\
ترميز التردد & تردد التحفيز من $4$ إلى $40$ نبضة/ث يحمل المسافة إلى المضرب & المقالة \\
اتجاه السُّلَّم & $4$ نبضات/ث حين تكون الكرة عند الجدار المقابل، وخطيًا حتى $40$ نبضة/ث حين تكون عند جدار المضرب: $f=4+36\,(1-x/L)$، حيث $x$ المسافة إلى جدار المضرب و$L$ عرض الملعب & المقالة \\
شكل النبضة & مستطيلة قصيرة ثنائية الطور: الطور الموجب أولًا، ثم السالب & المقالة \\
السعة & $75\mV$؛ اختيرت كي لا تُجبر الخلايا المثبَّطة على الإطلاق & المقالة \\
مدة الطورين & غير مذكورة؛ تؤخذ الإعدادات القياسية لنظام التحفيز ولا تُغيَّر في أثناء التجربة & اختيار \\
\bottomrule
\end{tabular}
\end{center}

\subsection*{قراءة المضرب}

في المقالة منطقتان حركيتان: النبضات في الأولى تحرّك المضرب إلى أعلى، وفي الثانية إلى أسفل؛ وتُوزن مساهمة المنطقتين لمراعاة اختلاف كثافة الخلايا واختلاف الفعالية~\cite{Kagan2022}. والصيغة الدقيقة غير مذكورة. وللإعادة نأخذ القاعدة~\eqref{eq:dishreadout}، $\Delta p=c\,(n_\uparrow-n_\downarrow)$ لكل نافذة $10\ms$، بوزن يسوّي بين المنطقتين بحسب متوسط فعاليتهما في الراحة (اختيار): يُقسم $n_\uparrow$ و$n_\downarrow$ على متوسط عدّ منطقتيهما في النافذة، مقيسًا قبل اللعب. ونختار المعامل $c$ بحيث يزيح فرقٌ بمقدار عدّ متوسط واحد المضربَ بنسبة $1\%$ من ارتفاع الملعب في النافذة (اختيار)؛ وعندئذ يقطع المضرب الملعب كله في $1\,\text{s}$ إذا رجحت إحدى المنطقتين رجحانًا ثابتًا.

ومواقع المنطقتين على الشريحة لم تُحدَّد في نص المقالة بإحداثيات؛ ليس هناك إلا مخطط. وللإعادة تكفي ثلاث مجموعات غير متداخلة من الأقطاب: ثمانية أقطاب حسية في الوسط، ومجموعة أقطاب تسجيل على كل جانب، إحداهما ”أعلى“ والأخرى ”أسفل“ (اختيار).

\subsection*{اللعبة}

أبعاد الملعب وطول المضرب وسرعة الكرة غير مذكورة في المقالة؛ لا يُذكر إلا أن الكرة تطير بسرعة ثابتة وترتد عن الجدران. وللإعادة (اختيار): ملعب $1\times1$، ومضرب على الجدار الأيمن طوله $0.2$ (إصابة عند $|y-p|<0.1$، كما في النموذج~\foreignlanguage{english}{\texttt{models/dishbrain\_model.py}})، والكرة تقطع الملعب في $1\,\text{s}$. والإخفاق من غير أي كرة مصدودة في الجولة يُعدّ ”إرسالًا ساحقًا“ (\foreignlanguage{english}{ace})؛ والجولة التي تزيد على ثلاث إصابات متتالية تُعدّ ”طويلة“ (المقالة).

\subsection*{التغذية الراجعة}

\begin{center}
\footnotesize
\setlength{\tabcolsep}{4pt}
\renewcommand{\arraystretch}{1.15}
\begin{tabular}{>{\raggedright}p{2.6cm}>{\raggedright}p{8.3cm}>{\raggedright\arraybackslash}p{1.6cm}}
\toprule
الحدث & ما يُقدَّم & المصدر \\
\midrule
إصابة & منبه متوقَّع: الأقطاب الحسية الثمانية كلها في آن واحد، $75\mV$، $100$ نبضة/ث، $100\ms$؛ ويُقطع تحفيز الكرة المعتاد خلال ذلك & المقالة \\
إخفاق & منبه غير متوقَّع: $150\mV$، $5$ نبضات/ث، $4\,\text{s}$، في لحظات عشوائية وعلى أقطاب مختارة عشوائيًا من الثمانية؛ ثم $4\,\text{s}$ بلا تحفيز، وتنطلق الكرة في اتجاه عشوائي & المقالة \\
قانون العشوائية & غير مذكور؛ وللإعادة: لحظات النبضات تدفق بواسوني بمتوسط تردد $5$ نبضات/ث & اختيار \\
\bottomrule
\end{tabular}
\end{center}

\subsection*{البروتوكول والتجارب الضابطة}

تدوم الجلسة الواحدة $20$ دقيقة؛ وتُقارن أول $5$ دقائق بآخر $15$ (المقالة). ومقاييس النتيجة ثلاثة: متوسط طول الجولة، ونسبة الإرسالات الساحقة، ونسبة الجولات الطويلة. وشروط التجربة (المقالة):
\begin{itemize}
\item \emph{المنبه}: الحلقة الكاملة، كما وُصفت أعلاه؛
\item \emph{الصمت}: بدل منبه الإصابة أو الإخفاق المدة نفسها بلا أي تحفيز، ثم تنطلق الكرة في اتجاه عشوائي؛
\item \emph{بلا تغذية راجعة}: بعد الإخفاق ترتد الكرة فحسب وتواصل طيرانها، ويستمر تحفيز موضع الكرة؛
\item \emph{الراحة}: اللعبة تجري والمضرب يتحرك بحسب الفعالية، لكن لا تحفيز أصلًا؛
\item \emph{الضوابط}: شريحة فيها وسط الاستنبات وحده؛ خلايا \foreignlanguage{english}{HEK293T}؛ ”مزرعة في الحاسوب“، أي محاكاة بدل الخلايا.
\end{itemize}
حجم العينة في السلسلة الرئيسية: $9$ مزارع من الفأر ($101$ جلسة)، و$11$ مزرعة بشرية ($138$)، و$6$ شرائح بوسط الاستنبات ($80$)، و$20$ مزرعة في الراحة ($42$)، و$3$ محاكيات ($38$)، والمجموع $399$ جلسة. وفي سلسلة التغذية الراجعة $486$ جلسة على مدى $3$ أيام بمعدل $3$ جلسات يوميًا، بالشروط ”المنبه“ و”الصمت“ و”بلا تغذية راجعة“ ($n=27$ و$17$ و$15$). ولم تظهر زيادة متوسط طول الجولة من أول $5$ دقائق إلى آخر $15$ إلا في المزارع الحية. وفي سلسلة التغذية الراجعة لم يظهر نمو ذو دلالة مع الزمن إلا في شرط ”المنبه“؛ وفي آخر الجلسة تفوّق شرط ”الصمت“ مع ذلك على ”بلا تغذية راجعة“، لكن بأثر أصغر، ولم يكن هناك تعلّم ثابت من جلسة إلى أخرى على مدى ثلاثة أيام~\cite{Kagan2022}.

\subsection*{دورة المنصة خطوةً خطوة}

كل $10\ms$ يفعل حاسوب المنصة ما يلي.
\begin{enumerate}
\item يقرأ عدد النبضات $n_\uparrow$ و$n_\downarrow$ في المنطقتين الحركيتين خلال النافذة المنقضية، ويقسمهما على متوسط عدّ المنطقة في الراحة.
\item يزيح المضرب بمقدار $\Delta p=c\,(n_\uparrow-n_\downarrow)$ ويحصره بين حدّي الملعب.
\item يحرّك الكرة خطوة واحدة؛ ويعكسها عن الجدران العلوي والسفلي والأيسر.
\item إذا بلغت الكرة الجدار الأيمن: عند $|y-p|<0.1$ فهي إصابة، ويزيد عدّاد الجولة واحدًا، ويُقدَّم منبه الإصابة ($100\ms$)؛ وإلا فهي إخفاق، وتُسجَّل الجولة، ويُقدَّم منبه الإخفاق ($4\,\text{s}$)، ثم توقف $4\,\text{s}$، وتنطلق الكرة من جديد.
\item وإلا يختار القطب $k$ بحسب الموضع الرأسي للكرة والتردد $f$ بحسب المسافة إلى جدار المضرب، ويقدّم إلى القطب $k$ نبضات ثنائية الطور بسعة $75\mV$ بالتردد $f$ حتى النافذة التالية.
\end{enumerate}
وهذا يكفي لبناء منصة متوافقة مع المنشورة، ولاختبار السؤال الرئيسي في الفصل: هل يعلّم المزرعةَ التوقُّعُ، أم مجرد الفرق بين الاستجابة للإصابة والاستجابة للإخفاق؟ ولهذا يحسن أن يُضاف إلى شروط المقالة الثلاثة شرط رابع ليس فيها: بعد الإخفاق منبه \emph{متوقَّع} بالمدة والطاقة نفسيهما اللتين للمنبه غير المتوقَّع. فإن تعلمت المزرعة مع ذلك، فالأمر في الفرق لا في التوقُّع.


\section{خلاصة}

\begin{table}[t]
\centering
\footnotesize
\setlength{\tabcolsep}{4pt}
\renewcommand{\arraystretch}{1.15}
\begin{tabular}{>{\raggedright}p{2.1cm}>{\raggedright}p{4.2cm}>{\raggedright}p{1.9cm}>{\raggedright\arraybackslash}p{4.6cm}}
\toprule
عقدة المنصة & كيف بُنيت & فصل الكتاب & ما هو معروف \\
\midrule
الدخل & المكان (أيّ الأقطاب الثمانية) والتردد ($4\text{--}40$ نبضة/ث) & \ref{sec:code}، \ref{sec:sensors} & يحدده المجرِّب \\
الخرج & الفرق بين عدّي المنطقتين في $10\ms$~\eqref{eq:dishreadout} & \ref{sec:glutcomputer}، \ref{sec:debugging} & يحدده المجرِّب؛ والضجيج وفق~\eqref{eq:rateerror} \\
المعلّم & منبه متوقَّع بعد الإصابة، وغير متوقَّع بعد الإخفاق؛ ولا دوبامين & \ref{sec:dofa} & نمو ذو دلالة داخل الجلسة مع التغذية الراجعة الكاملة فقط؛ ومع الصمت أثر أصغر؛ وغير ثابت من جلسة إلى أخرى \\
الآلية & التنبؤ بالدخل أو الخلط عند الإخفاق~\eqref{eq:dishdensity} & \ref{sec:dofa}، \ref{sec:recognition} & لم تُحدَّد؛ وكلتاهما تفسّر النتيجة \\
نظام الشبكة & الاقتراب من النظام الحرج في أثناء اللعب & \ref{sec:gaba} & صلته بجودة اللعب مقيسة \\
”الوعي“ & قول الباحثين & --- & لم يُبيَّن؛ ومطعون فيه \\
\bottomrule
\end{tabular}
\caption{\foreignlanguage{english}{DishBrain} بعين المهندس.}
\label{tab:dishbrain}
\end{table}

\foreignlanguage{english}{DishBrain} آلة من عصبونات حية في أبسط صورها: الدخل مرمَّز بترميز، والخرج مقروء بالعدّ، والمعلّم فرق بين الاستجابة للنجاح والاستجابة للفشل (الجدول~\ref{tab:dishbrain}). شبكة بلا عيون ولا عضلات ولا دوبامين تعلّمت اللعب قليلًا جدًا، وهذا وحده يجعلها منصة اختبار لأفكار الكتاب. وأيًّا كانت الآلية، تنبؤًا أو خلطًا أو شيئًا ثالثًا، فسيوجد الجواب كما ينصح الفصل~\ref{sec:debugging}: بالمسبار، وبإشارة الاختبار، وبتعطيل الأجزاء، على آلة نراها أخيرًا كاملة.

\section{تمارين}

\begin{exercise}[\lvlA]
\label{ex:22:1}
\emph{كم يلزم أن يُزاح العدّ.} تطلق المنطقتان معًا $R_\uparrow+R_\downarrow=2000$ نبضة في الثانية، والتدفقات بواسونية. (أ)~ما التشتت النسبي لعدّ منطقة واحدة في $10\ms$ عند $R=1000$ و$R=100$؟ (ب)~في $1\,\text{s}$ من طيران الكرة تتجمع الإزاحات~\eqref{eq:dishreadout}. ما أصغر $R_\uparrow-R_\downarrow$ يكون عندها متوسط الإزاحة ضعف التشتت؟
\end{exercise}

\begin{exercise}[\lvlA]
\label{ex:22:2}
\emph{كم جولة يلزم لرؤية التعلم.} الجولات مستقلة، ونسبة المصدود ثنائية الحدين. كم جولة يلزم في كل من الفترتين كي تتجاوز زيادة نسبة المصدود انحرافين معياريين للفرق: (أ)~من $0.20$ إلى $0.25$؛ (ب)~من $0.21$ إلى $0.38$، كما في النموذج؟
\end{exercise}

\begin{exercise}[\lvlB]
\label{ex:22:3}
\emph{اشتقاق~\eqref{eq:dishdensity}.} الضبط $\boldsymbol\psi$ يجري على مجموعة منتهية؛ عند الإخفاق (باحتمالية $P(\boldsymbol\psi)>0$) تنتقل الشبكة إلى $\boldsymbol\psi'$ باحتمالية متناظرة $K(\boldsymbol\psi,\boldsymbol\psi')$، وعند الإصابة تبقى. (أ)~برهن أن $\rho\propto1/P$ توزيع مستقر. (ب)~ما نسبة الإخفاقات مع ضبطين فيهما $P=0.9$ و$P=0.1$؟
\end{exercise}

\begin{exercise}[\lvlB]
\label{ex:22:4}
\emph{المنطقة الماصّة في النموذج.} يصل المضرب إلى $p=\min\bigl(1,\max(0,ay+b)\bigr)$، والكرة إلى الارتفاع $y\sim U(0,1)$، والإصابة إذا $|p-y|<h=0.1$. (أ)~برهن أن الشبكة تصدّ كل الكرات عند $|b|<h$ و$|a-1+b|<h$، وأوجد مساحة هذه المنطقة. (ب)~أوجد عدديًا متوسط نسبة المصدود عند $a\sim U(-1,1)$ و$b\sim U(0,1)$.
\end{exercise}

\begin{exercise}[\lvlB]
\label{ex:22:5}
\emph{محاكاة: سياج حول الضبوط.} في نسخة من \foreignlanguage{english}{\texttt{dishbrain\_model.py}} أعطِ $200$ مزرعة $20\,000$ جولة لكل منها. ما نسبة المصدود في آخر $500$؟ وماذا لو قُصّ الضبط بعد كل دفعة إلى المستطيل $a\in[-1,2]$، $b\in[-1,1]$؟ فسّر الفرق بالتمرين~\ref{ex:22:3}.
\end{exercise}

\begin{exercise}[\lvlB]
\label{ex:22:6}
\emph{تصميم ترميز الدخل.} تُصدّ الكرة إذا كان خطأ الارتفاع أقل من $h=0.1$؛ تقرأ الشبكة الدخل في $T=0.25\,\text{s}$، وحتى التردد $r$ يمكن تمييز نحو $\sqrt{rT}$ مستوى، والتحفيز لا يتجاوز $40$ نبضة في الثانية. (أ)~هل يكفي ترميز المكان على ثمانية أقطاب؟ (ب)~هل يمكن نقل الارتفاع بالتردد على قطب واحد؟
\end{exercise}

\begin{exercise}[\lvlC]
\label{ex:22:7}
\emph{بحث: تنبؤ أم خلط؟} عمّم النموذج على $d$ عددًا قابلًا للضبط ($p=\sum_{i<d}a_iy^i$). كيف يزداد مع $d$ عدد الجولات حتى نسبة مصدود $0.5$ عند الخلط وعند القاعدة ذات علامة الخطأ~\eqref{eq:perturbation}؟ وأيّ تجربة على المنصة تفصل بين الفرضيتين؟
\end{exercise}

\appendix
\renewcommand{\chaptername}{\appendixname}
\chapter{مواد لا تحدد نوع العصبون}
\label{sec:othermediators}

في الفصل~\ref{sec:neurontypes} فرزنا العصبونات بحسب ناقلها العصبي الرئيسي وحصلنا على عشرة أنواع (الجدول~\ref{tab:types}). وقد حان الوقت لإضافة التعقيدات. فالعصبون لا يقذف الناقل الرئيسي وحده، وفي الدماغ تعمل إلى جانب النواقل الرئيسية مواد أخرى. وهي تغيّر طريقة الشبكة في نقل الإشارات وتخزينها.

وإلى جانب ناقل العناوين وناقل البثّ (القسم~\ref{sec:buses}) هناك ناقل ثالث: \emph{القناة العكسية}. فبعض المواد لا يقذفها المرسل، بل \emph{المستقبِل}، فتعبر الشق راجعة إلى مخرج المرسل، وتغيّر كمية الناقل التي سيقذفها في المرة التالية. وهذا يشبه في شبكات الاتصال إشعار الاستلام الذي يستعمله المرسل لضبط إرساله. وبهذه القناة بالذات تستعين العصبونات حين تغيّر أوزان روابطها: فعبرها يعرف مخرج المرسل فعالية المستقبِل، أي العامل الثاني في قواعد التعلم في الفصل~\ref{sec:dofa}.

القائمة الكاملة للنواقل المعروفة طويلة: أكثر من مئة مادة، أغلبها سلاسل بروتينية قصيرة هي الببتيدات العصبية. لكنها كلها تندرج في بضعة أصناف، والمواد داخل كل صنف تتصرف تصرفًا متشابهًا. وقد جُمعت المواد التي لا تكون ناقلًا رئيسيًا في الجدول~\ref{tab:othermediators}، مع الأسئلة الثلاثة نفسها التي يطرحها المهندس: على أي ناقل تسير الإشارة، وهل تعمل بسرعة، وما علامتها المعتادة. وهي لا تحدد نوع العصبون: فإما أن تُقذف مع الناقل الرئيسي، وإما أن تقذفها خلايا غير عصبونية، وإما أن تسير في الاتجاه المعاكس.

\begin{table}[t]
\centering
\footnotesize
\setlength{\tabcolsep}{4pt}
\renewcommand{\arraystretch}{1.12}
\begin{tabular}{>{\raggedright}p{2.25cm}>{\raggedright}p{2.8cm}>{\raggedright}p{1.8cm}>{\raggedright}p{1.5cm}>{\centering}p{0.8cm}>{\raggedright\arraybackslash}p{4.3cm}}
\toprule
الصنف & المادة & الناقل & السرعة & العلامة & الدور في الحساب \\
\midrule
الأحماض الأمينية & دي-سيرين & موضعي & بطيء & و & مفتاح ثانٍ لمستقبل الغلوتامات: شرط ”و“ \\
\midrule
الأمينات الأحادية & الأمينات الأثرية & بثّ & بطيء & $\pm$ & ضبط دقيق لقنوات الأمينات الأحادية \\
\midrule
\multirow{2}{2.25cm}{البيورينات}
 & \foreignlanguage{english}{ATP} & عنواني & سريع وبطيء & $+$ & ناقل مرافق؛ إشارة بين العصبونات والدِّبق \\
 & الأدينوسين & حجمي & بطيء & $-$ & يتراكم مع الفعالية: عدّاد للحِمل~\cite{Dunwiddie2001} \\
\midrule
\multirow{8}{2.25cm}{الببتيدات العصبية (أكثر من مئة)}
 & الإنكيفالينات، الإندورفينات، الداينورفين & حجمي & بطيء & $-$ & كبت النقل \\
 & المادة \foreignlanguage{english}{P} & حجمي & بطيء & $+$ & تعزيز بطيء \\
 & الببتيد العصبي \foreignlanguage{english}{Y} & حجمي & بطيء & $-$ & تثبيط بطيء \\
 & السوماتوستاتين & حجمي & بطيء & $-$ & تثبيط بطيء، يرافق الغابا \\
 & الببتيد المعوي الفعال في الأوعية (\foreignlanguage{english}{VIP}) & حجمي & بطيء & $+$ & رفع التثبيط: تثبيط العصبونات المثبِّطة \\
 & الكوليسيستوكينين & حجمي & بطيء & $\pm$ & يرافق الغابا في بعض العصبونات المثبِّطة \\
 & الأوركسين & بثّ & بطيء & $+$ & مثبِّت لنظام اليقظة \\
 & الأوكسيتوسين، الفازوبريسين & بثّ & بطيء & $\pm$ & ضبط شامل للسلوك \\
\midrule
\multirow{2}{2.25cm}{الغازات}
 & أكسيد النيتريك \foreignlanguage{english}{NO} & عكسي، حجمي & بطيء & $\pm$ & يعبر الأغشية؛ إشارة عكسية عند تغيّر الأوزان~\cite{Garthwaite2008} \\
 & \foreignlanguage{english}{CO}، \foreignlanguage{english}{H$_2$S} & حجمي & بطيء & $\pm$ & دوره قليل الدراسة \\
\midrule
الدهون & الكانابينويدات الداخلية (الأنانداميد، \foreignlanguage{english}{2-AG}) & عكسي & بطيء & $-$ & المستقبِل يقلّل القذف عند المرسل: تغذية راجعة على الوزن~\cite{WilsonNicoll2001} \\
\bottomrule
\end{tabular}
\caption{مواد لا تحدد نوع العصبون. \emph{الناقل} و\emph{السرعة} و\emph{العلامة} كما في الجدول~\ref{tab:types}؛ ويضاف إليها: الناقل الحجمي، حيث تنتشر المادة في الحيز بين الخلايا حول مكان القذف؛ والعكسي، حيث تسير من المستقبِل راجعة إلى المرسل؛ والموضعي، حيث تعمل قرب مكان القذف. ”و“ إشارة إذن: لا تولّد تيارًا بذاتها، لكن من دونها لا يعمل مدخل آخر، كالمدخل الثاني لبوابة ”و“ (\foreignlanguage{english}{AND}) المنطقية. والببتيدات العصبية تُقذف دائمًا تقريبًا مع الناقل الرئيسي، وفي الغالب عند تردد نبضات مرتفع~\cite{vandenPol2012}.}
\label{tab:othermediators}
\end{table}

\chapter{الرموز}
\label{sec:notation}

يُرمز إلى أنواع العصبونات بالأحرف الأربعة الأولى من الاسم الإنجليزي لناقلها العصبي الرئيسي: \nt{GLUT}: الغلوتامات، \nt{GABA}: حمض غاما-أمينوبيوتيريك، \nt{GLYC}: الغلايسين، \nt{ACET}: الأسيتيل كولين، \nt{DOPA}: الدوبامين، \nt{SERO}: السيروتونين، \nt{HIST}: الهستامين، \nt{NORA}: النورأدرينالين، \nt{ADRE}: الأدرينالين، \nt{ASPA}: الأسبارتات (الجدول~\ref{tab:types}).

الحروف أقل من الكميات، والحرف نفسه يعني أحيانًا أشياء مختلفة في فصول مختلفة. في الجدول أدناه لكل معنى سطره؛ وفي العمود الأيسر الصيغة التي أُدخل فيها.

\begin{center}
\footnotesize
\setlength{\tabcolsep}{4pt}
\renewcommand{\arraystretch}{1.12}
\begin{longtable}{>{\raggedright}p{2.0cm}>{\raggedright}p{9.6cm}>{\raggedright\arraybackslash}p{2.4cm}}
\toprule
الرمز & المعنى & موضع إدخاله \\
\midrule
\endfirsthead
\toprule
الرمز & المعنى & موضع إدخاله \\
\midrule
\endhead
\bottomrule
\endfoot
$a$ & ميل خاصية النقل الخطية، $f(u)=au$ & \eqref{eq:feedback} \\
$a$ & مستوى \nt{ACET}: $0$ للقراءة، $1$ للكتابة & \eqref{eq:acetnet} \\
$a_i$, $a$ & تعب المجموعة: في مولّد الإيقاع؛ وفي أرجوحة النوم البطيء & \eqref{eq:halfcentre}, \eqref{eq:updown} \\
$a_S$, $a_P$ & حساسية إيقاع القلب للمسرّع وللمكبح & \eqref{eq:gasbrake} \\
$A(r)$, $A_0$ & استجابة المشبك للنبضة عند التردد $r$؛ وعند المشبك المستريح & \eqref{eq:depression} \\
$A_1$ & استجابة المستقبِل لحصة واحدة من الناقل & \eqref{eq:quantal} \\
$A$ & مساحة غشاء العصبون مع التغصنات & \eqref{eq:dendriteload} \\
$A_f$, $A_s$ & نسبة التصحيح التي تبقى حتى الحركة التالية، في العملية السريعة والبطيئة & \eqref{eq:tworate} \\
$b_i$ & التدفق الذاتي لناظمة الإيقاع & \eqref{eq:seroneuron} \\
$b$ & سرعة إتعاب العمل للمجموعة في مولّد الإيقاع أو في أرجوحة النوم & \eqref{eq:halfcentre}, \eqref{eq:updown} \\
$b$ & عدد البتات في عينة رقمنة واحدة & \eqref{eq:probedata} \\
$B_f$, $B_s$ & قوة تعلّم التصحيح من الخطأ، في العملية السريعة والبطيئة & \eqref{eq:tworate} \\
$c$ & تركيز الناقل في الحجم & \eqref{eq:seroneuron} \\
$c$ & معامل الارتباط بين ضجيج العصبونات المتجاورة & \eqref{eq:correlated} \\
$c$ & المعامل الذي يحرّك به فرقُ العدّين المضربَ & \eqref{eq:dishreadout} \\
$c_f$ & استجابة العصبون $C$: أكبر $s_{f,p}$ على كل المواضع & \eqref{eq:scstage} \\
$c_m$ & السعة النوعية للغشاء & \eqref{eq:dendriteload} \\
$C$ & كلفة الجهد: الفعل في زمن $\tau_a$ يكلّف $C/\tau_a$ & \eqref{eq:vigor} \\
$d'$ & قابلية التمييز بين دخلين: الفرق مقسومًا على الضجيج & \eqref{eq:gainsnr} \\
$d$, $d_j$ & التأخير على مسار الإشارة & \eqref{eq:glutneuron}, \eqref{eq:vstar} \\
$d$ & تأخير حلقة التغذية الراجعة مع العضلة & \eqref{eq:delaystable} \\
$D$ & معامل الانتشار & \eqref{eq:diffusionlength} \\
$D$ & المقام في ترددات شبكة \foreignlanguage{english}{\nt{GLUT}--\nt{GABA}} & \eqref{eq:isn} \\
$D$ & زمن انتظار المكافأة المؤجلة & \eqref{eq:patience} \\
$e$, $e_{ij}$ & أثر الأهلية في المشبك & \eqref{eq:threefactor} \\
$e$, $e_i$ & خطأ الخرج أو الحركة، الآتي على سلك الخطأ & \eqref{eq:lms}, \eqref{eq:adaptivefilter}, \eqref{eq:tworate} \\
$E_{\text{spk}}$ & طاقة نبضة واحدة مع عمل المشابك & \eqref{eq:energy} \\
$E$ & تردد مجموعة \nt{GLUT} & \eqref{eq:ei} \\
$E_E$ & جهد التوازن للمشبك المثير & \eqref{eq:shunt} \\
$f$, $F$ & خاصية النقل: التردد دالةً في الدخل & \eqref{eq:ratenet}, \eqref{eq:gain} \\
$f$, $f_0$ & تردد ضربات القلب؛ الإيقاع الذاتي لناظمة الإيقاع & \eqref{eq:gasbrake} \\
$f_s$ & تردد رقمنة قطب واحد & \eqref{eq:probedata} \\
$F(t)$, $F_1$ & قوة العضلة؛ قوة انتفاضة واحدة & \eqref{eq:forcesum} \\
$F_1$, $F_2$ & قوتا عضلتين تشدّان المفصل في اتجاهين متعاكسين & \eqref{eq:antagonists} \\
$g$ & حساسية المستقبل للتركيز & \eqref{eq:seroneuron} \\
$g_L$, $g_E$, $g_I$ & موصليات التسريب والإثارة والتثبيط & \eqref{eq:shunt} \\
$g$ & الكسب الذي يحدده \nt{NORA} & \eqref{eq:gain}, \eqref{eq:machine} \\
$g$ & التثبيط العام من \nt{GABA} في شبكة الذاكرة & \eqref{eq:acetnet} \\
$g$ & كسب بوابة الألم، المحدد من أعلى & \eqref{eq:gate} \\
$g$ & كسب مرحلة واحدة من الشلال الهرموني & \eqref{eq:cascade} \\
$g_j$ & المرشح $j$ بثابته الزمني الخاص عند دخل المرشح التكيفي & \eqref{eq:adaptivefilter} \\
$G$ & كسب حلقة التغذية الراجعة & \eqref{eq:feedback} \\
$G$ & سرعة التصحيح في حلقة ذات تأخير & \eqref{eq:delaystable} \\
$h$, $h_I$ & الدخل الخارجي لمجموعة \nt{GLUT} ولمجموعة \nt{GABA} & \eqref{eq:ei}, \eqref{eq:isn} \\
$h(t)$ & انتفاضة عضلية منفردة & \eqref{eq:twitch} \\
$H(t)$ & العتبة العليا لضغط النوم: ببلوغها تنام الآلة & \eqref{eq:sleeppressure} \\
$I$, $I_i$ & التدفق الخارجي إلى العصبون أو المجموعة & \eqref{eq:ratenet}, \eqref{eq:wta}, \eqref{eq:updown} \\
$I$ & تردد مجموعة \nt{GABA} & \eqref{eq:ei} \\
$I$ & سطوع المنبه أو شدته & \eqref{eq:photonnoise}, \eqref{eq:nakarushton} \\
$I$ & تيار الدخل الذي يشحن الغشاء & \eqref{eq:dendriteload} \\
$k$ & عدد النبضات في الرشقة & \eqref{eq:burst} \\
$k$ & كم مرة يجب أن تفوق الزيادةُ الضجيجَ & \eqref{eq:photonnoise} \\
$k$ & معامل التغذية الراجعة في الشلال الهرموني & \eqref{eq:cascade} \\
$k_s$ & القوة التي يضخ بها الألمُ التحسّسَ & \eqref{eq:sensitization} \\
$K$ & عدد سمات النتيجة & \eqref{eq:vectortd} \\
$K$ & صلابة المفصل & \eqref{eq:antagonists} \\
$K$ & كسب حلقة الشلال الهرموني، $K=g^3k$ & \eqref{eq:cascade}, \eqref{eq:cascadestable} \\
$K_\varphi$ & قوة مواءمة المولّد اليومي مع الإشارة الخارجية & \eqref{eq:pll} \\
$L$ & طول خط التأخير & \eqref{eq:itd} \\
$L$ & طول السلك من مستشعر الألم & \eqref{eq:paintime} \\
$L(t)$ & العتبة الدنيا لضغط النوم: ببلوغها تستيقظ الآلة & \eqref{eq:sleeppressure} \\
$M$ & مقياس الصلة السببية بين النذير والمكافأة & \eqref{eq:retro} \\
$M_k$ & توقُّع السمة $k$ & \eqref{eq:vectortd} \\
$M_{ik}$ & قوة التثبيط من عصبون \nt{GABA} رقم $k$ & \eqref{eq:machine} \\
$M$ & عزم الدوران في المفصل & \eqref{eq:antagonists} \\
$M$ & عدد خطوط الدخل إلى الخلّاطات & \eqref{eq:cover} \\
$n$, $\bar n$ & عدد النبضات في النافذة؛ ومتوسطه & \eqref{eq:rateerror} \\
$n$ & عدد مداخل العنصر العتبي & \eqref{eq:cover} \\
$n_\uparrow$, $n_\downarrow$ & عدد نبضات منطقتي ”أعلى“ و”أسفل“ في النافذة & \eqref{eq:dishreadout} \\
$N$ & عدد العصبونات & \eqref{eq:correlated}, \eqref{eq:energy} \\
$N$ & عدد أقطاب المسبار & \eqref{eq:probedata} \\
$N_s$ & عدد مواقع القذف بين عصبونين & \eqref{eq:quantal} \\
$p$ & احتمالية قذف الناقل لكل نبضة & \eqref{eq:burst} \\
$p$ & نسبة العصبونات النشطة في شبكة الذاكرة & \eqref{eq:acetnet} \\
$p$ & الاضطراب الذي يُتعلَّم تعويضه & \eqref{eq:tworate} \\
$P$ & القدرة المنفقة على الإشارات & \eqref{eq:energy}, \eqref{eq:dishpower} \\
$P$ & عدم اليقين في التقدير المخزَّن & \eqref{eq:kalman} \\
$P$ & فعالية ”المكبح“: أسلاك \nt{ACET} إلى القلب & \eqref{eq:gasbrake} \\
$P_{\max}$ & كم نمطًا عشوائيًا يستطيع العنصر العتبي أن يقسمها إلى صنفين & \eqref{eq:cover} \\
$P_{\text{miss}}$ & احتمالية الإخفاق & \eqref{eq:dishdensity} \\
$q$ & سرعة تغيّر العالم & \eqref{eq:kalman} \\
$q_k$ & تردد عصبون \nt{GABA} رقم $k$ & \eqref{eq:machine} \\
$q_0$ & الفعالية الخلفية للوسيط المثبِّط & \eqref{eq:gate} \\
$q$ & تردد الوسيط المثبِّط في بوابة الألم & \eqref{eq:gate} \\
$q$ & كم مرة تفوق الوحدة الحركية التالية سابقتها قوةً & \eqref{eq:sizeprinciple} \\
$r$, $r_i$ & تردد نبضات العصبون & \eqref{eq:ratenet} \\
$r$, $r_t$ & المكافأة (تدفق المكافأة) & \eqref{eq:rpe}, \eqref{eq:ctd} \\
$\mathbf r(t)$ & متجه استجابات الخزّان & \eqref{eq:reservoir} \\
$R$, $\bar R$ & المكافأة المتلقاة؛ وتوقُّعها أو متوسطها في وحدة الزمن & \eqref{eq:perturbation}, \eqref{eq:vigor} \\
$R$ & تباين الضجيج عند دخل المرشح & \eqref{eq:kalman} \\
$R$, $r_0$ & المكافأة المؤجلة والمكافأة الفورية & \eqref{eq:patience} \\
$R_{\max}$ & أقصى معدل للنقل، بت/ث & \eqref{eq:spikeentropy} \\
$r_{\max}$ & أعلى تردد للنبضات & \eqref{eq:gain}, \eqref{eq:nakarushton} \\
$R_{\text{raw}}$, $R_{\text{spk}}$ & تدفق بيانات المسبار: الخام والنبضات وحدها، بت/ث & \eqref{eq:probedata} \\
$s_j$ & علامة مخارج العصبون $j$ & \eqref{eq:dalesign} \\
$s_i$ & علامة مستقبل المستقبِل & \eqref{eq:seroneuron} \\
$s$, $\hat s$ & حالة العالم؛ وتقديرها & \eqref{eq:rpe}, \eqref{eq:kalman} \\
$s_{f,p}$ & استجابة العصبون $S$ للسمة $f$ في الموضع $p$ & \eqref{eq:scstage} \\
$S$ & فعالية ”المسرّع“: أسلاك \nt{NORA} إلى القلب & \eqref{eq:gasbrake} \\
$S$ & ضغط النوم & \eqref{eq:sleeppressure} \\
$t_1$ & لحظة النبضة الأولى & \eqref{eq:latency} \\
$t_k$ & لحظة النبضة رقم $k$ & \eqref{eq:forcesum} \\
$t_{f,i}$ & قالب السمة $f$ & \eqref{eq:scstage} \\
$T$ & نافذة عدّ النبضات؛ زمن تجميع الفوتونات & \eqref{eq:rateerror}, \eqref{eq:photonnoise} \\
$T_c$ & زمن صعود الانتفاضة العضلية & \eqref{eq:twitch} \\
$T_{\text{burst}}$ & الفاصل بين الرشقات الشبكية في المزرعة & \eqref{eq:burstinterval} \\
$u$ & الدخل الكلي للعصبون & \eqref{eq:gain} \\
$u$, $u(t)$ & أمر الدماغ: عند دخل المرشح التكيفي؛ وإلى الشلال الهرموني & \eqref{eq:adaptivefilter}, \eqref{eq:cascade} \\
$u(t)$ & دخل الخزّان & \eqref{eq:reservoir} \\
$U$ & مستوى الدخل الثابت & \eqref{eq:latency} \\
$U$ & نسبة مخزون الناقل التي تُقذف لكل نبضة & \eqref{eq:depression} \\
$v$ & سرعة الإشارة في السلك؛ سرعة حركة البقعة & \eqref{eq:itd}, \eqref{eq:vstar} \\
$V$ & الجهد على الغشاء & \eqref{eq:glutneuron}, \eqref{eq:shunt} \\
$V$ & توقُّع المكافأة المستقبلية & \eqref{eq:rpe}, \eqref{eq:ctd} \\
$w$, $w_{ij}$, $W_{ij}$ & وزن الرابط & \eqref{eq:glutneuron}, \eqref{eq:ratenet}, \eqref{eq:acetnet}, \eqref{eq:lms}, \eqref{eq:adaptivefilter} \\
$\mathbf w_k$ & أوزان مداخل العصبون $C$ رقم $k$ & \eqref{eq:traceinvariance} \\
$w_{q\mu}$, $w_{q\nu}$, $w_{z\mu}$ & أوزان الروابط في بوابة الألم & \eqref{eq:gate} \\
$W_{\text{out}}$ & أوزان القراءة الخطية للخزّان & \eqref{eq:reservoir} \\
$x$, $y$ & المداخل المنطقية والمخرج & \eqref{eq:dualrail}, \eqref{eq:veto} \\
$x$, $y$ & فعالية المرسل والمستقبِل & \eqref{eq:hebb} \\
$x$ & موضع الكاشف على خط التأخير & \eqref{eq:itd} \\
$x_i$ & الدخل من أعضاء الحس & \eqref{eq:acetnet} \\
$x_p$ & الدخل في الموضع $p$ & \eqref{eq:scstage} \\
$\mathbf x$ & مخارج العصبونات $S$، وهي دخل العصبون $C$ & \eqref{eq:traceinvariance} \\
$x_j$ & الدخل $j$ للمرشح التكيفي: الأمر بعد المرشح $g_j$ & \eqref{eq:lms}, \eqref{eq:adaptivefilter} \\
$x_f$, $x_s$ & تصحيحا العملية السريعة والبطيئة & \eqref{eq:tworate} \\
$x_1$, $x_2$, $x_3$ & المستويات في مراحل الشلال الهرموني الثلاث؛ $x_3$ الهرمون النهائي & \eqref{eq:cascade} \\
$x^*$ & نسبة مخزون الناقل المتعافي التي يبدأ عندها انهيار متسلسل جديد & \eqref{eq:burstinterval} \\
$y$ & دخل المرشح المشوَّش & \eqref{eq:kalman} \\
$y$, $\hat y$ & إشارة المستشعرات؛ وتنبؤ المرشح التكيفي بها & \eqref{eq:adaptivefilter} \\
$y_k$, $\bar y_k$ & فعالية العصبون $C$ رقم $k$؛ وأثرها & \eqref{eq:traceinvariance} \\
$y(t)$ & خرج قراءة الخزّان & \eqref{eq:reservoir} \\
$z$ & تردد خرج بوابة الألم & \eqref{eq:gate} \\
\midrule
$\alpha_i^\pm$ & وزن الأخطاء الموجبة والسالبة & \eqref{eq:asymmetric} \\
$\beta$ & قوة التثبيط المتبادل & \eqref{eq:wta} \\
$\beta$ & قوة تثبيط خرج بوابة الألم & \eqref{eq:gate} \\
$\gamma$ & معامل الخصم & \eqref{eq:rpe} \\
$\delta(\cdot)$ & دالة دلتا: قفزة لحظية & \eqref{eq:glutneuron} \\
$\delta$, $\delta_t$ & خطأ التنبؤ بالمكافأة & \eqref{eq:rpe}, \eqref{eq:ctd} \\
$\delta t$ & الفرق بين زمني وصول الصوت إلى الأذنين & \eqref{eq:itd} \\
$\Delta$, $\Delta_{\max}$ & الفاصل بين النبضات؛ نافذة التزامن & \eqref{eq:window} \\
$\Delta t$ & الدقة التي يميّز بها المستقبِل الزمن & \eqref{eq:spikeentropy} \\
$\Delta F$ & زيادة القوة من الوحدة الحركية التالية & \eqref{eq:sizeprinciple} \\
$\Delta I$ & الزيادة المميَّزة في السطوع & \eqref{eq:photonnoise} \\
$\Delta p$ & إزاحة المضرب في النافذة & \eqref{eq:dishreadout} \\
$\Delta u$ & الفرق بين دخلي المجموعتين & \eqref{eq:gainsnr} \\
$\Delta V$ & المقدار الذي يلزم رفع جهد الغشاء به & \eqref{eq:dendriteload} \\
$\Delta x$ & المسافة بين مستشعرين متجاورين & \eqref{eq:vstar} \\
$\varepsilon$ & الفرق النسبي بين الدخلين & \eqref{eq:wtagain} \\
$\eta$ & معدل التعلم & \eqref{eq:plasticity}, \eqref{eq:lms}, \eqref{eq:traceinvariance} \\
$\eta$ & الضجيج في الشبكة بعد المضخِّم & \eqref{eq:softmax} \\
$\mu$ & دخل اللمس إلى بوابة الألم & \eqref{eq:gate} \\
$\nu$ & دخل الألم & \eqref{eq:gate} \\
$\theta$ & عتبة الإطلاق & \eqref{eq:window}, \eqref{eq:scstage} \\
$\kappa$ & كمية الناقل لكل نبضة & \eqref{eq:seroneuron} \\
$\kappa_i$ & نسبة وزني الأخطاء الموجبة والسالبة & \eqref{eq:expectile} \\
$\kappa_m$ & صلة قوة العضلة بصلابتها & \eqref{eq:antagonists} \\
$\ell$ & المسافة التي ينتشر إليها الناقل & \eqref{eq:diffusionlength} \\
$\ell$ & ذراع العضلة التي تشدّ بها المفصل & \eqref{eq:antagonists} \\
$\xi$ & الارتجاف العشوائي لخرج العصبون & \eqref{eq:perturbation} \\
$\rho(\boldsymbol\psi)$ & نسبة الوقت الذي يُقضى في الضبط $\boldsymbol\psi$ & \eqref{eq:dishdensity} \\
$\sigma$ & التشتت، مقياس الضجيج & \eqref{eq:rateerror}, \eqref{eq:softmax} \\
$\sigma$ & السطوع الذي تبلغ عنده الاستجابة نصف أقصاها & \eqref{eq:nakarushton} \\
$\sigma_{\text{pre}}$, $\sigma_{\text{post}}$ & الضجيج قبل المضخِّم وبعده & \eqref{eq:gainsnr} \\
$\tau$ & الثابت الزمني للغشاء & \eqref{eq:glutneuron} \\
$\tau$ & الثابت الزمني لمرحلة الشلال الهرموني & \eqref{eq:cascade} \\
$\tau_{\text{eff}}$ & الثابت الزمني لمجموعة تثير نفسها & \eqref{eq:perfectint} \\
$\tau_c$ & عمر الناقل في الحجم & \eqref{eq:seroneuron} \\
$\tau_E$, $\tau_I$ & الثابتان الزمنيان لمجموعة \nt{GLUT} ومجموعة \nt{GABA} & \eqref{eq:ei} \\
$\tau_e$ & عمر أثر الأهلية & \eqref{eq:threefactor} \\
$\tau_{\text{tr}}$ & عمر أثر فعالية العصبون $C$ & \eqref{eq:traceinvariance} \\
$\tau_s$ & زمن عودة الحساسية بعد الإصابة & \eqref{eq:sensitization} \\
$\tau_r$, $\tau_d$ & زمن تراكم ضغط النوم في اليقظة وزمن تفريغه في النوم & \eqref{eq:sleeppressure} \\
$\tau_{\text{rec}}$ & زمن تعافي مخزون الناقل & \eqref{eq:depression} \\
$\tau_\gamma$ & أفق التخطيط & \eqref{eq:ctd} \\
$\tau_v$ & سرعة تحديث التوقُّع & \eqref{eq:ramp} \\
$\tau_a$ & زمن التعب في مولّد الإيقاع أو في أرجوحة النوم & \eqref{eq:halfcentre}, \eqref{eq:updown} \\
$\tau_a^*$ & أفضل زمن لأداء الفعل & \eqref{eq:vigor} \\
$\Phi$ & الطرف الأيمن لقاعدة التعلم & \eqref{eq:plasticity} \\
$\Phi$ & تدفق الفوتونات & \eqref{eq:photonnoise} \\
$\varphi_k$ & السمة $k$ للنتيجة المتلقاة & \eqref{eq:vectortd} \\
$\varphi$ & فرق الطور بين المولّد اليومي والإشارة الخارجية & \eqref{eq:pll} \\
$\boldsymbol\psi$ & ضبط الشبكة في نموذج \foreignlanguage{english}{DishBrain} & \eqref{eq:dishdensity} \\
$\omega$, $\omega_0$ & تردد الإشارة الخارجية (الضوء)؛ التردد الذاتي للمولّد اليومي & \eqref{eq:pll} \\
\end{longtable}
\end{center}

\chapter{الأجوبة}
\label{sec:answers}

أجوبة موجزة عن التمارين في آخر الفصول. أما الحلول الكاملة ففي دليل منفصل للمدرّسين.

\answer{ex:01:1}{(أ)~$8\cdot10^{10}$ من \nt{GLUT}، أي $10$ أضعاف سكان الأرض؛ وعصبونات البثّ $\approx1.9\cdot10^6$، أي مدينة بحجم فيينا. (ب)~$\approx8\cdot10^{14}$ نهاية، منها $\approx3\cdot10^{11}$ لعصبونات البثّ، بحصة $\approx4\cdot10^{-4}$، أي أكبر بـ$\approx16$ مرة من حصة هذه العصبونات ($2.2\cdot10^{-5}$).}
\answer{ex:01:2}{(أ)~$\tau_c\ln(c_0/c_*)\approx4.6\ms$. (ب)~يكون الخط حرًا عند $T\ge\tau_c\ln101$: حتى $\approx22$ قذفة في الثانية بدل $\approx220$.}
\answer{ex:01:3}{$V(s_k)=\gamma^{K-1-k}r$؛ $\delta_{\text{lamp}}=\gamma^Kr=re^{-T/\tau_\gamma}\approx0.60\,r$ لأي $\Delta$، و$\tau_\gamma\approx1.95\,\text{s}$.}
\answer{ex:01:4}{$n^*=93$، $47$، $25$، $12$: أي $n^*\approx5/\alpha$ عند قيم $\alpha$ الصغيرة.}
\answer{ex:01:5}{\nt{GLUT}: $1\to10\to10^4\to10^7$، أي $\approx10^7$ عصبون، و$4$ حلقات، و$4\ms$. \nt{DOPA}: $1\to10\to3.2\cdot10^5$، أي $\approx3\cdot10^5$ عصبون، و$3$ حلقات، أقل بـ$30$ مرة؛ وهي رتبة عدد عصبونات \nt{DOPA} نفسها.}
\answer{ex:01:6}{$\lambda=f_EJ_E-f_IJ_I$، ومنها $J_I=37.5$؛ نسبة التيارين $f_IJ_I/f_EJ_E=0.94$؛ ويحدث الانهيار المتسلسل ($\lambda\ge1$) إذا أُضعف التثبيط بـ$6.7\,\%$ فقط.}
\answer{ex:01:7}{$n\ge57$ عرضًا لكل تأخير في كل حالة. والهبوط نفسه قد ينتج مثلًا عن عدم دقة في التقدير الداخلي للزمن تتزايد مع $T$.}

\answer{ex:02:1}{صف طوله $3.2\,\text{mm}$ بخطوة $25\um$: $\approx129$ كاشفًا. ينطلق شريط عرضه $v\,\tau\ln2=1.7\,\text{mm}$، أي $\approx69$ كاشفًا، أكثر من نصف الصف؛ والاتجاه هو مركز الشريط.}
\answer{ex:02:2}{النافذة $-\tau\ln\frac{w_2}{\theta-w_1}\le\delta\le\tau\ln\frac{w_1}{\theta-w_2}$، أي $[-0.235,\,0.178]\ms$؛ والمركز $c=\frac\tau2\ln\frac{w_1(\theta-w_1)}{w_2(\theta-w_2)}=-28\,\mu\text{s}$. يزاح الاتجاه بـ$28\,\mu\text{s}$ نحو الأذن الأقوى، أي نحو ثلاث خطوات من خطوات الميز.}
\answer{ex:02:3}{يلزم وجود $s_i$ بحيث $s_iw_{ij}s_j\ge0$؛ وهي موجودة إذا وفقط إذا كانت الدورات زوجية. التثبيط المتبادل يؤول إليها (لا تذبذبات مستقرة)، وحلقة \foreignlanguage{english}{\nt{GLUT}--\nt{GABA}} لا تؤول (قد تتذبذب).}
\answer{ex:02:4}{ستة عصبونات $g_k=[x+y+z\ge k]$ و$\bar g_k=[\bar x+\bar y+\bar z\ge4-k]$، $k=1,2,3$؛ و$C=g_2$، $\bar C=\bar g_2$، $S=[g_1+\bar g_2+g_3\ge2]$، $\bar S=[\bar g_1+g_2+\bar g_3\ge2]$: ثمانية عصبونات. ومن \foreignlanguage{english}{AND} و\foreignlanguage{english}{OR}: $12$.}
\answer{ex:02:5}{$T_{\min}(0.6)=0.718\,\tau$؛ $I_c=0.225$، وقرب العتبة $T_{\min}\propto(I-I_c)^{-1/2}$.}
\answer{ex:02:6}{(أ)~$500$ حلقة، و$5\cdot10^4$ عصبون؛ التشتت $4.5\ms$ ($0.45\,\%$)، والخطأ النسبي $\propto T^{-1/2}$. (ب)~معامل عشوائي للتأخير مشترك بين كل الحلقات بتشتت $5\,\%$.}
\answer{ex:02:7}{الإنعاش مرة كل $\tau_F\ln2=0.69\,\text{s}$: $4.3$ نبضة في الثانية لكل عصبون مقابل $20$، أي أوفر بـ$4.6$ مرة ($4\cdot10^{-7}$ مقابل $2\cdot10^{-6}\,\text{J}$ عند $10^{-10}\,\text{J}$ للنبضة). القلّاب يفقد البت؛ والمكثف يحفظه إذا بدأ الإسكات في غضون $0.49\,\text{s}$ من الإنعاش (الاحتمال $0.71$).}

\answer{ex:03:1}{$35\to17.5\mV$ عند $g_E=g_L$؛ و$56\to40\mV$ عند $g_E=4g_L$، والطرح كان سيعطي $38.5\mV$: التحويل يقسم $g_E$ على $3$. تضيق النافذة إلى النصف (ينخفض $\tau_{\text{eff}}$ بعامل $2$).}
\answer{ex:03:2}{$\operatorname{tr}=\frac{w_{EE}-1}{\tau_E}-\frac1{\tau_I}$، $\det=\frac{1-w_{EE}+w_{EI}w_{IE}}{\tau_E\tau_I}$؛ وعلى الحد $\omega=\sqrt{\det}$، و$\tau_I=20\ms$، والدور $73\ms$.}
\answer{ex:03:3}{$d_c=\tau_E\arccos(-1/G)/\sqrt{G^2-1}$، والدور $2\pi\tau_E/\sqrt{G^2-1}\in(2d_c,4d_c)$. $G=2$: $d_c=12.1\ms$، والدور $36\ms$؛ $G=5$: $d_c=3.6\ms$، والدور $12.8\ms$، أي ضمن نطاق الإيقاعات السريعة تمامًا.}
\answer{ex:03:4}{(أ)~عند $I_2=\beta I_1=2$ صعودًا و$I_2=I_1/\beta=0.5$ نزولًا، بحلقة عرضها $1.5$. (ب)~بمقدار $\tau\ln10/(\beta-1)=2.3\tau$ لكل رتبة عشرية من $\varepsilon$.}
\answer{ex:03:5}{ثلاثة وسطاء، $d_k=5$ و$10$ و$15\ms$: يجب أن تغطي الحظور $15\ms$، بمعدل $5\ms$ لكل منها.}
\answer{ex:03:6}{$n+M+1=3+4+1=8$ عصبونات. واثنان: $\text{\nt{GABA}}=[x+y+z\ge2]$، و$\text{المخرج}=[x+y+z-2\,\text{\nt{GABA}}\ge1]$.}
\answer{ex:03:7}{$\binom84=70<100\le\binom94=126$: $9$ وسطاء (مبرهنة شبيرنر). ومن دون وصلات مباشرة: $100$، لأن رتبة $\beta(J-I)$ تساوي $n$.}
\answer{ex:03:8}{$h_c=\frac1{2w_{EE}}+\frac{w_{EI}w_{IE}-w_{EE}}{4w_{EE}^2}=\frac7{18}\approx0.39$؛ والنمذجة تعطي تغيّر العلامة بين $h=0.38$ و$0.40$.}

\answer{ex:04:1}{(أ)~$\approx1.1\cdot10^{11}\,\text{bit/J}$. (ب)~$\log_2\sqrt{10}\approx1.7$ بت مقابل $81$: أقل بـ$\approx50$ مرة.}
\answer{ex:04:2}{$r^*=(P_0/E_{\text{spk}})\ln\bigl(1/(r^*\Delta t)\bigr)$، و$P_0/E_{\text{spk}}=1.16\,\text{s}^{-1}$: $r^*\approx6$ نبضات في الثانية، و$7.4\cdot10^{10}\,\text{bit/J}$.}
\answer{ex:04:3}{$\mathcal I=\frac{N}{1+(N-1)c}=9.9$ (الحد $10$) مقابل $\mathcal I=\frac{N}{1-c}=1111$: فرق بـ$\approx110$ مرة؛ الارتباط ضار فقط حين يشبه الإشارة.}
\answer{ex:04:4}{$\theta\approx68.7$ و$19.9$ (بوحدات $w$)، و$m=19$ و$m=10$: يحتاج المستقبِل السريع إلى نحو نصف عدد المداخل المتزامنة، مع أن مدخله المتوسط أصغر بخمس مرات.}
\answer{ex:04:5}{$U\tau_{\text{rec}}\ge0.725\,\text{s}$ و$\tau_{\text{rec}}(1-2U)\le0.05\,\text{s}$، وهذا يتطلب $U\ge0.483$؛ أصغر $\tau_{\text{rec}}=0.725\,\text{s}$ عند $U=1$ (و$1.45\,\text{s}$ عند $U=0.5$).}
\answer{ex:04:6}{(أ)~$\theta\,e/(e-1)\approx1.58\,\theta$ لأي $\tau$ و$\theta$. (ب)~$14$ نبضة.}
\answer{ex:04:7}{(أ)~$1.62$ بت للكلمة: $0.77$ لعدد النبضات، و$0.84$ لمواضعها. (ب)~$1.13$ و$1.59$ بت: من $30$ كلمة يفقد التقدير نحو الثلث.}

\answer{ex:05:1}{$\ell\approx17\um$؛ و$\approx100$ يوم: البثّ عبر الدماغ يؤديه تفرّع السلك، أما الانتشار فلا يفعل إلا أن يبسط كل موقع تحرير على $\sim20\um$.}
\answer{ex:05:2}{$r=ab/(1+G)$، و$S=1/(1+G)$ تمامًا؛ $+1.7\,\%$ بدل $+20\,\%$.}
\answer{ex:05:3}{$T^*=1.21\,\text{s}$ عند $G=2$ و$0.19\,\text{s}$ عند $G=9$: عمليًا $G\sim2\text{--}4$، أي كبت بـ$3\text{--}5$ مرات.}
\answer{ex:05:4}{$\mathrm{CV}=1/\sqrt{2\lambda\tau_c}\approx9\cdot10^{-4}$ عند تردد كلي $\lambda$؛ والعصبون الواحد يحتاج إلى $\tau_c=12.5\,\text{s}$.}
\answer{ex:05:5}{$c\approx0.40$ في الحالتين؛ $\mathrm{CV}=0.050$ مقابل $0.158$، أي أقل بـ$\sqrt{10}$ مرة. تبقى ترددات العصبونات المفردة متناسبة مع $a_i$: لا يُضبط إلا المجموع.}
\answer{ex:05:6}{$N_1=\overline{A\vee B}$، $N_2=\overline{A\vee N_1}$، $N_3=\overline{B\vee N_1}$، $N_4=\overline{N_2\vee N_3}$، و$\mathrm{XOR}=N_5=\overline{N_4}$؛ كل تبديل يستغرق $\tau_c\ln2$، والمسار من $4$ بوابات يستغرق $2.8\,\text{s}$ لكل \foreignlanguage{english}{XOR}.}
\answer{ex:05:7}{(أ)~$\tau_\gamma=6/\ln3\approx5.5\,\text{s}$. (ب)~$\bar D<4\,\text{s}$ مقابل $D<2.2\,\text{s}$: عدم إمكان التنبؤ بالتأخير يزيد الصبر.}
\answer{ex:05:8}{$k_2=1/\tau_d=10\,\text{s}^{-1}$، $k_1=1/\tau_d-1/\tau_\gamma=9.5\,\text{s}^{-1}$؛ $\tau_\gamma=1/(k_2-mk_1)$: يتضاعف الأفق عند $+2.6\,\%$ (وعند $+5.3\,\%$ يصير لانهائيًا).}

\answer{ex:06:1}{$e=(1-e^{-T_a/\tau_e})e^{-d/\tau_e}$. (أ)~$0.110$ مقابل $4.5\cdot10^{-5}$، أي $\approx2.5\cdot10^3$ مرة. (ب)~$\tau_e^*=T_a/\ln(1+T_a/d)=0.59\,\text{s}$.}
\answer{ex:06:2}{المعدل $\nu_x\nu_y(A_+\tau_+-A_-\tau_-)=0.8A_+$ في الثانية، أي $\approx2900A_+$ في الساعة؛ $\tau_-=\tau_+/0.6=33\ms$.}
\answer{ex:06:3}{$\mathbf e_1$ عند $\mu^2+s_1^2>s_2^2$؛ هنا $1.01>0.25$، فيتعلم العصبون $\mathbf e_1$ مع أن التشتت على امتداد $\mathbf e_2$ أكبر بـ$25$ مرة: القاعدة تجد المتجه الرئيسي لمصفوفة الارتباطات لا التغايرات.}
\answer{ex:06:4}{$V=Re^{-(t_c+L-t)/\tau_\gamma}$ بين المصباح والعصير، فيكون $\dot V=V/\tau_\gamma$؛ ذروة مساحتها $Re^{-L/\tau_\gamma}$: $0.61R$ و$0.78R$.}
\answer{ex:06:5}{نسبة الإشارة إلى الضجيج $1/(N+1)$، وعدد المحاولات $\propto N+1$: $5\cdot10^5$ محاولة $\approx1$ شهر، و$5\cdot10^{11}$ محاولة $\approx8\cdot10^4$ سنة.}
\answer{ex:06:6}{(أ)~$k_2=1/\tau_d$، $k_1=1/\tau_d-1/\tau_\gamma$. (ب)~الخطأ الزائف $-(V/\tau_\gamma)\,\varepsilon/(1+\varepsilon)$، $\varepsilon=\tau_d/\tau_\gamma$، ومنه $\tau_d\le0.105\,\text{s}$.}
\answer{ex:06:7}{$0.46$، $0.90$، $0.99$، $0.95$، و$0.70$ عند $d=3\,\text{s}$. تقع النقاط على منحنى واحد بدلالة نسبة الأثر $F=(1-e^{-T_{\text{cue}}/\tau_e})e^{-d/\tau_e}$: تعلّم عند $F\gtrsim0.1$، واختيار عشوائي عند $F<10^{-2}$.}
\answer{ex:06:8}{$\eta^*=1/\bigl(2\sigma^2(N+2)\bigr)$، في $N+2$ محاولة؛ ومع $m$ عصبونات مرتعشة من $N$ في $N(m+2)/m\ge N+2$: التناثر لا يساعد.}

\answer{ex:07:1}{$V=\kappa p/\bigl(\kappa p+(1-\kappa)(1-p)\bigr)$: $0.059$، $0.20$، $0.50$؛ الوسيط يساوي $0$، أما النقطة~\eqref{eq:expectile} فتتغير مع $\kappa$ تغيرًا سلسًا.}
\answer{ex:07:2}{$p=1/3$، $x=0.75$، $\sigma_r=0.35$، $V(0.2)=0.083$.}
\answer{ex:07:3}{$V=\sqrt\kappa/(\sqrt\kappa+\sqrt{1-\kappa})$، من $0.25$ إلى $0.75$؛ التشتت $\propto\sqrt\eta$؛ وللقطرتين $V=0.25+0.5\kappa$، فيختلف التوزيعان بـ$0.05$ عند العصبونات الطرفية، ويلزم $\eta\lesssim0.01$.}
\answer{ex:07:4}{(أ)~$J^*=2\sqrt{C\bar R}$. (ب)~الزيادة $\bigl(k^{1/2}+k^{-1/2}\bigr)/2-1$: $6\,\%$ عند $k=2$ و$25\,\%$ عند $k=4$، فتكفي دقة ”في حدود الضعف“.}
\answer{ex:07:5}{ذروة مساحتها $R(e^{-1}-e^{-2})\approx0.23R$، أي ما يعادل $2.3\,\text{s}$ من الارتفاع التدريجي؛ وإذا كان الدوبامين يخبر بـ$V$ فلا ذروة، بل يقفز المستوى درجةً فيكبر $e$ مرة.}
\answer{ex:07:6}{الحدث يزيح الوزن $\propto A\tau_f/(\tau_e+\tau_f)$، والمستوى يعطي خطأ $s\tau_f$: $9\,\text{s}\le\tau_f\le17\,\text{s}$.}
\answer{ex:07:7}{في البداية $\Delta P=0.80$، $M=0.90$. (أ)~$\Delta P=0.74$ ($-8\,\%$)، $M=0.43$ ($-52\,\%$)؛ (ب)~$\Delta P=0.40$ ($-50\,\%$)، $M=0.80$ ($-11\,\%$): تفارق مزدوج.}
\answer{ex:07:8}{$N_{\text{eff}}\to2+\rho^2N\approx10^6$ محاولة: أقل بـ$10^4$ مرة مما مع سلك واحد، لكن تسربًا قدره $1\,\%$ ما زال يكلّف مليون محاولة.}

\answer{ex:08:1}{(أ)~$\approx1.7\cdot10^6$. (ب)~$g=\frac{0.9}{\sqrt{0.19}}\,\sigma_{\text{post}}/\sigma_{\text{pre}}\approx16.5$: الجواب لا يعتمد إلا على نسبة الضجيجين.}
\answer{ex:08:2}{(أ)~القيم الذاتية $-(1\pm g\beta)/\tau$؛ يخمد الفرق في $\tau/(1-g\beta)$: $40\ms$ عند $g\beta=0.5$ (يتضخم فرق المدخلين ثلاث مرات)، و$200\ms$ عند $0.9$ ($19$ مرة). (ب)~$0.905$ و$1.105$؛ وبين الحدين لا يفوز إلا الدخل الأقوى.}
\answer{ex:08:3}{$P(k)=e^{gu_k/\sigma}/\sum_le^{gu_l/\sigma}$، أي~\eqref{eq:softmax}؛ $g=10\ln9\approx22$.}
\answer{ex:08:4}{$\bar\Delta=1/3$، $g^*\approx3\ln(1/h)$: $21$ و$14$ و$9$ مقابل $16\text{--}23$ و$11$ و$8$ في النموذج.}
\answer{ex:08:5}{$g^*$ من $16\text{--}23$ عند $h=0.001$ إلى $6\text{--}8$ عند $h=0.2$؛ التقدير $3\ln(1/h)$ صحيح في حدود مرة ونصف عند $h\le0.05$؛ وعند $h\gtrsim\eta/10$ لا تلحق القاعدة $Q\leftarrow Q+\eta(R-Q)$ باستيعاب ما جلبه الاستكشاف، فيعلق $g^*$ قرب $8$.}
\answer{ex:08:6}{$T_p=\tau\ln9=44\ms$ عند $g_{\text{lo}}=0$ (النموذج $44\ms$)، و$70\ms$ عند $g_{\text{lo}}=0.25$: رشقة مدتها $2\text{--}3\tau$ مع كسب قريب من الصفر.}
\answer{ex:08:7}{(أ)~أفضل ثابت $0.925$ ($g=8$)، والعرّاف $0.929$ ($16/8$): لا يكاد يوجد ما يُكسب. (ب)~مثلًا حقب هادئة بمكافآت نادرة ($p\in[0,0.3]$) وحقب متقلبة مع $h=0.1$: $0.850$ مقابل $0.872$؛ ويستعيد الضبط نحو الربع عند $\eta=0.2$ وكل شيء تقريبًا عند $\eta=0.05$.}

\answer{ex:09:1}{(أ)~$P_\infty=0.2$، والذاكرة $20\,\text{s}$. (ب)~$P_\infty=1$، والذاكرة $1\,\text{s}$. (ج)~الذاكرة تتناسب مع سعة الضجيج: ست مرات.}
\answer{ex:09:2}{$P(t)=P_\infty\coth(t\sqrt{q/R})$. (أ)~$\frac12\sqrt{R/q}=5\,\text{s}$، أي نصف الذاكرة. (ب)~$t=\frac12\ln21\,\sqrt{R/q}\approx15\,\text{s}$: هذه هي المدة التي يجب أن يبقى فيها \nt{ACET} مرتفعًا.}
\answer{ex:09:3}{التباين $qT/2+R/(2T)$؛ $T^*=\sqrt{R/q}$، والتباين $\sqrt{qR}=P_\infty$، كما في المرشح~\eqref{eq:kalman}؛ ويزداد $(\kappa+1/\kappa)/2$ مرة: $1.25$ و$5.05$.}
\answer{ex:09:4}{$a>a_w=1-\theta/(mw)$: $0.15$ عند $w=0.041$، و$0.22$ عند $w=0.045$.}
\answer{ex:09:5}{تظهر الأخطاء عند $M\approx40\text{--}60$، وعند $M=80$ يوجد $1.9$ عصبون زائد، وعند $M=100$ تكون نسبة النمط $0.86$؛ تداخل الأنماط الأخرى تباينه $\approx(M-1)p(1-p)/N$، فيكون $M_{\max}\approx80$، و$M_{\max}\propto N/p$.}
\answer{ex:09:6}{مثلًا $\eta(a)=\eta\min\bigl(1,\max(0,(a-0.3)/0.6)\bigr)$. مع $\eta a$ دخلت العصبونات الغريبة الثلاثة كلها النمط؛ ومع $\eta(a)$ لم يدخل أي منها، والكتابة عند $a\ge0.9$ كما هي (النسبة $1.00$).}
\answer{ex:09:7}{(أ)~$40$ إعادة تشغيل؛ ومع $\eta(a)$ من التمرين~\ref{ex:09:6} أبدًا. (ب)~$200$ إعادة تشغيل، أي $10$ ليالٍ.}

\answer{ex:10:1}{الناقل $\approx10^{12}$ بت في الثانية ($11.4$ بت لكل نبضة)، والتحكم $\approx40$ بت في الثانية؛ $2.5\cdot10^{10}\,\text{s}$، أي نحو ثمانمئة سنة.}
\answer{ex:10:2}{$a>1/3$ (عند $a=0$ تنشأ تذبذبات متنامية $\approx67\,\text{Hz}$). لا: الاستقرار يحدده حاصل الضرب $g(1-a)$.}
\answer{ex:10:3}{المتوسط $\sigma^2R'x_j$، والتباين $(N+1)\,x_j^2\sigma^4R'^2$؛ النسبة $1/(N+1)$، ويلزم $n\approx z^2(N+1)\propto N$ محاولة.}
\answer{ex:10:4}{$L_v\in[32.9,\,47.9]\ms$: التأخير $35.4\ms$، والنافذة $\pm7.5\ms$؛ الربط الكاذب $2\Delta_{\max}r_ar_v\approx0.38$ في الثانية.}
\answer{ex:10:5}{$0.098$ مقابل $0.180$ عند أفضل $\alpha=0.2$: أقل بنسبة $45\,\%$.}
\answer{ex:10:6}{$n=\ln\bigl(\ln1.5/(0.6g)\bigr)/\ln(1-\alpha)$: $24$ و$4$؛ وعند $g=2$: $11$ و$2$.}
\answer{ex:10:7}{$S'=S\bigl(1-4\eta\sigma^2+4\eta^2\sigma^4(G+2)\bigr)$، وأفضل قيمة $\eta\sigma^2=1/(2(G+2))$؛ $\approx10^6$ محاولة، أي نحو شهر.}

\answer{ex:11:1}{(أ)~$T=k^2/\bigl(\Phi(\Delta I/I)^2\bigr)=9\,\text{s}$. (ب)~$45$ مستشعرًا، بقعة $\approx7\times7$، والدقة المكانية أسوأ بـ$\approx7$ مرات.}
\answer{ex:11:2}{(أ)~من $2$ إلى $3.3$ بت لكل نبضة، أي $22\text{--}34\,\%$ من الحد ($9.1\text{--}9.8$ بت). (ب)~$\log_2\binom{400}{20}\approx111$ بت؛ وفي المتوسط نوع مشترك واحد ($1$).}
\answer{ex:11:3}{(أ)~$\sigma/9\le I\le9\sigma$، أي $1.9$ رتبة. (ب)~$6$ و$7$.}
\answer{ex:11:4}{(أ)~$f<343/0.44\approx780\,\text{Hz}$. (ب)~$625$ نبضة، أي $125$ عصبونًا.}
\answer{ex:11:5}{$21$ بت لكل حدث، $4.8\cdot10^5$ حدث في الثانية، $7$ أحداث لكل بكسل من الحافة: $136$ بكسل في الثانية. الكاميرا العادية: $1.25$ إطار في الثانية.}
\answer{ex:11:6}{لوغاريتميًا $\approx1.0\,\text{s}$ في الاتجاهين (النظرية $0.96\tau$)؛ خطيًا $0.2\,\text{s}$ صعودًا و$3.0\,\text{s}$ هبوطًا (النظرية $0.18\tau$ و$3.0\tau$).}
\answer{ex:11:7}{(أ)~$\ln I$ موزع لوجستيًا بوسيط $\ln\sigma$ ومقياس $1$: الانتشار $\pi/\sqrt3$ باللوغاريتمات الطبيعية، أي $0.79$ رتبة. (ب)~$r=r_{\max}\Phi_I(I)^2$.}

\answer{ex:12:1}{$0.60$ (بلا ترابط كانت ستكون $0.18$)، والحد $\sqrt{c/(rT)}\approx0.58$: مئة عصبون لا تميّز إلا ”يوجد/لا يوجد“.}
\answer{ex:12:2}{(أ)~$10^8$ نبضة، $\approx10\,\text{mJ}$: $0.3\,\%$ من $3\,\text{J}$ للدماغ كله. (ب)~$3\,\text{J}$، أي أكثر بـ$\approx300$ مرة.}
\answer{ex:12:3}{(أ)~$\theta_A\in[2,3)$، $\theta_B\in[1,2)$؛ الشكل الآخر لا يعطي أكثر من $2$ و$1$. (ب)~$10\cdot96^2\approx9.2\cdot10^4$.}
\answer{ex:12:4}{(أ)~$T_d=0.85\,\tau_a=0.34\,\text{s}$. (ب)~$T_d\propto\tau_a$ ولا يتوقف على $I$: $\tau_a\approx3.5\,\text{s}$ (النمذجة: $3.1\,\text{s}$).}
\answer{ex:12:5}{$p\approx0.17$. تنخفض النسبة ببطء: $0.90$، $0.88$، $0.86$، $0.84$، $0.82$، $0.79$ عند $N=8$، $12$، $24$، $48$، $96$، $192$.}
\answer{ex:12:6}{التشغيلات العشرة كلها بلا خطأ عند $\tau_{\text{tr}}=20$ و$50$ و$200\ms$ (عند $30\ms$ فشل تشغيل واحد من عشرة)؛ عند $5\text{--}10\ms$ نحو $0.5$، وعند $0.5\text{--}2\,\text{s}$ تنخفض الجودة إلى $0.86$. الحد الأعلى: يجب أن يخمد الأثر خلال التوقف، $e^{-0.3\,\text{s}/\tau_{\text{tr}}}\ll1$.}
\answer{ex:12:7}{تأخير واحد لا يكفي: النافذة $\pm5\ms$ لا تغطي $\Delta x/v$ من $5$ إلى $20\ms$. فرعا تثبيط بـ$5<d_1\le10\ms$ و$15\le d_2\le d_1+10\ms$.}
\answer{ex:12:8}{\foreignlanguage{english}{\texttt{two\_stage}}: $\Delta=2$، تعادل من قلب واحد، وتغيّر الجواب من قلبين. عدّاد الآحاد: قلب واحد؛ $0.57$ عند $p=0.1$.}

\answer{ex:13:1}{(أ)~$q=100^{1/99}\approx1.048$، والخطوة $\approx4.8\%$؛ والمدى $\approx2.2\cdot10^3$ ($67\,\text{dB}$). (ب)~الخطوة $22f_1$، أي $220\%$ عند $10f_1$.}
\answer{ex:13:2}{$P_{\text{fail}}\approx1.8\cdot10^{-4}$ و$2.8\cdot10^{-16}$: نحو $300$ إخفاق في اليوم عند $S=2$، و$5\cdot10^{-10}$ في اليوم عند $S=4$.}
\answer{ex:13:3}{(أ)~$H(s)=eF_1T_c/(1+sT_c)^2$، مرشح من الرتبة الثانية بتردد قطع $1/(2\pi T_c)\approx3.2\,\text{Hz}$. (ب)~$r\approx22$ نبضة في الثانية (بالتوافقية الأولى $\approx20$).}
\answer{ex:13:4}{(أ)~عند $Gd=1/e$ يكون الجذر $s=-1/d$ مضاعفًا، ولا يقع الجذر الأيمن إلى يساره عند أي $G$ آخر. (ب)~$d=30\ms$: $G\approx12\,\text{s}^{-1}$، $\tau=30\ms$؛ $d=150\ms$: $G\approx2.5\,\text{s}^{-1}$، $\tau=150\ms$؛ ثابت الزمن يساوي التأخير.}
\answer{ex:13:5}{(أ)~$E\approx0.427$، $T\approx1.70\,\tau_a=0.68\,\text{s}$ (النموذج مع $\tau=20\ms$ يعطي $0.84\,\text{s}$). (ب)~المعادلة لا تحوي $I$؛ والإيقاع موجود عند $b>\beta-1$.}
\answer{ex:13:6}{عند $b=0.8$ لا إيقاع؛ وعند $b=1.05$، $1.2$، $1.5$، $2$، $4$ يكون $T\approx2.73$، $1.74$، $1.19$، $0.84$، $0.45\,\text{s}$؛ وعند أي $I$ يكون $T=0.840\,\text{s}$: الدخل يغيّر السعة فقط، لا الوتيرة.}
\answer{ex:13:7}{التعويض $s=u-a$ يعطي أكبر معدل خمود $a+1/d$ عند $G=e^{-1-ad}/d$؛ $a\ge33.3\,\text{s}^{-1}$، $G=e^{-2}/d\approx4.5\,\text{s}^{-1}$؛ وكلما اشتد الشد المشترك (زاد $a$) وجب أخذ $G$ أصغر.}
\answer{ex:13:8}{مسألة مفتوحة. مثلًا، عتبة في تغذية التعب $b[r_i-r_0]_+$ تعطي $I_{\min}=\beta b r_0/(b-\beta+1)\approx0.53$ عند $b=4$ و$r_0=0.2$، ودورًا من $1.8\,\text{s}$ عند $I=0.6$ إلى $0.52\,\text{s}$ عند $I=3$.}

\answer{ex:14:1}{(أ)~بمقدار $0.17\,\text{s}$ و$1.5\,\text{s}$. (ب)~من مسافة أقل من $11\,\text{cm}$.}
\answer{ex:14:2}{اللمس بلا ألم لا يمر عند أي $\mu$ ما دام $g\le\beta w_{q\mu}/w_{z\mu}=5$؛ واللمس $\mu=1$ يمر عند $g>6$: التحسّس القوي يجعل اللمسة العادية ألمًا.}
\answer{ex:14:3}{(أ)~$g^*=(1-k_sC)/(1-k_sA)$، مستقر عند $k_sA<1$. (ب)~بلا لمس عند $\nu\ge2$، ومع اللمس ابتداءً من $\nu\ge1.7$.}
\answer{ex:14:4}{$V(t)=-Pe^{-(T-t)/\tau_\gamma}$. (أ)~$-Pe^{-T/\tau_\gamma}\approx-0.37P$. (ب)~$+Pe^{-(T-t_e)/\tau_\gamma}\approx+0.61P$؛ يزداد مع $t_e$ حتى $+P$ ويعلّم تجنب الألم.}
\answer{ex:14:5}{عند $k_s=4$ لا انغلاق (الحالة المستقرة الثانية توجد عند $k_s\ge4.58$)؛ $T_{\min}\approx4.35$، $1.40$، $0.65$، $0.32\,\tau_s$ عند $k_s=4.6$، $5$، $6$، $8$.}
\answer{ex:14:6}{$w_{q\nu}=4/9\approx0.444$، $q_0=2/9\approx0.222$، $w_{q\mu}=49/45\approx1.089$؛ اللمس بلا ألم آمن حتى $g=49/9\approx5.4$.}
\answer{ex:14:7}{(أ)~العتبة $\nu_c(g)=\theta/g$ عند $\nu_c\ge q_0/w_{q\nu}$، وإلا $(\theta+\beta q_0)/(g+\beta w_{q\nu})$؛ $g^*\approx2.45$، $1.0$، $0.25$. (ب)~سؤال مفتوح.}

\answer{ex:15:1}{(أ)~$2\cdot10^5$ لكل عصبون، و$6\cdot10^{12}$ للدارة. (ب)~$\approx8$ بت في الثانية لكل سلك، فلا يقل الزمن عن $2.5\cdot10^4\,\text{s}\approx7$ ساعات.}
\answer{ex:15:2}{العامل $2$: الترابط بين الجارين $0.943$، و$\lambda$ من $4.18$ إلى $7.3\cdot10^{-4}$، والنسبة $\approx5.7\cdot10^3$؛ العامل $4$: الترابط $0.8$، والنسبة $\approx94$.}
\answer{ex:15:3}{(أ)~$\lambda=0.988$ و$0.808$: $82$ و$4.7$ حركة. (ب)~$x_s^*=0.690$، $x_f^*=0.172$، والمجموع $0.862$.}
\answer{ex:15:4}{(أ)~$116$ حركة. (ب)~$19$. والنموذج ذو العملية الواحدة لا يعطي أي توفير حين يكون المجموع صفرًا.}
\answer{ex:15:5}{(أ)~$G=50\,\text{s}^{-1}$ مقابل الحد $10.5\,\text{s}^{-1}$ بلا نموذج. (ب)~لا: عند $G=50\,\text{s}^{-1}$ يكون التأخير الحرج $d_c\approx103\ms$؛ وعند $d=150\ms$ يلزم $\hat k>0.445k$ (ولأي $G$ و$d$: $\hat k>k/2$).}
\answer{ex:15:6}{ينخفض الخطأ دون $0.5\%$ في $6\text{--}30\,\text{s}$، والأرضية $\approx(6\text{--}9)\cdot10^{-4}$ مقابل $7.5\cdot10^{-4}$ عند أفضل الأوزان؛ والأوزان عند $\eta=50$ لا تزال تختلف عن الأفضل بما يصل إلى $0.2$ على امتداد الاتجاهات السيئة التكييف لـ$C$ (التمرين~\ref{ex:15:2}).}
\answer{ex:15:7}{$\mathbf B=A\mathbf K$؛ $q_f/R\approx0.035$، $q_s/R\approx8.6\cdot10^{-4}$؛ وعند $4q_s$: $B_f\approx0.088$، $B_s\approx0.047$: العملية البطيئة تتعلم أسرع بأكثر من الضعف، والسريعة أبطأ.}

\answer{ex:16:1}{الدم دارة RC بـ$\tau=t_{1/2}/\ln2=14.4$ دقيقة؛ نسبة العظمى إلى الصغرى $2^{T/t_{1/2}}=64$؛ التوافقية الأساسية توهَّن إلى $0.55$؛ والنسبة $2$ عند $t_{1/2}=T=60$ دقيقة.}
\answer{ex:16:2}{$\varphi^*=\arcsin\bigl((\omega-\omega_0)/K_\varphi\bigr)\approx41$ دقيقة ($\omega-\omega_0\approx0.047$ راديان/يوم)، وزمن الاسترخاء $1/(K_\varphi\cos\varphi^*)\approx3.9$ يوم؛ وبعد إزاحة $+6$ و$-6$ ساعات: $7.7$ و$7.9$ يوم.}
\answer{ex:16:3}{$K_c(n)=\sec^n(\pi/n)$: $8$ و$2.37$، والخطأ $11\%$ و$30\%$؛ وعندما $n\to\infty$ يكون $K_c\to1$ والخطأ $\to50\%$.}
\answer{ex:16:4}{الأمران الأقصيان (المسرِّع $80$ والمكبح $0$) يعطيان $f=120$ بعد $0.76$ ثانية؛ و”المسرِّع وحده“ بعد $2.33$ ثانية، أي أطول بثلاثة أضعاف.}
\answer{ex:16:5}{$3\arctan(\omega\tau)+\omega d=\pi$، $K_c=(1+\omega^2\tau^2)^{3/2}$: $K_c=8$ و$3.54$ و$2.50$ و$1.76$، والدور $3.63$ و$5.46$ و$6.86$ و$9.27\,\tau$؛ عند $0.9K_c$ تخمد التذبذبات، وعند $1.1K_c$ تستقر على سعة ثابتة.}
\answer{ex:16:6}{$0.60\bar c$ و$0.87\bar c$ و$0.90\bar c$، والنهاية $0.91\bar c$؛ وعند تركيز ثابت تنعدم الاستجابة: هذا المتلقّي لا يقرأ إلا الدفعات.}
\answer{ex:16:7}{مسألة مفتوحة. مع التغذية الراجعة لا يوجد إيقاع إلا عند $K>8$، والدور متناسب مع $\tau$ المراحل ويكاد لا يعتمد على $K$ ($3.64\,\tau$ عند $K=8.5$، و$4.23\,\tau$ عند $K=40$): تغيير $K$ مرة ونصفًا يزيح الدور بنسبة $2\text{--}4\%$، وهي أقل من خطأ القياس. الحاسم هو فتح الحلقة (هرمون أخير ثابت عند مدخل المرحلة الأولى يقضي على هذا الإيقاع) أو الاستجابة لجرعة قصيرة إضافية من الهرمون في أطوار مختلفة من الدورة.}

\answer{ex:17:1}{$0.2\,\text{W}$ مقابل $81\,\text{\textmu W}$؛ والتدفق الخام يحتاج إلى $50\,\text{pJ}$ للبت على الأكثر.}
\answer{ex:17:2}{عند $c=0.1$: $5.6$ و$8.7$ و$9.2$ بت/ثانية، والنهاية $9.3$ بت/ثانية؛ وعند $c=0$ و$N=1000$: $65$ بت/ثانية.}
\answer{ex:17:3}{$B=\log_2N+P\log_2P+(1-P)\log_2\frac{1-P}{N-1}=4.87$ و$4.37$ و$3.29$ بت للحرف، أي $7.3$ و$6.6$ و$4.9$ بت/ثانية؛ ولبلوغ $10$ بت/ثانية يلزم $2.3$ حرف في الثانية.}
\answer{ex:17:4}{تباين الخطأ $2b/(Nm^2T)+\sigma_\xi^2$ لكل مركّبة؛ يتساوى الإسهامان عند $N\approx90$؛ والنهاية $\tfrac12\log_2(1+0.25/0.04)\approx1.4$ بت للمركّبة في النافذة.}
\answer{ex:17:5}{يتقارب الزوج عند $\eta_bd^2+\eta_db^2<2$، أي تقريبًا عند $\eta<1$: عند $0.8$ يتقارب، وعند $1.1$ تذبذبات مستقرة بدور $2$، وعند $1.5$ تذبذبات غير دورية، وعند $2$ تباعد؛ يجب الفصل بين سرعتي المتعلمَين.}
\answer{ex:17:6}{العتبة $\approx4.4\sigma$؛ $102\,\text{kbit/s}$، و$0.10\,\text{mW}$ مقابل $205\,\text{mW}$ للتدفق الخام، أي أقل بـ$2000$ مرة؛ ولا يتشبّع (أكثر من $3$ نبضات) إلا $5.7\cdot10^{-5}$ من العيّنات.}
\answer{ex:17:7}{مسألة مفتوحة. نحو $46$ بت/ثانية عند $\text{SNR}\approx0.9$ ونحو $330$ بت/ثانية مع الضجيج المستقل. إن كان العائق ضجيجًا يشبه الإشارة، فإن المعلومات المستخرجة تتشبع مع ازدياد عدد القنوات؛ وإن كان الجهل بالشفرة، فإنها تزداد مع مفكِّك أفضل (مثلًا مفكِّك يراعي توقيت النبضات).}

\answer{ex:18:1}{(أ)~$0.9949\le w\le1.0049$، أي $|1-w|\lesssim0.005$. (ب)~$K\ge400$ مشبك.}
\answer{ex:18:2}{$(\omega_R\tau_w)^2=\sqrt{\alpha(\alpha+2+2\varrho)}/\varrho-1$، $\varrho=\tau_m/\tau_w$؛ $f_R\approx11.1\,\text{Hz}$؛ $|Z(\omega_R)|/|Z(0)|=4.6$.}
\answer{ex:18:3}{(أ)~$f=1/\bigl(\tau\ln\frac{\mu}{\mu-\theta}\bigr)$؛ ولـ$1\,\text{Hz}$ يلزم $\mu/\theta-1\approx3.7\cdot10^{-44}$. (ب)~$T=\pi\tau/\sqrt\mu$، $f\approx3.2\,\text{Hz}$: $f\propto\sqrt\mu$.}
\answer{ex:18:4}{المكامل يعمل عند $f\lesssim17.7\,\text{Hz}$ (مرشح تمرير منخفض)، والرنّان في النطاق $5.5\text{--}22\,\text{Hz}$ فقط (مرشح تمرير نطاقي).}
\answer{ex:18:5}{(أ)~$T_{\min}=\frac{\tau}{w-1}\ln\frac{0.5}{0.3}\approx10.2\ms$. (ب)~$B>w-1-0.2=0.3$.}
\answer{ex:18:6}{الثابت الزمني للضبط $\tau/(\eta\langle r^2\rangle)$، $\eta=\tau\ln10/60\,\text{s}\approx3.8\cdot10^{-3}$؛ وعند $I\neq0$ ينزاح الوزن، فيجب أن تُطرح من $e$ نسخة الأمر $I/\tau$.}
\answer{ex:18:7}{مسألة مفتوحة. يتفوق الرنّان بـ$1.35$ مرة فقط عند $11\,\text{Hz}$ ويخسر $17$ مرة عند $1\,\text{Hz}$؛ والمكسب الرئيسي من إعادة التهيئة في الاختيار العتبي للنطاق (التمرين~\ref{ex:18:4}).}

\answer{ex:19:1}{(أ)~المشبك الواحد يعطي $6.7\mV$ ($R=66.7\,\text{M}\Omega$)، فبلوغ العتبة أصلًا يحتاج إلى $4$ على الأقل (الثلاثة لا تكفي إلا تقاربيًا)؛ و$8$ في $10\ms$؛ و$14$ في $5\ms$. (ب)~$0.1\ms\ll\tau_m$، وتصحيح التسريب $\sim0.25\%$.}
\answer{ex:19:2}{(أ)~$1.75\cdot10^{10}$ (ربع المازجات يستجيب لكل شيء)، و$4.4\cdot10^9$، و$0.06$ (عند $k=40$ لا يستجيب تقريبًا أبدًا). (ب)~بـ$10^3$ مرة: $2\cdot10^5$ مقابل $200$.}
\answer{ex:19:3}{(أ)~$C(2n,n)=2^{2n-1}$ بتناظر معاملات ذات الحدين. (ب)~$0.93$ و$0.09$؛ عرض الانتقال $\sim\sqrt{2n}$ في $P$، والعرض النسبي $\sim1/\sqrt n$.}
\answer{ex:19:4}{من تغصن واحد $V_{\text{soma}}<\kappa E_E<\theta$ مهما كان عدد المداخل؛ وعند $g_0=0.5g_L$ يلزم $n\ge3$ على كل تغصن.}
\answer{ex:19:5}{$I\ge C\sigma_V/\sigma_t=15\,\text{nA}$، أي $150$ مشبكًا متزامنًا، وهذا غير واقعي؛ أما العصبون الصغير فيكفيه $1\,\text{nA}$، ومع مشبك $4\,\text{nA}$ يكون الارتعاش $5\,\text{\textmu s}$.}
\answer{ex:19:6}{القمة عند الطرف القريب: $0.03\,\tau_m$ و$0.97$ ($L=0.2$)؛ $0.32\,\tau_m$ و$0.66$ ($L=1$)؛ $0.79\,\tau_m$ و$0.33$ ($L=2$). التقدير~\eqref{eq:dendriteload} بالسعة الكاملة لا يصح إلا عند $L\ll1$.}
\answer{ex:19:7}{مسألة مفتوحة. عند $m$ ثابت يزداد زمن الاستجابة خطيًا مع $P$ المحفوظ؛ وعند نسبة ثابتة $m=\varphi n$ يكفّ عن الاعتماد على $n$، فبطء العصبون الكبير نتيجة جمع مداخل كثيرة، لا نتيجة الحجم بحد ذاته.}

\answer{ex:20:1}{$t_{\text{w}}=\tau_r\ln\frac{1-L}{1-H}=16.8$ ساعة، $t_{\text{s}}=\tau_d\ln\frac HL=5.8$ ساعة، والدور $22.5$ ساعة؛ ويشدّ المولّدَ إلى $24$ ساعة التأرجحُ اليومي للعتبات، أي حلقة مقفلة الطور~\eqref{eq:pll}.}
\answer{ex:20:2}{(أ)~$L=0.111$؛ ومع $L=0.092$ كان النوم سيدوم $9.6$ ساعة. (ب)~بـ$22\%$ ($1/0.82=1.22$)، لا بـ$18\%$.}
\answer{ex:20:3}{$H=\frac{p-1}{p-1/q}=0.623$، $L=H/q=0.093$، حيث $p=e^{16/\tau_r}$، $q=e^{8/\tau_d}$. بعد الإيقاظ عقب $4$ ساعات ينزاح الطور إلى الأبد (النوم أبكر بـ$7.2$ ساعة)؛ ومع العتبات المتأرجحة يعود في يومين أو ثلاثة.}
\answer{ex:20:4}{(أ)~$r_\pm=\frac12\bigl(1\pm\sqrt{1-4k/w}\bigr)=0.947,\ 0.053$؛ $a_{\text{hi}}=3.51$، $a_{\text{lo}}=0.49$. (ب)~العمل $\approx0.33$ ثانية، والصمت $\approx0.59$ ثانية، والدور $\approx0.93$ ثانية، ونسبة العمل $0.36$.}
\answer{ex:20:5}{$a_{\text{hi}}/r_+<b<a_{\text{lo}}/r_-$: $3.71<b<9.20$. يخفض \nt{ACET} قيمة $b$ إلى ما دون $3.71$: فينتقل التقاطع إلى الفرع العلوي، وتستقر حالة العمل ($r\approx1$).}
\answer{ex:20:6}{$4.9$ و$6.8$ و$8.8$ و$5.5$ و$8.9$ ساعة عند $D=24$ و$32$ و$40$ و$48$ و$64$: المدة يحددها الطور اليومي لبدء النوم، لا الدَّين.}
\answer{ex:20:7}{بعد B يكون الخطأ على A في المتوسط $0.51$ (من $0.19$ إلى $1.08$ على $20$ مجموعة؛ وللاستجابة الصفرية $1$). ومع الخلط يكون الخطآن كلاهما أقل من $10^{-4}$.}
\answer{ex:20:8}{مسألة مفتوحة. التحجيم المنتظم يحفظ نسب الأوزان، لكنه لا يحفظ استجابات العصبون العتبي إلا إذا حُجّمت العتبة بالنسبة نفسها؛ أما الإعادات المخلوطة فتغيّر النسب. وثبات التماسات الكبيرة يناقض التحجيم المنتظم الخالص.}

\answer{ex:21:1}{(أ)~$T=\tau_{\text{rec}}\ln\frac{1}{1-x^*}$: $1.84\,\text{s}$ و$3.68\,\text{s}$. (ب)~من $3.36$ إلى $4.24\,\text{s}$: الحساسية $\dd T/\dd x^*=\tau_{\text{rec}}/(1-x^*)=80\,\text{s}$، وعند $x^*\to1$ تصير الرشقات الشبكية غير منتظمة.}
\answer{ex:21:2}{(أ)~$8.6\,\text{W}$، كما في~\eqref{eq:energy}. (ب)~$205\,\text{Mbit/s}$، أي $0.2\,\text{W}$، وهذا أكثر من قدرة المزرعة ($10^{-4}\,\text{W}$) بـ$2\cdot10^3$ مرة.}
\answer{ex:21:3}{$x_{\text{ss}}=1/(1+Ur_b\tau_{\text{rec}})<x^*$ عند $r_b>(1/x^*-1)/(U\tau_{\text{rec}})=0.28\,\text{Hz}$ لـ$x^*=0.9$ و$0.025\,\text{Hz}$ لـ$x^*=0.99$؛ وتنخفض قوة المشبك إلى $x_{\text{ss}}$، أي بنسبة $10\%$ و$1\%$.}
\answer{ex:21:4}{$u(t-k)$ متعامدة معيارية، و$m_k=\|P_Vu(t-k)\|^2$، حيث $P_V$ الإسقاط على الفضاء المولَّد بالمخارج ذي البعد $N$؛ وبمتراجحة بِسل $\sum_k m_k\le N$.}
\answer{ex:21:5}{$\mathrm{MC}\approx9\text{--}11$ مع $\tanh$ و$15\text{--}18$ للخزّان الخطي (ثلاث بذور عشوائية)، وكلاهما $\le30$: اللاخطية تُدفع ذاكرةً.}
\answer{ex:21:6}{(أ)~$n\le102$. (ب)~بـ$5.6$.}
\answer{ex:21:7}{مسألة مفتوحة. التجربة الضابطة الأساسية: جمّد القراءة وتحقق هل يبقى التحسن بعد توقف بلا تحفيز؛ فإن لم يبقَ فالذي تغيّر هو الحالة لا الأوزان.}

\answer{ex:22:1}{(أ)~$32\%$ و$100\%$. (ب)~$R_\uparrow-R_\downarrow\ge2\sqrt{2000\,\text{Hz}/1\,\text{s}}\approx89\,\text{Hz}$، أي $4.5\%$ فقط من التردد الكلي.}
\answer{ex:22:2}{(أ)~$n\ge4(p_1q_1+p_2q_2)/\Delta^2\approx560$. (ب)~$\approx56$.}
\answer{ex:22:3}{(أ)~التوازن التفصيلي: $\rho PK$ متناظرة. (ب)~نسبة الإخفاقات تساوي المتوسط التوافقي $1/\langle1/P\rangle=0.18$، مقابل $0.5$ بلا تغذية راجعة.}
\answer{ex:22:4}{(أ)~متوازي أضلاع مساحته $4h^2=0.04$ (ومع مراعاة حصر $p$ تبلغ عدديًا $0.045$). (ب)~$\approx0.18$.}
\answer{ex:22:5}{$0.49$: بعض المزارع يذهب إلى منطقة يكاد الإخفاق فيها يكون دائمًا، ويهيم فيها. ومع الحصر $1.00$: على مجموعة محدودة تتجمع الكثافة $\propto1/P$ في المنطقة الماصّة.}
\answer{ex:22:6}{(أ)~يلزم $\ge5$ مستويات؛ و$8$ أقطاب تكفي. (ب)~لا: يلزم $r_{\max}\ge25/T=100\,\text{Hz}$.}
\answer{ex:22:7}{مسألة مفتوحة. زمن البحث العشوائي يزداد مع $d$ تقريبًا كنسبة الحجوم $(L/h)^d$، أما البحث بعلامة الخطأ فخطيًا في $d$. والتجربة الفاصلة تجربة التبديل: بعد الإخفاق منبه متوقَّع لكنه قوي، وبعد الإصابة منبه غير متوقَّع لكنه ضعيف؛ ويلزم بضع عشرات من المزارع في كل مجموعة.}


\chapter{الملحق التجريبي}
\label{sec:supplement}

جمعنا هنا لكل فصل أقواله الرئيسية وما تستند إليه: التجربة التي قيس فيها ذلك، وما الذي قيس بالضبط، وأين نُشر. والحالة في العمود الأيسر تبيّن مدى رسوخ القول: \emph{مُقاس}: هناك تجربة مباشرة؛ \emph{نظرية}: نتيجة رياضية مشتقة في الفصل أو في دراسة كلاسيكية؛ \emph{نموذج}: نتيجة حساب بنموذج الكتاب (النصوص البرمجية في المجلد \foreignlanguage{english}{\texttt{models/}})؛ \emph{تقدير}: رقم من حيث رتبة المقدار بلا قياس مباشر؛ \emph{فرضية}: تفسير معقول لم تثبته التجربة بعد. وحيث لا تجربة، قيل ذلك صراحة.

\section{الفصل~\ref{sec:neurontypes}. أنواع العصبونات}

تستند أرقام الفصل إلى عدّ مباشر للخلايا في دماغ الإنسان حيثما وُجد هذا العدّ، وتستند قاعدة ”عصبون واحد، ناقل واحد“ إلى الأطالس الخلوية وإلى تجارب وجدت استثناءاتها أيضًا، ويستند دور \nt{DOPA} إلى تسجيلات النبضات، أما أدوار قنوات البثّ الأخرى فتستند إلى فرضيات.

{\footnotesize
\setlength{\tabcolsep}{3pt}\renewcommand{\arraystretch}{1.15}
\begin{longtable}{>{\raggedright}p{0.5cm}>{\raggedright}p{4.0cm}>{\raggedright}p{5.6cm}>{\raggedright}p{2.3cm}>{\raggedright\arraybackslash}p{1.7cm}}
\toprule
م & القضية & التجربة وما قيس & المصدر & الحالة \\
\midrule\endfirsthead
\toprule
م & القضية & التجربة وما قيس & المصدر & الحالة \\
\midrule\endhead
1 & في دماغ الإنسان نحو $8.6\cdot10^{10}$ عصبون، وكل عصبونات \nt{GLUT} تقريبًا عصبونات صغيرة في المخيخ (الجدول~\ref{tab:types}) & المجزِّئ المتماثل الخواص: يُذاب النسيج إلى معلّق متجانس من النوى، وتُعدّ النوى الحاملة للواسم العصبوني \foreignlanguage{english}{NeuN}. عند الرجال البالغين $86.1\pm8.1$ مليار عصبون؛ منها $19\,\%$ فقط في قشرة نصفي الكرة المخيين، والكتلة الرئيسية في المخيخ & \cite{Azevedo2009} & مُقاس \\
2 & لكل عصبون تقريبًا ناقل عصبي رئيسي واحد (القضية~\ref{prop:dale})، لكن ثمة استثناءات & أطلس ترنسكريبتومي لدماغ الفأر: نحو $4\cdot10^6$ خلية، و$5322$ عنقودًا؛ نُسب إلى كل عنقود ناقله بحسب جينات التخليق والتعبئة. والاستثناءات مقيسة مباشرة: عصبونات \nt{DOPA} في الشرائح تقذف الغابا أيضًا عبر الناقل الحويصلي نفسه وتثبط عصبونات المخطط؛ وفي جزء من محاور \nt{DOPA} يخرج الغلوتامات والدوبامين من نهايات مختلفة للسلك نفسه (المجهر الإلكتروني وعلم البصريات الوراثي) & \cite{Strata1999,Yao2023,Tritsch2012,Zhang2015} & مُقاس \\
3 & العمود كله في مصفوفة الأوزان بعلامة واحدة~\eqref{eq:dalesign} & استنتاج في الفصل من القضية~\ref{prop:dale} ومن أن مستقبلات الغلوتامات كلها تقريبًا مثيرة، ومستقبلات الغابا مثبطة. لا يوجد اختبار مباشر لقاعدة ”كل مستقبِلي العصبون الواحد يتلقون وزنًا بعلامة واحدة“ بوصفها قاعدة عامة؛ ومن الأمثلة المضادة القذف المشترك للغابا من عصبونات \nt{DOPA} (السطر~2) والمستقبِلون الحاملون لمستقبلات مختلفة العلامة (السطر~11) & استنتاج في الفصل & نظرية \\
4 & من ثلاثة أرباع إلى أربعة أخماس عصبونات القشرة من \nt{GLUT}، ومن الخُمس إلى الربع من \nt{GABA} & صبغ مناعي للغابا في أعمدة عرضها $50\um$ تخترق سماكة القشرة كلها في $10$ مناطق من قشرة القرد: في $8$ مناطق تشكل عصبونات الغابا نحو $25\,\%$ من كل العصبونات، وفي القشرة البصرية الأولية أقل من $20\,\%$ & \cite{Hendry1987} & مُقاس \\
5 & لعصبون \nt{GLUT} نحو $10^4$ مخرج & علم القياس المجسم عند التشريح: في القشرة الحديثة للرجال الشباب $1.64\cdot10^{14}$ مشبك (المجهر الإلكتروني، طريقة الشريحتين) ونحو $2.3\cdot10^{10}$ عصبون. والنسبة تعطي $\sim7\cdot10^3$ مشبك لكل عصبون؛ وفي المتوسط يساوي عدد المداخل عدد المخارج & \cite{Tang2001synapses,Pakkenberg1997} & مُقاس \\
6 & يتفرع مخرج عصبون \nt{DOPA} إلى $10^5\text{--}10^6$ نهاية؛ وهذا تقدير من مدى تغطية الهدف لا عدّ & وسم فيروسي لغشاء عصبونات \nt{DOPA} منفردة لدى الجرذ وإعادة بناء كاملة للمحور: العصبون الواحد يغطي من $0.45$ إلى $5.7\,\%$ (في المتوسط $2.7\,\%$) من حجم المخطط. قيست التغطية لا عدد النهايات: فهذا مستنتج من التغطية من حيث رتبة المقدار، ولا عدّ مباشرًا للنهايات. وأعداد النهايات لـ\nt{SERO} و\nt{NORA} تقديرات خشنة & \cite{Matsuda2009} & مُقاس \\
7 & عصبونات البثّ قليلة: \nt{DOPA} $\sim5\cdot10^5$ (المجموعة الرئيسية من مجموعتين، حدّ أدنى)، \nt{SERO} $\sim3\cdot10^5$ (مجموع عدّين جزئيين)، \nt{HIST} $\sim6\cdot10^4$ (الجدول~\ref{tab:types}، الشكل~\ref{fig:typepie}) & عدّ مجسم عند الإنسان: $550\,000$ عصبون مصطبغ في المادة السوداء (من دون السقيفة البطنية)؛ و$235\,000$ عصبون في نواة الرفاء الظهرية (كل عصبونات النواة)، و$125\,000$ عصبون سيروتونيني آخر في سقيفة الجسر؛ ونحو $64\,000$ عصبون هستاميني في الوطاء. لا عدّ مباشرًا لـ\nt{NORA} و\nt{GLYC} و\nt{ACET} و\nt{ADRE}: فهي في الجدول تقديرات خشنة لرتبة المقدار & \cite{Pakkenberg1991,Baker1990,Baker1991,Airaksinen1991} & مُقاس \\
8 & يقفز الغلوتامات في الشق إلى نحو ميلي مول ويهبط في $\sim1\ms$ & بإزاحة حاصر تنافسي سريع الانفصال عن مستقبلات \foreignlanguage{english}{NMDA} أثناء النقل المشبكي (مزرعة عصبونات الحُصين): الذروة $1.1\mM$، وثابت الهبوط $1.2\ms$ & \cite{Clements1992} & مُقاس \\
9 & في القشرة العاملة يكاد التياران المثير والمثبط يتوازنان، ويبقى العصبون قرب العتبة & النظرية: الشبكة ذات الوصلات القوية تبلغ النمط المتوازن من تلقاء نفسها. التجربة: في القشرة الجبهية الأمامية لابن مقرض في الكائن الحي، أثناء حالات ”الصعود“، يبقى جهد انعكاس التيار المشبكي قرب $-37\mV$ مئات من أجزاء الألف من الثانية، مع أن الموصلية تتغير بـ$\sim21\,\text{nS}$: الإثارة والتثبيط يرتفعان وينخفضان معًا & \cite{vanVreeswijk1996,Haider2006} & مُقاس \\
10 & عصبونات \nt{DOPA} تبلّغ عن خطأ التنبؤ بالمكافأة~\eqref{eq:rpe} (الشكل~\ref{fig:rpe}) & نبضات عصبونات \nt{DOPA} لدى القرد: رشقة للعصير غير المتوقع؛ وبعد التعلم رشقة للإشارة ولا استجابة للعصير المتنبأ به؛ وهبوط في التردد حين لا يُعطى العصير الموعود. وفي تجربة كمية يتناسب التردد مع الفرق بين المكافأة الحالية والمتوسط الموزون للمكافآت السابقة، لكن للأخطاء الموجبة فقط. أما المعادلة~\eqref{eq:rpe} نفسها فهي خوارزمية الفروق الزمنية & \cite{Schultz1997,Bayer2005,SuttonBarto2018} & مُقاس \\
11 & العلامة والسرعة يحددهما المستقبِل: المستقبلات مؤينة التوجه في أجزاء من جزء الألف من الثانية، وأيضية التوجه أبطأ بكثير (القضية~\ref{prop:receptor}) & نبضات سريعة من الغلوتامات على قطع من غشاء عصبونات الحُصين: التيار عبر المستقبلات مؤينة التوجه يرتفع (من $20$ إلى $80\,\%$) في $0.2\text{--}0.6\ms$ ويهبط في $2\text{--}3\ms$. والجهد المثبط البطيء للعصبونات الهرمية يزيله حاصر انتقائي للمستقبل أيضي التوجه للغابا. عائلتا مستقبلات الدوبامين متعاكستا التأثير؛ وفي المخطط تحتاج العصبونات ذات المستقبل \foreignlanguage{english}{D1} إلى الدوبامين لتقوية المشابك، وذات \foreignlanguage{english}{D2} لإضعافها & \cite{Colquhoun1992,DutarNicoll1988,Beaulieu2011,Shen2008} & مُقاس \\
12 & قناة البثّ ليست سلكًا واحدًا، بل حزمة من خطوط قليلة (القسم~\ref{sec:buses}) & وسم فيروسي-وراثي لعصبونات السيروتونين في نواة الرفاء الظهرية للفأر: العصبونات التي ترسل مخرجها إلى اللوزة وتلك التي ترسله إلى القشرة الجبهية تقع في مواضع مختلفة من النواة، وتتفرع إلى مناطق متتامة، وتستجيب استجابتين متعاكستين لمنبه مزعج. مراجعة: مجموعات فرعية من عصبونات الموضع الأزرق (\nt{NORA}) ترمّز عمليات مختلفة & \cite{Ren2018,Poe2020} & مُقاس \\
13 & \nt{NORA} يحدد الكسب $g$ & فرضية الكسب التكيفي؛ ونموذج شبكة ترفع فيه الكاتيكولامينات ميل منحنى النقل. لا قياس مباشرًا لـ”ازدياد $g$ مع النورأدرينالين“ في الشبكة كلها؛ وقد قيس أن قطر الحدقة يتبع نبضات عصبونات الموضع الأزرق لدى القرد & \cite{AstonJones2005,ServanSchreiber1990,Joshi2016} & فرضية \\
14 & \nt{ACET} يبدّل ”كتابة/قراءة“ ويحدد سرعة التعلم & فرضية. سندها: في شرائح القشرة الشمية للجرذ تضعف ناهضات الأسيتيل كولين ($100\,\mu\text{M}$) المشابك الداخلية ”المستحضِرة“ بأكثر من $60\,\%$، والمشابك الداخلة بأقل من $15\,\%$ & \cite{Hasselmo2006,HasselmoBower1992} & فرضية \\
15 & \nt{SERO} يحدد معامل الخصم $\gamma$؛ والعصبونات السيروتونينية تعمل بانتظام، ببضع نبضات في الثانية، وتكاد تصمت في إحدى مراحل النوم & فرضية ”السيروتونين $=\gamma$“. أقوى حجة: التنشيط البصري الوراثي لعصبونات السيروتونين في نواة الرفاء الظهرية للفأر يقلل عدد حالات التخلي المبكر عن انتظار مكافأة مؤجلة. وتردد النبضات والصمت في نوم حركة العين السريعة مقيسان (مراجعة للتسجيلات) & \cite{Doya2002,Miyazaki2014,JacobsAzmitia1992} & فرضية \\
\bottomrule
\end{longtable}}

\section{الفصل~\ref{sec:glutcomputer}. ماذا تحسب عصبونات \nt{GLUT}}

القضايا المنطقية في الفصل مبرهنات عن دارات ذات أوزان غير سالبة، مستنتجة في الفصل نفسه؛ وتؤكد التجربة نموذج العصبون، وأن الدارات المبنية وفق هذه المبرهنات (كاشف التزامن، وخط التأخير، والمجموعة الذاتية الاستمرار، والسلسلة) موجودة فعلًا في الدماغ.

{\footnotesize
\setlength{\tabcolsep}{3pt}\renewcommand{\arraystretch}{1.15}
\begin{longtable}{>{\raggedright}p{0.5cm}>{\raggedright}p{4.0cm}>{\raggedright}p{5.6cm}>{\raggedright}p{2.3cm}>{\raggedright\arraybackslash}p{1.7cm}}
\toprule
م & القضية & التجربة وما قيس & المصدر & الحالة \\
\midrule\endfirsthead
\toprule
م & القضية & التجربة وما قيس & المصدر & الحالة \\
\midrule\endhead
1 & العصبون دارة RC بعتبة وتصفير~\eqref{eq:glutneuron} & حُقن في جسم العصبونات الهرمية في قشرة الجرذ (شرائح) تيار ضجيجي يشبه التدفق المشبكي في الدماغ الحي. علاقة التردد المتوسط بمتوسط التيار وتشتته يصفها العصبون المكامل التسريبي إذا أضيف إليه تكيّف بطيء في التردد. أما مجالات $\tau$ والتأخيرات فمذكورة في الفصل بلا إحالة إلى تجربة & \cite{Rauch2003} & مُقاس \\
2 & النموذج~\eqref{eq:glutneuron} تبسيط: فروع المداخل نفسها تطلق نبضات محلية & تسجيل مباشر من تغصنات العصبونات الهرمية في القشرة البصرية للفأر في الكائن الحي: المنبه البصري يستدعي نبضات تغصنية محلية تعتمد على مستقبلات \foreignlanguage{english}{NMDA}؛ وكبتها يوسّع ضبط العصبون على الاتجاه & \cite{Smith2013} & مُقاس \\
3 & نافذة التزامن $\Delta_{\max}=\tau\ln\frac{w}{\theta-w}$~\eqref{eq:window}؛ وعند $\theta=1.5w$ تساوي $0.7\tau$ & نتيجة للنموذج~\eqref{eq:glutneuron}. يوجد قياس مباشر لنافذة جمع ضيقة في عصبونات ذات تثبيط سريع (الفصل~\ref{sec:gaba}، السطر~3 من جدوله) & استنتاج في الفصل & نظرية \\
4 & شبكة من عصبونات \nt{GLUT} وحدها لا تحسب إلا دوال رتيبة في المداخل (القضية~\ref{prop:monotone}) & البرهان في الفصل: تكرار من الصمت بطرف أيمن غير متناقص. هذه قضية عن النموذج؛ ولا تجربة تختبرها على شبكة حية بلا تثبيط & استنتاج في الفصل & نظرية \\
5 & أعضاء الحس تعطي زوجًا من الأسلاك: بعض الخلايا يستجيب لاشتداد المنبه، وبعضها لضعفه & الشبكية: منذ الخلايا ثنائية القطب تتفرع إشارة المخاريط إلى قناتي تشغيل وإطفاء. الحجب الدوائي لخلايا التشغيل ثنائية القطب لدى القرد: القناتان تسيران منفصلتين حتى القشرة؛ وبعد الحجب يسوء اكتشاف اشتداد الضوء، لا خفوته & \cite{Schiller1992} & مُقاس \\
6 & بالشيفرة ثنائية السلك تحسب شبكة من عصبونات \nt{GLUT} أي دالة منطقية لا تزامنيًا، بلا إيقاع مشترك، بثمن مضاعفة عدد العصبونات (القضية~\ref{prop:dualrail}، \eqref{eq:dualrail}، الشكل~\ref{fig:dualxor}) & الشيفرة ثنائية السلك وتحديد الجاهزية من الإشارة نفسها تقنية معيارية في الدارات اللاتزامنية غير الحساسة للتأخير. ودارة \foreignlanguage{english}{XOR} من ستة عصبونات مبنية في الفصل. ولا تجربة تُظهر أن الدماغ يحسب الدوال المنطقية بالشيفرة ثنائية السلك تحديدًا & \cite{Sparso2001} & نظرية \\
7 & أكبر فرق في زمن وصول الصوت إلى الأذنين عند الإنسان $\approx0.6\ms$، والدماغ يميّز نحو عشرة أجزاء من المليون من الثانية & مستمعون في تجربة بخيارين للإجابة: عتبة تمييز فرق الزمن (75\,\% إجابات صحيحة) هي $9\,\mu\text{s}$ لضجيج في النطاق $150\text{--}1700\,\text{Hz}$، و$11\,\mu\text{s}$ لنغمة $1\,\text{kHz}$، و$28\,\mu\text{s}$ لنقرة. والحد $0.6\ms$ حساب الفصل، $0.22\,\text{m}/343\,\text{m/s}$ & \cite{KlumppEady1956} & مُقاس \\
8 & لدى الطيور يتحول فرق الأزمنة إلى موضع عبر خطوط التأخير وكواشف التزامن~\eqref{eq:itd} (الشكل~\ref{fig:itd}) & بومة الحظائر: تدخل المحاور الآتية من الأذنين نواة كواشف التزامن من جهتين متقابلتين؛ ويتغير زمن وصول النبضات على طول المحور بانتظام مع العمق، ومدى التأخيرات عبر النواة ($160\,\mu\text{s}$ على $700\um$) يطابق مدى التأخيرات بين الأذنين & \cite{Carr1990} & مُقاس \\
9 & لم يُعثر لدى الثدييات على صف من التأخيرات؛ وتحل المهمة نفسها كواشف تزامن بتثبيط محسوب التوقيت من عصبونات \nt{GLYC} & تسجيل في الكائن الحي في الزيتونة العلوية الإنسية لدى الجربوع: الاستجابات لا تشكل خريطة اتجاهات؛ وإيصال الغليسين وحاصره عبر ماصة دقيقة يُظهر أن التثبيط \emph{الغليسيني} المحسوب التوقيت بدقة يحدد مجال التأخيرات العامل. التثبيط هنا من \nt{GLYC} لا من \nt{GABA} & \cite{Brand2002,Grothe2010} & مُقاس \\
10 & المجموعة ذات التغذية الراجعة القوية قلّاب؛ والفعالية الذاتية الاستمرار تدوم ثوانيَ بعد المنبه (الشكل~\ref{fig:bistable}) & القشرة الجبهية الأمامية للقرد في مهمة حركة عين مؤجلة نحو نقطة محفوظة (تأخير $1\text{--}6\,\text{s}$): $87$ من $170$ عصبونًا مرتبطًا بالمهمة تغيّر ترددها أثناء التأخير، و$79\,\%$ منها يتوقف ذلك فيها على اتجاه النقطة المحفوظة. أما أن هذه الفعالية تحفظها تغذية راجعة بحالتين مستقرتين فنظرية & \cite{Funahashi1989,Wang2001} & مُقاس \\
11 & يُكتب البت إذا دام الوارد أكثر من $\approx0.7\tau$ (الشكل~\ref{fig:recurrentgroup}ب) & حل المعادلة $\tau\,\dd r/\dd t=-r+f(wr+I)$ بطريقة أويلر عند $w=1$ و$I=0.6$؛ ولا برنامج في \texttt{models/}. لا تجربة & حساب في الفصل & نموذج \\
12 & يمكن حفظ الذاكرة في تيسير المشابك، بلا نبضات تقريبًا & النظرية: الكالسيوم المتبقي في النهايات يحفظ المكتوب نحو ثانية. السند: في القشرة الجبهية الأمامية الإنسية لابن مقرض (تسجيل من عدة عصبونات معًا) توجد شبكة فرعية من العصبونات الهرمية متصلة بمشابك تيسيرية ذات أثر لاحق طويل. ومراجعة لأرصاد الفترات ”الصامتة“ أثناء الاحتفاظ بالذاكرة. أي الطريقتين تستعمل القشرة ومتى: مسألة مفتوحة & \cite{Mongillo2008,WangMarkram2006,Stokes2015} & فرضية \\
13 & تُمحى الذاكرة بالتعب: ينضب مخزون الناقل العصبي & تسجيلات مزدوجة لعصبونات هرمية في القشرة: تهبط استجابة المشبك مع النبضات المتواترة، وسرعة الهبوط يحددها احتمال القذف. أما أن هذا بالذات ما يطفئ المجموعة الذاتية الاستمرار فلم يُبيَّن مباشرة & \cite{Tsodyks1997} & مُقاس \\
14 & سلسلة المجموعات مؤقت يطلقه حدث؛ لدى الطيور المغردة تنطلق العصبونات بالتتابع الصارم & تسجيل عصبونات معرّفة في المركز الصوتي الأعلى لعصفور الزيبرا في النوم وأثناء الغناء: كل عصبون يرسل مخرجه إلى النواة الحركية يطلق رشقة قصيرة واحدة في لحظة واحدة دقيقة من الأغنية، وتنطلق العصبونات متتابعة & \cite{Hahnloser2002} & مُقاس \\
\bottomrule
\end{longtable}}

\section{الفصل~\ref{sec:gaba}. \nt{GLUT} و\nt{GABA} معًا}

القضايا المنطقية والديناميكية في الفصل مستنتجة فيه نفسه من التحليل الخطي وقانون أوم؛ وتُظهر التجربة أن الدارات التي تقوم عليها، أي الحظر مع التأخير، والتثبيط الأمامي التالي للإثارة، والتثبيت بالتثبيط، وإيقاع حلقة \foreignlanguage{english}{\nt{GLUT}--\nt{GABA}}، تعمل فعلًا في الدماغ.

{\footnotesize
\setlength{\tabcolsep}{3pt}\renewcommand{\arraystretch}{1.15}
\begin{longtable}{>{\raggedright}p{0.5cm}>{\raggedright}p{4.0cm}>{\raggedright}p{5.6cm}>{\raggedright}p{2.3cm}>{\raggedright\arraybackslash}p{1.7cm}}
\toprule
م & القضية & التجربة وما قيس & المصدر & الحالة \\
\midrule\endfirsthead
\toprule
م & القضية & التجربة وما قيس & المصدر & الحالة \\
\midrule\endhead
1 & عصبونات \nt{GABA} في القشرة من الخُمس إلى الربع & صبغ مناعي للغابا في $10$ مناطق من قشرة القرد: نحو $25\,\%$ من العصبونات في $8$ مناطق، وأقل من $20\,\%$ في القشرة البصرية الأولية. ومراجعة لتنوع العصبونات المثبطة في القشرة & \cite{Hendry1987,Markram2004} & مُقاس \\
2 & ما يفعله التثبيط يحدده إلى حد بعيد موضع الاستقرار: التغصن، أو الجسم، أو موضع ولادة النبضة، أو نهاية سلك عصبون آخر & مراجعة: أصناف عصبونات \nt{GABA} في القشرة تختلف بهدفها على المستقبِل. تسجيل متزامن من جسم عصبون هرمي وتغصنه: التثبيط الأمامي أقوى بكثير على الجسم منه على التغصنات. والتثبيط على نهاية سلك عصبون آخر في الحبل الشوكي يقلل قذف الناقل العصبي لخط دخل واحد (مراجعة للقياسات) & \cite{Markram2004,Pouille2001,Rudomin1999} & مُقاس \\
3 & قلب العلامة عبر وسيط يكلف تأخيرًا واحدًا، نحو جزء من الألف من الثانية؛ والتثبيط التالي للإثارة يضيّق نافذة التزامن & العصبونات الهرمية في حُصين الجرذ في الشرائح: التثبيط عبر وسيط واحد يصل عقب الإثارة المباشرة بتأخير قصير، والنافذة التي يُجمع فيها مدخلان مثيران حتى العتبة أقل من $2\ms$. وعلى التغصنات، حيث هذا التثبيط أضعف، النافذة أوسع & \cite{Pouille2001} & مُقاس \\
4 & الحظر $a\wedge\bar b$~\eqref{eq:veto} مع \foreignlanguage{english}{OR} وواحد ثابت تامّ وظيفيًا (القضية~\ref{prop:gabacomplete}) & البرهان في الفصل. نتيجة كلاسيكية: شبكة من العناصر العتبية ذات المداخل المثبطة تحقق أي تعبير منطقي. قضية عن النموذج، ولا تجربة & \cite{McCullochPitts1943} & نظرية \\
5 & الحظر مع التأخير يعطي كاشفًا لاتجاه الحركة بسرعة مثلى $v^*=\Delta x/d$~\eqref{eq:vstar} & شبكية الأرنب: تنشأ الانتقائية للاتجاه بتثبيط يطفئ الإثارة عند الحركة في الاتجاه ”الصفري“. وأظهر التسجيل المنفصل للمدخلين المثير والمثبط للخلايا الانتقائية أن الخلايا عديمة المحاور النجمية (\nt{GABA}) تثبطها تثبيطًا غير متناظر، بالزوائد المتجهة إلى الجهة ”الصفرية“ فقط. والصيغة $v^*$ استنتاج الفصل & \cite{Barlow1965,Fried2002} & مُقاس \\
6 & التثبيط التحويلي يقسم الجهد ولا يطرح منه~\eqref{eq:shunt} (الشكل~\ref{fig:shunt}) & قانون أوم لمقسّم من ثلاث موصليات مع جهد اتزان للكلور قرب الراحة؛ لعصبون لا يطلق نبضات & استنتاج في الفصل & نظرية \\
7 & يؤثر التحويل في تردد النبضات بالطرح، أما القسمة فتنتج عن التثبيط مع الضجيج الخلفي المتوازن & النموذج: في العصبون المنطلق يكاد التيار عبر التحويل يكون ثابتًا، والأثر في التردد طرحي. التجربة: أُدخلت في العصبونات الهرمية للقشرة الحسية الجسدية لدى الجرذ (شرائح) موصليات خلفية مثيرة ومثبطة عبر المشبك الديناميكي؛ والنمو المتوازن المتزامن للاثنتين يقسم ميل الاستجابة دون إزاحة ملحوظة للتردد. مراجعة: القسمة على الفعالية الكلية توجد في مناطق كثيرة من الدماغ & \cite{HoltKoch1997,Chance2002,Carandini2012} & مُقاس \\
8 & عند $\beta>1$ يختار التثبيط المتبادل فائزًا واحدًا (القضية~\ref{prop:wta}، \eqref{eq:wta}) & القيم الذاتية $-(1\pm\beta)/\tau$ استنتاج في الفصل. والترددان $0.73$ و$0.53$ عند $\beta=0.5$ نقطة ثابتة $r_{1,2}=(I_{1,2}-\beta I_{2,1})/(1-\beta^2)$؛ ولا برنامج في \texttt{models/}. ولا يورد الفصل تجربة مباشرة & استنتاج في الفصل & نظرية \\
9 & تصفير الذاكرة بإشارة عبر \nt{GABA}؛ ومجموعتان بتثبيط متبادل تشكلان ماسكًا & ينتج من الشكل~\ref{fig:bistable} ومن القضية~\ref{prop:wta}. لا تجربة & استنتاج في الفصل & نظرية \\
10 & اتزان حلقة \foreignlanguage{english}{\nt{GLUT}--\nt{GABA}} مستقر إذا كان التثبيط قويًا بما يكفي وسريعًا بما يكفي (القضية~\ref{prop:ei}) & نموذج الجمهرتين~\eqref{eq:ei}؛ والشرطان (i) و(ii) هما محدد مصفوفته وأثرها، مستنتجان في الفصل & \cite{WilsonCowan1972} & نظرية \\
11 & عند $w_{EE}=1.5$ و$w_{EI}w_{IE}=2$ يحفظ التثبيط السريع ($\tau_I=2\ms$) المستوى $E^*=0.67h$، والبطيء ($30\ms$) يؤرجح التذبذبات (الشكل~\ref{fig:ei}) & حل عددي لـ~\eqref{eq:ei}، البرنامج \texttt{figures/make\_ei.py} & حساب & نموذج \\
12 & تعمل القشرة في نمط التثبيت بالتثبيط؛ وبصمته: ادفع عصبونات \nt{GABA} فيهبط ترددها~\eqref{eq:isn} & تنبأت النظرية بالاستجابة المتناقضة. التجربة: تسجيلات داخل الخلية في القشرة البصرية الأولية للقط: عند كبت الاستجابة بمنبه محيط لا تزداد الموصلية المثبطة بل تهبط مع المثيرة. والإثارة البصرية الوراثية لعصبونات \foreignlanguage{english}{PV} لدى الفأر في القشرة البصرية والحسية الجسدية والحركية تخفض تردد العصبونات المثبطة، حتى من دون منبه وتحت تخدير خفيف. والمعامل $-1/3$ بأرقام الشكل~\ref{fig:ei} استنتاج الفصل & \cite{Tsodyks1997isn,Ozeki2009,Sanzeni2020} & مُقاس \\
13 & الحلقة ”\nt{GLUT} يثير \nt{GABA}، و\nt{GABA} يثبط \nt{GLUT}“ تولّد إيقاعًا سريعًا دوره $12\text{--}50\ms$ & التحفيز البصري الوراثي لعصبونات \nt{GABA} السريعة في القشرة الحسية الجسدية للفأر بترددات $8\text{--}200\,\text{Hz}$ يقوّي انتقائيًا إيقاع غاما ($20\text{--}80\,\text{Hz}$، الدور $12\text{--}50\ms$)؛ وتحفيز عصبونات \nt{GLUT} لا يقوّي إلا الإيقاعات الأبطأ & \cite{Cardin2009,Buzsaki2012} & مُقاس \\
\bottomrule
\end{longtable}}

\section{الفصل~\ref{sec:code}. شفرة مورس}

التقديرات الحدية في الفصل من نظرية المعلومات والحساب؛ وما قيس هو أين ينشأ الضجيج في القناة، وما سرعة عمل الدماغ، وكم تكلف النبضة، أما السؤال عن الشيفرة التي تستعملها القشرة فعلًا فيبقى مفتوحًا.

{\footnotesize
\setlength{\tabcolsep}{3pt}\renewcommand{\arraystretch}{1.15}
\begin{longtable}{>{\raggedright}p{0.5cm}>{\raggedright}p{4.0cm}>{\raggedright}p{5.6cm}>{\raggedright}p{2.3cm}>{\raggedright\arraybackslash}p{1.7cm}}
\toprule
م & القضية & التجربة وما قيس & المصدر & الحالة \\
\midrule\endfirsthead
\toprule
م & القضية & التجربة وما قيس & المصدر & الحالة \\
\midrule\endhead
1 & ترددات عصبونات \nt{GLUT} في القشرة من أجزاء النبضة إلى عشرات النبضات في الثانية، والعصبونات تعمل معظم الوقت قرب الحد الأدنى & تسجيل بماصة ملاصقة للخلية (اختيار العصبون لا يتوقف على فعاليته) في القشرة السمعية لجرذ مستيقظ: في كل لحظة يعطي ترددًا فوق $20$ نبضة في الثانية أقل من $5\,\%$ من العصبونات؛ وفي بعض العصبونات يقل التردد الخلفي عن $0.01$ في الثانية؛ وتوزيع الترددات لوغاريتمي طبيعي & \cite{Hromadka2008} & مُقاس \\
2 & سعة القناة النبضية $R_{\max}\approx r\log_2\frac{e}{r\Delta t}$~\eqref{eq:spikeentropy}، وتكاد كلها تكون في توقيت النبضات (القضية~\ref{prop:capacity}) & إنتروبيا سلسلة ثنائية من خانات طولها $\Delta t$. أرقام الفصل: $8$ بتات للنبضة عند $r=10$ في الثانية و$\Delta t=1\ms$، ونحو $5$ بتات عند $\Delta t=10\ms$، وأقل من $2$ بت في الثانية لعدّاد النبضات & \cite{MacKay1952} & نظرية \\
3 & العصبونات الحسية تنقل من بت واحد إلى بضعة بتات للنبضة، وفي أفضل الحالات حتى نصف الحد & استعادة المنبه من نبضات عصبونات حسية معروفة المنبه؛ خلاصة القياسات في الكتاب & \cite{Rieke1997} & مُقاس \\
4 & مولّد النبضات دقيق؛ والتوقيت الدقيق تحدده التغيرات السريعة في المدخل & عصبونات قشرة الجرذ في الشرائح: مع تيار متكرر ذي تذبذبات سريعة تتكرر النبضات بدقة أفضل من $1\ms$، ومع تيار ثابت تتبعثر اللحظات & \cite{Mainen1995} & مُقاس \\
5 & المشبك غير موثوق: أكثر من نصف النبضات لا يصل إلى المستقبِل بسبب عشوائية القذف & التحفيز الأدنى لمداخل العصبونات الهرمية في الحُصين (شرائح): أكثر من نصف النبضات لا يستدعي استجابة. استُبعدت تقلبات عتبة التحفيز، وأعطال التوصيل عند تفرعات المحور، وانخفاض حرارة التجربة؛ فيبقى القذف الاحتمالي & \cite{AllenStevens1994} & مُقاس \\
6 & تدفق النبضات في القشرة الحية يبدو عشوائيًا: الفترات مبعثرة كما في تدفق بواسون، وتباين عدد النبضات قريب من متوسطه؛ وجزء من ”الضجيج“ رسالة عن حركات الحيوان & تسجيلات في المنطقتين البصريتين \foreignlanguage{english}{V1} و\foreignlanguage{english}{MT} لقرد مستيقظ: معامل تغير الفترات نحو $1$، ونسبة تباين عدد النبضات إلى متوسطه كما في عملية عشوائية؛ ونموذج يجمع مداخل صغيرة مستقلة يعطي تشتتًا أصغر بكثير. وفي القشرة البصرية للفأر يُتنبأ بجزء كبير من التشتت من تسجيل الفيديو لحركات الحيوان نفسه & \cite{Softky1993,Stringer2019} & مُقاس \\
7 & ترميز التردد بطيء: الخطأ $1/\sqrt{rT}$~\eqref{eq:rateerror}، والدماغ يتعرف على الصورة في $\sim150\ms$ & الصيغة إحصاء بواسون، استنتاج الفصل. التجربة: قرر أشخاص هل في صورة غير مألوفة معروضة $20\ms$ حيوان؛ تفترق جهود الدماغ لـ”نعم“ و”لا“ بعد نحو $150\ms$. وتقسيم هذا الزمن إلى نحو عشر مراحل من $10\text{--}20\ms$ تقدير الفصل لا قياس & \cite{Thorpe1996} & مُقاس \\
8 & التوسيط على مجموعة يحده الارتباط: لا ينخفض الخطأ بأكثر من $1/\sqrt c$ مرة~\eqref{eq:correlated} (القضية~\ref{prop:pooling}) & أزواج عصبونات من المنطقة \foreignlanguage{english}{MT} لدى القرد مسجلة معًا: معامل ارتباط أعداد النبضات $0.12$ في المتوسط، ويبين التحليل النظري أن حتى هذا الارتباط يحد ربح المجموعة. النظرية: لا يضع الحد إلا الجزء من الارتباط الذي يشبه الإشارة. ونموذج الموقف المقابل: تردد مجموعة من $50\text{--}100$ عصبون يُقرأ في فترة واحدة بين نبضتين، $10\text{--}50\ms$، ولا يُستعمل توقيت النبضات المفردة في القشرة & \cite{Zohary1994,MorenoBote2014,ShadlenNewsome1998} & مُقاس \\
9 & يمكن كتابة العدد في زمن النبضة الأولى~\eqref{eq:latency}، أو في ترتيب نبضات مجموعة، أو في طور الإيقاع المحلي & شبكية السمندل: البنية المكانية لصورة معروضة لحظة مكتوبة في التأخيرات النسبية للنبضات الأولى للعصبونات المخرجة، شبه مستقلة عن التباين. خلايا المكان في حُصين الجرذ: عند عبور موضعها ”الخاص“ تنزاح النبضات أبكر فأبكر في دورة إيقاع ثيتا ($7\text{--}12\,\text{Hz}$)، بإزاحة من $100^\circ$ إلى $355^\circ$. إلى أي حد تستعمل القشرة توقيت النبضات: سؤال مفتوح (السطر~8) & \cite{Gollisch2008,OKeefe1993} & مُقاس \\
10 & أي شيفرة تُقرأ يقرره الثابت الزمني للمستقبِل~\eqref{eq:reader} (القضية~\ref{prop:reader})؛ وفي القشرة النشطة العصبونات أقرب إلى كواشف التزامن & الصيغة~\eqref{eq:reader} استنتاج الفصل. مراجعة لتسجيلات في الكائن الحي: في القشرة النشطة موصلية الغشاء أعلى عدة مرات منها في الراحة، والثابت الزمني أقصر تبعًا لذلك. أما أن القشرة تبدّل الشيفرة بهذه الطريقة تحديدًا فلم يُبيَّن مباشرة & \cite{Destexhe2003} & فرضية \\
11 & المشبك الناضب ينقل تغيّر التردد، أي $\Delta r/r$ في جوهره~\eqref{eq:depression} (الشكل~\ref{fig:depression}) & تسجيلات مزدوجة لعصبونات هرمية في القشرة: سرعة النضوب يحددها احتمال القذف؛ مع النضوب البطيء تعكس الاستجابة التردد، ومع السريع التزامنات. نموذج مبني على النضوب المقيس في قشرة الجرذ: التغيرات النسبية المتساوية في التردد على مداخل سريعة وبطيئة تعطي استجابة متساوية، أي أن المشبك ينقل $\Delta r/r$. والأرقام $1.7\to2.2$ و$6.7$ و$90\ms$ حساب الفصل & \cite{Tsodyks1997,Abbott1997depression} & مُقاس \\
12 & الرشقة ”شَرطة“: احتمال ضياعها $(1-p)^k$~\eqref{eq:burst} صغير، والتيسير يقلله أكثر & مراجعة: في المشابك غير الموثوقة تضيع النبضات المفردة كثيرًا، أما الرشقات فتنتقل بموثوقية بفضل التيسير؛ والرشقة تستدعي تغيرات طويلة الأمد في المشابك. أما أن الدماغ يقرأ الرشقة رمزًا مستقلًا ففرضية & \cite{Lisman1997,AllenStevens1994} & فرضية \\
13 & عصبونات \nt{GABA} السريعة تحدد الإيقاع و”علامات الترقيم“؛ وصنف آخر يثبط المثبِّطات ويفتح البوابة (القضية~\ref{prop:punctuation}) & مراجعة لخصائص أصناف عصبونات \nt{GABA} في القشرة. علم البصريات الوراثي: الإثارة الإيقاعية لعصبونات \nt{GABA} السريعة تستدعي إيقاع غاما، والاستجابة للمنبه تتوقف على طور الدورة. في القشرة البصرية للفأر تثبط عصبونات \foreignlanguage{english}{PV} بعضها بعضًا، وعصبونات \foreignlanguage{english}{SST} كل الأصناف الأخرى، وعصبونات \foreignlanguage{english}{VIP} عصبونات \foreignlanguage{english}{SST} في الغالب. إثارة عصبونات \foreignlanguage{english}{VIP} لدى فأر مستيقظ ترفع التثبيط عن عصبونات \nt{GLUT}؛ وتشغّلها إشارة التعزيز. وتقسيم الأدوار قاعدةً عامة فرضية & \cite{Tremblay2016,Cardin2009,Pfeffer2013,Pi2013} & فرضية \\
14 & تحد الطاقة التردد المتوسط بقيمة من رتبة نبضة واحدة في الثانية~\eqref{eq:energy} & الميزانية الطاقية للمادة الرمادية لدى القوارض من بيانات تشريحية وفيزيولوجية: النبضات والتيارات المشبكية $47\,\%$ و$34\,\%$ من طاقة الإشارات؛ ولا يمكن أن ينشط معًا أكثر من $\sim15\,\%$ من العصبونات. ولقشرة الإنسان: لا يمكن أن ينشط بقوة معًا أكثر من $\sim1\,\%$ من العصبونات. والرقم $\sim10^9$ جزيء ATP للنبضة والقدرة $\sim10\,\text{W}$ للإشارات تقدير الفصل المبني على هذه الحسابات & \cite{Attwell2001,Lennie2003} & نظرية \\
15 & الشيفرة متناثرة ومضبوطة على أكبر عدد من البتات لكل جول؛ والقشرة تقايض الدقة بالطاقة (القضية~\ref{prop:economy}) & فرضية الشيفرة المقتصدة. السند: لدى الفئران عند تقييد الغذاء تهبط في القشرة البصرية موصلية مستقبلات \foreignlanguage{english}{AMPA} والإنفاق المشبكي من ATP ($-29\,\%$)، ويبقى تردد النبضات، ويتسع الضبط على الاتجاه بـ$32\,\%$ & \cite{Barlow1961,Padamsey2022} & فرضية \\
\bottomrule
\end{longtable}}

\section{الفصل~\ref{sec:serocomputer}. عصبونات \nt{SERO}}

خصائص عصبونات \nt{SERO} نفسها (الإيقاع، والتثبيط الذاتي، والإشارة بحسب المستقبل، والأنظمة الفرعية) مُقاسة مباشرة؛ وحسابات الفصل، أي المنظِّم والمرشّح والمنطق، تنتج من معادلاته؛ أما $\tau_c$ ومعامل انتشار السيروتونين وعدد مواقع التحرير فتقديرات؛ ودور الحلقة منظِّمًا للمستوى في الدماغ فرضية (ففي نواة الرفاء يكون التثبيط الذاتي نقطيًا ولا يعطي إلا توقفات قصيرة)، ودور \nt{SERO} أفقًا فرضية تدعمها تجارب سلوكية.

{\footnotesize
\setlength{\tabcolsep}{3pt}\renewcommand{\arraystretch}{1.15}
\begin{longtable}{>{\raggedright}p{0.5cm}>{\raggedright}p{4.0cm}>{\raggedright}p{5.6cm}>{\raggedright}p{2.3cm}>{\raggedright\arraybackslash}p{1.7cm}}
\toprule
م & القضية & التجربة وما قيس & المصدر & الحالة \\
\midrule\endfirsthead
\toprule
م & القضية & التجربة وما قيس & المصدر & الحالة \\
\midrule\endhead
1 & إشارة أثر السيروتونين يختارها مستقبِل المتلقي (القضية~\ref{prop:receptor}) & تُقسم مستقبلات السيروتونين إلى عائلات بحسب علم الأدوية والآلية: \foreignlanguage{english}{5-HT$_{1A}$} يفتح عبر بروتين \foreignlanguage{english}{G} قنوات البوتاسيوم ويثبّط الخلية، و\foreignlanguage{english}{5-HT$_{2A}$} و\foreignlanguage{english}{5-HT$_3$} المؤيني التوجه يثيران. الناقل الواحد يعطي استجابات متعاكسة في خلايا مختلفة. & \cite{BarnesSharp1999} & مُقاس \\
2 & عصبون \nt{SERO} في اليقظة يطلق بانتظام بضع نبضات في الثانية & تسجيل عصبونات نواة الرفاء الظهرية عند قطط تتحرك بحرية: تفريغ إيقاعي بطيء $2.8\pm0.2$ نبضة في الثانية في اليقظة الهادئة؛ في النوم البطيء ينخفض التردد بـ$34\text{--}68\,\%$، وفي النوم السريع بـ$98\,\%$. وتحت التخدير عند الجرذان $1.7\pm0.5$ نبضة في الثانية، بانتظام شديد. & \cite{TrulsonJacobs1979, GagnonParent2014, JacobsAzmitia1992} & مُقاس \\
3 & عصبون \nt{SERO} ناظمة إيقاع: لا يحتاج إلى دفعة لكل نبضة، لكنه يحتاج إلى دفق خلفي ثابت (في الشريحة نورأدرينالين عبر مستقبلات $\alpha_1$، وفي الكائن الحي من \nt{NORA}) & في شريحة الدماغ تفرّغ الخلايا ببطء وانتظام، وبعد كل نبضة فرط استقطاب لاحق مدته $200\text{--}800\ms$. الخلايا العاملة تلقائيًا في الشريحة أقل منها في الكائن الحي: الإيقاع المنتظم يحتاج إلى دخل نورأدرينالين مستمر عبر مستقبلات $\alpha_1$، وحجبها يخمده. & \cite{VandermaelenAghajanian1983} & مُقاس \\
4 & تكاد كل المستقبلات تكون أيضية التوجه؛ والاستجابة بطيئة: ترتفع وتهبط في غضون ثانية تقريبًا & كل عائلات المستقبلات ما عدا \foreignlanguage{english}{5-HT$_3$} مرتبطة ببروتينات \foreignlanguage{english}{G}. في الشريحة يعطي التحرير المستثار للسيروتونين عبر \foreignlanguage{english}{5-HT$_{1A}$} تيارًا مثبّطًا يرتفع ويهبط في غضون ثانية. & \cite{BarnesSharp1999, CourtneyFord2016} & مُقاس \\
5 & في مناطق الهدف يخرج السيروتونين إلى الحجم، لا إلى الشق وحده؛ وفي نواة الرفاء نفسها يجري تثبيط الجيران من نقطة إلى نقطة & قياس الفولتية الدوري السريع في الشرائح: التحرير لكل نبضة واحد في منطقة فيها مشابك وفي نواة الرفاء التي تكاد تخلو منها؛ ولا يعتمد على حجب الاسترداد والمستقبلات؛ وعدد الجزيئات لكل نهاية أقل بكثير من عدد بروتينات النقل والمستقبلات، والتركيز خارج الخلايا يطابق ألفة المستقبل الرئيسي. في نواة الرفاء يجري التثبيط عبر \foreignlanguage{english}{5-HT$_{1A}$} من نقطة إلى نقطة: بروتين النقل لا يدع السيروتونين يتراكم خارج الشق. & \cite{Bunin1998, CourtneyFord2016} & مُقاس \\
6 & التركيز وفق~\eqref{eq:seroneuron} يتناقص بسبب الضخ العائد؛ و$\tau_c$ من رتبة الثانية تقدير & قياس الفولتية في الدماغ الحي: السيروتونين خارج الخلايا يحدّه قبل كل شيء الاسترداد والتفكيك، لا التخليق. لا قياس مباشرًا لـ$\tau_c$ في الأعمال المستشهد بها؛ ورتبة المقدار تقدير الفصل. & \cite{Hashemi2012serotonin} & مُقاس؛ $\tau_c$ تقدير \\
7 & \foreignlanguage{english}{NOT} و\foreignlanguage{english}{NOR} بناظمة إيقاع واحدة؛ اكتمال منطق \nt{SERO} (القضية~\ref{prop:serocomplete}، الشكل~\ref{fig:serogates}) & ينتج من~\eqref{eq:seroneuron}: ناظمة الإيقاع القابلة للتثبيط بوابة \foreignlanguage{english}{NOR}، و\foreignlanguage{english}{NOR} كاملة وظيفيًا. لا توجد تجربة بُني فيها منطق من عصبونات \nt{SERO}؛ وشرط الأحجام المنفصلة لا يتحقق في الدماغ. & اشتقاق في الفصل & نظرية \\
8 & السيروتونين يثبّط عصبون \nt{SERO} نفسه عبر مستقبلات عليه & عند جرذان مخدّرة تثبّط ناهضات \foreignlanguage{english}{5-HT$_{1A}$}، وريديًا وبالرحلان الأيوني، تفريغ عصبونات الرفاء الظهرية بعلاقة الجرعة والاستجابة نفسها التي للسيروتونين؛ أما ناهضات \foreignlanguage{english}{5-HT$_{1B}$} فتكاد لا تؤثر. وفي الشريحة يعطي السيروتونين الداخلي المنشأ من الخلايا المجاورة عبر \foreignlanguage{english}{5-HT$_{1A}$} تيارًا وتوقفًا قصيرًا في التفريغ دون أن يغيّر التردد المتوسط. & \cite{Sprouse1987, CourtneyFord2016} & مُقاس \\
9 & في النموذج، الحلقة عبر الحجم المشترك منظِّم مستوى: التشتت والثابت الزمني يصغران بـ$1+G$ مرة~\eqref{eq:feedback}؛ وأن الحلقة تعمل هكذا في الدماغ فرضية & الصيغة الكلاسيكية للمضخّم بتغذية راجعة سالبة، وقد اشتُقت في الفصل من جديد. كسب الحلقة $G$ لنظام \nt{SERO} غير مقيس. والمقيس: في الشريحة يجري التثبيط الذاتي من نقطة إلى نقطة، ويعطي توقفًا قصيرًا ولا يغيّر التردد المتوسط. & \cite{Black1934feedback, CourtneyFord2016} & نظرية؛ وفي الدماغ فرضية \\
10 & التركيز متوسط متحرك للتردد، أي مرشّح تمرير منخفض~\eqref{eq:lowpass}، الشكل~\ref{fig:lowpass} & حل المعادلة الخطية~\eqref{eq:seroneuron}؛ والشكل حساب بهذه الصيغة عند $\tau_c=1\,\text{s}$. لا يوجد نموذج مستقل في \texttt{models/}. & اشتقاق في الفصل & نظرية \\
11 & تعرّج الشقوق يبطئ انتشار الجزيئة الصغيرة نحو $2.6$ مرة & طريقة المصدر النقطي: رحلان أيوني لرباعي ميثيل الأمونيوم وتسجيل تركيزه بقطب انتقائي للأيون في الشرائح وفي الدماغ الحي. التعرّج $\lambda\approx1.6$، أي أن $D$ الفعّال أصغر من الحر بـ$\lambda^2\approx2.6$ مرة؛ وحصة الحجم بين الخلايا نحو $0.2$. معامل انتشار السيروتونين نفسه في النسيج لم يُقَس في هذه الأعمال؛ وهو في الفصل تقدير. & \cite{Sykova2008} & مُقاس \\
12 & طول ”السلك“ المستقل $\ell\sim\sqrt{D\tau}\approx20\um$~\eqref{eq:diffusionlength}؛ وعدد الأسلاك محدود بعدد هذه المناطق (القضية~\ref{prop:volumewires}) & قانون الانتشار مع التعويض بتقديرات $D$ و$\tau$ (السطران~6 و~11)؛ والقضية~\ref{prop:volumewires} نتيجة. لا توجد تجربة مباشرة تعدّ القنوات المستقلة للنقل الحجمي. & اشتقاق في الفصل & نظرية \\
13 & خرج عصبون \nt{SERO} الواحد منتشر على مساحات هائلة من الدماغ؛ ومواقع التحرير تُقدَّر بـ$\sim10^5$ & إعادة بناء كاملة لعصبونات مفردة من الرفاء الظهرية عند الجرذ: كل المحاوير الصاعدة كثيرة التفرع، والطول الكلي لمحوار الخلية الواحدة يبلغ $18.7\,\text{cm}$، وفروع الخلية الواحدة تبلغ المخطط والقشرة الجبهية الأمامية واللوزة. عدد مواقع التحرير لكل خلية لم يُعَدّ مباشرة. & \cite{GagnonParent2014} & مُقاس؛ $10^5$ تقدير \\
14 & تنقسم عصبونات \nt{SERO} إلى بضعة أنظمة فرعية بمخارج واستجابات مختلفة & وسم جيني فيروسي عند الفئران: العصبونات الذاهبة إلى اللوزة وتلك الذاهبة إلى القشرة الجبهية تكاد لا تتقاطع في تفرعها، وتتلقى مداخل مختلفة، وتستجيب لمنبّه مزعج استجابتين متعاكستين؛ وتشغيل كل مجموعة وإيقافها يغيّران السلوك تغييرًا مختلفًا. & \cite{Ren2018} & مُقاس \\
15 & شرط الصبر $D<\tau_\gamma\ln(R/r_0)$~\eqref{eq:patience} مع الخصم الأسّي؛ ومع الخصم المقيس، الزائدي، $D<\tau_\gamma(R/r_0-1)$ & ينتج من الخصم $e^{-D/\tau_\gamma}$ أو $1/(1+D/\tau_\gamma)$؛ والاشتقاق في الفصل. قرود، اختيار على تأخيرات من رتبة الثواني: الخصم أقرب إلى المنحنى الزائدي منه إلى الأسّي؛ واستجابة عصبونات \nt{DOPA} للإشارة المنبئة بمكافأة مؤجلة تنخفض مع التأخير زائديًا أيضًا. & اشتقاق في الفصل؛ \cite{KobayashiSchultz2008} & نظرية \\
16 & \nt{SERO} يضبط الأفق $\tau_\gamma$ (القسم~\ref{sec:horizon}) & التشغيل البصري الجيني لعصبونات \nt{SERO} في الرفاء الظهرية عند الفئران يطيل انتظار المكافأة المؤجلة ويزيد المثابرة في البحث عن مكافأة صارت نادرة. وعند البشر يجعل خفضُ السيروتونين (نظام غذائي بلا تريبتوفان) خصمَ المكافأة المؤجلة أشد انحدارًا. كيف يتحول التركيز إلى $\tau_\gamma$ غير معروف. & \cite{Doya2002, Miyazaki2014, Lottem2018, Schweighofer2008} & فرضية \\
\bottomrule
\end{longtable}}

\section{الفصل~\ref{sec:dofa}. التعلّم مع \nt{DOPA}}

آليات المشبك (الدفعات، وكاشف التزامن، والكالسيوم، ونوافذ اللدونة) وسلوك الدوبامين خطأً للتنبؤ مقيسة مباشرة؛ وأن القواعد تقود إلى التدرج، وما ثمن البثّ، مبرهنتان؛ ونتيجة تعلّم الاختيار حساب بنموذج الكتاب؛ والدارة التي تحسب الخطأ فرضية.

{\footnotesize
\setlength{\tabcolsep}{3pt}\renewcommand{\arraystretch}{1.15}
\begin{longtable}{>{\raggedright}p{0.5cm}>{\raggedright}p{4.0cm}>{\raggedright}p{5.6cm}>{\raggedright}p{2.3cm}>{\raggedright\arraybackslash}p{1.7cm}}
\toprule
م & القضية & التجربة وما قيس & المصدر & الحالة \\
\midrule\endfirsthead
\toprule
م & القضية & التجربة وما قيس & المصدر & الحالة \\
\midrule\endhead
1 & لم يُعثر في الدماغ على الانتشار العكسي للخطأ بالشكل الذي هو عليه في الشبكات الاصطناعية & لا توجد تجربة مباشرة: هذه قضية عن غياب. تحلل المراجعات ما تتطلبه الخوارزمية (وصلات راجعة تكرر الأمامية، وطور مستقل) وما اقتُرح لها من تقريبات بيولوجية. & \cite{Lillicrap2020, WhittingtonBogacz2019} & فرضية \\
2 & يتفكك الوزن إلى عدد مواقع التحرير واحتمال التحرير والاستجابة للدفعة: $w=N_s\,p\,A_1$~\eqref{eq:quantal} & المشبك العصبي العضلي عند الضفدع مع تحرير مخفَّض: تتقلب استجابة المتلقي بدرجات هي مضاعفات الاستجابة المصغّرة التلقائية لدفعة واحدة؛ وعدد الدفعات لكل نبضة يخضع لقانون بواسون. & \cite{delCastillo1954} & مُقاس \\
3 & مستقبِل المتلقي كاشف تزامن؛ والكالسيوم يطلق تغيّر الوزن & تثبيت الجهد الرقعي (\foreignlanguage{english}{patch clamp}) لعصبونات الفأر: أيونات \foreignlanguage{english}{Mg$^{2+}$} تسدّ القناة التي يفتحها الغلوتامات عند كمون الراحة وتخرج منها عند إزالة الاستقطاب. شرائح الحُصين: خالب الكالسيوم في المتلقي يحجب التقوية الطويلة الأمد، وتحرير الكالسيوم بالضوء داخل الخلية يقوّي النقل بنفسه. & \cite{Nowak1984, Malenka1988calcium} & مُقاس \\
4 & القرار ينفذه أي جانب: يكبر $A_1$ عند المتلقي، ويتغير $p$ عند المرسِل & الغلوتامات المحرَّر بالضوء فوق شويكة واحدة يسبب نموًا سريعًا وثابتًا للشويكة وللتيار عبر مستقبلاتها؛ وترافق التقويةَ إدماجُ مستقبلات \foreignlanguage{english}{AMPA}. إزالة استقطاب المتلقي تكبت مؤقتًا التحرير عند المرسِل؛ وينقل الأثرَ القنبينوداتُ الداخلية التي تعبر الشق عائدة (أُظهر ذلك على مداخل \nt{GABA}). والاستنزاف القصير الأمد يعتمد على احتمال التحرير عند المرسِل. & \cite{Matsuzaki2004, Malinow2002, WilsonNicoll2001, Tsodyks1997} & مُقاس \\
5 & يقوى المشبك حين يكون المرسِل والمتلقي نشطين معًا~\eqref{eq:hebb} & أرنب مخدَّر: بعد سلاسل قصيرة من النبضات على مسار الدخل إلى الحُصين تبقى الاستجابة مقوّاة مدة من $30$ دقيقة إلى $10$ ساعات. وفي شريحة الحُصين لا تقوّي نبضات المتلقي إلا المشابك النشطة في اللحظة نفسها. & \cite{Bliss1973LTP, Kelso1986hebb} & مُقاس \\
6 & القاعدة~\eqref{eq:oja} تجد المتجه الذاتي الرئيسي لارتباطات المداخل (القضية~\ref{prop:oja}) & مبرهنة؛ والبرهان في الفصل. لا توجد تجربة تعلّم فيها عصبون المكوّن الرئيسي لمداخله؛ وآلية تقييد نمو الأوزان في المشبك لم تُحدَّد. & \cite{Oja1982} & نظرية \\
7 & هيب بالتوقيت: نافذة نحو $\pm20\ms$ (الشكل~\ref{fig:stdp})؛ ومن التسلسلات المتكررة تنمو السلاسل & أزواج من عصبونات الحُصين في المزرعة: نبضة المتلقي خلال $20\ms$ بعد نبضة المرسِل تعطي تقوية ثابتة، وخلال $20\ms$ قبلها تعطي إضعافًا؛ وكلاهما يعتمد على مستقبلات \foreignlanguage{english}{NMDA}. النسبة $A_-=0.6A_+$ في الشكل توضيحية. وأن السلاسل تنمو هكذا في الشبكة لم يُقَس مباشرة. & \cite{BiPoo1998} & مُقاس \\
8 & الأثر~\eqref{eq:threefactor}: الدوبامين الواصل بعد $1\text{--}2\,\text{s}$ من النشاط المشترك يقوّي الوصلة، وبعد ذلك لا (القضية~\ref{prop:threefactor}) & شرائح المخطط، تحفيز بصري جيني منفصل لمداخل الغلوتامات والدوبامين: لا تنمو الشويكة إلا إذا وصل الدوبامين بعد $0.3\text{--}2\,\text{s}$ من دخل الغلوتامات؛ والنافذة يحددها ارتفاع \foreignlanguage{english}{cAMP} السريع وتفككه. وفي القشرة يترك الإقران آثارًا منفصلة للتقوية والإضعاف، تحوّلها مستقبلات الأمينات الأحادية إلى تغيرات طويلة الأمد في نافذة من رتبة الثواني. & \cite{Yagishita2014, He2015eligibility} & مُقاس \\
9 & توجد نوافذ طولها ثوانٍ بلا دوبامين أيضًا: الإشارة الثالثة إثارة قوية للمتلقي & فأر حي، المنطقة \foreignlanguage{english}{CA1}: حدث هضبة واحد في المتلقي يقوّي المداخل الواصلة قبله وبعده بثوانٍ، وينشئ في عبور واحد حقلَ مكان. & \cite{Bittner2017} & مُقاس \\
10 & الخطوة المتوسطة لقاعدة العوامل الثلاثة تسير على تدرج المكافأة~\eqref{eq:perturbation} (القضية~\ref{prop:gradient}) & مبرهنة التدرج للتعلّم بالانحرافات العشوائية؛ واشتُقت في الفصل لضجيج صغير. & \cite{Williams1992} & نظرية \\
11 & الضجيج بحث: تتعلم الشبكة من انحرافاتها العشوائية & طيور الحسون الصغيرة: تعطيل النواة التي تخرج من دارة العقد القاعدية يزيل تنوّع الأغنية بحدة وبصورة عكوسة. الحسون البالغ: يُطبَّق تشويش على أداءات المقطع التي تقع طبقتها على جانب من الوسط، فتزيح الطيور الطبقة بسرعة إلى الجانب الآخر، أي تتعلم من تشتتها هي. & \cite{Olveczky2005, TumerBrainard2007} & مُقاس \\
12 & الاختيار بين اثنين يُتعلَّم عند $\tau_e=1\,\text{s}$ و$\delta=R-\bar R$؛ ولا يُتعلَّم عند $\tau_e=50\ms$؛ وبلا طرح تنمو الأوزان بلا حد، ومع معدل تعلّم كبير قد تثبّت الشبكة الاختيار الأسوأ & \texttt{models/dofa\_learning.py}، $\eta=0.05$، $500$ محاولة، $6$ بذور. $\tau_e=1\,\text{s}$: نسبة اختيار $A$ $0.65\text{--}0.79\to0.96\text{--}1.00$، والأوزان $0.85\text{--}1.33$. $\tau_e=50\ms$: $0.38\text{--}0.55$، والأوزان لا تتغير. $\delta=R$: وزن الفائز $860\text{--}1190$. $\delta=R$، $\eta=0.5$، $3$ بذور: واحدة ثبتت على $B$ ($0.04\to0$). يطبع البرنامج الحالات الأربع كلها؛ وإعادة الحساب طابقت الفصل. & \texttt{models/\allowbreak dofa\_\allowbreak learning.py} & نموذج \\
13 & زمن التعلّم بإشارة مشتركة واحدة يزداد متناسبًا مع عدد العصبونات الضاجّة (القضية~\ref{prop:broadcastcost}) & منحنيات التعلّم لشبكة خطية: اضطراب العصبونات أبطأ من التدرج الدقيق بعدد مرات يساوي عدد عصبونات الخرج، واضطراب الأوزان أبطأ فوق ذلك بعدد مرات يساوي عدد المداخل. & \cite{Werfel2005} & نظرية \\
14 & مخارج \nt{DOPA} تذهب قبل كل شيء إلى مناطق اختيار الأفعال & إعادة بناء كاملة لمحاوير عصبونات \nt{DOPA} سوداوية مخططية مفردة عند الجرذ: فروع الخلية الواحدة تغطي $0.45\text{--}5.7\,\%$ من حجم المخطط (بمتوسط $2.7\,\%$)، ولا تكاد توجد فروع خارج المخطط. & \cite{Matsuda2009} & مُقاس \\
15 & تتعلم القشرة التنبؤ بدخلها، ويُحسب الخطأ في مكانه & مراجعة: في القشرة الحسية استجابات لعدم التطابق بين الدخل المتوقع والواصل. أما أن هذا قاعدة تعلّم عامة للقشرة فغير مبرهن. & \cite{Keller2018} & فرضية \\
16 & يسلك الدوبامين سلوك الخطأ~\eqref{eq:ctd}: ذروة، وانتقال إلى الإشارة المنبئة، وهبوط (الشكل~\ref{fig:ctdcases}) & عصبونات \nt{DOPA} عند القرد: ذروة للعصير غير المتوقع؛ وبعد التعلّم ذروة للإشارة المنبئة ولا استجابة للعصير الموعود؛ وهبوط حين لا يصل العصير. الصيغة المتصلة~\eqref{eq:ctd} نظرية. & \cite{Schultz1997, Doya2000} & مُقاس \\
17 & دارة تحسب $\dd V/\dd t$ وتطرح التوقع & فئران، تسجيل وعلم بصري جيني في السقيفة البطنية: عصبونات \nt{DOPA} تطرح التوقع من المكافأة؛ وعصبونات \nt{GABA} المجاورة تثبّطها حين تكون المكافأة متوقعة، وإثارتها أو تثبيطها يغيّر الخطأ. أما كيف تُحسب المشتقة فلم يُبيَّن. & \cite{Eshel2015arithmetic} & فرضية \\
18 & عصبون \nt{DOPA} ناظمة إيقاع؛ والهبوطات أضحل من الذرى & في الشريحة تفرّغ عصبونات \nt{DOPA} تلقائيًا وبانتظام شديد، على كمون ناظم بطيء. وعند القرد يتبع التردد الخطأ الموجب كميًا، أما السالب فلا ينقله إلا ضعيفًا. & \cite{GraceOnn1989, Bayer2005} & مُقاس \\
19 & الفاعل والناقد (الشكل~\ref{fig:actorcritic}) & الخوارزمية من نظرية التعلّم المعزَّز. عند البشر (\foreignlanguage{english}{fMRI}) يسلك المخطط البطني سلوك الناقد في الاختيار ودونه، أما الظهري فلا يفعل ذلك إلا حين يجب اختيار الأفعال. & \cite{SuttonBarto2018, ODoherty2004actor} & فرضية \\
20 & إشارة $\delta$ في القاعدة معكوسة في العصبونات ذات العائلة الثانية من المستقبلات & شرائح المخطط: في العصبونات ذات مستقبلات \foreignlanguage{english}{D1} يعطي الإقران تقوية تحتاج إلى \foreignlanguage{english}{D1}؛ وفي ذات \foreignlanguage{english}{D2} إضعافًا يحتاج إلى \foreignlanguage{english}{D2}. & \cite{Shen2008} & مُقاس \\
\bottomrule
\end{longtable}}

\section{الفصل~\ref{sec:dofaalt}. نظريات أخرى في \nt{DOPA}}

كل واحدة من الصور الموسَّعة تستند إلى قياساتها المباشرة للدوبامين (تشتت نقاط الانقلاب، والاستجابات للمزعج والجديد، والارتفاع البطيء، والاستجابات لتغيّر الطعم)، لكن الصيغ التي تفسرها نظريات، وبين اثنتين منها تتناقض التجارب حتى الآن.

{\footnotesize
\setlength{\tabcolsep}{3pt}\renewcommand{\arraystretch}{1.15}
\begin{longtable}{>{\raggedright}p{0.5cm}>{\raggedright}p{4.0cm}>{\raggedright}p{5.6cm}>{\raggedright}p{2.3cm}>{\raggedright\arraybackslash}p{1.7cm}}
\toprule
م & القضية & التجربة وما قيس & المصدر & الحالة \\
\midrule\endfirsthead
\toprule
م & القضية & التجربة وما قيس & المصدر & الحالة \\
\midrule\endhead
1 & القاعدة غير المتناظرة~\eqref{eq:asymmetric} تتقارب إلى النقطة~\eqref{eq:expectile}؛ عند $\kappa=1/2$ هي المتوسط، ولقطرة باحتمال $1/2$ تكون $V=\kappa$ (الشكل~\ref{fig:expectiles}) & النقطة~\eqref{eq:expectile} هي ”التوقّعي“ (\foreignlanguage{english}{expectile})، أي حل مسألة المربعات الصغرى بأوزان مختلفة للانحرافات الموجبة والسالبة؛ ولمثال الفصل الاشتقاق في الفصل نفسه. & \cite{NeweyPowell1987}؛ اشتقاق في الفصل & نظرية \\
2 & لعصبونات \nt{DOPA} المختلفة نقاط انقلاب مختلفة، مرتبة بحسب اللاتناظر (القضية~\ref{prop:expectile}) & فئران، عصبونات \nt{DOPA} في السقيفة البطنية موسومة بصريًا جينيًا، ومكافآت مختلفة الحجم: النقطة التي تتحول فيها الذروة إلى هبوط تختلف من عصبون إلى آخر؛ وهي أعلى عند العصبونات ذات اللاتناظر الأكبر في الاستجابات؛ ومن استجابات المجموعة يُستعاد شكل توزيع المكافآت. أما أن هذا هو الوظيفة الرئيسية للتشتت ففرضية. & \cite{Dabney2020} & مُقاس \\
3 & جزء من عصبونات \nt{DOPA} يستجيب للمزعج بذروة & قرود: بعض عصبونات \nt{DOPA} يُستثار بالإشارة المنبئة بالمكافأة ويُثبَّط بالإشارة المنبئة بنفخة هواء في العين، وبعضها يُستثار بكلتيهما؛ والمجموعتان في أجزاء مختلفة من الدماغ المتوسط. & \cite{Matsumoto2009, BrombergMartin2010} & مُقاس \\
4 & القيمة المطلقة للخطأ أو الجِدّة تحدد معدل التعلّم & لا توجد في الأعمال المستشهد بها تجربة مباشرة تغيّر فيها إشارة الأهمية عند عصبونات \nt{DOPA} معدل التعلّم؛ والمراجعة تقترح دور إشارة ”انتبه، هذا مهم“. & \cite{BrombergMartin2010} & فرضية \\
5 & سلك مستقل لمحور الخطر & فئران: عصبونات \nt{DOPA} الذاهبة إلى الطرف الخلفي من المخطط تستجيب للمنبّهات الجديدة والقوية لا للمكافأة؛ وإتلافها يضعف تجنّب الشيء الجديد. & \cite{Menegas2018} & مُقاس \\
6 & زمن الفعل الأمثل $\tau_a^*=\sqrt{C/\bar R}$~\eqref{eq:vigor} & أصغر قيمة للمجموع $C/\tau_a+\bar R\tau_a$؛ والاشتقاق في الفصل. في النموذج الأصلي يحدد سرعةَ الأفعال ثمنُ الوقت، وهو يساوي متوسط المكافأة. & \cite{Niv2007}؛ اشتقاق في الفصل & نظرية \\
7 & المستوى البطيء للدوبامين يحمل $\bar R$ ويحدد سرعة الأفعال (القضية~\ref{prop:multiplex}) & جرذان، غسيل مجهري وقياس فولتية في النواة المتكئة: مستوى الدوبامين يتبع من دقيقة إلى دقيقة تواترَ المكافآت وسرعةَ الأفعال. وعند البشر يقوّي \foreignlanguage{english}{L-DOPA} اعتماد زمن الاستجابة على متوسط تواتر المكافآت. أما أن المستوى البطيء يساوي $\bar R$ تحديدًا فغير مبرهن. & \cite{Hamid2016, Beierholm2013vigor} & فرضية \\
8 & المستوى البطيء يتألف من مصدرين: التحرير يتغير عند المخارج دون تغيّر تردد النبضات، لكن الارتفاع البطيء موجود في التردد نفسه أيضًا & جرذان، مهمة واحدة: التحرير في النواة المتكئة يتبع التوقع المتغير للمكافأة، أما تردد نبضات عصبونات \nt{DOPA} المعرَّفة في السقيفة البطنية فلا. فئران في ممر افتراضي (السطر~10): الارتفاع التدريجي عند الاقتراب من المكافأة موجود في نبضات عصبونات \nt{DOPA} المعرَّفة أيضًا. & \cite{Mohebi2019, Kim2020} & مُقاس \\
9 & يرتفع الدوبامين تدريجيًا ما دام الحيوان يقترب من مكافأة بعيدة & جرذان في متاهة، قياس الفولتية في المخطط: يتزايد تركيز الدوبامين خلال ثوانٍ مع الاقتراب من الهدف، ويعتمد الانحدار على المسافة وعلى حجم المكافأة. & \cite{Howe2013ramp, Hamid2016} & مُقاس \\
10 & الارتفاع هو الخطأ $\delta$ مع تدقيق التقدير~\eqref{eq:ramp}؛ والقفزة في الممر تعطي ذروة (الشكل~\ref{fig:ramp}) & الصيغة والعدد $\delta=0.75\,V$ في الثانية اشتقاق في الفصل. التجربة: فئران في ممر افتراضي، ونقل لحظي إلى نقطة أقرب من الهدف؛ نبضات أجسام الخلايا، والكالسيوم في الأجسام والمخارج، والتركيز في المخطط تستجيب كخطأ لا كتوقع $V$. أي التفسيرين صحيح على كل المخارج، مسألة قيد النقاش. & \cite{Kim2020} & مُقاس \\
11 & النظر إلى الوراء: يخبر الدوبامين بمقياس الصلة السببية~\eqref{eq:retro} & فئران، تحرير الدوبامين في النواة المتكئة في تجارب تتنبأ فيها هذه النظرية والخطأ~\eqref{eq:ctd} بنتائج مختلفة: في عدد من التجارب تبع التحرير الأولى. & \cite{Jeong2022} & فرضية \\
12 & تنزاح استجابة الدوبامين أثناء التعلّم تدريجيًا من المكافأة إلى الإشارة المنبئة & فئران، إشارة مخارج عصبونات \nt{DOPA} في المخطط البطني على مدى التعلّم: تنتقل الذروة تدريجيًا من لحظة المكافأة إلى لحظة الإشارة المنبئة، كما يقتضي التعلّم وفق~\eqref{eq:ctd}. & \cite{Amo2022} & مُقاس \\
13 & عصبونات \nt{DOPA} تستجيب لتغيّر طعم المكافأة مع ثبات القيمة & جرذان، تسجيل عصبونات السقيفة البطنية: استبدال عصير بآخر مساوٍ له في القيمة يستثير استجابة، مع أن الخطأ~\eqref{eq:ctd} يساوي صفرًا. & \cite{Takahashi2017} & مُقاس \\
14 & الذرى الاصطناعية تعلّم الصلة بين منبّهين محايدين & جرذان، تجربة إقران مسبق لمنبّهين دون مكافأة: التشغيل البصري الجيني لعصبونات \nt{DOPA} عند الانتقال من الأول إلى الثاني ينشئ الصلة، وإيقافها يعيق تعلّمها. & \cite{Sharpe2017} & مُقاس \\
15 & متجه الأخطاء بحسب السمات~\eqref{eq:vectortd} & نظرية خطأ التنبؤ بالسمات؛ ولا يوجد اختبار مباشر للصيغة المتجهية. & \cite{Gardner2018} & فرضية \\
16 & عصبونات \nt{DOPA} المختلفة في المهمة الواحدة تحمل مقادير مختلفة & فئران في ممر افتراضي، تصوير ثنائي الفوتون لعصبونات \nt{DOPA}: بعضها مرتبط بالمكافأة، وبعضها بالسرعة والانعطاف، وبعضها بالإشارات المنبئة ودقة الاختيار. & \cite{Engelhard2019} & مُقاس \\
17 & حزمة أسلاك بدل سلك (القضية~\ref{prop:bundle}) & خلاصة السطور~\babelsublr{2--16}: كل تفكيك مقيس على حدة؛ ولا توجد تجربة تبيّن صورة موحدة تكون فيها كلها أجزاءً من إشارة واحدة. & --- & فرضية \\
\bottomrule
\end{longtable}}

\section{الفصل~\ref{sec:nora}. \nt{NORA}}

بنية نظام \nt{NORA}، ووضعاه، وصلته بالحدقة (لا عبر \nt{NORA} وحده)، وارتفاع نسبة الإشارة إلى الخلفية، وحرف U المقلوب في السلوك، كلها مقيسة مباشرة؛ والكسب الضربي $g$ نموذج؛ وأن الكسب يبدّل وضع الاختيار ويعمل مقلوبًا لدرجة الحرارة استنتاجان في الفصل؛ وفائدة ضبط الكسب حساب بنموذج الكتاب؛ أما ”\nt{NORA} مقبض الاستكشاف“ و”\nt{NORA} مقاطعة“ ففرضيتان.

{\footnotesize
\setlength{\tabcolsep}{3pt}\renewcommand{\arraystretch}{1.15}
\begin{longtable}{>{\raggedright}p{0.5cm}>{\raggedright}p{4.0cm}>{\raggedright}p{5.6cm}>{\raggedright}p{2.3cm}>{\raggedright\arraybackslash}p{1.7cm}}
\toprule
م & القضية & التجربة وما قيس & المصدر & الحالة \\
\midrule\endfirsthead
\toprule
م & القضية & التجربة وما قيس & المصدر & الحالة \\
\midrule\endhead
1 & عصبونات \nt{NORA} عند الإنسان بضع عشرات الآلاف، كلها تقريبًا في تجمّع واحد & مقاطع متسلسلة من جذع الدماغ البشري، صبغ بالتيروزين هيدروكسيلاز: $53\,900$ خلية في الموضع الأزرق، و$6\,260$ أخرى في النواة الفرعية المجاورة. & \cite{Baker1989LC} & مُقاس \\
2 & المخارج تتفرق في الدماغ كله تقريبًا، لكن أجزاء التجمع تخدم جزئيًا مناطق مختلفة & وسم جيني فيروسي عند الفئران: عصبونات \nt{NORA} في الموضع الأزرق تتلقى مداخل من مناطق واسعة في الدماغ وترسل مخارج إلى مناطق واسعة في الدماغ والنخاع الشوكي. وعند الجرذان: المجموعة الذاهبة إلى اللوزة تقوّي تعلّم الخوف، ومجموعة مستقلة ذاهبة إلى القشرة الجبهية الأمامية تقوّي إطفاءه. & \cite{Schwarz2015LC, Uematsu2017LC, Poe2020} & مُقاس \\
3 & الوضع الخلفي: نبضات منتظمة بتردد من أجزاء النبضة إلى بضع نبضات في الثانية، يتغير مع اليقظة & جرذان، تسجيل عصبونات الموضع الأزرق في دورة النوم واليقظة: التردد أعلى ما يكون في اليقظة، وأدنى في النوم البطيء، وقريب من الصفر في النوم السريع؛ وتغيرات التردد تسبق تغيّر الحالة. & \cite{AstonJonesBloom1981} & مُقاس \\
4 & الوضع الطوري: رشقة قصيرة لحدث مهم للمهمة، نحو لحظة القرار & قرود، مهمة البحث عن هدف نادر: كل عصبونات الموضع الأزرق تستجيب برشقة للهدف وحده، بعده بـ$\sim90\ms$ وقبل الاستجابة باليد بـ$\sim200\ms$؛ وعند الأداء السيئ تكون الرشقات أضعف. وفي مهمة الاختيار بين اثنين ترتبط الرشقة بالاستجابة أدق من ارتباطها بالمنبّه، وتظهر مع الاستجابات الصحيحة والخاطئة معًا. & \cite{AstonJones1994vigilance, Clayton2004LC} & مُقاس \\
5 & قطر الحدقة نقطة اختبار غير دقيقة لـ\nt{NORA}: فهو يتبع \nt{NORA} و\nt{ACET} ومناطق أخرى & قرود: نبضات الموضع الأزرق، حتى النبضات المفردة لخلية واحدة، تتنبأ باتساع الحدقة؛ وتوجد صلة مشابهة لكن متأخرة في الأكيمات وفي القشرة الحزامية. فئران: تقلبات الحدقة تتبع نشاط المخارج النورأدرينالية والكولينية في القشرة معًا. & \cite{Joshi2016, Reimer2016} & مُقاس \\
6 & النورأدرينالين يكبت النشاط الخلفي أكثر من المستثار؛ والكسب الضربي~\eqref{eq:gain} نموذج ينقل هذا الأثر & قرود يقظة، القشرة السمعية، رحلان أيوني للنورأدرينالين: يُكبت النشاط الخلفي للعصبونات أكثر من استجابتها لأصوات بني جنسها، فترتفع نسبة الإشارة إلى الضجيج. الدالة السينية الضربية~\eqref{eq:gain} مأخوذة من نموذج شبكي؛ والضرب الحرفي للانحدار غير مقيس. & \cite{Foote1975NE, ServanSchreiber1990} & مُقاس؛ $g$ نموذج \\
7 & الكسب يرفع $d'$ ضد الضجيج بعد المضخّم وحده~\eqref{eq:gainsnr} (القضية~\ref{prop:gainnoise}) & اشتقاق في الفصل؛ وفي النموذج الشبكي يحسّن الكسب كشف الإشارة. لا توجد تجربة تفصل الضجيج قبل المضخّم عن الضجيج بعده وتختبر هذا الفرق. & اشتقاق في الفصل؛ \cite{ServanSchreiber1990} & نظرية \\
8 & الحد $g\beta=1$ ينقل شبكة الاختيار من ”تتروّى“ إلى ”قررت“~\eqref{eq:wtagain} (القضية~\ref{prop:gainwta}، الشكل~\ref{fig:pitchfork}) & تحليل خطي لمجموعتين تثبّط إحداهما الأخرى؛ والاشتقاق في الفصل. لا توجد قياسات مباشرة لهذا الحد في الدماغ. & اشتقاق في الفصل & نظرية \\
9 & معامل الكسب مقلوب درجة حرارة الاختيار~\eqref{eq:softmax} & مع ضجيج لوجستي بعد المضخّم يكون احتمال الاختيار دالة \foreignlanguage{english}{softmax} مع $T=\sigma/g$؛ والاشتقاق في الفصل. & اشتقاق في الفصل & نظرية \\
10 & \nt{NORA} يحدد كم نستكشف: النشاط الخلفي المرتفع $g$ فعّال صغير (استكشاف)، والرشقة الطورية لحظة القرار $g$ كبير (قرار) & بشر: قبل الاختيار العشوائي تكون الحدقة القاعدية أوسع منها قبل اختيار أفضل خيار معروف؛ والأشخاص ذوو الحدقة الأوسع أميل إلى الاستكشاف. جرذان: الانتقال إلى وضع الاختيار العشوائي يتحكم فيه دخل \nt{NORA} إلى القشرة الحزامية الأمامية. أما أن الرشقة ترفع $g$ الفعّال فلم يُقَس مباشرة. & \cite{JepmaNieuwenhuis2011, Tervo2014stochastic, AstonJones2005} & فرضية \\
11 & المكافأة تعتمد على $g$ على شكل حرف U مقلوب؛ وأفضل $g$ أصغر في العالم السريع التغير (القضية~\ref{prop:explore}، الشكل~\ref{fig:invertedU}) & \texttt{models/nora\_explore.py}، $60$ محاكاة في كل منها $4000$ محاولة. تغيّر كل $1000$ محاولة: أفضل $g=16$، والنسبة $0.976$؛ وكل $20$: $g=8$، والنسبة $0.886$؛ وعند $g=0.5$: $0.77$. إعادة الحساب طابقت الفصل. & \texttt{models/\allowbreak nora\_\allowbreak explore.py} & نموذج \\
12 & حرف U المقلوب في السلوك: أفضل أداء عند نشاط خلفي معتدل مع رشقات واضحة & قرود، المهمة نفسها في السطر~4: في فترات الأداء الجيد يكون التردد الخلفي معتدلًا والرشقات للهدف واضحة؛ وعند التردد الخلفي المرتفع تغيب الرشقات وتكثر الاستجابات الخاطئة. وتغيرات النشاط الخلفي والمستثار تتبع تقلبات جودة الأداء. & \cite{Usher1999LC, AstonJones2005} & مُقاس \\
13 & \nt{NORA} يخبر بعدم اليقين غير المتوقع، و\nt{ACET} بالمتوقع & نظرية استدلال بايزي بنوعين من عدم اليقين؛ ولا توجد في الأعمال المستشهد بها تجربة مباشرة تفصل هاتين الإشارتين بحسب النواقل. & \cite{YuDayan2005} & فرضية \\
14 & بعد تغيّر حاد يتعلم الناس أسرع، وتتسع الحدقة & بشر، مهمة تنبؤ مع تغيرات مفاجئة: التغيرات القصيرة في الحدقة تتبع تقدير احتمال التغيّر، وتتنبأ مع الحدقة القاعدية بمدى تغيير المعطيات الجديدة للتقدير (معدل التعلّم)؛ وتغيير الحدقة بما لا يتصل بالضوء والمهمة يغيّر هذا المعدل. أما أن الحدقة تبيّن هنا \nt{NORA} تحديدًا فاستنتاج من المعطيات غير المباشرة في السطر~5. & \cite{Nassar2012} & مُقاس \\
15 & الرشقة الطورية تعمل مقاطعةً: تعيد الشبكة ضبط حالتها & لا توجد تجربة مباشرة تعيد فيها رشقة \nt{NORA} ضبط حالة الشبكة؛ والمقيس الرشقات عند الأحداث المهمة فقط (السطر~4). & \cite{BouretSara2005} & فرضية \\
16 & ضبط $g$ أو معدل التعلّم بحسب المفاجأة لا يكاد يفضل أفضل الإعدادات الثابتة & \texttt{models/nora\_explore.py}، عالم بحقب من $1000$ محاولة (تغيّر كل $1000$ وكل $20$). أفضل $g$ ثابت $=8$: $0.925$؛ وضبط $g$: حتى $0.927$؛ وضبط معدل التعلّم: حتى $0.926$؛ ومعدل ثابت $0.4$: $0.928$. إعادة الحساب طابقت الفصل. & \texttt{models/\allowbreak nora\_\allowbreak explore.py} & نموذج \\
\bottomrule
\end{longtable}}

\section{الفصل~\ref{sec:acet}. نضيف \nt{ACET}}

أفعال الأسيتيل كولين الثلاثة مقيسة في الشرائح وفي الدماغ الحي، والتعارض بين الكتابة والقراءة مبيَّن على نموذج الكتاب، أما أن \nt{ACET} يؤدي دور خط تمكين الكتابة ويخبر بعدم اليقين المتوقع ففرضية لها سند جزئي من التجربة.

{\footnotesize
\setlength{\tabcolsep}{3pt}\renewcommand{\arraystretch}{1.15}
\begin{longtable}{>{\raggedright}p{0.5cm}>{\raggedright}p{4.0cm}>{\raggedright}p{5.6cm}>{\raggedright}p{2.3cm}>{\raggedright\arraybackslash}p{1.7cm}}
\toprule
م & القضية & التجربة وما قيس & المصدر & الحالة \\
\midrule\endfirsthead
\toprule
م & القضية & التجربة وما قيس & المصدر & الحالة \\
\midrule\endhead
1 & عصبونات \nt{ACET} في الدماغ قليلة، مجموعة في بضع مجموعات، ومخارجها تتفرق في القشرة كلها: قناة بثّ & صبغ بأضداد لإنزيم تخليق الأسيتيل كولين ووسم راجع عند الجرذ: القشرة والحُصين واللوزة والبصلة الشمية والمهاد تتلقى الأسيتيل كولين أساسًا من ست مجموعات خلوية متراصة في الدماغ الأمامي القاعدي وفي الجذع & \cite{Mesulam1983} & مُقاس \\
2 & مستوى الأسيتيل كولين في القشرة مرتفع في اليقظة ومنخفض في النوم البطيء & غسيل مجهري في القشرة والحُصين عند قطط تتحرك بحرية مع تسجيل \foreignlanguage{english}{EEG}: تحرير الأسيتيل كولين في اليقظة الهادئة أعلى بـ$\sim100\%$، وفي اليقظة النشطة بـ$\sim175\%$، منه في النوم البطيء؛ وفي النوم السريع يكون في القشرة كما في اليقظة & \cite{Marrosu1995,Hasselmo1999} & مُقاس \\
3 & معنى الإشارة يحدده المستقبل: للأسيتيل كولين مستقبلات-قنوات سريعة، ومستقبلات بطيئة تطلق تفاعلات في الخلية (القضية~\ref{prop:receptor}) & مراجعة للقياسات: أثر الأسيتيل كولين في الاستثارية والنقل واللدونة يعتمد على موضع التحرير، ونوع المستقبل الفرعي (النيكوتينية قنوات أيونية، والمسكارينية عبر بروتينات \foreignlanguage{english}{G})، والخلية المتلقية & \cite{Picciotto2012} & مُقاس \\
4 & الفعل الأول: إضعاف الوصلات الداخلية دون المساس تقريبًا بالمداخل الخارجية (العامل $1-a$ في~\eqref{eq:acetnet}) & شرائح القشرة الشمية عند الجرذ: عند $100\,\mu\text{M}$ من ناهضات الأسيتيل كولين تنخفض كمونات الوصلات الداخلية (الطبقة \foreignlanguage{english}{Ib}) بأكثر من $60\%$، وكمونات الدخل الخارجي (الطبقة \foreignlanguage{english}{Ia}) بأقل من $15\%$، في الخلية نفسها؛ والكبت قبل مشبكي. شرائح الحُصين (\foreignlanguage{english}{CA1}): $89\%$ مقابل $40\%$ للمسار الداخلي ومسار الدخل & \cite{HasselmoBower1992,HasselmoSchnell1994} & مُقاس \\
5 & الفعل الثاني: تقوية الدخل (العامل $\frac{1+a}{2}$) & شرائح القشرة الحديثة: المستقبلات النيكوتينية تقوّي مشابك الدخل من المهاد ولا تؤثر في المشابك داخل القشرة؛ والمسكارينية تضعف النوعين. والأسيتيل كولين يرفع استثارية العصبونات (مراجعة) & \cite{Gil1997,Hasselmo2006} & مُقاس \\
6 & الفعل الثالث: تيسير تغيّر الأوزان (المعدل $\eta a$) & جرذ: التحفيز الكهربائي للنواة القاعدية (مصدر الأسيتيل كولين للقشرة) مقرونًا بنغمة يُحدث عند الحيوان البالغ إعادة تنظيم قوية لخريطة القشرة السمعية نحو الصوت المعروض & \cite{KilgardMerzenich1998,Hasselmo2006} & مُقاس \\
7 & قاعدة هيب للأنماط المتناثرة في المعادلة الثانية من~\eqref{eq:acetnet} تعطي ذاكرة ترابطية تُكمل النمط من جزئه & حساب سعة شبكة بأنماط متناثرة وقاعدة تغايرية: عند نسبة صغيرة $p$ من العصبونات النشطة تحفظ الشبكة أنماطًا أكثر بكثير مما عند $p=1/2$ & \cite{Tsodyks1988} & نظرية \\
8 & الكتابة عند $a$ منخفض تفسد الذاكرة: تكتب الشبكة ما تذكّرته لا ما رأته؛ وعند $a=0.1$ ليس الفساد أضعف منه عند $a=0$ & \texttt{models/\allowbreak acet\_memory.py}: $200$ عصبون، وخمسة أنماط في كل منها $20$، ونمط جديد يتقاسم $10$ عصبونات مع أحدها، و$10$ محاكاة. عند $a=0$ لا يُحفظ النمط الجديد، ويجرّ القديم $5.9$ عصبونًا غريبًا من $10$؛ وعند $a=0.1$ يُتذكَّر النمط الجديد ($0.97$)، لكن الاثنين يندمجان: الجديد يجرّ $7.3$ عصبونًا غريبًا، والقديم $7.9$؛ وبدءًا من $a=0.2$ أقل من واحد & اشتقاق في الفصل & نموذج \\
9 & القضية~\ref{prop:writeread}: لا يوجد مستوى $a$ تنجح عنده الكتابة والقراءة بالقدر نفسه (الشكل~\ref{fig:acetsim}) & البرنامج نفسه، ومعدل التعلّم $\eta a$: يُحفظ النمط الجديد كاملًا في عرض واحد عند $a\ge0.9$، و$0.8$ منه عند $a=0.75$، ولا يُحفظ عند $a\le0.5$. القراءة من نصف النمط: $1.0$ عند $a\le0.2$، و$0.96$ عند $a=0.3$، و$0.04$ عند $a=0.4$ & اشتقاق في الفصل & نموذج \\
10 & في الدماغ يبدّل سلك بثّ واحد، هو \nt{ACET}، بين الكتابة والقراءة & الفكرة مأخوذة من نماذج بُنيت على قياسات في الشرائح. وأكثر التجارب مباشرة عند الإنسان: حاجب المستقبلات المسكارينية سكوبولامين يسيء حفظ أزواج كلمات جديدة، وأشد ما يكون ذلك مع الأزواج المتداخلة مع القديمة، لكنه لا يسيء تذكّر الأزواج المحفوظة. لا يوجد قياس مباشر لوضعي الشبكة & \cite{HasselmoAndersonBower1992,Atri2004} & فرضية \\
11 & المرشح~\eqref{eq:kalman} أمثل؛ ووزن الدخل ومعدل التعلّم مقدار واحد $P/R$، وفي الحالة المستقرة $P=\sqrt{qR}$ والذاكرة $\sqrt{R/q}$ (القضية~\ref{prop:kalman}) & لا يتطلب تجربة: أمثلية مرشح كالمان-بيوسي نتيجة كلاسيكية في نظرية التقدير؛ والحالة المستقرة مشتقة في الفصل & \cite{KalmanBucy1961} & نظرية \\
12 & \nt{ACET} يخبر الشبكة بمقدار شبيه بـ$P$، أي عدم اليقين المتوقع & اقتُرح في نموذج احتمالي للانتباه. لا يوجد قياس مباشر لكون مستوى الأسيتيل كولين يتبع عدم اليقين في التقدير & \cite{YuDayan2005} & فرضية \\
13 & جزء من عصبونات \nt{ACET} يستجيب للمكافأة والعقاب في نحو عشرين ميلي ثانية، وبقوة تزداد كلما كان التعزيز أقل توقعًا & فأر، عصبونات كولينية في الدماغ الأمامي القاعدي معرَّفة بصريًا جينيًا: استجابة للمكافأة والعقاب بعد $18\pm3\ms$، يزداد حجمها مع المفاجأة في التعزيز، والاستجابات متشابهة في عصبونات نواتين & \cite{Hangya2015} & مُقاس \\
14 & \nt{NORA} يلاحظ التغيّر غير المتوقع في العالم، و\nt{ACET} يُبقي الشبكة في وضع الكتابة & تقسيم العمل مقترح في نموذج. وبصورة غير مباشرة: عند الإنسان يتتبع قطر الحدقة المرتبط بـ\nt{NORA} التغيرات ومعدل التعلّم. لا يوجد اختبار مباشر للزوج \foreignlanguage{english}{\nt{NORA}--\nt{ACET}} & \cite{YuDayan2005,Nassar2012} & فرضية \\
15 & في النوم البطيء تعيد الشبكة في وضع القراءة تشغيل ما كُتب نهارًا، وهذه الإعادات تنقله إلى الذاكرة الطويلة الأمد & تسجيلات مجموعات عصبونات الحُصين عند الجرذ: الخلايا التي عملت معًا نهارًا تنطلق معًا أكثر في النوم البطيء التالي، وبالترتيب نفسه، وهذا مقيس. وللنقل: كبت تموجات الحُصين بعد التعلّم يسيء الذاكرة المكانية، وإطالة التموجات التلقائية يحسّنها؛ أما أن السبب هو المستوى المنخفض لـ\nt{ACET} فغير مبرهن & \cite{WilsonMcNaughton1994,Skaggs1996,Girardeau2009,FernandezRuiz2019,Hasselmo1999} & فرضية \\
\bottomrule
\end{longtable}}

\section{الفصل~\ref{sec:synthesis}. الآلة كاملة}

لا يعرض فصل الخلاصة تجارب جديدة: كل سطر فيه يستند إلى الفصل الذي نوقش فيه الادعاء، وإلى المصادر نفسها؛ الراسخ هو ناقل \nt{GLUT} والتحكم في التدفق عبر \nt{GABA}، أما أدوار قنوات البث وعملها المشترك ففرضية في معظمها.

{\footnotesize
\setlength{\tabcolsep}{3pt}\renewcommand{\arraystretch}{1.15}
\begin{longtable}{>{\raggedright}p{0.5cm}>{\raggedright}p{4.0cm}>{\raggedright}p{5.6cm}>{\raggedright}p{2.3cm}>{\raggedright\arraybackslash}p{1.7cm}}
\toprule
م & القضية & التجربة وما قيس & المصدر & الحالة \\
\midrule\endfirsthead
\toprule
م & القضية & التجربة وما قيس & المصدر & الحالة \\
\midrule\endhead
1 & لـ$8.6\cdot10^{10}$ عصبون أربعة أنواع فقط من إشارات التحكم~\eqref{eq:machine} & عدّ نوى الخلايا في متجانس الدماغ (المجزِّئ المتساوي الخواص): لدى الرجل البالغ $86.1\pm8.1$ مليار عصبون & \cite{Azevedo2009} & مُقاس \\
2 & القضية~\ref{prop:architecture}، الطبقتان الأوليان: تسير البيانات على ناقل العناوين \nt{GLUT}، ويحدد المرسلُ علامةَ الوصلة، وتوجّه عصبونات \nt{GABA} التدفق وتطبّعه (الفصلان~\ref{sec:neurontypes} و\ref{sec:gaba}) & ناقل عصبي رئيسي واحد لكل عصبون (مراجعة للقياسات)؛ التثبيط الأمامي يضيّق النافذة التي يجمع فيها العصبون الهرمي مداخله؛ والقسمة على النشاط الكلي مقيسة في مناطق كثيرة & \cite{Strata1999,Pouille2001,Carandini2012} & مُقاس \\
3 & العناصر غير موثوقة: يفقد المشبك معظم النبضات (الجدول~\ref{tab:computer}، الفصل~\ref{sec:code}) & شرائح الحُصين، تنبيه أدنى: أكثر من نصف النبضات لا يستدعي استجابة في \foreignlanguage{english}{CA1}؛ استُبعد إخفاق العتبة والتوصيل على المحور، والسبب هو التحرير الاحتمالي للناقل العصبي & \cite{AllenStevens1994} & مُقاس \\
4 & يعمل الدماغ على $\sim20$ واطًا، بمعدل نبضة في الثانية تقريبًا لكل عصبون (الفصل~\ref{sec:code})؛ ولم يبنِ المهندسون بعدُ آلة كهذه & التقدير~\eqref{eq:energy} من التكاليف المقيسة: نحو $10^9$ جزيء \foreignlanguage{english}{ATP} لكل نبضة مع عمل المشابك (قشرة القوارض)؛ والتقدير المفصّل للقشرة يعطي ترددًا أدنى. حالة التقنية بحسب مراجعة & \cite{Attwell2001,Lennie2003,MehonicKenyon2022} & نظرية \\
5 & تعلّم معظم المشابك حاصل ضرب أثر محلي في عدد مشترك واحد $\delta$ (المعادلة الثانية في~\eqref{eq:machine}، الفصل~\ref{sec:dofa}) & عصبونات \nt{DOPA} لدى القرد تستجيب لخطأ التنبؤ بالمكافأة؛ وفي شرائح الجسم المخطط يكبّر الدوبامين الشويكة فقط إذا وصل بعد $0.3\text{--}2\,\text{s}$ من الدخل الغلوتاماتي، فالأثر مقيس & \cite{Schultz1997,Yagishita2014} & مُقاس \\
6 & سلك خطأ مستقل إلى كل مخرج لا يوجد إلا في دارة تعلّم الحركات (الفصل~\ref{sec:motorlearning}) & المخيخ: تزامن إشارة الليف المتسلق مع الدخل يُضعف هذا الدخل؛ ولدى القرد تحمل النبضات المعقدة لخلايا بوركنجي اتجاه خطأ الحركة وتستدعي التصحيح & \cite{Ito2001,Herzfeld2018} & مُقاس \\
7 & كل قناة بث حزمة من عدة خطوط (القضية~\ref{prop:bundle}) & في \nt{DOPA}: في المهمة الواحدة تحمل عصبونات مختلفة المكافأة والحركة وسمات المنبّه؛ والخطأ السريع مختلف عن المستوى البطيء للدوبامين. أما في \nt{SERO} و\nt{NORA} و\nt{ACET} فهذا التوزيع أقل دراسة & \cite{Engelhard2019,Hamid2016,Mohebi2019} & مُقاس \\
8 & \nt{SERO}: منظّم المستوى، وبحسب الفرضية الأفق $\tau_\gamma$ (الفصل~\ref{sec:serocomputer}) & التثبيط الذاتي: ناهضات مستقبلاته \foreignlanguage{english}{5-HT$_{1A}$} تثبّط عصبونات السيروتونين في نوى الرفاء، وهذا مُقاس. الأفق مقترح نظريًا؛ وأقوى تجربة أن التنشيط البصري الوراثي لهذه العصبونات لدى الفأر يطيل انتظار المكافأة المؤجلة & \cite{Sprouse1987,Doya2002,Miyazaki2014} & فرضية \\
9 & \nt{NORA}: الكسب $g$ يختار بين البحث والحسم (الفصل~\ref{sec:nora}) & ارتفاع نسبة الإشارة إلى الضجيج: لدى القرد اليقظ يكبت النورأدرينالين في القشرة السمعية النشاطَ الخلفي أكثر مما يكبت الاستجابة لأصوات أبناء جنسه، وهذا مُقاس؛ والكسب الضربي نموذج؛ والدور في الاختيار بين البحث والحسم نظرية & \cite{Foote1975NE,ServanSchreiber1990,AstonJones2005} & فرضية \\
10 & \nt{ACET}: يبدّل بين الكتابة والقراءة (الفصل~\ref{sec:acet}) & الأفعال الثلاثة (إضعاف الوصلات الداخلية، وتقوية الدخل، وتيسير اللدونة) مقيسة؛ وتبديل الأنماط تدعمه بصورة غير مباشرة تجربة السكوبولامين لدى الإنسان & \cite{HasselmoBower1992,Gil1997,Atri2004} & فرضية \\
11 & الصيغة~\eqref{eq:machine} تصف الآلة بكاملها & هي خلاصة نماذج الفصول من~\ref{sec:glutcomputer} إلى~\ref{sec:acet}، وكل جزء منها يستند إلى فصله؛ ولم تُختبر تجريبيًا بوصفها كلًّا & استنتاج في الفصل & فرضية \\
12 & تعمل المقابض منسَّقةً بلا تحكم مشترك: الخطأ، المفاجأة، الكتابة، العودة (الشكل~\ref{fig:timeline}) & لا توجد تجربة: المخطط نوعي. حلقات منفردة: نظرية تقسيم العمل بين \nt{NORA} و\nt{ACET}، وارتباط حدقة العين بسرعة التعلم لدى الإنسان & \cite{YuDayan2005,Nassar2012} & فرضية \\
13 & التعلم بإشارة مشتركة واحدة يبطؤ كلما كثرت العصبونات المؤثرة في النتيجة (القضية~\ref{prop:broadcastcost}) & حساب منحنيات التعلم لشبكة خطية: التعلم باضطراب المخارج أبطأ من الانحدار التدرجي بعدد مرات يساوي عدد المخارج & \cite{Werfel2005} & نظرية \\
14 & كيف تضبط القشرة الدارات المتعددة الطبقات مجهول؛ والمرشحون: رشقات النبضات، والحسابات في فروع العصبون، والتعلم الذاتي بالتنبؤ & الرشقات حاملةً للخطأ نموذج؛ ولم تستطع تكرار استجابة نموذج مفصّل لعصبون قشري إلا شبكة من $5\text{--}8$ طبقات، وهذا نموذج؛ والنبضات الكالسيومية في فروع عصبونات قشرة الإنسان التي تعطي ”أو الحصرية“ مقيسة في الشرائح؛ والتعلم الذاتي مراجعة للأدلة & \cite{Payeur2021,Beniaguev2021,Gidon2020,Keller2018} & فرضية \\
\bottomrule
\end{longtable}}

\section{الفصل~\ref{sec:sensors}. مداخل الآلة}

بنية المستشعرات مقيسة مباشرة، بتسجيل تيار الخلايا المفردة واهتزازات القوقعة وبعدّ خلايا الشبكية؛ وحد الحساسية وصيغة ضجيج الطلقات يتبعان من الفيزياء؛ أما أن شفرات المستشعرات مضبوطة على أكبر قدر من المعلومات ففرضية لها سند قوي.

{\footnotesize
\setlength{\tabcolsep}{3pt}\renewcommand{\arraystretch}{1.15}
\begin{longtable}{>{\raggedright}p{0.5cm}>{\raggedright}p{4.0cm}>{\raggedright}p{5.6cm}>{\raggedright}p{2.3cm}>{\raggedright\arraybackslash}p{1.7cm}}
\toprule
م & القضية & التجربة وما قيس & المصدر & الحالة \\
\midrule\endfirsthead
\toprule
م & القضية & التجربة وما قيس & المصدر & الحالة \\
\midrule\endhead
1 & المستشعر محوِّل: المستقبِل يعطي تيارًا، ومخرج الخلايا الحسية في العين والأذن غلوتامات على ناقل \nt{GLUT}؛ والخلايا الأولى لا تُصدر نبضات (الشكل~\ref{fig:transducer}) & العصي والمخاريط لدى السلحفاة تطلق حمضًا أمينيًا مثيرًا: تلتقطه قنوات غلوتاماتية على رقعة غشاء مفصولة. الخلايا الشعرية الداخلية في قوقعة الجرذ: التيارات في نهايات الألياف العصبية تمر عبر مستقبِلات \foreignlanguage{english}{AMPA} الغلوتاماتية، ويزداد تردد التحرير مع إزالة الاستقطاب المتدرجة للخلية حتى $150\,\text{Hz}$ & \cite{CopenhagenJahr1989,GlowatzkiFuchs2002} & مُقاس \\
2 & مستشعر الضوء في الظلام يستجيب لفوتون واحد بتيار نحو بيكوأمبير & قطب ماصّ على القطعة الخارجية لعصية العلجوم: الاستجابات للومضات الخافتة مكمّاة، والحدث $\approx1\,\text{pA}$ ($5\%$ من التيار المظلم) بانتشار $0.2\,\text{pA}$ وكفاءة $0.5$. والإنسان يُبلغ عن وصول فوتون واحد أكثر مما تقتضيه المصادفة & \cite{Baylor1979,Tinsley2016} & مُقاس \\
3 & مستشعر الصوت يلاحظ إزاحة أقل من نانومتر & طريقة موسباور، الغشاء القاعدي لقوقعة خنزير غينيا عند موضع $16\text{--}19\,\text{kHz}$: عند عتبة استجابة العصب السمعي سرعة الغشاء نحو $0.04\,\text{mm/s}$، أي إزاحة $\approx0.35\,\text{nm}$ (التحويل من عندنا). المقيس هو الغشاء لا حزمة المستشعر & \cite{Sellick1982,Hudspeth1989} & مُقاس \\
4 & الدقة في الظلام يحدّها ضجيج الطلقات للفوتونات~\eqref{eq:photonnoise}؛ وفي الضوء الخافت يطول التجميع & الصيغة تتبع من إحصاء بواسون. في العصية يطابق الضجيجُ في الضوء الخافت مجموعَ أحداث كمومية مستقلة؛ وعلى خلفية ساطعة تكون الاستجابة للومضة أصغر وأقصر. أما رقم ”أعشار الثانية“ فلم يُختبر هنا على حدة & \cite{Baylor1979} & نظرية \\
5 & مدى الدخل عشر رتب للضوء واثنتا عشرة للصوت، أما تردد النبضات فأقل من ثلاث رتب & للصوت: تزداد استجابة الغشاء القاعدي عند التردد المميِّز بميل يصل إلى $0.2\,\text{dB/dB}$ في النطاق $40\text{--}80\,\text{dB}$، والانضغاط ظاهر فوق $100\,\text{dB}$ أيضًا. أما أعداد الرتب نفسها فتقديرات معيارية، بلا مصدر في الفصل & \cite{Ruggero1997} & مُقاس \\
6 & استجابة المستشعر~\eqref{eq:nakarushton}، و$\sigma$ تتبع متوسط السطوع فتزيح المنحنى على المحور اللوغاريتمي (الشكل~\ref{fig:adaptation}) & طُبّقت الصيغة أول مرة على كمونات \foreignlanguage{english}{S} في شبكية السمك. مخاريط السلحفاة: الاستجابات تخضع لـ$V=I V_m/(I+\sigma)$ وتغطي $\sim3.5$ رتبة من السطوع؛ وبعد $2\text{--}3$ دقائق من الخلفية ينزاح المنحنى أفقيًا محافظًا على عرضه. الصلة بالقسمة على النشاط الكلي مراجعة & \cite{NakaRushton1966,NormannPerlman1979,Carandini2012} & مُقاس \\
7 & المستشعرات تنقل التغيرات، وشفراتها مضبوطة على أكبر قدر من المعلومات لكل نبضة (القضية~\ref{prop:sensors}) & المبدأ مقترح نظريًا. وأقوى دليل: لدى الذبابة تطابق خاصية ”التباين $\leftarrow$ الاستجابة“ لعصبونات الطبقة الأولى التوزيعَ التراكمي لتباينات المشاهد الطبيعية، وبذلك تُبلغ أكبر كمية من المعلومات & \cite{Barlow1961,Laughlin1981} & فرضية \\
8 & مستشعرات التشغيل والإطفاء شفرة سلكين؛ وكاميرا الأحداث تكررها في السيليكون & مراجعة للتجارب: قناتا التشغيل والإطفاء في الشبكية تنفصلان منذ المشبك الأول ويمكن إيقاف كل منهما على حدة. الكاميرا: $128\times128$ بكسل بلا إطارات، مدى $120\,\text{dB}$، وتأخير $15\,\text{\textmu s}$ & \cite{Schiller1992,Lichtsteiner2008,Gallego2022} & مُقاس \\
9 & في العين $\sim10^8$ مستشعر ضوء و$\sim10^6$ عصبون مخرج: ضغط بنحو $\sim100$ مرة & عدّ على شبكيات بشرية كاملة: في المتوسط $92$ مليون عصية و$4.6$ مليون مخروط؛ وخلايا المخرج (العقدية) من $0.7$ إلى $1.5$ مليون في عيون مختلفة & \cite{Curcio1990,CurcioAllen1990} & مُقاس \\
10 & لدى الفأر أكثر من ثلاثين قناة مخرج للعين & الفأر: تصوير كالسيومي ثنائي الفوتون لأكثر من $11\,000$ خلية مخرج وتجميع الاستجابات، فظهر أكثر من $30$ نوعًا وظيفيًا بفارق واضح. لدى الإنسان لم يُقَس العدد & \cite{Baden2016} & مُقاس \\
11 & عين خنزير غينيا تنقل $\sim875\,000$ بت في الثانية؛ والتحويل إلى عين الإنسان يعطي $\sim10^7$ بت في الثانية & خنزير غينيا، مصفوفة متعددة الأقطاب، حركة في مشاهد طبيعية: $6\text{--}13$ بت في الثانية لكل خلية، و$\sim875\,000$ بت في الثانية لنحو $\sim10^5$ خلية مخرج. وللإنسان ($\sim10^6$ خلية) $\sim10^7$ بت في الثانية، وهو تحويل المؤلفين & \cite{Koch2006} & مُقاس \\
12 & الأذن بنك مرشحات تمرير النطاق بخريطة ترددات شبه لوغاريتمية (الشكل~\ref{fig:filterbank}) & موضع أكبر اهتزاز للغشاء القاعدي يتوقف على التردد؛ ودالة ”التردد--الموضع“ شبه أُسية، وتصف بشكل واحد قواقع الإنسان والحيوانات الحية & \cite{Greenwood1990,Robles2001} & مُقاس \\
13 & المرنانات فعّالة: بعض المستشعرات يضخّم الاهتزازات الضعيفة آلاف المرات في السعة ($66\text{--}76\,\text{dB}$)، ومع ازدياد الشدة ينخفض الكسب & قياس الاهتزاز بالليزر عند قاعدة قوقعة الشنشيلة: للنغمات الضعيفة عند التردد المميِّز كسب $66\text{--}76\,\text{dB}$ نسبةً إلى الركاب؛ والموت يخفض الحساسية بـ$60\text{--}81\,\text{dB}$ (بـ$10^3\text{--}10^4$ مرة في السعة) ويزيل الانضغاط. مصدر القوة الخلايا الشعرية الخارجية & \cite{Ruggero1997,Robles2001} & مُقاس \\
14 & مضخّم الأذن موضوع عند حافة الإثارة الذاتية & حساب: قرب نقطة تفرّع هوبف لا مفر من انضغاط المدى والتوليف الحاد والنغمات التركيبية، وهي كل اللاخطيات المعروفة في السمع. أما أن القوقعة مضبوطة هكذا بالذات فلم يُبرهن مباشرة & \cite{Eguiluz2000} & فرضية \\
15 & في الترددات المنخفضة تسير النبضات في طور مع الصوت (لدى خنزير غينيا يضعف التقيد بدءًا من $\sim600\,\text{Hz}$ ويبقى ملحوظًا حتى $\sim3.5\,\text{kHz}$)، وبها يتحدد الاتجاه؛ وفي العالية تبقى شفرة الموضع & العصب السمعي لخنزير غينيا: يبدأ التقيد بالطور بالانخفاض قرب $600\,\text{Hz}$ ويختفي فوق $3.5\,\text{kHz}$؛ ولدى القط والقرد الحد أعلى بنحو أوكتاف. فرق زمن الوصول إلى الأذنين مراجعة & \cite{PalmerRussell1986,Grothe2010} & مُقاس \\
16 & بقية المستشعرات (الجدول~\ref{tab:sensors}): للشم $\sim400$ نوع من المستقبِلات، نوع لكل عصبون، وشفرة تركيبية؛ والذوق جزئيًا عبر \foreignlanguage{english}{ATP} والسيروتونين؛ واللمس مستشعرات سريعة التعوّد وبطيئته؛ والدفء والبرد منفصلان & الفأر: المستقبِل الواحد يتعرف على روائح كثيرة، والرائحة الواحدة على مستقبِلات كثيرة، والروائح المختلفة على مجموعات مختلفة. بلا مستقبِلات \foreignlanguage{english}{ATP} لا يستجيب عصب الذوق؛ والبراعم الذوقية تطلق السيروتونين. تسجيل ألياف مفردة من جلد يد الإنسان: أربعة أنواع بحسب سرعة التعوّد. وقناة البرد مختلفة عن قنوات الحرارة & \cite{Mombaerts2004,Malnic1999,Finger2005,Huang2005,JohanssonVallbo1979,McKemy2002} & مُقاس \\
\bottomrule
\end{longtable}}

\section{الفصل~\ref{sec:recognition}. التعرّف}

سرعة التعرّف، وبنية المراحل الأولى، والقراءة الخطية للجواب، والتعلم من استمرارية الزمن مقيسة لدى الإنسان والقرد؛ واختزال السلسلة كلها إلى عمليتين والتعلم من الفيديو مختبَران على نماذج الكتاب؛ أما تفسيرات الخدع فتتراوح بين الراسخ الأسس والمختلف فيه.

{\footnotesize
\setlength{\tabcolsep}{3pt}\renewcommand{\arraystretch}{1.15}
\begin{longtable}{>{\raggedright}p{0.5cm}>{\raggedright}p{4.0cm}>{\raggedright}p{5.6cm}>{\raggedright}p{2.3cm}>{\raggedright\arraybackslash}p{1.7cm}}
\toprule
م & القضية & التجربة وما قيس & المصدر & الحالة \\
\midrule\endfirsthead
\toprule
م & القضية & التجربة وما قيس & المصدر & الحالة \\
\midrule\endhead
1 & المقارنة بالنموذج بكسلًا ببكسل مع الإزاحة لا تصيب بأفضل من رمي عملة؛ وزوج المراحل~\eqref{eq:scstage} يتعرف على الشكل في أي مكان & \foreignlanguage{english}{\texttt{models/\allowbreak recognition\_models.py}}، ”شبكية“ من $24$ نقطة، $4000$ محاولة: النموذج $0.49$، $0.51$، $0.52$ عند نسبة نقاط مقلوبة $0$، $0.1$، $0.2$؛ وزوج $S$ و$C$: $1.00$، $0.86$، $0.70$ & استنتاج في الفصل & نموذج \\
2 & يُتعرَّف على حيوان في صورة في $\sim150\ms$، بمرور واحد عبر نحو عشر مراحل من $10\text{--}20\ms$ لكل منها & الإنسان، مهمة ”هل يوجد حيوان أم لا“ على صور معروضة $20\ms$: الكمونات المستثارة للجوابين تفترق بعد $\sim150\ms$. القرد: زمن الاستجابة الأولى يزداد من مرحلة إلى أخرى (\foreignlanguage{english}{V1} أولًا، ثم \foreignlanguage{english}{V2} و\foreignlanguage{english}{V4}). عدد المراحل و”صفر نبضات أو نبضة واحدة لكل مرحلة“ استنتاج من هذه الأرقام & \cite{Thorpe1996,Schmolesky1998} & مُقاس \\
3 & المرحلة $S$ ”و“ على الأجزاء، والمرحلة $C$ الأكبر على المواضع (القضية~\ref{prop:conveyor}) & القط، القشرة البصرية الأولى: الخلايا البسيطة تستجيب لحافة باتجاه معيّن في مكان معيّن، والمعقدة للاتجاه نفسه في منطقة أوسع. تسجيل داخل خلوي للخلايا المعقدة: الاستجابة لشريطين قريبة لدى خلايا كثيرة من أكبر الاستجابتين، وفي المتوسط أقرب ما تكون إلى عملية \foreignlanguage{english}{max}. الأكبر عبر التثبيط المتبادل هو القضية~\ref{prop:wta}، والقسمة على النشاط الكلي مراجعة & \cite{HubelWiesel1962,Lampl2004,Carandini2012} & مُقاس \\
4 & كلما صعدنا في السلسلة كبرت التراكيب وتعقدت، وقلّ اعتماد الجواب على الموضع & القرد، تبسيط المنبهات إلى ”السمة الحرجة“: خلايا \foreignlanguage{english}{V4} والقشرة الصدغية السفلية الخلفية تستجيب لتراكيب معقدة من الشكل واللون والنسيج في مجال صغير؛ وفي القشرة الصدغية السفلية الأمامية المجال أكبر. وتناوب $S$ و$C$ في النموذج يعيد إنتاج هذه الصورة & \cite{KobatakeTanaka1994,Riesenhuber1999} & مُقاس \\
5 & الشبكات الالتفافية تكرر هذا التناوب وتتنبأ باستجابات المراحل العليا أفضل من غيرها؛ والشيء يُقرأ خطيًا من ترددات مجموعة من العصبونات & شبكة مدرَّبة على التعرّف بلا مواءمة مع الدماغ: الطبقة العليا تتنبأ باستجابات عصبونات القشرة الصدغية السفلية لدى القرد، والوسطى باستجابات \foreignlanguage{english}{V4}. مصنِّف خطي على نشاط $\sim100$ خلية مختارة عشوائيًا في نوافذ تبدأ من $12.5\ms$ يتعرف على الشيء وفئته مع اختلاف الموضع والحجم & \cite{Fukushima1980,Yamins2014,Hung2005} & مُقاس \\
6 & قاعدة هيب مع أثر~\eqref{eq:traceinvariance} تعلّم اللاتغيّر من فيديو بلا أسماء إذا لم يكن الأثر أقصر من $\sim20\ms$ & فكرة السمات البطيئة والقاعدة مع أثر نظرية. \foreignlanguage{english}{\texttt{models/\allowbreak recognition\_models.py}}، عصبونا $C$، $16$ مدخلًا، $10$ تشغيلات: التوقف بين الأشياء $0.3\,\text{s}$. عند $\tau_{\text{tr}}=5\ms$ التطابق مع الشكل $0.52$ في المتوسط، وعند $10\ms$ $0.56$؛ وعند $20$ و$50$ و$200\ms$ $1.00$ في كل التشغيلات (عند $30\ms$ يخطئ تشغيل واحد من عشرة) & \cite{Foldiak1991,WiskottSejnowski2002} & نموذج \\
7 & في الدماغ يُتعلَّم اللاتغيّر من استمرارية الزمن & القرد: استُبدل الشيء خفيةً أثناء قفزة العين إلى مكان معيّن من المجال؛ فتغيّر لاتغيّرُ الموضع لدى عصبونات القشرة الصدغية السفلية وفق الاستبدال، تغيرًا دالًّا بعد ساعة. أما أنها القاعدة الرئيسية لتعلم البصر ففرضية & \cite{LiDiCarlo2008} & مُقاس \\
8 & الحركة: كواشف الاتجاه، ثم زوج ”و“/”أو“ نفسه على مناطق كبيرة & كواشف الاتجاه في شبكية الأرنب؛ وفي المنطقة \foreignlanguage{english}{MT} لدى القرد كل العصبونات الـ$110$ المسجلة انتقائية للاتجاه، والاستجابة للحركة أقوى منها في \foreignlanguage{english}{V1} & \cite{Barlow1965,Albright1984} & مُقاس \\
9 & المراحل العليا ترسل التنبؤ إلى أسفل، والخطأ يصعد إلى أعلى & النموذج يفسر التأثيرات خارج المجال الكلاسيكي في القشرة البصرية الأولى؛ وبعض المشاهدات يتفق معه (مراجعة)، ولم يُثبت بوصفه البنية العامة لخط الأنابيب & \cite{RaoBallard1999,Keller2018} & فرضية \\
10 & الصور الصعبة تتطلب وصلات راجعة وعدة دورات & القرد: في الصور ”الصعبة“ يتكوّن القرار في القشرة الصدغية السفلية متأخرًا بنحو $\sim30\ms$؛ وهذه الاستجابات المتأخرة تتنبأ بها الشبكة الأمامية تنبؤًا ضعيفًا، والشبكات ذات الوصلات الراجعة أو الأعمق تنبؤًا أفضل & \cite{Kar2019} & مُقاس \\
11 & تزييف غير ملحوظ في البكسلات يخدع الشبكة؛ والبصر البشري أمتن بكثير، لكنه ليس محصنًا & على الشبكات: اضطراب صغير مختار خصيصًا يغيّر الجواب؛ ومثال ”الباندا $\leftarrow$ الجِبّون“ من العمل الثاني. ولدى الإنسان عند العرض القصير تزيح التزييفاتُ التي تنتقل جيدًا بين الشبكات الاختيارَ & \cite{Szegedy2014,Goodfellow2015,Elsayed2018} & مُقاس \\
12 & خدع التوقع: الدخل غير الموثوق يتنحى للتوقع، فالباهت يتحرك ”أبطأ“، والفستان يُرى على نحو مختلف (القسم~\ref{sec:illusions}) & المحزّزات الأقل تباينًا تبدو أبطأ حركة. استطلاع لـ$1401$ شخص: يُرى الفستان في ثلاثة أشكال ثابتة، وتلميح الإضاءة يبدّل اللون؛ ولدى $\sim13\,000$ مشاهد يحسم الأمرَ افتراضُ الإضاءة. والنموذج البايزي مع توقع الحركات البطيئة يفسر عددًا من خدع الحركة & \cite{Thompson1982,LaferSousa2015,Wallisch2017,Weiss2002,Purves2011} & فرضية \\
13 & ضبط الكسب: الشلال، والميل، والتباين؛ وشبكة هيرمان ليست تثبيط الجيران & كواشف الاتجاه في شبكية الأرنب تخفض ترددها عند الحركة الطويلة، وبعد التوقف تهبط دون الخلفية. خدعة الميل يعيد إنتاجها نموذج بالقسمة على نشاط الجيران. وتعديلات شبكة هيرمان التي لا تغيّر استجابة عصبونات مخرج الشبكية تزيل البقع، فرُفض التفسير بتثبيط الجيران & \cite{BarlowHill1963,Schwartz2009,SchillerCarvey2005} & مُقاس \\
14 & تعويض التأخيرات: الشيء المتحرك يبدو متقدمًا على الومضة & الخدعة مقيسة. والتجارب النفسية الفيزيائية تناقض كلًّا من مدّ الحركة إلى الأمام وفرق التأخيرات؛ واقتُرح تفسير بأثر رجعي، بحسب الأحداث خلال $\sim80\ms$ بعد الومضة. والخلاف لم يُحسم & \cite{Nijhawan1994,EaglemanSejnowski2000} & فرضية \\
15 & ”الثعابين“ الساكنة تدور: فرق التأخيرات~\eqref{eq:latency} تقرؤه كواشف الاتجاه & الخدعة تعمل بلا لون أيضًا؛ وأزواج العناصر المتجاورة تولّدها منفردة؛ وعصبونات قشرة القرد الانتقائية للاتجاه تعطي استجابة موجَّهة لهذه الأزواج الساكنة. ولدى الإنسان يطلق الدورانَ القفزاتُ الدقيقة للعين والرمش & \cite{Conway2005,OteroMillan2012} & مُقاس \\
16 & الصور المزدوجة: تثبيط متبادل، وتعب، وضجيج؛ والتبدلات يسببها الضجيج في الغالب & نموذج مجموعات عصبونات متنافسة: التبدلات فيه يسببها \emph{الضجيج}، ولا تحدث بدونه؛ ويفسر ازدياد تكرار التبدلات مع شدة المنبه. والتعب هنا دوره أصغر & \cite{MorenoBote2007} & فرضية \\
17 & بعض الخدع نتيجة للمسألة التي تتعلمها الآلة & الشبكة المدرَّبة على التنبؤ بالإطار التالي من الفيديو ”ترى“ دوران الحلقات الساكنة ولا تراه في الصور الضابطة؛ والشبكات المخصصة لتنقية الصور وزيادة حدتها تنخدع بخدع اللون والسطوع كالإنسان & \cite{Watanabe2018,GomezVilla2020} & مُقاس \\
\bottomrule
\end{longtable}}

\section{الفصل~\ref{sec:muscles}. المخرج إلى العضلات}

بنية المخرج، أي الوحدات الحركية، وهامش أمان المشبك، والتشغيل بحسب الحجم، وزوج العضلات، وحلقة التمدد، ومولّدات الإيقاع، مقيسة مباشرة؛ وشرط استقرار الحلقة ذات التأخير مبرهنة؛ والنموذج الداخلي للجسم فرضية راسخة الأسس؛ وأرقام الأشكال حساب بنموذج الكتاب.

{\footnotesize
\setlength{\tabcolsep}{3pt}\renewcommand{\arraystretch}{1.15}
\begin{longtable}{>{\raggedright}p{0.5cm}>{\raggedright}p{4.0cm}>{\raggedright}p{5.6cm}>{\raggedright}p{2.3cm}>{\raggedright\arraybackslash}p{1.7cm}}
\toprule
م & القضية & التجربة وما قيس & المصدر & الحالة \\
\midrule\endfirsthead
\toprule
م & القضية & التجربة وما قيس & المصدر & الحالة \\
\midrule\endhead
1 & الوحدة الحركية: عصبون \nt{ACET} واحد يتحكم في مجموعة من الألياف، من بضع مئات إلى نحو ألفين؛ وفي العضلات ذات الضبط الدقيق تكون الوحدات أصغر & عضلات الإنسان، عدّ المحاور الحركية والألياف العضلية: في المتوسط نحو $340$ ليفًا لكل عصبون في العضلة بين العظام في اليد، ونحو $410$ في العضلة العضدية الكعبرية، ونحو $1930$ في عضلة الساق التوأمية الإنسية. ولا يذكر الفصل عددًا دقيقًا لأصغر الوحدات. & \cite{Feinstein1955mu} & مُقاس \\
2 & المشبك على العضلة يطلق من الناقل العصبي أضعاف ما يلزم لإطلاق الليف & مراجعة للقياسات على المشابك العصبية العضلية: هامش الأمان (نسبة الناقل المُطلَق إلى العتبي) لدى الثدييات البالغة عادة $3\text{--}5$؛ ولدى الإنسان يقوم الهامش على طيات غشاء المتلقي أكثر مما يقوم على حجم الإطلاق. & \cite{WoodSlater2001} & مُقاس \\
3 & النفضة~\eqref{eq:twitch} تتصاعد خلال زمن $T_c$ من رتبة $50\ms$ & وحدات حركية مفردة في يد الإنسان عند جهد إرادي، والقوة مُتوسَّطة على نبضات الوحدة: زمن التصاعد $30\text{--}100\ms$، وهو أقل من $70\ms$ لدى $80\,\%$ من الوحدات. والدالة $(t/T_c)e^{1-t/T_c}$ تقريب من نموذج لمجمع الوحدات. & \cite{MilnerBrown1973rec, Fuglevand1993} & مُقاس \\
4 & مجموع النفضات~\eqref{eq:forcesum}: عند $5$ و$15$ و$40$ نبضة في الثانية تكون القوة $0.2\text{--}1.1F_1$ و$1.8\text{--}2.2F_1$ ونحو $5.4F_1$ (الشكل~\ref{fig:tetanus}) & حساب بـ\foreignlanguage{english}{\texttt{models/\allowbreak muscle\_\allowbreak models.py}}، $T_c=50\ms$: $0.21\text{--}1.10$ و$1.76\text{--}2.19$ و$5.28\text{--}5.49\,F_1$ في آخر $0.3\,\text{s}$. والجمع الخطي تبسيط: لا تشبّع للعضلة الحقيقية في النموذج. & \foreignlanguage{english}{\texttt{models/\allowbreak muscle\_\allowbreak models.py}} & نموذج \\
5 & تُشغَّل الوحدات بحسب الحجم؛ وأضعفها أضعف من أقواها مئة مرة & الجذور البطنية لدى القط: مع دخل منعكس متزايد تفرغ أولًا العصبونات ذات النبضات الصغيرة في الجذر، أي الصغيرة. العضلة بين العظام في يد الإنسان: قوة نفضة الوحدات من $0.1$ إلى $10\,\text{g}$، وتزداد شبه خطيًا مع القوة التي تُشغَّل عندها الوحدة. & \cite{Henneman1957, MilnerBrown1973rec} & مُقاس \\
6 & خطوة القوة تتناسب مع القوة، $\Delta F/F\approx q-1$~\eqref{eq:sizeprinciple}: فاصلة عائمة & مستنتج في الفصل من المتوالية الهندسية للقوى. ويتفق مع التجربة: عند الجهود الكبيرة يُشغَّل لزيادة القوة نفسها عدد أقل بكثير من الوحدات الجديدة. & استنتاج في الفصل؛ \cite{MilnerBrown1973rec} & نظرية \\
7 & انتشار القوة يزداد بما يتناسب مع القوة نفسها & جهد إرادي متساوي القياس لعضلة الإبهام: الانحراف المعياري للقوة يزداد خطيًا مع متوسط القوة؛ وعند الجهد نفسه المستثار بالتنبيه الكهربائي يبقى الانتشار ثابتًا. نموذج المجمع: الازدياد الخطي يعطيه تشغيل الوحدات بترتيب القوة. & \cite{Jones2002sdn} & مُقاس \\
8 & نعومة الحركات نتيجة للضجيج المعتمد على الإشارة & نموذج: المسار الذي يقلل انتشار النقطة النهائية تحت هذا الضجيج يعيد إنتاج المسارات المقيسة للعين واليد. ولا توجد تجربة مباشرة تميّز هذا السبب من غيره. & \cite{HarrisWolpert1998} & فرضية \\
9 & الفرق بين قوتي زوج العضلات يحدد العزم، ومجموعهما الصلابة~\eqref{eq:antagonists} & نظرية التحكم في الممانعة بالشد المشترك. ذراع الإنسان: صلابة كل مفصل، مقيسة بدفعات صغيرة، تتناسب تقريبًا مع العزم فيه؛ والنسب المختلفة للشد المشترك تغيّر شكل قطع الصلابة الناقص لليد واتجاهه. & \cite{Hogan1984, GomiOsu1998stiff} & مُقاس \\
10 & مستشعر التمدد يثير عصبون \nt{ACET} الخاص به، ويطفئ المضادة عبر وسيط مثبِّط (الشكل~\ref{fig:antagonists}) & الحبل الشوكي لدى القط: وسيط واحد تثيره ألياف مستشعرات التمدد يعطي كمونات تثبيطية أحادية المشبك في عصبونات \nt{ACET} للعضلة المضادة، بتأخير مشبكي $0.28\text{--}0.42\ms$. ولم يُحدَّد في هذه التجربة الناقلُ العصبي للوسيط (الغليسين). & \cite{Jankowska1972Ia} & مُقاس \\
11 & تأخير الحلقة: الشوكية $25\text{--}30\ms$، وعبر القشرة أطول بمرة ونصف إلى ثلاث مرات، وعبر العين $100\text{--}200\ms$ & ذراع الإنسان، دفعة للمفصل: استجابة عضلية قصيرة بعد $20\text{--}45\ms$، وطويلة بعد $45\text{--}105\ms$، وإرادية بعد $120\text{--}180\ms$. إزاحة الموضع المرئي للإصبع أثناء الحركة: التصحيح بعد $160\ms$ في المتوسط. & \cite{Pruszynski2008rapid, SaundersKnill2003vis} & مُقاس \\
12 & $\dd x/\dd t=-Gx(t-d)$ مستقر عند $Gd<\pi/2$~\eqref{eq:delaystable}؛ وعلى الحد التردد $1/(4d)$ & مبرهنة عن جذور معادلة ذات تأخير. تحقق بـ\foreignlanguage{english}{\texttt{models/\allowbreak muscle\_\allowbreak models.py}}، $d=30\ms$: عند $3\,\text{s}$ السعة $3\cdot10^{-6}$ عند $G=40$، و$0.13$ عند $G=50$، و$38$ عند $G=55$ في الثانية (الحد $52$). & \cite{Hayes1950} & نظرية \\
13 & الرعشة الفسيولوجية $8\text{--}12\,\text{Hz}$؛ وحلقة التمدد أحد مصادرها & مراجعة للقياسات: الرعشة الدقيقة في نطاق نحو $10\,\text{Hz}$ موجودة لدى الأصحاء؛ ومنشؤها مختلط: تذبذبات مركزية، وخصائص تفريغ الوحدات، والرنين الميكانيكي، ورنين الحلقة المنعكسة، ويغلب أحدها أو غيره بحسب الظروف. & \cite{McAuleyMarsden2000} & فرضية \\
14 & يتنبأ الدماغ بنتيجة الأمر بالنموذج الداخلي للجسم & يمسك الإنسان ثقلًا بالقرص ويحرك ذراعه: مع الحمل القصوري واللزج والمرن تتغير قوة القبض متزامنة مع الحمل، لا بتأخير الحلقة، أي أن الحمل متنبَّأ به. وتجمع المراجعة تجارب كهذه؛ ولا ملاحظة مباشرة للنموذج نفسه في العصبونات. & \cite{FlanaganWing1997grip, WolpertGhahramani2000} & فرضية \\
15 & إيقاع المشي والتنفس والمضغ تولّده شبكات محلية عند دخل ثابت & مراجعة للتجارب: الجهاز العصبي المعزول (الحبل الشوكي، العقد) يُصدر نمطًا حركيًا إيقاعيًا بلا دخل إيقاعي وبلا تغذية راجعة من العضلات. & \cite{MarderBucher2001} & مُقاس \\
16 & التثبيط المتبادل مع التعب~\eqref{eq:halfcentre} يعطي إيقاعًا بدور $0.84\,\text{s}\approx2\tau_a$ (الشكل~\ref{fig:halfcentre}) & مبرهنة: الشبكة ذات التثبيط المتبادل والتكيّف التي لا حالة مستقرة لها عند دخل ثابت تتذبذب حتمًا. حساب بـ\foreignlanguage{english}{\texttt{models/\allowbreak muscle\_\allowbreak models.py}} ($I=1$، $\beta=2$، $b=2$، $\tau=20\ms$، $\tau_a=0.4\,\text{s}$): الدور $0.84\,\text{s}$. أما أن المولّدات الحقيقية مبنية هكذا بالذات فلم يُبيَّن. & \cite{Matsuoka1985osc} & نموذج \\
\bottomrule
\end{longtable}}

\section{الفصل~\ref{sec:pain}. الألم}

مستشعرات الألم، والسلكان، والبوابة في الحبل الشوكي، ومقبض الصوت الهابط، والتحسّس، والاستيلاء على الانتباه مقيسة مباشرة؛ وصيغتا البوابة والتحسّس نموذجان مبسّطان من الكتاب؛ أما القاعدة التي تختار بها الآلة علوّ صوت الإنذار ففرضية.

{\footnotesize
\setlength{\tabcolsep}{3pt}\renewcommand{\arraystretch}{1.15}
\begin{longtable}{>{\raggedright}p{0.5cm}>{\raggedright}p{4.0cm}>{\raggedright}p{5.6cm}>{\raggedright}p{2.3cm}>{\raggedright\arraybackslash}p{1.7cm}}
\toprule
م & القضية & التجربة وما قيس & المصدر & الحالة \\
\midrule\endfirsthead
\toprule
م & القضية & التجربة وما قيس & المصدر & الحالة \\
\midrule\endhead
1 & عتبة عالية: عتبات التسخين لمستشعرات الألم المختلفة من $41^\circ$ إلى $46\text{--}48^\circ$ & تسجيل ألياف مفردة. جلد القرد: العتبة الوسيطة للمستشعرات البطيئة (\foreignlanguage{english}{C}) $41^\circ$، وللمستشعرات السريعة من النوع الثاني $46^\circ$، ومن النوع الأول فوق $53^\circ$. جلد الإنسان: مستشعرات \foreignlanguage{english}{C} التي تستجيب للضغط أيضًا $40.7^\circ$، والتي لا تستجيب له $48^\circ$. & \cite{Treede1995heat, Weidner1999cfib} & مُقاس \\
2 & مستشعرات الألم تتعوّد أضعف بكثير من مستشعرات اللمس، وتصبح بعد إصابة خفيفة أشد حساسية؛ وضعف الاستجابة بعد تسخين قوي مقيس لدى نوع واحد & حرق جلدي خفيف ($50^\circ$، $100\,\text{s}$): بعد $5\text{--}10$ دقائق تنخفض عتبة الألم لدى البشر وعتبة مستشعرات \foreignlanguage{english}{C} لدى القرد بـ$2\text{--}6^\circ$؛ ولا يوجد هذا الانزياح لدى مستشعرات الجلد الأخرى. ولدى المستشعرات السريعة من النوع الثاني تُكبت الاستجابة بعد التسخين القوي. ومقارنة التعوّد بمستشعرات اللمس تقدير عام، لا تسنده هنا تجربة منفصلة. & \cite{LaMotte1982hyper, Treede1995heat} & مُقاس \\
3 & سلكان: السريع $5\text{--}30\,\text{m/s}$، والبطيء $0.5\text{--}2\,\text{m/s}$؛ وزمن الوصول $t=L/v$~\eqref{eq:paintime} & سرعة التوصيل: مستشعرات الألم السريعة لدى القرد في المتوسط $15$ و$25\,\text{m/s}$ (نوعان)؛ ومستشعرات \foreignlanguage{english}{C} لدى الإنسان $0.8\text{--}1.0\,\text{m/s}$. وعند $L=1\,\text{m}$ يكون ذلك عشرات الملّي ثانية ونحو ثانية. & \cite{Treede1995heat, Weidner1999cfib} & مُقاس \\
4 & يصل الألم مرتين: أولًا السريع الدقيق، ثم الكليل الطويل & زمن رد الفعل لتسخين ساعد الإنسان: من $1100$ إلى $400\ms$ مع ارتفاع الحرارة؛ والتسخين القوي للجلد المشعر يعطي ألمًا مزدوجًا، ولجلد الراحة لا. الألم الأول لا تحمله إلا ألياف أسرع من $6\,\text{m/s}$، ويوافقه بحسب الكمون المستشعرات السريعة من النوع الثاني، وهي غائبة في الراحة. وتقسيم الأدوار (”أين“ و”ما مدى السوء“) تأويل. & \cite{CampbellLaMotte1983first, Treede1995heat} & مُقاس \\
5 & البوابة: عصبون المخرج في الحبل الشوكي يتلقى الألم واللمس، والوسيط المثبِّط يثيره اللمس (الشكل~\ref{fig:gate}) & مخطط البوابة مقترح نظريًا. فئران، وسم جيني وإزالة مجموعات من العصبونات: العصبونات المثيرة الحاملة للسوماتوستاتين لازمة للألم الميكانيكي؛ والعصبونات المثبِّطة الحاملة للداينورفين تمنع دخل مستشعرات اللمس من الوصول إليها، وبدونها تستدعي اللمسة الخفيفة ألمًا. & \cite{MelzackWall1965, Duan2014} & مُقاس \\
6 & اللمس يخفت الألم & الإنسان، نبضات ليزر مؤلمة على اليد: اللمس المتزامن يخفض تقدير الألم والقدرة على تمييز طاقة النبضات؛ ويضعف الأثر مع المسافة بين نقطة اللمس ونقطة الألم داخل القطعة الجلدية الواحدة. & \cite{Mancini2014touch} & مُقاس \\
7 & افتراض نظرية البوابة: الألم القوي يطفئ بنفسه الوسيط المثبِّط، فتنفتح البوابة على مصراعيها & معلَّم في الفصل افتراضًا للنظرية الأصلية. والعمل على الفئران يصف كيف تمنع البوابة اللمس من دخول قناة الألم؛ ولم يُبيَّن فيه إطفاء الوسيط بالألم. & \cite{MelzackWall1965} & فرضية \\
8 & البوابة~\eqref{eq:gate}: العتبة $0.18$ بلا لمس، و$0.86$ مع اللمس، و$0.11$ عند $g=2$؛ وعند $\nu=1$ يخفض اللمس المخرج من $1$ إلى $0.25$، وعند $\nu=2$ لا يؤثر؛ واللمس نفسه لا يمر (الشكل~\ref{fig:gatecurves}) & الحساب بـ\foreignlanguage{english}{\texttt{models/\allowbreak pain\_\allowbreak gate.py}} بأوزان النص يعيد إنتاج كل هذه الأرقام. & \foreignlanguage{english}{\texttt{models/\allowbreak pain\_\allowbreak gate.py}} & نموذج \\
9 & المسارات الهابطة من جذع الدماغ تغيّر كسب البوابة في الاتجاهين & الجرذ، النخاع المستطيل، تسجيل عند تسخين الذيل: بعض العصبونات تفرغ قبل السحب (”تشغيل“)، وبعضها يصمت (”إطفاء“)؛ والتنبيه الضعيف لهذه المنطقة يكبت المنعكس. مراجعة: الدارات نفسها تيسّر نقل الألم وتثبطه، والأفيونيات تعمل عبرها. & \cite{Fields1983rvm, Fields2004} & مُقاس \\
10 & توقع الراحة يخفت الألم عبر القناة الأفيونية & أناس بألم حاد: لدى من خفّض توقعُ الراحة ألمَهم أعاد حجبُ المستقبِلات الأفيونية الألمَ إلى مستوى الآخرين؛ ولدى الآخرين لم يغيّر الحجب الألم. & \cite{Levine1978} & مُقاس \\
11 & يختار الدماغ علوّ صوت الإنذار بحسب جدوى المقاطعة & لا توجد تجربة قيست فيها قاعدة اختيار علو الصوت. & --- & فرضية \\
12 & التحسّس: بعد الإصابة يزداد مخرج البوابة وتنخفض العتبة، جزئيًا بفعل الحبل الشوكي؛ ويدوم حتى بعد توقف المنبّه المؤذي & الجرذ: بعد إصابة الجلد تنخفض عتبة منعكس الثني وتزداد الاستجابة؛ ويبيّن التحليل الكهربائي الفسيولوجي أن جزءًا من الزيادة ينشأ في الحبل الشوكي نفسه. والمنبّه المؤذي يُحدث اضطرابات طويلة في الحساسية تدوم بعد المنبّه نفسه. ولا يذكر الفصل زمن الانحسار الدقيق. & \cite{Woolf1983} & مُقاس \\
13 & التحسّس~\eqref{eq:sensitization} تغذية راجعة موجبة؛ ومع حلقة قوية قد ينغلق الإنذار على نفسه بعد الالتئام & ينتج من النموذج~\eqref{eq:sensitization} (تمرين في آخر الفصل). ولا توجد تجربة تبيّن أن الألم المستمر بعد الالتئام حلقة منغلقة من هذا النوع بالذات. & استنتاج في الفصل & فرضية \\
14 & الألم مكافأة سالبة، والتعلم منه يتبع خطأ الفروق الزمنية & تصوير بالرنين، بشر، سلاسل من نُذر الألم: نشاط الجسم المخطط البطني والقشرة الجزيرية الأمامية يطابق إشارات التعلم بالفروق الزمنية. القرد: بعض عصبونات \nt{DOPA} يُثار بنفخة هواء مزعجة وبنذيرها، وبعضها يُثبَّط. & \cite{Seymour2004, Matsumoto2009} & مُقاس \\
15 & الألم يقاطع العمل الجاري & بشر، منبّه كهربائي مؤلم واختبارات صوتية: الاستجابة لاختبار بعد $100\ms$ من بدء الألم أبطأ بوضوح، وبعد $1.5\,\text{s}$ لا فرق عن الضابط. وتجمع المراجعة تجارب كهذه في نموذج للاستيلاء على الانتباه. & \cite{Crombez1996disrupt, EcclestonCrombez1999} & مُقاس \\
16 & السحب يُغلق في الحبل الشوكي & الجرذ: عصبونات الطبقات العميقة من القرن الخلفي تكرر النمط المكاني لمنعكس السحب لعضلات منفردة؛ ولا يمكن إثارتها بالتنبيه المعاكس الاتجاه من الجزء العنقي، أي أنها وسطاء شوكية. & \cite{Schouenborg1995wr} & مُقاس \\
\bottomrule
\end{longtable}}

\section{الفصل~\ref{sec:motorlearning}. تعلّم الحركة}

قاعدة المربعات الصغرى وفائدة سلك الخطأ المستقل مبرهنتان؛ وبنية الدارة ذات سلك الخطأ، والإضعاف عند التزامن، وضبط الأثر على التأخير، والأثر اللاحق، والعودة التلقائية للمهارة مقيسة؛ أما أن الدارة كلها تعمل تعلمًا موجَّهًا وتبني نموذجًا داخليًا ففرضية رائدة؛ وأرقام نموذج العمليتين حساب بنموذج الكتاب.

{\footnotesize
\setlength{\tabcolsep}{3pt}\renewcommand{\arraystretch}{1.15}
\begin{longtable}{>{\raggedright}p{0.5cm}>{\raggedright}p{4.0cm}>{\raggedright}p{5.6cm}>{\raggedright}p{2.3cm}>{\raggedright\arraybackslash}p{1.7cm}}
\toprule
م & القضية & التجربة وما قيس & المصدر & الحالة \\
\midrule\endfirsthead
\toprule
م & القضية & التجربة وما قيس & المصدر & الحالة \\
\midrule\endhead
1 & القاعدة~\eqref{eq:lms} تسير في المتوسط مع تدرج مربع الخطأ & نتيجة من نظرية المرشحات التكيّفية: التصحيح المتناسب مع الدخل وخطأ المخرج انحدار تدرجي عشوائي لمتوسط مربع الخطأ. & \cite{WidrowHoff1960} & نظرية \\
2 & مع سلك خطأ لكل مخرج لا تتوقف السرعة على عدد المخارج؛ ومع عدد مشترك واحد تنخفض مع عدد العصبونات الضاجّة (القضية~\ref{prop:teachers}) & منحنيات التعلم لشبكة خطية: التعلم باضطرابات عشوائية للعصبونات أبطأ من التدرج الدقيق بعدد مرات يساوي عدد العصبونات المضطربة. & \cite{Werfel2005} & نظرية \\
3 & الدارة ذات سلك الخطأ تعمل تعلمًا موجَّهًا (القضية~\ref{prop:errorwire}) & نظريات مستنتجة من بنية الدارة. تجارب الأسطر من 7 إلى 10 تتفق معها؛ أما أن الدارة كلها تعمل هكذا بالذات فلم يُبرهن. & \cite{Marr1969, Albus1971} & فرضية \\
4 & طبقة التوسيع: نحو $7\cdot10^{10}$ عصبون \nt{GLUT} صغير، أكثر مما في بقية الدماغ كله & عدّ النوى في معلّق من النسيج المذاب: في مخيخ الإنسان $69\cdot10^9$ عصبونًا من أصل $86\cdot10^9$ في الدماغ كله. القياس المجسّم: $101\cdot10^9$ عصبون صغير في مخيخ رجال مسنين. & \cite{Azevedo2009, Andersen1992cereb} & مُقاس \\
5 & رفع بُعد الدخل يجعل العلاقة المطلوبة خطية & مبرهنة عدد التقسيمات: المجموع الموزون لـ$n$ مدخلًا مع عتبة ينفذ أي تصنيف لنحو $2n$ متجهًا في وضع عام. أما أن المخيخ يستعمل هذا بالذات فجزء من فرضية السطر~3. & \cite{Cover1965} & نظرية \\
6 & نحو $3\cdot10^7$ عصبون \nt{GABA} مخرج، لكل منها نحو $10^5$ مدخل & القياس المجسّم لمخيخ الإنسان: $30.5\cdot10^6$ عصبون مخرج. المجهر الإلكتروني، الجرذ: نحو $1.75\cdot10^5$ مدخل من طبقة التوسيع لكل عصبون مخرج. والناقل العصبي المثبِّط لعصبونات المخرج في المراجعة. & \cite{Andersen1992cereb, NapperHarvey1988, Ito2001} & مُقاس \\
7 & سلك خطأ واحد بالضبط لكل عصبون مخرج، نحو مرة في الثانية، وفي كل مرة ومضة قوية & مراجعة للتسجيلات: كل عصبون مخرج يتلقى مدخل \nt{GLUT} متسلقًا واحدًا، يستدعي نبضة معقدة بتردد من رتبة $1\,\text{Hz}$. & \cite{Ito2001} & مُقاس \\
8 & احتمال الومضة يتوقف على اتجاه خطأ الحركة & القرد، الجزء المحرك للعين في المخيخ، تكيّف الرمشات السريعة: اتجاه الخطأ البصري يحدد احتمال النبضة المعقدة، وحجمه توقيتها؛ والومضة تغيّر النبضات البسيطة في المحاولة التالية وتزيح الحركة على امتداد الخطأ المفضل للخلية. & \cite{Herzfeld2018} & مُقاس \\
9 & تزامن ومضة الخطأ مع مدخل نشط يُضعف المدخل ($-\eta e_i x_j$) & مراجعة للتجارب: التنبيه المشترك لمدخل من طبقة التوسيع ولسلك الخطأ يُحدث إضعافًا طويل الأمد لهذا المدخل؛ وتنبيه المدخل وحده لا يُحدثه. & \cite{Ito2001} & مُقاس \\
10 & طول الأثر مضبوط على تأخير الخطأ: عند تأخير $100\ms$ يضعف أكثر ما يضعف المدخل الذي كان نشطًا قبله بـ$100\ms$ & شرائح وفئران حية: في الجزء المسؤول عن تعلم حركة العين يكون الإضعاف مضبوطًا بدقة على فاصل نحو $120\ms$ بين المدخل والخطأ، وهو زمن وصول إشارة الخطأ في هذه المهمة؛ وفي جزء آخر تُضبط خلايا مختلفة على فواصل مختلفة. & \cite{Suvrathan2016} & مُقاس \\
11 & الدارة مرشح تكيّفي~\eqref{eq:adaptivefilter} يتعلم نموذجًا داخليًا للجسم & نموذج مستنتج من بنية الدارة. ولا توجد تجربة مباشرة قيس فيها المرشح المتعلَّم بوصفه دالة تحويل الجسم. & \cite{Fujita1982} & فرضية \\
12 & التنبؤ يطرح نتائج الحركات الذاتية، ولهذا لا ندغدغ أنفسنا & بشر: اللمسة الذاتية تبدو أقل دغدغة من اللمسة نفسها من الغير؛ وفي التصوير تستجيب القشرة الحسية الجسدية للمسة الذاتية أضعف، ويقل نشاط المخيخ حين تلمس الحركة الجلد. والتفسير بطرح التنبؤ فرضية. & \cite{Blakemore1998} & فرضية \\
13 & مع قوة خارجية يتناقص الخطأ، وبعد إزالتها تخطئ اليد في الاتجاه الآخر & بشر يمدّون أيديهم بمقبض روبوت في حقل قوة: تعود المسارات مستقيمة؛ وبعد إزالة الحقل يكون الأثر اللاحق شبه صورة مرآوية للتشوهات الأولى؛ وينتقل أيضًا إلى أجزاء من الفضاء لم يُتدرَّب فيها. & \cite{ShadmehrMussaIvaldi1994} & مُقاس \\
14 & تتعلم عمليتان~\eqref{eq:tworate}؛ وبعد الاضطراب المعاكس تعود المهارة بنفسها & بشر، حقل قوة، ثم حقل معاكس قصير ومحاولات مع تثبيت الخطأ: يعود التصحيح بنفسه إلى القديم؛ ونموذج العمليتين يعيد إنتاج ذلك، ونموذج العملية الواحدة لا. & \cite{Smith2006} & مُقاس \\
15 & أرقام الشكل~\ref{fig:tworate}: $0.84$ (البطيئة $0.63$) خلال $200$ حركة؛ $6$ معاكسة؛ السريعة $-0.53$ والبطيئة $0.46$؛ عودة إلى $0.37$ خلال $40$ حركة & حساب بـ\foreignlanguage{english}{\texttt{models/\allowbreak motor\_\allowbreak learning.py}}، $A_f=0.92$، $B_f=0.1$، $A_s=0.996$، $B_s=0.02$: $0.840$ ($0.634$ و$0.205$)، $6$ حركات، $-0.531$ و$0.457$، وقمة $0.371$ عند الحركة الـ$40$. & \foreignlanguage{english}{\texttt{models/\allowbreak motor\_\allowbreak learning.py}} & نموذج \\
16 & تلزم عمليتان لأن الجسم يتغير بسرعات مختلفة & نظرية: التقدير الأمثل مع اضطرابات بأعمار متعددة يعطي عدة مرشحات بثوابت زمنية مختلفة، ويفسر العودة التلقائية وتسارع إعادة التعلم. & \cite{Kording2007} & فرضية \\
\bottomrule
\end{longtable}}

\section{الفصل~\ref{sec:chemoutput}. المخرج الكيميائي}

السلكان إلى القلب واختلاف سرعتيهما، والتغذية الراجعة السالبة في سلاسل الهرمونات، والشفرة بتردد الدفعات، والمولّد اليومي وضبطه بالضوء، كلها مُقاسة مباشرة؛ والحد $K<8$ مبرهنة من نظرية التحكم مشتقة في الفصل؛ أما النبضات التي تولّدها التغذية الراجعة للسلسلة نفسها ففرضية تدعمها تجربة قوية.

{\footnotesize
\setlength{\tabcolsep}{3pt}\renewcommand{\arraystretch}{1.15}
\begin{longtable}{>{\raggedright}p{0.5cm}>{\raggedright}p{4.0cm}>{\raggedright}p{5.6cm}>{\raggedright}p{2.3cm}>{\raggedright\arraybackslash}p{1.7cm}}
\toprule
م & القضية & التجربة وما قيس & المصدر & الحالة \\
\midrule\endfirsthead
\toprule
م & القضية & التجربة وما قيس & المصدر & الحالة \\
\midrule\endhead
1 & ناظمة إيقاع القلب تنبض وحدها نحو $100$ مرة في الدقيقة؛ في الراحة يكون المكبح مضغوطًا، والتردد عادةً نحو $60\text{--}70$ & بشر، إيقاف السلكين إلى القلب إيقافًا تامًا: التردد الذاتي أعلى من تردد الراحة ويتناقص مع العمر، وهو عند البالغين نحو $100$ ضربة في الدقيقة. إذن يغلب المكبح في الراحة. تردد الراحة $60\text{--}70$ قيمة نموذجية لم يُستشهد لها بمصدر مستقل. & \cite{JoseCollison1970ihr} & مُقاس \\
2 & المسرِّع \nt{NORA} والمكبح \nt{ACET} مرشحا تمرير منخفض، والمكبح أسرع~\eqref{eq:gasbrake} & كلب، تحفيز معدَّل التردد للعصب المبهم أو للعصب الودي القلبي ودالة التحويل إلى تردد الأذين: كلا الطريقين مرشح تمرير منخفض، لكن تردد القطع في الطريق الودي أدنى بكثير وتضاف إليه تأخيرة صافية نحو $1.7\,\text{s}$. الخصائص تعتمد على متوسط التوتر، فالصيغة الخطية~\eqref{eq:gasbrake} تقريب. & \cite{Berger1989} & مُقاس \\
3 & المكبح سريع لأن \nt{ACET} يفتح قناة البوتاسيوم دون مرسال ثانٍ، أما \nt{NORA} فعبر سلسلة تفاعلات & رقع غشائية من خلايا الأذين: قناة البوتاسيوم التي يفتحها الأسيتيل كولين تفتحها الوحدتان الفرعيتان $\beta\gamma$ لبروتين \foreignlanguage{english}{G} المرتبط بالمستقبل، دون مرسال ثانٍ داخل الخلية. طريق \nt{NORA} عبر \foreignlanguage{english}{cAMP} لم يُستشهد له هنا. & \cite{Logothetis1987gbg} & مُقاس \\
4 & يدور الدم في الجسم في نحو دقيقة؛ ويصل \nt{ADRE} من مجرى الدم إلى كل الأعضاء التي تملك المستقبل & تقدير عام لرتبة المقدار؛ وهكذا صيغ في الفصل، ولم يُستشهد بمصدر. & --- & تقدير \\
5 & تحيط بسلسلة الهرمونات تغذية راجعة سالبة: الهرمون الأخير يثبّط المراحل الأولى & مراجعة للتجارب: هرمونات قشرة الكظر تثبّط إفراز الغدة النخامية لـ\foreignlanguage{english}{ACTH} وإفراز الوطاء للهرمون المطلِق في ثلاثة نطاقات زمنية: سريع ومتوسط وبطيء. & \cite{KellerWood1984fb} & مُقاس \\
6 & ثلاث مراحل متماثلة مستقرة عند $K<8$~\eqref{eq:cascadestable}، والدور عند الحد $2\pi\tau/\sqrt3\approx3.6\tau$ & مشتق في الفصل: عند التردد $\sqrt3/\tau$ تعطي المراحل الثلاث إزاحة طور $180^\circ$ وتوهينًا بـ$8$ أضعاف. & اشتقاق في الفصل & نظرية \\
7 & خطأ المستوى $1/(1+K)$، لذلك لا يثبّت هذا المنظِّم المستوى بدقة أفضل من نحو $10\,\%$ (القضية~\ref{prop:cascade}) & نتيجة للسطر~6 في التغذية الراجعة التناسبية. المحور الحقيقي فيه تغذية راجعة على عدة مراحل وفي عدة نطاقات زمنية (السطر~5)، فالحد يخص النموذج~\eqref{eq:cascade} ولم يُقَس. & اشتقاق في الفصل & نظرية \\
8 & عند $K=4$ و$7$ يستقر المستوى، وعند $K=12$ نبضات منتظمة بدور $3.7\text{--}3.9\,\tau$ (الشكل~\ref{fig:cascade}) & حساب بـ\texttt{models/\allowbreak chem\_\allowbreak models.py}: عند $K=12$ الأدوار المتتالية $3.9$ و$3.7$ و$3.8$ و$3.7\,\tau$؛ وعند $K=4$ السعة في نهاية الحساب $0.003$. & \texttt{models/\allowbreak chem\_\allowbreak models.py} & نموذج \\
9 & يخرج هرمون الإجهاد نبضاتٍ نحو مرة في الساعة؛ وتولّد النبضاتِ التغذيةُ الراجعة نفسها، دون مولّد مستقل & جرذان، تسريب مستمر للهرمون المطلِق: يتذبذب \foreignlanguage{english}{ACTH} والكورتيكوستيرون بتردد فيزيولوجي يقارب مرة في الساعة، فلا حاجة إلى مدخل إيقاعي من الوطاء للنبضات. ونموذج النخامية والكظر بتغذية راجعة متأخرة يعطي مثل هذه التذبذبات. & \cite{Walker2012glc, Walker2010} & فرضية \\
10 & المتلقّي يقرأ تردد الدفعات لا المستوى & قرود بلا إفراز ذاتي للهرمون المطلِق لموجّهات الغدد التناسلية: التسريب المستمر لا يعيد إفراز هرمونات النخامية، والدفعات مرة كل ساعة تعيده؛ والانتقال إلى التسريب المستمر يخمد الاستجابة من جديد. تعب المستقبل بوصفه مرشح تمرير عالٍ تفسير. & \cite{Belchetz1978} & مُقاس \\
11 & الترددات المختلفة للدفعات تؤثر تأثيرات مختلفة في خلايا مختلفة من الغدة نفسها & القرود نفسها: زيادة تواتر الدفعات إلى $2\text{--}5$ في الساعة تخفض هرموني النخامية كليهما؛ وتقليله إلى دفعة كل $3$ ساعات يخفض أحدهما (\foreignlanguage{english}{LH}) ويرفع الآخر (\foreignlanguage{english}{FSH})، فالنسبة بينهما يحددها التردد. & \cite{Wildt1981gnrh} & مُقاس \\
12 & يولّد الإيقاعَ اليومي تغذيةٌ راجعة سالبة داخل الخلية بتأخر عدة ساعات & مراجعة: بروتينات جينات الساعة تثبّط بتأخر نسخ جيناتها هي؛ والحلقة من النسخ والترجمة تعطي دورًا يقارب يومًا. & \cite{TakahashiJS2017} & مُقاس \\
13 & المولّد الرئيسي مجموعة من عشرات آلاف العصبونات تزامن بقية الجسم & مراجعة: النواة فوق التصالبة هي المولّد اليومي الرئيسي؛ عصبوناتها المفردة في المزرعة مولّدات مستقلة، والشبكة تزامنها؛ وعند الفأر نحو $10^4$ عصبون في كل جانب. & \cite{Welsh2010scn} & مُقاس \\
14 & الدور الذاتي عند الإنسان $24.18$ ساعة & بشر في ضوء خافت وبجدول منفصل عن اليوم: دور إيقاعات الميلاتونين والحرارة والكورتيزول $24.18$ ساعة في المتوسط عند الشباب والمسنين، بتشتت ضيق. & \cite{Czeisler1999} & مُقاس \\
15 & الضوء يضبط المولّد كحلقة مقفلة الطور~\eqref{eq:pll}؛ والإزاحة اللازمة نحو $11$ دقيقة في اليوم & بشر، ضوء ساطع مفرد مدته $6.7$ ساعة في أوقات مختلفة من اليوم: منحنى استجابة طورية من النوع الأول بسعة $5$ ساعات، تأخير قبل الحد الأدنى لحرارة الجسم وتقديم بعده. المعادلة~\eqref{eq:pll} نموذج معياري؛ و$11$ دقيقة $=0.18$ ساعة حساب بسيط. & \cite{Khalsa2003prc} & مُقاس \\
16 & بلا ضوء لا تزامن، وينجرف الإيقاع إلى دوره الذاتي & مكفوفون كليًا لا يدركون الضوء ويعيشون في المجتمع العادي: عند نحو نصفهم تسير إيقاعات الميلاتونين والكورتيزول بدورها الخاص الذي لا يطابق اليوم. & \cite{Sack1992blind} & مُقاس \\
\bottomrule
\end{longtable}}

\section{الفصل~\ref{sec:debugging}. تصحيح أخطاء الآلة}

ما يسمعه القطب، وانخفاض أبعاد النشاط، وعمل الواجهات القائمة على المصفوفات الصلبة، والتعلم المشترك للدماغ والمفكِّك، والنقاط البصرية الناتجة عن التحفيز، كلها مُقاسة في تجارب محكَّمة؛ وتقديرات تدفقات البيانات حساب في الفصل؛ أما الجهاز المزروع في البشر من \foreignlanguage{english}{Neuralink}، فبقدر ما أمكن التحقق، نشرت بياناته الشركة نفسها في الغالب.

{\footnotesize
\setlength{\tabcolsep}{3pt}\renewcommand{\arraystretch}{1.15}
\begin{longtable}{>{\raggedright}p{0.5cm}>{\raggedright}p{4.0cm}>{\raggedright}p{5.6cm}>{\raggedright}p{2.3cm}>{\raggedright\arraybackslash}p{1.7cm}}
\toprule
م & القضية & التجربة وما قيس & المصدر & الحالة \\
\midrule\endfirsthead
\toprule
م & القضية & التجربة وما قيس & المصدر & الحالة \\
\midrule\endhead
1 & يسمع القطب نبضات العصبونات في نصف قطر من رتبة مئة ميكرومتر، عادةً من عصبون واحد إلى بضعة عصبونات؛ وتُميَّز بشكلها & جرذ، الحقل \foreignlanguage{english}{CA1}، تسجيل داخل خلوي وخارج خلوي متزامن: شكل النبضة خارج الخلية يكرر عرض النبضة داخلها وسعتها. تقدير: كان يمكن للقطب الرباعي أن يميز $60\text{--}100$ عصبون، لكن ما يُستخرج فعلًا نحو ستة. & \cite{Henze2000extra} & مُقاس \\
2 & في الدماغ $8.6\cdot10^{10}$ عصبون؛ والمسبار يسمع نحو جزء من مئة مليون & عدّ النوى في معلّق من النسيج المذاب: $86\cdot10^9$ عصبون في دماغ الإنسان. والنسبة حساب في الفصل. & \cite{Azevedo2009} & مُقاس \\
3 & نشاط مئات العصبونات في القشرة الحركية ينحصر في عشرة أو عشرين متغيرًا مشتركًا & مراجعة لتسجيلات في القشرة الحركية عند القردة: نشاط المجموعة توصفه متغيرات كامنة قليلة مشتركة بين عصبونات كثيرة. & \cite{Gallego2017} & مُقاس \\
4 & ضجيج الجيران متشارك، ومئة عصبون تحمل تقريبًا ما تحمله ألف (القضية~\ref{prop:pooling}) & القشرة البصرية عند القرد: معامل ارتباط ضجيج العصبونات المتجاورة نحو $0.1$، والمكسب من المتوسط يتشبع عند مئة عصبون. ونظرية الارتباطات المقيِّدة للمعلومات. & \cite{Zohary1994, MorenoBote2014} & مُقاس \\
5 & التدفق الخام $2\cdot10^8$ بت/ثانية مقابل $8\cdot10^4$ بت/ثانية في النبضات~\eqref{eq:probedata} & حساب في الفصل من سعة تدفق النبضات~\eqref{eq:spikeentropy}. ومعاملات الرقمنة واقعية: رقاقة \foreignlanguage{english}{Neuralink} بتردد $19.3\,\text{kHz}$ و$10$ بتات لكل قناة. & اشتقاق في الفصل؛ \cite{Rieke1997, Musk2019} & نظرية \\
6 & المنظومة الأولى لـ\foreignlanguage{english}{Neuralink}: الرقاقات تضخّم وترقمن، والتدفق الخام يخرج عبر كابل، وبرنامج في الخارج يستخرج النبضات؛ أما الاستخراج على الرقاقة والغلاف المغلق والاتصال اللاسلكي والطاقة بلا أسلاك فللجهاز المخصص للبشر، بحسب ما أعلنته الشركة & مقالة الشركة: التدفق الخام الكامل عبر كابل \foreignlanguage{english}{USB-C}، وبرنامج يستخرج النبضات في الزمن الحقيقي خارج الرقاقة؛ والضغط على متن الجهاز والطاقة والإرسال بلا أسلاك ذُكرت خطةً للأجهزة السريرية. وعن الجهاز المزروع في البشر لا يوجد إلا إعلانات الشركة. وضرورة استخراج النبضات على الرقاقة استنتاج الفصل من~\eqref{eq:probedata}. & \cite{Musk2019}؛ تقارير \foreignlanguage{english}{Neuralink}، لا مقالة محكَّمة & مُقاس (مقالة الشركة)؛ وللبشر تقرير الشركة \\
7 & حتى $3072$ قطبًا على $96$ خيطًا، $32$ لكل خيط؛ يغرس الخيوطَ روبوتٌ يتجنب الأوعية & مقالة الشركة: $3072$ قطبًا على $96$ خيطًا، $32$ لكل خيط؛ يغرس الروبوت $6$ خيوط ($192$ قطبًا) في الدقيقة؛ وحتى $70\,\%$ من الأقطاب لدى جرذان بمصفوفة مزروعة لمدة طويلة تسمع نبضات. & \cite{Musk2019} & مُقاس (مقالة الشركة) \\
8 & عند البشر نحو ألف قناة على $64$ خيطًا؛ بعض الخيوط انزاح؛ والمؤشر نحو $10$ بت/ثانية & إعلانات الشركة فقط. لم يعثر البحث في \foreignlanguage{english}{PubMed} على مقالة محكَّمة بنتائج على البشر؛ وهذا يتفق مع التحفظ في نص الفصل. & تقارير \foreignlanguage{english}{Neuralink}؛ لا مقالة محكَّمة & مُقاس (تقرير الشركة) \\
9 & واجهة على مصفوفة من مئة قطب تتحكم في المؤشر & إنسان مشلول، $96$ قطبًا في القشرة الحركية الأولية: نية تحريك الذراع تغيّر النبضات بعد ثلاث سنوات من الإصابة؛ ويعطي المفكِّك ”مؤشرًا عصبيًا“ (البريد، التلفاز)، والتحكم في يد اصطناعية وفي ذراع روبوتية. & \cite{Hochberg2006} & مُقاس \\
10 & الكتابة ”باليد المتخيَّلة“: نحو $90$ حرفًا في الدقيقة، أقل من $8$ بت/ثانية & إنسان مشلول اليد، محاولات الكتابة باليد، ومفكِّك تكراري: $90$ حرفًا في الدقيقة بدقة $94.1\,\%$ في الزمن الحقيقي، وأكثر من $99\,\%$ بعد التصحيح التلقائي. والتقدير بالبتات حساب في الفصل. & \cite{Willett2021} & مُقاس \\
11 & محاولات الكلام تتحول إلى نص بسرعة نحو $60$ كلمة في الدقيقة & إنسان فقد الكلام الواضح، مصفوفات في القشرة: $62$ كلمة في الدقيقة؛ و$9.1\,\%$ أخطاء في الكلمات بمعجم من $50$ كلمة، و$23.8\,\%$ بمعجم من $125\,000$. & \cite{Willett2023} & مُقاس \\
12 & تستخرج الواجهة جزءًا صغيرًا من السعة: عشرات البتات في الثانية للعصبون الواحد، وآلاف للمئة (القضية~\ref{prop:probe}) & الحد الأعلى~\eqref{eq:spikeentropy} نظرية؛ والمقارنة بالسطرين~8 و10 استنتاج في الفصل. أما أن عنق الزجاجة هو انخفاض الأبعاد والجهل بالشفرة لا السلك، فلم يُختبر بتجربة مباشرة. & \cite{Rieke1997} & نظرية \\
13 & الدماغ والمفكِّك يتعلم كل منهما من الآخر & قرود تتحكم في المؤشر؛ المفكِّك يتكيف في حلقة مغلقة، والعصبونات يُعاد تعلّمها في الوقت نفسه: تزداد الدقة، وتبقى المهارة ويقل تأثرها بحركات اليد نفسها. ونصيحة الفصل بين سرعتي التعلم قاعدة هندسية، ولا تجربة مستقلة لها في الفصل. & \cite{Orsborn2014coad} & مُقاس \\
14 & التيار عبر القطب يثير العصبونات والأسلاك المحيطة دون تمييز للنوع & قشرة الفأر والقط، تسجيل الكالسيوم بالفوتونين: حتى التيار الضعيف يثير عصبونات متفرقة منتشرة على مسافة تصل إلى المليمترات، عبر إثارة المحاوير مباشرة في حجم قطره عشرات الميكرومترات. أي أن الإثارة غير انتقائية، لكنها ليست متصلة. & \cite{Histed2009stim} & مُقاس \\
15 & تحفيز القشرة البصرية يضيء نقاطًا؛ وحتى $16$ نقطة معًا تتيح التعرف على بعض الحروف وحدود الأشياء & إنسان كفيف، $96$ قطبًا في القشرة البصرية لمدة $6$ أشهر: عتبة القطب الواحد $66.8\pm36.5\,\text{\textmu A}$؛ والتحفيز المتزامن للمجموعات (حتى $16$ قطبًا، وهو حد المحفِّز) يخفض العتبة ويعطي أنماطًا قابلة للتمييز، ويُتعرَّف على بعض الحروف وحدود الأشياء. & \cite{Fernandez2021} & مُقاس \\
16 & الاتجاه الثاني لـ\foreignlanguage{english}{Neuralink} جهاز يحقن إشارات في القشرة البصرية & إعلانات الشركة عن التطوير فقط؛ لم يُعثر على بيانات محكَّمة عن عمله. & تقارير \foreignlanguage{english}{Neuralink} & فرضية \\
17 & تراكيز نواقل البث العصبية يمكن قياسها بمستشعر كيميائي في النسيج & شرائح دماغ الجرذ، قياس الفولتية الدوري السريع على ليف كربوني: قيس إطلاق السيروتونين وإعادة امتصاصه؛ والمادة تخرج خارج حدود المشبك. & \cite{Bunin1998} & مُقاس \\
\bottomrule
\end{longtable}}

\section{الفصل~\ref{sec:eetypes}. العصبونات أجهزةً}

أوضاع عمل الخلايا المفردة مُقاسة مباشرة، بالتيار عبر القطب وتسجيل الجهد؛ ورياضيات المكامل والتقسيم إلى صنفين نظرية كلاسيكية؛ أما استعمال الدماغ لإعادة تهيئة الأوضاع في الحساب فيبقى فرضية.

{\footnotesize
\setlength{\tabcolsep}{3pt}\renewcommand{\arraystretch}{1.15}
\begin{longtable}{>{\raggedright}p{0.5cm}>{\raggedright}p{4.0cm}>{\raggedright}p{5.6cm}>{\raggedright}p{2.3cm}>{\raggedright\arraybackslash}p{1.7cm}}
\toprule
م & القضية & التجربة وما قيس & المصدر & الحالة \\
\midrule\endfirsthead
\toprule
م & القضية & التجربة وما قيس & المصدر & الحالة \\
\midrule\endhead
1 & الرنّان: تيار بطيء يعاكس تغير الجهد يجعل العصبون مرشح تمرير نطاقي قمته عند بضعة هرتزات إلى عشراتها (القسم~\ref{sec:resonator}) & في الشرائح حُقن في عصبونات \nt{GLUT} تيار جيبي يزداد تردده تدريجيًا وقيست الممانعة $|Z(f)|$: القمة عند $2\text{--}5\,\text{Hz}$ في $33^\circ$C (ونحو $7\,\text{Hz}$ في $38^\circ$C)؛ ويصنع الرنين تياران بطيئان، بوتاسيومي وكاتيوني، ويقويه تيار صوديوم مستمر. وتعرض المراجعة الحيلة نفسها لأصناف كثيرة من الخلايا & \cite{HuVervaekeStorm2002,HutcheonYarom2000} & مُقاس \\
2 & التيار البطيء يتصرف كالمحاثة، ويشكّل مع سعة الغشاء دارة رنين & قيست الاستجابة دون العتبية لمحوار على التيار ووُصفت بمعادلات الموصليات بعد تقريبها خطيًا؛ والدارة المكافئة تحتوي محاثة ”ظاهراتية“ تعتمد قيمتها على الجهد & \cite{Mauro1970} & نظرية \\
3 & الثابت الزمني لمجموعة بتغذية راجعة $\tau_{\text{eff}}=\tau/(1-w)$~\eqref{eq:perfectint}؛ عند $\tau=0.1$ ثانية و$w=0.995$ يكون $20$ ثانية & مشتق في الفصل من المعادلة الخطية للمجموعة؛ واقتُرح آليةً لحفظ موضع العين في عمل نظري & اشتقاق في الفصل؛ \cite{Seung1996} & نظرية \\
4 & مكامل موضع العين يحفظ المدخل بالشبكة لا بخاصية خلية مفردة & تسجيل داخل خلوي عند سمكة في العقدة التي تحفظ موضع العين: بعد دوران سريع يتغير تردد العصبون درجةً ويبقى؛ والتيار القصير في خلية واحدة لا يسبب إلا إزاحة عابرة، أما درجات الجهد فتبقى حتى دون العتبة، إذ تحفظها المداخل المشبكية & \cite{Aksay2001} & مُقاس \\
5 & المكامل غير التسريبي يحتاج إلى ضبط دائم بالتعلم؛ وعند اختلال الضبط إما ينسى وإما ينجرف إلى الإشباع & دُرّبت سمكة بتغذية راجعة بصرية تتناسب مع $\pm$ موضع العين: صار الحفظ غير مستقر أو تسريبيًا، وهبط الثابت الزمني للعصبونات بأكثر من رتبة، إلى $<1$ ثانية؛ وأعادت التغذية الراجعة البصرية العادية الاستقرار تدريجيًا & \cite{Major2004} & مُقاس \\
6 & العصبون ثنائي الاستقرار: مدخل قصير يشغّل هضبة طويلة، والتثبيط يطفئها؛ والسيروتونين يشغّل هذه الخاصية (القسم~\ref{sec:bistable}) & تسجيل داخل خلوي لعصبونات \nt{ACET} التي تتحكم في العضلات عند القط: إثارة مشبكية قصيرة تنقل العصبون إلى عمل طويل، وتثبيط قصير يعيد ضبطه؛ وعند القط مقطوع الحبل الشوكي تظهر الحالة بعد إعطاء سلائف السيروتونين & \cite{Hounsgaard1988} & مُقاس \\
7 & صنفان: المكاملات (يرتفع التردد من الصفر) والرنّانات (عند العتبة تردد منتهٍ فورًا) & الصنفان مشتقان من نظرية التشعبات. والتجربة: في شرائح القشرة تبدأ عصبونات \nt{GLUT} العمل بتردد صغير كيفما شئنا (النوع~1)، وعصبونات \nt{GABA} السريعة تبدأ بقفزة من تردد منتهٍ (النوع~2) & \cite{Izhikevich2007,Tateno2004} & مُقاس \\
8 & العصبون الواحد مكامل أو كاشف تزامن: التثبيط يقصّر نافذة الجمع & في شرائح الحُصين النافذة التي تُجمع فيها المداخل المثيرة حتى العتبة أقصر من $2\,\text{ms}$؛ ويقصّرها التثبيط الأمامي الذي يصل فور الإثارة؛ وفي التغصنات، حيث التثبيط أضعف، النافذة أوسع & \cite{Pouille2001} & مُقاس \\
9 & خط تأخير من المحاوير (الجدول~\ref{tab:eetypes}) & قيست عند البومة الصيادة التأخيرات في المحاوير الذاهبة إلى كواشف التزامن: اختلاف طول المحوار يعطي تأخيرًا مختلفًا، والكواشف مضبوطة على فرق زمن وصول الصوت & \cite{Carr1990} & مُقاس \\
10 & ناظمة إيقاع في اليقظة وخلية صامتة في النوم & عصبونات \nt{SERO} تعمل بانتظام تقريبًا نهارًا، وتتباطأ في النوم البطيء، وتكاد تصمت في النوم السريع (التفاصيل في الفصل~\ref{sec:sleep}) & \cite{JacobsAzmitia1992,McGintyHarper1976} & مُقاس \\
11 & يحدد الجهازَ القنواتُ الأيونية وقنواتُ البث التي تضبطها (القضية~\ref{prop:eetypes}) & ثبت على شبكات صغيرة: المواد المنظِّمة تغيّر الخصائص الذاتية للعصبونات وقوة المشابك، فتُصدر الدارة نفسها إيقاعات مختلفة & \cite{Marder2012} & مُقاس \\
12 & الدماغ يستعمل إعادة تهيئة الأوضاع في الحساب (القضية~\ref{prop:eetypes}) & لا توجد تجربة مباشرة كانت فيها إعادة تهيئة وضع الخلية لازمة لحل مسألة؛ ثمة فقط تجارب تغيّر فيها المنظِّمات مخرج الدارة & \cite{Marder2012} & فرضية \\
\bottomrule
\end{longtable}}

\section{الفصل~\ref{sec:dendrites}. عدد التغصنات}

بنية الأصناف وأعداد المداخل مُقاسة تشريحيًا وبالتسجيلات الكهربائية؛ والتقدير~\eqref{eq:dendriteload} فيزياء بسيطة للمكثف، أما التفسير الحسابي لعدد المداخل فيستند إلى النظرية والنماذج.

{\footnotesize
\setlength{\tabcolsep}{3pt}\renewcommand{\arraystretch}{1.15}
\begin{longtable}{>{\raggedright}p{0.5cm}>{\raggedright}p{4.0cm}>{\raggedright}p{5.6cm}>{\raggedright}p{2.3cm}>{\raggedright\arraybackslash}p{1.7cm}}
\toprule
م & القضية & التجربة وما قيس & المصدر & الحالة \\
\midrule\endfirsthead
\toprule
م & القضية & التجربة وما قيس & المصدر & الحالة \\
\midrule\endhead
1 & السعة النوعية للغشاء $c_m\approx1\,\text{\textmu F/cm}^2$ & سُحب غشاء من جسم العصبون إلى الماصّة، وطُبقت درجة جهد، ومن انحدار التيار السعوي وُجدت السعة الكلية ثم قُسمت على المساحة: $0.9\,\text{\textmu F/cm}^2$ في ثلاثة أصناف من العصبونات & \cite{Gentet2000} & مُقاس \\
2 & زمن الشحن $t\approx c_mA\Delta V/I$~\eqref{eq:dendriteload}: العصبون الكبير يحتاج إلى عشرات المداخل المتزامنة، والصغير يكفيه مشبك كبير واحد & تقدير بقانون شحن المكثف؛ والأعداد ($20$ و$300\,\text{pF}$، و$0.1$ وبضعة $\text{nA}$) رتب مقادير & اشتقاق في الفصل & نظرية \\
3 & مشبك ضخم واحد يعطي تيارًا من بضعة نانوأمبيرات ويبلغ بالعصبون الصغير العتبة فورًا & تسجيل متزامن للنهاية وللمتلقي في شريحة: تيار المشبك الواحد $2\text{--}13\,\text{nA}$، وكل نبضة داخلة تسبب استجابة فوق العتبة، بتأخر نحو $0.4\,\text{ms}$ عند $36^\circ$C؛ والمتلقي يُخرج \nt{GLYC} & \cite{Borst1995,BorstSoria2012} & مُقاس \\
4 & صفر تغصنات: في العصبونات الحسية للمس والألم تذهب النبضة من المستشعر إلى الحبل الشوكي متجاوزة جسم الخلية & البنية (محوار ينقسم إلى اثنين عند الجسم) ثابتة تشريحيًا. وأن التوصيل لا يتطلب قابلية الجسم للإثارة بُيّن على نموذج مُتحقَّق منه لهذا العصبون: من دون جسم قابل للإثارة تعبر النبضة التفرع بالموثوقية نفسها & \cite{AmirDevor2003} & نظرية \\
5 & تغصن واحد: عصبون الشم يحمل نوعًا واحدًا من المستقبلات وينقل خط دخل واحدًا & مراجعة لتجارب على جينات المستقبلات: كل عصبون شمي يعبّر عن جين مستقبل واحد من عائلة كبيرة & \cite{Mombaerts2004} & مُقاس \\
6 & تغصنان: في كواشف اتجاه الصوت تحط المداخل من الأذنين اليسرى واليمنى على تغصنين مختلفين & تشريح وتسجيلات عند الثدييات جمعتها مراجعة: مداخل الأذنين مفرّقة على تغصنين متعاكسين & \cite{Grothe2010} & مُقاس \\
7 & تفريق المداخل على تغصنين يجعل بوابة AND أكثر صرامة & بُيّن على نموذج: الإشباع~\eqref{eq:shunt} على تغصن واحد يمنع مداخل جانب واحد من بلوغ العتبة، فيتحسن كشف التزامن. ولا توجد تجربة مباشرة غُيّر فيها توزيع المداخل & \cite{AgmonSnir1998} & نظرية \\
8 & نحو $7\cdot10^{10}$ عصبون \nt{GLUT} صغير في دارة التعلم الحركي، لكل منها $\sim4$ مداخل & عدّ الخلايا بالتجزئة المتجانسة الخواص: في هذا الجزء من دماغ الإنسان نحو $69$ مليار عصبون. وتسجيل خلية واحدة كهذه عند جرذ حي: عند اللمس تصل رشقات من تيارات منفصلة من مداخل قليلة، وتطلق الخلية حين تُجمع & \cite{Azevedo2009,Chadderton2004} & مُقاس \\
9 & العنصر العتبي ذو $n$ مدخلًا يفصل ما لا يزيد على $P_{\max}\approx2n$ نمطًا عشوائيًا~\eqref{eq:cover}؛ وللعصبون المخرج في دارة التعلم الحركي الحد $\sim2\cdot10^5$، والتقدير الواقعي $\sim4\cdot10^4$ تقابل & الحد نتيجة كلاسيكية في الهندسة التوافقية. والتقدير الواقعي نظرية السعة بأوزان موجبة فقط مع هامش للضجيج، مطبقةً على العدد المُقاس لمداخل هذا العصبون & \cite{Cover1965,Brunel2004} & نظرية \\
10 & المازجات ذات المداخل القليلة لازمة لتوسيع المدخل؛ و”الأربعة“ قريبة من الأمثل & نظرية: أبعاد التمثيل الذي تعطيه طبقة المازجات تبلغ أقصاها تقريبًا عند عدد المداخل المشاهَد تشريحيًا؛ وفرضية التوسيع الأصلية في الأعمال الكلاسيكية & \cite{Marr1969,Albus1971,LitwinKumar2017} & فرضية \\
11 & الأشجار الكبيرة: $10^4\text{--}10^5$ مدخل & عدّ المشابك في صور المجهر الإلكتروني: $\approx1.7\cdot10^5$ مشبك على عصبون \nt{GABA} مخرج واحد في دارة التعلم الحركي عند الجرذ؛ ونحو $30\,000$ مدخل مثير و$1\,700$ مثبِّط على عصبون \nt{GLUT} في الحُصين & \cite{NapperHarvey1988,Megias2001} & مُقاس \\
12 & فروع الشجرة الكبيرة دارات فرعية مستقلة لكل منها عتبتها؛ والعصبون شبكة صغيرة متعددة الطبقات & تحفيز نقطي لموضعين على تغصنات رفيعة: المداخل على فرع واحد تُجمع وفق منحنى سيني، وعلى فرعين مختلفين خطيًا. وفي تغصنات عصبونات قشرة الإنسان وُجدت نبضات كالسيوم متدرجة تتيح لخلية واحدة حل مسألة غير قابلة للفصل الخطي. ونموذج: محاكاة عصبون واحد تحتاج إلى شبكة عميقة & \cite{Polsky2004,Gidon2020,Beniaguev2021} & مُقاس \\
13 & التغصنات مخارج: كل فرع من عصبون \nt{GABA} في الشبكية كاشف اتجاه مستقل & تسجيل الكالسيوم بالفوتونين في فروع منفردة: الإشارة في الفرع انتقائية للحركة من جسم الخلية نحو الطرف، أما جهد الجسم فلا & \cite{Euler2002} & مُقاس \\
14 & عدد المداخل يتبع المهمة (القضية~\ref{prop:dendrites}) & بنية الأصناف مُقاسة (الأسطر \babelsublr{3--13})؛ أما أن عدد المداخل اختير لأجل السرعة أو التصنيف فلم يُبيَّن إلا بالنظرية والنماذج لأصناف منفردة & \cite{LitwinKumar2017,AgmonSnir1998} & فرضية \\
\bottomrule
\end{longtable}}

\section{الفصل~\ref{sec:sleep}. النوم}

أوضاع النوم والإعادات وتناقص المشابك مُقاسة بتسجيلات العصبونات والديال المجهري والمجهر الإلكتروني؛ وجدول النوم والأرجحة البطيئة توصفهما نماذج الكتاب، أما غرض النوم ففرضيات في الغالب.

{\footnotesize
\setlength{\tabcolsep}{3pt}\renewcommand{\arraystretch}{1.15}
\begin{longtable}{>{\raggedright}p{0.5cm}>{\raggedright}p{4.0cm}>{\raggedright}p{5.6cm}>{\raggedright}p{2.3cm}>{\raggedright\arraybackslash}p{1.7cm}}
\toprule
م & القضية & التجربة وما قيس & المصدر & الحالة \\
\midrule\endfirsthead
\toprule
م & القضية & التجربة وما قيس & المصدر & الحالة \\
\midrule\endhead
1 & في النوم لا تصمت عصبونات القشرة؛ الذي يتغير هو وضع العمل & تسجيلات طويلة لعصبونات القشرة عند الجرذان: في النوم البطيء يبقى التردد غير صفري، وتتناوب فترات العمل مع فترات الصمت؛ وبعد يقظة طويلة يكون التردد أعلى في كل الحالات، وبعد النوم أدنى & \cite{Vyazovskiy2009} & مُقاس \\
2 & \nt{ACET} مرتفع نهارًا وفي النوم السريع، ومنخفض في النوم البطيء (الجدول~\ref{tab:sleepmodes}) & ديال مجهري في القشرة عند قطط تتحرك بحرية: في اليقظة يكون إطلاق الأسيتيل كولين أعلى بـ$100\text{--}175\%$ منه في النوم البطيء، وفي النوم السريع يقارب مستوى اليقظة الهادئة & \cite{Marrosu1995} & مُقاس \\
3 & \nt{NORA} و\nt{SERO} يهدآن في النوم البطيء ويكادان يصمتان في النوم السريع & تسجيل عصبونات \nt{NORA} منفردة عند الجرذان وعصبونات \nt{SERO} عند القطط بحسب مراحل النوم: التردد أقصى في اليقظة، وأدنى في النوم البطيء، ويكاد يكون صفرًا في النوم السريع & \cite{AstonJonesBloom1981,McGintyHarper1976} & مُقاس \\
4 & في النوم السريع تُطفأ العضلات بالتثبيط عبر \nt{GABA} و\nt{GLYC} & عند الجرذان قيس نشاط عصبونات \nt{ACET} لعضلة المضغ وأُعطيت الحاصرات مباشرة إليها: لا يزول الشلل إلا إذا حُجبت مستقبلات \nt{GABA} بنوعيها ومستقبلات \nt{GLYC} معًا & \cite{BrooksPeever2012} & مُقاس \\
5 & ضغط النوم $S$ مكثف بعتبتين~\eqref{eq:sleeppressure}، $\tau_r=18.2$ ساعة، $\tau_d=4.2$ ساعة & نموذج العمليتين؛ الثوابت الزمنية مستخرجة من كيفية ازدياد قدرة الموجات البطيئة في تخطيط الدماغ مع اليقظة وتناقصها في النوم، وشكل العتبات من أوقات الاستيقاظ التلقائي & \cite{Borbely1982,Daan1984} & نظرية \\
6 & تصل الآلة بنفسها إلى $16.1$ ساعة يقظة و$8.0$ ساعات نوم (الشكل~\ref{fig:twoprocess}) & حساب وفق~\eqref{eq:sleeppressure} بعتبات متأرجحة، البرنامج \texttt{models/sleep\_models.py}: $16.1$ ساعة يقظة و$7.9\text{--}8.0$ ساعات نوم & --- & نموذج \\
7 & بعد $40$ ساعة بلا نوم $S=0.90$، لكن النوم لا يدوم إلا $8.8$ ساعة: الدَّين يُسدَّد بالعمق لا بالطول & النموذج: \texttt{models/sleep\_models.py} يعطي $S=0.90$ و$8.8$ ساعة. والتجربة: بعد $40.5$ ساعة بلا نوم عند ثمانية أشخاص كانت قدرة الموجات البطيئة في تخطيط الدماغ في ليلة التعافي الأولى أعلى من المعتاد بدلالة إحصائية & \cite{Borbely1981} & نموذج \\
8 & $S$ فيزيائيًا هو الأدينوزين & ديال مجهري عند القطط: يرتفع الأدينوزين خارج الخلايا في إحدى المناطق المتحكمة في اليقظة خلال اليقظة، ويواصل الارتفاع مع الحرمان من النوم، وينخفض أثناء النوم. أما أن $S$ هو الأدينوزين وحده فلم يُبيَّن & \cite{PorkkaHeiskanen1997,Dunwiddie2001} & فرضية \\
9 & في النوم البطيء تتأرجح القشرة: عمل لبضع مئات من $\text{ms}$ ثم صمت، أقل من مرة في الثانية ($<1\,\text{Hz}$، وتحت التخدير غالبًا $0.3\text{--}0.4\,\text{Hz}$) & تسجيل داخل خلوي عند قطط مخدّرة: إزالة استقطاب مدتها $0.8\text{--}1.5$ ثانية وفرط استقطاب طويل، والإيقاع $<1\,\text{Hz}$، وغالبًا $0.3\text{--}0.4\,\text{Hz}$؛ ويُرى التناوب نفسه بين العمل والصمت في النوم الطبيعي (السطر~1) & \cite{Steriade1993,Vyazovskiy2009} & مُقاس \\
10 & الأرجحة مجموعة بتغذية راجعة وتعب~\eqref{eq:updown}؛ عند $b=5$ الدور $1.1$ ثانية (قيمة النموذج) والعمل نحو $35\%$ من الوقت، وعند $b=1$ العمل منتظم (الشكل~\ref{fig:updown}) & حساب، البرنامج \texttt{models/sleep\_models.py}: الدور $1.11$ ثانية، ونسبة وقت العمل $0.35$ (المتوسط على عدد صحيح من الأدوار)؛ وعند $b=1$ نسبة العمل $1.0$. دور النموذج أقصر من المُقاس (السطر~9) & --- & نموذج \\
11 & \nt{ACET} يطفئ الأرجحة وينقل القشرة إلى العمل المنتظم & تحفيز نواة عصبونات \nt{ACET} عند القط حجب الإيقاع البطيء في $79\%$ من عصبونات القشرة ونقلها إلى العمل المنتظم، أساسًا باختفاء فرط الاستقطاب الطويل. والأسيتيل كولين يثبّط تيارات البوتاسيوم البطيئة التي تصنع التعب & \cite{SteriadeAmzica1993,McCormick1992} & مُقاس \\
12 & ومضات الإيقاع $7\text{--}15\,\text{Hz}$ تولّدها حلقة \foreignlanguage{english}{\nt{GLUT}--\nt{GABA}} بين القشرة وعقدة الترحيل & في شرائح عقدة الترحيل البصرية عند النمس تنشأ الومضات من رشقات ارتدادية لخلايا الترحيل ردًا على تثبيط من نواة \nt{GABA} مجاورة تثيرها هي نفسها & \cite{vonKrosigk1993,SteriadeMcCormickSejnowski1993} & مُقاس \\
13 & إعادات اليوم تجري مسرَّعة: في الحُصين أسرع بنحو $20$ مرة، وفي القشرة بـ$6\text{--}7$ مرات & في النوم البطيء تجري التسلسلات المتكررة لعصبونات الحُصين مضغوطة زمنيًا. بعد الجري على مسار تكررت تسلسلات عصبونات المكان بالترتيب نفسه في ومضات قصيرة من $\sim100\,\text{ms}$، مضغوطة نحو $20$ مرة؛ وفي القشرة الضغط $6\text{--}7$ مرات & \cite{Nadasdy1999,LeeWilson2002,Euston2007} & مُقاس \\
14 & تعزيز تزامن الإعادات مع القشرة يحسّن الذاكرة (علاقة سببية) & عند الجرذان بعد التعلم عُزّز ارتباط الإعادات بالموجات البطيئة وومضات الإيقاع في القشرة بتحفيز مقيَّد بالإعادات: في اليوم التالي كانت الذاكرة عالية، وعند الضوابط بمستوى الصدفة & \cite{Maingret2016} & مُقاس \\
15 & غرض الإعادات هو التعلم المخلوط للذاكرة البطيئة & نظرية نظامي التعلم وحسابات على شبكات اصطناعية: التعلم المتتابع يفسد القديم، والمخلوط لا. ولا توجد تجربة مباشرة على الدماغ تبيّن أن الإعادات لازمة لهذا بالذات & \cite{McClelland1995} & فرضية \\
16 & بعد النوم تتناقص المشابك تناسبيًا مع حجمها، والكبيرة تكاد لا تتغير & مجهر إلكتروني ثلاثي الأبعاد لـ$6920$ مشبكًا في قشرة الفأر: مساحة التماس بعد النوم أصغر بـ$\sim18\%$ مما بعد اليقظة، والتناقص متناسب مع الحجم، والمشابك الكبيرة الثابتة لا تتغير & \cite{deVivo2017} & مُقاس \\
17 & المهمة الرئيسية للنوم إعادة تطبيع الأوزان & فرضية الاتزان الداخلي المشبكي؛ والسطران 1 و16 يتفقان معها، لكنهما لا يثبتان أنها الوظيفة الرئيسية & \cite{TononiCirelli2014} & فرضية \\
18 & النوم السريع اختبار لنموذج العالم على منصة & الوضع (الجدول~\ref{tab:sleepmodes}) مُقاس؛ أما أن الشبكة تشغّل نموذج العالم فافتراض نظري بلا تجربة & \cite{Hobson2009} & فرضية \\
19 & غسل الدماغ في النوم: السؤال مفتوح & تجربة: في النوم يكون الحيز بين الخلايا أوسع بـ$60\%$ والتنظيف أسرع. وأخرى بطريقة قياس مختلفة: في النوم وتحت التخدير يكون التنظيف أبطأ بوضوح. النتائج متناقضة & \cite{Xie2013,Miao2024} & فرضية \\
20 & النوم الموضعي: مناطق من قشرة حيوان مستيقظ تسقط لحظات قصيرة في الصمت & عند الجرذان بعد يقظة طويلة تصمت عصبونات منطقة منفردة من القشرة لحظات قصيرة، كما في النوم البطيء، مع موجة بطيئة في تخطيط الدماغ الموضعي؛ وفي هذه اللحظات يخطئ الجرذ في المهمة أكثر & \cite{Vyazovskiy2011} & مُقاس \\
\bottomrule
\end{longtable}}

\section{الفصل~\ref{sec:livingmachine}. آلة من عصبونات حية}

سلوك المزارع على مصفوفات الأقطاب مُقاس في تجارب مباشرة بالتسجيل والتحفيز؛ وآلية الرشقات الشبكية تستند إلى نموذج استنزاف المشابك وإلى تجارب على مزارع دقيقة؛ أما الأقوال عن الدوبامين في الطبق فتستند إلى تجارب منفردة، يَعدّ الباحثون أنفسهم بعضها مجرد عرض توضيحي.

{\footnotesize
\setlength{\tabcolsep}{3pt}\renewcommand{\arraystretch}{1.15}
\begin{longtable}{>{\raggedright}p{0.5cm}>{\raggedright}p{4.0cm}>{\raggedright}p{5.6cm}>{\raggedright}p{2.3cm}>{\raggedright\arraybackslash}p{1.7cm}}
\toprule
م & القضية & التجربة وما قيس & المصدر & الحالة \\
\midrule\endfirsthead
\toprule
م & القضية & التجربة وما قيس & المصدر & الحالة \\
\midrule\endhead
1 & المزرعة على الشريحة: أغلبها \nt{GLUT}، ونسبة أصغر \nt{GABA} تتوقف على المستحضر؛ في مزارع الحُصين نحو $6\%$ & شبكة من آلاف العصبونات على مصفوفة أقطاب تجربة قياسية (السطر~3). نسبة عصبونات \nt{GABA} في مزارع الحُصين عُدّت بالصبغ الخاص بـ\foreignlanguage{english}{GABA}: نحو $6\%$ من العصبونات. ولا نذكر رقمًا متحققًا منه لمزارع القشرة & \cite{Benson1994} & مُقاس \\
2 & العضيّات: نسيج عصبي ثلاثي الأبعاد قطره نحو مليمتر من خلايا جذعية بشرية & استُنبتت من الخلايا الجذعية عضيّات ثلاثية الأبعاد فيها عدة مناطق دماغية، منها قشرة بطبقات من الخلايا السَّلَفية وعصبونات ناضجة & \cite{Lancaster2013} & مُقاس \\
3 & بلا دخل تشتعل المزرعة برشقات شبكية فواصلها من ثانية إلى دقائق، بحسب المزرعة & سُجّلت $58$ مزرعة قشرية كثافتها من $3\,000$ إلى $50\,000$ عصبون خمسة أسابيع ($963$ تسجيلًا مدة كل منها نصف ساعة): بعد أسبوعين تغلب في كل المزارع تقريبًا رشقات الشبكة كلها؛ والفواصل بين الرشقات الشبكية من $1$ إلى $300\,\text{s}$ وتتوقف بشدة على المستحضر & \cite{Wagenaar2006} & مُقاس \\
4 & تنتهي الرشقة الشبكية باستنزاف الناقل، والفاصل $T_{\text{burst}}\approx\tau_{\text{rec}}\ln\frac{1}{1-x^*}$~\eqref{eq:burstinterval}؛ و$2\text{--}4\,\text{s}$ تقدير النموذج مع $\tau_{\text{rec}}=0.8\,\text{s}$، والتعافي المقيس أبطأ & الصيغة مشتقة في الفصل. النموذج: شبكة بمشابك تُستنزف تولّد من تلقاء نفسها رشقات تشمل الشبكة كلها. وتجربة على مزارع دقيقة من $4\text{--}30$ عصبونًا: مدة الرشقة الشبكية تتوقف على الزمن المنقضي منذ الرشقة الشبكية السابقة، واستنزاف حويصلات الناقل يلغي الرشقات الشبكية؛ واستنتاج الباحثين أن الرشقة الشبكية يحدّها مخزون الناقل. والتعافي هناك عشرات الثواني، وهو ما يعطي بالصيغة نفسها فواصل من عشرات الثواني والدقائق & استنتاج في الفصل؛ \cite{Tsodyks2000,CohenSegal2011,Tsodyks1997} & نظرية \\
5 & ”معلّم يصمت“: يُوقف التحفيز حين تظهر الاستجابة المطلوبة، فتثبت الاستجابة & حُفّزت المزرعة بتردد $0.3\text{--}1\,\text{Hz}$ حتى تظهر الاستجابة المختارة بعد $50\pm10\ms$ من المنبه، ثم أُوقف التحفيز $5$ دقائق؛ وبعد عدة دورات صارت الاستجابة المطلوبة تُستثار فورًا؛ ورُسمت منحنيات التعلم & \cite{Shahaf2001} & مُقاس \\
6 & المزرعة تلعب التنس؛ ونمو السلاسل ذو الدلالة داخل الجلسة مع التغذية الراجعة الكاملة فقط & التفاصيل في الفصل~\ref{sec:dishbrain} وفي ملحقه: مع الصمت بدل المنبه تلعب المزرعة في آخر الجلسة أيضًا أفضل مما بلا تغذية راجعة، لكن بأثر أصغر؛ والتعلم من جلسة إلى أخرى غير ثابت & \cite{Kagan2022} & مُقاس \\
7 & أن المزرعة تتعلم التنبؤ بدخلها فرضية؛ والعبارات الرنانة عن ”الوعي“ مطعون فيها & مبدأ الطاقة الحرة اقتراح نظري؛ ولا تجربة مباشرة تميّزه من تفسيرات أبسط. وقد أشار ثلاثون باحثًا في مقالة ردّ إلى أن تغيّر السلوك داخل الحلقة لا يدل على ذكاء ولا على إحساس & \cite{Friston2010,Balci2023} & فرضية \\
8 & المزرعة تقود جسمًا افتراضيًا نحو الهدف بنمط منبهات مختار & كان برنامج يختار بنفسه أنماط التحفيز التدريبية بحسب النتيجة الجارية؛ وفي عشرات الدقائق بلغت الشبكة الحالات المطلوبة، ودامت اللدونة $80$ دقيقة بعد ساعتين من التدريب. والمنبهات نفسها، إذا قُدّمت بلا صلة بالسلوك، لم تؤثر & \cite{Bakkum2008} & مُقاس \\
9 & الخزّان: لا يُدرَّب إلا القارئ الخطي $y=W_{\text{out}}\mathbf r$~\eqref{eq:reservoir} & نظرية الحوسبة بالخزّان: أي نظام لاخطي ذي ذاكرة متلاشية مع مخرج خطي مدرَّب & \cite{Maass2002,JaegerHaas2004} & نظرية \\
10 & العضيّة خزّانًا تميّز المتكلمين بدقة نحو $78\%$ (بحسب الباحثين)؛ القارئ يتعلم بالخطأ، وروابط العضيّة تتغير بلا معلّم & قُدّمت أصوات عدة متكلمين إلى العضيّة أنماطًا من التيار على مصفوفة الأقطاب؛ وميّز القارئ المدرَّب المتكلمَ بدقة نحو $78\%$، كما يذكر الباحثون. ويذكرون أيضًا أن الاتصالية الوظيفية للعضيّة تغيّرت في أثناء التدريب، ويسمّون ذلك تعلّمًا بلا معلّم في العضيّة نفسها & \cite{Cai2023} & مُقاس \\
11 & مليون عصبون ينفق على النبضات نحو $10^{-4}\,\text{W}$~\eqref{eq:dishpower} & تقدير: كلفة النبضة $10^{-10}\,\text{J}$ من ميزانية طاقة القشرة~\eqref{eq:energy} مضروبة في عدد النبضات؛ ولا قياس مباشر لقدرة المزرعة & استنتاج في الفصل؛ \cite{Attwell2001} & نظرية \\
12 & دوبامين بومضة ضوء: $1\mM$ من الدوبامين المقيَّد، و$240$ تكرارًا، وازداد عدد النبضات المستثارة & على منصة العضيّات البعيدة حُرّر الدوبامين من قفص حساس للضوء ($1\mM$) بالأشعة فوق البنفسجية ($365\,\text{nm}$، $0.8\,\text{s}$) حين أثار المنبه نبضات أكثر؛ وفي $240$ خطوة (نحو $5$ ساعات) ازداد عدد النبضات إجمالًا. ويكتب الباحثون أن الصلة بالدوبامين لا تثبت بتجربة واحدة؛ ولا تجارب ضابطة & \cite{Jordan2024} & فرضية \\
13 & الدوبامين من خلايا حية يغيّر وزن ترانزستور عضوي تغييرًا طويلًا وقابلًا للعكس & استُنبتت خلايا تفرز الدوبامين فوق عنصر عصبي الشكل عضوي؛ تأكسد الدوبامين عند القطب يغيّر موصلية القناة؛ وقد بُيّن تغيّر الوزن طويل الأمد وإعادته & \cite{Keene2020} & مُقاس \\
14 & توجد عضيّات بعصبونات \nt{DOPA} عاملة؛ ولم يُستعمل دوبامينها معلّمًا & في عضيّات من خلايا جذعية بشرية وُجدت عصبونات \nt{DOPA} ناضجة نشطة كهربائيًا، وإنتاج للدوبامين. ولم نجد في المنشورات تجارب يعلّم فيها هذا الدوبامين شبكة أخرى & \cite{Jo2016} & مُقاس \\
15 & عضيّة توازن عمودًا: منبهات التعليم من البرنامج أفضل من العشوائية، بلا تثبيت، عبر مشابك \nt{GLUT} & عضيّات من قشرة الفأر في حلقة مغلقة مع مسألة ”العمود على العربة“: في أغلب العضيّات أعطت الإشارات التي اختارها التعلم المعزَّز نتيجة أفضل من العشوائية أو من غياب الإشارة؛ وبعد $45$ دقيقة من الراحة لم يبقَ التحسن؛ وحجب مستقبلات \foreignlanguage{english}{AMPA} و\foreignlanguage{english}{NMDA} ألغاه & \cite{Robbins2026} & مُقاس \\
16 & بلا سجلات تحكم لا تستطيع الشبكة على المنصة أن تقرر بنفسها ما يُعدّ نجاحًا (القضية~\ref{prop:livingmachine}) & في كل تجارب الأسطر \babelsublr{5--15} يحدد إشارةَ التعليم المجرِّبُ أو البرنامج؛ ولا تجربة وضعت فيها المزرعة بنفسها معيارًا للنجاح، وهذا استنتاج من الغياب لا تجربة & --- & فرضية \\
\bottomrule
\end{longtable}}

\section{الفصل~\ref{sec:dishbrain}. \foreignlanguage{english}{DishBrain}}

بنية المنصة ونتائج اللعب مأخوذة من دراسة تجريبية واحدة ومن متابعتها؛ وصيغ الضجيج والانجراف نحو الضبوط الجيدة مشتقة في الفصل ومختبرة بنموذج الكتاب؛ أما آلية التعلم في الطبق فلم تُحدَّد.

{\footnotesize
\setlength{\tabcolsep}{3pt}\renewcommand{\arraystretch}{1.15}
\begin{longtable}{>{\raggedright}p{0.5cm}>{\raggedright}p{4.0cm}>{\raggedright}p{5.6cm}>{\raggedright}p{2.3cm}>{\raggedright\arraybackslash}p{1.7cm}}
\toprule
م & القضية & التجربة وما قيس & المصدر & الحالة \\
\midrule\endfirsthead
\toprule
م & القضية & التجربة وما قيس & المصدر & الحالة \\
\midrule\endhead
1 & المنصة: نحو $800\,000$ خلية من قشرة الفأر أو نحو $10^6$ خلية بشرية على الشريحة؛ $26\,000$ قطب على $8\,\text{mm}^2$، وتسجيل حتى $1024$، وتحفيز عبر $8$ & طرائق الدراسة: شريحة \foreignlanguage{english}{MaxOne}، $26\,000$ قطب من البلاتين على $8\,\text{mm}^2$، وتسجيل $1024$ قناة بـ$20\,000$ عينة في الثانية؛ ويمكن تقنيًا تحفيز حتى $32$ قطبًا، ونجح التحفيز المستقل عبر $8$؛ واستُعملت مزارع الفأر من نحو $10$ أيام & \cite{Kagan2022} & مُقاس \\
2 & حلقة بنبضة ساعة $10\ms$: نبضات المنطقتين الحركيتين تحرّك المضرب (الشكل~\ref{fig:dishloop}) & تُعدّ النبضات في نوافذ مدة كل منها $10\ms$؛ والتأخير من النبضة إلى المنبه المجيب نحو $5\ms$، يحدده مخزن العتاد المؤقت & \cite{Kagan2022} & مُقاس \\
3 & الدخل: ترميز المكان (أيّ الأقطاب الثمانية) وترميز التردد $4\text{--}40$ نبضة/ث & في التجربة الرئيسية يزيد تردد التحفيز خطيًا من $4\,\text{Hz}$ حين تكون الكرة عند الجدار البعيد إلى $40\,\text{Hz}$ عند المضرب؛ وسعة نبضات الكرة $75\mV$؛ والنبضات مستطيلة ثنائية الطور & \cite{Kagan2022} & مُقاس \\
4 & الخرج $\Delta p=c\,(n_\uparrow-n_\downarrow)$~\eqref{eq:dishreadout}، وضجيج العدّ في النافذة $1/\sqrt{RT}$ & قاعدة الفرق بين المنطقتين موصوفة في الدراسة (بأوزان للمنطقتين بحسب الفعالية)، والصيغة الدقيقة غير معطاة. وتقدير الضجيج نتيجة للعدّ البواسوني، مشتق في الفصل & \cite{Kagan2022}؛ استنتاج في الفصل & نظرية \\
5 & المعلّم: بعد الإصابة منبه متوقَّع $75\mV$، $100$ نبضة/ث، $0.1\,\text{s}$؛ وبعد الإخفاق منبه غير متوقَّع، $150\mV$، $5$ نبضات/ث، $4\,\text{s}$ في أماكن ولحظات عشوائية، ثم $4\,\text{s}$ توقف & الطرائق: بعد الإصابة $75\mV$، $100\,\text{Hz}$، $100\ms$ على الأقطاب الثمانية كلها معًا؛ وبعد الإخفاق $150\mV$، $5\,\text{Hz}$ على أقطاب عشوائية في لحظات عشوائية مدة $4\,\text{s}$، ثم $4\,\text{s}$ بلا تحفيز. ولا دوبامين في المزرعة & \cite{Kagan2022} & مُقاس \\
6 & الشبكة تتصرف بحيث يصير دخلها متوقَّعًا & مبدأ الطاقة الحرة اقتراح نظري اشتق منه الباحثون البروتوكول؛ والتجربة تتفق معه، لكنها لا تميّزه من آليات أبسط (السطران~\babelsublr{7--8}) & \cite{Friston2010,Kagan2022} & فرضية \\
7 & إذا كان الفشل يهزّ الشبكة والنجاح لا يهزّها، فالوقت في الضبط $\rho\propto1/P_{\text{miss}}$~\eqref{eq:dishdensity} (القضية~\ref{prop:dishscramble}) & ينتج من توازن التدفقات في سلسلة ماركوف بدفعات متناظرة، مشتق في الفصل. ولا اختبار مباشر على مزرعة & استنتاج في الفصل & نظرية \\
8 & نموذج ”الخلط عند الإخفاق“ يتعلم: نسبة الإصابات $0.21\to0.38$؛ ومع دفعة بعد كل جولة $\approx0.12$، وبلا دفعات $0.19$ (الشكل~\ref{fig:dishmodel}) & حساب، النص البرمجي \foreignlanguage{english}{\texttt{models/dishbrain\_model.py}}، $40$ مزرعة نموذجية بـ$2000$ جولة لكل منها: $0.21\to0.38$؛ $0.13\to0.11$؛ $0.19\to0.19$ & --- & نموذج \\
9 & سلاسل الكرات المصدودة تطول في المزارع الحية ذات الحلقة الكاملة فقط؛ والضوابط لا تتعلم & $399$ جلسة مدة كل منها $20$ دقيقة: $9$ مزارع من الفأر ($101$ جلسة)، و$11$ بشرية ($138$)، و$6$ شرائح بوسط الاستنبات ($80$)، و$20$ مزرعة في الراحة ($42$)، و$3$ محاكيات ($38$). متوسط طول الجولة في آخر $15$ دقيقة أكبر منه في أول $5$ في المزارع الحية فقط؛ ولا أثر في الخلايا غير العصبونية (\foreignlanguage{english}{HEK293T}) ولا في الشرائح بوسط الاستنبات & \cite{Kagan2022} & مُقاس \\
10 & النمو ذو الدلالة داخل الجلسة مع التغذية الراجعة الكاملة فقط؛ ومع الصمت بدل المنبه يكون اللعب في آخر الجلسة أفضل مما بلا تغذية راجعة، لكن الأثر أصغر & $486$ جلسة على مدى $3$ أيام: ”المنبه“ ($n=27$)، و”الصمت“ ($17$)، و”بلا تغذية راجعة“ ($15$). النمو مع الزمن ذو دلالة في شرط ”المنبه“ فقط؛ وفي آخر الجلسة تفوّق شرط ”الصمت“ على ”بلا تغذية راجعة“ بأثر أصغر & \cite{Kagan2022} & مُقاس \\
11 & الأثر صغير؛ وهل تتعلم الشبكة تعلمًا دائمًا مسألة مفتوحة & داخل الجلسة التحسن متين، لكن الباحثين أنفسهم يقولون إن التعلم بين الجلسات على مدى أيام لم يُلاحَظ ثابتًا & \cite{Kagan2022} & مُقاس \\
12 & في أثناء اللعب تقترب الشبكة من النظام الحرج، وكلما اقتربت كان اللعب أفضل & تحليل الانهيارات المتسلسلة للفعالية في المزارع نفسها: الدخل المنظَّم يقرّب الشبكة من النظام الحرج، وجودة اللعب ترتبط بالقرب منه؛ لكن الحرجية وحدها، بلا معلومات عن عواقب الأفعال، لا تكفي للتعلم & \cite{Habibollahi2023} & مُقاس \\
13 & ”وعي“ المزرعة لم يُبيَّن & مقالة ردّ لثلاثين باحثًا: تغيّر السلوك في حلقة مغلقة لا يثبت ذكاءً ولا إحساسًا؛ ولا تجربة تبيّن ذلك & \cite{Balci2023} & فرضية \\
14 & التجربة الضابطة الرابعة المقترحة: منبه متوقَّع بعد الإخفاق بالمدة والطاقة نفسيهما (القسم~\ref{sec:dishspec}) & لم تُجرَ تجربة كهذه؛ هذا اقتراح الكتاب & --- & فرضية \\
\bottomrule
\end{longtable}}


\begin{thebibliography}{99}
% English entries: typeset left to right
\begin{otherlanguage}{english}
\bibitem{Azevedo2009} F. A. C. Azevedo et al., \emph{Equal numbers of neuronal and nonneuronal cells make the human brain an isometrically scaled-up primate brain}, J. Comp. Neurol. \textbf{513}, 532 (2009).
\bibitem{Schultz1997} W. Schultz, P. Dayan, and P. R. Montague, \emph{A neural substrate of prediction and reward}, Science \textbf{275}, 1593 (1997).
\bibitem{SuttonBarto2018} R. S. Sutton and A. G. Barto, \emph{Reinforcement Learning: An Introduction}, 2nd ed. (MIT Press, Cambridge, MA, 2018).
\bibitem{Doya2002} K. Doya, \emph{Metalearning and neuromodulation}, Neural Networks \textbf{15}, 495 (2002).
\bibitem{Strata1999} P. Strata and R. Harvey, \emph{Dale's principle}, Brain Res. Bull. \textbf{50}, 349 (1999).
\bibitem{vanVreeswijk1996} C. van Vreeswijk and H. Sompolinsky, \emph{Chaos in neuronal networks with balanced excitatory and inhibitory activity}, Science \textbf{274}, 1724 (1996).
\bibitem{AstonJones2005} G. Aston-Jones and J. D. Cohen, \emph{An integrative theory of locus coeruleus-norepinephrine function: adaptive gain and optimal performance}, Annu. Rev. Neurosci. \textbf{28}, 403 (2005).
\bibitem{Beaulieu2011} J.-M. Beaulieu and R. R. Gainetdinov, \emph{The physiology, signaling, and pharmacology of dopamine receptors}, Pharmacol. Rev. \textbf{63}, 182 (2011).
\bibitem{Sparso2001} J. Spars{\o} and S. Furber (eds.), \emph{Principles of Asynchronous Circuit Design: A Systems Perspective} (Kluwer, Boston, 2001).
\bibitem{Carr1990} C. E. Carr and M. Konishi, \emph{A circuit for detection of interaural time differences in the brain stem of the barn owl}, J. Neurosci. \textbf{10}, 3227 (1990).
\bibitem{Wang2001} X.-J. Wang, \emph{Synaptic reverberation underlying mnemonic persistent activity}, Trends Neurosci. \textbf{24}, 455 (2001).
\bibitem{Hahnloser2002} R. H. R. Hahnloser, A. A. Kozhevnikov, and M. S. Fee, \emph{An ultra-sparse code underlies the generation of neural sequences in a songbird}, Nature \textbf{419}, 65 (2002).
\bibitem{Barlow1965} H. B. Barlow and W. R. Levick, \emph{The mechanism of directionally selective units in rabbit's retina}, J. Physiol. \textbf{178}, 477 (1965).
\bibitem{Carandini2012} M. Carandini and D. J. Heeger, \emph{Normalization as a canonical neural computation}, Nat. Rev. Neurosci. \textbf{13}, 51 (2012).
\bibitem{Pouille2001} F. Pouille and M. Scanziani, \emph{Enforcement of temporal fidelity in pyramidal cells by somatic feed-forward inhibition}, Science \textbf{293}, 1159 (2001).
\bibitem{Tsodyks1997isn} M. V. Tsodyks, W. E. Skaggs, T. J. Sejnowski, and B. L. McNaughton, \emph{Paradoxical effects of external modulation of inhibitory interneurons}, J. Neurosci. \textbf{17}, 4382 (1997).
\bibitem{Buzsaki2012} G. Buzs\'aki and X.-J. Wang, \emph{Mechanisms of gamma oscillations}, Annu. Rev. Neurosci. \textbf{35}, 203 (2012).
\bibitem{Rieke1997} F. Rieke, D. Warland, R. de Ruyter van Steveninck, and W. Bialek, \emph{Spikes: Exploring the Neural Code} (MIT Press, Cambridge, MA, 1997).
\bibitem{Mainen1995} Z. F. Mainen and T. J. Sejnowski, \emph{Reliability of spike timing in neocortical neurons}, Science \textbf{268}, 1503 (1995).
\bibitem{Softky1993} W. R. Softky and C. Koch, \emph{The highly irregular firing of cortical cells is inconsistent with temporal integration of random EPSPs}, J. Neurosci. \textbf{13}, 334 (1993).
\bibitem{Thorpe1996} S. Thorpe, D. Fize, and C. Marlot, \emph{Speed of processing in the human visual system}, Nature \textbf{381}, 520 (1996).
\bibitem{Zohary1994} E. Zohary, M. N. Shadlen, and W. T. Newsome, \emph{Correlated neuronal discharge rate and its implications for psychophysical performance}, Nature \textbf{370}, 140 (1994).
\bibitem{MorenoBote2014} R. Moreno-Bote et al., \emph{Information-limiting correlations}, Nat. Neurosci. \textbf{17}, 1410 (2014).
\bibitem{Gollisch2008} T. Gollisch and M. Meister, \emph{Rapid neural coding in the retina with relative spike latencies}, Science \textbf{319}, 1108 (2008).
\bibitem{OKeefe1993} J. O'Keefe and M. L. Recce, \emph{Phase relationship between hippocampal place units and the EEG theta rhythm}, Hippocampus \textbf{3}, 317 (1993).
\bibitem{Destexhe2003} A. Destexhe, M. Rudolph, and D. Par\'e, \emph{The high-conductance state of neocortical neurons in vivo}, Nat. Rev. Neurosci. \textbf{4}, 739 (2003).
\bibitem{Tsodyks1997} M. V. Tsodyks and H. Markram, \emph{The neural code between neocortical pyramidal neurons depends on neurotransmitter release probability}, Proc. Natl. Acad. Sci. USA \textbf{94}, 719 (1997).
\bibitem{Lisman1997} J. E. Lisman, \emph{Bursts as a unit of neural information: making unreliable synapses reliable}, Trends Neurosci. \textbf{20}, 38 (1997).
\bibitem{Attwell2001} D. Attwell and S. B. Laughlin, \emph{An energy budget for signaling in the grey matter of the brain}, J. Cereb. Blood Flow Metab. \textbf{21}, 1133 (2001).
\bibitem{Lennie2003} P. Lennie, \emph{The cost of cortical computation}, Curr. Biol. \textbf{13}, 493 (2003).
\bibitem{Yagishita2014} S. Yagishita et al., \emph{A critical time window for dopamine actions on the structural plasticity of dendritic spines}, Science \textbf{345}, 1616 (2014).
\bibitem{Werfel2005} J. Werfel, X. Xie, and H. S. Seung, \emph{Learning curves for stochastic gradient descent in linear feedforward networks}, Neural Comput. \textbf{17}, 2699 (2005).
\bibitem{Doya2000} K. Doya, \emph{Reinforcement learning in continuous time and space}, Neural Comput. \textbf{12}, 219 (2000).
\bibitem{Oja1982} E. Oja, \emph{Simplified neuron model as a principal component analyzer}, J. Math. Biol. \textbf{15}, 267 (1982).
\bibitem{BiPoo1998} G.-Q. Bi and M.-M. Poo, \emph{Synaptic modifications in cultured hippocampal neurons: dependence on spike timing, synaptic strength, and postsynaptic cell type}, J. Neurosci. \textbf{18}, 10464 (1998).
\bibitem{Williams1992} R. J. Williams, \emph{Simple statistical gradient-following algorithms for connectionist reinforcement learning}, Mach. Learn. \textbf{8}, 229 (1992).
\bibitem{Bayer2005} H. M. Bayer and P. W. Glimcher, \emph{Midbrain dopamine neurons encode a quantitative reward prediction error signal}, Neuron \textbf{47}, 129 (2005).
\bibitem{Shen2008} W. Shen, M. Flajolet, P. Greengard, and D. J. Surmeier, \emph{Dichotomous dopaminergic control of striatal synaptic plasticity}, Science \textbf{321}, 848 (2008).
\bibitem{Dabney2020} W. Dabney, Z. Kurth-Nelson, N. Uchida, C. K. Starkweather, D. Hassabis, R. Munos, and M. Botvinick, \emph{A distributional code for value in dopamine-based reinforcement learning}, Nature \textbf{577}, 671 (2020).
\bibitem{Matsumoto2009} M. Matsumoto and O. Hikosaka, \emph{Two types of dopamine neuron distinctly convey positive and negative motivational signals}, Nature \textbf{459}, 837 (2009).
\bibitem{BrombergMartin2010} E. S. Bromberg-Martin, M. Matsumoto, and O. Hikosaka, \emph{Dopamine in motivational control: rewarding, aversive, and alerting}, Neuron \textbf{68}, 815 (2010).
\bibitem{Menegas2018} W. Menegas, K. Akiti, R. Amo, N. Uchida, and M. Watabe-Uchida, \emph{Dopamine neurons projecting to the posterior striatum reinforce avoidance of threatening stimuli}, Nat. Neurosci. \textbf{21}, 1421 (2018).
\bibitem{Niv2007} Y. Niv, N. D. Daw, D. Joel, and P. Dayan, \emph{Tonic dopamine: opportunity costs and the control of response vigor}, Psychopharmacology \textbf{191}, 507 (2007).
\bibitem{Mohebi2019} A. Mohebi et al., \emph{Dissociable dopamine dynamics for learning and motivation}, Nature \textbf{570}, 65 (2019).
\bibitem{Hamid2016} A. A. Hamid et al., \emph{Mesolimbic dopamine signals the value of work}, Nat. Neurosci. \textbf{19}, 117 (2016).
\bibitem{Kim2020} H. R. Kim, A. N. Malik et al., \emph{A unified framework for dopamine signals across timescales}, Cell \textbf{183}, 1600 (2020).
\bibitem{Jeong2022} H. Jeong et al., \emph{Mesolimbic dopamine release conveys causal associations}, Science \textbf{378}, eabq6740 (2022).
\bibitem{Amo2022} R. Amo et al., \emph{A gradual temporal shift of dopamine responses mirrors the progression of temporal difference error in machine learning}, Nat. Neurosci. \textbf{25}, 1082 (2022).
\bibitem{Takahashi2017} Y. K. Takahashi et al., \emph{Dopamine neurons respond to errors in the prediction of sensory features of expected rewards}, Neuron \textbf{95}, 1395 (2017).
\bibitem{Sharpe2017} M. J. Sharpe et al., \emph{Dopamine transients are sufficient and necessary for acquisition of model-based associations}, Nat. Neurosci. \textbf{20}, 735 (2017).
\bibitem{Gardner2018} M. P. H. Gardner, G. Schoenbaum, and S. J. Gershman, \emph{Rethinking dopamine as generalized prediction error}, Proc. R. Soc. B \textbf{285}, 20181645 (2018).
\bibitem{Engelhard2019} B. Engelhard et al., \emph{Specialized coding of sensory, motor and cognitive variables in VTA dopamine neurons}, Nature \textbf{570}, 509 (2019).
\bibitem{Clements1992} J. D. Clements, R. A. J. Lester, G. Tong, C. E. Jahr, and G. L. Westbrook, \emph{The time course of glutamate in the synaptic cleft}, Science \textbf{258}, 1498 (1992).
\bibitem{Hasselmo2006} M. E. Hasselmo, \emph{The role of acetylcholine in learning and memory}, Curr. Opin. Neurobiol. \textbf{16}, 710 (2006).
\bibitem{JacobsAzmitia1992} B. L. Jacobs and E. C. Azmitia, \emph{Structure and function of the brain serotonin system}, Physiol. Rev. \textbf{72}, 165 (1992).
\bibitem{Schiller1992} P. H. Schiller, \emph{The ON and OFF channels of the visual system}, Trends Neurosci. \textbf{15}, 86 (1992).
\bibitem{Grothe2010} B. Grothe, M. Pecka, and D. McAlpine, \emph{Mechanisms of sound localization in mammals}, Physiol. Rev. \textbf{90}, 983 (2010).
\bibitem{Markram2004} H. Markram et al., \emph{Interneurons of the neocortical inhibitory system}, Nat. Rev. Neurosci. \textbf{5}, 793 (2004).
\bibitem{Joshi2016} S. Joshi, Y. Li, R. M. Kalwani, and J. I. Gold, \emph{Relationships between pupil diameter and neuronal activity in the locus coeruleus, colliculi, and cingulate cortex}, Neuron \textbf{89}, 221 (2016).
\bibitem{ServanSchreiber1990} D. Servan-Schreiber, H. Printz, and J. D. Cohen, \emph{A network model of catecholamine effects: gain, signal-to-noise ratio, and behavior}, Science \textbf{249}, 892 (1990).
\bibitem{YuDayan2005} A. J. Yu and P. Dayan, \emph{Uncertainty, neuromodulation, and attention}, Neuron \textbf{46}, 681 (2005).
\bibitem{Nassar2012} M. R. Nassar et al., \emph{Rational regulation of learning dynamics by pupil-linked arousal systems}, Nat. Neurosci. \textbf{15}, 1040 (2012).
\bibitem{BouretSara2005} S. Bouret and S. J. Sara, \emph{Network reset: a simplified overarching theory of locus coeruleus noradrenaline function}, Trends Neurosci. \textbf{28}, 574 (2005).
\bibitem{Hebb1949} D. O. Hebb, \emph{The Organization of Behavior: A Neuropsychological Theory} (Wiley, New York, 1949).
\bibitem{MacKay1952} D. M. MacKay and W. S. McCulloch, \emph{The limiting information capacity of a neuronal link}, Bull. Math. Biophys. \textbf{14}, 127 (1952).
\bibitem{WilsonCowan1972} H. R. Wilson and J. D. Cowan, \emph{Excitatory and inhibitory interactions in localized populations of model neurons}, Biophys. J. \textbf{12}, 1 (1972).
\bibitem{AllenStevens1994} C. Allen and C. F. Stevens, \emph{An evaluation of causes for unreliability of synaptic transmission}, Proc. Natl. Acad. Sci. USA \textbf{91}, 10380 (1994).
\bibitem{Barlow1961} H. B. Barlow, \emph{Possible principles underlying the transformations of sensory messages}, in \emph{Sensory Communication}, ed. W. A. Rosenblith (MIT Press, Cambridge, MA, 1961), pp.~217--234.
\bibitem{Cardin2009} J. A. Cardin et al., \emph{Driving fast-spiking cells induces gamma rhythm and controls sensory responses}, Nature \textbf{459}, 663 (2009).
\bibitem{Lillicrap2020} T. P. Lillicrap, A. Santoro, L. Marris, C. J. Akerman, and G. Hinton, \emph{Backpropagation and the brain}, Nat. Rev. Neurosci. \textbf{21}, 335 (2020).
\bibitem{Fremaux2016} N. Fr\'emaux and W. Gerstner, \emph{Neuromodulated spike-timing-dependent plasticity, and theory of three-factor learning rules}, Front. Neural Circuits \textbf{9}, 85 (2016).
\bibitem{Haider2006} B. Haider, A. Duque, A. R. Hasenstaub, and D. A. McCormick, \emph{Neocortical network activity in vivo is generated through a dynamic balance of excitation and inhibition}, J. Neurosci. \textbf{26}, 4535 (2006).
\bibitem{HoltKoch1997} G. R. Holt and C. Koch, \emph{Shunting inhibition does not have a divisive effect on firing rates}, Neural Comput. \textbf{9}, 1001 (1997).
\bibitem{Chance2002} F. S. Chance, L. F. Abbott, and A. D. Reyes, \emph{Gain modulation from background synaptic input}, Neuron \textbf{35}, 773 (2002).
\bibitem{Ozeki2009} H. Ozeki, I. M. Finn, E. S. Schaffer, K. D. Miller, and D. Ferster, \emph{Inhibitory stabilization of the cortical network underlies visual surround suppression}, Neuron \textbf{62}, 578 (2009).
\bibitem{HasselmoBower1992} M. E. Hasselmo and J. M. Bower, \emph{Cholinergic suppression specific to intrinsic not afferent fiber synapses in rat piriform (olfactory) cortex}, J. Neurophysiol. \textbf{67}, 1222 (1992).
\bibitem{HasselmoAndersonBower1992} M. E. Hasselmo, B. P. Anderson, and J. M. Bower, \emph{Cholinergic modulation of cortical associative memory function}, J. Neurophysiol. \textbf{67}, 1230 (1992).
\bibitem{Hasselmo1999} M. E. Hasselmo, \emph{Neuromodulation: acetylcholine and memory consolidation}, Trends Cogn. Sci. \textbf{3}, 351 (1999).
\bibitem{Tsodyks1988} M. V. Tsodyks and M. V. Feigel'man, \emph{The enhanced storage capacity in neural networks with low activity level}, Europhys. Lett. \textbf{6}, 101 (1988).
\bibitem{KalmanBucy1961} R. E. Kalman and R. S. Bucy, \emph{New results in linear filtering and prediction theory}, J. Basic Eng. \textbf{83}, 95 (1961).
\bibitem{WilsonMcNaughton1994} M. A. Wilson and B. L. McNaughton, \emph{Reactivation of hippocampal ensemble memories during sleep}, Science \textbf{265}, 676 (1994).
\bibitem{BarnesSharp1999} N. M. Barnes and T. Sharp, \emph{A review of central 5-HT receptors and their function}, Neuropharmacology \textbf{38}, 1083 (1999).
\bibitem{Bunin1998} M. A. Bunin and R. M. Wightman, \emph{Quantitative evaluation of 5-hydroxytryptamine (serotonin) neuronal release and uptake: an investigation of extrasynaptic transmission}, J. Neurosci. \textbf{18}, 4854 (1998).
\bibitem{Sprouse1987} J. S. Sprouse and G. K. Aghajanian, \emph{Electrophysiological responses of serotoninergic dorsal raphe neurons to 5-HT$_{1A}$ and 5-HT$_{1B}$ agonists}, Synapse \textbf{1}, 3 (1987).
\bibitem{Sykova2008} E. Syková and C. Nicholson, \emph{Diffusion in brain extracellular space}, Physiol. Rev. \textbf{88}, 1277 (2008).
\bibitem{WilsonNicoll2001} R. I. Wilson and R. A. Nicoll, \emph{Endogenous cannabinoids mediate retrograde signalling at hippocampal synapses}, Nature \textbf{410}, 588 (2001).
\bibitem{Garthwaite2008} J. Garthwaite, \emph{Concepts of neural nitric oxide-mediated transmission}, Eur. J. Neurosci. \textbf{27}, 2783 (2008).
\bibitem{vandenPol2012} A. N. van den Pol, \emph{Neuropeptide transmission in brain circuits}, Neuron \textbf{76}, 98 (2012).
\bibitem{Dunwiddie2001} T. V. Dunwiddie and S. A. Masino, \emph{The role and regulation of adenosine in the central nervous system}, Annu. Rev. Neurosci. \textbf{24}, 31 (2001).
\bibitem{Skaggs1996} W. E. Skaggs and B. L. McNaughton, \emph{Replay of neuronal firing sequences in rat hippocampus during sleep following spatial experience}, Science \textbf{271}, 1870 (1996).
\bibitem{Baylor1979} D. A. Baylor, T. D. Lamb, and K.-W. Yau, \emph{Responses of retinal rods to single photons}, J. Physiol. \textbf{288}, 613 (1979).
\bibitem{Hudspeth1989} A. J. Hudspeth, \emph{How the ear's works work}, Nature \textbf{341}, 397 (1989).
\bibitem{NakaRushton1966} K. I. Naka and W. A. H. Rushton, \emph{S-potentials from colour units in the retina of fish (Cyprinidae)}, J. Physiol. \textbf{185}, 536 (1966).
\bibitem{Lichtsteiner2008} P. Lichtsteiner, C. Posch, and T. Delbruck, \emph{A $128\times128$ 120 dB 15 $\mu$s latency asynchronous temporal contrast vision sensor}, IEEE J. Solid-State Circuits \textbf{43}, 566 (2008).
\bibitem{Koch2006} K. Koch et al., \emph{How much the eye tells the brain}, Curr. Biol. \textbf{16}, 1428 (2006).
\bibitem{Robles2001} L. Robles and M. A. Ruggero, \emph{Mechanics of the mammalian cochlea}, Physiol. Rev. \textbf{81}, 1305 (2001).
\bibitem{Mombaerts2004} P. Mombaerts, \emph{Genes and ligands for odorant, vomeronasal and taste receptors}, Nat. Rev. Neurosci. \textbf{5}, 263 (2004).
\bibitem{WoodSlater2001} S. J. Wood and C. R. Slater, \emph{Safety factor at the neuromuscular junction}, Prog. Neurobiol. \textbf{64}, 393 (2001).
\bibitem{Henneman1957} E. Henneman, \emph{Relation between size of neurons and their susceptibility to discharge}, Science \textbf{126}, 1345 (1957).
\bibitem{HarrisWolpert1998} C. M. Harris and D. M. Wolpert, \emph{Signal-dependent noise determines motor planning}, Nature \textbf{394}, 780 (1998).
\bibitem{Hogan1984} N. Hogan, \emph{Adaptive control of mechanical impedance by coactivation of antagonist muscles}, IEEE Trans. Autom. Control \textbf{29}, 681 (1984).
\bibitem{McAuleyMarsden2000} J. H. McAuley and C. D. Marsden, \emph{Physiological and pathological tremors and rhythmic central motor control}, Brain \textbf{123}, 1545 (2000).
\bibitem{WolpertGhahramani2000} D. M. Wolpert and Z. Ghahramani, \emph{Computational principles of movement neuroscience}, Nat. Neurosci. \textbf{3}, 1212 (2000).
\bibitem{MarderBucher2001} E. Marder and D. Bucher, \emph{Central pattern generators and the control of rhythmic movements}, Curr. Biol. \textbf{11}, R986 (2001).
\bibitem{Miyazaki2014} K. W. Miyazaki et al., \emph{Optogenetic activation of dorsal raphe serotonin neurons enhances patience for future rewards}, Curr. Biol. \textbf{24}, 2033 (2014).
\bibitem{Fuglevand1993} A. J. Fuglevand, D. A. Winter, and A. E. Patla, \emph{Models of recruitment and rate coding organization in motor-unit pools}, J. Neurophysiol. \textbf{70}, 2470 (1993).
\bibitem{Curcio1990} C. A. Curcio, K. R. Sloan, R. E. Kalina, and A. E. Hendrickson, \emph{Human photoreceptor topography}, J. Comp. Neurol. \textbf{292}, 497 (1990).
\bibitem{Hayes1950} N. D. Hayes, \emph{Roots of the transcendental equation associated with a certain difference-differential equation}, J. London Math. Soc. \textbf{25}, 226 (1950).
\bibitem{MelzackWall1965} R. Melzack and P. D. Wall, \emph{Pain mechanisms: a new theory}, Science \textbf{150}, 971 (1965).
\bibitem{Fields2004} H. Fields, \emph{State-dependent opioid control of pain}, Nat. Rev. Neurosci. \textbf{5}, 565 (2004).
\bibitem{Levine1978} J. D. Levine, N. C. Gordon, and H. L. Fields, \emph{The mechanism of placebo analgesia}, Lancet \textbf{312}, 654 (1978).
\bibitem{Woolf1983} C. J. Woolf, \emph{Evidence for a central component of post-injury pain hypersensitivity}, Nature \textbf{306}, 686 (1983).
\bibitem{Seymour2004} B. Seymour et al., \emph{Temporal difference models describe higher-order learning in humans}, Nature \textbf{429}, 664 (2004).
\bibitem{EcclestonCrombez1999} C. Eccleston and G. Crombez, \emph{Pain demands attention: a cognitive-affective model of the interruptive function of pain}, Psychol. Bull. \textbf{125}, 356 (1999).
\bibitem{WidrowHoff1960} B. Widrow and M. E. Hoff, \emph{Adaptive switching circuits}, IRE WESCON Convention Record, part 4, 96 (1960).
\bibitem{Marr1969} D. Marr, \emph{A theory of cerebellar cortex}, J. Physiol. \textbf{202}, 437 (1969).
\bibitem{Albus1971} J. S. Albus, \emph{A theory of cerebellar function}, Math. Biosci. \textbf{10}, 25 (1971).
\bibitem{Cover1965} T. M. Cover, \emph{Geometrical and statistical properties of systems of linear inequalities with applications in pattern recognition}, IEEE Trans. Electron. Comput. \textbf{EC-14}, 326 (1965).
\bibitem{Ito2001} M. Ito, \emph{Cerebellar long-term depression: characterization, signal transduction, and functional roles}, Physiol. Rev. \textbf{81}, 1143 (2001).
\bibitem{Suvrathan2016} A. Suvrathan, H. L. Payne, and J. L. Raymond, \emph{Timing rules for synaptic plasticity matched to behavioral function}, Neuron \textbf{92}, 959 (2016).
\bibitem{Fujita1982} M. Fujita, \emph{Adaptive filter model of the cerebellum}, Biol. Cybern. \textbf{45}, 195 (1982).
\bibitem{Blakemore1998} S.-J. Blakemore, D. M. Wolpert, and C. D. Frith, \emph{Central cancellation of self-produced tickle sensation}, Nat. Neurosci. \textbf{1}, 635 (1998).
\bibitem{Smith2006} M. A. Smith, A. Ghazizadeh, and R. Shadmehr, \emph{Interacting adaptive processes with different timescales underlie short-term motor learning}, PLoS Biol. \textbf{4}, e179 (2006).
\bibitem{Kording2007} K. P. K\"ording, J. B. Tenenbaum, and R. Shadmehr, \emph{The dynamics of memory as a consequence of optimal adaptation to a changing body}, Nat. Neurosci. \textbf{10}, 779 (2007).
\bibitem{Yao2023} Z. Yao et al., \emph{A high-resolution transcriptomic and spatial atlas of cell types in the whole mouse brain}, Nature \textbf{624}, 317 (2023).
\bibitem{Sanzeni2020} A. Sanzeni, B. Akitake, H. C. Goldbach, C. E. Leedy, N. Brunel, and M. H. Histed, \emph{Inhibition stabilization is a widespread property of cortical networks}, eLife \textbf{9}, e54875 (2020).
\bibitem{Padamsey2022} Z. Padamsey, D. Katsanevaki, N. Dupuy, and N. L. Rochefort, \emph{Neocortex saves energy by reducing coding precision during food scarcity}, Neuron \textbf{110}, 280 (2022).
\bibitem{Stringer2019} C. Stringer, M. Pachitariu, N. Steinmetz, C. B. Reddy, M. Carandini, and K. D. Harris, \emph{Spontaneous behaviors drive multidimensional, brainwide activity}, Science \textbf{364}, eaav7893 (2019).
\bibitem{Bittner2017} K. C. Bittner, A. D. Milstein, C. Grienberger, S. Romani, and J. C. Magee, \emph{Behavioral time scale synaptic plasticity underlies CA1 place fields}, Science \textbf{357}, 1033 (2017).
\bibitem{WhittingtonBogacz2019} J. C. R. Whittington and R. Bogacz, \emph{Theories of error back-propagation in the brain}, Trends Cogn. Sci. \textbf{23}, 235 (2019).
\bibitem{Reimer2016} J. Reimer et al., \emph{Pupil fluctuations track rapid changes in adrenergic and cholinergic activity in cortex}, Nat. Commun. \textbf{7}, 13289 (2016).
\bibitem{Beniaguev2021} D. Beniaguev, I. Segev, and M. London, \emph{Single cortical neurons as deep artificial neural networks}, Neuron \textbf{109}, 2727 (2021).
\bibitem{Gidon2020} A. Gidon et al., \emph{Dendritic action potentials and computation in human layer 2/3 cortical neurons}, Science \textbf{367}, 83 (2020).
\bibitem{Gallego2022} G. Gallego et al., \emph{Event-based vision: a survey}, IEEE Trans. Pattern Anal. Mach. Intell. \textbf{44}, 154 (2022).
\bibitem{Berger1989} R. D. Berger, J. P. Saul, and R. J. Cohen, \emph{Transfer function analysis of autonomic regulation. I. Canine atrial rate response}, Am. J. Physiol. \textbf{256}, H142 (1989).
\bibitem{Walker2010} J. J. Walker, J. R. Terry, and S. L. Lightman, \emph{Origin of ultradian pulsatility in the hypothalamic--pituitary--adrenal axis}, Proc. R. Soc. B \textbf{277}, 1627 (2010).
\bibitem{Belchetz1978} P. E. Belchetz, T. M. Plant, Y. Nakai, E. J. Keogh, and E. Knobil, \emph{Hypophysial responses to continuous and intermittent delivery of hypothalamic gonadotropin-releasing hormone}, Science \textbf{202}, 631 (1978).
\bibitem{Czeisler1999} C. A. Czeisler et al., \emph{Stability, precision, and near-24-hour period of the human circadian pacemaker}, Science \textbf{284}, 2177 (1999).
\bibitem{TakahashiJS2017} J. S. Takahashi, \emph{Transcriptional architecture of the mammalian circadian clock}, Nat. Rev. Genet. \textbf{18}, 164 (2017).
\bibitem{Musk2019} E. Musk and Neuralink, \emph{An integrated brain-machine interface platform with thousands of channels}, J. Med. Internet Res. \textbf{21}, e16194 (2019).
\bibitem{Hochberg2006} L. R. Hochberg et al., \emph{Neuronal ensemble control of prosthetic devices by a human with tetraplegia}, Nature \textbf{442}, 164 (2006).
\bibitem{Willett2021} F. R. Willett et al., \emph{High-performance brain-to-text communication via handwriting}, Nature \textbf{593}, 249 (2021).
\bibitem{Willett2023} F. R. Willett et al., \emph{A high-performance speech neuroprosthesis}, Nature \textbf{620}, 1031 (2023).
\bibitem{Gallego2017} J. A. Gallego, M. G. Perich, L. E. Miller, and S. A. Solla, \emph{Neural manifolds for the control of movement}, Neuron \textbf{94}, 978 (2017).
\bibitem{Fernandez2021} E. Fern\'andez et al., \emph{Visual percepts evoked with an intracortical 96-channel microelectrode array inserted in human occipital cortex}, J. Clin. Invest. \textbf{131}, e151331 (2021).
\bibitem{Tritsch2012} N. X. Tritsch, J. B. Ding, and B. L. Sabatini, \emph{Dopaminergic neurons inhibit striatal output through non-canonical release of GABA}, Nature \textbf{490}, 262 (2012).
\bibitem{Zhang2015} S. Zhang et al., \emph{Dopaminergic and glutamatergic microdomains in a subset of rodent mesoaccumbens axons}, Nat. Neurosci. \textbf{18}, 386 (2015).
\bibitem{Mongillo2008} G. Mongillo, O. Barak, and M. Tsodyks, \emph{Synaptic theory of working memory}, Science \textbf{319}, 1543 (2008).
\bibitem{Stokes2015} M. G. Stokes, \emph{`Activity-silent' working memory in prefrontal cortex: a dynamic coding framework}, Trends Cogn. Sci. \textbf{19}, 394 (2015).
\bibitem{Smith2013} S. L. Smith, I. T. Smith, T. Branco, and M. H\"ausser, \emph{Dendritic spikes enhance stimulus selectivity in cortical neurons in vivo}, Nature \textbf{503}, 115 (2013).
\bibitem{Baden2016} T. Baden et al., \emph{The functional diversity of retinal ganglion cells in the mouse}, Nature \textbf{529}, 345 (2016).
\bibitem{Lottem2018} E. Lottem et al., \emph{Activation of serotonin neurons promotes active persistence in a probabilistic foraging task}, Nat. Commun. \textbf{9}, 1000 (2018).
\bibitem{Payeur2021} A. Payeur, J. Guerguiev, F. Zenke, B. A. Richards, and R. Naud, \emph{Burst-dependent synaptic plasticity can coordinate learning in hierarchical circuits}, Nat. Neurosci. \textbf{24}, 1010 (2021).
\bibitem{MehonicKenyon2022} A. Mehonic and A. J. Kenyon, \emph{Brain-inspired computing needs a master plan}, Nature \textbf{604}, 255 (2022).
\bibitem{Poe2020} G. R. Poe et al., \emph{Locus coeruleus: a new look at the blue spot}, Nat. Rev. Neurosci. \textbf{21}, 644 (2020).
\bibitem{Hangya2015} B. Hangya, S. P. Ranade, M. Lorenc, and A. Kepecs, \emph{Central cholinergic neurons are rapidly recruited by reinforcement feedback}, Cell \textbf{162}, 1155 (2015).
\bibitem{FernandezRuiz2019} A. Fern\'andez-Ruiz et al., \emph{Long-duration hippocampal sharp wave ripples improve memory}, Science \textbf{364}, 1082 (2019).
\bibitem{Herzfeld2018} D. J. Herzfeld, Y. Kojima, R. Soetedjo, and R. Shadmehr, \emph{Encoding of error and learning to correct that error by the Purkinje cells of the cerebellum}, Nat. Neurosci. \textbf{21}, 736 (2018).
\bibitem{Duan2014} B. Duan et al., \emph{Identification of spinal circuits transmitting and gating mechanical pain}, Cell \textbf{159}, 1417 (2014).
\bibitem{Tremblay2016} R. Tremblay, S. Lee, and B. Rudy, \emph{GABAergic interneurons in the neocortex: from cellular properties to circuits}, Neuron \textbf{91}, 260 (2016).
\bibitem{Ren2018} J. Ren et al., \emph{Anatomically defined and functionally distinct dorsal raphe serotonin sub-systems}, Cell \textbf{175}, 472 (2018).
\bibitem{Pfeffer2013} C. K. Pfeffer, M. Xue, M. He, Z. J. Huang, and M. Scanziani, \emph{Inhibition of inhibition in visual cortex: the logic of connections between molecularly distinct interneurons}, Nat. Neurosci. \textbf{16}, 1068 (2013).
\bibitem{Pi2013} H.-J. Pi et al., \emph{Cortical interneurons that specialize in disinhibitory control}, Nature \textbf{503}, 521 (2013).
\bibitem{Keller2018} G. B. Keller and T. D. Mrsic-Flogel, \emph{Predictive processing: a canonical cortical computation}, Neuron \textbf{100}, 424 (2018).
\bibitem{HubelWiesel1962} D. H. Hubel and T. N. Wiesel, \emph{Receptive fields, binocular interaction and functional architecture in the cat's visual cortex}, J. Physiol. \textbf{160}, 106 (1962).
\bibitem{Riesenhuber1999} M. Riesenhuber and T. Poggio, \emph{Hierarchical models of object recognition in cortex}, Nat. Neurosci. \textbf{2}, 1019 (1999).
\bibitem{Fukushima1980} K. Fukushima, \emph{Neocognitron: a self-organizing neural network model for a mechanism of pattern recognition unaffected by shift in position}, Biol. Cybern. \textbf{36}, 193 (1980).
\bibitem{Yamins2014} D. L. K. Yamins et al., \emph{Performance-optimized hierarchical models predict neural responses in higher visual cortex}, Proc. Natl. Acad. Sci. USA \textbf{111}, 8619 (2014).
\bibitem{Hung2005} C. P. Hung, G. Kreiman, T. Poggio, and J. J. DiCarlo, \emph{Fast readout of object identity from macaque inferior temporal cortex}, Science \textbf{310}, 863 (2005).
\bibitem{WiskottSejnowski2002} L. Wiskott and T. J. Sejnowski, \emph{Slow feature analysis: unsupervised learning of invariances}, Neural Comput. \textbf{14}, 715 (2002).
\bibitem{Foldiak1991} P. F\"oldi\'ak, \emph{Learning invariance from transformation sequences}, Neural Comput. \textbf{3}, 194 (1991).
\bibitem{LiDiCarlo2008} N. Li and J. J. DiCarlo, \emph{Unsupervised natural experience rapidly alters invariant object representation in visual cortex}, Science \textbf{321}, 1502 (2008).
\bibitem{RaoBallard1999} R. P. N. Rao and D. H. Ballard, \emph{Predictive coding in the visual cortex: a functional interpretation of some extra-classical receptive-field effects}, Nat. Neurosci. \textbf{2}, 79 (1999).
\bibitem{Kar2019} K. Kar, J. Kubilius, K. Schmidt, E. B. Issa, and J. J. DiCarlo, \emph{Evidence that recurrent circuits are critical to the ventral stream's execution of core object recognition behavior}, Nat. Neurosci. \textbf{22}, 974 (2019).
\bibitem{Szegedy2014} C. Szegedy et al., \emph{Intriguing properties of neural networks}, in \emph{Proc. ICLR} (2014), arXiv:1312.6199.
\bibitem{Weiss2002} Y. Weiss, E. P. Simoncelli, and E. H. Adelson, \emph{Motion illusions as optimal percepts}, Nat. Neurosci. \textbf{5}, 598 (2002).
\bibitem{Purves2011} D. Purves, W. T. Wojtach, and R. B. Lotto, \emph{Understanding vision in wholly empirical terms}, Proc. Natl. Acad. Sci. USA \textbf{108}, Suppl.~3, 15588 (2011).
\bibitem{LaferSousa2015} R. Lafer-Sousa, K. L. Hermann, and B. R. Conway, \emph{Striking individual differences in color perception uncovered by `the dress' photograph}, Curr. Biol. \textbf{25}, R545 (2015).
\bibitem{Wallisch2017} P. Wallisch, \emph{Illumination assumptions account for individual differences in the perceptual interpretation of a profoundly ambiguous stimulus in the color domain: ``The dress''}, J. Vis. \textbf{17}(4), 5 (2017).
\bibitem{Schwartz2009} O. Schwartz, T. J. Sejnowski, and P. Dayan, \emph{Perceptual organization in the tilt illusion}, J. Vis. \textbf{9}(4), 19 (2009).
\bibitem{SchillerCarvey2005} P. H. Schiller and C. E. Carvey, \emph{The Hermann grid illusion revisited}, Perception \textbf{34}, 1375 (2005).
\bibitem{Nijhawan1994} R. Nijhawan, \emph{Motion extrapolation in catching}, Nature \textbf{370}, 256 (1994).
\bibitem{Conway2005} B. R. Conway, A. Kitaoka, A. Yazdanbakhsh, C. C. Pack, and M. S. Livingstone, \emph{Neural basis for a powerful static motion illusion}, J. Neurosci. \textbf{25}, 5651 (2005).
\bibitem{OteroMillan2012} J. Otero-Millan, S. L. Macknik, and S. Martinez-Conde, \emph{Microsaccades and blinks trigger illusory rotation in the ``rotating snakes'' illusion}, J. Neurosci. \textbf{32}, 6043 (2012).
\bibitem{MorenoBote2007} R. Moreno-Bote, J. Rinzel, and N. Rubin, \emph{Noise-induced alternations in an attractor network model of perceptual bistability}, J. Neurophysiol. \textbf{98}, 1125 (2007).
\bibitem{Watanabe2018} E. Watanabe, A. Kitaoka, K. Sakamoto, M. Yasugi, and K. Tanaka, \emph{Illusory motion reproduced by deep neural networks trained for prediction}, Front. Psychol. \textbf{9}, 345 (2018).
\bibitem{GomezVilla2020} A. G\'omez-Villa, A. Mart\'in, J. Vazquez-Corral, M. Bertalm\'io, and J. Malo, \emph{Color illusions also deceive CNNs for low-level vision tasks: analysis and implications}, Vision Res. \textbf{176}, 156 (2020).
\bibitem{delCastillo1954} J. del Castillo and B. Katz, \emph{Quantal components of the end-plate potential}, J. Physiol. \textbf{124}, 560 (1954).
\bibitem{Nowak1984} L. Nowak, P. Bregestovski, P. Ascher, A. Herbet, and A. Prochiantz, \emph{Magnesium gates glutamate-activated channels in mouse central neurones}, Nature \textbf{307}, 462 (1984).
\bibitem{Malinow2002} R. Malinow and R. C. Malenka, \emph{AMPA receptor trafficking and synaptic plasticity}, Annu. Rev. Neurosci. \textbf{25}, 103 (2002).
\bibitem{Matsuzaki2004} M. Matsuzaki, N. Honkura, G. C. R. Ellis-Davies, and H. Kasai, \emph{Structural basis of long-term potentiation in single dendritic spines}, Nature \textbf{429}, 761 (2004).
\bibitem{Rudomin1999} P. Rudomin and R. F. Schmidt, \emph{Presynaptic inhibition in the vertebrate spinal cord revisited}, Exp. Brain Res. \textbf{129}, 1 (1999).
\bibitem{HutcheonYarom2000} B. Hutcheon and Y. Yarom, \emph{Resonance, oscillation and the intrinsic frequency preferences of neurons}, Trends Neurosci. \textbf{23}, 216 (2000).
\bibitem{Seung1996} H. S. Seung, \emph{How the brain keeps the eyes still}, Proc. Natl. Acad. Sci. USA \textbf{93}, 13339 (1996).
\bibitem{Hounsgaard1988} J. Hounsgaard, H. Hultborn, B. Jespersen, and O. Kiehn, \emph{Bistability of alpha-motoneurones in the decerebrate cat and in the acute spinal cat after intravenous 5-hydroxytryptophan}, J. Physiol. \textbf{405}, 345 (1988).
\bibitem{Izhikevich2007} E. M. Izhikevich, \emph{Dynamical Systems in Neuroscience: The Geometry of Excitability and Bursting} (MIT Press, Cambridge, MA, 2007).
\bibitem{AgmonSnir1998} H. Agmon-Snir, C. E. Carr, and J. Rinzel, \emph{The role of dendrites in auditory coincidence detection}, Nature \textbf{393}, 268 (1998).
\bibitem{BorstSoria2012} J. G. G. Borst and J. Soria van Hoeve, \emph{The calyx of Held synapse: from model synapse to auditory relay}, Annu. Rev. Physiol. \textbf{74}, 199 (2012).
\bibitem{Euler2002} T. Euler, P. B. Detwiler, and W. Denk, \emph{Directionally selective calcium signals in dendrites of starburst amacrine cells}, Nature \textbf{418}, 845 (2002).
\bibitem{AstonJonesBloom1981} G. Aston-Jones and F. E. Bloom, \emph{Activity of norepinephrine-containing locus coeruleus neurons in behaving rats anticipates fluctuations in the sleep-waking cycle}, J. Neurosci. \textbf{1}, 876 (1981).
\bibitem{BrooksPeever2012} P. L. Brooks and J. H. Peever, \emph{Identification of the transmitter and receptor mechanisms responsible for REM sleep paralysis}, J. Neurosci. \textbf{32}, 9785 (2012).
\bibitem{Borbely1982} A. A. Borb\'ely, \emph{A two process model of sleep regulation}, Hum. Neurobiol. \textbf{1}, 195 (1982).
\bibitem{PorkkaHeiskanen1997} T. Porkka-Heiskanen et al., \emph{Adenosine: a mediator of the sleep-inducing effects of prolonged wakefulness}, Science \textbf{276}, 1265 (1997).
\bibitem{Steriade1993} M. Steriade, A. N\'u\~nez, and F. Amzica, \emph{A novel slow ($<1$ Hz) oscillation of neocortical neurons in vivo: depolarizing and hyperpolarizing components}, J. Neurosci. \textbf{13}, 3252 (1993).
\bibitem{McCormick1992} D. A. McCormick, \emph{Neurotransmitter actions in the thalamus and cerebral cortex and their role in neuromodulation of thalamocortical activity}, Prog. Neurobiol. \textbf{39}, 337 (1992).
\bibitem{SteriadeMcCormickSejnowski1993} M. Steriade, D. A. McCormick, and T. J. Sejnowski, \emph{Thalamocortical oscillations in the sleeping and aroused brain}, Science \textbf{262}, 679 (1993).
\bibitem{Nadasdy1999} Z. N\'adasdy et al., \emph{Replay and time compression of recurring spike sequences in the hippocampus}, J. Neurosci. \textbf{19}, 9497 (1999).
\bibitem{Maingret2016} N. Maingret et al., \emph{Hippocampo-cortical coupling mediates memory consolidation during sleep}, Nat. Neurosci. \textbf{19}, 959 (2016).
\bibitem{McClelland1995} J. L. McClelland, B. L. McNaughton, and R. C. O'Reilly, \emph{Why there are complementary learning systems in the hippocampus and neocortex}, Psychol. Rev. \textbf{102}, 419 (1995).
\bibitem{TononiCirelli2014} G. Tononi and C. Cirelli, \emph{Sleep and the price of plasticity: from synaptic and cellular homeostasis to memory consolidation and integration}, Neuron \textbf{81}, 12 (2014).
\bibitem{deVivo2017} L. de Vivo et al., \emph{Ultrastructural evidence for synaptic scaling across the wake/sleep cycle}, Science \textbf{355}, 507 (2017).
\bibitem{Xie2013} L. Xie et al., \emph{Sleep drives metabolite clearance from the adult brain}, Science \textbf{342}, 373 (2013).
\bibitem{Miao2024} A. Miao et al., \emph{Brain clearance is reduced during sleep and anesthesia}, Nat. Neurosci. \textbf{27}, 1046 (2024).
\bibitem{Vyazovskiy2011} V. V. Vyazovskiy et al., \emph{Local sleep in awake rats}, Nature \textbf{472}, 443 (2011).
\bibitem{Lancaster2013} M. A. Lancaster et al., \emph{Cerebral organoids model human brain development and microcephaly}, Nature \textbf{501}, 373 (2013).
\bibitem{Jordan2024} F. D. Jordan et al., \emph{Open and remotely accessible Neuroplatform for research in wetware computing}, Front. Artif. Intell. \textbf{7}, 1376042 (2024).
\bibitem{Smirnova2023} L. Smirnova et al., \emph{Organoid intelligence (OI): the new frontier in biocomputing and intelligence-in-a-dish}, Front. Sci. \textbf{1}, 1017235 (2023).
\bibitem{Wagenaar2006} D. A. Wagenaar, J. Pine, and S. M. Potter, \emph{An extremely rich repertoire of bursting patterns during the development of cortical cultures}, BMC Neurosci. \textbf{7}, 11 (2006).
\bibitem{Shahaf2001} G. Shahaf and S. Marom, \emph{Learning in networks of cortical neurons}, J. Neurosci. \textbf{21}, 8782 (2001).
\bibitem{Kagan2022} B. J. Kagan et al., \emph{In vitro neurons learn and exhibit sentience when embodied in a simulated game-world}, Neuron \textbf{110}, 3952 (2022).
\bibitem{Friston2010} K. Friston, \emph{The free-energy principle: a unified brain theory?}, Nat. Rev. Neurosci. \textbf{11}, 127 (2010).
\bibitem{Balci2023} F. Balci et al., \emph{A response to claims of emergent intelligence and sentience in a dish}, Neuron \textbf{111}, 604 (2023).
\bibitem{Bakkum2008} D. J. Bakkum, Z. C. Chao, and S. M. Potter, \emph{Spatio-temporal electrical stimuli shape behavior of an embodied cortical network in a goal-directed learning task}, J. Neural Eng. \textbf{5}, 310 (2008).
\bibitem{Maass2002} W. Maass, T. Natschl\"ager, and H. Markram, \emph{Real-time computing without stable states: a new framework for neural computation based on perturbations}, Neural Comput. \textbf{14}, 2531 (2002).
\bibitem{JaegerHaas2004} H. Jaeger and H. Haas, \emph{Harnessing nonlinearity: predicting chaotic systems and saving energy in wireless communication}, Science \textbf{304}, 78 (2004).
\bibitem{Cai2023} H. Cai et al., \emph{Brain organoid reservoir computing for artificial intelligence}, Nat. Electron. \textbf{6}, 1032 (2023).
\bibitem{Habibollahi2023} F. Habibollahi, B. J. Kagan, A. N. Burkitt, and C. French, \emph{Critical dynamics arise during structured information presentation within embodied in vitro neuronal networks}, Nat. Commun. \textbf{14}, 5287 (2023).
\bibitem{Robbins2026} A. Robbins et al., \emph{Goal-directed learning in cortical organoids}, Cell Rep. \textbf{45}, 116984 (2026).
\bibitem{Keene2020} S. T. Keene et al., \emph{A biohybrid synapse with neurotransmitter-mediated plasticity}, Nat. Mater. \textbf{19}, 969 (2020).
\bibitem{Jo2016} J. Jo et al., \emph{Midbrain-like organoids from human pluripotent stem cells contain functional dopaminergic and neuromelanin-producing neurons}, Cell Stem Cell \textbf{19}, 248 (2016).
\bibitem{Hendry1987} S. H. Hendry, H. D. Schwark, E. G. Jones, and J. Yan, \emph{Numbers and proportions of GABA-immunoreactive neurons in different areas of monkey cerebral cortex}, J. Neurosci. \textbf{7}, 1503 (1987).
\bibitem{Tang2001synapses} Y. Tang, J. R. Nyengaard, D. M. G. De Groot, and H. J. G. Gundersen, \emph{Total regional and global number of synapses in the human brain neocortex}, Synapse \textbf{41}, 258 (2001).
\bibitem{Pakkenberg1997} B. Pakkenberg and H. J. G. Gundersen, \emph{Neocortical neuron number in humans: effect of sex and age}, J. Comp. Neurol. \textbf{384}, 312 (1997).
\bibitem{Matsuda2009} W. Matsuda, T. Furuta, K. C. Nakamura, H. Hioki, F. Fujiyama, R. Arai, and T. Kaneko, \emph{Single nigrostriatal dopaminergic neurons form widely spread and highly dense axonal arborizations in the neostriatum}, J. Neurosci. \textbf{29}, 444 (2009).
\bibitem{Pakkenberg1991} B. Pakkenberg, A. M{\o}ller, H. J. G. Gundersen, A. Mouritzen Dam, and H. Pakkenberg, \emph{The absolute number of nerve cells in substantia nigra in normal subjects and in patients with Parkinson's disease estimated with an unbiased stereological method}, J. Neurol. Neurosurg. Psychiatry \textbf{54}, 30 (1991).
\bibitem{Baker1990} K. G. Baker, G. M. Halliday, and I. T\"ork, \emph{Cytoarchitecture of the human dorsal raphe nucleus}, J. Comp. Neurol. \textbf{301}, 147 (1990).
\bibitem{Baker1991} K. G. Baker, G. M. Halliday, P. Halasz, J.-P. Hornung, L. B. Geffen, R. G. H. Cotton, and I. T\"ork, \emph{Cytoarchitecture of serotonin-synthesizing neurons in the pontine tegmentum of the human brain}, Synapse \textbf{7}, 301 (1991).
\bibitem{Airaksinen1991} M. S. Airaksinen, A. Paetau, L. Palj\"arvi, K. Reinikainen, P. Riekkinen, R. Suomalainen, and P. Panula, \emph{Histamine neurons in human hypothalamus: anatomy in normal and Alzheimer diseased brains}, Neuroscience \textbf{44}, 465 (1991).
\bibitem{Colquhoun1992} D. Colquhoun, P. Jonas, and B. Sakmann, \emph{Action of brief pulses of glutamate on AMPA/kainate receptors in patches from different neurones of rat hippocampal slices}, J. Physiol. \textbf{458}, 261 (1992).
\bibitem{DutarNicoll1988} P. Dutar and R. A. Nicoll, \emph{A physiological role for GABA$_B$ receptors in the central nervous system}, Nature \textbf{332}, 156 (1988).

\bibitem{Rauch2003} A. Rauch, G. La Camera, H.-R. L\"uscher, W. Senn, and S. Fusi, \emph{Neocortical pyramidal cells respond as integrate-and-fire neurons to in vivo-like input currents}, J. Neurophysiol. \textbf{90}, 1598 (2003).
\bibitem{KlumppEady1956} R. G. Klumpp and H. R. Eady, \emph{Some measurements of interaural time difference thresholds}, J. Acoust. Soc. Am. \textbf{28}, 859 (1956).
\bibitem{Brand2002} A. Brand, O. Behrend, T. Marquardt, D. McAlpine, and B. Grothe, \emph{Precise inhibition is essential for microsecond interaural time difference coding}, Nature \textbf{417}, 543 (2002).
\bibitem{Funahashi1989} S. Funahashi, C. J. Bruce, and P. S. Goldman-Rakic, \emph{Mnemonic coding of visual space in the monkey's dorsolateral prefrontal cortex}, J. Neurophysiol. \textbf{61}, 331 (1989).
\bibitem{WangMarkram2006} Y. Wang, H. Markram, P. H. Goodman, T. K. Berger, J. Ma, and P. S. Goldman-Rakic, \emph{Heterogeneity in the pyramidal network of the medial prefrontal cortex}, Nat. Neurosci. \textbf{9}, 534 (2006).

\bibitem{McCullochPitts1943} W. S. McCulloch and W. Pitts, \emph{A logical calculus of the ideas immanent in nervous activity}, Bull. Math. Biophys. \textbf{5}, 115 (1943).
\bibitem{Fried2002} S. I. Fried, T. A. M\"unch, and F. S. Werblin, \emph{Mechanisms and circuitry underlying directional selectivity in the retina}, Nature \textbf{420}, 411 (2002).

\bibitem{Hromadka2008} T. Hrom\'adka, M. R. DeWeese, and A. M. Zador, \emph{Sparse representation of sounds in the unanesthetized auditory cortex}, PLoS Biol. \textbf{6}, e16 (2008).
\bibitem{Abbott1997depression} L. F. Abbott, J. A. Varela, K. Sen, and S. B. Nelson, \emph{Synaptic depression and cortical gain control}, Science \textbf{275}, 221 (1997).
\bibitem{ShadlenNewsome1998} M. N. Shadlen and W. T. Newsome, \emph{The variable discharge of cortical neurons: implications for connectivity, computation, and information coding}, J. Neurosci. \textbf{18}, 3870 (1998).

\bibitem{TrulsonJacobs1979} M. E. Trulson and B. L. Jacobs, \emph{Raphe unit activity in freely moving cats: correlation with level of behavioral arousal}, Brain Res. \textbf{163}, 135 (1979).
\bibitem{GagnonParent2014} D. Gagnon and M. Parent, \emph{Distribution of VGLUT3 in highly collateralized axons from the rat dorsal raphe nucleus as revealed by single-neuron reconstructions}, PLoS One \textbf{9}, e87709 (2014).
\bibitem{VandermaelenAghajanian1983} C. P. Vandermaelen and G. K. Aghajanian, \emph{Electrophysiological and pharmacological characterization of serotonergic dorsal raphe neurons recorded extracellularly and intracellularly in rat brain slices}, Brain Res. \textbf{289}, 109 (1983).
\bibitem{CourtneyFord2016} N. A. Courtney and C. P. Ford, \emph{Mechanisms of 5-HT$_{1A}$ receptor-mediated transmission in dorsal raphe serotonin neurons}, J. Physiol. \textbf{594}, 953 (2016).
\bibitem{Hashemi2012serotonin} P. Hashemi et al., \emph{Brain dopamine and serotonin differ in regulation and its consequences}, Proc. Natl. Acad. Sci. USA \textbf{109}, 11510 (2012).
\bibitem{Black1934feedback} H. S. Black, \emph{Stabilized feedback amplifiers}, Bell Syst. Tech. J. \textbf{13}, 1 (1934).
\bibitem{KobayashiSchultz2008} S. Kobayashi and W. Schultz, \emph{Influence of reward delays on responses of dopamine neurons}, J. Neurosci. \textbf{28}, 7837 (2008).
\bibitem{Schweighofer2008} N. Schweighofer et al., \emph{Low-serotonin levels increase delayed reward discounting in humans}, J. Neurosci. \textbf{28}, 4528 (2008).

\bibitem{Malenka1988calcium} R. C. Malenka, J. A. Kauer, R. S. Zucker, and R. A. Nicoll, \emph{Postsynaptic calcium is sufficient for potentiation of hippocampal synaptic transmission}, Science \textbf{242}, 81 (1988).
\bibitem{Bliss1973LTP} T. V. P. Bliss and T. L{\o}mo, \emph{Long-lasting potentiation of synaptic transmission in the dentate area of the anaesthetized rabbit following stimulation of the perforant path}, J. Physiol. \textbf{232}, 331 (1973).
\bibitem{Kelso1986hebb} S. R. Kelso, A. H. Ganong, and T. H. Brown, \emph{Hebbian synapses in hippocampus}, Proc. Natl. Acad. Sci. USA \textbf{83}, 5326 (1986).
\bibitem{He2015eligibility} K. He et al., \emph{Distinct eligibility traces for LTP and LTD in cortical synapses}, Neuron \textbf{88}, 528 (2015).
\bibitem{Olveczky2005} B. P. \"Olveczky, A. S. Andalman, and M. S. Fee, \emph{Vocal experimentation in the juvenile songbird requires a basal ganglia circuit}, PLoS Biol. \textbf{3}, e153 (2005).
\bibitem{TumerBrainard2007} E. C. Tumer and M. S. Brainard, \emph{Performance variability enables adaptive plasticity of `crystallized' adult birdsong}, Nature \textbf{450}, 1240 (2007).
\bibitem{Eshel2015arithmetic} N. Eshel et al., \emph{Arithmetic and local circuitry underlying dopamine prediction errors}, Nature \textbf{525}, 243 (2015).
\bibitem{GraceOnn1989} A. A. Grace and S.-P. Onn, \emph{Morphology and electrophysiological properties of immunocytochemically identified rat dopamine neurons recorded in vitro}, J. Neurosci. \textbf{9}, 3463 (1989).
\bibitem{ODoherty2004actor} J. O'Doherty et al., \emph{Dissociable roles of ventral and dorsal striatum in instrumental conditioning}, Science \textbf{304}, 452 (2004).

\bibitem{NeweyPowell1987} W. K. Newey and J. L. Powell, \emph{Asymmetric least squares estimation and testing}, Econometrica \textbf{55}, 819 (1987).
\bibitem{Beierholm2013vigor} U. Beierholm et al., \emph{Dopamine modulates reward-related vigor}, Neuropsychopharmacology \textbf{38}, 1495 (2013).
\bibitem{Howe2013ramp} M. W. Howe, P. L. Tierney, S. G. Sandberg, P. E. M. Phillips, and A. M. Graybiel, \emph{Prolonged dopamine signalling in striatum signals proximity and value of distant rewards}, Nature \textbf{500}, 575 (2013).

\bibitem{Baker1989LC} K. G. Baker, I. T\"ork, J.-P. Hornung, and P. Halasz, \emph{The human locus coeruleus complex: an immunohistochemical and three dimensional reconstruction study}, Exp. Brain Res. \textbf{77}, 257 (1989).
\bibitem{Schwarz2015LC} L. A. Schwarz et al., \emph{Viral-genetic tracing of the input--output organization of a central noradrenaline circuit}, Nature \textbf{524}, 88 (2015).
\bibitem{Uematsu2017LC} A. Uematsu et al., \emph{Modular organization of the brainstem noradrenaline system coordinates opposing learning states}, Nat. Neurosci. \textbf{20}, 1602 (2017).
\bibitem{AstonJones1994vigilance} G. Aston-Jones, J. Rajkowski, P. Kubiak, and T. Alexinsky, \emph{Locus coeruleus neurons in monkey are selectively activated by attended cues in a vigilance task}, J. Neurosci. \textbf{14}, 4467 (1994).
\bibitem{Clayton2004LC} E. C. Clayton, J. Rajkowski, J. D. Cohen, and G. Aston-Jones, \emph{Phasic activation of monkey locus ceruleus neurons by simple decisions in a forced-choice task}, J. Neurosci. \textbf{24}, 9914 (2004).
\bibitem{Foote1975NE} S. L. Foote, R. Freedman, and A. P. Oliver, \emph{Effects of putative neurotransmitters on neuronal activity in monkey auditory cortex}, Brain Res. \textbf{86}, 229 (1975).
\bibitem{JepmaNieuwenhuis2011} M. Jepma and S. Nieuwenhuis, \emph{Pupil diameter predicts changes in the exploration--exploitation trade-off: evidence for the adaptive gain theory}, J. Cogn. Neurosci. \textbf{23}, 1587 (2011).
\bibitem{Tervo2014stochastic} D. G. R. Tervo et al., \emph{Behavioral variability through stochastic choice and its gating by anterior cingulate cortex}, Cell \textbf{159}, 21 (2014).
\bibitem{Usher1999LC} M. Usher, J. D. Cohen, D. Servan-Schreiber, J. Rajkowski, and G. Aston-Jones, \emph{The role of locus coeruleus in the regulation of cognitive performance}, Science \textbf{283}, 549 (1999).

\bibitem{Mesulam1983} M.-M. Mesulam, E. J. Mufson, B. H. Wainer, and A. I. Levey, \emph{Central cholinergic pathways in the rat: an overview based on an alternative nomenclature (Ch1--Ch6)}, Neuroscience \textbf{10}, 1185 (1983).
\bibitem{Marrosu1995} F. Marrosu et al., \emph{Microdialysis measurement of cortical and hippocampal acetylcholine release during sleep-wake cycle in freely moving cats}, Brain Res. \textbf{671}, 329 (1995).
\bibitem{Picciotto2012} M. R. Picciotto, M. J. Higley, and Y. S. Mineur, \emph{Acetylcholine as a neuromodulator: cholinergic signaling shapes nervous system function and behavior}, Neuron \textbf{76}, 116 (2012).
\bibitem{HasselmoSchnell1994} M. E. Hasselmo and E. Schnell, \emph{Laminar selectivity of the cholinergic suppression of synaptic transmission in rat hippocampal region CA1: computational modeling and brain slice physiology}, J. Neurosci. \textbf{14}, 3898 (1994).
\bibitem{Gil1997} Z. Gil, B. W. Connors, and Y. Amitai, \emph{Differential regulation of neocortical synapses by neuromodulators and activity}, Neuron \textbf{19}, 679 (1997).
\bibitem{KilgardMerzenich1998} M. P. Kilgard and M. M. Merzenich, \emph{Cortical map reorganization enabled by nucleus basalis activity}, Science \textbf{279}, 1714 (1998).
\bibitem{Atri2004} A. Atri et al., \emph{Blockade of central cholinergic receptors impairs new learning and increases proactive interference in a word paired-associate memory task}, Behav. Neurosci. \textbf{118}, 223 (2004).
\bibitem{Girardeau2009} G. Girardeau, K. Benchenane, S. I. Wiener, G. Buzs\'aki, and M. B. Zugaro, \emph{Selective suppression of hippocampal ripples impairs spatial memory}, Nat. Neurosci. \textbf{12}, 1222 (2009).

\input{supplement/bib_10}
\bibitem{CopenhagenJahr1989} D. R. Copenhagen and C. E. Jahr, \emph{Release of endogenous excitatory amino acids from turtle photoreceptors}, Nature \textbf{341}, 536 (1989).
\bibitem{GlowatzkiFuchs2002} E. Glowatzki and P. A. Fuchs, \emph{Transmitter release at the hair cell ribbon synapse}, Nat. Neurosci. \textbf{5}, 147 (2002).
\bibitem{Tinsley2016} J. N. Tinsley et al., \emph{Direct detection of a single photon by humans}, Nat. Commun. \textbf{7}, 12172 (2016).
\bibitem{Sellick1982} P. M. Sellick, R. Patuzzi, and B. M. Johnstone, \emph{Measurement of basilar membrane motion in the guinea pig using the M\"ossbauer technique}, J. Acoust. Soc. Am. \textbf{72}, 131 (1982).
\bibitem{Ruggero1997} M. A. Ruggero, N. C. Rich, A. Recio, S. S. Narayan, and L. Robles, \emph{Basilar-membrane responses to tones at the base of the chinchilla cochlea}, J. Acoust. Soc. Am. \textbf{101}, 2151 (1997).
\bibitem{NormannPerlman1979} R. A. Normann and I. Perlman, \emph{The effects of background illumination on the photoresponses of red and green cones}, J. Physiol. \textbf{286}, 491 (1979).
\bibitem{Laughlin1981} S. Laughlin, \emph{A simple coding procedure enhances a neuron's information capacity}, Z. Naturforsch. C \textbf{36}, 910 (1981).
\bibitem{CurcioAllen1990} C. A. Curcio and K. A. Allen, \emph{Topography of ganglion cells in human retina}, J. Comp. Neurol. \textbf{300}, 5 (1990).
\bibitem{Greenwood1990} D. D. Greenwood, \emph{A cochlear frequency-position function for several species --- 29 years later}, J. Acoust. Soc. Am. \textbf{87}, 2592 (1990).
\bibitem{Eguiluz2000} V. M. Egu\'iluz, M. Ospeck, Y. Choe, A. J. Hudspeth, and M. O. Magnasco, \emph{Essential nonlinearities in hearing}, Phys. Rev. Lett. \textbf{84}, 5232 (2000).
\bibitem{PalmerRussell1986} A. R. Palmer and I. J. Russell, \emph{Phase-locking in the cochlear nerve of the guinea-pig and its relation to the receptor potential of inner hair-cells}, Hear. Res. \textbf{24}, 1 (1986).
\bibitem{Malnic1999} B. Malnic, J. Hirono, T. Sato, and L. B. Buck, \emph{Combinatorial receptor codes for odors}, Cell \textbf{96}, 713 (1999).
\bibitem{Finger2005} T. E. Finger et al., \emph{ATP signaling is crucial for communication from taste buds to gustatory nerves}, Science \textbf{310}, 1495 (2005).
\bibitem{Huang2005} Y.-J. Huang et al., \emph{Mouse taste buds use serotonin as a neurotransmitter}, J. Neurosci. \textbf{25}, 843 (2005).
\bibitem{JohanssonVallbo1979} R. S. Johansson and A. B. Vallbo, \emph{Tactile sensibility in the human hand: relative and absolute densities of four types of mechanoreceptive units in glabrous skin}, J. Physiol. \textbf{286}, 283 (1979).
\bibitem{McKemy2002} D. D. McKemy, W. M. Neuhausser, and D. Julius, \emph{Identification of a cold receptor reveals a general role for TRP channels in thermosensation}, Nature \textbf{416}, 52 (2002).

\bibitem{Schmolesky1998} M. T. Schmolesky et al., \emph{Signal timing across the macaque visual system}, J. Neurophysiol. \textbf{79}, 3272 (1998).
\bibitem{Lampl2004} I. Lampl, D. Ferster, T. Poggio, and M. Riesenhuber, \emph{Intracellular measurements of spatial integration and the MAX operation in complex cells of the cat primary visual cortex}, J. Neurophysiol. \textbf{92}, 2704 (2004).
\bibitem{KobatakeTanaka1994} E. Kobatake and K. Tanaka, \emph{Neuronal selectivities to complex object features in the ventral visual pathway of the macaque cerebral cortex}, J. Neurophysiol. \textbf{71}, 856 (1994).
\bibitem{Albright1984} T. D. Albright, \emph{Direction and orientation selectivity of neurons in visual area MT of the macaque}, J. Neurophysiol. \textbf{52}, 1106 (1984).
\bibitem{Goodfellow2015} I. J. Goodfellow, J. Shlens, and C. Szegedy, \emph{Explaining and harnessing adversarial examples}, in \emph{Proc. ICLR} (2015), arXiv:1412.6572.
\bibitem{Elsayed2018} G. F. Elsayed et al., \emph{Adversarial examples that fool both computer vision and time-limited humans}, in \emph{Advances in Neural Information Processing Systems 31} (2018), arXiv:1802.08195.
\bibitem{Thompson1982} P. Thompson, \emph{Perceived rate of movement depends on contrast}, Vision Res. \textbf{22}, 377 (1982).
\bibitem{BarlowHill1963} H. B. Barlow and R. M. Hill, \emph{Evidence for a physiological explanation of the waterfall phenomenon and figural after-effects}, Nature \textbf{200}, 1345 (1963).
\bibitem{EaglemanSejnowski2000} D. M. Eagleman and T. J. Sejnowski, \emph{Motion integration and postdiction in visual awareness}, Science \textbf{287}, 2036 (2000).

\bibitem{Feinstein1955mu} B. Feinstein, B. Lindeg{\aa}rd, E. Nyman, and G. Wohlfart, \emph{Morphologic studies of motor units in normal human muscles}, Acta Anat. \textbf{23}, 127 (1955).
\bibitem{MilnerBrown1973rec} H. S. Milner-Brown, R. B. Stein, and R. Yemm, \emph{The orderly recruitment of human motor units during voluntary isometric contractions}, J. Physiol. \textbf{230}, 359 (1973).
\bibitem{Jones2002sdn} K. E. Jones, A. F. de C. Hamilton, and D. M. Wolpert, \emph{Sources of signal-dependent noise during isometric force production}, J. Neurophysiol. \textbf{88}, 1533 (2002).
\bibitem{GomiOsu1998stiff} H. Gomi and R. Osu, \emph{Task-dependent viscoelasticity of human multijoint arm and its spatial characteristics for interaction with environments}, J. Neurosci. \textbf{18}, 8965 (1998).
\bibitem{Jankowska1972Ia} E. Jankowska and W. J. Roberts, \emph{Synaptic actions of single interneurones mediating reciprocal Ia inhibition of motoneurones}, J. Physiol. \textbf{222}, 623 (1972).
\bibitem{Pruszynski2008rapid} J. A. Pruszynski, I. Kurtzer, and S. H. Scott, \emph{Rapid motor responses are appropriately tuned to the metrics of a visuospatial task}, J. Neurophysiol. \textbf{100}, 224 (2008).
\bibitem{SaundersKnill2003vis} J. A. Saunders and D. C. Knill, \emph{Humans use continuous visual feedback from the hand to control fast reaching movements}, Exp. Brain Res. \textbf{152}, 341 (2003).
\bibitem{FlanaganWing1997grip} J. R. Flanagan and A. M. Wing, \emph{The role of internal models in motion planning and control: evidence from grip force adjustments during movements of hand-held loads}, J. Neurosci. \textbf{17}, 1519 (1997).
\bibitem{Matsuoka1985osc} K. Matsuoka, \emph{Sustained oscillations generated by mutually inhibiting neurons with adaptation}, Biol. Cybern. \textbf{52}, 367 (1985).

\bibitem{Treede1995heat} R.-D. Treede, R. A. Meyer, S. N. Raja, and J. N. Campbell, \emph{Evidence for two different heat transduction mechanisms in nociceptive primary afferents innervating monkey skin}, J. Physiol. \textbf{483}, 747 (1995).
\bibitem{Weidner1999cfib} C. Weidner et al., \emph{Functional attributes discriminating mechano-insensitive and mechano-responsive C nociceptors in human skin}, J. Neurosci. \textbf{19}, 10184 (1999).
\bibitem{CampbellLaMotte1983first} J. N. Campbell and R. H. LaMotte, \emph{Latency to detection of first pain}, Brain Res. \textbf{266}, 203 (1983).
\bibitem{LaMotte1982hyper} R. H. LaMotte, J. G. Thalhammer, H. E. Torebj{\"o}rk, and C. J. Robinson, \emph{Peripheral neural mechanisms of cutaneous hyperalgesia following mild injury by heat}, J. Neurosci. \textbf{2}, 765 (1982).
\bibitem{Mancini2014touch} F. Mancini, T. Nash, G. D. Iannetti, and P. Haggard, \emph{Pain relief by touch: a quantitative approach}, Pain \textbf{155}, 635 (2014).
\bibitem{Fields1983rvm} H. L. Fields, J. Bry, I. Hentall, and G. Zorman, \emph{The activity of neurons in the rostral medulla of the rat during withdrawal from noxious heat}, J. Neurosci. \textbf{3}, 2545 (1983).
\bibitem{Crombez1996disrupt} G. Crombez, C. Eccleston, F. Baeyens, and P. Eelen, \emph{The disruptive nature of pain: an experimental investigation}, Behav. Res. Ther. \textbf{34}, 911 (1996).
\bibitem{Schouenborg1995wr} J. Schouenborg, H.-R. Weng, J. Kalliom{\"a}ki, and H. Holmberg, \emph{A survey of spinal dorsal horn neurones encoding the spatial organization of withdrawal reflexes in the rat}, Exp. Brain Res. \textbf{106}, 19 (1995).

\bibitem{Andersen1992cereb} B. B. Andersen, L. Korbo, and B. Pakkenberg, \emph{A quantitative study of the human cerebellum with unbiased stereological techniques}, J. Comp. Neurol. \textbf{326}, 549 (1992).
\bibitem{ShadmehrMussaIvaldi1994} R. Shadmehr and F. A. Mussa-Ivaldi, \emph{Adaptive representation of dynamics during learning of a motor task}, J. Neurosci. \textbf{14}, 3208 (1994).

\bibitem{JoseCollison1970ihr} A. D. Jose and D. Collison, \emph{The normal range and determinants of the intrinsic heart rate in man}, Cardiovasc. Res. \textbf{4}, 160 (1970).
\bibitem{Logothetis1987gbg} D. E. Logothetis, Y. Kurachi, J. Galper, E. J. Neer, and D. E. Clapham, \emph{The $\beta\gamma$ subunits of GTP-binding proteins activate the muscarinic K$^+$ channel in heart}, Nature \textbf{325}, 321 (1987).
\bibitem{KellerWood1984fb} M. E. Keller-Wood and M. F. Dallman, \emph{Corticosteroid inhibition of ACTH secretion}, Endocr. Rev. \textbf{5}, 1 (1984).
\bibitem{Walker2012glc} J. J. Walker et al., \emph{The origin of glucocorticoid hormone oscillations}, PLoS Biol. \textbf{10}, e1001341 (2012).
\bibitem{Wildt1981gnrh} L. Wildt et al., \emph{Frequency and amplitude of gonadotropin-releasing hormone stimulation and gonadotropin secretion in the rhesus monkey}, Endocrinology \textbf{109}, 376 (1981).
\bibitem{Welsh2010scn} D. K. Welsh, J. S. Takahashi, and S. A. Kay, \emph{Suprachiasmatic nucleus: cell autonomy and network properties}, Annu. Rev. Physiol. \textbf{72}, 551 (2010).
\bibitem{Khalsa2003prc} S. B. S. Khalsa, M. E. Jewett, C. Cajochen, and C. A. Czeisler, \emph{A phase response curve to single bright light pulses in human subjects}, J. Physiol. \textbf{549}, 945 (2003).
\bibitem{Sack1992blind} R. L. Sack, A. J. Lewy, M. L. Blood, L. D. Keith, and H. Nakagawa, \emph{Circadian rhythm abnormalities in totally blind people: incidence and clinical significance}, J. Clin. Endocrinol. Metab. \textbf{75}, 127 (1992).

\bibitem{Henze2000extra} D. A. Henze et al., \emph{Intracellular features predicted by extracellular recordings in the hippocampus in vivo}, J. Neurophysiol. \textbf{84}, 390 (2000).
\bibitem{Histed2009stim} M. H. Histed, V. Bonin, and R. C. Reid, \emph{Direct activation of sparse, distributed populations of cortical neurons by electrical microstimulation}, Neuron \textbf{63}, 508 (2009).
\bibitem{Orsborn2014coad} A. L. Orsborn et al., \emph{Closed-loop decoder adaptation shapes neural plasticity for skillful neuroprosthetic control}, Neuron \textbf{82}, 1380 (2014).

\bibitem{HuVervaekeStorm2002} H. Hu, K. Vervaeke, and J. F. Storm, \emph{Two forms of electrical resonance at theta frequencies, generated by M-current, h-current and persistent Na$^+$ current in rat hippocampal pyramidal cells}, J. Physiol. \textbf{545}, 783 (2002).
\bibitem{Mauro1970} A. Mauro, F. Conti, F. Dodge, and R. Schor, \emph{Subthreshold behavior and phenomenological impedance of the squid giant axon}, J. Gen. Physiol. \textbf{55}, 497 (1970).
\bibitem{Aksay2001} E. Aksay et al., \emph{In vivo intracellular recording and perturbation of persistent activity in a neural integrator}, Nat. Neurosci. \textbf{4}, 184 (2001).
\bibitem{Major2004} G. Major et al., \emph{Plasticity and tuning of the time course of analog persistent firing in a neural integrator}, Proc. Natl. Acad. Sci. USA \textbf{101}, 7745 (2004).
\bibitem{Tateno2004} T. Tateno, A. Harsch, and H. P. C. Robinson, \emph{Threshold firing frequency-current relationships of neurons in rat somatosensory cortex: type 1 and type 2 dynamics}, J. Neurophysiol. \textbf{92}, 2283 (2004).
\bibitem{Marder2012} E. Marder, \emph{Neuromodulation of neuronal circuits: back to the future}, Neuron \textbf{76}, 1 (2012).
\bibitem{McGintyHarper1976} D. J. McGinty and R. M. Harper, \emph{Dorsal raphe neurons: depression of firing during sleep in cats}, Brain Res. \textbf{101}, 569 (1976).

\bibitem{Gentet2000} L. J. Gentet, G. J. Stuart, and J. D. Clements, \emph{Direct measurement of specific membrane capacitance in neurons}, Biophys. J. \textbf{79}, 314 (2000).
\bibitem{Borst1995} J. G. G. Borst, F. Helmchen, and B. Sakmann, \emph{Pre- and postsynaptic whole-cell recordings in the medial nucleus of the trapezoid body of the rat}, J. Physiol. \textbf{489}, 825 (1995).
\bibitem{AmirDevor2003} R. Amir and M. Devor, \emph{Electrical excitability of the soma of sensory neurons is required for spike invasion of the soma, but not for through-conduction}, Biophys. J. \textbf{84}, 2181 (2003).
\bibitem{Chadderton2004} P. Chadderton, T. W. Margrie, and M. H\"ausser, \emph{Integration of quanta in cerebellar granule cells during sensory processing}, Nature \textbf{428}, 856 (2004).
\bibitem{LitwinKumar2017} A. Litwin-Kumar et al., \emph{Optimal degrees of synaptic connectivity}, Neuron \textbf{93}, 1153 (2017).
\bibitem{NapperHarvey1988} R. M. A. Napper and R. J. Harvey, \emph{Number of parallel fiber synapses on an individual Purkinje cell in the cerebellum of the rat}, J. Comp. Neurol. \textbf{274}, 168 (1988).
\bibitem{Megias2001} M. Meg\'{\i}as, Z. Emri, T. F. Freund, and A. I. Guly\'as, \emph{Total number and distribution of inhibitory and excitatory synapses on hippocampal CA1 pyramidal cells}, Neuroscience \textbf{102}, 527 (2001).
\bibitem{Polsky2004} A. Polsky, B. W. Mel, and J. Schiller, \emph{Computational subunits in thin dendrites of pyramidal cells}, Nat. Neurosci. \textbf{7}, 621 (2004).
\bibitem{Brunel2004} N. Brunel et al., \emph{Optimal information storage and the distribution of synaptic weights: perceptron versus Purkinje cell}, Neuron \textbf{43}, 745 (2004).

\bibitem{Vyazovskiy2009} V. V. Vyazovskiy et al., \emph{Cortical firing and sleep homeostasis}, Neuron \textbf{63}, 865 (2009).
\bibitem{Daan1984} S. Daan, D. G. Beersma, and A. A. Borb\'ely, \emph{Timing of human sleep: recovery process gated by a circadian pacemaker}, Am. J. Physiol. \textbf{246}, R161 (1984).
\bibitem{Borbely1981} A. A. Borb\'ely et al., \emph{Sleep deprivation: effect on sleep stages and EEG power density in man}, Electroencephalogr. Clin. Neurophysiol. \textbf{51}, 483 (1981).
\bibitem{SteriadeAmzica1993} M. Steriade, F. Amzica, and A. Nu\~nez, \emph{Cholinergic and noradrenergic modulation of the slow ($\approx0.3$ Hz) oscillation in neocortical cells}, J. Neurophysiol. \textbf{70}, 1385 (1993).
\bibitem{vonKrosigk1993} M. von Krosigk, T. Bal, and D. A. McCormick, \emph{Cellular mechanisms of a synchronized oscillation in the thalamus}, Science \textbf{261}, 361 (1993).
\bibitem{LeeWilson2002} A. K. Lee and M. A. Wilson, \emph{Memory of sequential experience in the hippocampus during slow wave sleep}, Neuron \textbf{36}, 1183 (2002).
\bibitem{Euston2007} D. R. Euston, M. Tatsuno, and B. L. McNaughton, \emph{Fast-forward playback of recent memory sequences in prefrontal cortex during sleep}, Science \textbf{318}, 1147 (2007).
\bibitem{Hobson2009} J. A. Hobson, \emph{REM sleep and dreaming: towards a theory of protoconsciousness}, Nat. Rev. Neurosci. \textbf{10}, 803 (2009).

\bibitem{CohenSegal2011} D. Cohen and M. Segal, \emph{Network bursts in hippocampal microcultures are terminated by exhaustion of vesicle pools}, J. Neurophysiol. \textbf{106}, 2314 (2011).
\bibitem{Tsodyks2000} M. Tsodyks, A. Uziel, and H. Markram, \emph{Synchrony generation in recurrent networks with frequency-dependent synapses}, J. Neurosci. \textbf{20}, RC50 (2000).
\bibitem{Benson1994} D. L. Benson, F. H. Watkins, O. Steward, and G. Banker, \emph{Characterization of GABAergic neurons in hippocampal cell cultures}, J. Neurocytol. \textbf{23}, 279 (1994).

\input{supplement/bib_22}
\end{otherlanguage}
\end{thebibliography}


\end{document}
